% SPDX-FileCopyrightText: 2024-2026 Fabian Januszewski
% SPDX-License-Identifier: CC-BY-4.0

\documentclass[11pt]{article}
% VERSION 1.25.3 (slot 26, 2026-11-10: SIGNED BY THE OWNER 2026-11-10 ("I sign those title pages."); the text of 1.25.3-pre.1, unchanged: header, date and release record only -- no statement, label, number or tag moves; the title page names no sign-off)
% VERSION 1.25.3-pre.1 (slot 26, 2026-11-10, pre-release for owner signature, NOT signed): title page only, no change of mathematical content (public release HACS 0.1.1): the \date line names only the version, the date and the revision history; the words 'signed by the owner 2026-11-09' leave the title page (the sign-off record stays in the release record of Appendix app:history). No statement, label, number or tag moves.
% VERSION 1.25.2 (slot 26, 2026-11-09: SIGNED BY THE OWNER 2026-11-09 ("I sign."); the text of 1.25.2-pre.1, unchanged: header, date and release record only -- no statement, label, number or tag moves; referee LAYOUT_signed_referee.md ACCEPT)
% VERSION 1.25.2-pre.1 (slot 26, 2026-11-09, pre-release for owner re-signature, NOT signed): paths only, no change of content (repository layout, owner ruling 2026-11-09): one line names src/ModArithmetic/ENGINE_API.md (old -> new: tools/layout/mapping.tsv); no statement, label, number or tag moves.
% VERSION 1.25.1 (slot 26, 2026-11-09: SIGNED BY THE OWNER 2026-11-09; the text of 1.25.1-pre.1, unchanged: header, date and release record only -- no statement, label, number or tag moves)
% VERSION 1.25.1-pre.1 (slot 26, 2026-11-09, pre-release, NOT signed): no change of mathematical content: machine names and private paths removed for the public release (owner ruling 2026-11-09).
\usepackage{amsmath, amssymb, amsthm}
\usepackage{geometry}
\geometry{a4paper, margin=0.8in}
\usepackage{hyperref}
\usepackage{booktabs}
\usepackage{array}
\usepackage{enumitem}
\setlist{nosep}
\usepackage{longtable}
\emergencystretch=3em
\setcounter{tocdepth}{1}

% Theorem environments (same style as hecke_action.tex)
\newtheorem{theorem}{Theorem}[section]
\newtheorem{lemma}[theorem]{Lemma}
\newtheorem{proposition}[theorem]{Proposition}
\newtheorem{corollary}[theorem]{Corollary}
\newtheorem{claim}[theorem]{Claim}
\theoremstyle{remark}
\newtheorem{definition}[theorem]{Definition}
\newtheorem{remark}[theorem]{Remark}
\newtheorem{algorithm}[theorem]{Algorithm}
\newtheorem{certificate}[theorem]{Certificate}
\newtheorem{prediction}[theorem]{Prediction}
\newtheorem{checkitem}[theorem]{Check}

% Macros shared with hecke_action.tex (copied, not \input: that file is frozen)
\newcommand{\F}{\mathbb{F}}
\newcommand{\Z}{\mathbb{Z}}
\newcommand{\Q}{\mathbb{Q}}
\newcommand{\C}{\mathbb{C}}
\newcommand{\GL}{\mathrm{GL}}
\newcommand{\SL}{\mathrm{SL}}
\newcommand{\Sym}{\mathrm{Sym}}
\newcommand{\End}{\mathrm{End}}
\newcommand{\Hom}{\mathrm{Hom}}
\newcommand{\Ga}{\Gamma}
\newcommand{\De}{\Delta}
\newcommand{\Fl}{\F_\ell}
\newcommand{\Flr}{\F_{\ell^r}}
\newcommand{\smat}[4]{\begin{pmatrix} #1 & #2 \\ #3 & #4 \end{pmatrix}}
\DeclareMathOperator{\res}{res}
\DeclareMathOperator{\cor}{cor}
\DeclareMathOperator{\rank}{rank}
\DeclareMathOperator{\im}{im}
\DeclareMathOperator{\soc}{soc}
\DeclareMathOperator{\rad}{rad}
\DeclareMathOperator{\Ext}{Ext}
\DeclareMathOperator{\Ann}{Ann}
\DeclareMathOperator{\MaxSpec}{MaxSpec}
\DeclareMathOperator{\Frob}{Frob}
\DeclareMathOperator{\num}{num}
\DeclareMathOperator{\coker}{coker}
\DeclareMathOperator{\Fitt}{Fitt}
% Additional macros for this document
\newcommand{\T}{\mathbb{T}}
\newcommand{\m}{\mathfrak{m}}
\newcommand{\Htot}{H_{\mathrm{tot}}}
\newcommand{\Mn}{M_n}
\newcommand{\Dn}{D_n}
\newcommand{\Hc}{H_c}
\newcommand{\Hpar}{H_{\partial}}
\newcommand{\HA}[1]{[C,~\texttt{#1}]}      % citation into gl2_contract.tex by label
\newcommand{\Mnote}[1]{[M~\S#1]}            % owner's mathematics note
\newcommand{\Snote}[1]{[S~Stage~#1]}        % owner's strategy note
% STATUS tags for claims (Appendix A indexes them)
\newcommand{\PROVED}{\textsf{[proved here]}}
\newcommand{\STD}{\textsf{[standard, reference]}}
\newcommand{\CONJ}{\textsf{[conjectural]}}
\newcommand{\VERIFY}{\textsf{[to be verified computationally]}}
% v1.19.0-pre.1 (slot-19 WP-2): two tags of the companion reports, and citations into [U] and [L37]
\newcommand{\CMPD}{\textsf{[computed and compared]}}   % computed by a named package, two routes; proof-grade, not a proof here
\newcommand{\SPEC}{\textsf{[specification]}}           % an object or interface fixed for the code; no claim
\newcommand{\VERIFYloc}{\textsf{[locator not verified, rule~A2]}} % a literature locator not checked against the text (v1.19.0-pre.1)
\newcommand{\Ucite}[1]{[U,~\texttt{#1}]}               % docs/latex/hida_hecke_algebra.tex, by label
\newcommand{\Lcite}[1]{[L37,~\texttt{#1}]}             % docs/latex/l37_loewy4_example.tex, by label (not binding)

\newenvironment{decisionbox}
  {\par\medskip\noindent\begin{quote}\textbf{DECISION.}\ }
  {\end{quote}\medskip}

\title{Hecke Module Extensions in $H^\ast(\SL_2(\Z);\Sym^n\Fl^2)$\\[1ex]
\large Computational classification of all extension classes, in all computable degrees, over $\Fl$ and over $\Flr$ --- a concise mathematical specification}
\author{Fabian Januszewski\\ Universit\"at Paderborn\\ \texttt{fabian.januszewski@math.uni-paderborn.de}}
\date{Version 1.25.3 --- 2026-11-10 (revision history: Appendix~\ref{app:history})}

\begin{document}
\maketitle

\begin{abstract}
This document is the programme specification that accompanies the binding mathematical contract
\texttt{gl2\_contract.tex}, cited as [C]. In [C]'s notation and normalisations it states what it means to
compute and classify \emph{all} extension classes of Hecke modules occurring in the cohomology groups the
project computes, and how to do it with deterministic certificates at every step. The question which
$T_q$ are needed is reduced to a certified local statement, one maximal ideal at a time: the operators
computed must generate the local Hecke algebra $\T_\m$. The $\T_\m$-modules $H^\bullet_\m$ are recorded by
exhibiting them, with the invariants that make two records comparable (Loewy and socle series, active
tangent directions, Jordan types, positions of the canonical submodules): the classification data
\ref{D1}--\ref{D7}. Taking the six cohomology groups and the sequences between them from [C], it treats
the Hecke algebra acting, with its separation and generation bounds; primary decomposition and the local
algebras; the classification data and the predictions; the algorithms stage by stage, each with its
certificate; scalar extension to $\Flr$; the output record schema, over $\Fl$ and $\Z/\ell^{k}$; first
test cases; and an interface for $L$-values, with its Eisenstein instance, measures and truncated
$\ell$-adic $L$-functions. Every claim carries a STATUS tag; the tags are indexed in Appendix~\ref{app:claims}.
\end{abstract}

\tableofcontents

%======================================================================
\setcounter{section}{-1}
\section{Conventions and standing assumptions}\label{sec:conv}
%======================================================================

\subsection{Notation dictionary}
The two source notes use $p$ for the coefficient prime and $k$ for the symmetric-power
degree. This document follows the contract [C] instead, because the implementation does:

\begin{center}
\begin{tabular}{@{}lll@{}}
\toprule
object & this document / [C] & [M], [S] \\
\midrule
coefficient prime (characteristic of the field) & $\ell$ (\texttt{ell}) & $p$ \\
Hecke prime (index of $T_q$) & $q$, or $p$ (\texttt{p\_hecke}) & $q$ \\
symmetric-power degree & $n$ (\texttt{n\_sym}) & $k$ \\
modular weight & $k:=n+2$ & $k+2$ \\
coefficient module & $\Mn:=\Sym^n\Fl^2$ & $M=M_k(\F_p)$ \\
\bottomrule
\end{tabular}
\end{center}

\noindent Every formula quoted from [M] is rewritten with $k\mapsto n$, $p\mapsto\ell$.

\subsection{The standing conventions of [C]}
\noindent\textbf{[C] is \texttt{gl2\_contract.tex}}, the unified binding mathematical contract
(this revision cites it at \textbf{v1.7.3}, 2026-09-22 --- the slot-6 close, status only, on top of
v1.7.2; before that slot 5's $X$ revision v1.7.0,
its referee pass v1.7.1 and the slot close v1.7.2, which is status only, on top of slot 4's
$X$ revision v1.6.0, which resolves the eight items [C] v1.5.2 deferred, its referee pass
v1.6.1 and slot 4's close v1.6.2, and the slot-3 chain v1.5.0 / v1.5.1 /
v1.5.2; \emph{none of the eight adds, renames or removes a label}, so every label cited here
exists there and the v1.3.3 labels are all preserved. What v1.6.0 changes that this
document depends on: \HA{obl:coneexact}(i) is a statement on the stored equivariant-cochain basis with
the flat identities recorded, (iv)'s explicit $\im=\ker$ check is degree-generic, (v) and
\HA{eq:conedims}'s first identity no longer need $H^2_{\mathrm{eq}}=0$, and
\HA{sec:recordschema} gains item~6, the extension record and its six rules --- which
\S\ref{sec:schema} of this document implements. What v1.7.0 changes: \HA{obl:coneexact}(iii)'s
$\ell\in\{2,3\}$ carve-out on the restriction half is \emph{dropped}, the identity asserted at
every $\ell$ with the proof written in, and \HA{sec:recordschema}'s items are restated for the
record schema \texttt{hecke\_ext/v2}, which \S\ref{sec:schema} now names),
which supersedes \texttt{hecke\_action.tex} (the document [C] named through v1.2.2) and
\texttt{monodromy\_v5.tex} (cited here as [v5]); both are frozen. Every label this document
cites as \HA{$\cdot$} is defined there.
\begin{itemize}
\item $\Ga=\SL_2(\Z)$, $\De=\{g\in M_2(\Z):\det g>0\}$; $T=\smat1101$, $S=\smat0{-1}10$,
  $ST$, $TS$ as in \HA{sec:model}. $\Ga_\infty=\langle T,-I\rangle$.
\item $\rho_n(A)=\Sym^n(A)$ in the \emph{column} convention $A\cdot x=ax+cy$, $A\cdot y=bx+dy$,
  basis $u_j=x^{n-j}y^j$, no determinant twist, extended multiplicatively to the monoid $\De$
  \HA{def:rhon}, \HA{prop:symconv}. This coincides with the action of \Mnote{0}.
\item $\ell\ge5$ throughout (hypothesis (H3) of \HA{sec:conventions}), so that
  $H^\bullet(K,R;L)\cong H^\bullet(\Ga;\Mn)$ \HA{rmk:orbifold}. $n$ is \emph{even}: for odd $n$
  the central relation forces $H^\bullet=0$ \HA{rmk:negI}, and the pipeline returns empty
  matrices.
\item $T_p=T_{\Ga\,\mathrm{diag}(1,p)\Ga}$ is the operator of \HA{def:hecke}, assembled by
  \HA{alg:assembly}, normalised by $T_p|_{H^0(\Sym^0)}=p+1$, the Eisenstein eigenvalue
  $1+p^{n+1}$ on $H^1$, and the recursion
  $T_{p^{r+1}}=T_pT_{p^r}-p^{n+1}T_{p^{r-1}}$ for the \emph{full} operators \HA{thm:relations}
  ($T_{p^2}=T_{\De^{\mathrm{prim}}_{p^2}}+p^nI$). The bad-prime operator is $U_\ell:=T_\ell$
  \HA{cor:ulrecursion}. Prime indices suffice: $T_{p^r}\in\Fl[T_p]$ and $T_{MN}=T_MT_N$.
\item The \emph{convention warning} of \Mnote{4} (operator vs.\ its transpose) is settled by
  [C]: the oracles $a_2(\De)=-24$ at $n=10$ and $1+p^{n+1}$ on the Eisenstein class are part of
  the acceptance gate (\texttt{HeckeOracleValidation}, \texttt{HeckeEndToEndValidation}).
  Consequently the cuspidal eigenvalues on $H^1$ are $a_q(f)$ and the boundary characters of
  \S\ref{sec:boundary} are as stated there, not swapped.
\end{itemize}

\subsection{Sym versus divided powers}\label{sec:symdiv}
For $n\ge\ell$ the $\Ga$-module $\Mn$ is not self-dual; its dual is the divided-power module
$\Dn=(\Sym^n\Fl^2)^\vee$ \Mnote{0}. The pipeline computes $\Mn$ today; $\Dn$ enters through
the duality \emph{predictions} of \HA{cl:pd} until it is \emph{computed} --- as a first-class
contragredient local system $\rho_n^\vee(g)=\rho_n(\operatorname{adj}g)^{T}$ with the pairing map
at cochain and cohomology level, \HA{def:contragredient}, \HA{prop:pairing} --- in phase E7 of
[X] v1.5.1 (slots 6--7; owner decision OQ4, decided). The adjointness
$\langle A_c(T_p)x,y\rangle=\langle x,A(\Dn)(T_p)y\rangle$ of \HA{cl:pd}(2) is then the
acceptance gate for $X_3$ at $\ell\ge5$ (owner decision OP3), and the socle$\leftrightarrow$Loewy
reversal of Prediction~\ref{pr:dual} is asserted by \texttt{HeckeDualityValidation} (v1.4.0;
through v1.3.3 this sentence read ``$\Dn$ enters solely through the predictions'' and cited Open
Question~4 as open).

\subsection{What ``certified'' means}
Repository policy (\texttt{FpPoly.h}, \texttt{HeckeModuleDecomposition.h}): randomisation may
\emph{find} an object (a factor, a witness, a cyclic vector) but never \emph{certifies} it; every
reported statement is re-derived by a deterministic, unconditional check that a third party can
repeat from the emitted data alone. Every stage below names its certificate.

%======================================================================
\section{The cohomologies computable now, and the sequences relating them}\label{sec:coh}
%======================================================================

\subsection{The six groups, and what [C] proves about them}\label{sec:imported}
\emph{Recalled, not claimed here.} For $M=\Mn$ the \textbf{six groups} are
\begin{gather*}
X_1=H^0(\Ga;M),\qquad X_2=H^0(\partial;M),\qquad X_3=H^1_c(\Ga;M),\\
X_4=H^1(\Ga;M),\qquad X_5=H^1(\partial;M),\qquad X_6=H^2_c(\Ga;M),
\end{gather*}
with $\Htot:=\bigoplus_iX_i$, $d:=\dim\Htot$, $H^\bullet(\partial;M):=H^\bullet(\Ga_\infty;M)$
and $H^\bullet_c$ the cohomology of the (shifted) cone of the restriction $r^\bullet$. This is
\HA{eq:sixgroups} of \HA{sec:six} verbatim; the two defining pointers are \HA{def:bdrygroup}
for the boundary groups and \HA{def:pair} with \HA{lem:conerelative}(3) for the compactly
supported ones. The six-term exact sequence relating them is \HA{eq:les}, its dimension
bookkeeping \HA{cl:les}(2), and $H^1_!:=\ker r^1=\im(X_3\to X_4)$ is defined there. [C] writes
the restriction maps $r^0,r^1$ with the same letter and the $T$-invariant
$N=y^\ell-x^{\ell-1}y$ as $\mathrm{Nm}$, reserving $N=p^\nu$ for the Hecke level (its notation
collision~C18).

\paragraph{Nine statements moved into [C], and why this document no longer states them.}
The amalgam presentation, the dimension of $H^1$ with its Dickson correction, the torsion
origin of the Dickson classes, the boundary groups, the six-term sequence, the boundary
characters and Poincar\'e duality are \emph{facts about the objects} [C] constructs, not part
of the classification programme this document specifies. The scope cut that moved them is
recorded in [C] as erratum~E18 (\HA{err:e18}). Until v1.3 this document kept its own copies,
so nine labels were defined in both places and a citation of the form \texttt{[C,~cl:pd]} did
not resolve to a unique statement; v1.3 removes the copies. Nothing was lost: every one of the
nine is preserved in [C] under the same label, five of them now carrying a proof this document
did not give --- the amalgam presentation, the Bockstein, the boundary groups, the six-term
sequence, and clauses~(1) and~(3) of the boundary characters.

\begin{center}\footnotesize
\begin{tabular}{@{}>{\raggedright\arraybackslash}p{7.8cm}>{\raggedright\arraybackslash}p{5.2cm}@{}}
\toprule
statement, as this document stated it through v1.2.2 & now cited as \\
\midrule
\S1.1 ``Six groups'' --- the six $X_i$, $\Htot$, $d$ & \HA{sec:six}, \HA{eq:sixgroups} \\
Claim ``Amalgam presentation'' & \HA{cl:amalgam} \\
Claim ``Dimension of $H^1$'' & \HA{cl:dimH1} \\
Claim ``Torsion origin of the Dickson classes'' & \HA{cl:bockstein} \\
Claim ``Boundary groups'' & \HA{cl:boundary} \\
Claim ``The sequence (3.1)'' & \HA{cl:les} \\
the displayed sequence of that claim & \HA{eq:les} \\
Claim ``Boundary eigenvalues'' & \HA{cl:chi} \\
Claim ``Poincar\'e duality'' & \HA{cl:pd} \\
\bottomrule
\end{tabular}
\end{center}

\paragraph{Four of the nine are corrected in [C]; the v1.2.2 forms are superseded.}
A reader holding an older PDF of this document should read these four in [C] and not here.
\begin{enumerate}
\item \textbf{The Dickson-locus sentence of the boundary-eigenvalue claim is false}
  ([C] erratum~E21, \HA{rmk:chicoincide}). Through v1.2.2 this document concluded
  ``$\chi_1\equiv\chi_2$ at every Dickson locus''. On the Dickson locus
  $n=a(\ell+1)+b(\ell^2-\ell)$ one has $n\equiv2a\pmod{\ell-1}$, so the criterion is
  $(\ell-1)\mid(n-2)$, \emph{not} always: it holds at the first Dickson class $n=\ell+1$ ---
  this document's own example $(n,\ell)=(6,5)$ --- and fails already at $(12,5)$, where
  $\chi_1(2)\equiv3$ and $\chi_2(2)\equiv2$ in $\F_5$. What survives is the use it was put to:
  the first Dickson class never separates the left- from the right-coset indexing, so the
  separation comes from the Eisenstein eigenvalue $1+q^{n+1}$ (\HA{prop:bdrychar}(3)), which
  is how \ref{add:bdry} settles the convention.
\item \textbf{``$H^i(\partial;M)=0$ for $i\ge2$'' is false at $\ell=2$}
  (\HA{rmk:bdryrepairs}(i), \HA{lem:ell2unbounded}). The hypothesis that carries it is
  $2\in R^\times$, and \HA{cl:boundary} now states it. This is the same $\ell=2$ unboundedness
  this document records in \ref{add:cone}; under its standing hypothesis $\ell\ge5$ nothing here is
  affected.
\item \textbf{The hypotheses of the six-term sequence changed in both directions}
  (\HA{rmk:sixscope}). $H^0_c=0$ is \emph{unconditional} --- \HA{lem:conerelative}(4) proves
  it with no hypothesis on $\ell$ and none on $n$ --- so [C] does not list it, where this
  document did; and ``$n$ even'' is needed for the six groups to be non-zero, not for
  exactness. But $H^2(\Ga;M)=0$ is \emph{kept}, and is not removable: \HA{lem:h2torsion}(2)
  exhibits $(\ell,n)=(2,0)$, a pair with $n$ even, at which $H^q(\Ga;\F_2)\ne0$ for every
  $q\ge0$.
\item \textbf{Poincar\'e duality holds in the adjointness form only} (\HA{rmk:pdscope}(1)).
  This document's ``$T_q$ self-adjoint once $\Dn$ carries the adjugate $\De$-structure
  $g\mapsto\det(g)g^{-1}$'' is true of the adjointness identity
  $\langle\rho_n(g)v,\varphi\rangle=\langle v,\rho_n^\vee(\det(g)g^{-1})\varphi\rangle$, which
  is exact, integral and carries no determinant factor, and \emph{false} of the diagonal one
  $\langle\rho_n(g)v,\rho_n^\vee(g)\varphi\rangle=(\det g)^n\langle v,\varphi\rangle$, which
  carries the twist $q^{\,n}$ for $g=\mathrm{diag}(1,q)$. \HA{cl:pd}(2) is stated in the
  former, and it is that form Prediction~\ref{pr:dual} and identity~5 of
  \S\ref{sec:identities} rest on. [C] also carries the claim as \textsc{open} (owner Open
  Question~4): its coefficient-level input \HA{prop:pairing} is proved and its cohomological
  content is not.
\end{enumerate}

\subsection{What the code computes today}\label{sec:today}
The following is a statement about the tree at commit \texttt{d54c589} (branch
\texttt{matrix-engine}), checked against \texttt{Hecke/README.md} and the headers, and agreeing
with the independent audit \texttt{docs/implementation/\allowbreak hecke\_extensions\_\allowbreak capability\_inventory.md}.
\begin{itemize}
\item $X_1,X_4$ and $H^2(\Ga;M)$ ($=0$, asserted as a $0\times0$ block) on the 7-vertex model,
  with bases, projectors and the Hecke matrices $A^{(q)}(T_p)$ for $q\in\{0,1,2\}$ and every
  prime $p$ (including $p=\ell$), plus $T_{p^r},T_n,T(p,p)$ (\texttt{HeckeOperatorFamily}).
  This is the whole \emph{level (b)} of the plan.
\item $X_2,X_5$ as equivariant cohomology of the circle with monodromy $\rho_n(T)$
  (\texttt{LocalSystemPresets::\allowbreak CreateS1Generic}), and the restriction maps $r^0,r^1$ as
  matrices between the stored bases (\texttt{Equivariant\allowbreak Cochain\allowbreak Complex\allowbreak Morphism\allowbreak OnK}, exercised
  by \texttt{Test\allowbreak SL2Z\_\allowbreak To\_\allowbreak S1Cusp\_\allowbreak Restriction\allowbreak Morphism}). \textbf{No Hecke action is implemented
  on $X_2,X_5$}: \texttt{HeckeOperatorOnK} is bound to the 7-vertex model
  (\texttt{RequireModel(\allowbreak SL2Z7VertexModel())}).
\item \textbf{$X_3,X_6$ are not computable yet}: the mapping cone exists only for the
  non-equivariant \texttt{CochainComplex} (\texttt{Simplicial\allowbreak Cohomology.h},
  \texttt{Cone\allowbreak Complex()}), not for \texttt{Equivariant\allowbreak Cochain\allowbreak Complex\allowbreak Morphism\allowbreak OnK}. The algebraic
  cone for equivariant complexes is \HA{def:pair} and \HA{lem:conerelative}, and its Hecke
  action \HA{def:coneaction}, \HA{thm:conewelldef} (\HA{sec:cone}); building it is a mechanical
  addition. ([v5] \S8 carried the construction before the scope cut and is now frozen.)
\item The involution $\iota=\mathrm{diag}(-1,1)$ \Mnote{0} has \textbf{no production
  implementation}; it is built, and [C] \HA{obl:iota} is discharged, by the two validations
  \texttt{HeckeIotaTauSymmetryValidation} and \texttt{HeckeIotaSignValidation}.
\item Level (c): the primary decomposition of $H^q$ under one operator or a commuting family,
  with certified irreducible factors, generalized-Jordan partitions, residue degrees
  (\S\ref{sec:primary}), and the JSON report. No local algebra, tangent space, Loewy or socle
  series, commutant, or $\Flr$ arithmetic exists: Stages 3--5 of \S\ref{sec:alg} are entirely new
code (inventory \S1); the substrate they need (\texttt{Matrix} products / RREF / kernel,
\texttt{SolveColumnsExact}, \texttt{FpPolyT}) is present.
\end{itemize}

\noindent\textbf{Addendum, v1.3.3 (2026-09-16, the close of slot 1 of the extension stage;
branch \texttt{hecke-extensions}, $77$ registered validations).} The statement above is kept as
the record of \texttt{d54c589}; what slot 1 changed, with no golden value moved:
\begin{itemize}
\item \textbf{A Hecke action on $X_2,X_5$ now exists}: \texttt{Hecke/HeckeBoundaryOperator}
  computes $\Theta^\partial_g$ of [C] \HA{def:bdryop} on the circle model from the certified
  $\Gamma$-level $\tau_\partial$ of \texttt{Hecke/HeckeBoundaryDatum} ([C] \HA{def:taubdry}),
  with [C] \HA{obl:taubdry} discharged in full by \texttt{HeckeBoundaryTauValidation} and
  \texttt{HeckeBoundaryOperatorValidation}; the record fields \texttt{boundary\_model = circle}
  and \texttt{boundary\_is\_group\_cohomology} $=(\ell\ne2)$ say which object is computed. The
  operator is not yet integrated into the report of \S\ref{sec:schema} (E4b, slot 2), and the
  groups \texttt{H0\_bd}, \texttt{H1\_bd} stay \texttt{unavailable} in the record until then.
\item \textbf{$X_3,X_6$ are still not computable}: no $\eta$, no equivariant cone.
\item \textbf{The locality certificate of Certificate~\ref{sec:localcert} exists as a read-only
  function} (\texttt{hecke::CertifyLocality}, called by nothing in \texttt{compute()} until E0b):
  on the whole golden grid ($236$ decompositions, $348$ blocks) every block certifies local with
  the shipped witness, so E0b's fix moves no decomposition hash (\S\ref{sec:gap}'s defect is
  real --- $48$ constructed non-local families that the shipped predicate accepts are rejected
  --- but does not occur on the grid); and the $(160,691)$ probe verdict is at Claim~\ref{cl:bk}.
\item \textbf{The Sturm bound $B'(k,\ell)$ is computed} (\texttt{HeckeExt/HeckeSturmBound}, the
  first unit of the extension layer \texttt{HeckeExt/}; Claim~\ref{cl:bprime}), and
  \textbf{the identity certificate} of the plan's \S1.3 in its slot-1 subset
  (\texttt{HeckeExt/HeckeIdentityCertificate}: the $\dim H^0$-robust Eichler--Shimura count,
  commutativity, $d_2=0$, the $q^2$ recursion against the independent double-coset route) runs
  on ten pipeline points (\texttt{HeckeIdentityValidation}).
\item \textbf{A reach limit of $U_\ell$, recorded}: at $n=16$ the exact-integer coset
  construction of \texttt{HeckeCosetTables} completes for $U_\ell$ at
  $\ell\in\{691,1009,2003,3617\}$ ($1$--$34$\,s, $0.1$--$0.7$\,GB) and overflows its
  \texttt{int64} tier (``exact multiply overflow'') at $\ell\ge5003$, independently of $n$; the
  plan's $(16,43867)$ RSS point was a placeholder-model prediction that cannot be run. The two
  measured points $(128,691)$: $2717$\,MB and $(256,131)$: $1677$\,MB fit
  $\mathrm{rss}=(\ell+1)(c_0+c_1(n+1)^2)$ with $c_1\approx1.8$--$1.9\cdot10^{-4}$ and a $c_0$
  the two fits disagree on by $4.4\times$; the constants are re-fitted before E1b freezes them.
\end{itemize}

\noindent\textbf{Addendum, v1.4.0 (2026-09-16, slot 2 of the extension stage, as reported by its
packages WP-1, WP-2, WP-4, WP-5 and WP-7).} What slot 2 built, and the one golden re-baseline
it made:
\begin{itemize}
\item \textbf{E0b --- the locality certificate runs inside the refinement loop} (WP-1;
  \texttt{HeckeModuleDecomposition::ComputeMultiOp} = \texttt{RefineToFixedPoint} +
  \texttt{RunLocalityStage} + \texttt{FinalizeChangeOfBasis}): one candidate sequence --- the
  single-operator witnesses with $\deg\pi_j=f$ first, the all-ones combination, then the
  fixed-seed combinations --- capped by the budget; a candidate with $\ge2$ distinct irreducible
  factors \emph{splits} the block at the parent's position in \texttt{CanonicalLess} order; the
  first clause-(1) candidate fixes \texttt{residue\_witness} exactly as before (the shipped
  semantics, byte for byte), so on the golden grid, where candidate~$0$ certifies, no block list
  moves. The mathematical settings layer \texttt{Hecke/HeckeSettings} ([X] \S1.5, owner decision
  R1) exists with the key \texttt{locality} (\texttt{budget=<N>}, default $64$, or
  \texttt{strict}); \texttt{strict} parses and validates and \emph{throws}
  \texttt{not\_implemented\_in\_this\_build} until slot~3 (WP-6), never degrading to budget mode
  (slot-2 owner confirmation~3). \texttt{hecke\_report/v2} (WP-2): the per-block keys
  \texttt{local}, \texttt{witness\_poly}, \texttt{coefficient\_field} $=\{$\texttt{w\_hat\_exponent},
  \texttt{P}$\}$ with \texttt{P} keyed by prime, \texttt{level\_c.locality} $=\{$\texttt{mode},
  \texttt{budget}, \texttt{source}, \texttt{complete}, \texttt{splits}$\}$, the corrected
  multiplicity note, and the four per-complex degree fields of \S\ref{sec:schema} --- the stage's
  \emph{one} re-baseline, the $384$ report texts and nothing else, explained key path by key path
  by \texttt{golden\_explain.py}.
\item \textbf{E1a --- the local algebra and certificates (a), (b)} (WP-4, WP-5;
  \texttt{HeckeExt/HeckeLocalAlgebra}, \texttt{HeckeExt/HeckeGenerationCertificates}, namespace
  \texttt{hecke\_ext}): Algorithm~\ref{alg:stage3} as implemented (its ``as implemented''
  paragraph), Certificate~\ref{cert:cyclic} with the canonical generator,
  Certificate~\ref{cert:comm} with the socle bound, the structured system on a cyclic-most
  element and the resumable \texttt{CommutantState} (R3 instance~3), the cross-block evaluation
  test \texttt{SameMaximalIdeal}, and the pipeline-free map \texttt{generation\_verdict} of the
  default configuration; \texttt{HeckeExtAlgebraValidation} on constructed truths ($0.10$\,s).
  The real-block arms are E1b's.
\item \textbf{E4b --- $H^1_!$, Harder and the diagnostics} (WP-7;
  \texttt{Hecke/HeckeBoundaryDiagnostics}, \texttt{HeckeExt/HeckeBoundary},
  \texttt{HeckeBoundaryParabolicValidation}, record \texttt{hecke\_boundary\_parabolic/v1} with
  the per-complex degree fields): $H^1_!=\ker r^1$ computed on $315$ $(\text{cell},p,\nu)$
  pairs; Harder's criterion asserted on the Kummer list in the form ``$r^1=0$ on the joint
  eigenspace'' (Claim~\ref{cl:harder}); \texttt{cochain\_commutes} and \texttt{eta\_zero} both
  false on $315/315$ (\ref{add:cone}); and the two refutations of this revision
  (\S\ref{sec:tests}, \ref{D6}). $X_2,X_5$ are still not in the record of \S\ref{sec:schema}
  (E2).
\item \textbf{Still not computable}: $X_3,X_6$ (the $\eta$-table is WP-8, the first part of
  E5), the production $\iota$ (E3, slot~3), $\Dn$ (E7). The $U_\ell$ reach repair (E1b-0) and the
  RSS constants are WP-3's and are not recorded here.
\end{itemize}

\noindent\textbf{Addendum, v1.4.2 (2026-09-17, slot 3 of the extension stage, as reported by its
packages WP-1, WP-2, WP-3, WP-5 and WP-6).} What slot 3 built; it made \textbf{no} golden
re-baseline and no golden emission:
\begin{itemize}
\item \textbf{E1b-0 --- the $U_\ell$ reach} (WP-1): the overflow at $\ell\ge5003$ was in
  [C] \HA{alg:tau} step~1's fundamental-domain reduction, not in the coset construction; step~1
  now escalates per operation to an \texttt{\_\_int128} tier, bit-identically wherever the
  \texttt{int64} tier completes, and the reach is $\ell\le23170$, past which the $\tau$-table is
  refused by its own certification ([C] \HA{lem:bounds}, \HA{thm:backend}). The RSS model is
  frozen: $\mathrm{rss}=(\ell+1)(c_0+c_1(n+1)^2)$\,MB with $c_0=0.123152$,
  $c_1=2.24985\cdot10^{-4}$ (\texttt{u\_ell\_rss/v2-2026-09-17}, $16$ measured points, residuals
  $+17.9\,\%/-15.9\,\%$); at $M=2048$ it admits $U_\ell$ for $n\le804$ at $\ell\le13$, $n\le110$
  at $\ell=691$, $n\le18$ at $\ell=10007$, and the ten Kummer verdicts are unchanged (seven
  admitted, three refused). The setting \texttt{u\_ell} is registered.
\item \textbf{E1b --- generation certificate (c) on real blocks} (WP-2,
  \texttt{HeckeExt/HeckeSturmSweep}): Certificate~\ref{cert:sturm} sweeps every prime
  $q\le B'(k,\ell)$ by a membership test in $A_S$; at $(10,691)$ the algebra is
  $\F_{691}[\varepsilon]/\varepsilon^2$, (a) fails, (b) is \texttt{excluded\_by\_socle\_bound}
  with the naive $\dim C=5$, and (c) is \texttt{sturm($B'=2$)}, \texttt{final = true}; at
  $(n,\ell)=(22,5),(26,7)$ the $\{2,3\}$ sweep is rejected and the $B'$ sweep accepted; both
  unavailable modes run on a forced cap with the real $A^{(1)}(T_q)$; Stage~1 consults the
  frozen $U_\ell$ admission. The settings are registered as \texttt{sturm} and
  \texttt{sturm\_unavailable}, which is how \S\ref{sec:schema} names them since v1.5.0; the
  default configuration never produces \texttt{"heuristic"}.
\item \textbf{E0b-S --- strict locality mode} (WP-3): Certificate~\ref{sec:localcert}'s strict
  mode as specified, the covering example verified, the golden grid identical under it. The
  implementation reports a strict-mode witness as a matrix and checks that it commutes with
  the family; it asserts $\dim_{\Fl}\bar A=\operatorname{lcm}_p\deg\pi_p$ on a local block.
  \textbf{Both are now statements of this document} (v1.5.0): the hypothesis a
  \emph{reported} witness must satisfy is stated in Certificate~\ref{sec:localcert}, and
  the sharpness identity is Claim~\ref{cl:sharpbound}. With them, the mode-dependence of
  the sub-block order after a split (Certificate~\ref{sec:localcert}, closing paragraph)
  and the two \texttt{phase\_reached} keys of \S\ref{sec:schema}.
\item \textbf{E5-2, E5-3 --- the cone and its Hecke action} (WP-5, WP-6;
  \texttt{EquivariantMappingCone}, \texttt{Hecke/HeckeConeAction}): $H^\bullet_c$ with the long
  exact sequence and its maps in stored bases, and the operator on it induced by $\Phi_g$ with
  the derived $h=\tilde\eta$ ([C] \HA{def:coneaction}); every clause of [C]
  \HA{obl:coneexact} but the cross-routes (vii) is evaluated and \HA{obl:eta} is discharged. At
  $(10,691)$: $\dim H^q_c=[0,3,0]$, $\rank r^1=1$, $\dim H^1_!=2$; at the ten Kummer cells
  $\dim H^1_c=1+2\dim S_{n+2}$ and $H^2_c=0$. So $X_3,X_6$ are \emph{computable} on the model
  complexes; they are not yet in the record of \S\ref{sec:schema} (E2).
\item \textbf{Still not computable}: the production $\iota$ (E3, moved to slot~4), $\Dn$ (E7);
  the record \texttt{hecke\_ext/v1} (E2, from slot~4). \emph{Read as of the slot-3 close}
  (v1.6.1, the referee pass): this list is a snapshot of that date and it is the last
  addendum in this section, so a reader arriving here should not take it for the present
  state --- slot~4 built the production $\iota$ (E3) and wrote the record, and slot~5 moved
  it to \texttt{hecke\_ext/v2} with the boundary, cone and per-sign objects; \emph{only} $\Dn$
  (E7) is still not computable. The version block above carries what slots~4 and~5 built; the
  slot-4 and slot-5 addenda are the close's. \emph{(v1.7.2, the slot-6 close:)} $\Dn$ is
  computable too, with its Hecke action and the pairings (slot-6 WP-1, WP-3), so nothing on
  this list remains.
\end{itemize}

\begin{decisionbox}
``All cohomologies the project can compute at this stage'' means, with Hecke action:
$X_1$, $X_4$ (and the vanishing $H^2$). The classification programme of
\S\ref{sec:alg} is stated for all six $X_i$; its Stages 0--1 list the three additions that
make the other four available (boundary Hecke action, equivariant cone, $\iota$), all of
which are formula-level. Two of them need no new mathematics; the cone needs a \emph{definition},
now fixed: the Hecke action on $X_3,X_6$ is the one induced by $\Phi_g$ from the certified triple
$(\tau,\tau_\partial,\eta)$ over $\Z$ with the cochain homotopy \emph{derived} (\ref{add:cone}), which is
what makes the operator on $X_3$ well defined at all --- with a freely chosen cochain homotopy it is
not (\ref{add:cone}, third bullet). Until the three additions land, the record of
\S\ref{sec:schema} carries the four missing groups as \texttt{null} with an explicit
\texttt{"unavailable"} marker, and the extension classes ``in the cone'' are reported as
\emph{predicted} from $r^1$ (\HA{cl:les}), not computed.
\end{decisionbox}

\subsection{Boundary characters}\label{sec:boundary}
The characters $\chi_a(q):=q^a+q^{\,n+1-a}$ by which the Hecke algebra acts on $X_2$ and $X_5$,
their eigenvectors $x^{n-b\ell}\mathrm{Nm}^b$ on $X_2=M^T$, the Eisenstein value
$\chi_0(q)=1+q^{n+1}$, the right-coset variant $\chi_{b+1}$ that settles the convention hedge of
\Mnote{4}, and the coincidence criterion $\chi_a\equiv\chi_{a'}$ iff $a\equiv a'$ or
$a\equiv n+1-a'\pmod{\ell-1}$ are \HA{cl:chi}, with \HA{prop:bdrychar} for the eigenvector
computation. The Dickson-locus corollary this document drew from the criterion through v1.2.2 is
\emph{false} and is corrected in \HA{rmk:chicoincide}; see item~1 of \S\ref{sec:imported}.

\begin{claim}[Harder's denominator, Ribet's converse]\label{cl:harder}\STD\
For $n<\ell-1$: $r^1=0$ on the Eisenstein class of $H^1(\Ga;\Mn)$ iff
$\ell\mid\num(B_{k}/k)$, $k=n+2$ (Harder \cite{harder3}; Ribet \cite{ribet}). Then the reduced
Eisenstein class lies in $H^1_!$ and $H^1_{\m_{\mathrm{Eis}}}$ is a non-split $\T$-module ---
the prototype of the extensions classified here \Mnote{5}. All cases with $n<\ell-1$ and
$k\le26$:
$(k,\ell)$ one of $(12,691)$, $(16,3617)$, $(18,43867)$, $(20,283)$, $(20,617)$, $(22,131)$,
$(22,593)$, $(24,103)$, $(24,2294797)$, $(26,657931)$ \cite{washington} ($B_{24}=-236364091/2730$,
$236364091=103\cdot2294797$).
\emph{The mod-$\ell$ reading (v1.4.0).} ``The Eisenstein class'' is the integral class --- the
reduction of the Eisenstein eigenline of $M_k(\Z_\ell)$ into $H^1(\Ga;\Mn(\Fl))$ --- and the
\STD{} tag covers that reading. Mod $\ell$ it is \emph{not} a Hecke datum at a Kummer prime: the
joint eigenspace $E$ of the $T_q$ for the eigenvalue $1+q^{n+1}$ on $H^1$ is two-dimensional
there (the socle of the non-split block, \ref{D6}), so the statement a mod-$\ell$ computation can
check, and does check, is that $r^1$ vanishes on \emph{all} of $E$; the state
\texttt{inconclusive} is reserved for $0<\rank(r^1E)<\dim E$ (not observed).
\textbf{Verified 2026-09-16} (\texttt{HeckeBoundaryParabolicValidation}, slot-2 WP-7):
$\rank(r^1E)=0$ on all ten cells of the list, and the converse $\rank(r^1E)=1$ on the $27$
non-Kummer cells with $n<\ell-1$ of its grid --- $37/37$ consistent, $0$ inconclusive. What this
does \emph{not} say is $H^1_!=H^1$: see \S\ref{sec:tests}.

\emph{The criterion's hypothesis and the divisibility predicate's domain are different
things} (v1.5.0; slot-4 WP-7 finding~4). This claim is stated for $n<\ell-1$, i.e.\
$k=n+2\le\ell$; the \emph{question} ``does $\ell$ divide $\num(B_k/k)$'' is meaningful for
every even $k\ge2$ and every $\ell$, and a predicate answering it must therefore be total
where the criterion is not. A naive evaluation of $B_k\bmod\ell$ by the $\Fl$-recursion
$\sum_{j}\binom{k+1}{j}B_j=0$ is valid only for $k<\ell-1$, and two provable branches
cover the rest (v1.5.1, correcting v1.5.0, which named only the two instances $k=\ell-1$
and $k=\ell$ and so left every even $k$ with $(\ell-1)\mid k$ and $k>\ell-1$ --- at
$\ell\in\{2,3\}$ that is \emph{every} even $k$ --- with no branch at all).
\textbf{Whenever $(\ell-1)\mid k$}, von Staudt--Clausen puts $\ell$ into the
\emph{denominator} of $B_k$ (indeed $\ell\,B_{k}\equiv-1\pmod\ell$), so $\ell\nmid\num(B_k)$
and a fortiori $\ell\nmid\num(B_k/k)$, which divides $\num(B_k)$: the answer is
\texttt{false} outright and no recursion is run. That branch subsumes $k=\ell-1$, the
instance at the top of this claim's own range; and since $\ell-1\in\{1,2\}$ divides every
even $k$, it is the branch that decides \emph{every} even $k$ at $\ell\in\{2,3\}$ --- in
particular this claim's own edge cell $(n,\ell)=(0,2)$, where $k=2=\ell$ and
$\num(B_2/2)=\num(1/12)=1$, so the answer is \texttt{false}. (The remaining degenerate case
belongs to \emph{odd} $k\ge3$, which is outside this claim: there $B_k=0$, the numerator
vanishes and every prime divides it, a degenerate \texttt{true}.) \textbf{Otherwise}
$k\not\equiv0\pmod{\ell-1}$, and the Kummer congruence $B_{k}/k\equiv
B_{k'}/k'\pmod\ell$ for $k\equiv k'\pmod{\ell-1}$ reduces
the question to $k'=k\bmod(\ell-1)<\ell-1$, where the recursion is valid. \textbf{Conformance}: a shipped predicate must handle those branches explicitly
and say which one it took, and an implementation that pins a table of $\num(B_k/k)$ for
small $k$ (\texttt{HeckeExt/HeckeBoundary}, $k\le26$) is answering the same question over a
finite range, not a different one. Implemented as
\texttt{HeckeExt/HeckePredictions}' \texttt{bernoulli\_mod} / \texttt{kummer\_divides}
(slot-4 WP-7, with \texttt{HeckePredictionsValidation}); the rewiring of the pinned table to
it is E6's, because it changes what \texttt{harder\_applicable} records for $k>26$.
\end{claim}

%======================================================================
\section{The Hecke algebra acting}\label{sec:hecke}
%======================================================================

\subsection{Definitions}
Let $\T_\infty\subseteq\End_{\Fl}(\Htot)$ be the $\Fl$-algebra generated by $\{T_q\}_{q\text{ prime}}$
(all $q$, $q=\ell$ included), and for a finite set $S$ of primes $\T_S\subseteq\T_\infty$ the
subalgebra generated by $\{T_q\}_{q\in S}$. Both are commutative artinian
(\HA{obl:algebra}; finite-dimensional), so
\[
\T_\infty=\prod_{\m\in\MaxSpec\T_\infty}\T_{\m},\qquad \Htot=\bigoplus_\m H_{\m},\qquad
H_\m:=\{v:\m^Nv=0,\ N\gg0\},
\]
with residue fields $\T_\infty/\m\cong\Flr$, $r=r_\m$. The \emph{character} of $\m$ is the
Galois orbit of $(a_q(\m))_q\in\Flr^{\,\mathbb P}$, $a_q(\m):=T_q\bmod\m$. Equivariance of
\HA{eq:les} makes each $X_i$ a $\T_\infty$-module and $X_{i,\m}:=X_i\cap H_\m$.

\begin{claim}[Locality of extensions]\label{cl:local}\PROVED\
$\Ext^1_{\T_\infty}(A,B)=0$ for $\T_\infty$-modules $A,B$ supported at different maximal
ideals. Hence every non-split extension of Hecke modules occurring in any $X_i$, or between
terms of \HA{eq:les}, lives inside a single $H_\m$, and $H_\m$ is a module over the local
ring $\T_\m$.
\end{claim}
\begin{proof}
$\T_\infty=\prod\T_\m$ is a product of rings; the idempotent $e_\m$ acts as $1$ on $A$ and $0$
on $B$, so any extension $0\to B\to E\to A\to0$ splits via $e_\m E$.
\end{proof}

\begin{claim}[Generation criterion]\label{cl:gen}\PROVED\
$\T_S=\T_\infty$ iff (a) the natural map $\MaxSpec\T_\infty\to\MaxSpec\T_S$ is a bijection
(``$S$ separates''), and (b) for every $\m$ the images of $T_q-a_q(\m)$, $q\in S$, span
$\m\T_\m/\m^2\T_\m$ over $\T_\infty/\m$ (here $a_q(\m)$ is read inside $\T_\m$ through its
coefficient field, Remark~\ref{rmk:coef}). When (b) holds at $\m$, the $\T_S$-module structure of
$H_\m$ \emph{is} its $\T_\infty$-module structure, and $\dim_{\Flr}\m/\m^2$ is the minimal
number of operators that can generate $\T_\m$ over its residue field.
\end{claim}
\begin{proof}
$\T_S\subseteq\T_\infty$ is a finite (integral) extension of artinian rings. If (a) holds,
$\T_S=\prod\T_{S,\m'}$ with $\m'=\m\cap\T_S$ and $\T_{S,\m'}\to\T_\m$ local. Nakayama for the
finite $\T_{S,\m'}$-module $\T_\m$: surjectivity iff $\T_{S,\m'}\to\T_\m/\m'\T_\m$ is onto; since
$\T_\m/\m'\T_\m$ is local with residue field $\T_\m/\m$ and (b) says $\m'\T_\m+\m^2=\m$,
Nakayama again gives $\m'\T_\m=\m$ and then $\T_\m/\m'\T_\m=\T_\infty/\m$, which is the image of
$\T_{S,\m'}$ iff the residue field extension is generated by the $a_q$, $q\in S$ --- part of (a)
(the residue field of $\m'$ equals that of $\m$ exactly when $\MaxSpec$ maps bijectively with
the same residue degrees; see Remark~\ref{rmk:resdeg}). Conversely $\T_S=\T_\infty$ trivially
gives (a),(b).
\end{proof}

\begin{remark}\label{rmk:resdeg}
Condition (a) contains a residue-degree clause: $S$ must separate $\m$ from its own Galois
conjugates' images, i.e.\ $\Flr=\Fl(a_q(\m):q\in S)$. This is exactly what the certified
\texttt{residue\_degree} witness of Stage S11 (\S\ref{sec:primary}) establishes.
\end{remark}

\subsection{Separation}\label{sec:sep}
\begin{claim}[Finiteness of the anemic spectrum]\label{cl:sigma}\STD\
For fixed $\ell\ge5$ the set $\Sigma(\ell)$ of anemic eigensystems $\{a_q\}_{q\ne\ell}$ occurring in
$H^1(\Ga;\Mn)$ for \emph{some} $n$ is finite: every level-1 mod-$\ell$ eigensystem is a
$\theta$-twist of one in weight $\le\ell+1$ (Serre \cite{serre}, Jochnowitz \cite[Thm.~4.1, p.~276]{jochnowitz}, finiteness
Ash--Stevens \cite[Thm.~2.2, p.~204]{ash_stevens}; both locators verified on page images, slot-17 R17-4, v1.17.0-pre.5); equivalently $\Sigma(\ell)$ is the set of semisimple
$\bar\rho:G_{\Q,\{\ell\}}\to\GL_2(\bar\Fl)$ that are modular of level 1. Together with the
Eisenstein-type characters $\chi_a$, $a\in\Z/(\ell-1)$ (\HA{cl:chi}), this is the whole
list of characters of $\T_\infty$ on $\Htot$ for every $n$.
\end{claim}

\begin{claim}[Separating set]\label{cl:sep}\PROVED\ (given Claim~\ref{cl:sigma})
There is a finite set $S_{\mathrm{sep}}(\ell)$ of primes $q\ne\ell$, depending only on $\ell$,
such that $\{T_q\}_{q\in S_{\mathrm{sep}}(\ell)}$ separates all pairs in
$\Sigma(\ell)\cup\{\chi_a\}$ \emph{as Galois orbits}. It is computable once per $\ell$ from the
weights $n\le\ell-1$ and their $\theta$-twists (\S\ref{sec:atlas}).
\end{claim}
\begin{proof}
Finitely many distinct systems; two distinct systems differ at some $q\ne\ell$; take the union.
\end{proof}

\begin{claim}[$T_\ell$ is indispensable]\label{cl:Tl}\STD\
No $\ell$-independent set of primes $q\ne\ell$ separates the \emph{full} spectrum: for a
normalised eigenform $f$ of weight $w$ with $a_\ell(f)\neq0$, the system of $\theta^{\ell-1}f$
(weight $w+\ell^2-1$, $a_\ell=0$) coexists in that weight with $fE_{\ell-1}^{\ell+1}$
($a_\ell=a_\ell(f)$) and agrees with it at every $q\ne\ell$ \Mnote{6.1}. If $a_\ell(f)=0$ the two
full systems coincide: multiplicity $\ge2$, and in weight $\ell w$ the span of $f,\,f|V$ carries
$U_\ell=\smat0100$, a non-split self-extension. Hence $T_\ell$ (i.e.\ $U_\ell$, computed by the
same pipeline \HA{cor:ulrecursion}) is always in $S$.
\end{claim}

\subsection{Generation bounds}\label{sec:genbounds}
\begin{claim}[Sturm]\label{cl:sturm}\STD\ (a, first sentence) / \CONJ\ ($B=B(k)$ in the second) / \PROVED\ (b)
Let $k=n+2$ and $B(k):=\lfloor k/12\rfloor+1$.
\begin{enumerate}[label=(\alph*)]
\item The $\Z$-algebra $\T_\Z(S_k)$ generated by all $T_m$ on $S_k(\Z)$ is spanned as a
  $\Z$-module by $T_m$, $m\le\lfloor k/12\rfloor$ (Sturm \cite{sturm}; Stein \cite[\S9.4]{stein};
  verified on the Miller basis for $12\le k\le36$). Hence, by the pairing $(T,g)\mapsto a_1(Tg)$, which identifies $\T_\Q(M_k)$ with
  $M_k(\Q)^\vee$ (a form with $a_m=0$ for all $m\ge1$ is constant, hence $0$), $\T_\Z(M_k)\otimes\Q$ is spanned by $T_m$, $m\le B$, for
  every $B$ such that no nonzero $g\in M_k$ has $a_1(g)=\dots=a_B(g)=0$; such a $B$ exists. That $B=B(k)$ suffices is equivalent to
  $a_d(f_0)\ne0$, resp.\ (for $k\equiv2\bmod12$, where $B(k)=d+1$) to $(a_d(f_0),a_{d+1}(f_0))\ne(0,0)$, for the Miller element $f_0=1+O(q^{d})$, $d=\dim M_k$; this is
  \CONJ, verified (exact rational arithmetic on $E_4$, $E_6$, $\De$) for even $4\le k\le70$ (v1.24.0-pre.2: ``equivalent'' corrected for $k\equiv2\bmod12$ and the range
  extended, HeckeZ referee A7-m6, N8; v1.24.0-pre.1 erratum, owner ruling Q9 of 2026-11-04: the earlier argument via $\De g$ assumed $a_0(g)=0$; for
  $g=1+196560q^2+\dots\in M_{12}$, $\De g=q-24q^2+\dots$ vanishes to order $1$ only).
\item The $\Z$-module statement for $M_k$ --- $\T_\Z(M_k)$ spanned by $T_m$, $m\le B(k)$ --- is
  \emph{false}, and false at primes $\ge5$. On the Miller basis of $M_k(\Z)$ the index
  $[\T_\Z(M_k):\langle T_m:m\le B(k)\rangle_\Z]$ is divisible by $5^2$ at $k=24$, by $5\cdot7$ at
  $k=28$, by $5\cdot7\cdot131$ at $k=34$ (lower bounds $400$, $2240$, $36680$ from
  $\langle T_m:m\le3B+3\rangle$; recomputed for this revision), and the reductions of
  $T_1,\dots,T_{B(k)}$ span only $2$ of the $3$ dimensions of the image of $\T_\Z(M_k)\otimes\Fl$ in $\End M_k(\Fl)$
  at $(k,\ell)=(24,5),(28,7),(34,131)$. The obstruction is the constant term: the pairing of (a)
  sees $a_m$, $m\ge1$, only, and Sturm's theorem mod $\ell$ for $M_k$ needs $a_0$. At $(12,691)$
  the statement happens to hold (the index $3$ is prime to $691$), which is why the test case
  of \S\ref{sec:tests} does not expose it.
\end{enumerate}
\end{claim}

\begin{definition}[The computed bound $B'(k,\ell)$]\label{def:bprime}
For $B\ge2$ let $W_B:=\Z[T_q:q\le B\text{ prime}]\subseteq\T_\Z(M_k)\subseteq\End_\Z M_k(\Z)$ (the
$\Z$-subalgebra; a finitely generated $\Z$-module) and $\T^{\mathrm{sat}}:=\T_\Q(M_k)\cap\End_\Z M_k(\Z)$,
the saturation (computable as the saturation of $\langle T_m:m\le B\rangle_\Z$ for any $B$ at which this span has rank $\dim M_k$; by Claim~\ref{cl:sturm}(a) such $B$ exist, and $B=B(k)$ works in the verified range; v1.24.0-pre.1 erratum).
$B'(k,\ell)$ is the least $B\le B_{\max}$ such that the $\Fl$-algebra generated in
$\End M_k(\Fl)$ by the reductions $T_q\bmod\ell$, $q\le B$ prime, computed on the Miller basis,
has dimension $\dim M_k$. The cap $B_{\max}$ is carried by the mathematical setting
\texttt{sturm} (as \texttt{cap=<B>}; v1.5.0, the registered spelling --- through v1.4.2 this
document wrote \texttt{sturm.cap}, which no registry has)
(default $12\,B(k)$; [X] \S1.5, slot-1 owner confirmation~6), and the outcome ---
\texttt{found($B'$)} or \texttt{none\_at\_cap} --- is recorded either way together with the
resumable state of [X] \S1.3.1 R3 (\texttt{budget\_history}; \texttt{resume(cap)} refuses a lower
cap and is a no-op on a complete state). At $n=0$ ($k=2$, $M_2=0$) no bound is defined: the
record carries the outcome \texttt{not\_applicable} at record level --- a record-layer state
that E2 implements --- while the implementation's \texttt{SturmOutcome} stays three-valued
$\{\texttt{found},\texttt{none\_at\_cap},\texttt{in\_progress}\}$ (v1.4.0; slot-1 WP-3 \S7
item~7, slot-2 plan \S3.16).
\end{definition}

\begin{claim}[$B'$ certifies generation over $\Z_\ell$]\label{cl:bprime}\PROVED\
If $B'=B'(k,\ell)$ exists then $W_{B'}+\ell\,\T_\Z(M_k)=\T_\Z(M_k)$: every $T_q$ is a polynomial
with $\Z$-coefficients in the $T_{q'}$, $q'\le B'$ prime, plus $\ell$ times an element of
$\T_\Z(M_k)$. It exists iff $\ell\nmid[\T^{\mathrm{sat}}:\T_\Z(M_k)]$.
\end{claim}
\begin{proof}
$\T^{\mathrm{sat}}$ is saturated, so $\T^{\mathrm{sat}}\otimes\Fl\hookrightarrow\End M_k(\Fl)$ and
the dimension condition says $W_{B'}+\ell\T^{\mathrm{sat}}=\T^{\mathrm{sat}}$; by Nakayama
$W_{B'}\otimes\Z_\ell=\T^{\mathrm{sat}}\otimes\Z_\ell\supseteq\T_\Z\otimes\Z_\ell\supseteq W_{B'}\otimes\Z_\ell$,
so $\T_\Z/W_{B'}$ is finite of order prime to $\ell$ and $W_{B'}+\ell\T_\Z=\T_\Z$. Conversely the
dimension condition forces $\ell\nmid[\T^{\mathrm{sat}}:\T_\Z]$, and if $\ell\nmid[\T^{\mathrm{sat}}:\T_\Z]$
then, since $W_B\uparrow\T_\Z$, it holds for $B$ large. Cost: $\dim M_k\times B'$ integer
$q$-expansion coefficients and one algebra closure mod $\ell$ --- independent of the cohomology
pipeline. Computed values (this revision): $B'=2$ at $(12,691)$ and at $(34,131)$ ($\F_{131}[T_2]$
is already everything, although $\langle T_1,T_2,T_3\rangle_\Z$ is not); $B'=5$ at
$(24,5),(30,5),(32,5),(36,5)$; $B'=7$ at $(28,7)$. At $(24,5)$ and $(28,7)$ the $\Fl$-algebra
generated by $T_2,T_3$ has dimension $2<3$: the primes $\le B(k)$ do \emph{not} generate.
\textbf{Implemented} (v1.3.3; slot 1, \texttt{HeckeExt/HeckeSturmBound}, pipeline-free ---
$q$-expansions of $E_4$, $E_6$ and $\Delta=q\prod(1-q^n)^{24}$ reduced mod $\ell$ as they are
formed, so $\ell\in\{2,3\}$ work too; the cap $12\,B(k)$ as the setting \texttt{sturm}; a
certificate with \texttt{verify()}; the resumable state of the implementation plan's R3):
\texttt{HeckeSturmBoundValidation} pins every value above and, in addition, $B'=2$ on all ten
Kummer cells of Claim~\ref{cl:harder}'s list and at $(162,691)$, $(200,691)$ and
$(258,691)$ (the last recorded, $1.4$\,ms, the slowest case of the list), $B'=5$ at
$(60,5)$, $B'=3$ at $(36,7)$, and the proper subalgebra $\dim\langle T_2,T_3\rangle=2<3$ at
$(24,5)$ and $(28,7)$.
\end{proof}

\begin{claim}[What Sturm certifies mod $\ell$]\label{cl:sturmmodl}\PROVED\ (from \HA{cl:bockstein} and Claim~\ref{cl:bprime}, given a recorded $B'(k,\ell)$)
The $\Z$-algebras generated by the $T_m$ on $M_k(\Z)$ and on $H^1(\Ga;\Mn(\Z))/\mathrm{tors}$
are canonically isomorphic (both torsion-free with the same relations after $\otimes\Q$,
Eichler--Shimura). Let $q$ be a prime and write $T_q=F_q(T_{q'}:q'\le B')+\ell t'$ with
$F_q\in\Z[X_{q'}]$, $t'\in\T_\Z(M_k)$ (Claim~\ref{cl:bprime}). Then $t:=T_q-F_q(T_{q'})$ acts on
$H^1(\Ga;\Mn(\Fl))$ with image inside the image of $\mathrm{tors}\,H^1(\Ga;\Mn(\Z))$, which is
$\cong H^0(\Ga;\Mn(\Fl))$ (the Dickson part). Hence $t=0$ on every $H^1_\m$ whose character is
not that of a Dickson class, and on those $\m$ the operators $\{T_q\}_{q\le B'\text{ prime}}$
generate $\T_\m|_{H^1}$. On Dickson-type $\m$, on $X_2,X_3,X_5,X_6$, and when $B'(k,\ell)$ does
not exist, Sturm gives no certificate. In the record this is the tri-state \texttt{final} of
\S\ref{sec:schema} (v1.4.0): \texttt{sturm(B\textquotesingle)} on a non-Dickson $\m$ of $X_4$
makes the ideal \texttt{final = true}; off that range the sweep is a heuristic and the ideal is
\texttt{final = false} with reason \texttt{sturm\_heuristic} in the default configuration, or
\texttt{final = "heuristic"} under \texttt{final\_mode = sturm\_heuristic}; a missing $B'$ is
reason \texttt{sturm\_unavailable}, or \texttt{"heuristic"} under
\texttt{sturm\_unavailable = heuristic} ([X] \S1.3.1 R3.2). With $B(k)$ in place of $B'$ the conclusion fails inside
these hypotheses: at $(n,\ell)=(22,5)$ and $(26,7)$ there are no Dickson classes, and $T_2,T_3$
($B(k)=3$) generate a proper subalgebra of $\T_\Z(M_k)\otimes\Fl$ (Claim~\ref{cl:bprime}).
\end{claim}
\begin{proof}
$t$ preserves $H^1(\Ga;\Mn(\Z))$ and acts as $\ell t'$ on the torsion-free quotient, so
$t\,H^1(\Z)\subseteq\ell H^1(\Z)+\mathrm{tors}$; reduce mod $\ell$
($H^1(\Fl)=H^1(\Z)\otimes\Fl$, \HA{cl:bockstein}). $(\mathrm{tors}\otimes\Fl)_\m=0$ unless
$\m$ is Dickson-type.
\end{proof}

\noindent
(v1.15.1-pre.4; slot-15 WP-19) \emph{What the dimension condition of Definition~\ref{def:bprime} certifies.} By
Claim~\ref{cl:mkfaithful}, $\T_\Z(M_k)\otimes\Fl$ acts on $M_k(\Fl)$ with image of dimension $\dim M_k$ at every $\ell\ge5$.
So an $\Fl$-algebra generated by $T_q\bmod\ell$ that reaches dimension $\dim M_k$ \emph{is} that image. $\T_\Z(M_k)\otimes\Z_\ell$ is
saturated, and $B'(k,\ell)$ exists at every $\ell\ge5$ (Claim~\ref{cl:bprime}'s criterion
$\ell\nmid[\T^{\mathrm{sat}}:\T_\Z(M_k)]$ always holds); \texttt{none\_at\_cap} can still be recorded when the cap is too small.
The statement of this claim is unchanged.

\begin{claim}[The forms' Hecke algebra mod $\ell$ is faithful of full dimension on $M_k(\Fl)$]\label{cl:mkfaithful}
\PROVED\ (v1.15.1-pre.4, slot-15 WP-19, proof in \texttt{docs/implementation/slot15/WP-19\_forms\_proofs.md} \S1; refereed in the
WP-4 window) Let $\ell\ge5$, $k\ge4$ even, $M_k(\Fl):=M_k(\Z)\otimes\Fl$, $S_k(\Fl):=S_k(\Z)\otimes\Fl$, and let $\T(X)$ be the
$\Fl$-algebra generated by all $T_m$ in $\End(X)$.
\begin{enumerate}[label=(\alph*)]
\item $\dim\T(M_k(\Fl))=\dim M_k$ and $M_k(\Fl)\cong\Hom(\T(M_k(\Fl)),\Fl)$ as $\T$-modules; hence $\dim\T_\m=\dim M_k(\Fl)_\m$ for every
  maximal ideal $\m$, and $\T_\Z(M_k)\otimes\Fl\to\End M_k(\Fl)$ is injective. The same holds for $S_k$ at every prime $\ell$.
\item If $(\ell-1)\mid k$, then $M_k(\Fl)=\Fl\cdot1\oplus S_k(\Fl)$ as $\T$-modules, and the system $m\mapsto\sigma_{k-1}(m)\bmod\ell$ of the
  constant form $1=E_{\ell-1}^{k/(\ell-1)}$ occurs in no cusp form mod~$\ell$; so $\T(M_k(\Fl))=\Fl\times\T(S_k(\Fl))$.
\end{enumerate}
\end{claim}
\begin{proof}
The pairing $(t,f)\mapsto a_1(tf)$ is $\T$-balanced, and $a_1(T_mf)=a_m(f)$. Its kernel on the form side consists of the constants in
$X$, and its kernel on the algebra side is $0$ when $X$ has no non-zero constant. A non-zero constant lies in $M_k(\Fl)$ iff
$(\ell-1)\mid k$: if $f\equiv c\neq0$, then $fE_4\equiv cE_4$, and Serre's weight congruence \cite{serre} (LNM~350, n$^\circ$~1.3, Th\'eor\`eme~1, p.~198, verified (slot-16 WP-8): for $p\ge3$ and $f,f'$ with rational coefficients of weights $k,k'$, $f\ne0$ and $v_p(f-f')\ge v_p(f)+m$ force $k\equiv k'\pmod{(p-1)p^{m-1}}$) with $m=1$
gives $(\ell-1)\mid k$. Conversely $E_{\ell-1}\equiv1$ (von Staudt--Clausen). So the pairing is perfect for $S_k$, and for $M_k$ when
$(\ell-1)\nmid k$, and (a) follows by localising the pairing.

For (b), $1$ is a $T_m$-eigenvector with eigenvalue $\sigma_{k-1}(m)$ (read on $a_0$), and $f-a_0(f)\cdot1\in S_k(\Fl)$ (Miller basis). Suppose a
cusp form mod~$\ell$ carries that system. Then it has $a_m=\sigma_{k-1}(m)$ for $m\ge1$. Lift it to $\tilde F\in S_k(\Z)$ and put
$c:=-B_k/2k$, which has $v_\ell(c)=-1-v_\ell(k)$ by von Staudt--Clausen. Then $h:=(G_k-\tilde F)/c\in M_k(\Q)$ satisfies
$h\equiv1\pmod{\ell^{2+v_\ell(k)}}$, and $hE_4\equiv E_4$ gives, by the same theorem, $(\ell-1)\ell^{1+v_\ell(k)}\mid k$, which is
absurd. So the idempotent of the constant's system lies in $\T(M_k(\Fl))$ and projects onto $\Fl\cdot1$, and (a) follows in this
case as well.

$\ell\ge5$ enters through Serre's theorem (its modulus for $p\ge5$) and through $E_{\ell-1}$ being a modular form of level one.
\end{proof}

\begin{claim}[F9: the Eichler--Shimura count on a block, with its torsion term]\label{cl:formsF9}
\PROVED\ (v1.15.1-pre.4, slot-15 WP-19, proof in \texttt{WP-19\_forms\_proofs.md} \S2; refereed in the WP-4 window) Let $\ell\ge5$ and
$n>0$ even, $k=n+2$. Let $P'$ be a finite set of primes, where $T_\ell$ is the level-one operator (the pipeline's $U_\ell$). Let
$\lambda\in\Fl^{P'}$, and let $X_\lambda$ be the joint generalised $\lambda$-eigenspace. Then
\[
  \dim H^1(\Ga;\Mn)_\lambda=\dim S_k(\Fl)_\lambda+\dim M_k(\Fl)_\lambda+\dim H^0(\Ga;\Mn)_\lambda .
\]
By \HA{prop:heckeH0}(1),(2), $\dim H^0(\Ga;\Mn)_\lambda$ is $0$ if $\ell\in P'$ and $\lambda_\ell\neq0$. Otherwise it is the number of
$(a,b)$ with $a(\ell+1)+b(\ell^2-\ell)=n$ and $(1+p)p^a\equiv\lambda_p$ for all $p\in P'\setminus\{\ell\}$. If $\{T_p\}_{p\in P}$,
$P\supseteq P'$, generates $\T(S_k(\Fl))$, resp.\ $\T(M_k(\Fl))$, the first two terms equal the dimensions of the algebras
$\Fl[T_p|_{X_\lambda}:p\in P]$ (Claim~\ref{cl:mkfaithful}). \texttt{HeckeFormsHeckeAlgebraValidation}'s F9 identity
$D=\dim\T(S_k)_\m+\dim\T(M_k)_\m$ is therefore exact on every block with $\dim H^0(\Ga;\Mn)_\lambda=0$ (the torsion guard), and is off
by exactly that term otherwise.
\end{claim}
\begin{proof}
$H^1(\Ga;\Mn)=\mathbb L/\ell$ is an extension of $\mathbb L'/\ell$ by $\mathbb L_{\mathrm{tors}}/\ell$ \HA{prop:cohforms}(1). Localising at $\lambda$
is exact. $\dim(\mathbb L_{\mathrm{tors}}/\ell)_\lambda=\dim\mathbb L[\ell]_\lambda=\dim H^0(\Ga;\Mn)_\lambda$, by counting cyclic summands and by
the Hecke-equivariant Bockstein \HA{cl:bockstein}(2),(3). For a lattice, $\dim(L/\ell)_\lambda$ counts the $\bar\Q_\ell$-systems reducing to
$\lambda$. By \HA{lem:heckerankone}'s proof, $\mathbb L'\otimes\bar\Q_\ell\cong M_k\oplus\overline{S_k}$, semisimply, so the free term is
$2s+e=s+(s+e)$, where $s$ counts the cuspidal and $e$ the Eisenstein systems over $\lambda$; and $S_k(\Z_\ell)$, $M_k(\Z_\ell)$ give $s$ and
$s+e$. On $H^0$, $T_p$ is diagonal on Dickson monomials and $U_\ell$ is nilpotent \HA{prop:heckeH0}.
\end{proof}


\begin{claim}[Bella\"iche--Khare]\label{cl:bk}\STD\ (recalled; to be checked against the paper before use as a default)
For $\ell\ge5$ and an irreducible unobstructed modular $\bar\rho$ of level 1, the completed
cuspidal anemic Hecke algebra $\T_{\bar\rho}$ is isomorphic to $\Fl[[x,y]]$, with $x,y$
realisable as $T_q-a_q$, $T_{q'}-a_{q'}$ for suitable primes $q,q'$ (Bella\"iche--Khare
\cite{bk}; Deo \cite{deo} for $U_\ell$ and higher level; Bella\"iche \cite{eigenbook}). Every
fixed-weight $\T_{n,\m}$ is a quotient, so two well-chosen operators away from $\ell$ plus
$T_\ell$ generate $\T_{n,\m}$ for all $n$ simultaneously; which pair works is a tangent-space
condition checkable in any one weight. For $n\gg\ell$ no single $T_q$ suffices mod $\ell$
(finitely many values in $\Flr$ against $\sim n/6$ eigensystems); for small $n$ one operator
can ($T_2$ alone generates at $(10,691)$, Claim~\ref{cl:ramanujan}). The campaign row
$(n,\ell)=(160,691)$ is a \emph{candidate} example (\S\ref{sec:tests}): $T_5$ alone
has a $4{:}1$ block where $T_2$ and $T_3$ have $2{:}1,2{:}1,2{:}1,4{:}2$; whether that $4{:}1$
is two $\Fl$-rational systems merged at $a_5$ or is $T_2$'s $4{:}2$ block with an
$\Fl$-rational $a_5$ was not determined by the single-operator rows of
\texttt{results\_table.tsv}. \textbf{Verified computationally on 2026-09-16} (v1.3.3; the
slot-1 probe \texttt{HeckeLocalityProbe}, one \texttt{ComputeMultiOp} run on
$\{T_2,T_3,T_5\}$ with the locality certificate on every block): the simultaneous
decomposition has seven blocks
$2{:}1,\,1{:}1,\,2{:}1,\,2{:}1,\,4{:}2,\,6{:}3,\,10{:}5$ (dimension${:}$residue degree, in
that order), every one local, \texttt{certify() = 0}; the $4{:}2$ block is a genuine
residue-degree-$2$ system. \textbf{Which alternative holds} (v1.7.0; measured by
\texttt{HeckeExtSeparationValidation}, its \texttt{SEPARATION\_SINGLE} line, slot-5 WP-8): the
\emph{second}. Run with $\{T_2\}$, $\{T_3\}$, $\{T_5\}$ and $\{T_2,T_3\}$ alone, every family
gives seven blocks with the same dimension multiset $\{1,2,2,2,4,6,10\}$; $T_2$ alone, $T_3$
alone and $\{T_2,T_3\}$ each give the family's list exactly, and $T_5$ alone gives
$4{:}1,2{:}1,2{:}1,2{:}1,1{:}1,6{:}3,10{:}5$ --- its $4{:}1$ is the family's $4{:}2$ block,
with $a_5$ rational on it, seen with a coarser invariant. It is not two
$\F_{691}$-rational systems merged at $a_5$ and split by the refinement, as v1.3.3--v1.6.2
said: that would give eight blocks, and every family gives seven. (v1.7.1: the validation
compares the $(\dim{:}f)$ lists and the block counts, not subspaces; that the partitions
\emph{coincide} follows from the counts, since the family's decomposition refines each single
operator's and every run has seven blocks.) So the row is evidence for
the refinement logic, but it is \emph{not} an instance of this claim's ``no single $T_q$
suffices'': \textbf{the block partition at $(160,691)$, with its residue degrees, is already
$T_2$'s (and $T_3$'s); $T_5$ alone gives the same partition and under-reports one residue
degree.} (Through v1.6.2 this sentence read ``as
this claim predicts for $n\gg\ell$, no single $T_q$ separates everything while $\{T_2,T_3\}$
does'', which the measurement refutes; v1.7.0's ``no single $T_q$ certifies the residue degrees
that $\{T_2,T_3\}$ does'' is refuted by the same line, $T_2$ alone giving the family's residue
degrees; the general statement for $n\gg\ell$ above is a
counting argument and is not affected.) The row does not by itself verify the isomorphism
$\T_{\bar\rho}\cong\Fl[[x,y]]$, which stays \STD{} (recalled).
\end{claim}

\begin{claim}[$T_\ell$ nilpotent on $\theta$-images]\label{cl:Tlnil}\textsf{[(1): proved in [C]; (2): standard, reference]} (v1.12.0-pre.1)
\begin{enumerate}
\item $U_\ell$ kills $\theta$-images over $\Fl$: on $H^1$ the image of
  $L\cdot:H^1(\Ga;M_{n-\ell-1})\to H^1(\Ga;\Mn)$ (multiplication by the Dickson invariant,
  $\SL_2(\Fl)$-equivariant, Hecke-equivariant up to a $\det$-twist \Mnote{2}) is killed by
  $U_\ell=T_\ell$, i.e.\ $A(U_\ell)_n\circ\theta=0$, for every $n\ge\ell+1$ and every prime
  $\ell$ --- \emph{already on cochains}. \emph{Pointer} (v1.12.0-pre.1): this sentence is
  \HA{prop:thetatwist} (in \HA{sec:dicksontheta}, [C] v1.11.0-pre.1, slot-11 WP-2), proved there in
  the code's conventions (\HA{def:rhon}, \HA{H2}):
  every coset representative of $\Ga\,\mathrm{diag}(1,\ell)\,\Ga$ has determinant $\ell$, so the
  Hecke operator $U_\ell$ on $H^\bullet(\Ga;M_{n-\ell-1}\otimes\det)$ is $\ell\cdot A(U_\ell)_{n-\ell-1}=0$
  over $\Fl$, and $\theta$ intertwines the twisted operator with $A(U_\ell)_n$. This document keeps
  no proof of its own; where the two differ, [C] wins. The computational check is (6e) of
  \S\ref{sec:identities} (slot-11 WP-6); that $\theta$-images are non-ordinary is
  \HA{cor:thetanonordinary}. \emph{Not} claimed: that $U_\ell$ vanishes on all of $H^0$ ---
  it does not in general (\HA{prop:heckeH0}(4); (6f) of \S\ref{sec:identities}). \emph{Through
  v1.11.1} the sentence read ``$U_\ell$ kills
  $\theta$-images; \ldots\ is killed by $T_\ell$ after the twist'' and carried \STD.
\item For $n\ge\ell+1$, $T_\ell$ has large kernel and typically nilpotent Jordan blocks on
  $H_\m$; its Jordan type is a basic extension invariant \Mnote{6.3}. \STD{} (unchanged).
\end{enumerate}
\end{claim}

%======================================================================
\section{Primary decomposition and the local algebras \texorpdfstring{$\T_\m$}{T\_m}}\label{sec:primary}
%======================================================================

\subsection{What exists: Stages S10--S11}
\texttt{HeckeModuleDecomposition} (header comment is the binding definition) computes, for a
degree $q$ and a commuting family $\{A^{(q)}(T_{p_j})\}_j$ over $\Fl$:
\begin{enumerate}[label=(\roman*)]
\item the common refinement of the single-operator primary decompositions: blocks
  $V=\{v:\pi_j(T_{p_j})^{N}v=0\ \forall j\}$ with $\pi_j\in\Fl[X]$ irreducible;
\item per block and operator, the minimal polynomial $\pi_j^{e_j}$, the characteristic
  polynomial, and the generalized-Jordan partition of $\dim V/\deg\pi_j$ from
  $\rank\pi_j(T_{p_j})^r$;
\item the residue degree $f:=\operatorname{lcm}_j\deg\pi_j$, certified as
  $[\T_S/\m:\Fl]$ by a \emph{witness} $w=\sum_jc_jT_{p_j}$ whose minimal polynomial on the block
  has squarefree part a single irreducible of degree exactly $f$ (deterministic-seed search;
  reported as \texttt{residue\_degree\_lower\_bound} if no witness is found);
\item an adapted basis $U$, and \texttt{certify()}: all of the above re-derived
  deterministically (Berkowitz charpoly, Frobenius/gcd irreducibility proof, product identity,
  block-diagonality of every $U^{-1}A_jU$, nilpotency exponents, partitions from ranks).
\end{enumerate}
Factorisation is Cantor--Zassenhaus with a fixed seed, always followed by
\texttt{CertifyFactorization}.

\subsection{A gap: per-operator refinement is not \texorpdfstring{$\MaxSpec\T_S$}{MaxSpec T\_S}}\label{sec:gap}
\begin{claim}\label{cl:gap}\PROVED\
The common refinement (i) can leave a block on which $\T_S$ is \emph{not} local. It is local iff
the residue-degree witness (iii) exists with $\deg\pi_w=f$ \emph{and} $\dim_{\Fl}(\T_S|_V/\rad)=f$.
\end{claim}
\begin{proof}
Example: $\ell$ arbitrary, $\alpha\in\F_{\ell^2}\setminus\Fl$, $V=\F_{\ell^2}\oplus\F_{\ell^2}$,
$t_1=(\alpha,\alpha^\ell)$, $t_2=(\alpha,\alpha)$ acting diagonally. Each $t_j$ has minimal
polynomial the irreducible $\pi=\min(\alpha)$, so (i) does not split $V$, yet
$\Fl[t_1,t_2]\ni t_1-t_2=(0,\alpha^\ell-\alpha)$ is a non-nilpotent zero-divisor: two maximal
ideals, characters $(\alpha,\alpha)$ and $(\alpha^\ell,\alpha)$ --- not Galois conjugate. In
general $\T_S|_V/\rad$ is a product of fields each containing a root of every $\pi_j$; it is a
single field iff its $\Fl$-dimension is $f$; equivalently, iff some (then a generic) element
$w$ has irreducible squarefree minimal polynomial $\pi_w$ of degree $f$ \emph{and} every $t_j$ is
congruent modulo nilpotents to a polynomial in $w$ --- the two clauses of
Certificate~\ref{sec:localcert}; the first alone is not sufficient (in the example
$t_2=(\alpha,\alpha)$ has irreducible squarefree minimal polynomial of degree $f=2$). Until
slot 2 (E0b) the code accepted a single
operator with $\deg\pi_j=f$ as the witness (\texttt{HeckeModuleDecomposition.cpp:687--692}), so in
the example it reported $V$ as one block with \texttt{residue\_degree\_certified=true}: the witness
search certified a \emph{primitive element of $\T_S/\rad$ assuming locality}, not locality itself
(independently found in \texttt{hecke\_extensions\_capability\_inventory.md} \S2.3.4). Since E0b
(slot-2 WP-1, 2026-09-16) the refinement loop enforces both clauses of
Certificate~\ref{sec:localcert} and the example splits; the $48$ constructed non-local families
of \texttt{HeckeLocalityValidation}, which a frozen copy of the old predicate accepts, are all
split and certified, and no block of the golden grid moved (every one was local, slot 1).
\end{proof}

\begin{decisionbox}
A block is a maximal-ideal component of $\T_S$ only when certified local
(Certificate~\ref{sec:localcert}); Stage~2 of \S\ref{sec:alg} refines by the primary
decomposition of the witnesses $w=\sum c_jT_{p_j}$ themselves until that certificate passes,
within a bounded number of attempts (the budget mode, the default); the strict mode of
Certificate~\ref{sec:localcert} continues to termination.
\end{decisionbox}

\subsection{The local algebra and its invariants}\label{sec:localalg}
Fix a certified block $V=H_\m$ (over $\Fl$; $\dim_{\Fl}V=d_\m$) and generators
$t_j:=T_{p_j}|_V$, $j\in S$. Let $A_S:=\Fl[t_j:j\in S]\subseteq\End_{\Fl}(V)$ (computed as an
$\Fl$-subspace of $\End(V)$, \S\ref{sec:alg} Stage 3). Its invariants:
\begin{itemize}
\item $\dim_{\Fl}A_S$; the radical $\rad A_S=\{t\in A_S:t\text{ nilpotent}\}=\m_S$ (the unique
  maximal ideal once local), $A_S/\m_S\cong\Flr$;
\item the \emph{embedding dimension} $\dim_{\Flr}\m_S/\m_S^2$ (Claim~\ref{cl:gen}(b));
\item the nilpotency index $\nu$ ($\m_S^\nu=0\ne\m_S^{\nu-1}$), i.e.\ the Loewy length of
  $A_S$; the socle $\soc A_S=\Ann_{A_S}(\m_S)$ and whether $A_S$ is Gorenstein
  ($\dim_{\Flr}\soc A_S=1$);
\item the Hilbert function $h(i)=\dim_{\Flr}\m_S^i/\m_S^{i+1}$.
\end{itemize}
All are $\Fl$-linear algebra on $A_S$ (\S\ref{sec:field}: no field extension is needed; the
$\Flr$-dimensions are the $\Fl$-dimensions divided by $r$).

\begin{remark}[The coefficient field]\label{rmk:coef}
A local artinian $\Fl$-algebra $A$ with residue field $\F=\Flr$ contains a unique subfield $K$
mapping isomorphically onto $\F$ (Hensel; $\F/\Fl$ is separable): $K=\{x\in A:x^{\ell^r}=x\}$,
the Teichm\"uller representatives. Every $x\in A$ is uniquely $x=s+n$ with $s\in K$, $n$
nilpotent, and $s=x^{\ell^{rN}}$ once $\ell^{rN}$ exceeds the nilpotency index.
Claims~\ref{cl:gen}(b), \ref{cl:yoneda} and \ref{cl:split} read ``$T_q-a_q(\m)$'' and
``$\F$-vector space'' through $K$; Stage~3 constructs it explicitly.
\end{remark}

\begin{claim}[Residue field acts through $A_S$]\label{cl:resfield}\PROVED\
Let $w\in A_S$ be a residue witness ($\pi_w$ irreducible of degree $r$, $\min_V(w)=\pi_w^{e}$)
and put $\hat w:=w^{\ell^{N}}$ with $N$ a multiple of $r$ and $\ell^N\ge D$. Then
$\pi_w(\hat w)=\pi_w(w)^{\ell^N}=0$, $\hat w\equiv w\pmod{\m_S}$, and $\Fl[\hat w]=K\cong\Fl[X]/(\pi_w)$
is the coefficient field of $A_S$. For every $j$ there is a unique $P_j\in\Fl[X]_{<r}$ with
$n_j:=t_j-P_j(\hat w)$ nilpotent, namely the one with $P_j(\hat w)$ the $K$-component of $t_j$
(for $j=w$: $n_w=w-\hat w$); the tuple $(P_j(\bar X))_j\in\Flr^{S}$ is the character $a(\m)$ in
the coordinates $\Flr=\Fl[X]/(\pi_w)$. In general $P_j(w)\notin K$ and $w$ itself must not be
used: for $A_S=\F[\varepsilon]/(\varepsilon^2)$, $t=\alpha+\varepsilon$, $w=t$, one has
$t-P(w)=0$ with the \emph{unreduced} $P=X$ although the nilpotent part of $t$ is
$\varepsilon\neq0$; with $\hat w=\alpha$ one gets $n_t=\varepsilon$. (The character is always
the reduced representative $P_j=X\bmod\pi_w\in\Fl[X]_{<r}$: here $r=1$ and $P_t=\alpha$, so
$P_t(w)=P_t(\hat w)=\alpha$ and the two recipes agree --- the substitution of $w$ for $\hat w$
is wrong only at $r\ge2$, where $P_j(w)\notin K$ in general. Both facts are pinned by
\texttt{HeckeLocalityValidation}, slot 1: the trap arm at $r\ge2$, and the $r=1$ twin arm on
which the analytic and the matrix path agree verdict field for verdict field.)
\end{claim}
\begin{proof}
$A_S/\m_S$ is a field of degree $r$ containing $\Fl[\bar w]$ of degree $r$, so $\bar w$ generates
it; $\hat w$ has the same residue ($r\mid N$) and satisfies $\pi_w(\hat w)=0$ because Frobenius
is a ring endomorphism of $A_S$ fixing $\Fl$, so $\Fl[\hat w]$ is a field of degree $r$, i.e.\
$K$. Uniqueness of $P_j$: $K\to\F$ is an isomorphism and $\bar t_j=\overline{P_j(\hat w)}$.
\end{proof}

%======================================================================
\section{What ``all extension classes'' means}\label{sec:ext}
%======================================================================

\subsection{Simple modules and \texorpdfstring{$\Ext^1$ over $\T_\m$}{Ext1 over T\_m}}
Over the local artinian ring $\T_\m$ with residue field $\F:=\Flr$ there is exactly one simple
module, $\F$ itself ($\dim_{\Fl}\F=r$). Every finite $\T_\m$-module $X$ has a composition
series with all factors $\F$; its length is $\dim_{\Fl}X/r$. Over $\Fl$ (before scalar
extension) $H_\m$ is $\T_\m$-primary with all Jordan--H\"older factors isomorphic to $\F$.

\begin{claim}[$\Ext^1$ between simples]\label{cl:ext1}\PROVED\
$\Ext^1_{\T_\m}(\F,\F)\cong\Hom_{\F}(\m/\m^2,\F)$, of $\F$-dimension the embedding dimension.
An extension $0\to\F\to E\to\F\to0$ is classified by the $\F$-linear map
$\m/\m^2\to\Hom_\F(\F,\F)=\F$, $t\mapsto(\bar b\mapsto t\,b)$.
\end{claim}
\begin{proof}
Apply $\Hom_{\T_\m}(-,\F)$ to $0\to\m\to\T_\m\to\F\to0$: $\Hom(\T_\m,\F)\to\Hom(\m,\F)\to
\Ext^1(\F,\F)\to0$, and $\Hom_{\T_\m}(\m,\F)=\Hom_\F(\m/\m^2,\F)$; the first map is zero
because $\m\cdot\F=0$. The description of the class is the standard one for a two-step
module: $E/\m E$ and $\m E$ are $\F$-spaces and $\m$ acts through $\m/\m^2$.
\end{proof}

\begin{claim}[Yoneda versus module-theoretic data]\label{cl:yoneda}\PROVED\
For a two-step $\T_\m$-module $E$ ($\m^2E=0$) with $A:=\m E\subseteq E$, $B:=E/\m E$ (both
$\F$-spaces through the coefficient field, Remark~\ref{rmk:coef}), the Yoneda class of $0\to A\to E\to B\to0$ in $\Ext^1_{\T_\m}(B,A)=\Hom_\F(\m/\m^2,\Hom_\F(B,A))$ is
the \emph{structure map} $\theta_E:\m/\m^2\to\Hom_\F(B,A)$, $t\mapsto(b\mapsto tb)$; it is a
complete invariant of the isomorphism class of $E$ up to $\GL(A)\times\GL(B)$, and $E$ is
determined by $(\dim B,\dim A,\theta_E)$. The image of $\theta_E$ is the space of \emph{active
directions}; $\rank\theta_E=0$ iff $E$ is semisimple.
\end{claim}
\begin{proof}
$E$ is the $\T_\m$-module $B\oplus A$ with $t\cdot(b,a)=(0,\theta_E(t)b)$; isomorphisms of such
modules that are the identity on the filtration are exactly pairs $(g_B,g_A)$ conjugating
$\theta_E$; the Yoneda class is by definition the obstruction to splitting the same data.
\end{proof}

\subsection{The classification data (binding)}\label{sec:classdata}
For a longer Loewy length the isomorphism type of a $\T_\m$-module is \emph{not} captured by a
finite list of numbers: for embedding dimension $\ge2$ the module category of
$\F[x,y]/(x,y)^\nu$ is of wild representation type for $\nu\ge3$ (Drozd \cite{drozd}; Ringel
\cite{ringel}; $\nu=2$ is the tame Kronecker case) \STD. ``Classifying all extension classes'' is therefore taken to mean the following
\emph{certified record} per $(n,\ell,\m)$, which determines the modules up to isomorphism by
exhibiting them, together with the invariants that make two records comparable:

\begin{enumerate}[label=(D\arabic*)]
\item\label{D1} $\T_\m$ (in the form of the matrices of the certified generators
  $t_j$ on $H_\m$, an $\Fl$-basis of $A_S$, and the residue witness), its invariants of
  \S\ref{sec:localalg}, and the certificate that $A_S=\T_\m$ (Stage~4, \S\ref{sec:alg}).
\item\label{D2} For each of the six $X_i$: the Loewy (radical) series
  $X_{i,\m}\supseteq\m X_{i,\m}\supseteq\m^2X_{i,\m}\supseteq\cdots$ and the socle series
  $0\subseteq\soc X\subseteq\soc^2X\subseteq\cdots$ ($\soc^{j+1}X/\soc^jX=\soc(X/\soc^jX)$),
  with layer dimensions in units of $r$; $X$ is semisimple iff $\m X=0$.
\item\label{D3} For each generator $t_j$ and each $X_i$: the Jordan type of $t_j-a_{p_j}(\m)$
  over $\F$, i.e.\ the partition read from the nullities of $\pi_j(t_j)^s$ on $X_{i,\m}$ divided
  by $f$ (Claim~\ref{cl:split}: the S10/S11 generalized-Jordan partition with every multiplicity
  divided by $f/\deg\pi_j$), and, when
  $\ell$-adic lifts are relevant, the same for $T_\ell$ specifically (Claim~\ref{cl:Tlnil}).
\item\label{D4} For every consecutive two-step piece $\m^iX/\m^{i+2}X$ the structure map
  $\theta$ from $\m/\m^2$ to $\Hom_\F(\m^iX/\m^{i+1}X,\ \m^{i+1}X/\m^{i+2}X)$ ($\F$-linear, $\F$
  acting through the coefficient field), its rank and its matrices; the subset of $S$ whose
  $n_j=t_j-P_j(\hat w)$ (Claim~\ref{cl:resfield}) span the active directions \emph{over $\F$} is
  the minimal set needed to \emph{see} the extensions at $\m$.
\item\label{D5} Cyclicity and freeness: is $X_{i,\m}$ cyclic over $\T_\m$ (a quotient of
  $\T_\m$), is it free of rank one (``multiplicity one''), and the minimal number of generators
  $\dim_\F X/\m X$.
\item\label{D6} The position of the canonical submodules relative to the Loewy layers: $H^1_!$,
  the \emph{Eisenstein line} at $\m$ --- defined as $\soc(X^+_{4,\m})$ when that socle is
  one-dimensional over $\F$ and undefined otherwise (the reductions of the integral Eisenstein
  and cuspidal eigenlines are not mod-$\ell$ Hecke data: they reduce into the same joint
  eigenspace $E=\soc X_{4,\m}$ of the $T_q$, \S\ref{sec:tests}. \emph{Sharpened, v1.4.0}:
  through v1.3.3 this read ``at $(10,691)$ both reduce to the same line''; measured on the
  unsplit $H^1$, $E$ is \emph{two-dimensional} at $(10,691)$ and at every one of the ten Kummer
  cells --- the socle of the non-split block, $\varepsilon\T_\m\oplus\F$ in the notation of
  Claim~\ref{cl:ramanujan}, one line per $\iota$-sign expected --- and a line on every
  off-Dickson non-Kummer cell with $k\le26$ ($43$ cells)
  (\texttt{HeckeBoundaryParabolicValidation}, slot-2 WP-7, 2026-09-16); so without the sign
  split no ``Eisenstein line'' exists at a Kummer cell, and $\soc(X^+_{4,\m})$ is a line there
  only after E3 (\VERIFY{} per sign). Only a mod-$\ell^2$ run separates the two reductions: the
  $\Z/\ell^2$-based definition of E8 is this datum's \emph{refinement}, never a replacement, and
  the mod-$\ell$ socle semantics stays the published one until an E8 run exists --- slot-1 owner
  confirmation~9(a); v1.9.0: an E8 run exists since slot~8, but only in the validations and the
  golden harness, never in a record, so the published semantics is unchanged; what the run
  measured at $(10,691)$ is recorded in \S\ref{sec:ring}~(d)) --- the Dickson
  classes ($X_1$ and $r^0X_1\subseteq X_2$),
  $L\cdot H^1(M_{n-\ell-1})$ when $n\ge\ell+1$, and the images/kernels of the maps of
  \HA{eq:les}; for each: a Loewy layer, a sum of layers, or skew.
\item\label{D7} The list of \emph{predicted} extensions (\S\ref{sec:expect}) with a
  confirmed/refuted flag, and the list of unpredicted non-split pieces.
\end{enumerate}
``$X_i$ carries a non-trivial Hecke module extension at $\m$'' $:\Leftrightarrow\m X_{i,\m}\ne0$.

\subsection{What to expect}\label{sec:expect}
\begin{prediction}[Eisenstein--cuspidal]\label{pr:eiscusp}\STD\ (Ribet \cite{ribet}, Mazur \cite{mazur}; Kummer congruences)
\emph{Reworded at v1.12.0.} Let $k=n+2\not\equiv0\pmod{\ell-1}$ and let $\chi_0$ be the system of $E_k$'s
own class: the anemic eigenvalues $1+p^{k-1}$ ($p\ne\ell$) \emph{together with} the
$U_\ell$-eigenvalue~$1$. Then $\chi_0\equiv$ a cuspidal system at $\m$ iff $\ell\mid\num(B_k/k)$, i.e.\
(Kummer) iff $\ell\mid\num(B_w/w)$ for the $w\equiv k\pmod{\ell-1}$ with $2<w\le\ell-1$; the ``usual
exceptional residues'' are $w\equiv0,2\pmod{\ell-1}$, where no such $w$ exists (at $w\equiv2$ the answer
is ``never'', $B_2/2$ being an $\ell$-unit at $\ell\ge5$). Then $r^1$ vanishes on the relevant class and
$X_{4,\m}$ is non-split. For $n<\ell-1$ this is Claim~\ref{cl:harder}; the Galois side is the class of
Ribet's lemma in $H^1(\Q_{\{\ell\}},\omega^{-(n+1)})$. \emph{What the prediction does not say.} Through
v1.12.0-pre.2 it read ``$\chi_a\equiv$ a cuspidal system at $\m$ iff $\ell\mid\num(B_w/w)$ for some
$w\equiv n+2-2a\pmod{\ell-1}$, $2<w\le\ell-1$ (up to the usual exceptional residues)'', for every
twist~$a$ and on the anemic system. Read literally that is false: at $(10,7)$, $a=1$,
$\chi_1(p)=p+p^{10}\equiv p+p^{4}\equiv\tau(p)\bmod7$ (Ramanujan's $\tau(n)\equiv n\sigma_3(n)\bmod7$), a
cuspidal system, while $w\equiv4$ and $B_4/4=-1/120$ is a $7$-unit; for $a\ge1$ the $\theta$-twist of an
Eisenstein system is a mod-$\ell$ cusp form at every $\ell$ (\HA{rmk:thetaclassical}), so no Bernoulli
condition governs it. At $a=0$ the anemic system alone does not suffice either: the congruence census
(\S\ref{sec:census}, edition~1) has $58$ non-reduced \texttt{mixed} rows at the regular primes
$7,11,13$ carrying $E_k$'s anemic system with $U_\ell$-eigenvalue~$0$ (the system of $\theta^{\ell-1}E_k$ modulo $\ell$, whose $a_\ell$
is~$0$), and $592$ \texttt{mixed} rows at those primes in all. (Found by slot-11 WP-19, confirmed by
its referee WP-20.) The \STD{} tag is for the statement as reworded.
\emph{Which ``cuspidal'' (v1.12.1, the slot-11 WP-15 referee's m-6).} ``A cuspidal system'' above means
the system of a characteristic-zero cusp eigenform (as the $(198,691)$ item of \S\ref{sec:tests} and the
examples document read it); under that reading both directions are the classical ones (Ribet, Mazur).
If it is read as ``the system of a mod-$\ell$ cusp form'', the direction ($\Rightarrow$) still holds, but
the ``iff'' is then the Serre / Swinnerton-Dyer constant-term criterion ($E_k\bmod\ell$ with vanishing
constant term is a mod-$\ell$ cusp form of weight $k$, and lifts by Deligne--Serre with $a_\ell\equiv1$),
not the Ribet--Mazur statement. The truth of the claim does not depend on the reading.
\end{prediction}

\begin{prediction}[Cuspidal--cuspidal]\label{pr:cuspcusp}\CONJ\ (expected on the evidence of \Mnote{7}; to be confirmed case by case)
Congruent forms in one weight arise from $E_{\ell-1}$-multiplication, from $\theta^{\ell-1}$
(Claim~\ref{cl:Tl}), and accidentally. Multiplicity $\ge2$ forces the split/non-split
question; $T_\ell$ is usually the active direction, but anemic nilpotents occur (large
$\T_\m$).
\end{prediction}

\begin{prediction}[Boundary and cone]\label{pr:cone}\CONJ\
When $(\ell-1)\mid n$ or more generally $\chi_a\equiv\chi_{a'}$, the cone can glue
$H^0(\partial)$- and $H^1(\partial)$-type classes; when a Dickson class has the character of a
cusp form or of another boundary class, $X_3$ and $X_4$ can glue them. These are the extension
classes that are ``in the cone'' rather than in $H^1$; exactness of \HA{eq:les} determines
$X_{3,\m}$ only up to extension, the cone determines it. \emph{Which} cone module is meant is
therefore part of the prediction, and it is fixed by \ref{add:cone}: the $\T_\m$-module structure on
$X_{3,\m}$ is the one induced by $\Phi_g(a,b)=(\Theta_g a,\Theta^\partial_g b+ha)$ with $h$ derived
from the certified pair homotopy $\eta$. It is \emph{not} determined by exactness plus a freely
chosen cochain homotopy: the two would differ by an arbitrary element of $\Hom(\ker r^1,\coker r^0)$,
which is non-zero in exactly the situations this prediction is about --- at $(10,691)$ it has
dimension $3$ while $\dim X_3=4$ (\ref{add:cone}, third bullet). The gluing is therefore a statement about
a well-defined module, and the computation reports it as such.
\end{prediction}

\begin{prediction}[Duality]\label{pr:dual}\STD\ (\HA{cl:pd})
The Loewy layers of $H^{2-i}(\Ga;\Dn)_\m$ are the socle layers of $X_{i,\m}=H^i_c(\Ga;\Mn)_\m$
in reverse order, and vice versa ($\rad^i(V^\vee)=(\soc^iV)^\perp$; both series are in
\ref{D2}, so nothing new is computed --- but ``reversed Loewy series'' would fail on any
non-rigid module); a module cyclic for $\Mn$ need not be cyclic for $\Dn$. Computed in E7
(slots 6--7 of [X] v1.5.1, scheduled; asserted at $(10,691)$ and on the Kummer cells by
\texttt{HeckeDualityValidation}, v1.4.0).
\end{prediction}

\begin{prediction}[Galois interpretation]\label{pr:galois}\CONJ\ (Bella\"iche--Chenevier \cite{bc})
Non-split $\T_\m$-extensions between a cuspidal and an Eisenstein-type factor correspond to
non-zero classes in $H^1(\Q_{\{\ell\}},\bar\rho_1\otimes\bar\rho_2^\vee)$; the computed
active-direction ranks are to be compared with those Galois cohomology dimensions.
\end{prediction}

\begin{prediction}[Concentration for $n\gg\ell$]\label{pr:conc}\STD\ (from Claim~\ref{cl:sigma})
$\dim H^1\sim n/6$ while $\#\Sigma(\ell)$ is bounded: for fixed $\ell$ the module is
dominated by a bounded number of $\m$ with $\dim H_\m\to\infty$; the invariants of interest
become the embedding dimension and Loewy length of $\T_\m$ as functions of $n$ (compare the
Krull-dimension-2 picture of Claim~\ref{cl:bk}).
\end{prediction}

\begin{prediction}[The lattice is nearly split on congruence-free classes (R); formerly C-F9]\label{pr:latticesplit}\CONJ\
(restated a fifth time at v1.25.0-pre.1, \textbf{owner ruling of 2026-11-06}: on irreducible classes the hypothesis ``not dihedral'' is replaced by
``of one weight''; the v1.22.0-pre.1 form is \textbf{refuted on cohomology} at a non-ordinary class at $\ell=59$, see below
(\texttt{docs/implementation/slot25/R59\_verification.md}); it is the statement below with ``when $\bar\rho_\m$ is irreducible, of one weight'' read as
``not dihedral'' and without the definition of \emph{of one weight}, and is not repeated.
Restated a fourth time at v1.22.0-pre.1 (dihedral classes excluded; owner ruling of 2026-10-09 on census finding F33-2,
\texttt{docs/implementation/slot22/owner\_ruling\_F33-2\_1009.md}, wording of the slot-22 Fable referee \texttt{docs/implementation/slot22/F33-2\_referee.md}~\S3.3):
the v1.17.0-pre.4 form is \textbf{refuted on cohomology} at the dihedral classes at $\ell=23$ and $31$, see below; it is the statement below with
``congruence-free and not dihedral'' read as ``congruence-free'' and without the definition of \emph{dihedral}, and is not repeated.
Restated a third time at v1.17.0-pre.4, owner ruling of 2026-10-02 on finding~B-1 of the slot-17 Fable referee F1,
\texttt{docs/implementation/slot17/WP-3\_referee\_F1.md}: the v1.17.0-pre.2 form, asserted on \emph{Kummer-free} classes, is \textbf{contradicted on the
forms side} at irreducible congruence classes, the regular prime $\ell=139$ among them, see below. Restated a second time at v1.17.0-pre.2, owner ruling
O-R17-2a of 2026-10-02, before the referee: the v1.16.0 form, which excluded only the Kummer--Frobenius locus, is \textbf{refuted as stated} on cohomology at
$(608,547)$ and contradicted on the forms side at three further primes by slot-17 R17-2, see below. The label is kept. Below the new form stand, as published
and marked contradicted or refuted, the v1.17.0-pre.2 form, the v1.16.0 form (restated v1.15.1-pre.7: slot-15 WP-12,
proposal~(E-1) of \texttt{docs/implementation/slot15/WP-12\_U\_E\_proposals.md}; the Fable referee \texttt{docs/implementation/slot15/WP-4\_referee.md}~\S8;
owner ruling S15-O2 of 2026-10-01) with its evidence, and the v1.14.0 statement with its evidence.)
Let $\ell\ge5$, $n\ge1$ even, $k=n+2$, $\Lambda:=H^1(\Ga;\Sym^n\Z_\ell)$, and $\m$ a maximal ideal of $\T$ on $\Lambda$ with residue field $\Fl$, not
of Dickson type; write $a_q\in\Fl$ for the eigenvalue of $T_q$ at $\m$, and $\bar\rho_\m$ for its semisimple residual Galois representation. When
$\bar\rho_\m\cong\omega^i\oplus\omega^j$ is reducible, let $d\in[1,\ell-2]$ with $d\equiv j-i\pmod{\ell-1}$ ($d$ is odd) and call $w_1=d+1$, $w_2=\ell-d$
(so $w_1+w_2=\ell+1$, both even and, $\m$ not being of Dickson type, both in $[4,\ell-3]$) the \emph{seeds} of~$\m$; the pair $\{w_1,w_2\}$ does not depend on
the order of $i$, $j$. Call $\m$ a \emph{Kummer class} if $\bar\rho_\m$ is reducible and $\ell$ divides $B_{w_1}$ or $B_{w_2}$, and \emph{Kummer-free}
otherwise --- that is, if $\bar\rho_\m$ is irreducible, or reducible with both seeds $\ell$-regular. Call $\m$ \emph{congruence-free} if for every even $w$
with $2\le w\le\ell+1$ and every $0\le j\le\ell-2$ the generalised eigenspace of the system $(p^{-j}a_p)_{p\ne\ell}$ in $M_w(\Fl)$ has dimension at most~$1$.
On a reducible non-Dickson class this is Kummer-free (at a seed $w$ the eigenspace has dimension $1+g_w$, below); on an irreducible class it says that no two
level-$1$ eigenforms of a weight $\le\ell+1$ in the class are congruent.
Call $\m$ \emph{dihedral} if $\bar\rho_\m$ is irreducible and $\bar\rho_\m\cong\bar\rho_\m\otimes\omega^{(\ell-1)/2}$
(equivalently, by Brauer--Nesbitt, $a_p=0$ for every prime $p\ne\ell$ with $\bigl(\tfrac{p}{\ell}\bigr)=-1$; equivalently
$\bar\rho_\m$ is induced from $\Q(\sqrt{\ell^{*}})$, $\ell^{*}=(-1)^{(\ell-1)/2}\ell$, i.e.\ $\bar\rho_\m|_{G_{\Q(\zeta_\ell)}}$ is
reducible, the case the Taylor--Wiles hypothesis excludes). At level~$1$ the dihedral classes are, for $\ell\equiv3\pmod4$
with $h(-\ell)>1$, the twists $\omega^{j}\otimes\bar\rho_f$ of the classes of the cusp forms $f$ of weight $(\ell+1)/2$ with
$\bar\rho_f|_{I_\ell}\cong1\oplus\omega^{(\ell-1)/2}$ --- $f$ is its own companion, the test above sees its line at the twists
$j$ and $j+(\ell-1)/2$ --- and there are none at $\ell\equiv1\pmod4$ (parity: $\Q(\sqrt\ell)$ is real and, $\ell$ prime, has $h^{+}=h$, so an induced class unramified outside $\ell$ is even). A reducible class is never called dihedral
here, although $\omega^{i}\oplus\omega^{i+(\ell-1)/2}$ is its own quadratic twist: on those classes the statement is kept
(evidence below). (v1.22.0-pre.1.)
Call an irreducible $\m$ \emph{of one weight} if there is exactly one pair $(w,j)$, $w$ even, $2\le w\le\ell+1$, $0\le j\le\ell-2$, for which the
generalised eigenspace of the system $(p^{-j}a_p)_{p\ne\ell}$ in $M_w(\Fl)$ is non-zero --- the pairs and the eigenspaces of the test that defines
\emph{congruence-free} above; so an irreducible class that is congruence-free and of one weight has, over all pairs, exactly one non-zero eigenspace,
and it is a line. An irreducible class of one weight is \emph{not dihedral}: since $p^{(\ell-1)/2}\equiv\bigl(\tfrac{p}{\ell}\bigr)\pmod\ell$, the
system at $(w,j+\tfrac{\ell-1}{2})$ (twist read modulo $\ell-1$) is $\bigl(\bigl(\tfrac p\ell\bigr)p^{-j}a_p\bigr)_{p\ne\ell}$, which on a dihedral class
(where $a_p=0$ whenever $\bigl(\tfrac p\ell\bigr)=-1$) equals the system at $(w,j)$; so a dihedral class has either no such pair or at least two.
The new hypothesis therefore implies the old one, and the statement below asserts less than the v1.22.0-pre.1 form. Not of one weight are, besides the
dihedral classes (pairs $(w,j)$, $(w,j+\tfrac{\ell-1}{2})$, $w=\tfrac{\ell+1}2$, above), every non-ordinary class (pairs $(k_0,j)$ and
$(\ell+3-k_0,j+k_0-2)$, by the $\theta$-cycle; \STD, literature, checked at the eight $\Fl$-rational non-ordinary classes at the census primes
$\ell\le199$ in R59~\S6) and every ordinary class with a companion form (pairs $(k_0,j)$ and $(\ell+1-k_0,j+k_0-1)$, Gross's companion forms;
\STD, literature, no instance other than the dihedral ones computed). The expected meaning (the weight part of Serre's conjecture; not verified
here): an irreducible $\m$ is of one weight exactly when $\bar\rho_\m|_{I_\ell}$ is a non-split extension of two characters, ordinary without companion.
(v1.25.0-pre.1, owner ruling of 2026-11-06.)
If $\m$ is congruence-free and, when $\bar\rho_\m$ is irreducible, of one weight, then
\[
  (T_q-a_q)(T_{q'}-a_{q'})\,\Lambda_\m\subseteq\ell\,\Lambda_\m\qquad\text{for all primes }q,q'\quad(T_\ell=U_\ell);
\]
equivalently (as for the v1.14.0 statement below), the local algebra $A$ of $\T$ on $V=H^1(\Ga;\Mn)_\m$ has $\m_A^2=0$. Nothing is asserted on the other
classes. (Regularity of $\ell$ no longer implies the hypothesis: the v1.17.0-pre.2 sentence ``at a regular prime every $\m$ is Kummer-free, so in particular
the statement holds at every regular~$\ell$'' is withdrawn, contradicted at $\ell=139$ below.)

\emph{Where the excluded locus sits} (v1.17.0-pre.2; the bookkeeping of slot-17 R17-1, \texttt{docs/implementation/slot17/R17-1\_report.md}~\S\S2--5, not
binding). Fix an irregular pair $(\ell,w_K)$, $\ell\mid B_{w_K}$, and put $w_O=\ell+1-w_K$. Its Kummer classes are $\omega^a(1\oplus\omega^{w_O-1})=
\omega^b(1\oplus\omega^{w_K-1})$; R17-1 splits the cuspidal ones into the \emph{standard} family $a\in[1,w_K-2]$ and the \emph{swapped} family
$b\in[1,w_O-2]$, besides the two Eisenstein ends $a=0$ (the class of $E_{w_O}$) and $b=0$ (that of $E_{w_K}$). A Kummer--Frobenius ideal of the v1.16.0 form
(seed $w=w_O$, twist $0\le t\le J(w)=w_K-1$) is a Kummer class from its onset on, of the standard family ($t=a$) or at one of the two Eisenstein ends
($t=0$, or $t=J(w)$, the class $b=0$); the swapped family is \emph{never} Kummer--Frobenius, since that would need $\ell\mid B_{w_O}$ as
well, a double divisibility (R17-1 Lemma~3.1, proved there; none for $\ell<2\cdot10^4$, by scan; since v1.18.0-pre.1 Claim~\ref{cl:seedtwist} of this document). By R17-1's Theorem~4.1 (proved there from literature inputs, which are VERIFY there;
not a statement of [C] --- since this document now argues from it, it is to become a statement of [C] or of this document in slot~18; v1.17.0-pre.3, the
referee F1's m-E5; \textbf{since v1.18.0-pre.1 it is Claim~\ref{cl:dimformula}} of this document, with its inputs located and tagged there), $\dim V=m_1(1+2g_{w_1})+m_2(1+2g_{w_2})$ on a reducible non-Dickson class, with $m_1,m_2$ the Jordan--H\"older multiplicities of the
two Serre weights in $\Sym^n\Fl^2$ and $g_w$ the dimension of the generalised eigenspace of $E_w$'s system ($a_p=1+p^{w-1}$) in $S_w(\Fl)$, which is $0$ unless $\ell\mid B_w$. So on a reducible
congruence-free (equivalently, Kummer-free) class the coefficients $1+2g$ never exceed~$1$ ($\dim V=m_1+m_2$): on reducible classes the restated form asserts
$\m_A^2=0$ exactly where the Kummer coefficient does not enter. Its irreducible analogue --- two congruent level-$1$ eigenforms of one weight $\le\ell+1$ ---
is what the hypothesis now also excludes (v1.17.0-pre.4).

\emph{Contradicted on the forms side: the v1.17.0-pre.2 form} (v1.17.0-pre.4; the slot-17 referee F1, finding~B-1, reproduced with tracked tooling by
slot-17 R17-2, \texttt{docs/implementation/slot17/R17-2\_report.md}~\S7: \texttt{r172/r172\_forms.py --seed}, \texttt{r172/r172\_seedcheck.py}; computed by
script, forms side only, nothing here is cohomology). The v1.17.0-pre.2 form asserted (R) on \emph{Kummer-free} classes (residual representation irreducible, or
reducible with both seeds $\ell$-regular), and added that it therefore holds at every regular~$\ell$. At $(\ell,k_0)=(67,32)$, $(131,28)$, $(139,28)$ the space
$S_{k_0}$ has dimension~$2$ and is one Galois orbit over a quadratic Hecke field in which $\ell$ ramifies (the discriminant of the $T_2$-charpoly is
$2^6\cdot3^2\cdot67\cdot273067$, resp.\ $2^6\cdot3^6\cdot131\cdot139$); the two eigenforms are congruent modulo the prime above~$\ell$, confirmed over~$\Z$
($(T_p-\lambda_p)^2=0\ne T_p-\lambda_p$ on $S_{k_0}(\Fl)$ for every $p\le23$), and no $p^a+p^b$ fits the system, so the class is irreducible, hence Kummer-free;
$67\mid B_{58}$ and $131\mid B_{22}$, while $139$ is \textbf{regular}. At the $\theta^j$-twists $k=\ell k_0+j(\ell+1)$, $j=1,2$, i.e.\ the six cells
$n=2210$, $2278$ ($\ell=67$), $3798$, $3930$ ($\ell=131$), $4030$, $4170$ ($\ell=139$), the generalised eigenspace $W\subseteq S_k(\Fl)$ has $\dim W=4$,
Hilbert function $[1,1,1,1]$, $\dim A=\dim W$ (generated; the same with the anemic set $\{2,3,5,7\}$ and with $p\le43$), $T_2$ of Jordan type~$[4]$ and
$T_\ell$ nilpotent: $\T^{\mathrm f}_\m/\ell\cong\Fl[x]/x^4$. Six further pairs, $(73,40)$, $(71,54)$, $(89,68)$, $(181,38)$, $(227,46)$, $(349,38)$, each
first confirmed over~$\Z$, show the same $\Fl[x]/x^4$ at both onset cells ($12$ cells). By \HA{prop:cohforms}(5) these blocks are torsion-free and $A$ is the
quotient of $\T^{\mathrm f}_\m/\ell$ by the saturation defect (\HA{rmk:sattor}), so the pre.2 form fails there unless $\dim I\ge2$; $A$ itself was not
computed (the cheapest cohomology candidate, $(2210,67)$ --- prediction $\dim V=8$, $A\cong\F_{67}[x]/x^4$, saturated --- is cluster-only; predict RSS first).
Since $139$ is regular, the sentence ``(R) holds at every regular~$\ell$'' (slot~15, \texttt{WP-4\_referee.md}~\S8, kept in the v1.16.0 and pre.2 forms) is
contradicted on the forms side, and so is the v1.16.0 form at this second, independent locus (an irreducible class is never Kummer--Frobenius). None of these
classes is congruence-free (at $w=k_0$, $j=0$ the eigenspace has dimension~$2$), so none is a counterexample to the restated form.

\emph{Refuted on cohomology: the v1.17.0-pre.4 form, at dihedral classes} (v1.22.0-pre.1; census~3.3 finding F33-2,
\texttt{docs/implementation/slot22/C33\_followups\_report.md}, the Fable referee \texttt{F33-2\_referee.md}; computed).
At $\ell=23$ the class of $\Delta$ and at $\ell=31$ that of the weight-$16$ form are dihedral by $\Q(\sqrt{-\ell})$:
$\bar\rho\cong\bar\rho\otimes\omega^{(\ell-1)/2}$, $\bar\rho|_{I_\ell}\cong1\oplus\omega^{(\ell-1)/2}$, the form its own
companion. Every twist is congruence-free in the sense above (the only non-zero eigenspaces are the form's own line at
$(w,j)=((\ell+1)/2,0)$ and $((\ell+1)/2,(\ell-1)/2)$, each of dimension~$1$), and the v1.17.0-pre.4 form asserted
$\m_A^2=0$ there. Census~3.3 (\texttt{rows33.tsv}, C1a local run, cells certified, $24$ cells recomputed byte for byte)
reads $A\cong\Fl[x]/x^{3}$, Hilbert function $[1,1,1]$, on $161$ rows: $140$ at $\ell=23$ ($n\ge298$, $\dim V=8$ or $12$,
module Loewy layers $[4,2,2]$ and $[6,4,2]$) and $21$ at $\ell=31$ ($n\ge526$, $\dim V=8$), at every twist $s$, with
$(T_2-a_2)^{2}\ne0$ on $V$ on each (recomputed from the records' matrices); the operators of the primes inert in
$\Q(\sqrt{-\ell})$ and $U_\ell$ lie in $\m_A^{2}$. The forms side agrees: at $(298,23)$, $(540,23)$ and $(526,31)$,
$\T^{\mathrm f}_\m/\ell\cong\Fl[x]/x^{4}$ on the four-dimensional generalised eigenspace (by script, modulo $\ell^{6}$:
every eigenform of the class is ramified at $\ell$, Newton slope $\tfrac12$ at the split primes and $1$ at the inert
ones, $\T^{\mathrm f}_\m$ monogenic), and $A$ is its quotient by a saturation defect of length $1$ (\HA{prop:cohforms}(5));
at $(540,23)$ the defect has length~$2$ and $A\cong\Fl[x]/x^{2}$, so the seven rows $n=540+24t$ (among the $768$ consistent dihedral rows, not among the $161$) hold the statement only
through the defect. The onset is $n=\ell(\ell+3)/2-1$ ($298$, $526$); for $\ell=59,71,79,83$ it lies beyond $n\le700$.
No non-dihedral class in census~3.3 reaches Loewy length~$3$ (the referee's \S4: of the $100\,025$ non-Dickson, non-dihedral
rows of residue degree~$1$, Loewy length~$1$ on $85\,123$, $2$ on $14\,818$, $4$ on the $84$ Kummer--Frobenius rows at
$(\ell,w_K)=(37,32)$, which are not congruence-free, and $3$ on none). The hypothesis now
excludes dihedral classes ($929$ rows of residue degree~$1$ in census~3.3: the $161$ above, and $768$ consistent with $\m_A^2=0$,
below their onsets); the reducible self-twist classes ($\omega^{i}\oplus\omega^{i+(\ell-1)/2}$, both seeds
$(\ell+1)/2$) stay inside: $1770$ rows at $\ell=7,11,19,23,31,43,\dots$, $\dim V$ up to $24$, all with $\m_A^{2}=0$.
Whether the true obstruction is the self-twist or, more generally, two cuspidal Serre weights cannot be decided on census~3.3
(no ordinary class with a cuspidal companion other than the self-twist ones occurs at level~$1$ and $\ell\le103$; the referee's
\S\S2.2, 2.6): a non-dihedral class of that shape with $\m_A^2\ne0$ would refute the present form. (v1.25.0-pre.1: one does, at $\ell=59$, a
non-ordinary class --- the slot-22 scan took ``ordinary'' as $a_\ell\ne0$ and did not look at non-ordinary classes, every one of which has that shape;
next paragraph.)

\emph{Refuted on cohomology: the v1.22.0-pre.1 form, at a non-ordinary class} (v1.25.0-pre.1, owner ruling of 2026-11-06; census~4.0 rows found by the
gallery referee, \texttt{docs/implementation/slot25/PG140\_referee.md}~R1, verified independently from scratch, the data source
\texttt{docs/implementation/slot25/R59\_verification.md}, scripts and outputs in its \texttt{r59/}; computed). At $\ell=59$ the class of $f_{16}=\Delta E_4$,
system $(a_2,a_3,a_5,a_7,a_{59})\equiv(39,15,13,14,0)$, is irreducible (no $p^{i}+p^{j}$ fits the system at the $11$ primes $p\le31$), so neither of Dickson
type nor a Kummer class; \emph{not dihedral} ($59\equiv3\bmod8$, so $\bigl(\tfrac{2}{59}\bigr)=-1$, while $a_2=39\ne0$; also $a_{11}=41\ne0$); and
\emph{congruence-free}: over all even $w\le60$ and all $58$ twists the eigenspace in $M_w(\F_{59})$ is non-zero at exactly the two pairs $(16,0)$ and
$(46,14)$, a line at both (exact at this $\ell$: $T_2,T_3,T_5$ span the Hecke algebra of $M_w(\F_{59})$ for $w\le60$). It is \emph{non-ordinary},
$a_{59}\equiv0$, and so has two cuspidal Serre weights, $16$ and $46=\ell+3-16$: it is not of one weight. Census~4.0 (main section) has four Loewy-$3$ rows
of the class, residue degree~$1$, \texttt{final}, \texttt{sturm(B$'$)}: $n=884$ and $944$ ($\theta^0$, resp.\ $\theta^1$; $\dim V=6$,
$A\cong\F_{59}[x]/x^{3}$, module Loewy layers $[2,2,2]$) and $n=942$, $1000$ ($\theta^0$; $\dim V=8$, Hilbert function $(1,2,1)$, layers $[4,2,2]$).
Recomputed with the shared binary from the Hecke matrices on $H^1$: $(T_2-a_2)^2\ne0$ on $V$ at all four cells (ranks of $T_2-a_2$, its square and cube
$4,2,0$; block dimensions $6$, $8$, $6$, $8$), and at $n=884$ by a second, independent linear-algebra route on the dumped matrices. On the forms side
$\T^{\mathrm f}_\m/\ell$ has the census Hilbert function at all four cells, so the saturation defect is $0$ there (\HA{prop:cohforms}(5)); the blocks are
torsion-free ($\dim V=2\dim W$). So the v1.22.0-pre.1 form, which asserted $\m_A^2=0$ on this congruence-free, non-dihedral class, fails; nothing proved
in this document and no data sentence becomes false (the census-3.3 sentences above are scoped to $n\le700$; the onset here is
$n=\ell(k_0-1)-1=884$). The forms side repeats the pattern at a non-ordinary class with large image, $\ell=79$, weight~$38$ (companion weight $44$):
$\T^{\mathrm f}_\m/\ell$ of Hilbert function $(1,1,1)$ at $n=2922=79\cdot37-1$ and $(1,2,1)$ at $n=3000$ (by script; cohomology not computed), so the
obstruction is non-ordinarity, not the $S_4$-type image of the class at $59$. \emph{Consistency of the restated form} (R59~\S7.1): of the $672$ rows of
census~4.0 (main and b1; section b2 has none) with Loewy length $\ge3$, all of residue degree~$1$, the $418$ dihedral rows at $\ell=23$, $31$ and the $4$
rows at $59$ are not of one weight, the $235$ Kummer rows at $\ell=37$, $67$, $547$ and the $15$ rows at $\ell=139$ ($\theta^{j}(f_{138})$, a congruence
in weight $138$) are not congruence-free; and the other irreducible, non-dihedral rows of residue degree~$1$ with $\dim V\ge6$ (where (R) has content)
are $5805$ rows at $11\le\ell\le31$, $23$ ordinary classes (of $\Delta$, $f_{16}$, $f_{18}$, $f_{20}$, $f_{22}$, $f_{24}$, $f_{26}$) each at exactly one pair
$(w,j)$, hence of one weight, with Loewy length~$1$ on $2008$ and $2$ on $3797$. So no row of census~4.0 contradicts the restated form, and within the
census it newly excludes, among the rows with $\dim V\ge6$, only the four at $\ell=59$. The narrower alternative ``congruence-free, not dihedral and
ordinary'' (R59~\S7.3) would keep the ordinary classes with a companion other than themselves inside, the shape warned about above, of which no instance
has data; the owner chose ``of one weight''.

\emph{Refuted as stated: the v1.16.0 form} (v1.17.0-pre.2; slot-17 R17-2, \texttt{docs/implementation/slot17/R17-2\_report.md}; every item at a Kummer class
that is \emph{not} Kummer--Frobenius, so a counterexample to the v1.16.0 form and none to the restated one):
\begin{enumerate}
\item \emph{On cohomology} (computed; R17-2~\S6.1). At $\ell=547$, $w_K=486$ has $g_{486}=2$ (two cusp forms of weight $486$ congruent to $E_{486}$ modulo
  $547$, while $v_{547}(B_{486})=1$; R17-1~\S4, by script), so every cuspidal Kummer class at $(547,486)$ already has rank $1+g=3$ in its first window. The
  class $a=1$ at $n=608$ ($k=610$, $a_p\equiv p+p^{62}$) lies far below its Kummer--Frobenius onset $k\ge547\cdot62+548=34462$.
  \texttt{HeckeSaturationDefectValidation} with \texttt{HECKE\_SAT\_DEFECT\_CELLS="608:547"} (the slot-17 shared binary; operators $T_2,T_3,T_5,T_7,U_{547}$;
  $e_{\max}=0$; 23\,min, $3.7$\,GiB): on the block of the system $(a_2,a_3,a_5,a_7,a_{547})=(17,185,281,442,0)$, $\dim V=6$, $\rank\T_\m=3$, the census algebra
  $A$ and $\T_\m/547$ both have Hilbert function $(1,1,1)$, so $A\cong\F_{547}[x]/x^3$; the saturation defect is $I_\m=0$ ($\T_\m$ saturated, \HA{prop:satdefect}),
  and the \texttt{stable} verdict is certified at $K=2$ (\HA{cor:stablecert}). Hence $\m_A^2\ne0$ on $V$. The forms side reads the same algebra
  ($\T^{\mathrm f}_\m/547\cong\F_{547}[x]/x^3$, $\dim W=3$, by script; the run's own forms cross-route agrees), as \HA{prop:cohforms}(5) requires on a saturated
  torsion-free block.
\item \emph{On the forms side only} (computed by script, R17-2~\S\S2--3: the algebra generated by the $T_p$, $p\le23$, on the generalised eigenspace $W$ of
  $S_k(\Fl)$, with $\dim A=\dim W$ as its generation certificate, unchanged with the anemic set $\{2,3,5,7\}$ and with $p\le43$): $\T^{\mathrm f}_\m/\ell\cong
  \Fl[x]/x^4$ at the swapped classes $b=1,\dots,4$ at $\ell=37$ ($k=1222$, $1260$, $1298$, $1336$; $n=1220,\dots,1334$), $b=1,2$ at $\ell=103$ ($k=2576$, $2680$)
  and $b=1,2$ at $\ell=691$ ($k=8984$, $9676$; the Kummer seed $12$ of $\Delta$); and $\T^{\mathrm f}_\m/547\cong\F_{547}[x]/x^3$ at $(1156,547)$ (class $a=2$).
  On these torsion-free blocks \HA{prop:cohforms}(5) makes $A$ the quotient of $\T^{\mathrm f}_\m/\ell$ by the saturation defect, so the census-algebra form fails
  there unless $\dim I\ge2$ ($\ge1$ at $547$), \HA{rmk:sattor} (v1.17.0-pre.3 wording, the referee F1's m-E4). $A$ itself was not computed at these cells ($(1220,37)$ needs an exclusive memory window, $(2574,103)$
  and $(8982,691)$ are beyond local reach); every saturation defect measured so far at the Kummer classes has length $\le1$ (elsewhere defects are
  longer: length $2$ at $(540,23)$ above; $\delta$ up to $190$ in census~4.0; v1.24.0-rc.2), and every Loewy-$4$ Kummer block measured is saturated, so
  the expectation (not a claim) is that cohomology agrees.
\end{enumerate}

\emph{What held at the Kummer classes tested} (now outside the assertion; recorded as data). The Eisenstein ends: $\m^2=0$ on the cusp side of the $b=0$
classes $(1182,37)$ ($E_{32}$), $(484,547)$ ($E_{486}$) and $(8290,691)$ ($\Delta$) (forms side; the Eisenstein part needs cohomology). The pre-jump
cells of R17-1~\S5.2 on cohomology (R17-2~\S6.2, console route with $U_\ell$): $(738,67)$, $(944,59)$, $(1184,37)$ have module Loewy layers $[3,1]$ and
$A\cong\Fl[x]/x^2$ ($(670,67)$, the census control, reproduced); $(738,67)$ is saturated and \texttt{stable} certified.

\emph{Observed on the Kummer locus} (not asserted here; a candidate for a later restatement of Prediction~\ref{pr:kfonset}, slot~18 --- taken up at v1.18.0-pre.1
as Prediction~\ref{pr:kfonset}~(B), (C) and Prediction~\ref{pr:depth}, both \CONJ). Every cuspidal Kummer-class block
computed below the second growth of the multiplicities ($(m_1,m_2)\le(2,2)$) is curvilinear of the largest length its rank allows, $A$ (or, where only the forms side
was computed, $\T^{\mathrm f}_\m/\ell$) $\cong\Fl[x]/x^{\dim V/2}$: the $24$ Kummer--Frobenius onset cells at $37$ and the slot-17 onset cells at $67$, $59$, $101$
(Prediction~\ref{pr:kfonset}); the swapped cells and the $g=2$ first windows of item~2 and of $(608,547)$; and, at rank~$2$, the Kummer rows of edition~3 with
$\dim V=4$ (modules $[2,2]$ and the pre-jump $[3,1]$). At the second growth $(3,3)$ it fails on the forms side ($(1,2,2,1)$ at $k=1412$, $1450$, $\ell=37$;
Prediction~\ref{pr:kfonset}).

\emph{Evidence for the restated form.} Slot~17 tested only Kummer classes and irreducible congruence classes, none of them congruence-free, and adds none. The
evidence is the v1.16.0 form's table below, read on its congruence-free rows: at $\ell\in\{7,11,13\}$ every row is congruence-free ($M_w$, $w\le\ell+1$, has
dimension $\le2$, and its one cusp form $\Delta$, at $w=12$, is irreducible modulo $11$ and~$13$); at the other primes the rows mix congruence-free classes,
Kummer classes and irreducible congruence classes, every row with Loewy length $\ge3$ is a Kummer class (the $18$ Kummer--Frobenius records and six weight-shift
cells at~$37$), and the irreducible congruence classes with $\ell\le107$ ($\ell=67$, $71$, $73$, $89$) have their onsets at $n\ge2210$, beyond every row, so
passing from Kummer-free to congruence-free loses no evidence (v1.17.0-pre.4; the referee F1's B-1, item~4). No row in the table violates the restated form. The $161$ dihedral rows of census~3.3 (above) violate the
v1.17.0-pre.4 form, which did not exclude dihedral classes; they lie outside the present form (v1.22.0-pre.1, the
slot-22 Fable referee \texttt{docs/implementation/slot22/F33-2\_referee.md}, owner ruling of 2026-10-09). (v1.17.0-pre.3, the referee F1's m-E3.) On a torsion-free block with $\dim V\le4$, $\m_A^2=0$ is automatic:
$\dim A\le\rank\T_\m=\lceil\dim V/2\rceil\le2$ (\HA{lem:heckerankone}, \HA{prop:cohforms}(5)); so the $120+229$ below-onset Kummer rows were never evidence, and
the content of the prediction lies at $\dim V\ge5$ (the rows at $\ell\le13$ with $\dim V$ up to $26$, and the three-copy rows at~$13$). Whether it holds on every
congruence-free class is open; the C++ readings keep \emph{reporting}, never gating.

\medskip\noindent\textbf{The v1.16.0 form (restated v1.15.1-pre.7; refuted as stated at v1.17.0-pre.2 by $(608,547)$, above; kept as published, with its
evidence and its refutation of the v1.14.0 statement).} The definition of \emph{Kummer--Frobenius} made here is kept: Prediction~\ref{pr:kfonset} uses it.
Let $\ell\ge5$, $n\ge1$ even, $k=n+2$, $\Lambda:=H^1(\Ga;\Sym^n\Z_\ell)$, and $\m$ a maximal ideal of $\T$ on $\Lambda$ with residue field
$\Fl$, not of Dickson type; write $a_q\in\Fl$ for the eigenvalue of $T_q$ at $\m$. Call $\m$ \emph{Kummer--Frobenius} if
\begin{itemize}
\item its residual representation is $\omega^t(1\oplus\omega^{w-1})$ for an even $w$ with $4\le w\le\ell-3$;
\item $\ell$ divides the companion Bernoulli number $B_{\ell+1-w}$;
\item $k\ge\ell w+t(\ell+1)$, with $0\le t\le J(w)$, where $J(w)\equiv-w\pmod\ell$ (the twist $\theta^tE_w$ has not dropped filtration).
\end{itemize}
If $\m$ is not Kummer--Frobenius, then
\[
  (T_q-a_q)(T_{q'}-a_{q'})\,\Lambda_\m\subseteq\ell\,\Lambda_\m\qquad\text{for all primes }q,q'\quad(T_\ell=U_\ell).
\]
Equivalently, the local algebra $A$ of $\T$ on $V=H^1(\Ga;\Mn)_\m$ has $\m_A^2=0$. In particular this holds at every regular
$\ell$ (the referee's form~(R), \texttt{WP-4\_referee.md}~\S8).

\emph{Refuted as stated} (v1.15.1-pre.7). The v1.14.0 statement (no Kummer--Frobenius exclusion) fails at $\ell=37$, by three
independent routes:
(1)~\emph{certified $\Fl$ records}: the ideals $\operatorname{eis}(a,a+5)$ (residual $\omega^a(1\oplus\omega^5)$, $37\,\|\,B_{32}$)
of the slot-14 C1a records (first-cluster job \texttt{34567994}, staged data set \texttt{np\_dataset\_c15}, not published) at
$n=654,656,\dots,676$ ($a=1,\dots,12$) and $n=690,\dots,700$ ($a=1,\dots,6$) --- $18$ records, each \texttt{verify()}$=0$,
\texttt{final = true} by \texttt{sturm($B'$)}, \texttt{sturm\_proof\_range = true} --- have \texttt{loewy\_A}$\;=4$, embedding
dimension~$1$, Hilbert function $[1,1,1,1]$, i.e.\ $A\cong\F_{37}[x]/x^4$, and module Loewy layers $[3,2,2,1]$
($V\cong A\oplus A/x^3\oplus A/x$, $\dim V=8$); since $\m_A$ is generated by the $T_q-a_q$, some product
$(T_q-a_q)(T_{q'}-a_{q'})$ is non-zero on $V$, whatever the full generator set (referee \S8, point~2);
(2)~\emph{the forms side} (slot-15 WP-12~\S1.2): $\T^{\mathrm f}_\m/\ell\cong\F_{37}[x]/x^4$ on $S_k(\F_{37})$, $k=656,\dots,702$,
consistent with \HA{prop:cohforms}(5) on these torsion-free blocks (none of the $18$ is of Dickson type; the Dickson locus in the
window is $n=684$);
(3)~\emph{the lattice mod~$37^2$} (slot-15 WP-13, \texttt{HeckeCensusLatticeValidation} at $(654,37)$,
\texttt{docs/implementation/slot15/WP-13\_report.md}): ideal~$8$ --- the system $a_{2,3,5,7}\equiv(29,29,16,33)$, $U_{37}\equiv0$,
$\dim V=8$, the cluster record's ideal~$8$ --- prints \texttt{CENSUS\_LATTICE 654 37 8 \dots{} status=failed cf9\_lattice=false
cf9\_field=false first\_fail=2,2 m\_nonzero=true lattice\_rank=8 loewy\_V=4}: $(T_2-29)^2\Lambda_\m\not\subseteq37\Lambda_\m$.
The failing product does not involve $U_{37}$, so the refutation does not depend on $U_\ell$; ideal~$1$ of the same cell
($\operatorname{eis}(2,5)$, not Kummer--Frobenius) holds. All $18$ ideals are Kummer--Frobenius ($w=6$, $t=a$), so none is a
counterexample to the restated form.

\emph{Evidence for the restated form} (of v1.15.1-pre.7, i.e.\ the v1.16.0 form; every item at a non-Kummer--Frobenius $\m$; no item against it until
slot~17):
\begin{center}\small
\begin{tabular}{@{}>{\raggedright\arraybackslash}p{3.3cm}>{\raggedright\arraybackslash}p{2.2cm}>{\raggedright\arraybackslash}p{2.2cm}>{\raggedright\arraybackslash}p{2.4cm}>{\raggedright\arraybackslash}p{4.3cm}@{}}
\toprule
source & $\ell$ & range & rows / blocks & max Loewy of $A$ ($A_1$) \\
\midrule
census edition~2 (slot-13 WP-9 lattice test) & $7,11,13,691$ & $n\le300$ & $654$ non-reduced rows &
  $2$; C-F9 held $654/654$ (at $691$ the Kummer class $B_{12}$ has onset $k\ge691\cdot680$, far beyond) \\
WP-M1 route (first cluster, \texttt{34551683}) & $7,11,13$ & $n\le200$ / $150$ & $1233$ non-reduced blocks & $\m_A^2=0$ on all \\
staged np data set (slot-14 WP-C15) & $37,\dots,103$ & $n\le700$ (partial) & $2008$ non-reduced rows &
  $2$ everywhere except the $18$ Kummer--Frobenius rows at $37$; this includes the $120$ non-reduced Kummer-predicted rows
  (Loewy~$2$) below their onset at $37,59,101,103$ ($14,70,19,17$; onsets $k\ge37\cdot6+38t$, $59\cdot16$, $101\cdot34$,
  $103\cdot80$) \\
census edition~3, its complete C1b block (v1.16.1; first-cluster \texttt{34568172} and RB1--RB4 \texttt{34573807}--\texttt{10}, slot-15
  WP-C18, WP-15) & $61,67,\dots,103$ (ten primes) & every even $n\le700$, $3500/3500$ cells & $62\,113$ rows, $2165$ non-reduced &
  $2$ at every prime (embedding dimension $\le1$, every block Gorenstein); no Kummer--Frobenius ideal of Loewy $\ge3$ exists there
  (the only Kummer--Frobenius ideals in range are the two $t=0$ ideals at $(668,67)$, Loewy $2$ and $1$), so $0$ violations of~(R) \\
weight-shift $A_1$ (WP-L0930, \texttt{a1\_loewy\_scan.py}) & $17,19,29,37,47,59,71,79,83,107,131$ & $n\le360$ / $200$ / $180$ &
  $1788$ blocks with a recorded $A_1$ & $2$, except \emph{exactly} the six Kummer--Frobenius cells $(258,37)$, $(294,37)$,
  $(330,37)$ ($\operatorname{eis}(1,6)$), $(296,37)$, $(332,37)$ ($\operatorname{eis}(2,7)$), $(334,37)$ ($\operatorname{eis}(3,8)$),
  each $[1,1,1,1]$ \\
C-F9 lattice at $(654,37)$, slot-15 WP-13 & $37$ & $1$ cell & ideal~$1$ ($\operatorname{eis}(2,5)$, not K--F): held; ideal~$8$
  (K--F): failed & --- \\
\bottomrule
\end{tabular}
\end{center}
\emph{Process note.} The slot-13 and slot-14 evidence for the v1.14.0 statement (the $654/654$ and the $1233/1233$) lay entirely
below every Kummer--Frobenius onset, so it never tested the refuted case. \emph{The open part}: the case $t=0$ (the untwisted
$E_w$ class) is not covered by the data --- there the block is the Eisenstein ideal itself and F9's (Mazur--Wiles) shape applies;
the next tests are $(736,67)$ ($67\mid B_{58}$, $w=10$, $\operatorname{eis}(1,10)$) and $(372,37)$ ($t=4$). The remaining $17$
Loewy-$4$ cells at $\ell=37$ ($n=656,\dots,700$) are lattice-test candidates for the cluster ([X] v1.20.1, WP-C19). (v1.16.1) The rows of the table are, from $\ell=17$ on, the rows of census edition~3 (\S\ref{sec:census}); there (R) has
$0$ violations on all $80\,293$ rows present, and the C-F9 lattice test held on all $155$ evaluable cheap new-prime rows ($9$ rows of
residue degree $>1$ not evaluated). Whether the
restated form holds at every regular prime is open. The C++ readings (\texttt{cf9\_holds()},
\texttt{HeckeCensusLatticeValidation}'s summary, \texttt{LOEWY\_RING\_CF9\_FINDING}) keep \emph{reporting}, never gating, as before.

\medskip\noindent\textbf{The v1.14.0 statement (refuted as stated, v1.15.1-pre.7; kept as published, with its evidence).}
(v1.14.0-pre.2; candidate~C-F9 of
the slot-12 finality proposal \texttt{docs/implementation/slot12/WP-12\_finality\_proposal.md} v2 \S7, tested by
slot-13 WP-9, owner answer Q8)
Let $\ell\ge5$, $n\ge1$, $\Lambda:=H^1(\Ga;\Sym^n\Z_\ell)$, and $\m$ a maximal ideal of the Hecke algebra $\T$ acting
on $\Lambda$ with residue field $\Fl$; write $a_q\in\Fl$ for the eigenvalue of $T_q$ at $\m$. Then
\[
  (T_q-a_q)(T_{q'}-a_{q'})\,\Lambda_\m\subseteq\ell\,\Lambda_\m\qquad\text{for all primes }q,q'
\]
(with $T_\ell=U_\ell$). \emph{Equivalently}: since $H^2(\Ga;\Sym^n\Z_\ell)[\ell]=0$ (\HA{cl:bockstein}(1)),
$\Lambda/\ell=H^1(\Ga;\Mn)$ Hecke-equivariantly, so the statement is that the local algebra $A$ of $\T$ on
$V=H^1(\Ga;\Mn)_\m$ has $\m_A^2=0$, i.e.\ Loewy length $\le2$ --- in contrast with $\T_\m/\ell$ acting on $M_k(\Fl)$,
whose Loewy length reaches~$6$ on the same census (slot-12 WP-12, the forms side). (v1.14.0-pre.3, referee
finding~M4.) The mod-$\ell$ form needs no hypothesis ``$\m$ not of Dickson type'' for the \emph{equivalence} (the
identification $\Lambda/\ell=H^1(\Ga;\Mn)$ holds for every $\m$ at $\ell\ge5$). The \emph{evidence}, however, is
entirely non-Dickson: every tested row is \texttt{mixed} or \texttt{cuspidal}, has a free lattice and
\texttt{sturm\_proof\_range = true}, and every Dickson-type block computed so far is reduced ($\m_A=0$, where the
statement is trivial: the $340$ Dickson-type census rows and the $100$ Dickson-type blocks of slot-13 WP-11's scan).
On a non-reduced Dickson-type $\m$ the prediction therefore has \emph{zero} evidence --- an untested extrapolation.
(v1.15.0, slot-14 WP-6, \texttt{HeckeRingLoewyValidation}'s route over $\Z/\ell^{K}$: the $206$ Dickson-type (torsion) blocks
read on $348$ cells --- $101$ of them in the window $\ell=7$, $n\le200$; $\ell\in\{11,13\}$, $n\le150$ --- are all reduced mod~$\ell$, so
the prediction holds trivially on them; a non-reduced Dickson-type block has not occurred, and the extrapolation stays untested.)
\emph{Evidence} (computed; census edition~2, \S\ref{sec:census}): on all $654$ non-reduced $H^1$ rows
($\ell\in\{7,11,13,691\}$, $n\le300$, $\dim V\le26$; $128$ / $298$ / $226$ / $2$ by~$\ell$) the statement holds, and
fails on none; it was tested twice, independently --- on the lattice mod~$\ell^2$ by the module route
($H^1(\Ga;\Sym^n\Z/\ell^2)$, the $\m$-part cut out by Fitting projectors, never by the $\Fl$ decomposition) and on $V$
over~$\Fl$ --- and the two agree on every row (slot-13 WP-9, \texttt{HeckeCensusLatticeValidation}, its
all-rows mode \texttt{HECKE\_CENSUS\_LATTICE\_ALL=1}; the TSV \texttt{docs/implementation/slot13/WP-9\_census\_lattice.tsv}).
(v1.15.0) \emph{A third, independent route}: slot-14 WP-M1's scan (the $\Z/\ell^2$ module route and joint idempotents, with no
$\Fl$ decomposition; \texttt{docs/implementation/slot14/WP-M1\_loewy\_note.md}) re-derived $A$ on $1229$ $\Fl$-rational blocks in $250$
cells ($\ell=7$, even $n\le200$; $\ell\in\{11,13\}$, even $n\le150$). The $259$ non-reduced blocks among them are the census's rows in
that window, and all have $\m_A^2=0$. The $\Delta$-block of $(400,7)$, outside the census window, also has $\m_A^2=0$. The C++
route over $\Z/\ell^{K}$ (slot-14 WP-6, \texttt{HeckeRingLoewyValidation}; \texttt{docs/implementation/slot14/WP-6\_report.md}~\S3)
re-derived that window block for block and $98$ further cells ($348$ cells): the statement held on $452/452$ non-reduced blocks.
(v1.15.1-pre.3; P12's cluster half, owner sign-off 2026-09-30) \emph{Beyond the census window}: the WP-M1 route on the first cluster (job
\texttt{34551683}: $\ell=7$, $n=402\ldots604$; $\ell\in\{11,13\}$, $n=302\ldots556$) found $\m_A^2=0$ on all $1233$ non-reduced
$\Fl$-rational blocks of its completed cells ($98$ / $384$ / $751$ by~$\ell$; the one further non-reduced block, recovered from the
killed cell $(604,7)$, also has it), and on none fails; the data set is \texttt{docs/matrix-engine/results/noctua2\_section\_m/}
(published 2026-09-30). The forms side ($q$-expansions, all $T_q$ up to the Sturm bound) re-derived $\T_\m/\ell$ on all $343$ blocks of
the WP-M1 scans with a $\Z/\ell^{K}$ algebra and at $(400,7)$, $(402,7)$, $(500,7)$, $(604,7)$ (slot-14 WP-7,
\texttt{docs/implementation/slot14/WP-7\_referee.md}~\S3). The tag stays \CONJ.
\emph{What the prediction is not}: the Hecke algebra over $\Z/\ell^{K}$, reduced afterwards ($A_K\otimes\Fl$ on the $\m$-part of $H^1(\Ga;\Sym^n\Z/\ell^{K})$), has Loewy
length $3$--$6$ on census rows and exactly~$8$ at $(400,7)$. The prediction therefore concerns the reduction of the \emph{module}, not
of the algebra.
The pairs tested are $q,q'\in S=\{2,3,5,7\}\cup\{\ell\}$; since every row of edition~2 is \texttt{final = true}
($A_S=\T_\infty|_V$, \S\ref{sec:certgen}), $\m_{A_S}^2=0$ is the statement for all $q,q'$. The $2250$ reduced $H^1$
rows satisfy it trivially ($\m_A=0$). \emph{What does not explain it.} The candidate mechanism of the proposal ---
the Jordan--H\"older filtration of $\Sym^n(\Fl^2)$ by Serre weights $F(a,b)=\Sym^a\otimes\det^b$, which does give $V$ a
Hecke filtration with scalar graded pieces --- does not bound the Loewy length: at most two \emph{distinct} weights
occur on any row, the Loewy length of $V$ is~$2$ on \emph{every} row while the number of pieces with multiplicity
ranges over $1$--$26$, and ``Loewy length $\le$ number of distinct weights'' \textbf{fails on $89$ rows} (the $87$
cuspidal rows at $\ell=13$, where the weight of $\Delta$ occurs two or three times --- at $(192,13)$ and $(194,13)$ all
three copies lie in~$V$ --- and the two $\ell=691$ rows, where $\Sym^n$ is irreducible). Why extensions between repeated
copies of one weight never nest deeper is open. \emph{Not asserted}: any statement at $\ell\in\{2,3\}$, at residue
degree $>1$, or on $X_i$ other than $X_4$; no code gates this prediction beyond the census test itself.
\emph{Evidence (slot~19, v1.19.0-pre.5)}: at $(1220,37)$ nine congruence-free hidden blocks with $A_1=\F_{37}$ (consistent; $\m_A^{2}=0$ trivially) and one Kummer block with
$\m_{A_1}^{2}\ne0$ (outside the statement; Remark~\ref{rmk:frobkummer}). No change of the statement.
\end{prediction}

\begin{prediction}[The Kummer onset and its window: the standard and the swapped family, and $g=2$]\label{pr:kfonset}\CONJ\ (v1.15.1-pre.7; companion of the restated
Prediction~\ref{pr:latticesplit}, slot-15 WP-12~(E-1); \textbf{restated at v1.17.0-pre.1} with the window bound proposed by the $\ell=37$ report
\texttt{docs/latex/l37\_loewy4\_example.tex} (not binding; its \texttt{rmk:kfonset}, \texttt{conj:window}, \texttt{prop:secondwindow}), for the slot-17
referee and the owner's ruling~O3 of 2026-10-01. \textbf{Generalised at v1.18.0-pre.1} (hecke-extensions slot~18, WP-2; the slot-17 WP-2 flag and the slot-17
referee F1's B-2 and m-E5; the owner's ruling~O3 of 2026-10-03): part~(A) is the v1.17.0 statement, unchanged in content and reworded only through
Claim~\ref{cl:seedtwist}; parts~(B), the swapped family, and~(C), $g=2$, are new, both \CONJ; the title read ``The Kummer--Frobenius onset, and its window''
through v1.17.0. The label is kept; the v1.15.1-pre.7 statement stands below, marked superseded.)
Let $(\ell,w_K)$ be an irregular pair, $w_O:=\ell+1-w_K$, $g:=g_{w_K}$ (Claim~\ref{cl:dimformula}), and $\m$ a cuspidal Kummer class of the pair, presented
by Claim~\ref{cl:seedtwist}(2) as $\bar\rho_\m\cong\omega^{t}(1\oplus\omega^{w-1})$ with $w\in\{w_O,w_K\}$ and $1\le t\le J(w)-1$: $w=w_O$ is the standard family
(there $\m$ is Kummer--Frobenius in the sense of Prediction~\ref{pr:latticesplit}'s v1.16.0 form from its onset on, $\ell$ dividing the companion $B_{\ell+1-w}$),
$w=w_K$ the swapped family ($\ell\mid B_w$ itself; never Kummer--Frobenius). Write $A$ for the census algebra on $V=H^1(\Ga;\Mn)_\m$, $k=n+2$ running through the
class.
\begin{enumerate}
\item[(A)] \emph{(The standard family, $g=1$; the v1.17.0 statement.)} If $w=w_O$ and $g=1$ (exactly one cusp form of weight $w_K=\ell+1-w$ congruent to
  $E_{w_K}$ modulo~$\ell$; this holds at $154$ of the $155$ irregular pairs with $\ell<2000$ and fails at $(547,486)$, slot-17 R17-1~\S4; hypothesis added at
  v1.17.0-pre.3, the slot-17 referee F1's B-2), then $A\cong\Fl[x]/x^4$ for every $k$ in the class of $\m$ with
  \[
    \ell w+t(\ell+1)\;\le\;k\;<\;w+t(\ell+1)+(\ell^{2}-1),
  \]
  and the Loewy length of $A$ rises from at most $2$ to $4$ at $k=\ell w+t(\ell+1)$.
\item[(B)] \emph{(The swapped family, $g=1$; new at v1.18.0-pre.1.)} If $w=w_K$ and $g=1$, the same: $A\cong\Fl[x]/x^4$ for every $k$ in the class with
  $\ell w_K+t(\ell+1)\le k<w_K+t(\ell+1)+(\ell^{2}-1)$, and the Loewy length of $A$ rises from at most $2$ to $4$ at $k=\ell w_K+t(\ell+1)$ (the \emph{swapped
  meeting}; at the same twist~$t$ it is the cheaper of the two meetings of a pair exactly when $w_K<w_O$, e.g.\ at $103$ and $691$; v1.18.0-pre.2, the referee F1's
m-E7: pre.1 omitted ``at the same twist'').
\item[(C)] \emph{($g=2$, the first window; new at v1.18.0-pre.1.)} If $g=2$, in either family, then $A\cong\Fl[x]/x^3$ for every $k$ in the class with
  $w+t(\ell+1)\le k<\ell w+t(\ell+1)$ (the first window: there $(m_O,m_K)=(1,1)$, $\dim V=1+(1+2g)=6$ and $\rank\T_\m=3$ by Claim~\ref{cl:dimformula}).
\end{enumerate}
In (A)--(C) ``rises from at most~$2$'' and the bounds $\dim A\le\rank\T_\m$ are \PROVED{} by Claim~\ref{cl:dimformula} (Remark~\ref{rmk:dimwindows}); what is
conjectural is that $A$ is curvilinear of the largest length the rank allows, which includes that $\T_\m$ is saturated there (\HA{prop:cohforms}(5),
\HA{rmk:sattor}). \emph{Nothing is asserted}: for $k\ge w+t(\ell+1)+(\ell^{2}-1)$ in (A), (B) (the multiplicities grow a second time there, $\dim V=12$; the
forms-side algebra has Hilbert function $(1,2,2,1)$ at every such cell computed, in both families); in (C) from the meeting $k=\ell w+t(\ell+1)$ on
($\dim V=4+4g=12$, $\rank\T_\m=6$; the meeting is Prediction~\ref{pr:depth}'s case $(m,r_0)=(2,3)$, computed $\F_{547}[x]/x^6$ on the forms side at
$k=34462$, the meeting of the class $a=1$, by the slot-18 referee F1, \texttt{docs/implementation/slot18/wp3/}; v1.18.0-pre.2, pre.1 read ``no datum''); at $g\ge3$ (no such pair is known; $g=1$ at $154$ of the $155$
pairs below $2000$, $g=2$ only at $(547,486)$); for the two end twists $t=0$ (the class of $E_w$: the Eisenstein ideal itself) and $t=J(w)=\ell-w$ (the
residual representation is then $1\oplus\omega^{\ell-w}$, the class of $E_{\ell+1-w}$, the other Eisenstein end; v1.17.0-pre.2 wording, pre.1 read ``the swapped
configuration''); the definition of Kummer--Frobenius keeps its cap $0\le t\le J(w)$ for Prediction~\ref{pr:latticesplit}. Prediction~\ref{pr:depth} states
the curvilinearity in the form of the Hecke algebra $\T_\m/\ell$, for these families and for the irreducible congruence classes.
\emph{What changed at v1.18.0-pre.1, and why.} (B): the swapped family meets at $\ell w_K+t(\ell+1)$, and the slot-17 forms-side data there
(Prediction~\ref{pr:latticesplit}, ``Refuted as stated'', item~2) are curvilinear $\Fl[x]/x^4$ at every onset cell computed; slot-18 WP-2 tested the upper end
(evidence below). (C): at $g=2$ the coefficient $1+2g=5$ enters already in the first window, so Loewy length~$3$ appears without any doubling of Serre
weights; $(608,547)$ is the cohomology datum (Prediction~\ref{pr:latticesplit}, item~1). Nothing is extrapolated beyond a configuration with a datum: (C) is
not asserted at the meeting, where it would read $\Fl[x]/x^6$ (Prediction~\ref{pr:depth} states that case separately; v1.18.0-pre.2: it now has one
forms-side datum, $(34460,547)$, $\F_{547}[x]/x^6$, the referee F1's M-2; (C)'s claim stays the first window). (A)'s text is
unchanged; its hypothesis ``Kummer--Frobenius with seed $w$ and twist $t$, $1\le t\le J(w)-1$'' is now the case $w=w_O$ of Claim~\ref{cl:seedtwist}(2).
\emph{What changed, and why} (v1.17.0-pre.1). (a)~\emph{The upper end.} ``From $k=\ell w+t(\ell+1)$ on'' is withdrawn. On the forms side the unbounded
form is false at $\ell=37$: at $k=1412$ ($t=1$) and $k=1450$ ($t=2$) the algebra generated on the block $W$ of $S_k(\F_{37})$ by the $T_p$, $p\le23$, has
dimension $6=\dim W$ and Hilbert function $(1,2,2,1)$, so it is $\T^{\mathrm f}_\m/37$, and it is not curvilinear (the report's
\texttt{prop:secondwindow}, by script). For the cohomological $A$ itself the point is open: at $(1410,37)$ the report's dimension formula gives
$\dim V=12$ on a torsion-free block, so $\rank\T_\m=6$ (\HA{lem:heckerankone}) and $A$ is a quotient of $\T_\m/37$ (\HA{prop:cohforms}(5)); either $A$ has
the Hilbert function $(1,2,2,1)$ or $\T_\m$ is not saturated there, and it was not computed (the report's open problem \texttt{pb:secondwindow}). The
new upper bound is where the Jordan--H\"older multiplicities of $\Sym^{n}\F_\ell^{2}$ grow a second time: at $\ell=37$, $w=6$ the pair $(m_1,m_2)$ of the
report's \texttt{thm:dimformula} becomes $(3,3)$ at $k=6+38t+1368$ (its \texttt{comp:jh}, exact, $t=1,2,3$; one period $\ell^{2}-1=1368$ in~$n$).
(b)~\emph{The end twists.} At $\ell=37$, for $t=J(6)=31$ the dimension formula gives $\dim V=3$, $7$, $11$ from $k=32$, $1184$, $1400$ --- never the $8$
of the window --- and for $t=0$ the block is the Eisenstein ideal of $E_6$; neither was evidence for the prediction. (c)~\emph{At $\ell\ne37$} the window's
upper end is computed on the multiplicity side and extrapolated only in the shape of the algebra (v1.17.0-pre.3, the referee F1's m-E1; pre.2 read ``an
extrapolation of the $\ell=37$ pattern''): slot-17 R17-1's committed tables (\texttt{docs/implementation/slot17/r171/classes/}) give the multiplicities
$(2,2)$ at $k=\ell w+t(\ell+1)$ and the next event $(3,3)$ at exactly $k=w+t(\ell+1)+\ell^{2}-1$ on $1492$ of $1492$ standard cuspidal classes of its $13$ pairs
(v1.18.0-pre.1 correction: pre.3 through v1.17.0 read ``$2092$ of $2092$'', the referee F1's count; the standard classes of the $13$ pairs number
$\sum(w_K-2)=1492$, recounted by \texttt{docs/implementation/slot18/wp2/window\_events.py}, which also gives $1906$ of $1906$ for the swapped family,
Remark~\ref{rmk:dimwindows});
that $A$ is $\Fl[x]/x^4$ throughout the window, and stops being so at its upper end, is the $\ell=37$ pattern extended; at pre.1 no other irregular prime had a computed
Loewy-$4$ cell, and at $\ell=61\ldots103$ the first window lies beyond $n=700$ (below). (v1.17.0-pre.2) Slot-17 R17-2 tested the window's lower end at $67$,
$59$ and $101$ (below); its upper end was then still untested off~$37$. (v1.17.0-pre.5) On the forms side the window is now tested off~$37$ up to its last
cell: at $(3496,59)$ and $(4498,67)$ ($k=3498$, $4500$, $t=1$; the next cell of the class, $k+\ell-1$, is the upper bound itself) the block has rank~$4$ and
$\T^{\mathrm f}_\m/\ell$ is curvilinear, $\Fl[x]/x^4$ (by script, the $\ell=37$ report's referee F4,
\texttt{docs/implementation/slot17/r173/f4\_edeg\_out.txt}); the cells beyond the upper bound were not computed off~$37$. (d)~\emph{Scope} (v1.17.0-pre.2). The prediction concerns the Kummer--Frobenius (standard)
family only, at $g_{\ell+1-w}=1$. The swapped family and every class at a pair with $g\ge2$ (v1.17.0-pre.3; pre.2 read ``the $g\ge2$ first windows''), also Kummer classes, give curvilinear algebras too (Prediction~\ref{pr:latticesplit}, ``Observed
on the Kummer locus''); they are outside this prediction, and its generalisation to them is a slot-18 item. (v1.18.0-pre.1: that generalisation is now parts~(B)
and~(C); this paragraph is the v1.17.0 scope of part~(A), kept as published.)
\emph{Evidence for the restated form} (all at $\ell=37$, $w=6$, inside the window, $1\le t\le30$): the $24$ onset cells below (cohomology, census and
weight-shift data); on the forms side $\T^{\mathrm f}_\m/37\cong\F_{37}[x]/x^4$ at $k=260,\dots,728$ ($t=1$) and at $k=1376$ ($t=1$, inside the window,
$36$ below its upper end), by script. In the window $\dim V=8$ is proved in the ranges computed (the report's \texttt{cor:onset}, from the Serre-weight
multiplicities), so there $A\cong\F_{37}[x]/x^4$ is equivalent to $\T_\m/37\cong\F_{37}[x]/x^4$ with $\T_\m$ saturated.
(v1.17.0-pre.2) \emph{Off $\ell=37$}, at the window's lower end (slot-17 R17-2, \texttt{docs/implementation/slot17/R17-2\_report.md}~\S2, \S6.2): on cohomology
at $(736,67)$ ($67\mid B_{58}$, $w=10$, $t=1$, $k=738$): $\dim V=8$, module Loewy layers $[3,2,2,1]$, $A\cong\F_{67}[x]/x^4$ with $U_{67}$ in the operator set,
$\T_\m$ saturated ($I_\m=0$) and the \texttt{stable} verdict certified at the system $(21,25,45,40,0)$ (\texttt{HeckeSaturationDefectValidation}); on the forms
side, by script, $\T^{\mathrm f}_\m/\ell\cong\Fl[x]/x^4$ ($\dim A=\dim W=4$) at $(736,67)$, $(804,67)$ ($t=2$), $(1002,59)$ ($59\mid B_{44}$, $w=16$, $t=1$),
$(1062,59)$ ($t=2$) and $(3534,101)$ ($101\mid B_{68}$, $w=34$, $t=1$). Each was predicted before it was run (R17-1~\S7); none failed.
(v1.18.0-pre.1) \emph{The C-F9 lattice at the window's cells, $\ell=37$} (slot-17 \texttt{docs/implementation/slot17/LAT37\_local\_report.md}, the env-armed
\texttt{HeckeCensusLatticeValidation} on the October records, $U_{37}$ in the operator set; \texttt{refused:0}, \texttt{fallback:0}): at $(700,37)$, $(698,37)$
and $(696,37)$ exactly one ideal per cell fails the lattice test, the Kummer--Frobenius ideal of dimension~$8$ and Loewy length~$4$ (\texttt{first\_fail=2,2},
\texttt{cf9\_lattice = cf9\_field = false}, module layers of the $(656,37)$ shape), and every ideal of dimension~$4$ holds, as at $(654,37)$ and $(656,37)$;
$(694,37)$, $(692,37)$, $(690,37)$ run in slot~18 (WP-L37). (At $(696,37)$ the validation's own (J1)/(J2) self-checks failed for an environmental reason, a
removed build tree; the lattice rows are unaffected, report \S1.)

\emph{Evidence for part~(B)} (v1.18.0-pre.1; forms side only, $\T^{\mathrm f}_\m/\ell$ on the block $W\subseteq S_k(\Fl)$ with the generation certificate
$\dim A=\dim W$; nothing here is cohomology). At the swapped meeting: $\F_{37}[x]/x^4$ at $b=1,\dots,4$ ($k=1222$, $1260$, $1298$, $1336$), $\F_{103}[x]/x^4$ at
$b=1,2$ ($k=2576$, $2680$), $\F_{691}[x]/x^4$ at $b=1,2$ ($k=8984$, $9676$) (slot-17 R17-2~\S\S2--3). At the window's last cell and one cell past it
(slot-18 WP-2, \texttt{docs/implementation/slot18/wp2/}, by script, \texttt{r172\_forms.py} unchanged; the predictions committed before the run, \texttt{predictions.txt};
$11/11$ held): $\Fl[x]/x^4$ at $(1400,37)$ ($b=1$), $(1438,37)$ ($b=2$) and $(10632,103)$ ($b=1$), each the last cell below
$w_K+t(\ell+1)+\ell^{2}-1$; and $\dim W=6$, Hilbert function $(1,2,2,1)$, not curvilinear, at the next cell $(1436,37)$, $(1474,37)$, $(10734,103)$, exactly as
in the standard family at $k=1412$, $1450$. (v1.18.0-pre.2) The slot-18 referee F1 (\texttt{docs/implementation/slot18/wp3/}, the same script, predictions
written before each run, \texttt{cells\_f1.txt}; $5/5$ held here): $\Fl[x]/x^4$ at $(1364,37)$ ($b=1$, mid-window), $(1476,37)$ ($b=3$, last cell),
$(1514,37)$ ($b=4$, last cell) and $(10736,103)$ ($b=2$, last cell); $\dim W=6$, $(1,2,2,1)$, not curvilinear, at $(1512,37)$ ($b=3$, first cell past the window).
On the multiplicity side the window of~(B) is computed on all $1906$ swapped classes of R17-1's $13$ pairs
(Remark~\ref{rmk:dimwindows}). The census algebra $A$ was not computed at any swapped cell: the cheapest, $(1220,37)$ (the $\ell=37$ report's \texttt{pb:swappedcoh}), needs an
exclusive memory window (slot-18 R18-3, $\approx20$~GB); $(2574,103)$ and $(8982,691)$ are beyond local reach. So (B) is \CONJ{} with the saturation of
$\T_\m$ at these blocks as its untested half.
\emph{Evidence for part~(C)} (v1.18.0-pre.1). On cohomology: $(608,547)$, the class $a=1$ of the standard family at the first cell of its first window,
$A\cong\F_{547}[x]/x^3$, $I_\m=0$, \texttt{stable} certified (Prediction~\ref{pr:latticesplit}, item~1). On the forms side, $\F_{547}[x]/x^3$ ($\dim W=3$,
generated): $(1156,547)$ ($a=2$, R17-2), and (slot-18 WP-2, predicted before the run) $(1154,547)$, $(1700,547)$, $(2246,547)$ (the second to fourth cells of the
class $a=1$) and, in the \emph{swapped} family, $(1032,547)$ and $(1578,547)$ (the first two cells of $b=1$; its first window runs from $k=1034$ to the swapped
meeting $k=266390$). (v1.18.0-pre.2, the referee F1's M-2; \texttt{docs/implementation/slot18/wp3/}, predicted before the run, $3/3$ held here; operators
$T_2,\dots,T_{11}$, no $T_\ell$): $\F_{547}[x]/x^3$, generated, at $(16988,547)$ (the $31$st cell of $a=1$), at $(33914,547)$ (its $62$nd and \emph{last} cell,
$k=33916$) and, swapped, at $(11952,547)$ (the $21$st cell of $b=1$). The first window of $a=1$ ($62$ cells, $610\le k<34462$) is now sampled at its cells
$1$--$4$, $31$ and $62$, so (C) is stated over the whole first window on evidence at its start, middle and end (pre.1: ``the data cover their first two to four
cells''); the swapped window of $b=1$ (from $k=1034$ to $266390$) is sampled at its cells $1$, $2$ and $21$ only. If a cell inside a first window at
$547$ read a non-curvilinear algebra, (C) would be restricted to what was computed. At the meeting $(34460,547)$, $k=34462$, the same run gives
$\dim W=6=\rank\T_\m$, Hilbert function $(1,1,1,1,1,1)$, generated: $\F_{547}[x]/x^6$ --- the datum of Prediction~\ref{pr:depth} at $(m,r_0)=(2,3)$, outside~(C).

\medskip\noindent\textbf{The v1.15.1-pre.7 statement (superseded at v1.17.0-pre.1 by the window form above; kept as published, with its evidence).}
On a Kummer--Frobenius $\m$ (Prediction~\ref{pr:latticesplit}), $A\cong\Fl[x]/x^4$ from $k=\ell w+t(\ell+1)$ on: the Loewy length of
$A$ rises from $2$ to $4$ at exactly that weight.
\emph{Evidence} (computed, not certified as a law): the weight-shift data predict Loewy length~$4$ at $\ell=37$ for
$\operatorname{eis}(t,t+5)$ exactly at $n=220+38t+36m$, $m\ge0$, $t\ge1$. In the scanned range $n\le356$ that gives the six cells
$(258,37)$, $(294,37)$, $(330,37)$, $(296,37)$, $(332,37)$, $(334,37)$; all six are observed, and no other cell has Loewy length
$\ge3$. The cluster window $n=648\ldots700$ gives $t=1,\dots,12$ and $t=1,\dots,6$, the $18$ records of
Prediction~\ref{pr:latticesplit}. So the onset has $24/24$ hits and $0$ false positives; it was predicted (WP-12 pass~1, \S1.2:
$(258,37)$, $(294,37)$ as ``the cheap test'') before the weight-shift scan was read.
(v1.16.1) \emph{At $\ell=61\ldots103$} (census edition~3's complete C1b block, slot-15 WP-C18~\S3.2): the first Loewy-$4$ cell of
every Kummer class there lies beyond $n=700$ ($n=736$ at $67\mid B_{58}$, $n=3534$ at $101\mid B_{68}$, $n=8342$ at
$103\mid B_{24}$), so the law predicts no Loewy-$4$ cell and none is observed; the $229$ Kummer-predicted rows lie below their
onset and have Loewy $\le2$. This is consistent with the law but \textbf{does not test its positive half}; the first cell that would
is $(736,67)$, beyond the C1 grid. (v1.17.0-pre.2, two corrections.) $(736,67)$ was tested in slot~17 and held (the evidence of the restated form, above).
And $n=8342$ is the first \emph{Kummer--Frobenius} Loewy-$4$ candidate at $103$, not the first Loewy-$4$ cell there: the swapped $(2574,103)$ comes first
(slot-17 R17-1~\S5.2) and reads $\F_{103}[x]/x^4$ on the forms side (R17-2~\S2), outside this prediction and a counterexample to the v1.16.0 form of
Prediction~\ref{pr:latticesplit}.
\emph{Mechanism} (\CONJ, WP-12~\S1.2): at $k=\ell w+t(\ell+1)$ the Frobenius direction of the seed $E_w$ enters the block. In the
split case it adds a line (module Loewy $[3,1]$, still $\m_A^2=0$); in the Kummer case it doubles the Ribet-congruent block to a
curvilinear $\Fl[x]/x^4$. Why the doubling is curvilinear is open: Bella\"iche--Khare's $A_{\bar\rho}\cong\Fl[[x,y]]$
(unobstructed, Vandiver; \cite{bk}) bounds only the embedding dimension, by~$2$. (v1.17.0-pre.1) At the window's upper end the multiplicities grow a
second time and the block's rank rises from $4$ to $6$; on the forms side the algebra there has embedding dimension~$2$ (at $\ell=37$, $t=1,2$; the
report's \texttt{prop:secondwindow}). Why the window closes exactly
there, and whether the census algebra follows the forms algebra beyond it, is open.
\end{prediction}

\begin{remark}[The split\,/\,ramified switch is not obstructedness]\label{rmk:i8switch}
\PROVED{} (from \cite{bk}; v1.16.3, owner ruling~O4 of 2026-10-01, slot-15 WP-12 proposal~(E-3) of
\texttt{docs/implementation/slot15/WP-12\_U\_E\_proposals.md}; the reading is \texttt{docs/implementation/slot15/WP-12\_report.md}~\S5
and~\S8.4(e).)
Slot-14 WP-C14's I-8 dichotomy for a twisted Eisenstein ideal $\operatorname{eis}(i,j)$, $0<i<j$, with Bernoulli seeds $B_{w_1}$,
$B_{w_2}$, $w_1=j-i+1$, $w_2=\ell+i-j$ (the \emph{split} case: both seeds are $\ell$-units, two $\ell$-adically rational eigenforms,
module Loewy $[3,1]$ from $k_2$ on, \S\ref{sec:census}, data label \texttt{twisted\_eis\_split}; the \emph{ramified} case: $\ell$
divides a seed, one ramified Galois orbit --- the Kummer case, compare Prediction~\ref{pr:kfonset}) is \textbf{not} Bella\"iche--Khare obstructedness (\cite{bk} Thm.~22~(i), (iii)). For
$\bar\rho=\omega^i\oplus\omega^j$, $a=j-i$, one has $a+1=w_1$ and $\ell-a=w_2$, so Thm.~22~(i) is exactly ``both seeds are units''; but
Thm.~22~(iii) makes $\bar\rho$ unobstructed whenever Vandiver's conjecture holds at~$\ell$, the Kummer case included (the Kummer class lies
in the odd eigenspace, \cite{ribet}, while the obstruction lives in even ones), and Vandiver is verified for every $\ell<2^{31}$
(Hart--Harvey--Ong \cite{hho}). So both cases are unobstructed, $A_{\bar\rho}\cong\Fl[[x,y]]$ in both, and the dichotomy is a statement
about $\T_{\bar\rho}$ in characteristic~$0$ (one ramified orbit against two rational forms), not about $A_{\bar\rho}$.
\end{remark}

\begin{prediction}[The Frobenius layer (F1)]\label{pr:frobeniuslayer}\CONJ\ (v1.16.3, owner ruling~O4 of 2026-10-01: slot-15 WP-12
proposal~(E-2) of \texttt{docs/implementation/slot15/WP-12\_U\_E\_proposals.md}, accepted as a \CONJ; the law~L-F1 of
\texttt{docs/implementation/slot15/WP-12\_report.md}~\S2.2)
Let $\ell\ge5$ and let $f$ be an ordinary level-one eigenform of weight $k_0\le\ell$ (v1.20.0-pre.7: $k_0\le\ell+1$ through pre.6; the boundary weight $k_0=\ell+1$ is excluded after slot-20 WP-SW, below), $\bar\rho_f$ its residual representation. Let
$J(k_0)$ be the least $J\ge0$ with $k_0+J\equiv0\pmod\ell$, so that $\theta^jf$ keeps filtration $k_0+j(\ell+1)$ for $j\le J(k_0)$. For
$1\le j\le J(k_0)$, the block of $\bar\rho_f\otimes\omega^j$ in $H^1(\Ga;\Sym^n\Fl^2)$ goes from $D=2$ to $D=4$ at
\[
  n_V=\ell k_0+j(\ell+1)-2,
\]
and it is \emph{hidden} ($A_1=\Fl$, $A_K/\ell\cong\Fl[x]/x^2$) from $n_V$ on.

\emph{Why} (\PROVED{} given two per-cell facts --- $D(n_V)=4$, and the two blocks at $n_0$ are reduced; WP-12~\S2.2).
(i)~At $n_0=\ell k_0-2$ the $\bar\rho_f$-part of $H^1$ is $B_{\mathrm{ord}}\oplus B_V$, separated by $U_\ell$ ($a_\ell\ne0$ on
$B_{\mathrm{ord}}$, $a_\ell=0$ on $B_V$, the $V(f)=f^\ell$ direction); each has $D=2$ and reduced algebra $\Fl$ (the per-cell fact).
(ii)~$\theta_L^j$ maps $H^1(\Ga;M_{n_0})$ to $H^1(\Ga;M_{n_V})$ with $T_p\theta_L=p\,\theta_LT_p$ and $U_\ell\theta_L=0$
(\HA{thm:thetales}(2)), injectively off the Eisenstein line (\HA{thm:thetales}(3)); so $\theta_L^j(B_{\mathrm{ord}}\oplus B_V)$ is a
$4$-dimensional Hecke submodule of the $\bar\rho_f\otimes\omega^j$ block on which every $T_p$ acts as the scalar $p^ja_p(f)$ and $U_\ell$
as~$0$. (iii)~$D(n_V)=4$ (the per-cell fact) makes it the whole block; hence $A_1=\Fl$. (iv)~The block is torsion-free (non-Dickson), so
\HA{lem:heckerankone} gives a local $\Z_\ell$-algebra of rank $\lceil4/2\rceil=2$, whose reduction is local of dimension~$2$, i.e.\
$\Fl[x]/x^2$; by \HA{cor:stablecert} this is $A_K/\ell$. (v)~So $A_1\subsetneq A_K/\ell$, a saturation defect
(\HA{prop:saturation}(2)): the congruence is hidden.
\emph{What stays conjectural}: the onset --- that nothing enters the block earlier, so that the doubling happens exactly at $n_V$ --- and
that the $V$-block first appears at weight exactly $\ell k_0$ with $V(f)=f^\ell$ (compare \HA{rmk:thetaforms}(1)). Both held on every
tested cell. \emph{Why $j\le J(k_0)$ only}: for $j>J(k_0)$, $\theta^jf$ has dropped filtration, appears early, and at $k_0+j(\ell+1)$
the block gains $\theta_L^j(B_{\mathrm{ord}})$ glued non-semisimply (not hidden); identifying the dropped form with a companion-form
twist needs Gross's companion-form theorem and is not attempted.

\emph{Evidence} (computed, WP-12~\S2.3, \texttt{law\_f1.py} on the weight-shift data of WP-L0930): at $\ell=17$, $33$ irreducible chains
with the first doubling as predicted and $0$ failures, $17$ first hidden rungs as predicted (the three $\theta^j(\Delta E_6)$ at
$n=322,340,358$ included) and $0$ failures; at $\ell=19$, $49$ and $9$ (plus $n=362,366$ beyond $360$), $0$ failures; at
$\ell=29,37,47$ the doublings $14,15,2$ as predicted and no hidden rung predicted or observed in range; at $\ell=59\ldots199$ none
predicted, none observed; $0$ spurious hidden blocks at every $\ell$. The predictions of WP-12~\S2.4 ($\ell=17$, $n=362\ldots500$;
$\ell=19$; at $\ell=37$ no hidden block below $n=480$) were written before any further data and are untested.
\emph{Evidence (slot~16, 2026-10-01)}: consistent, no qualification of the statement. Slot-16 WP-5
(\texttt{HeckeSaturationDefectValidation}, \texttt{docs/implementation/slot16/WP-5\_report.md}~\S3) read the hidden-chain cells $(220,17)$
($\theta\Delta$), $(238,17)$ ($\theta^2\Delta$), $(246,19)$ ($\theta\Delta$): $D=4$, $\operatorname{rank}\T=2$, $A_1=\Fl$,
$\T/\ell=\Fl[x]/x^2$, $\T^{\mathrm{sat}}/\T=\Z/\ell$, the defect ideal $I=\m$ principal, $3/3$ held; its default arm~(H) gates the
law's $j=1$ onset at $\ell=13$, the cell $(168,13)$ (slot-16 WP-3/WP-5 audit). Slot-16 WP-11
(\texttt{docs/implementation/slot16/WP-11\_report.md}) read $(220,17)$ and $(246,19)$ with $a_\ell=0$, $\dim A_1=1$ and the \texttt{stable}
verdict certified ($\T_\m/\ell$ of Hilbert function $[1,1]$), and records that there $\Lambda_\m\not\cong\T_\m^2$: a hidden block is free
of rank~$4$ over a rank-$2$ algebra with a saturation defect --- a fact beside this prediction, not against it.
\emph{Evidence (slot~19, 2026-10-04)}: consistent, no qualification. At $(1220,37)$ (slot-19 WP-RC's \texttt{HeckeSaturationDefectValidation} run, read by slot-19 sat1220,
\texttt{docs/implementation/slot19/sat1220\_analysis.md}, predictions committed first) all nine $\F_{37}$-rational irreducible rungs $n_V\le1220$ of the weight class are hidden
($D=4$, $A_1=\F_{37}$, $I=\m$) and there is no spurious hidden block: the first test at $\ell=37$ beyond WP-12's range, $0$ failures (\CMPD; two rungs at $n_V=1184$, two at
$n_V=1220$). The tenth block with $\dim I=1$ there is the rung of a Kummer class, Remark~\ref{rmk:frobkummer}.

\emph{Evidence (slot~20, forms side; v1.20.0-pre.3, scope qualified at v1.20.0-pre.4)}: consistent with the transport step~(ii) of \emph{Why}; \textbf{outside the
statement's hypotheses} --- the family-(c) seed is a non-reduced block of two congruent eigenforms ($\T/\ell=\Fl[\epsilon]/\epsilon^{2}$), not a reduced $D=2$ pair, and
the onset block is predicted saturated, not hidden (\Lcite{pb:irrcoh}). At the family-(c) onsets $(2210,67)$, $(2278,67)$, $(3798,131)$, $(3930,131)$
($j=1,2,1,2$) the $\theta^{j}$-image of the seed space $S_{k_0}(\Fl)$ is $\ker x^{2}$ of the curvilinear block, so the first-order Hecke data of the doubled block is the
$\theta^{j}$-twist of the seed's $\mathcal O_K/\m^{2}$: $T_p\mapsto p^{j}(a_p\bmod\m^{2})$, $t_1(p)=2^{-j}p^{j}s(p)$ (\PROVED{} given the curvilinearity, \CMPD{} at all $1228$
primes $p<10^{4}$ at each cell; slot-20 R20-2~\S3, \texttt{docs/implementation/slot20/R20-2\_frobenius\_shadows.md}; \Lcite{sec:fieldsirr}). Neither the onset nor the hidden conclusion is tested by this; the first-order data is non-scalar, as a non-reduced source demands.
\emph{Evidence and an observation (slot~20, cohomology; v1.20.0-pre.6)}: slot-20 WP-SW's three cells $(1258,37)$, $(1296,37)$, $(1334,37)$, read with slot-19 sat1220's
\texttt{lf1\_rungs.py}: at each, the $9$ irreducible rungs with $n_V\le n$ are present and hidden ($D=4$, $A_1=\F_{37}$), among them the onset $n_V=1334$ of $k_0=32$, $j=4$ exactly at $(1334,37)$, and
no predicted rung is missing (\CMPD). One hidden $D=4$ block at $(1334,37)$ is not a predicted rung: $\bar\rho_f\otimes\omega$, $f$ the weight-$38$ ($=\ell+1$) form with
$(a_2,a_3,a_5,a_7,a_{37})\equiv(0,1,0,36,1)$, $J(38)=36$, predicted onset $n_V=1442$. Both reduced $D=2$ copies of $\bar\rho_f$ ($a_{37}=1$ and $a_{37}=0$) are present already at $n=1296$ (in the role of \emph{Why}~(i)'s $n_0=\ell k_0-2=1404$), and
$1296+(\ell+1)=1334$: \emph{Why} (ii)--(v) with $n=1296$ in the role of $n_0$ explains the block (\PROVED{} given the per-cell facts read in the logs). So at the boundary weight $k_0=\ell+1$ the statement as written through v1.20.0-pre.6 ($k_0\le\ell+1$) is \textbf{refuted at this one class}: the
doubling is not at $n_V=1442$ but at $1334$, because the $V$-copy of $\bar\rho_f$ is present at weight $1298=\ell k_0-3(\ell-1)$, not first at
$\ell k_0=1406$; the mechanism (ii)--(v) holds. The hypothesis is therefore restricted to $k_0\le\ell$ (v1.20.0-pre.7, F1$'$ M-E2), where every tested cell is
consistent; the onset at $k_0=\ell+1$ is open (why the $a_\ell=0$ copy appears early is not examined; \Lcite{app:ribetdata}). Tag unchanged, \CONJ.
\end{prediction}

\begin{remark}[The Frobenius layer on a Kummer class (v1.19.0-pre.5)]\label{rmk:frobkummer}
\CMPD{} at one cell, \CONJ{} in general (slot-19 sat1220, \texttt{docs/implementation/slot19/sat1220\_analysis.md}~\S\S1--3). At the rung $n_V=\ell w_K+j(\ell+1)-2$ of a Kummer class
($\ell\mid B_{w_K}$, Ribet's form ordinary; $k_0=w_K$ in Prediction~\ref{pr:frobeniuslayer}) the block is the swapped family's $D=8$ block (the swapped onset
$k=\ell w_K+b(\ell+1)$ of Prediction~\ref{pr:kfonset}(B) \emph{is} this rung, $j=b$), and the $\theta$-layer hides not all of $\m$ but $\m^{3}$: $\T_\m/\ell\cong\Fl[x]/x^{4}$,
$A_1\cong\Fl[x]/x^{3}$, $\T^{\mathrm{sat}}_\m/\T_\m\cong\Z/\ell$, $I_\m=\m^{3}$. \emph{Computed} at $(1220,37)=(37\cdot32+38-2,37)$, the rung of $\theta^{1}$ on the weight-$32$ Kummer class,
system $(9,19,14,16,0)$: torsion-free, \texttt{stable} certified at $K=2$, $\T_\m=\Z_{37}[T_p:p\in\{2,3,5,7,37\}]$ the full Hecke algebra of the block (the $S$-primes generate
$\T^{\mathrm f}_\m/37$, Nakayama; so no other generating set changes $I_\m$), $\T_\m/37\cong\F_{37}[x]/x^{4}$ on the forms side, $A_1\cong\F_{37}[x]/x^{3}$, radical layers $[4,3,1]$,
$T_2-9$ of Jordan type $[3,2,2,1]$; an independent route on $H^{1}(\Ga;\Sym^{1220}\F_{37})$ agrees. At v1.19.0-pre.5 not yet verified: that $\theta V_{\mathrm{src}}$ is the socle of $V$. \emph{Computed at v1.20.0-pre.1} (\CMPD{} as a run, one route; slot-20 WP-RC20 window TS,
\texttt{docs/implementation/slot20/TS\_launch.md}, the console script \texttt{theta\_socle\_1220\_37.oc} of sat1220, $V_{\mathrm{src}}$ the block of $(1182,37)$ at $a_p=1+p^{31}$): $\dim V=8$,
$\dim V_{\mathrm{src}}=7$, $\theta$ is injective on $V_{\mathrm{src}}$, $\theta V_{\mathrm{src}}=\ker(U_{37}|V)$ of dimension~$7$ ($\mathrm{rank}\,U_{37}|V=1$), $\theta V_{\mathrm{src}}$ is \emph{not} killed by
$\m$ ($T_p-a_p$ has rank~$3$ on it for $p=2,3,5,7$), and $\dim\mathrm{soc}(V)=4$: \textbf{$\theta V_{\mathrm{src}}$ is not the socle of $V$}, so the identification of the defect's
$\theta$-layer with the socle (sat1220~\S3's reading) is refuted. The facts $\dim I_\m=1$, $I_\m=\m^{3}$, $A_1\cong\F_{37}[x]/x^{3}$ above are not touched; a mechanism is open. This is consistent with Prediction~\ref{pr:frobeniuslayer} ``Why (ii)'' and with \HA{thm:thetales}(2),(3),(5) (v1.20.0-pre.2, F1 m-E6): $\theta_L$ is injective
on $H^{1}(M_{1182})$ ($\mathrm{rank}\,\delta=0$) and $a_p=1+p^{31}\ne1+p^{-1}$; $U_\ell\theta_L=0$ gives $\theta V_{\mathrm{src}}\subseteq\ker U_{37}$, equality by dimension; $T_p\theta_L=p\,\theta_LT_p$
carries the non-semisimple part of $V_{\mathrm{src}}$ ($T_p-a_p$ of rank~$3$) into $V$. (ii)'s scalar $T_p$ rests on \emph{reduced} $D=2$ source blocks, which the $7$-dimensional
$V_{\mathrm{src}}$ of a reducible $\bar\rho$ is not; the prediction's hypotheses do not hold at this cell, and no statement of it moves. \emph{Conjectured}: the same ($\dim I_\m=1$, $I_\m=\m^{3}$, $A_1\cong\Fl[x]/x^{3}$) at every swapped cell $n=\ell w_K+b(\ell+1)-2$, $1\le b\le J(w_K)$ (the cells run or planned: at $37$, $b=1,\dots,4$ of $J(32)=5$, the fifth $n=1372$ predicted
\texttt{a20b9531} and running; at $103$, $b=1,2$ of $J(24)=79$, $(2574,103)$, $(2678,103)$; at $691$, $b=1,2$ of $J(12)=679$, $(8982,691)$, $(9674,691)$); slot-20 WP-SW runs $b=2,3,4$ at $37$ ($n=1258$, $1296$, $1334$; predictions committed first), the two $103$ cells go to a cluster kit
and the $691$ cells are recorded as untested (owner ruling O10 of 2026-10-04). \emph{v1.20.0-pre.6}: WP-SW \textbf{confirms} it at $b=2,3,4$ (\CMPD; $D=8$, \texttt{stable} certified at $K=2$,
$\T/37=\F_{37}[x]/x^{4}$, $A_1=\F_{37}[x]/x^{3}$, $\dim I_\m=1$, $I_\m=\m^{3}$ at each; \Lcite{conj:swappedsat}), so at $37$ at the first four of the five cells of the range; \CONJ{} at $b=5$ and beyond $37$. \emph{1.20.0}: \textbf{confirmed at $b=5$} as well, at the boundary cell
$(1372,37)$ ($b=J(32)$; \CMPD, \texttt{docs/implementation/slot20/SW1372\_report.md}, prediction \texttt{a20b9531} first): the block of system $(33,22,18,10,0)$ has $D=8$, \texttt{stable} certified
at $K=2$, $\T/37=\F_{37}[x]/x^{4}$, $A_1=\F_{37}[x]/x^{3}$, $\dim I_\m=1$, $I_\m=\m^{3}$; so at $37$ the conjecture holds at all five cells of its range ($b=1,\dots,5$), and it stays \CONJ{} beyond $37$. Prediction~\ref{pr:latticesplit}~(R) is not touched: the class is not congruence-free.
\end{remark}

\subsection{The dimension of a reducible block at every \texorpdfstring{$\ell$}{l}}\label{sec:dimblock}
(New at v1.18.0-pre.1, hecke-extensions slot~18, WP-2. Slot-17 R17-1's Lemmas~3.1, 3.2 and Theorem~4.1,
\texttt{docs/implementation/slot17/R17-1\_report.md}~\S\S3--4, become statements of this document, as the slot-17 referee F1 asked (m-E5, ``since [E] now
argues from Theorem~4.1''); the $\ell=37$ report's proposal \texttt{conj:depth} (\texttt{docs/latex/l37\_loewy4\_example.tex} 0.2.0-pre.2, not binding) becomes
Prediction~\ref{pr:depth}. Placed at the end of \S\ref{sec:ext}, so that no earlier statement number moves.) Throughout $\ell\ge5$. A level-$1$ residual
representation $\omega^{i}\oplus\omega^{j}$ has the \emph{seeds} $w_1=d+1$, $w_2=\ell-d$ of Prediction~\ref{pr:latticesplit} ($d\in[1,\ell-2]$, $d\equiv j-i\pmod{\ell-1}$,
odd); it is of Dickson type iff a seed equals~$2$, and otherwise both seeds lie in $[4,\ell-3]$. For an even $w\ge4$, $J(w)\in[0,\ell-1]$ is defined by
$J(w)\equiv-w\pmod\ell$, as in Prediction~\ref{pr:latticesplit}.

\begin{claim}[Two seeds per class; the standard and the swapped family]\label{cl:seedtwist}\PROVED\
\begin{enumerate}
\item Let $4\le w,w'\le\ell-3$ be even. Then $1\oplus\omega^{w-1}$ and $1\oplus\omega^{w'-1}$ agree up to a twist by a power of~$\omega$ iff $w'\in\{w,\ell+1-w\}$.
\item Let $(\ell,w_K)$ be an \emph{irregular pair} ($\ell\mid B_{w_K}$, $w_K$ even, $4\le w_K\le\ell-3$) and $w_O:=\ell+1-w_K$; then $J(w_O)=w_K-1$ and
  $J(w_K)=w_O-1$. The classes with seeds $\{w_O,w_K\}$ --- the \emph{Kummer classes} of the pair --- are $\omega^{a}(1\oplus\omega^{w_O-1})=\omega^{b}(1\oplus\omega^{w_K-1})$,
  $a\in\Z/(\ell-1)$, $b\equiv a+w_O-1$, and each is exactly one of: the class $a=0$ (that of $E_{w_O}$); the class $b=0$ (that of $E_{w_K}$); a class of the
  \emph{standard family}, $1\le a\le w_K-2$; a class of the \emph{swapped family}, $1\le b\le w_O-2$. Equivalently, a Kummer class other than the two Eisenstein
  ends has exactly one presentation $\omega^{t}(1\oplus\omega^{w-1})$ with $w\in\{w_O,w_K\}$ and $1\le t\le J(w)-1$: with $w=w_O$, $t=a$ in the standard family
  (where $\ell$ divides the companion $B_{\ell+1-w}$, the Kummer--Frobenius situation of Prediction~\ref{pr:latticesplit}'s v1.16.0 form) and with $w=w_K$,
  $t=b$ in the swapped family (where $\ell$ divides $B_w$ itself).
\item A class is a Kummer class of two different irregular pairs only if $\ell\mid B_w$ and $\ell\mid B_{\ell+1-w}$ for one of its seeds~$w$ (a \emph{double
  divisibility}).
\end{enumerate}
\end{claim}
\begin{proof}
(1) Characters of the cyclic group $(\Z/\ell)^{\times}$ are determined by their exponents modulo $\ell-1$, so $\omega^{t}\oplus\omega^{t+w'-1}\cong1\oplus\omega^{w-1}$
iff $\{t,t+w'-1\}\equiv\{0,w-1\}$: either $t\equiv0$ and $w'\equiv w$, or $t\equiv w-1$ and $w'-1\equiv-(w-1)$, i.e.\ $w'\equiv\ell+1-w$; as $w,w',\ell+1-w$ all lie
in $[4,\ell-3]$, the congruences are equalities. (2) $J(w_O)=\ell-w_O=w_K-1$, $J(w_K)=w_O-1$. Take $a\in[0,\ell-2]$: for $1\le a\le w_K-2$, $b=a+w_O-1\in[w_O,\ell-2]$;
for $a=w_K-1$, $b=\ell-1\equiv0$; for $w_K\le a\le\ell-2$, $b\equiv a+w_O-\ell\in[1,w_O-2]$. So the four sets partition $\Z/(\ell-1)$, and by~(1) the presentations
of a Kummer class with a seed in $[4,\ell-3]$ are exactly the two given. (3) By~(1) the seeds of a class determine the pair up to $w\leftrightarrow\ell+1-w$.
\end{proof}
\noindent By script (slot-17 R17-1~\S3, three independent routes to $B_w\bmod\ell$): no double divisibility for $\ell<2\cdot10^{4}$ ($900$ irregular primes,
$1153$ pairs); so at every reachable cell each Kummer class belongs to one pair. A heuristic (not a statement) expects the first one far beyond.

\begin{claim}[The self-mirrored seed is regular]\label{cl:selfmirror}\PROVED\ (given the \STD\ input~(i) of the proof)
If $\ell\equiv3\pmod4$ and $\ell>3$, then $\ell\nmid B_{(\ell+1)/2}$: the seed $w=(\ell+1)/2$, the one with $w=\ell+1-w$, is never irregular.
\end{claim}
\begin{proof}
Let $\chi=(\cdot/\ell)=\omega^{(\ell-1)/2}$, odd, and $m:=(\ell+1)/2$, even, $4\le m\le\ell-3$ for $\ell\ge7$ (at $\ell=7$, $m=4$). (i) \STD: the class number
formula $h(-\ell)=-B_{1,\chi}=-\frac1\ell\sum_{a=1}^{\ell-1}a\,\chi(a)$ for $\ell>3$ (Washington \cite[Ch.~4]{washington}; the theorem number is \textbf{not
verified}, rule~A2). (ii) $B_{1,\omega^{m-1}}\equiv B_m/m\pmod\ell$, \cite[Cor.~5.15, p.~61]{washington1st} (verified on the page image, slot-17 R17-3; the
$\ell=37$ report's \texttt{app:verify}); here $\omega^{m-1}=\chi$, so $h(-\ell)\equiv-B_m/m\pmod\ell$. (iii) $0<h(-\ell)<\ell$: pairing $a$ with $\ell-a$
($\chi(\ell-a)=-\chi(a)$), $\sum_{a=1}^{\ell-1}a\chi(a)=\sum_{a=1}^{(\ell-1)/2}\chi(a)(2a-\ell)$, of absolute value at most $\sum_{a\le(\ell-1)/2}(\ell-2a)=(\ell-1)^{2}/4$,
so $h(-\ell)\le(\ell-1)^{2}/(4\ell)<\ell$. Hence $\ell\nmid B_m/m$, and $\ell\nmid m$, so $\ell\nmid B_m$.
\end{proof}
\noindent By script (R17-1~\S3): the congruence and the bound on all $154$ primes $7\le\ell<2000$ with $\ell\equiv3\pmod4$.

\begin{claim}[The dimension of a reducible block, at every $\ell$]\label{cl:dimformula}\PROVED\ (given the \STD\ inputs named in the proof; two of their
locators are \textbf{not verified}, see the status paragraph below)
Let $\ell\ge5$, $n\ge2$ even, $k=n+2$, and $\m$ a maximal ideal of $\T$ on $H^1(\Ga;\Mn)$ with residue field~$\Fl$ whose residual representation
$\bar\rho_\m\cong\omega^{i}\oplus\omega^{j}$ is reducible and not of Dickson type, with seeds $w_1=d+1$, $w_2=\ell-d$. For an even $w$, $4\le w\le\ell-3$, let
$g_w$ be the dimension of the generalised eigenspace of the system $a_p=1+p^{w-1}$ ($p\ne\ell$) in $S_w(\Fl)$; $g_w=0$ unless $\ell\mid B_w$. Then
\[
  \dim_{\Fl}H^1(\Ga;\Mn)_\m\;=\;m_1(n)\,(1+2g_{w_1})\;+\;m_2(n)\,(1+2g_{w_2}),
\]
\[
  m_1(n):=[\Mn:\Sym^{w_1-2}\Fl^{2}\otimes{\det}^{i}],\qquad m_2(n):=[\Mn:\Sym^{w_2-2}\Fl^{2}\otimes{\det}^{i+w_1-1}],
\]
the Jordan--H\"older multiplicities of $\Mn=\Sym^{n}\Fl^{2}$ as a $\GL_2(\Fl)$-module (${\det}^{s}$ acting on $T_p$ through $p^{s}$, the convention of
Claim~\ref{cl:glover}). If moreover $\m$ is \emph{cuspidal} --- $\bar\rho_\m\not\cong1\oplus\omega^{k-1}$, i.e.\ $E_k$'s system is not in~$\m$ --- then the block
$\Lambda_\m=H^1(\Ga;\Sym^{n}\Z_\ell)_\m$ is torsion-free, $\dim H^1(\Ga;\Mn)_\m$ is even, $\rank\T_\m=\frac12\dim H^1(\Ga;\Mn)_\m$, and the census algebra
$A$ of the block has $\operatorname{Loewy}(A)\le\dim A\le\frac12\dim H^1(\Ga;\Mn)_\m$.
\end{claim}
\begin{proof}
(The $\ell=37$ report's \texttt{thm:dimformula}, steps~(1)--(3), with $37$ replaced by $\ell$ and its computer input~(ii) replaced by~(c) below.)
(a) \emph{Additivity.} For $p\ne\ell$ the matrices of determinant $p$ in the double coset of $T_p$ act on $\Mn$ through their reductions in $\GL_2(\Fl)$, so every
$\GL_2(\Fl)$-submodule carries the anemic Hecke action and the long exact sequences of $\Ga$-cohomology of a composition series of $\Mn$ are Hecke-equivariant
(\STD: \cite[Lemma~1.1.1, p.~194]{ash_stevens}, verified on the page image, slot-17 R17-4). $H^{q}(\Ga;N)=0$ for $q\ge2$ and every $\Fl[\Ga]$-module $N$ (the
amalgam argument of \HA{cl:amalgam}(2),(3), which uses nothing about the coefficient module, $6\in\Fl^{\times}$), and localisation at $\m$ is exact; so
$\dim H^1(\Mn)_\m-\dim H^0(\Mn)_\m=\sum_\sigma[\Mn:\sigma]\bigl(\dim H^1(\Ga;\sigma)_\m-\dim H^0(\Ga;\sigma)_\m\bigr)$ over the irreducible
$\sigma=\Sym^{r}\Fl^{2}\otimes{\det}^{s}$, $0\le r\le\ell-1$, $s\in\Z/(\ell-1)$.
(b) \emph{$H^0$ at $\m$.} $H^0(\Ga;\sigma)=\sigma^{\SL_2(\Fl)}$ is non-zero only for $r=0$, where $T_p$ acts by $(1+p)p^{s}$, the Dickson system
$\omega^{s}(1\oplus\omega)$; so $H^0(\Ga;\sigma)_\m=0$, and, by left exactness along the composition series, $H^0(\Ga;\Mn)_\m=0$.
(c) \emph{$H^1(\Ga;\sigma)_\m$.} As groups $H^1(\Ga;\Sym^{r}\otimes{\det}^{s})=H^1(\Ga;\Sym^{r}\Fl^{2})$, with $T_p$ multiplied by $p^{s}$. For $r=0$ it is $0$
($\Ga^{\mathrm{ab}}\cong\Z/12$); for $r$ odd, $-I$ acts by $-1$ and it is $0$. For even $2\le r\le\ell-1$, $\Sym^{r}\Fl^{2}$ is irreducible and non-trivial, so
$H^0(\Ga;\Sym^{r}\Fl^{2})=0$, $H^1(\Ga;\Sym^{r}\Z_\ell^{2})$ is torsion-free and reduces onto $H^1(\Ga;\Sym^{r}\Fl^{2})$ (\HA{cl:bockstein}); by
Eichler--Shimura with its Hecke equivariance (\HA{lem:heckerankone}'s proof: \cite[Thm.~12.2.2, p.~108; Prop.~12.4.10, p.~117]{diamond_im}, verified, slot-17 L17-5)
the $\m'$-part of the reduction has dimension the number of characteristic-$0$ eigensystems of $M_{r+2}$ reducing to $\m'$, the Eisenstein one counted once and
each cuspidal one twice. The Eisenstein system of $E_{r+2}$ twisted by $\omega^{s}$ is $\omega^{s}\oplus\omega^{s+r+1}$; it equals $\bar\rho_\m$ iff
$(r,s)=(w_1-2,i)$ or $(w_2-2,i+w_1-1)$ (the computation of Claim~\ref{cl:seedtwist}(1); the solution $r=\ell-1$ needs $d\in\{1,\ell-2\}$, both of Dickson type; v1.18.0-pre.2, the referee F1's m-E2,
pre.1 read ``$d=1$''). A cuspidal system of
weight $w=r+2\le\ell+1$ twisted into $\bar\rho_\m$ has reducible residual representation, and \emph{a level-$1$ eigenform $f$ of weight $2\le w\le\ell+1$ with
reducible $\bar\rho_f$ has $\bar\rho_f^{\mathrm{ss}}\cong1\oplus\omega^{w-1}$}: if $f$ is ordinary, $\bar\rho_f|_{I_\ell}\sim\bigl(\begin{smallmatrix}\omega^{w-1}&\ast\\0&1\end{smallmatrix}\bigr)$
(Deligne, \STD: \cite[Thm.~2.1.4, p.~561]{wiles_ordinary}, verified, slot-16 L16-3), and a level-$1$ $\bar\rho^{\mathrm{ss}}$ is a sum of powers of $\omega$ (a character
of $G_\Q$ unramified outside $\ell$ with values in $\bar\F_\ell^{\times}$ is a power of $\omega$, \STD: class field theory over $\Q$), hence $\bar\rho_f^{\mathrm{ss}}\cong1\oplus\omega^{w-1}$;
if $a_\ell(f)\equiv0$, $\bar\rho_f|_{I_\ell}\cong\omega_2^{w-1}\oplus\omega_2^{\ell(w-1)}$ (Fontaine, \STD: \cite[Thm.~2.6]{edixhoven}; statement and attribution
verified in a secondary source, the theorem number pending, see the status paragraph; for $w\le\ell$ also \cite[Thm.~4.5 with Prop.~4.4(1)]{berger_li_zhu}), irreducible since $\ell+1\nmid w-1$ for $2\le w\le\ell+1$, a contradiction. So the cuspidal systems
twisted into $\bar\rho_\m$ are those of the cusp forms of weight $w\in\{w_1,w_2\}$ congruent to $E_w$, $g_w$ of them with multiplicity (the generalised eigenspace
in $S_w(\Fl)$ counts the characteristic-$0$ systems reducing to it, $S_w(\Z_\ell)$ being a Hecke-stable lattice; this count, and the count on $H^1$ above, use the
semisimplicity of the Hecke action in characteristic~$0$ and multiplicity one at level~$1$, \STD: the inputs of \HA{lem:heckerankone}, \cite[\S\S3.5--3.6,
pp.~54--55]{diamond_im}, verified, slot-17 L17-5, with Serre's Cours Ch.~VII \S5 and Diamond--Shurman Thm.~5.5.4 as cited there; v1.18.0-pre.2, the referee
F1's m-E2); at $w=\ell+1$ the ordinary case is of Dickson
type. That $g_w=0$ unless $\ell\mid B_w$ is the Eisenstein-ideal congruence (\STD, Prediction~\ref{pr:eiscusp}). Summing, $\dim H^1(\Ga;\sigma)_\m=1+2g_{w}$ for
$\sigma$ the Serre weight at the seed $w$, and $0$ for every other~$\sigma$.
For a cuspidal $\m$: the torsion of $H^1(\Ga;\Sym^{n}\Z_\ell)$ meets only the systems of \HA{prop:heckeH0} (the remark after \HA{prop:cohforms}), which are of
Dickson type, so $\Lambda_\m$ is torsion-free of rank $f_\m=\dim H^1(\Ga;\Mn)_\m$ (\HA{cl:bockstein}); $\m$ contains no Eisenstein system of weight~$k$, so $f_\m$
is even and $\rank\T_\m=f_\m/2$ by \HA{lem:heckerankone}; $A$ is a quotient of $\T_\m/\ell$ (\HA{prop:cohforms}(5)), so $\dim A\le f_\m/2$.
(A corollary, v1.18.0-pre.2, the referee F1's m-E2(iii): on a cuspidal Kummer class $f_\m=m_O+m_K(1+2g)\equiv m_O+m_K\pmod2$ is even, so $m_O\equiv m_K\pmod2$ ---
why the two seeds' multiplicities change together in Remark~\ref{rmk:dimwindows}.)
\end{proof}
\emph{Status of the inputs} (the owner's ruling~O3 of 2026-10-03: \PROVED{} only with the literature inputs \STD{} and the unverified ones flagged, as
\HA{lem:heckerankone}). Verified on page images: Ash--Stevens Lemma~1.1.1, Diamond--Im Thm.~12.2.2 / Prop.~12.4.10 and \S\S3.5--3.6, Wiles Thm.~2.1.4.
\textbf{Edixhoven Thm.~2.6: attribution verified in a secondary source, locator pending} (v1.18.0-pre.2, slot-18 R18-4,
\texttt{docs/implementation/slot18/R18-4\_literature.md}~\S2; pre.1 read ``not verified (the number read from the text layer, slot-17 R17-3)''): the statement
(ordinary: Deligne; supersingular, $k\le\ell+1$: Fontaine, $\bar\rho_f|_{I_\ell}\cong\psi^{k-1}\oplus\psi'^{k-1}$ with the two fundamental characters of level~$2$) and
its attribution, with the proof in \cite{edixhoven}, are confirmed in Edixhoven's Boston survey \emph{Serre's conjecture} (printed p.~19, the section on
``Theorem~4.5 of [15]''); the number ``Thm.~2.6'' is not confirmed, the paper itself not being in the library (an optional scan item for the owner). A second
locator, for $2\le w\le\ell$: \cite[Thm.~4.5, p.~12, with Prop.~4.4(1), p.~11]{berger_li_zhu} (read on the arXiv copy, slot-17 R17-4) gives
$\bar\rho_f|_{G_{\Q_\ell}}\cong\mathrm{ind}\,\omega_2^{w-1}$ when $v_\ell(a_\ell)>\lfloor(w-2)/(\ell-1)\rfloor$, which for $w\le\ell$ is $a_\ell\equiv0$ (at $w=\ell+1$ it
needs $v_\ell(a_\ell)\ge2$, so Fontaine's theorem stays the input there). The class-field-theory sentence is textbook and carries no locator. If the
Edixhoven input failed, the right-hand side would be a \emph{lower bound}, with equality for every $\m$ into which no non-ordinary ($a_\ell\equiv0$) level-$1$
eigenform of weight $\le\ell+1$ with reducible $\bar\rho_f$ twists (at its Serre weight $\Sym^{w'-2}\otimes{\det}^{s}$, which need not be a seed); the census
check below ($21\,690/21\,690$) is evidence that none does in the ranges computed (v1.18.0-pre.2, the referee F1's m-E1: pre.1 placed the possible extra
contribution at the seeds). The multiplicities $m_1,m_2$ are computed by Glover's recursion
(Claim~\ref{cl:glover}; slot-17 R17-1's \texttt{glover\_jh.py}, checked against exact Brauer characters on every semisimple class for $n\le1500$ at
$\ell\in\{17,37,67\}$); the recursion is read off \cite{glover} ((4.1), p.~439; (4.2), p.~440; (5.7), p.~454; (6.4), p.~461; held since 2026-10-05, read in the slot-21 literature pass,
\texttt{docs/implementation/slot21/missing\_literature.md}~\S7; v1.21.0-pre.1, in place of ``whose locator is not in the library''); a closed multiplicity formula is \emph{not} in Glover's paper.
\emph{Evidence} (census edition~3 against the formula, slot-17 R17-1~\S4, \texttt{census\_check.py}): every reducible non-Dickson $H^1$ class of every cell,
$21\,690$ of $21\,690$ (cell, class) pairs at $\ell\in\{7,\dots,103,691\}$, $0$ missing; and every cohomological and forms-side dimension the project measured at
$\ell=37$ (the report's \texttt{cor:onset}).

\begin{remark}[The windows on the multiplicity side]\label{rmk:dimwindows}
\PROVED{} (from Claim~\ref{cl:dimformula}) and computed. For a cuspidal Kummer class presented as $\omega^{t}(1\oplus\omega^{w-1})$ by
Claim~\ref{cl:seedtwist}(2), $g:=g_{w_K}$, and $(m_O,m_K)$ the multiplicities at the seeds $w_O$, $w_K$: $\dim V=m_O+m_K(1+2g)$ and $\rank\T_\m=\frac12\dim V$.
By script on slot-17 R17-1's class tables (\texttt{docs/implementation/slot17/r171/classes/}, $13$ irregular pairs, $\ell\le691$; recounted by
\texttt{docs/implementation/slot18/wp2/window\_events.py}), on every one of the $1492$ standard and $1906$ swapped cuspidal classes: $(m_O,m_K)$ becomes $(1,1)$ at
$k=w+t(\ell+1)$, $(2,2)$ at $k=\ell w+t(\ell+1)$ (the \emph{meeting}) and $(3,3)$ at $k=w+t(\ell+1)+\ell^{2}-1$, with no other change in between
(v1.18.0-pre.2, the referee F1's m-E3: \texttt{window\_events.py} tests the $(2,2)$ and $(3,3)$ events only; the exact event sequence --- these three events and
nothing else up to the upper end, and $\dim V=m_O+m_K(1+2g)$ at every event --- is checked by \texttt{docs/implementation/slot18/wp3/window\_check\_f1.py},
$1492/1492$ and $1906/1906$). So in the
\emph{first window} $w+t(\ell+1)\le k<\ell w+t(\ell+1)$, $\dim V=2+2g$ and $\rank\T_\m=1+g$; in the \emph{second}, $\ell w+t(\ell+1)\le k<w+t(\ell+1)+\ell^{2}-1$,
$\dim V=4+4g$ and $\rank\T_\m=2+2g$; then $6+6g$. At a regular class ($g_{w_1}=g_{w_2}=0$) the same windows give rank $1$, $2$, $3$. These are statements
about dimensions in the ranges computed (the class tables run to $k_{\max}=\ell\max(w)+\ell^{2}+2\ell$), not about the shape of~$A$.
\end{remark}

\begin{prediction}[The depth of a curvilinear block]\label{pr:depth}\CONJ\ (new at v1.18.0-pre.1; the $\ell=37$ report's \texttt{conj:depth} (0.2.0-pre.2, its \texttt{sec:fieldsproposals}),
taken into this document by slot-18 WP-2 under the owner's ruling~O3 of 2026-10-03)
Let $\m$ be a cuspidal, non-Dickson block of $H^1(\Ga;\Mn)$ with residue field~$\Fl$, and $\T_\m$ the Hecke algebra on the torsion-free $\Lambda_\m$ (so
$\T_\m/\ell\cong\T^{\mathrm f}_\m/\ell$ on the forms side, \HA{prop:cohforms}). Suppose one of:
\begin{enumerate}
\item[(a)] $\m$ is a Kummer class of an irregular pair $(\ell,w_K)$ at a weight where $(m_O,m_K)=(m,m)$ (Remark~\ref{rmk:dimwindows}), with \emph{seed rank}
  $r_0:=1+g_{w_K}$ (the standard and swapped families with $g=1$, and the $g=2$ classes at $547$);
\item[(b)] $\bar\rho_\m\cong\bar\rho_f\otimes\omega^{j}$, $j\in\{1,2\}$, at $k=\ell k_0+j(\ell+1)$ (the Frobenius layer of Prediction~\ref{pr:frobeniuslayer}), where
  $f$ is an ordinary level-$1$ eigenform of weight $k_0\le\ell+1$ whose Hecke algebra on $S_{k_0}$, completed at $f$'s system, is the ring of integers of a
  ramified quadratic extension of $\Q_\ell$ (two Galois-conjugate eigenforms congruent modulo~$\ell$); set $r_0:=2$, $m:=2$.
\end{enumerate}
If $m\le2$, then $\T_\m/\ell\cong\Fl[x]/x^{m r_0}$. In~(a), $\rank\T_\m=m\,r_0$ is \PROVED{} (Remark~\ref{rmk:dimwindows}), so the prediction is that
$\T_\m/\ell$ is curvilinear of the largest length its rank allows; in~(b), $\rank\T_\m=4$ is computed at every cell below, not proved here. In~(b) the
block is torsion-free, so that $\T_\m/\ell\cong\T^{\mathrm f}_\m/\ell$, by the $H^0$-argument of Claim~\ref{cl:dimformula}'s proof (the torsion meets only Dickson-type
systems), which does not use reducibility (v1.18.0-pre.2, the referee F1's m-E5).
\end{prediction}
\noindent\emph{Evidence} (forms side unless said; generated, $\dim A=\dim W$). The $\ell=37$ report's Table~\texttt{tab:families}: $57$ curvilinear blocks ---
$29$ standard Kummer, $8$ swapped, $2$ with $g=2$, $18$ irreducible --- $26$ of them also on cohomology (the $24$ cells at $37$, $(736,67)$, $(608,547)$, each with
$\T_\m$ saturated). Slot-18 WP-2 (Prediction~\ref{pr:kfonset}'s evidence, \texttt{docs/implementation/slot18/wp2/}): $8$ more --- $(1400,37)$, $(1438,37)$,
$(10632,103)$ at $(m,r_0)=(2,2)$, swapped; $(1154,547)$, $(1700,547)$, $(2246,547)$, $(1032,547)$, $(1578,547)$ at $(m,r_0)=(1,3)$, standard and swapped. At
$(m,r_0)=(1,2)$ the bound $\dim\T_\m/\ell\le2$ makes the statement automatic. (v1.18.0-pre.2) The slot-18 referee F1 (\texttt{docs/implementation/slot18/wp3/},
predicted before each run, $9/9$ held): $8$ more --- $(1364,37)$, $(1476,37)$, $(1514,37)$, $(10736,103)$ at $(m,r_0)=(2,2)$, swapped; $(16988,547)$, $(33914,547)$,
$(11952,547)$ at $(m,r_0)=(1,3)$; and $(34460,547)$, $k=34462$, the meeting of the class $a=1$, at $(m,r_0)=(2,3)$: $\dim W=6=\rank\T_\m$, Hilbert function
$(1,1,1,1,1,1)$, generated, $\F_{547}[x]/x^6$. In all $73$ curvilinear blocks on the forms side, $26$ of them on cohomology. \emph{Why $m\le2$}: at $m=3$ it
fails, Hilbert function $(1,2,2,1)$ at $k=1412$, $1450$ (standard, $37$) and $1438$, $1476$, $1514$ (swapped, $37$), $10736$ (swapped, $103$). The $(2,3)$ datum
separates the two readings of the depth law: the hypothesis is $m\le2$, not $m r_0\le4$ --- $(m,r_0)=(2,3)$ is curvilinear of length~$6$ while $(3,2)$ is
$(1,2,2,1)$ (pre.1 read ``\emph{Untested}: $(m,r_0)=(2,3)$ \ldots\ an extrapolation, kept as such''; the referee F1's M-2). \emph{Untested}: a congruence between two Galois orbits rather than a conjugate pair; any cell of~(b)
on cohomology (the cheapest is $(2210,67)$, cluster only). Whether $\T_\m$ is saturated at every block of (a), (b) --- so that the census algebra $A$ equals
$\T_\m/\ell$ --- is the report's \texttt{q:satfam}: measured yes at the $26$ cohomology cells; Prediction~\ref{pr:kfonset} asserts it, with the
curvilinearity, in its windows. \emph{Not taken into this document}: the report's \texttt{conj:law} (the reducible quotient modulo~$\ell^{2}$, $28/28$
cells, unexplained). Slot-18 R18-1 tested it at scale (\texttt{docs/implementation/slot18/R18-1\_report.md}, final; predictions committed
before the grid, \texttt{r181/predictions.md}): $601$ of $715$ predicted cells run at $8$ primes. (C-1) \emph{as stated} is \textbf{refuted} at all $50$
second-window cells at $\ell=37$ ($k=1412,\dots,2708$, $a=1,2$; the block is not curvilinear and every $s\in\Z/37$ is admitted for a generating set
of $\T_\m$ --- the $T_p$, $p<100$, at $19$ cells, with $T_{37}$ added at the $31$ cells $a=1$, $k\ge1628$, where the $T_p$ alone do not generate
(the report's referee F4, M1; re-run \texttt{docs/implementation/slot18/r185/}) --- so $\T_\m/I_{\mathrm{red}}$ is not cyclic; one of them, $(1412,37)$,
was seen in a smoke run before the predictions commit). The restricted form
(C-1$'$) --- the law wherever $\T_\m/I_{\mathrm{red}}$ is cyclic --- held at $548/548$ cells at $7$ primes ($37,41,43,59,67,101,103$), $\ell$-adically
(modulo $\ell^{e-1}$) at $20/20$, vacuous at the $3$ $g=2$ cells at $547$; (C-2), (C-3) held at $433/433$. The report [L37] v0.3.0-pre.2 (refereed by F4 in pre.1, repairs applied) carries (C-1$'$)
as its conjecture. A pointer only, no \CONJ{} taken in (R18-1's final figures replace F1's preliminary ones, $503/50/3$ of $556$, in this pre.2
text; slot-18 R18-5).

\subsection{The block \texorpdfstring{$(258,37)$}{(258,37)} exactly: the Hecke algebra, its over-orders, the lattice}\label{sec:block258}
(New at v1.19.0-pre.1, hecke-extensions slot~19, WP-2; owner proposal of 2026-10-03, \texttt{docs/implementation/slot18/lead\_carry.md}~C4. The statements
are slot-19 R19-1's (\texttt{docs/implementation/slot19/R19-1\_report.md}~\S8, predictions committed before any computation, $10/10$ held); the checkable
item~\ref{chk:latcmp} is the gate of slot-19 WP-LC, which merges only after this text is at pre.1 (rule~22). Placed at the end of \S\ref{sec:ext}, so that no
earlier number moves.) \textbf{Setting.} $(n,\ell)=(258,37)$, $k=260$, $\m=\m_1$ the class $\mathrm{eis}(1,6)=\theta E_6\equiv\theta^{6}E_{32}$ --- the first cell of
Prediction~\ref{pr:kfonset}'s Kummer--Frobenius onset and of the $\ell=37$ report's table (\Lcite{thm:example}). $\Lambda:=H^{1}(\Ga;\Mn\otimes\Z_{37})$,
$\Lambda_\m$ its $\m$-part: free of rank~$8$ over $\Z_{37}$, torsion-free, $\Lambda_\m/37=V$ (\Lcite{thm:example}(2)). $\T_\m\subseteq\End(\Lambda_\m)$ is the
$\Z_{37}$-algebra generated by \emph{all} $T_p$ (with $T_{37}=U_{37}$), $\T^{\mathrm f}_\m$ the same on $S_{260}(\Z_{37})_\m$; $t:=T_2-29$; $A:=\F_{37}[x]/x^{4}$.
\textbf{Tags} (here only): \CMPD{} = computed by slot-19 R19-1 from the exact integral Hecke matrices of $S_{260}$, by two independent routes (python-flint and
in-house), at $37$-adic precision $37^{32}$, far above every valuation read --- proof-grade for the statements as written, but not a proof in this document.

\begin{claim}[The Hecke algebra of the block]\label{cl:258alg}
\begin{enumerate}
\item \CMPD{} $\T^{\mathrm f}_\m=\Z_{37}[t]\cong\Z_{37}[X]/(P)$, $P\equiv X^{4}\pmod{37}$, the valuations of the coefficients of $P$ (low to high) $(2,2,2,1,0)$,
  $v(\operatorname{disc}P)=6$: every $T_n|_{S_m}$, $n\le21$ (a $\Z$-basis of the Hecke algebra of $S_{260}$ by $a_1$-duality with the Miller basis) and $n=37$,
  is a polynomial in $t$ with $\Z_{37}$-coefficients. So $\Z_{37}[t]$ is the \emph{full} Hecke algebra of the block, not an anemic one.
\item \PROVED{} (from \HA{prop:cohforms}(2),(3) and \Lcite{thm:example}(2)) $\T_\m\cong\T^{\mathrm f}_\m$, $T_p\mapsto T_p$ for every $p$ ($p=37$ included). In
  particular $\T_\m=\Z_{37}[t]$ and $\T_\m$ is Gorenstein (monogenic).
\item \CMPD{} $P=P_1P_2$ with $P_1$, $P_2$ Eisenstein of degree~$2$; $U_{37}$ has slopes $1$ on $K_1$ and $5$ on $K_2$ ($K_i:=\Q_{37}[X]/(P_i)$, ramified quadratic);
  the normalisation is $\tilde\T:=\mathcal O_{K_1}\times\mathcal O_{K_2}$, $\mathcal O_{K_i}=\Z_{37}[\rho_i]$; $v(\operatorname{Res}(P_1,P_2))=2$; all six pairwise distances of the four
  roots are $1/2$, so every congruence between two of the four eigenforms has depth exactly $1/2$ (it holds modulo $\pi$, not modulo $\pi^{2}$; attained at $a_2$).
\item \CMPD{} $\tilde\T/\T_\m\cong\F_{37}[\varepsilon]/\varepsilon^{2}$ as a $\T_\m$-module (cyclic, generated by the class of the idempotent $e_1$: $37e_1\in\T_\m$,
  $te_1\notin\T_\m$; Smith form $(37,37)$, not $37^{2}$), and $\tilde\T/\T_\m\cong\Z_{37}[X]/(P_1,P_2)$, the inter-orbit congruence module. The forms congruence module
  $S_\m/(S_\m^{K_1}\oplus S_\m^{K_2})\cong(\Z/37)^{2}$; the $\mathcal O$-span of the four eigenforms has index $37^{3}$ in $S_\m\otimes\mathcal O$.
\item \CMPD{} The over-orders of $\T_\m$ in $\tilde\T$ are exactly $\T_\m\subsetneq R\subsetneq\tilde\T$ with $R:=\T_\m+\Z_{37}\,te_1$ (of the $38$ lattices of index~$37$ between $\T_\m$
  and $\tilde\T$ exactly one is a ring). In fact, since $\tilde\T/\T_\m\cong A/x^{2}$ is cyclic by~(4), exactly one of the $38$ is even a $\T_\m$-submodule --- the socle line
  $R=\T_\m+\Z_{37}te_1$ --- so the over-order list $\{\T_\m,R,\tilde\T\}$ follows from (4) and the single check that $R$ is a ring (referee F1, slot~19, \texttt{f1\_block258.py}). $R$ is local and Gorenstein ($\dim\soc R/37R=1$), $R=(\m:\m)$ and $\m=tR\cong R$. As $A$-modules ($x\mapsto t$):
  $\T_\m/37\cong A$, $R/37R\cong A/x^{3}\oplus A/x$, $\tilde\T/37\tilde\T\cong A/x^{2}\oplus A/x^{2}$. Every over-order of $\T_\m$ (and of its two quotients
  $\mathcal O_{K_i}$) is Gorenstein: $\T_\m$ is a \emph{Bass order}.
\end{enumerate}
\end{claim}

\begin{proof}[Proof of (2)]
$\Lambda$ is torsion-free (\Lcite{thm:example}(2)), so $\T^{\mathrm{coh}}=\T^{\mathbb L}$ on the block, and \HA{prop:cohforms}(3) gives $\T^{\mathbb L}\cong\T^{\mathrm f}$ with
$T_p\mapsto T_p$; localise at $\m$. $\T_\m\cong\Z_{37}[X]/(P)$ is a complete intersection, hence Gorenstein. (R19-1 recorded ``$\T_\m\cong\T^{\mathrm f}_\m$: forms side only'';
this is the proof.)
\end{proof}

\begin{claim}[The lattice]\label{cl:258lattice}
\PROVED{} given Claim~\ref{cl:258alg}(5) (\CMPD), \Lcite{thm:example}(1) ($V\cong A\oplus A/x^{3}\oplus A/x$, \CMPD{} at the cell) and \emph{Gorenstein ideal
theory} (over a reduced one-dimensional Gorenstein order $S$, a fractional ideal $I$ with $(I:I)=S$ is invertible; \STD, \cite{bass_ubiquity}, \S6, locator not verified,
rule~A2 --- the paper is not in the project library (slot-19 referee F4, m1) and nothing here rests on it --- and proved in the proof below, v1.19.0-pre.2; an invertible ideal of a semilocal ring is principal). Let $\Lambda_\m=\Lambda_\m^{+}\oplus\Lambda_\m^{-}$ be the $\iota$-eigenspace decomposition
(\HA{prop:iotasplit}; $\T_\m$-stable by \HA{prop:iotacommutes}). Then one of $\Lambda_\m^{\pm}$ is $\cong\T_\m$ and the other is $\cong\m\cong R$; so
\[
  \Lambda_\m\;\cong\;\T_\m\oplus\m\qquad\text{as $\T_\m$-modules.}
\]
The argument does not decide which sign is the free one; the measurement does: \CMPD{} (slot-19 WP-LC, \texttt{HeckeLatticeModuleValidation},
Check~\ref{chk:latcmp}(d) at $M=1,\dots,4$; v1.19.0-pre.3) the \emph{plus} sign is free,
\[
  \Lambda_\m^{+}\cong\T_\m,\qquad\Lambda_\m^{-}\cong\m
\]
(``$+$'' the Eisenstein side of \HA{def:iota}). Here $\Lambda_\m$ is the $H^{1}$ lattice. On the modular-symbol lattice $H^{1}_c(\Ga;\Mn\otimes\Z_{37})_\m$ the free sign is the
other one ($-$, the geometric-even sign; slot-19 R19-8, re-derived by isign), because $\m$ is residually Eisenstein
($T_p\equiv p+p^{6}\bmod 37$, $p=2,3,5$). The natural map $H^{1}_{c,\m}\to H^{1}_\m$ is injective with a cokernel of
order $37$ in each sign (\CMPD: one cyclic factor per sign by isign (K); the order by R19-9 \S4, P9, the dual index $[H^{1}(D)_\m:H^{1}_{!}(D)_\m]=37$, reproduced by slot-19 WP-BIL), so the two lattices are not isomorphic as $\T_\m[\iota]$-modules (\CMPD, isign). So the statement is \PROVED{} given \CMPD{} inputs, and the sign is \CMPD.
\end{claim}

\begin{proof}
By Eichler--Shimura (\HA{prop:cohforms}, its per-sign multiplicity-one statement in~(5) --- ``one of the $\pm$-parts \dots{} free of rank one'' --- and \Lcite{prop:cuspidal}(2): no Eisenstein system in the block in characteristic~$0$ (the block is residually Eisenstein)) each $\Lambda_\m^{\pm}\otimes\Q$ is free of rank one over
$\T_\m\otimes\Q=K_1\times K_2$, so $\Lambda_\m^{\pm}$ is isomorphic to a fractional ideal $I^{\pm}$ of $\T_\m$. Its multiplier ring $S^{\pm}:=(I^{\pm}:I^{\pm})$ is an order
containing $\T_\m$, hence contained in $\tilde\T$: $S^{\pm}\in\{\T_\m,R,\tilde\T\}$ (Claim~\ref{cl:258alg}(5)), all Gorenstein, so $I^{\pm}$ is invertible over $S^{\pm}$ and
$I^{\pm}\cong S^{\pm}$. (The Gorenstein input, for $S=S^{\pm}$ local, i.e.\ $\T_\m$ or $R$: fractional ideals over a one-dimensional Gorenstein local ring are reflexive,
$(S:(S:I))=I$; with $I^{\ast}:=(S:I)$, $(S:II^{\ast})=((S:I^{\ast}):I)=(I:I)=S$. If $II^{\ast}\subsetneq S$ then $II^{\ast}\subseteq\mathfrak m_S$ and
$(S:II^{\ast})\supseteq(S:\mathfrak m_S)\supsetneq S$, as $\mathrm{Ext}^{1}_S(k,S)\neq0$ in depth~$1$ --- a contradiction; so $II^{\ast}=S$, $I$ is invertible, hence principal.
For $S=\tilde\T$, a product of DVRs, every fractional ideal is principal anyway.) Reducing modulo~$37$, $V=V^{+}\oplus V^{-}$ with $V^{\pm}\cong S^{\pm}/37S^{\pm}\in\{A,\ A/x^{3}\oplus A/x,\ A/x^{2}\oplus A/x^{2}\}$. As the Jordan type of $x$ on a finite $A=\F_{37}[x]/x^{4}$-module is unique, the only pair whose sum is $A\oplus A/x^{3}\oplus A/x$ is $\{A,\ A/x^{3}\oplus A/x\}$, i.e.\ $\{S^{+},S^{-}\}=\{\T_\m,R\}$; and
$R\cong tR=\m$ ($t$ is a non-zero-divisor). (R19-1's route, Bass's structure theorem for torsion-free modules over a Bass order \cite{bass_ubiquity}~\S7, reaches
the same conclusion without the $\iota$-split; the split makes the weaker input above enough. R19-1's input is a literature locator that stays unchecked --- the
paper is not in the project library --- and this route does not need it; slot-19 referee F4, m1.)
\end{proof}

\begin{remark}[An independent check already measured]\label{rmk:258Urank}
$U_{37}\bmod 37$ has rank $2$ on $A$ and rank $1$ on $A/x^{3}\oplus A/x$ (\CMPD, R19-1~\S5), so Claim~\ref{cl:258lattice} predicts $\rank U_{37}|_V=3$ --- the value measured
on cohomology (\Lcite{tab:cells}; slot-18 R18-3). The excluded alternatives give $4$ ($\T_\m^{2}$) and $2$ ($\T_\m\oplus\tilde\T$).
\end{remark}

\begin{checkitem}[The lattice comparison at $(258,37)$; the gate of slot-19 WP-LC]\label{chk:latcmp}
For $M=1,2,3,4$ let $H_M:=H^{1}(\Ga;\Mn\otimes\Z/37^{M})_\m$, computed on the $\Z/\ell^{K}$ route (\HA{sec:saturation}; \S\ref{sec:ring}: Howell and Smith forms over
$\Z/37^{M}$), with the operators $t=T_2-29$ and $U_{37}$, and $\iota^{\ast}$. Write $\langle d_1,\dots,d_r\rangle_M$ for the Smith valuations of an operator on $H_M$
(the invariants $\min(d_i,M)$, so that its cokernel is $\bigoplus_i\Z/37^{\min(d_i,M)}$). Then:
\begin{enumerate}
\item[(a)] \PROVED{} $H_M\cong\Lambda_\m/37^{M}$: free of rank~$8$ over $\Z/37^{M}$ (torsion-freeness, \Lcite{thm:example}(2); $H^{2}(\Ga;-)$ has no $\ell$-torsion for
  $\ell\ge5$, \HA{lem:h2torsion}; so $H^{1}(\Ga;\Mn\otimes\Z/37^{M})=\Lambda/37^{M}$, and localisation at $\m$ commutes).
\item[(b)] $\dim_{\F_{37}}H_M/(37,t)H_M=3$ (the minimal number of $\T_\m$-generators: $1$ for $\T_\m$, $2$ for $\m$).
\item[(c)] $t$: $\langle0,0,0,0,0,1,1,2\rangle_M$; \quad $U_{37}$: $\langle0,0,0,1,5,6,6,6\rangle_M$.
\item[(d)] Per sign: $H_M=H_M^{+}\oplus H_M^{-}$, each free of rank~$4$; on one sign $\dim H/(37,t)H=1$, $t$: $\langle0,0,0,2\rangle_M$, $U_{37}$: $\langle0,0,6,6\rangle_M$
  (the free summand $\T_\m$); on the other $\dim H/(37,t)H=2$, $t$: $\langle0,0,1,1\rangle_M$, $U_{37}$: $\langle0,1,5,6\rangle_M$ (the summand $\m$). The sign of the
  free summand is \emph{reported}, not predicted.
\item[(e)] $P(t)=0$ on $H_M$, and $T_p=Q_p(t)$ on $H_M$ for every prime $p<100$ ($p=37$ included), with $P$ and the $Q_p$ the forms-side data of R19-1 modulo $37^{6}$
  (\texttt{docs/implementation/slot19/r191/data/universal\_eigenform\_n1000\_mod37e6.txt}: $T_n=c_0+c_1t+c_2t^{2}+c_3t^{3}$).
\end{enumerate}
\textbf{Status.} (a) \PROVED. (b)--(e) are \PROVED{} given the \CMPD{} inputs of Claims~\ref{cl:258alg}, \ref{cl:258lattice} (the Smith valuations of $t$ and $U_{37}$ on each of
$\T_\m$, $R$, $\tilde\T$ are R19-1's, \S5): (b)--(d) are the invariants of $(\T_\m\oplus\m)/37^{M}$, (e) is Claim~\ref{cl:258alg}(2). So a failure of~(a) or~(e) is a
defect of the code; a failure of (b)--(d) with (a), (e) holding refutes a \CMPD{} input or \Lcite{thm:example}(1), and is a finding. \textbf{Discrimination.}
$\T_\m\oplus\tilde\T$ shares (b) and $t$'s valuations but has $U_{37}$: $\langle0,0,1,1,5,5,6,6\rangle_M$ (so already $\rank U_{37}|_{H_1}=2$, against~$3$); $\T_\m^{2}$ has
(b)~$=2$; so (b)--(d) at $M\ge2$ separate the three candidates. \textbf{Cost} (plan): $\le6$\,GB at $M=4$, predicted from $1.0$\,GB at $M=1$.
\textbf{Run} (v1.19.0-pre.3). \CMPD{} (slot-19 WP-LC, \texttt{HeckeLatticeModuleValidation}; predictions committed before the run,
\texttt{docs/implementation/slot19/wplc/predictions.md}~\S1, outcomes \S3): (a)--(e) \emph{held} at $M=1,\dots,4$ --- rank~$8$, free; $\mu=3$; $t$:
$\langle0^{5},1,1,2\rangle_M$, $U_{37}$: $\langle0^{3},1,5,6,6,6\rangle_M$; the free summand on the plus sign (Claim~\ref{cl:258lattice}); (e) at all $25$ primes $p<100$
(the full arm), on the block and on each sign. The planted negative controls $\T_\m\oplus\tilde\T$ ($U_{37}$) and $\T_\m^{2}$ ($\mu$) are rejected; independent runs
at $K=2,3$ agree with the $K=4$ run reduced. Measured cost: $\approx22$\,min and $2.9$\,GB (full arm), $\approx2$\,min and $0.8$\,GB (default arm).
\end{checkitem}

\begin{remark}[The $\ell=7$ controls of WP-LC: per-sign freeness fails]\label{rmk:latcmpseven}
\CMPD{} (slot-19 WP-LC, the controls of \texttt{HeckeLatticeModuleValidation}; cells chosen by a forms-side scan of every even $n\le110$,
\texttt{docs/implementation/slot19/wplc/l7\_pred.py}, and predicted before the run, \texttt{wplc/predictions.md}~\S2; outcomes \S3; the over-order identification at
$(48,7)$ closed by the slot-19 referee F4, m5; v1.19.0-pre.3). Three non-Eisenstein blocks at $\ell=7$, written by $(n,7)$ and the residual system
$(a_2,a_3,a_5,a_7)$: $(48,7)(1,0,0,0)$, rank~$2$, Loewy length~$2$, $\T_\m=\Z_7[t]$; $(90,7)(3,3,1,0)$, rank~$2$, Loewy~$2$, $\T_\m=\Z_7[t]$; $(90,7)(1,0,0,0)$, rank~$4$,
Loewy~$3$, $\T_\m$ not monogenic in $t$. Every block is free over $\Z/7^{M}$ --- measured (all cyclic-decomposition exponents equal $K$), and also the argument of
Check~\ref{chk:latcmp}(a), whose universal-coefficient step \HA{lem:h2torsion} needs only $\ell\ge5$ (F4, m7); rank, sign ranks, $P(t)=0$ and $T_p=Q_p(t)$ for $p\le13$
hold, the items forced by \HA{prop:cohforms}. The prediction ``$\Lambda_\m^{\pm}\cong\T_\m$ on both signs'' is \textbf{refuted on five of the six signs}:
\begin{enumerate}
\item $(48,7)(1,0,0,0)$: minus $\cong\T_\m$ (free); plus $\cong\Z_7[t/7]$. With $s=t/7$, $s^{2}+as+b=0$, $v(a)=v(b)=0$, $v(a^{2}-4b)=2$, so $\Z_7[s]$ is \emph{not} maximal: the
  over-orders form the chain $\T_\m\subset\Z_7[t/7]\subset\mathcal O_{\max}$, of index~$7$ at each step, $\mathcal O_{\max}=\Z_7[(s+a/2)/7]$ étale. The Smith valuations of $t$
  and $\mu$ do not separate $\Z_7[t/7]$ from $\mathcal O_{\max}$ (both $\langle1,1\rangle$, $\mu=2$), but $U_7$ does: $\langle1,5\rangle$, $\langle2,4\rangle$, $\langle3,3\rangle$ on
  $\T_\m$, $\Z_7[t/7]$, $\mathcal O_{\max}$ (F4, from WP-LC's $P$ and $Q_7$). The plus reading $\langle2,4\rangle_4$ is $\Z_7[t/7]$ and nothing else; and a lattice over a
  quadratic order is isomorphic to its multiplier ring (quadratic orders are Gorenstein; the argument in the proof of Claim~\ref{cl:258lattice}), so this identifies
  the plus sign. The same shape as $(258,37)$: one sign free, the other the next over-order.
\item $(90,7)(3,3,1,0)$: both signs $\cong\Z_7[t/7]$, which is here the maximal order ($v(\mathrm{disc})=0$, $\cong\Z_7\times\Z_7$; $t$: $\langle1,2\rangle$, $\mu=2$); no
  identification gap.
\item $(90,7)(1,0,0,0)$: plus consistent with $\Z_7[t/7]$ ($t$: $\langle1,1,1,2\rangle$, $\mu=4$, $U_7$: $\langle2,4,4,4\rangle_4$); minus none of the $\Z_7[t/7^{j}]$ ($t$:
  $\langle0,0,2,3\rangle_M$, $\mu=2$, $U_7$: $\langle1,4,4,4\rangle_4$). No candidate was enumerated here; the readings are pinned, the order is not named.
\end{enumerate}
\emph{Scope} (F4, m6). This is per-sign freeness of $H^{1}(\Ga;\Mn\otimes\Z_7)_\m$ over $\T_\m$ failing at $n=48,90$ ($\ell=7\le n$), outside the range $n<\ell-1$ of
Claim~\ref{cl:lveis}(2) (Harder) and of the Fontaine--Laffaille comparison: a \CMPD{} statement about three blocks (five signs), not a theorem about the lattice for
$\ell<k-1$. It is qualitatively expected; no literature pointer is offered yet. The validation pins the observed order per sign and prints
\texttt{prediction\_free=refuted}; that is not a re-pin of an item of Check~\ref{chk:latcmp}. Adding $\mathcal O_{\max}$ to the code's candidate list is a later code
item (F4, m5).
\end{remark}

\begin{remark}[What (258,37) says about the depth law and the next cells]\label{rmk:258next}
\CMPD{} (R19-1~\S6) the universal eigenform $\sum_nT_nq^{n}\in(\T_\m/37^{6})[[q]]$ to $n\le1000$, the four eigenforms in $\mathcal O_{K_1}\times\mathcal O_{K_2}$ modulo
$37^{6}$ and the $\Q(\alpha)$-eigenform ($\alpha=a_2$, degree $21$, Maeda) to the Sturm bound, with Ribet's form $g\in S_{32}$ ($a_n(g)\equiv\sigma_{31}(n)$ modulo $37$, not
modulo $37^{2}$: depth $1=v(B_{32}/32)$) and $\theta^{6}g\equiv\theta E_6\equiv f_i\pmod\pi$. The Galois side of the two seeds is slot-18 R18-2's pair of Kummer
generators, $X^{37}-\eta_{32}$ (Ribet's $H$) and $X^{37}-\eta_6$ (\texttt{docs/implementation/slot18/R18-2\_report.md}~\S1; [L37]~\texttt{comp:kummergen}); the
number-field layer that reproduces them in C++ is \S\ref{sec:lvnf}. \emph{Not claimed}: the cohomology lattice at the other $23$ cells of [L37]'s table
(the same argument applies wherever Claim~\ref{cl:258alg}(5)'s over-order data are computed); the small Loewy-$2$\,/\,$3$ controls at $\ell=7$ that the plan asks
for were not chosen by R19-1; WP-LC chose and ran them (Remark~\ref{rmk:latcmpseven}, v1.19.0-pre.3).
\end{remark}

\noindent\emph{Code map (\textbf{built}, slot-19 WP-LC; planned at v1.19.0-pre.1, contract first; v1.19.0-pre.3).} \texttt{HeckeExt/HeckeLatticeModule} and
\texttt{HeckeLatticeModuleValidation} (registered \texttt{standard}, moved to \texttt{slow} by the owner's ruling of 2026-10-03; $301$ checks; read-only env arm
\texttt{HECKE\_LATTICE\_MODULE\_FULL=1}: all $p<100$ and independent $K=2,3$ runs, $585$ checks): $H_M$ on the shipped $\Z/\ell^{K}$ module route
(\texttt{EquivariantCohomologyOnK}) over $\Z/37^{M}$, with $\iota^{\ast}$, the Howell\,/\,Smith forms of \S\ref{sec:ring}, $M=1\dots4$; Check~\ref{chk:latcmp} (a)--(e) at $(258,37)$; the
eigenform data of (e) read from R19-1's tracked file; the Loewy-$2$\,/\,$3$ controls at $\ell=7$ chosen by WP-LC and predicted before the run (rule~20; Remark~\ref{rmk:latcmpseven}). Value-free: no
golden line, digest, record key or census verdict changes.

%======================================================================
\section{The algorithms, stage by stage, with certificates}\label{sec:alg}
%======================================================================

Throughout, $d_q=\dim H^q$, $d=\dim\Htot=O(n)$, $D:=d_\m=\dim_{\Fl}H_\m$, $k:=\#S$, and costs are
in $\Fl$-multiplications. Everything is exact linear algebra over $\Fl$ on the existing engine
(\texttt{engine::elim}, \texttt{FpPoly}); no floating point, no field extension
(\S\ref{sec:field}).

\subsection*{Stage 0 --- preconditions and the three additions}
Available: bases of $X_1,X_4$, projectors $P_q$, $A^{(q)}(T_p)$ for all $p$
(\S\ref{sec:today}). Required additions, in order of value:
\begin{enumerate}[label=(A\arabic*)]
\item\label{add:bdry} \emph{Boundary Hecke action} on $X_2,X_5$: the formula of \HA{def:hecke}
  with $\Ga_\infty$ in place of $\Ga$. The relevant set is
  $B_q:=\{\delta\in\Ga\,\mathrm{diag}(1,q)\Ga:\ \delta\text{ upper triangular}\}$ (the elements
  fixing the cusp; $B_q/\Ga_\infty\to\Ga\,\mathrm{diag}(1,q)\Ga/\Ga$ is a bijection by column
  Hermite normal form), and the decomposition to use is the analogue of the \emph{left} cosets
  $\Ga g\Ga=\bigsqcup_i\delta_i\Ga$ of \HA{def:hecke}:
  $B_q=\delta_\infty\Ga_\infty\sqcup\bigsqcup_{j=0}^{q-1}\delta_j\Ga_\infty$ with
  $\delta_\infty=\mathrm{diag}(1,q)$, $\delta_j=\smat{q}{j}{0}{1}$, so that
  $T\delta_\infty=\delta_\infty T^q$, $T\delta_j=\delta_{j+1}$, $T\delta_{q-1}=\delta_0T$. On the
  circle model the operator is $c\mapsto\sum_i\rho_n(\delta_i)\,c(\tau_\infty(\cdot))$ with a
  one-dimensional $\tau$ (a $\Ga_\infty$-twisted chain map on the line), certified by
  $\partial\tau=\tau\partial$ exactly as \HA{alg:tau}, and by \HA{cl:chi} as oracle.
  \emph{Status (v1.4.0; [C] v1.4.1).} This one-dimensional $\tau_\infty$ is the
  $\Ga_\infty$-level twisted datum of \HA{obl:taubdry}(v.c), a validation route; the normative
  operator is \HA{eq:bdrychain} from the canonical $\tau_\partial$ of \HA{alg:taubdry} (cited,
  not re-described here; shipped as \texttt{Hecke/HeckeBoundaryDatum}). Clause (v) of
  \HA{obl:taubdry} is three named comparisons --- (v.a) degree $0$ at every $\nu$, against
  $\sum_{\delta\in B_N/\Ga_\infty}\rho_n(\delta)$ on $M^T$; (v.b) degree $1$ at $\nu=1$, against
  \HA{eq:bdrygroupform} in the coordinate $\Phi$ of \HA{prop:bdrychar}(3),
  $[\varphi]\mapsto[\varphi(\langle v_0,Tv_0\rangle)]$, representative map and lift on that one
  orientation (the opposite orientation gives a different form); (v.c) every computed degree at
  $\nu=1$, the datum above --- so at $\nu\ge2$ clause (v) is \emph{degree $0$ alone}, $H^1$
  there being covered by (vii) and the operator's own certificates; the general-$\nu$
  double-coset decomposition of $B_N$ is recorded in [C] as \emph{not stated}. Through v1.3.3
  this item read clause (v) as one comparison in every degree at every $\nu$.
  \emph{Scope of the circle model} (owner decision on the verdict's point~5, 2026-09-15): the
  boundary group is $\Ga_\infty=\langle T\rangle\times\langle-I\rangle$ and the boundary cohomology
  is $H^\bullet(\partial;M):=H^\bullet(\Ga_\infty;M)$ as defined in \HA{sec:six}; the circle
  model is $C^\bullet_{\mathrm{eq}}(K_\partial,\mathrm{Rel}_\partial;L)$ with
  $\mathrm{Rel}_\partial$ the cusp relation $f(v_3)=\rho_n(T)f(v_0)$ \emph{together with the
  $-I$ self-relations} $f(\sigma)=\rho_n(-I)f(\sigma)$ (\HA{prop:circlemodel}, v1.4.0; through
  v1.3.3 this document, like [C], called it ``the single relation''). It computes the boundary
  equivariant complex $C^\bullet_\partial$ for \emph{every} $n$ and every coefficient ring
  (\HA{prop:circlemodel}(4)): the single-relation description holds exactly when
  $\rho_n(-I)=\mathrm{id}$, i.e.\ $n$ even or $2=0$, and for $n$ odd with $2$ invertible the
  circle model is the zero complex (\HA{rmk:bdryoddn}) --- what the code computes. Its
  cohomology is $H^\bullet(\langle T\rangle;M)$ for $n$ even, and agrees with
  $H^\bullet(\Ga_\infty;M)$ \emph{exactly when}
  $2\in R^\times$ --- for $n$ even $-I$ acts trivially, so
  $H^\bullet(\Ga_\infty;M)=H^\bullet(\langle T\rangle;M)^{-I}=H^\bullet(\langle T\rangle;M)$ when $2$
  is invertible. The circle model is therefore the specialisation of the boundary datum at
  $2\in R^\times$, as the 7-vertex complex is the specialisation of the orbifold model at
  $6\in R^\times$; at $\ell=2$ they part company, and substantively so (closing paragraph
  of \ref{add:cone}). Under the standing hypothesis $\ell\ge5$ of \S\ref{sec:conv} nothing in this
  document is affected.
  \textbf{Not} the upper-triangular right-coset representatives $\beta_j=\smat1j0q$,
  $\beta_\infty=\mathrm{diag}(q,1)$ of \HA{prop:cosets}(1): they represent
  $\Ga_\infty\backslash B_q$, and their sum
  $M_\beta:=\sum_j\rho_n(\beta_j)+\rho_n(\beta_\infty)$ is a well-defined endomorphism of $M^T$
  \emph{only} --- the monomials $x^{n-b\ell}\mathrm{Nm}^b$ are an eigenbasis with eigenvalue
  $\chi_{b+1}$ in place of $\chi_b$ --- not asserted to be induced by any cochain map of
  $C^\bullet_\partial$ (a sum over $\Ga_\infty\backslash B_q$ is not $\Ga_\infty$-equivariant as
  a formula on cochains), and built by the implementation as the fibre-level \emph{negative
  control} \texttt{beta\_control\_fiber\_operator} alone (\HA{prop:bdrychar}(2), v1.4.0; through
  v1.3.3 this document called it ``the $\det$-twisted (transposed) operator''). Equivariance
  $r^qA^{(q)}=A^{(q)}_\partial r^q$ of the two actions is \HA{obl:taubdry}(vii): asserted at
  $\ell\ge5$ ($864$ $(\text{cell},q)$ pairs, slot 1), and at $\ell\in\{2,3\}$, where the model
  complexes compute the orbifold model and not group cohomology, a \textsc{comp} \emph{record}
  only (\HA{rmk:conerestrictsmallell}: $216/216$, consistent with a pair homotopy over $\Z$,
  claimed nowhere, and a failure elsewhere would contradict nothing); the negative control
  breaks it on $H^1$ exactly when $\chi_0(p)\ne\chi_1(p)$ (\HA{cl:chi}(3)), which on $H^0$ of
  the circle model it cannot. The Eisenstein eigenvalue $1+q^{n+1}$ on $[y^n]\in X_5$, read in
  the coordinate $\Phi$, is the oracle.
\item\label{add:cone} \emph{Equivariant mapping cone} of $r^\bullet$: the algebraic cone of
  \HA{def:pair} and \HA{lem:conerelative} (re-homed into [C] from the now-frozen [v5] \S8) for
  \texttt{Equivariant\allowbreak Cochain\allowbreak Complex\allowbreak Morphism\allowbreak OnK}, giving $X_3,X_6$ with bases and the maps of
  \HA{eq:les} as matrices. $r^\bullet$ is Hecke-equivariant on cochains only \emph{up to
  homotopy}, so $(\Theta_g,\Theta_g^\partial)$ is not a cochain map of the cone in general.
  \textbf{Open Question~2 is settled} (verdict
  \texttt{docs/implementation/\allowbreak hecke\_extensions\_\allowbreak OQ2\_\allowbreak verdict.md},
  Appendix~\ref{app:oq2}); the binding text is amendment C-A3 of the implementation plan, which is
  the mechanism by which [C] gains it. In outline:
  \begin{itemize}
  \item \emph{The datum is a triple over $\Z$}, one level below the cochains, where $\tau$ lives:
    the bulk $\tau$ of \HA{def:tau}; a $g$-twisted chain datum $\tau_\partial$ on
    $\partial X:=\Ga\cdot i(K_\partial)$, the $\Ga$-orbit of the cusp line, with the same
    $\Ga''=\Ga_0(q)$, the same $s_i$ and the same $\delta_i$ as the bulk; and a \emph{pair homotopy}
    $\eta:C_\ast(\partial X)\to C_{\ast+1}(X)$ satisfying (T1) and
    $\partial\eta+\eta\partial=j\tau_\partial-\tau j$, certified exactly like $\tau$ (\HA{obl:tau}).
    $\tau_\partial$ and $\eta$ exist over $\Z$ for \emph{every} $\ell$: every simplex of
    $\partial X$ has stabiliser $\{\pm I\}$, which acts trivially, so no stabiliser order is ever
    inverted --- unlike \HA{lem:tauexists}, which needs $\ell\ge5$ for the elliptic stabilisers.
    The bulk $\tau$ is the canonical table of \HA{alg:tau}, integral exactly when no fill used the
    mod-$\ell$ fallback. \PROVED\
    The canonical $\tau_\partial$ is the table of \HA{alg:taubdry} (v1.4.0; shipped since
    slot~1): the nearest-line-vertex rule with ties to the larger real part, over $\Z$ with
    coefficients in $\{0,\pm1\}$ for every $\ell$; other admissible choices are homotopic
    (\HA{lem:taubdryexists}(3)). It is \emph{not} $\iota$-symmetric (\HA{rmk:taubdrynoiota};
    \ref{add:iota}).
    The $\Ga$-level $\tau_\partial$ on $\partial X$ is the \emph{normative} boundary datum, and the
    one-dimensional $\tau_\infty$ of \ref{add:bdry} is its Shapiro transcription and a validation route
    only (owner decision, 2026-09-15; Appendix~\ref{app:oq2}).
  \item \emph{The cochain homotopy is derived, never chosen freely}: by the coset transport of
    \HA{thm:welldef}(b),
    $(hc)(\sigma):=\sum_i\rho_n(\delta_i)\,\varphi_c\bigl(\eta(s_i\cdot i(\sigma))\bigr)$, giving
    $h^i:C^i_{\mathrm{eq}}(K)\to C^{i-1}_{\mathrm{eq}}(K_\partial)$ with
    $dh+hd=\Theta_g^\partial r-r\Theta_g$, and the cone map
    $\Phi_g(a,b)=(\Theta_g a,\ \Theta_g^\partial b+ha)$ on
    $C(r)^i=C^{i+1}_{\mathrm{eq}}(K)\oplus C^i_{\mathrm{eq}}(K_\partial)$ with $D(a,b)=(-da,\,db+ra)$.
    Every map of \HA{eq:les} is then equivariant automatically. \PROVED\
    The sign convention is fixed once and for all (owner decision, 2026-09-15;
    Appendix~\ref{app:oq2}): $\partial\eta+\eta\partial=j\tau_\partial-\tau j$, so $h=\tilde\eta$
    carries \emph{no} sign. The identity $dh+hd=\Theta^\partial_gr-r\Theta_g$ and the cone map are
    not moved by that choice --- they are fixed by $D$.
  \item \emph{The ambiguity, and why it is not one here.} Two cochain homotopies satisfying that
    identity differ by a degree-$(-1)$ chain map $k$, and the two cone maps differ on $X_3=H^1_c$ by
    the composite $X_3\to X_4\xrightarrow{\,k_\ast\,}X_2\to X_3$, i.e.\ by an element of
    $\Hom(\ker r^1,\coker r^0)$ --- at $(n,\ell)=(10,691)$ a space of dimension $3$. For a
    \emph{free} cochain-level $h$ \emph{every} element of that $\Hom$ occurs \emph{over a
    field} (the realisation argument splits the cochain groups, so it is stated for
    $R=\F_\ell$ and not claimed over $\Z/\ell^k$ or $\Z/m$), so the operator on $X_3$ would be
    genuinely undetermined and $X_3$ would be lost as a Hecke module. \PROVED\
    For an $\eta$-derived $h$ it vanishes identically: since $X$ is contractible and
    two-dimensional (\STD, Borel--Serre), so that $Z_1(X)=B_1(X)$ and $Z_2(X)=0$, two pair
    homotopies differ by $\partial M-M\partial$ for a \emph{unique} $M:C_0(\partial X)\to C_2(X)$,
    and
    $\Psi(a,b):=(0,-\widetilde M a)$ is then an integral cone homotopy between the two $\Phi_g$.
    Hence $T_q$ on $H^\bullet_c$ is independent of $\eta$ over $\Z$ for every $\ell$; independent of
    $(\tau,\tau_\partial)$, of the $s_i$ and of $g$ when $6$ is invertible in the coefficient ring;
    and equal there to the operator of the Hecke correspondence on the Borel--Serre pair, so the
    Hecke relations hold on $H^\bullet_c$ on the nose. \PROVED
  \item \emph{What does not decide it.} The on-the-nose case is exactly $\eta\equiv0$, i.e.\
    $\tau j=j\tau_\partial$ (a strict map of pairs; one can be constructed on this model, but the
    shipped $\tau$ is not strict --- measured on every cell, v1.4.0: \texttt{cochain\_commutes}
    and \texttt{eta\_zero} are both false on $315/315$ $(\text{cell},p,\nu)$ pairs of
    \texttt{HeckeBoundaryParabolicValidation}'s grid, the cusp-support walk finding interior
    support on $15$ of the $21$ boundary slots of the canonical $\tau$ at $p=2$ ($10$ of $12$
    vertex slots, $5$ of $9$ edge slots), so $h=0$ is admissible nowhere and the $\eta$-table of
    E5 is genuinely needed; the slot counts [C] \HA{rmk:conestrict}(1) quotes are those of the
    pre-re-baseline table, a contract finding for the next status bump). The entrywise cochain identity $r\Theta_g=\Theta_g^\partial r$
    alone does \emph{not} license $h=0$ --- it can hold while $\tau j\ne j\tau_\partial$, and then
    $h_\eta$ is a chain map that need not induce zero on $X_3$ (whether it actually does for some
    $(\tau,\tau_\partial)$ is undecided, \CONJ) --- so it is recorded as a diagnostic and gates
    nothing. It
    is $\tau$, not $r^\bullet$, that one could change to obtain strict commutation. \PROVED
  \end{itemize}
  All of \ref{add:cone} transfers verbatim to a free $\Ga$-model of the pair (the Borel construction
  $(X,\partial X)\times E\Ga$, or an isotropy-resolution model), whose equivariant cohomology is the
  orbifold cohomology at $\ell\in\{2,3\}$ and of which the 7-vertex complex is the specialisation
  where $6$ is invertible; there existence and homotopy uniqueness of $\tau,\tau_\partial,\eta$ hold
  over $\Z$ for every $\ell$. \STD\
  That general model enters the binding contract \emph{now, as definitions}, with its construction
  deferred (owner decision, 2026-09-15; Appendix~\ref{app:oq2}) --- which is what makes
  ``the 7-vertex complex is a specialisation'' a theorem rather than a slogan.

  \emph{The boundary group, and what it costs at $\ell=2$.} The boundary group is
  $\Ga_\infty=\langle T\rangle\times\langle-I\rangle$ and
  $H^\bullet(\partial;M):=H^\bullet(\Ga_\infty;M)$ (\HA{sec:six}; owner decision, 2026-09-15,
  \emph{amending} the verdict, which made no recommendation on this point). Since
  $\Ga_\infty\cong\Z\times\Z/2$ and $\langle T\rangle\cong\Z$ has cohomological dimension~$1$,
  Lyndon--Hochschild--Serre reduces to
  \[
  0\longrightarrow H^{q-1}(\Z/2;M)_T\longrightarrow H^q(\Ga_\infty;M)\longrightarrow
  H^q(\Z/2;M)^T\longrightarrow0 .
  \]
  When $2\in R^\times$ the terms with $q\ge1$ vanish and this returns $H^0=M^T$, $H^1=M_T$,
  $H^q=0$ for $q\ge2$: the circle model of \ref{add:bdry}. At $\ell=2$, however, $-I$ acts trivially and
  $H^q(\Z/2;M)\cong M$ for \emph{every} $q\ge0$, so $H^q(\Ga_\infty;M)\ne0$ for every $q\ge0$ ---
  $H^\bullet(\partial;M)$ is \emph{unbounded above in degree}. The bulk is no better (the
  $12$-torsion argument of \HA{cl:amalgam} needs $\ell\ge5$), and hence neither is the cone:
  at $\ell=2$ the sequence \HA{eq:les} is a \emph{window} of a long exact sequence that does not
  terminate, and $X_2,X_3,X_5,X_6$ are the degree-$\le2$ members of infinite families. What the
  implementation is to compute there --- a degree bound and a truncation policy, and what
  ``exactness of \HA{eq:les}'' asserts at the top of the truncation --- is \emph{decided in the
  contract} ([C] v1.3.0, \HA{open:ell2}, policy \HA{oc1a}: $q_{\max}$ a recorded run parameter,
  default $2$, cohomology reported only in degrees $\le q_{\max}$), and its construction is
  deferred with the free $\Ga$-model above; nothing is decided here. Under this document's standing hypothesis $\ell\ge5$ (\S\ref{sec:conv}) none of
  this is in force.
\item\label{add:iota} \emph{The involution $\iota$}: the 7-vertex model is symmetric under
  $z\mapsto-\bar z$ ($0\leftrightarrow3$, $1\leftrightarrow2$, $4\leftrightarrow5$, $6$ fixed),
  $\iota T\iota^{-1}=T^{-1}$, $\iota S\iota^{-1}=S^{-1}$, $\iota(ST)\iota^{-1}=(TS)^{-1}$, and
  $\rho_n(\iota)=\mathrm{diag}((-1)^{n-j})_j$ (note $\iota\notin\De$; this is a separate
  automorphism, not a Hecke operator). Implementable as an
  \texttt{Equivariant\allowbreak Cochain\allowbreak Complex\allowbreak Morphism\allowbreak OnK} with a simplicial \emph{automorphism} and
  orientation signs; $[\iota,T_q]=0$ is a runtime assertion. \emph{The boundary datum and
  $\iota$ (v1.4.0).} The bulk $\tau$ is $\iota$-symmetric by construction (\HA{alg:tau} step
  (1b) breaks ties through the $\iota_K$-fixed vertex $v_6$; \HA{obl:iota}); the boundary datum
  $\tau_\partial$ is \emph{not}, and at $p=2$ cannot be made so by any $T$-equivariant vertex
  rule, since no vertex of the $\infty$-line is $\iota$-fixed (\HA{rmk:taubdrynoiota},
  [C] v1.4.1; the conjugate datum is admissible and homotopic, so $\Theta^\partial_g$ on
  cohomology is unaffected, and nothing records or asserts the symmetry). If the sign split of
  $X_2,X_5$ or of the cone ever needs $\iota\tau_\partial\iota=\tau_\partial$ at chain level,
  that is a \emph{contract addition} (owner close-out decision (a) of slot~1), not an
  assumption of this document. $\Htot=H^+\oplus H^-$ halves every
  matrix and, as \S\ref{sec:tests} shows, is what makes the Ramanujan block cyclic on each
  sign. \emph{The sign (owner decision OQ3, resolved 2026-09-16; [C] v1.2.0
  \HA{def:iota}(iv), \HA{rmk:iotasign}).} The involution of [C] is
  $\iota^{\ast}:=-\iota^{\ast}_{\mathrm{geo}}$, the \emph{negative} of the geometric pullback
  $(\iota^{\ast}_{\mathrm{geo}}f)(\sigma)=\rho_n(\iota)f(\iota_K\sigma)$, so that $H^+$ is
  \emph{by definition} the eigenspace containing the Eisenstein classes: the v1.1.1 run of
  \texttt{HeckeIotaSignValidation} had found them in the $(-1)$-eigenspace of the geometric
  involution on every cell of a $77$-cell grid, and the owner negated the definition rather
  than the record. Consequences to keep in mind when reading $\pm$ here: constants are odd
  ($\iota^{\ast}=-1$ on $H^0(\Ga;\F_\ell)$), $H^0(\partial)$ is odd and $H^1(\partial)$ even
  for $\ell>n$, $\iota^{\ast}$ is the geometric involution twisted by $\det$ and, for $n$
  even, not a pullback along any lift of $z\mapsto-\bar z$, and at $\ell=2$ the two coincide. Every
  $H^\pm$ of this document is in terms of $\iota^{\ast}$ so defined
  (\texttt{docs/implementation/oq3\_resolution.md}).
\end{enumerate}
Recommended once: two independent models already run in parallel in
\texttt{HeckeCrossCheckValidation} --- the amalgam / \texttt{BarCocycleH1} route and the 7-vertex
complex, the latter standing in for the $\Ga(2)$/$S_3$-invariants model of \Mnote{1}, which does
not exist in the tree (inventory) --- and the identities of \S\ref{sec:identities} become one
validation.

\subsection*{Stage 1 --- base operators}
Compute $T_q$ for $S_0=\{2,3,5,7,\ell\}$ (and $\iota$) on all available $X_i$. Cost per
$T_q$: the $\tau$-table has $(q+1)\cdot N_\bullet$ entries and $\Theta^{(q)}$ is applied to $d_q$
basis cocycles, i.e.\ $O\bigl((q+1)\,n^2\,d_q\bigr)$ multiplications through the tape for
the $\Theta$-application alone, \emph{plus the fixed cost of the cohomology run} that produces
the bases, which dominates the measured wall ($T_2$ at $n=128$ in $0.12$\,s against $0.09$\,s
for the cohomology alone; at $n=1000$, $14.4$\,s against $14.8$\,s, \texttt{CLAUDE.md}; v1.7.0:
the $n=1000$ pair has the whole \emph{below} its part, so those two walls differ by less than
their run-to-run noise --- it supports ``dominates'' and must not be read as a split).
Run immediately: $[T_q,T_{q'}]=0$, $[T_q,\iota]=0$, equivariance of every map in \HA{eq:les},
the boundary characters of \HA{cl:chi}, \HA{cl:dimH1}. Failures are convention or
coset bugs, not mathematics.

\subsection*{Stage 2 --- maximal ideals}
\begin{algorithm}\label{alg:stage2}
Per $X_i$ (and per sign once \ref{add:iota} exists):
\begin{enumerate}
\item Run \texttt{HeckeModuleDecomposition::ComputeMultiOp} on $\{A(T_q)\}_{q\in S_0}$ per
  $X_i$ (per sign): blocks $V$ with per-operator $\pi_q^{e_q}$, partitions, lower bound $f$.
  Cost: Berkowitz $O(d^4)$ per operator plus factorisation; negligible against Stage 1.
\item \emph{Generic refinement} (new, Claim~\ref{cl:gap}): for each block, for the
  deterministic witness sequence $w_c=\sum_jc_jt_j$ ($c$ from the fixed-seed search already in
  S11), compute the primary decomposition of $w_c$ on $V$ and split $V$ accordingly; repeat
  until every block passes the locality certificate below, or until a fixed attempt budget
  (the setting \texttt{locality = budget=<N>}, default $N=64$ candidates per block, counting
  the single-operator witnesses, the all-ones combination and the seeded combinations as
  \emph{one} sequence with one attempt index) is exhausted, in which case the block is recorded
  with \texttt{"local": false}, its \texttt{residue\_degree\_lower\_bound}, and
  \texttt{final = false} with reason \texttt{budget\_exhausted:locality} and
  \texttt{resumable.complete = false} --- the loop never spins. (v1.4.0. The reason is not
  \texttt{not\_local}: a block \emph{proved} non-local splits and never survives the loop, so a
  survivor is a block whose locality is \emph{unknown under the budget}; slot-2 WP-1 \S7
  item~3.) Under \texttt{locality = strict} (owner decision OQ16; never the default, never run
  by the golden harness, byte-identical to budget mode wherever budget mode certifies) the loop
  continues after the $N$ linear witnesses with the terminating route of
  Certificate~\ref{sec:localcert} and emits \texttt{local = true} on every block,
  \texttt{local = false} only by interruption at a step boundary (reason \texttt{interrupted},
  resumable). Strict mode is specified here, parsed since slot~2 and \textbf{implemented
  since slot~3} (v1.4.2; slot-3 WP-3, \texttt{HeckeModuleDecomposition::StrictLocalityRoute}):
  it no longer throws \texttt{not\_implemented\_in\_this\_build}, and re-running the whole
  golden grid under it reaches the terminating route on no block.
\item Read off the character $a(\m)=(P_q(\bar X))_{q\in S_0}\in(\Fl[X]/(\pi_w))^{S_0}$
  (Claim~\ref{cl:resfield}) and
  compare with the atlas (Stage 6): label \emph{cuspidal} (lowest-weight representative and
  $\theta$-twist), \emph{Eisenstein-type} (all $a$ with $\chi_a\equiv a(\m)$), or \emph{both}.
  If the atlas certifies that $S_0$ separates $\Sigma(\ell)\cup\{\chi_a\}$ the decomposition is
  complete on the anemic side; otherwise add $S_{\mathrm{sep}}(\ell)$ to $S_0$ and repeat.
  \emph{In the default run} (v1.8.0-pre.1; slot-7 owner ruling~3b~(B)) $\Sigma(\ell)$ exists
  only where it is built in-process, $\ell\in\{5,7,11,13\}$: there $S_{\mathrm{sep}}(\ell)$ is
  added to $S_0$ (a run from the start on $S_0\cup S_{\mathrm{sep}}(\ell)$ is that repeat) and
  the label is read off. At every other $\ell$ this step does not run, no stored atlas is read,
  and the record says so --- \texttt{s\_sep} and every \texttt{label} are the placeholder
  \texttt{\{"status": "atlas\_unavailable"\}} of \S\ref{sec:schema} --- never a silent
  omission and never a label from a different run.
\end{enumerate}
\end{algorithm}

\begin{certificate}[Locality]\label{sec:localcert}\PROVED\
A block $V$ with generators $t_j$ is $\m$-primary for $\T_S$ with $[\T_S/\m:\Fl]=f$ iff there is
$w\in A_S$ with $\pi_w:=$ squarefree part of $\min_V(w)$ irreducible of degree $f$, and, with
$\hat w:=w^{\ell^N}$ ($f\mid N$, $\ell^N\ge D$; Claim~\ref{cl:resfield}), for every $j$ a
polynomial $P_j\in\Fl[X]_{<f}$ with $n_j:=t_j-P_j(\hat w)$ nilpotent on $V$. Verification:
$\mathtt{CertifyIrreducible}(\pi_w)$; $\hat w$ by $N$ successive $\ell$-th powers; $P_j$
\emph{deterministically}: take any $0\ne v\in U:=\ker\pi_j(t_j)\cap\ker\pi_w(w)$ (an
$A_S$-submodule containing $\soc V$, hence non-zero when $V$ is local). On $U$ both $t_j$ and
$w=\hat w$ are semisimple and commute, so if $V$ is local, uniqueness of the Jordan--Chevalley
decomposition of $t_j|_U=P_j(\hat w)|_U+n_j|_U$ forces $n_j|_U=0$, i.e.\ $t_jv=P_j(\hat w)v$: a
consistent $D\times f$ linear system with a unique solution ($v,wv,\dots,w^{f-1}v$ are
independent because $\Fl[w]|_{\ker\pi_w(w)}$ is a field). Then verify $n_j^{D}=0$ on $V$ by
squaring, $O(D^3\log D)$ per $j$. Inconsistency or non-nilpotency proves that $V$ is not local:
refine (Algorithm~\ref{alg:stage2}(2)). Sufficiency: $A_S/\rad$ is generated by $\bar w$ hence
is $\Fl[X]/(\pi_w)$, a field. Necessity: a generic element of a local $A_S$ is such a $w$; more
precisely the \emph{linear} witnesses $\sum_jc_jt_j$, $c\in\Fl^{k}$, contain one whenever $f$ has
at most $\ell$ distinct prime factors (the residues $\bar t_j$ generate $\Flr$ as a field, so
their $\Fl$-span lies in no maximal subfield, and an $\Fl$-space is not a union of $\le\ell$
proper subspaces --- automatic for $\ell\ge5$, $f<30030$), and they separate two
non-Galois-conjugate characters of residue degrees $f,f'$ whenever $\min(f,f')\le\ell$ (the
same covering argument over the Galois orbit). Outside these ranges the fixed-seed search may
fail, which is why Algorithm~\ref{alg:stage2}(2) is bounded and reports \texttt{"local": false}
rather than assuming.

\emph{Strict mode (v1.4.0; [X] \S3.1, transcribed; owner decision OQ16).} When the $N$ linear
witnesses fail, locality is decided and a witness found by linear algebra with no search, in
three steps on $A_S$ computed as in Algorithm~\ref{alg:stage3}(1).
(1) \emph{The radical.} On a commutative $\Fl$-algebra the Frobenius $x\mapsto x^\ell$ is
$\Fl$-linear, so $\rad A_S=\ker(\Frob^N)$ with $\ell^N\ge D$, the kernel of an explicit
$\dim A_S\times\dim A_S$ matrix ($N$ successive Frobenius matrices, each column
$b_i\mapsto b_i^\ell$ by square-and-multiply on $D\times D$ matrices). The trace form is not
usable here: it vanishes identically on $\Fl[x]/(x^\ell)$, a local algebra of dimension $\ell$.
(2) \emph{The Berlekamp subalgebra decides and splits.} $\bar A:=A_S/\rad A_S\cong\prod_{i=1}^m\F_{\ell^{f_i}}$
with $m$ the number of maximal ideals of $A_S$ on $V$, and the Frobenius-fixed subalgebra
$B:=\ker(\Frob_{\bar A}-\mathrm{id})=\Fl^m$ has $\dim_{\Fl}B=m$. So $V$ is local iff $\dim B=1$.
If $\dim B\ge2$, the first non-scalar RREF basis element $\bar b$ of $B$ (a fixed choice) lifts to
$w\in A_S$ whose semisimple part has $\ge2$ distinct eigenvalues in $\Fl$, and the primary
decomposition of $w$ on $V$ splits $V$ into $\ge2$ \emph{non-zero} blocks --- $V$ is faithful
over $A_S$, so no primitive idempotent of $\bar A$ lifts to an element acting as $0$ ---
inserted at the parent's position in \texttt{CanonicalLess} order of the eigenvalues; recurse.
\emph{Termination}: every split strictly reduces $\dim V$. No witness is searched here.
(3) \emph{A residue witness over the whole residue field.} With $\dim B=1$, $\bar A=\F_{\ell^f}$
is a field, and a field is never the union of its proper subfields
($\sum_{r\mid f\text{ prime}}\ell^{f/r}\le\omega(f)\,\ell^{f/2}<\ell^f$ for $f\ge2$, $\ell\ge2$),
so an element of $A_S$ whose residue generates $\bar A$ exists although the $\Fl$-span of the
$\bar t_j$ may lie in a union of proper subfields; it is drawn from $\bar A$'s own coordinate
vectors (lifted to $A_S$ by the same coordinates) by the fixed-seed sequence, each candidate
tested by $\deg\min_{\bar A}(\bar w)=f$ --- an $f\times f$ Krylov computation, failure
probability $\le\omega(f)\,\ell^{-f/2}$, expected trials $\approx1$ --- with the lexicographic
enumeration of the $\ell^f$ coordinate vectors as the deterministic, terminating fallback; then
$\pi_w$, $N$, $\hat w$, the $P_j$ and the nilpotency checks exactly as above, which succeed
because $V$ is local. Strict mode changes \emph{how the witness is found}, never what is
certified: the certificate is the same. The covering example is $(\ell,f)=(2,30)$:
$V=\F_{2^{30}}$ as a $30$-dimensional $\F_2$-space, $t_1=a\cdot$, $t_2=b\cdot$ with
$a\in\F_{2^{15}}$, $b=a+c\in\F_{2^{10}}$, $c\in\F_{2^6}$, so that $\F_2(a,b)=\F_{2^{30}}$ and
$A_S$ is a field ($f=30=2\cdot3\cdot5$, three prime factors $>\ell$), while the three non-zero
linear witnesses $a,b,a+b$ have residue degrees $15,10,6$: budget mode emits
\texttt{local = false} with \texttt{residue\_degree\_lower\_bound = 30} (the correct budget-mode
verdict) and strict mode \texttt{local = true}, $f=30$ ([X] \S3.2; the strict-mode arms of
\texttt{HeckeLocalityValidation}, slot 3 --- \emph{verified} (v1.4.2, slot-3 WP-3): built from
a certified degree-$30$ irreducible, budget mode \texttt{local = false} with bound $30$, strict
mode \texttt{local = true}, $f=30$, \texttt{phase\_reached = residue\_search}, in $0.10$\,s). Cost $O(\dim A_S\cdot D^3\cdot N\log\ell)\le
O(D^5\log D\log\ell)$ per block, entered only after $N$ failed linear witnesses --- on no
computed cell so far. What strict mode records: \texttt{level\_c.locality} $=\{$\texttt{mode:
strict}, \texttt{budget}, \texttt{phase\_reached} $\in\{$\texttt{linear}, \texttt{berlekamp},
\texttt{residue\_search}$\}\}$ and \texttt{local = true} on every block; the golden
\texttt{decomp\_*} hashes pin the budget mode, and strict mode must reproduce them bit for bit.
\emph{Per block and per run} (v1.5.0; slot-3 WP-3 finding~6): \texttt{phase\_reached} is the
phase at which \emph{that} block's verdict was reached, and the run's high-water mark over
all blocks --- split parents included, whose phase belongs to no surviving block --- is the
separate key \texttt{phase\_reached\_high\_water} of \S\ref{sec:schema}.

\emph{What a \textbf{reported} witness must satisfy (v1.5.0; slot-3 WP-3 finding~3).} The
sufficiency argument above reads ``$A_S/\rad$ is generated by $\bar w$'', which presumes
$w\in A_S$. That is automatic for a \emph{linear} witness $w=\sum_jc_jt_j$, which a record
reports as the coefficient list $(c_j)$ and a verifier re-forms inside $A_S$; it is
\textbf{not} checkable from a witness reported as a \emph{matrix}, which strict mode's
step~(3) produces (a polynomial in the $t_j$, not a linear combination, so no coefficient
list can carry it). A verifier working from the reported data alone therefore checks the
weaker hypothesis \textbf{$w$ commutes with $\{t_j\}_{j\in S}$}, and that hypothesis still
yields locality: with $\pi_w$ irreducible of degree $f$, $K:=\Fl[\hat w]\cong\F_{\ell^f}$
is a field, the $n_j=t_j-P_j(\hat w)$ are nilpotent, and they commute with $\hat w$ and
with each other (each $t_j$ does, and $w$ does by hypothesis), so $A_S\subseteq K[n_j:j]$,
a local ring whose maximal ideal is its nil radical; hence $A_S$ meets that ideal in
nilpotents only and $A_S/\rad A_S$ embeds in $K$ as a finite domain, i.e.\ is a field, and
$A_S$ is local. \textbf{Conformance}: a record reporting a matrix witness must carry it as
such and the commutation is a certificate clause of its own; a record reporting a linear
witness keeps the coefficient list. Implemented as \texttt{PrimaryBlock::witness\_matrix}
with \texttt{certify()} code~30 (slot-3 WP-3), empty in budget mode, so no published value
moves.

\emph{The sub-block order after a split is per mode, not intrinsic} (v1.5.0; slot-3 WP-3
finding~2; owner confirmation of 2026-09-21, item~1(c)). Both modes insert the sub-blocks
at the parent's position in \texttt{CanonicalLess} order \emph{of the splitting element's
factors}, and both are deterministic; but the two modes split by \emph{different} elements
--- a linear witness and the Berlekamp lift of strict mode's step~(2) above ---
so the correspondence factor $\leftrightarrow$ maximal ideal, and with it the resulting
\emph{order} of the list, may differ (measured: $9$ of $9$ cells agree as an unordered
partition with bases entrywise equal, $6$ of $9$ also in order). \textbf{What is
mode-independent is the partition, not the list.} A statement that two modes produce ``the
same blocks'' is therefore a statement about the partition; a \texttt{decomp\_*\_U} hash,
which reads the list, is a function of the mode and is golden only at the default. An
intrinsic order --- by each sub-block's own $\pi_w$, say --- was considered and
\textbf{declined}: it would change the split path of \emph{both} modes, and the default
mode's off-grid order with it, for no consumer that exists.
\end{certificate}

\begin{claim}[The residue degree of a local block, and the sharpness of the budget-mode bound]\label{cl:sharpbound}\PROVED
Let $V$ be a block on which $A_S=\Fl[t_p:p\in S]$ is local, and let $\pi_p$ be the
squarefree part of $\min_V(t_p)$. Then
\[
  \dim_{\Fl}\bar A\;=\;\operatorname{lcm}_{p\in S}\deg\pi_p ,
  \qquad \bar A:=A_S/\rad A_S .
\]
Consequently the residue degree $f=[\T_\m/\m:\Fl]$ of a \emph{local} block equals
$\operatorname{lcm}_p\deg\pi_p$, which is exactly the quantity budget mode reports as
\texttt{residue\_degree\_lower\_bound} when its witness family is exhausted: \textbf{on a
block that is in fact local, that lower bound is sharp}.
\end{claim}

\begin{proof}
$A_S$ local means $\bar A$ is a field, and it is generated over $\Fl$ by the residues
$\bar t_p$, so $\bar A=\Fl(\bar t_p:p\in S)$ is the compositum inside $\bar A$ of the
subfields $\Fl(\bar t_p)$. The minimal polynomial of $\bar t_p$ over $\Fl$ is $\pi_p$: the
squarefree part generates the radical of $(\min_V(t_p))$, so $\pi_p(t_p)$ is nilpotent and
$\pi_p(\bar t_p)=0$, while $\pi_p$ is irreducible in $\Fl[X]$ --- if
$\min_V(t_p)=\prod_if_i^{e_i}$ had two distinct $f_i$, the Chinese remainder theorem would
make $\Fl[t_p]\cong\prod_i\Fl[X]/(f_i^{e_i})$ and so exhibit a non-trivial idempotent
\emph{in $A_S$}, which a local ring does not have (v1.5.1: through v1.5.0 this step read
``a proper factorisation would give a non-trivial idempotent decomposition of the field
$\Fl[\bar t_p]$'', which it does not --- inside a field a proper factorisation of $\pi_p$
says only that $\min(\bar t_p)$ is a proper divisor of $\pi_p$, i.e.\ the very alternative
to be excluded; the idempotent lives one level down, in $A_S$) --- so no proper divisor
annihilates $\bar t_p$. Hence
$\Fl(\bar t_p)=\F_{\ell^{\deg\pi_p}}$. Inside a finite field the compositum of
$\F_{\ell^{a}}$ and $\F_{\ell^{b}}$ is $\F_{\ell^{\operatorname{lcm}(a,b)}}$, and the
statement follows by induction on $|S|$. For the consequence: $\bar A=\T_\m/\m$ when the
block is $\m$-primary for $\T_S$ (Claim~\ref{cl:resfield}), and the bound budget mode
emits on an exhausted family is $\operatorname{lcm}_p\deg\pi_p$ by construction --- an
\emph{equality} exactly when the block is local, and a \emph{lower} bound on every residue
degree of a block that survives the refinement loop, because such a block has every
$\pi_p$ \textbf{irreducible} (a candidate with two distinct irreducible factors splits the
block, Algorithm~\ref{alg:stage2}(2)), whence $\deg\pi_p$ divides the degree $d_i$ of every
factor $K_i$ of $\bar A=\prod_iK_i$ and so does $\operatorname{lcm}_p\deg\pi_p$.
(v1.5.1, correcting v1.5.0's ``a lower bound in general, because a non-local block's
$\bar A$ is a product and the lcm bounds only its largest factor'': without the
irreducibility of the $\pi_p$ that is false in both directions --- for $\bar A=\Fl\times\Fl$
with $S=\{p\}$ and $t_p=(0,1)$ the lcm is $\deg X(X-1)=2$ while both residue degrees are
$1$, and for $\bar A=\F_{\ell^{12}}\times\F_{\ell^{6}}$ with $t_p$ of local degrees $4,2$
and $t_q$ of local degree $3$ in both it is $\operatorname{lcm}(6,3)=6<12$.)
(Slot-3 WP-3 finding~4, where the implementation asserts it. It is also why
the covering example of Certificate~\ref{sec:localcert} has bound $30$ and not merely
$\ge30$.)
\end{proof}

\begin{certificate}[Exhaustiveness of $\MaxSpec$]\PROVED\
$\sum_\m\dim H_\m=d$ (all blocks certified local, pairwise distinct characters up to Galois,
$U$ invertible and every $U^{-1}A(T_q)U$ block-diagonal): this is \texttt{certify()} plus the
locality certificate. Completeness against \emph{uncomputed} operators is Stage 4.
\end{certificate}

\subsection*{Stage 3 --- the local algebra and its tangent space}
\begin{algorithm}\label{alg:stage3}
For each $\m$ (matrices of size $D$):
\begin{enumerate}
\item $A_S$ as an $\Fl$-subspace of $\End(V)\cong\Fl^{D^2}$: start from $\{I\}$, multiply the
  current basis by the $t_j$, row-reduce, until stable. By commutativity the products needed are
  monomials $\prod t_j^{a_j}$ of total degree $<\nu\le D$; cost
  $O(k\cdot\dim A_S\cdot D^3)$ products and an RREF on $\dim A_S\times D^2$. Bound:
  $\dim A_S\le D$ once $V$ is known cyclic (Stage 4a), $\le D^2/4+1$ in general (Schur).
\item The coefficient field $K=\Fl[\hat w]$, $\hat w=w^{\ell^N}$, and
  $n_j:=t_j-P_j(\hat w)$ for \emph{every} $j$, the witness included ($n_w=w-\hat w$;
  Claim~\ref{cl:resfield}); then $A_S=K[n_j]_j$ and $\m_S=(n_j)_j\subseteq A_S$ as an ideal
  (span of the generators times the basis of $A_S$; no separate $\pi_w(w)$ generator), $\m_S^2$,
  $\m_S^i$; $\nu$; $\soc A_S=\Ann_{A_S}(\m_S)$ by a linear system; Gorenstein iff
  $\dim_{\Fl}\soc A_S=f$.
\item Tangent-space generation test: the images $\bar n_j$ in $\m_S/\m_S^2$, which is a
  $K$-vector space through $\hat w$; record which subset of $S$ spans it \emph{over $K$} (the
  $K$-span of $\bar n_j$ is the $\Fl$-span of $\hat w^i\bar n_j$, $0\le i<f$), and
  $\dim_{\Flr}\m_S/\m_S^2=\dim_{\Fl}(\m_S/\m_S^2)/f$. Over $\Fl$ one direction would count as $f$,
  and with $P_j(w)$ in place of $P_j(\hat w)$ the witness's own direction is lost
  (Claim~\ref{cl:resfield}). Expectation (Claim~\ref{cl:bk}):
  $\le2$ anemic directions plus possibly the $T_\ell$ direction at cuspidal $\m$; Eisenstein-type
  and the dominant $\m$ for $n\gg\ell$ are where surprises would show.
\end{enumerate}
\end{algorithm}
\noindent\emph{As implemented} (v1.4.0; slot-2 WP-4, \texttt{HeckeExt/HeckeLocalAlgebra},
\texttt{ComputeLocalAlgebra}, namespace \texttt{hecke\_ext}): $\hat w$ and the $P_j$ are
\emph{consumed} from the locality certificate on its matrix path --- the $k=1$ analytic path of
[X] \S3.1 never forms $\hat w$, so Stage~3 forces the matrix path --- with $\pi_w(\hat w)=0$
re-asserted as an identity (Frobenius is a ring endomorphism, so its failure is a bug, not a
finding) and $N$ re-derived; $\m_S=(n_j)_j$ with no separate $\pi_w(w)$ generator, since
$n_w=w-\hat w=\sum_jc_jn_j$ is automatic, and $\m_S^{i+1}=\sum_j\m_S^in_j$ (only the non-zero
powers stored); the tangent test over $K$ is the $\Fl$-span of $\{\hat w^i\bar n_j:0\le i<f\}$
modulo $\m_S^2$ in one row reduction, \texttt{minimal\_active} greedy in \emph{ascending} $p$
(hence deterministic), $\mathrm{embdim}=\dim_{\Fl}(\m_S/\m_S^2)/f$; the output carries both the
$\dim A_S\times D^2$ coordinate matrix and the basis as $D\times D$ matrices, so membership is
one solve and a certificate compares entrywise (the RREF of a span is unique). A non-local input
is \emph{reported}, never thrown --- E0b's loop is what acts on it. Stage~2's cross-block
evaluation test is \texttt{SameMaximalIdeal}: blocks $B,B'$ carry the same $\m$ iff, with $B$'s
own $c$, $\pi_w$ and $\{P_j\}$ and $N'$ the least multiple of $f$ with $\ell^{N'}\ge\dim B'$
(taken for the \emph{tested} block), $\pi_w(w'|_{B'})$ and every $t_j|_{B'}-P_j(\hat w')$ are
nilpotent; it is what the exhaustiveness certificate of E2 calls for pairwise distinctness.

\subsection*{Stage 4 --- certifying $A_S=\T_\m$}\label{sec:certgen}
\begin{certificate}[(a) Cyclicity]\label{cert:cyclic}\PROVED\
If $\dim_{\Fl}V/\m_SV=f$ (Nakayama: $V$ is cyclic over $A_S$; deterministic --- no random
vectors needed; \emph{some} $v\notin\m_SV$ generates, and every such $v$ does) then
$\End_{A_S}(V)=A_S|_V$, and since every
$T_q$ commutes with $A_S$ and preserves $V$ (it commutes with the polynomials in $A_S$ that cut
$V$ out), $\T_\m=A_S|_V$: \emph{no further operator can change the module structure at $\m$}.
Cost $O(kD^3)$. Cyclicity fails exactly at multiplicity $>1$ (e.g.\ every cuspidal $\m$ on
$X_4$ without the $\iota$-split, \HA{cl:dimH1}: holomorphic and antiholomorphic copies),
which is why \ref{add:iota} is on the list.

\emph{The generator is canonical} (v1.4.1; stated here rather than only recorded, which is what
v1.4.0's implementation note asked for). Which generator a certificate \emph{reports} is not
free: a record is verified from the data it carries, so the choice must be a function of the
data. \textbf{The generator of this certificate is the least standard basis vector of $V$ that
lies outside $\m_SV$} --- least in the fixed index order of the stored basis of $V$, which
exists because $\m_SV\subsetneq V$ by Nakayama. It is canonical given $(V,\m_S)$ and the stored
basis, it needs no search beyond one pass over $D$ vectors, and a verifier recomputes it and
compares. Any other $v\notin\m_SV$ generates equally well and the certificate is \emph{true} for
it; but a record emitting one is non-conforming, and \texttt{verify()} rejects it
(clause~13 of its fifteen). This is strictly stronger than ``any $v\notin\m_SV$'' and is
deliberately so: it is what makes two runs on the same block emit the same certificate.
\emph{Implemented} (v1.4.0; slot-2 WP-4, \texttt{HeckeExt/HeckeLocalAlgebra}).
\end{certificate}
\begin{proof}
$V=A_Sv$; an $A_S$-endomorphism $\phi$ is determined by $\phi(v)=av$, $a\in A_S$, and then
$\phi=a\cdot$. Commutation and invariance as stated.
\end{proof}

\begin{certificate}[(b) Commutant bound]\label{cert:comm}\PROVED\
$C:=\End_{A_S}(V)=\{X:Xt_j=t_jX\ \forall j\}$ is the kernel of a linear map
$\Fl^{D^2}\to\Fl^{kD^2}$. Then $A_S\subseteq\T_\m\subseteq C$ and $\T_\m$ is commutative. If
$\dim C=\dim A_S$, done. Otherwise every further $T_q$ lies in the commutative subalgebra of $C$
generated by $A_S$ and the $T_{q'}$ already known; testing one new $T_q$ is a membership test in
the current $A_S$ (Stage 4c). Cost: naive $O(kD^6)$; with $w$ in rational canonical form
$V\cong\bigoplus_i\Fl[X]/(g_i)$, $\End_{\Fl[w]}(V)$ has the explicit basis of dimension
$\sum_{i,i'}\deg\gcd(g_i,g_{i'})$ and only the centralisers of the remaining $t_j$ inside it
need solving. \emph{The sharp count (v1.4.0).} The element put in rational canonical form is
not the witness $w$ (which may be semisimple, hence scalar on the block, and then gives the
naive $D^2$ unknowns) but a \emph{cyclic-most} $a\in A_S$: the first of the RREF basis elements
$b_1,\dots,b_{\dim A_S}$ and their prefix sums whose minimal polynomial on $V$ has maximal
degree --- a choice that affects the cost, never the verdict. $A_S$ being local, $a$ is
$\pi_a$-primary with invariant factors $g_i=\pi_a^{s_i}$, $i=1,\dots,m$, and the basis has
exactly
\[
u=\sum_{i,i'}\deg\gcd(g_i,g_{i'})=\deg\pi_a\sum_{i,i'}\min(s_i,s_{i'})\;\le\;mD
\]
elements, so the $k$ systems $Xt_j=t_jX$ have $u\le mD$ unknowns instead of $D^2$, cost
$O(kuD^3)$ to form and $O(kD^2u^2)$ to reduce; $u$ can reach $D^2$ only when $m\approx D$
(through v1.3.3 this read ``$O(m^2D)$ unknowns'' and carried a \VERIFY{} on the cost, discharged
by slot-2 WP-5, which implements the count and pins $\dim C$ entrywise against the naive
commutant). The system is budgeted (\texttt{kCommutantBudget}, a named constant; formation gated
up front, rows charged one by one) and resumable from the row reached
(\texttt{CommutantState}, R3 instance~3); an unaffordable shape ends \texttt{over\_budget} and
never falls back to the naive route.

\emph{The socle bound (v1.4.0; underspecified through v1.3.3 --- this certificate carried no
statement of \emph{when} (b) can succeed; [X] \S2.2 Stage~4(b), slot-2 WP-5).} Let
$t:=\dim_\F V/\m_SV$ (the minimal number of generators, \ref{D5}) and $s:=\dim_\F\soc V$,
$\soc V=\{v:\m_Sv=0\}=\bigcap_j\ker n_j$, both in units of $f$. Then
\[
\dim_{\Fl}C\;\ge\;\dim_{\Fl}A_S+f\,t\,s-\dim_{\Fl}\soc A_S ,
\]
so \textbf{$\dim C=\dim A_S$ requires $f\,t\,s\le\dim_{\Fl}\soc A_S$}, and
$f\,t\,s>\dim_{\Fl}\soc A_S$ is a \emph{proof} that (b) fails (outcome
\texttt{excluded\_by\_socle\_bound}, the four numbers recorded). \emph{Proof.} Every
$\Fl$-linear map $V\twoheadrightarrow V/\m_SV\to\soc V\hookrightarrow V$ whose middle map is
$K$-linear is $A_S$-linear, both ends being $A_S/\m_S$-modules, so
$C\supseteq\Hom_K(V/\m_SV,\soc V)$, of $\Fl$-dimension $fts$; and
$A_S\cap\Hom_K(V/\m_SV,\soc V)\subseteq\{a\in A_S:a\,\m_SV=0\}=\soc A_S$ because $V$ is faithful.
For cyclic $V$ ($t=1$, $V\cong A_S$) the bound is an equality and (a) already certifies; for
$t\ge2$ it forces $\dim_\F\soc A_S\ge2$, i.e.\ $A_S$ \emph{not} Gorenstein. Hence (b) is never
attempted on a field, on any $\F[\varepsilon]/(\varepsilon^2)$ or on any Gorenstein $A_S$ ---
in particular not at $(10,691)$, where $t=s=2$, $f=1$, $\dim\soc A_S=1$ give
$\dim C\ge2+4-1=5$, the exact value, and not on any Kummer block --- and it can beat (a) only on
co-cyclic-like modules over a non-Gorenstein $A_S$, e.g.\ $V=A_S^\vee$ for
$A_S=\F[x,y]/(x,y)^2$ ($t=2$, $s=1$, $\dim\soc A_S=2$, $\End_{A_S}(A_S^\vee)=A_S$), which
\texttt{HeckeExtAlgebraValidation} pins. Where (b) is not attempted or is \texttt{over\_budget},
the ideal falls through to (c); a computed $\dim C>\dim A_S$ is the same proof that (b) fails.
In the default configuration an ideal that reaches closure by no route is emitted
\texttt{final = false} with reason \texttt{closure\_unreachable} (the socle bound or a computed
$\dim C>\dim A_S$, with (a) failing and (c) off its proof range) or
\texttt{budget\_exhausted:commutant}, never withheld (owner decision OQ8). \texttt{verify()}
re-derives $\dim C$ from the stored $(g_i)$ and the change of basis $R$, a quantity that is
basis-independent: the rational canonical basis is \emph{not} pinned, any valid $R$ verifies,
and a tamper must break the companion relation (an image column of $R$), not perturb $R$
(slot-2 WP-5 \S4).
\end{certificate}

\begin{certificate}[(c) Sturm sweep --- fallback]\label{cert:sturm}\PROVED\ \emph{given a recorded $B'(k,\ell)$, on the range of Claim~\ref{cl:sturmmodl}}
First compute $B'(k,\ell)$ (Definition~\ref{def:bprime}: Hecke matrices on the Miller basis of
$M_k(\Z)$ reduced mod $\ell$, then the algebra closure --- independent of the cohomology
pipeline and trivial in cost) and record it with its certificate (the cap, the closure
dimension reached, $\dim M_k$). If it does not exist below the cap (\texttt{none\_at\_cap};
equivalently $\ell\mid[\T^{\mathrm{sat}}:\T_\Z(M_k)]$ once the cap is large enough), (c) is
\emph{unavailable}: in the default configuration \texttt{sturm\_unavailable = nonfinal} the
ideal is \texttt{final = false} with reason \texttt{sturm\_unavailable}; under
\texttt{sturm\_unavailable = heuristic} the sweep runs to the cap anyway and the ideal is
\texttt{final = "heuristic"}, with the primes actually swept in \texttt{sturm\_primes} and the
mode in \texttt{sturm.mode} --- never in the default configuration (v1.4.0; slot-1 owner
confirmation~6, [X] \S1.3.1 R3.2).
Otherwise, for all primes $q\le B'(k,\ell)$ and $q\in S_{\mathrm{sep}}(\ell)\cup\{\ell\}$
(Claim~\ref{cl:sep}), compute $T_q$ (Stage 1 cost each; embarrassingly parallel over $q$ and
over $\tau$-table entries --- the natural GPU stage), restrict to $V$ (a projection already at
hand), and test $T_q|_V\in A_S$ by a linear solve against the basis of $A_S$
($O(\dim A_S\cdot D^2)$). This \emph{proves} $A_S=\T_\m$ on $X_{4,\m}$ for every
non-Dickson-type $\m$ (Claim~\ref{cl:sturmmodl}); on Dickson-type $\m$ and on $X_2,X_3,X_5,X_6$
the sweep proves nothing, whatever its margin, and the record says so (v1.4.0):
\texttt{"certificate": "sturm(B\textquotesingle)"} versus \texttt{"cyclic"} /
\texttt{"commutant"}, and \texttt{final = false} with reason \texttt{sturm\_heuristic} under the
default \texttt{final\_mode = closure}, or \texttt{final = "heuristic"} with
\texttt{sturm\_primes} under the user-selected \texttt{final\_mode = sturm\_heuristic}. A
$U_\ell$ omitted by the coset budget (\texttt{u\_ell.reason = u\_ell\_over\_budget}) never makes
a proof-range record non-final: the proof list is the primes $\le B'$. With
$B(k)=\lfloor k/12\rfloor+1$ in place of $B'$ the sweep would certify $A_S=\T_\m$ \emph{falsely}
at $(n,\ell)=(22,5)$ and $(26,7)$, on the pipeline's own $H^1$ (\S\ref{sec:tests}). Number of
operators $\approx B'/\ln B'$.
\end{certificate}

\begin{certificate}[(d0) Dickson --- the $X_1$ blocks]\label{cert:dickson}\PROVED\ \emph{(v1.13.0;
Proposition~D of the slot-12 finality proposal, \texttt{docs/implementation/slot12/WP-12\_finality\_proposal.md}
v2 \S3.1, owner-approved 2026-09-26)}
Let $\ell\ge5$, $n\ge1$ even, $V\subseteq X_1=H^0(\Ga;\Mn)$ a certified local $\T_S$-block (Stage~2)
\textbf{with $\ell\in S$}, and let $\Lambda(n,\ell)$ be the set of classes $a\bmod(\ell-1)$ of the
Dickson monomials $L^aD^b$ of degree $n=a(\ell+1)+b(\ell^2-\ell)$. If every $t_q=T_q|_V$,
$q\in S\setminus\{\ell\}$, is a scalar $\lambda_q$ and \emph{exactly one} $a\in\Lambda(n,\ell)$ satisfies
$(1+q)q^{a}=\lambda_q$ for all $q\in S\setminus\{\ell\}$, then $\T_\infty|_V=\Fl[U_\ell|_V]=A_S$, i.e.\
$A_S=\T_\m$ on $V$ in the sense of Claim~\ref{cl:gen}(b). In particular this holds whenever $S$ contains
a prime $q$ that is a primitive root mod $\ell$ with $q\not\equiv-1$ (at $\ell=5,11,13$ the prime $2$, at
$\ell=7$ the prime $3$, all in $S_0$).
The record (\texttt{hecke\_ext/v4}) writes \texttt{certificate\_routes.d0} $\in\{$\texttt{passed},
\texttt{failed:<class>}, \texttt{not\_applicable:u\_ell}, \texttt{not\_run}$\}$: \texttt{failed:<a>} names
the second class of $\Lambda(n,\ell)$ the $S$-eigenvalues do not separate from the first;
\texttt{not\_applicable:u\_ell} is written when $U_\ell$ was not computed ($\ell\notin S$, e.g.\
\texttt{u\_ell\_over\_budget} or \texttt{u\_ell = off}); \texttt{not\_run} when (a)--(c) already passed or
the block is not an $X_1$ block of the unsplit record. A passed (d0) names the certificate
\texttt{"dickson"}; a failed one leaves the verdict to the reasons (a)--(c) produced. \texttt{verify()}
re-derives (d0) from the record alone: the stored $t_p$ over \texttt{operator\_primes}, $n$ and $\ell$.
\end{certificate}
\begin{proof}
By \HA{prop:heckeH0}(1) (no hypothesis on the datum) the $L^aD^b$ span $X_1=M_n^{\Ga}$ and are
simultaneous eigenvectors of every $T_p$, $p\ne\ell$, with eigenvalue $(1+p)p^{a}$, which depends on
$a\bmod(\ell-1)$ only; so $X_1=\bigoplus_{c\in\Lambda(n,\ell)}X_1^{(c)}$ (the class spans), every $T_p$,
$p\ne\ell$, acts on $X_1^{(c)}$ by the scalar $(1+p)p^{c}$, and the $T_S$-primary decomposition of $X_1$ is
the grouping of the $X_1^{(c)}$ by their $S$-eigenvalue systems. A block $V$ whose system $(\lambda_q)$ is
attained by exactly one $c$ is therefore contained in $X_1^{(c)}$, so every $T_p|_V$, $p\ne\ell$, is the
scalar $(1+p)p^{c}$. $\T_\infty$ is generated by the $T_p$ over all primes and $T_\ell=U_\ell$
(\HA{def:operators}; the diamonds $T(p,p)$ act by the scalar $p^{n}$), so $\T_\infty|_V=\Fl[U_\ell|_V]$,
and $U_\ell|_V\in A_S$ because $\ell\in S$. If $q\in S$ is a primitive root with $1+q\ne0$ in $\Fl$,
$(1+q)q^{a}=(1+q)q^{a'}$ forces $a\equiv a'\pmod{\ell-1}$, so the uniqueness hypothesis holds. \emph{No
boundary map, no $r^0$ and no transcription enter.}
\end{proof}

\begin{certificate}[(d1) Boundary --- the $X_4$ blocks]\label{cert:boundary}\PROVED\ \emph{(v1.13.0;
Proposition~F$'$ of the slot-12 finality proposal v2 \S3.4; unconditional since [C] v1.12.0 proves
Proposition~T as \HA{prop:bdrytranscription} and Lemma~B$'$ as \HA{lem:bdrytorus}; owner-approved
2026-09-26)}
Let $\ell\ge5$, $n\ge1$ even, $V=X_{4,\m}$ a certified local $\T_S$-block of $X_4=H^1(\Ga;\Mn)$
\textbf{with $\ell\in S$}, $A_S$ its local algebra, and $R(\ell)$ the least prime $p\ne\ell$ of every class
of $(\Z/\ell)^{\times}$ (key \texttt{boundary\_reps}, in class order $1,\dots,\ell-1$). If
\textup{(H1)} $r^1|_V$ is injective ($\rank(r^1V)=\dim V$) and \textup{(H2)} $t_p=T_p|_V\in A_S$ for every
$p\in R(\ell)$, then $\T_\infty|_V=A_S$.
The record writes \texttt{certificate\_routes.d1} $\in\{$\texttt{passed}, \texttt{not\_injective},
\texttt{failed:<p>}, \texttt{not\_run:budget} (v1.13.1), \texttt{not\_applicable:u\_ell}, \texttt{not\_run}$\}$ ---
\texttt{failed:<p>} the first $p\in R(\ell)$ with $t_p\notin A_S$ --- and \texttt{boundary\_reps} $=R(\ell)$ exactly
when the route reached its representatives (\texttt{passed} or \texttt{failed:<p>}) (the
representatives are \emph{tested}, never generators: \texttt{generators\_certified} stays the Stage-4
$S_\m$). A passed (d1) names the certificate \texttt{"boundary"}. \emph{No shortcut}: all $\ell-1$ classes
are tested, also at $\ell=5$. The driver computes $r^1$ from the boundary pair
(\texttt{SL2ZBoundaryPair}, the restriction alone unless \texttt{with\_cone} builds the cone on it), and
the representatives outside $S$ by the anemic-pool operators, streamed beyond $q\le31$ (from $\ell=11$
on). Its certificate object re-derives (H1) from the stored block basis and $r^1$, and (H2) by the
independent algebra-closure route of \texttt{HeckeAlgebraInvariants}, and the driver refuses a record
whose (d1) does not re-derive.
\emph{The stream's cost cap (v1.13.1; slot-12 FIN audit G1; [X] \S1.3.1 R3).} After (H1) and before any
representative is computed, the route checks
$\sum_{p\in R(\ell),\,p>31}\mathrm{PredictUEllRssMb}(n,p)\le4\cdot B_{u_\ell}$ --- the frozen $u_\ell$ RSS model
(\texttt{u\_ell\_rss/v2}) applied to each streamed $T_p$ ($p+1$ cosets) against four times the \texttt{u\_ell}
budget in force; over the cap it writes \texttt{not\_run:budget} and computes nothing. No settings key: the
cap is fixed and tied to \texttt{u\_ell}, so the route is resumed by re-running with a larger \texttt{u\_ell}
budget. Every census cell ($\ell\le13$, $n\le200$) is within it at the default $2048$\,MB (worst
$(200,13)$: $\approx2.5$\,GB of $8$\,GB); $(10,691)$ ($\approx391$\,GB predicted) and every $\ell=10007$ cell are
not. \emph{The model is kept (v1.15.1-pre.2; owner ruling Q-D, 2026-09-29).} \texttt{u\_ell\_rss/v2-2026-09-17} is the pre-lever fit; since
slot-14 Q8 the $U_\ell$ levers are on by default and the measured RSS is $2.7$--$25\times$ below it, so \texttt{u\_ell\_over\_budget} verdicts are
conservative (slot-14 WP-C10). The owner ruled the frozen model stays: a re-fit would be a value change (verdicts, records, hashes) and must bump the model
version. \emph{What \texttt{verify()} re-derives from the record} (v1.14.0-pre.1; audit G2 repaired --- it
supersedes v1.13.1's ``structure only'', which was a recorded deviation from the slot-12 proposal v2 \S3.4
item~3): (d0) in full (code~31), the structure of (d1) (code~30: the closed list, the group,
\texttt{boundary\_reps} $=R(\ell)$, the cost cap from $(n,\ell,$\texttt{u\_ell.budget\_mb}$)$) \emph{and}, on an
ideal with \texttt{d1 = passed}, (d1) itself (code~32) from the record's own data. There the ideal carries
\texttt{boundary\_data} $=\{$\texttt{rank\_r1}, \texttt{r1}, \texttt{t}$\}$ (\S\ref{sec:schema} item~(j)):
\texttt{r1} $=r^1|_V=r^1V$, an $h^1(\partial)\times D$ matrix over $\Fl$ ($D=\dim V$; rows in the stored basis
of $H^1(\partial)$, columns in the block basis of $V$), \texttt{rank\_r1} its rank, and \texttt{t} the
$\ell-1$ matrices $t_p=T_p|_V$ ($D\times D$), $p\in R(\ell)$ in class order, in the \emph{same} block basis as
the ideal's stored \texttt{t\_p} ($p\in S$). Code~32 refuses: \texttt{boundary\_data} present iff
\texttt{d1 = passed} violated; a shape that is not $h\times D$ / $D\times D$ or keys that are not
\texttt{boundary\_reps}; an entry outside $[0,\ell)$; $\rank(\texttt{r1})\ne$ \texttt{rank\_r1} or
\texttt{rank\_r1} $\ne D$ (H1); a $t_p$ with $p\in R(\ell)\cap S$ differing from the stored \texttt{t\_p}; a
$t_p$ not commuting with the stored generators; and a $t_p$ outside $A_S$ (H2), decided as in the certificate
object: $\dim\Fl[t_q\,(q\in S),\,t_p]=\dim\Fl[t_q\,(q\in S)]$ by the independent algebra-closure route of
\texttt{HeckeAlgebraInvariants} over the record's stored generators. The $t_p$ are \emph{tested}, never
generators; nothing in the record is recomputed from the pipeline. \emph{Size}: only on passed ideals ---
the census's $382$ passed rows have $D\le26$ and $\ell-1\le12$. What the record still does not carry is
$V$ and the ambient $A^{(1)}(T_p)$: that $t_p$ is \emph{the} restriction of $T_p$ to $V$ and $r^1|_V$
\emph{the} restriction of $r^1$ is the in-process certificate's clause (code~4 / code~2 there) and the
driver's refusal, exactly as the stored \texttt{t\_p} of (d0) are taken as the driver's restrictions.
\end{certificate}
\begin{proof}
$V$ is $\T_\infty$-stable (\HA{obl:algebra}, \HA{thm:relations}). The restriction intertwines:
$r^1A(T_g)=A^{\partial}(T_g)r^1$ for every $g$, by \HA{thm:conewelldef}(a) for a Hecke datum of the pair,
which exists by \HA{lem:tauexists}(2) and \HA{lem:etaexists} ($6\in\Fl^{\times}$ as $\ell\ge5$). Let
$p\ne\ell$ be prime and $p'\in R(\ell)$ with $p'\equiv p\pmod\ell$. By \HA{lem:bdrytorus}(2) in degree~$1$
--- which rests on \HA{prop:bdrytranscription}, the Shapiro transcription of the pipeline's
$T^{\partial}_p$ on $H^1(C^{\bullet}_{\partial})$ to the group form \HA{eq:bdrygroupform} --- the boundary
operators satisfy $T^{\partial}_{p}=T^{\partial}_{p'}$. Hence for $v\in V$,
$r^1(T_p-T_{p'})v=(T^{\partial}_p-T^{\partial}_{p'})r^1v=0$, and (H1) gives $T_p|_V=T_{p'}|_V\in A_S$ by (H2).
$T_\ell=U_\ell\in A_S$ as $\ell\in S$, and the diamonds are scalars. So every generator of $\T_\infty$
restricts into $A_S$, i.e.\ $\T_\infty|_V=A_S$. \emph{Injectivity, not Eisenstein-ness, is the
hypothesis}: on the \texttt{mixed} block $(34,7)$ ideal~$2$ ($\dim V=4$, $\rank r^1=2$) every
representative lies in $A_S$ and the boundary law $t_q=t_{q'}$ fails; (d1) refuses it
(\texttt{not\_injective}), which \texttt{HeckeBoundaryFinalityValidation} arm~(X) pins.
\end{proof}

\noindent\emph{Scope and order (v1.13.0).} (d0) and (d1) certify \emph{exactly} the predicate of
(a)--(c) --- Claim~\ref{cl:gen}(b) at $\m$, $\T_\infty|_V=A_S$ --- not a weaker ``the ideal does not
split'', and they run \emph{only} on the blocks where (a)--(c) did not pass. They reach precisely the
range (c) cannot (Claim~\ref{cl:sturmmodl}'s Dickson-type clause, and $X_1$, which
Certificate~\ref{cert:sturm}'s heuristic list does not name): on census editions~1 and~2 every
\texttt{sturm\_heuristic} row is discharged --- editions~1 / 2 hold $270$ / $376$ such rows
($184$ / $184$ at $\ell=5$, of them $36$ on $X_1$; $68$ / $150$ at $7$; $13$ / $37$ at $11$; $5$ / $5$ at
$13$), all certified by the validation's independent computation (slot-12 WP-14, arm~(H)) and by the
driver (S12-FIN, arm~(R) by default, arm~(H) on all $376$). A block of $X_4$ on which $r^1$ is not injective \emph{and} (c) fails would stay non-final. None occurs in census
editions~1--3 (edition~3: $80\,293$ rows, all final; its $3302$ non-injective $H^1$ rows are all off the Dickson
range, $S_V=0$, and final by (c)). None occurs in the mixed-finality scans either: $\ell=5$, $n\le400$; $\ell=7$,
$n\le500$; $\ell\in\{11,13\}$, $n\le340$ on the Dickson cells; $\ell\in\{17,\dots,43\}$ on the Dickson cells with
$n\le200$ (v1.18.0; slot-18 WP-19c's $64$-cell probe). There, no block has $\ker r^1|_V\ne0\ne S_V$. On such a block, slot-13 WP-11's Proposition~M (refereed,
\PROVED\ from \HA{cl:bockstein}, \HA{thm:conewelldef}, \HA{lem:bdrytorus} and Claim~\ref{cl:sturmmodl}'s proof)
would certify $\T_\infty|_V=A_S$. It does so under the hypotheses (H2)--(H5), by gluing the Sturm law modulo the
torsion image $S_V$ to the boundary law modulo $K=\ker r^1|_V$. It is \emph{specified only}: certificate (d2) lands
with the next record event, and only once a row needs it.

Practical order: (a) on every $\m$; (b) on the survivors the socle bound admits; (c) only where
(b) leaves room and only on those $\m$; (d0) / (d1) (v1.13.0) only on the $X_1$ / $X_4$ blocks
(a)--(c) leave non-final. In the regime $n\gg\ell$ the dominant $\m$ will need (c). Record per $\m$
the certified generating set $S_\m$, the certificate type, and $\dim\m/\m^2$. \emph{The answer
to ``which $T_q$ do we need for $(n,\ell)$'' is $\bigcup_\m S_\m$ with its certificates.}

\subsection*{Stage 5 --- module structure at $\m$}
\begin{algorithm}\label{alg:stage5}
With $\T_\m=A_S$ certified, for each $X_i$ and $V=X_{i,\m}$:
\begin{enumerate}
\item Loewy series: $\m^{j+1}V=\sum_a n_a\,\m^jV$ over a basis $\{n_a\}$ of $\m_S$
  (spans and RREF, $O(\dim\m_S\cdot D^3)$); socle series:
  $\soc V=\bigcap_a\ker n_a$ (one stacked kernel), $\soc^{j+1}V$ the preimage of
  $\soc(V/\soc^jV)$. Layer dimensions divided by $f$.
\item Jordan type of each $n_j$ on $V$: the S10 partition machinery, re-used.
\item Structure maps $\theta$ of \ref{D4} for each consecutive pair of layers: the
  $\Fl$-linear map $\m_S/\m_S^2\to\Hom_{\Fl}(\text{top},\text{bottom})$ induced by
  $n\mapsto(\bar b\mapsto n b)$; its $\Flr$-rank is its $\Fl$-rank divided by $f$
  (Claim~\ref{cl:frank}). Active directions and the minimal subset of $S_\m$ spanning them.
\item Canonical submodules \ref{D6}: for $W$ each of $H^1_!$, $\soc(X^+_{4,\m})$ (the
  Eisenstein line, when defined), $X_1$, $r^0X_1$, $L\cdot H^1(M_{n-\ell-1})$, and the kernels
  and images of the maps of \HA{eq:les}, compute
  $\dim(W\cap\m^jV)$ and the image of $W$ in each layer. ($L\cdot$ needs the $(n-\ell-1,\ell)$ run,
  Hecke-equivariant up to the known twist \Mnote{2}.)
\item Cyclic/free \ref{D5}: cyclic iff $\dim V/\m_SV=f$; free of rank one iff moreover
  $\dim V=\dim A_S$. For $X_4$ versus $X_3$, compare with the duality prediction for $\Dn$
  (Prediction~\ref{pr:dual}).
\end{enumerate}
\end{algorithm}

\begin{claim}\label{cl:frank}\PROVED\
If $\phi:U\to W$ is $\Flr$-linear between $\Flr$-spaces, its rank as an $\Fl$-linear map is
$r\cdot\rank_{\Flr}\phi$; all $\Flr$-dimensions of \ref{D2}--\ref{D5} are therefore computable
over $\Fl$ as $\Fl$-dimensions divided by $r=f$, once $V$ is known $\m$-primary with residue
field $\Flr$ (the $\Flr$-structure on the layers is the one through $A_S/\m_S$).
\end{claim}
\begin{proof}
$\im\phi$ is an $\Flr$-subspace; $\dim_{\Fl}=r\dim_{\Flr}$.
\end{proof}

\subsection*{Stage 6 --- the per-$\ell$ atlas}\label{sec:atlas}
Once per $\ell$: run Stages 1--5 for all even $n\le\ell-1$ (where $\Mn$ is irreducible, no Dickson
classes, $d=O(\ell)$), collect the anemic systems with their $a_\ell$ and Galois-orbit structure,
close under $\theta$-twist $a_q\mapsto q^ja_q$, $j\in\Z/(\ell-1)$; this is $\Sigma(\ell)$
(Claim~\ref{cl:sigma}; weight $\ell+1$ must be included: $\De\bmod11$ is a cuspidal level-1
system of weight $12=\ell+1$ that is neither Eisenstein-congruent nor a twist of a lower
weight). Greedily choose $S_{\mathrm{sep}}(\ell)$ (typically 2--4 primes). For
each cuspidal $\bar\rho\in\Sigma(\ell)$, from a weight with multiplicity $\ge2$
($n=w-2+\ell^2-1$, or any weight in the $E_{\ell-1}$-tower) compute $\dim\m/\m^2$ and a generating
pair of anemic operators (Claim~\ref{cl:bk}); store it as the default $S_\m$ (Stage 4 still
certifies). Attach to each $(n,\ell)$ the predictions of \S\ref{sec:expect} before computing;
every record confirms or refutes each, and flags the unpredicted.

\noindent\emph{The atlas against the literature} (v1.15.1-pre.1, a citation and a check, not a claim of this document).
Citro--Ghitza \cite{citro_ghitza} enumerate the level-one eigensystems mod~$p$ one by one over $\overline{\F}_p$: $(p-1)^2/4$
Eisenstein ones (their Cor.~13) and, for every $p<2000$, the number of the others (the table of their appendix). $\Sigma(\ell)$
stores Galois orbits of anemic systems instead, so the comparison is: Eisenstein orbits (all $\Fl$-rational) against
$(\ell-1)^2/4$, and the sum of the residue degrees $r$ of the cuspidal orbits against the table. They agree wherever both exist:
$\Sigma(5),\Sigma(7),\Sigma(11),\Sigma(13)$ hold $4,9,25,36$ Eisenstein systems and $0,0,10,12$ cuspidal ones (slot-6 WP-8; the
table's rows $p=11,13$ read $10,12$), and the stored $\Sigma(691)$ (\texttt{hecke\_atlas/v1}, FNV \texttt{0x73c7210047a81b97})
holds $119\,025=690^2/4$ Eisenstein orbits and $883\,200$ cuspidal orbits whose degrees sum to $6\,762\,000$, the table's row
$p=691$ (read from the store on 2026-09-28). The same row's Eisenstein--cuspidal weights, $12$ and $200$, are the two irregular
classes of $691$.

\subsection*{Stage 7 --- scaling in $n$ for fixed $\ell$}
Run $n$ in residue classes mod $(\ell-1)$; the $E_{\ell-1}$-tower and the $L$-maps connect
$n$, $n+\ell-1$, $n+\ell+1$, so records are comparable $\m$ by $\m$ (regression test; detection of
where an extension is born). Memory: dense $\Fl$ matrices of size $d\times d$, $d\approx n/3$
for all six groups and both signs together ($\dim H^1\approx n/6$); for $n\sim10^4$ that is $\sim10^7$ entries per operator --- but this estimate applies
to level (b) only. Stage 1 works on cochains ($\dim C^1=11(n+1)$, peak RSS $\approx n^2$: $1.4$\,GB
at $n=1000$), so $n\sim10^4$ is out of reach of the present model \emph{on the reference
machine} (inventory \S1, Stage 7). \emph{Reframed (v1.4.0; owner decision OQ12):} $n\le1000$ is
a hardware-bounded limit there, not a design limit --- the target hardware includes HPC nodes
with the memory to push it, the model is kept as the blueprint for higher-rank groups, and a
Manin-sized $H^1$ model (the amalgam presentation of \HA{cl:amalgam}, seeded by
\texttt{BarCocycleH1}) is a downgraded follow-up, not the direction (through v1.3.3 this read
``the practical range today is $n\le1000$, and beyond it a Manin-sized $H^1$ model is the way
forward''). The first extension campaign is bounded at $n\le200$ (slot-1 owner
confirmation~9(b); E6). The
$\tau$-evaluation of Stage 1 dominates (measured exponent $\approx2.05$ in $n$ for $T_2$,
\texttt{CLAUDE.md}); the decomposition layer on legacy \texttt{Matrix} (Berkowitz $O(d^4)$) is fine
to $d\approx200$.

\begin{claim}[The Glover factors of $\Sym^n$, and the convention that fixes them]\label{cl:glover}
\STD{} (1) / \PROVED{} (2), (3), (4) (v1.21.0-pre.5: (3) \STD{} $\to$ \PROVED).
Let $\ell$ be a prime, $G=\GL_2(\Fl)$, $B\subset G$ the upper-triangular Borel and
$\chi_1,\chi_2$ its two diagonal characters. \textbf{The convention} (v1.5.0; slot-4 WP-7
finding~3, and the one thing item~6 of \S\ref{sec:identities} left free): \emph{$\Sym^n$
means the standard symmetric power} $\Mn=\Sym^n(\Fl^{2})$ of \HA{def:rhon} --- the module
on which $\mathrm{diag}(\alpha,\beta)$ has eigenvalues $\alpha^{\,n-t}\beta^{\,t}$,
$0\le t\le n$ --- and \emph{not} its contragredient. This is not a nicety:
$\Sym^n(V)^{\ast}\cong\Sym^n(V)\otimes\det^{-n}$ (\S\ref{sec:symdiv},
\HA{sec:contragredient}), so the two readings give the same $n_i$ and the same congruence
while shifting every individual $j_i$ by $-n$. \textbf{The sentence that fixes it}, in the
form an implementation can quote verbatim: \emph{$\Sym^n$ is the symmetric power of the \textbf{standard}
representation --- the module $\rho_n$ of \HA{def:rhon} acts on, in which
$\mathrm{diag}(a,d)$ has $\mathrm{entry}(t,t)=a^{\,n-t}d^{\,t}$ --- and not of its dual; the
dual would shift every $\det$-twist by $-n$.}

\textbf{And this convention \emph{gates} phase~E7} (v1.6.0; slot-4 WP-7 finding~3, the
consequence it did not draw). E7's whole subject is the contragredient local system $\Dn$ and
$\rho_n^{\vee}$, and the first thing E7 does with the factor multiset is compare the residue
characters of $H^1(\Dn)$ against $\theta$-twisted systems --- which is item~6 of
\S\ref{sec:identities}, and which E6's atlas will evaluate. Under the two readings that
comparison differs by $\theta^{-n}$ on \emph{every} factor. Both readings are internally
consistent, so \textbf{a green validation suite cannot detect the discrepancy}: every
structural invariant an implementation can check --- the $n_i$, the dimension count
$\sum_i(n_i+1)=n+1$, the central-character congruence, even the Brauer character against a
matrix $\rho_n$, which is computed in whichever convention the matrices are in --- is
convention-independent or self-consistent. A wrong reading would therefore surface as an E7
\emph{finding}, indistinguishable from a mathematical result. E7 may not start against an
unfixed convention: \textbf{[E] v1.6.0 is a precondition of E7's start}, and this paragraph
is the thing it is a precondition for.
\begin{enumerate}
\item The simple $\Fl[G]$-modules are $L(r,j)=\Sym^{r}\otimes\det^{j}$ with
  $0\le r\le\ell-1$ and $j\in\Z/(\ell-1)$ (for $\ell=2$ the twist group is trivial), and
  the central character forces $r+2j\equiv n\pmod{\ell-1}$ on every Jordan--H\"older factor
  of $\Mn$~\cite{glover} ((3.1), p.~433; (6.2), p.~459; in Glover's notation $V_m\otimes D^{n}$ is $\Sym^{m-1}\otimes\det^{n}$).
\item With $\theta:=x^{\ell}y-xy^{\ell}=\prod_{[a:b]\in\mathbb{P}^{1}(\Fl)}(bx-ay)$, which
  satisfies $g\cdot\theta=\det(g)\,\theta$, multiplication by $\theta$ is an injection of
  $G$-modules $\Sym^{\,n-\ell-1}\otimes\det\hookrightarrow\Sym^{n}$, and its cokernel has
  dimension $(n+1)-(n-\ell)=\ell+1$ and is the principal series
  $\mathrm{Ind}_B^G(\chi_1^{\,n}\chi_2^{\,0})$ --- restriction of a degree-$n$ form to each
  of the $\ell+1$ lines of $\mathbb{P}^{1}(\Fl)$.
\item $\mathrm{Ind}_B^G(\chi_1^{\,s}\chi_2^{\,t}):=\Fl[G]\otimes_{\Fl[B]}\chi_1^{\,s}\chi_2^{\,t}$, where
  $\chi_1^{\,s}\chi_2^{\,t}\bigl(\begin{smallmatrix}a&b\\0&d\end{smallmatrix}\bigr)=a^{s}d^{t}$, has the two Jordan--H\"older factors
  $\Sym^{r}\otimes\det^{t}$ and $\Sym^{\,\ell-1-r}\otimes\det^{s}$ with
  $r=(s-t)\bmod(\ell-1)$, the trivial/Steinberg pair when $r=0$ (v1.21.0-pre.5: \PROVED{} below; \cite{glover} states no
  induced-module result --- its (4.1), p.~439, (4.2), p.~440, (5.7), p.~454, (6.4), p.~461, describe $\Sym^{n}$ by explicit bases ---
  so the earlier citation for this item is withdrawn).
\item Hence the multiset of Jordan--H\"older factors of $\Mn$, written as pairs
  $(n_i,j_i)$, is computed by the terminating recursion
  \[
    \mathrm{JH}(n)=\begin{cases}
      \emptyset & n<0,\\[2pt]
      \{(n,0)\} & 0\le n\le\ell-1,\\[2pt]
      \{(r,0),\ (\ell-1-r,\,n)\}\ \cup\ \bigl(\mathrm{JH}(n-\ell-1)\otimes\det\bigr),
        & n\ge\ell,\ \ r:=n\bmod(\ell-1),
    \end{cases}
  \]
  in $\lfloor (n+1)/(\ell+1)\rfloor$ steps, with at most
  $2\lfloor (n+1)/(\ell+1)\rfloor+1$ raw
  factors before equal pairs are merged, and $\sum_i(n_i+1)=n+1$ counted with multiplicity.
  (v1.5.1, corrected from $\lfloor n/(\ell+1)\rfloor+1$ steps and
  $2\lfloor n/(\ell+1)\rfloor+1$ factors, which undercounts: at $n=\ell$ the recursion emits
  two raw factors against a bound of $1$, and at $n=2\ell+1$ four against $3$.)
\end{enumerate}
\end{claim}

\begin{proof}
(2) The product formula for $\theta$ and $g\cdot\theta=\det(g)\theta$ are classical; $\theta$
is a non-zero divisor in $\Fl[x,y]$, so multiplication by it is injective, and it raises
degree by $\ell+1$ and twists the $G$-action by $\det$, which is the stated map. The
cokernel is $\Fl[x,y]_n/\theta\,\Fl[x,y]_{n-\ell-1}$; since $\theta$ is the product of the
$\ell+1$ distinct linear forms cutting out the lines of $\mathbb{P}^{1}(\Fl)$, a form of
degree $n$ lies in $(\theta)$ iff it vanishes on all of them, so the quotient is the image
of the restriction map into the $\ell+1$ one-dimensional spaces, i.e.\ the induced module
with the stated character. (3) Write $\tilde\chi:=\chi_1^{\,s}\chi_2^{\,t}$, $q:=\ell-1$, $V:=\Sym^{r}\otimes\det^{t}\oplus\Sym^{\,q-r}\otimes\det^{s}$;
both summands are simple by~(1) ($0\le r\le q$). Two finite-dimensional $\Fl[G]$-modules with equal Brauer characters
have the same Jordan--H\"older factors (the Brauer characters of the simple modules are linearly independent; $\Fl$ is a
splitting field by~(1)), so it suffices to compare Brauer characters on the $\ell$-regular, i.e.\ semisimple, elements $g$,
with Teichm\"uller lifts $\alpha,\beta\in\mu_{\ell^2-1}$ of the eigenvalues of $g$. Frobenius' formula
$\chi_{\mathrm{Ind}}(g)=|B|^{-1}\sum_{x\in G,\ x^{-1}gx\in B}\tilde\chi(x^{-1}gx)$, valid for Brauer characters, gives:
$(\ell+1)\alpha^{s+t}$ at a central $g$ ($\alpha=\beta$); $\alpha^{s}\beta^{t}+\alpha^{t}\beta^{s}$ at a split $g\sim\mathrm{diag}(\alpha,\beta)$,
$\alpha\ne\beta$ (the conjugates of $g$ in $B$ have diagonal $(\alpha,\beta)$ or $(\beta,\alpha)$, each for $|B|$ values of $x$); and $0$ at a
non-split $g$, which has no conjugate in $B$. On $V$, with $\chi_{\Sym^{m}}(g)=\sum_{i=0}^{m}\alpha^{m-i}\beta^{i}$:
at a central $g$, $(r+1)\alpha^{r+2t}+(q-r+1)\alpha^{q-r+2s}=(\ell+1)\alpha^{s+t}$, since $\alpha^{q}=1$ and $r\equiv s-t$;
at a split $g$, $\alpha,\beta\in\mu_q$ and $\alpha^{\ell-r}=\alpha^{1-r}$, so
\[
(\alpha-\beta)\chi_V(g)=(\alpha\beta)^{t}\bigl(\alpha^{r+1}-\beta^{r+1}\bigr)+(\alpha\beta)^{s}\bigl(\alpha^{1-r}-\beta^{1-r}\bigr)
=(\alpha\beta)^{t}\bigl(\alpha^{r+1}-\beta^{r+1}+\alpha\beta^{r}-\alpha^{r}\beta\bigr)=(\alpha-\beta)(\alpha\beta)^{t}(\alpha^{r}+\beta^{r}),
\]
using $(\alpha\beta)^{s}=(\alpha\beta)^{t+r}$; hence $\chi_V(g)=\alpha^{t+r}\beta^{t}+\alpha^{t}\beta^{t+r}=\alpha^{s}\beta^{t}+\alpha^{t}\beta^{s}$;
at a non-split $g$, $\beta=\alpha^{\ell}$, $\alpha=\beta^{\ell}$ and $N:=\alpha\beta\in\mu_q$, so $\alpha^{\ell-r}-\beta^{\ell-r}=\beta\alpha^{-r}-\alpha\beta^{-r}=-N^{-r}(\alpha^{r+1}-\beta^{r+1})$
and $(\alpha-\beta)\chi_V(g)=(N^{t}-N^{s-r})(\alpha^{r+1}-\beta^{r+1})=0$. The two Brauer characters agree on every semisimple class,
which proves~(3); for $r=0$ the factors are $\det^{t}$ and $\Sym^{q}\otimes\det^{t}$ (as $s\equiv t$), the twisted trivial/Steinberg pair.
Item~(3) is used below only through these factors, so the realisation of the cokernel in~(2) enters (4) only through its Brauer
character, $\alpha^{n}+\beta^{n}$ at a split $g$ ($=(\alpha^{n+1}-\beta^{n+1}-\alpha\beta(\alpha^{n-\ell}-\beta^{n-\ell}))/(\alpha-\beta)$ with $\alpha^{q}=\beta^{q}=1$), which is
$\chi_{\mathrm{Ind}}$ at $(s,t)=(n,0)$, and $0$ at a non-split $g$ likewise. \emph{Checked} (v1.21.0-pre.5; own Python, exact in $\Z[x]/\Phi_{\ell^2-1}$, the
Frobenius sum taken over the explicit group): every $(s,t)\in(\Z/q)^2$ on every semisimple class at $\ell\in\{2,3,5,7,11,13\}$, and the multiplicities of
$\Mn$, solved over all $\ell(\ell-1)$ simples and verified exactly, equal to the recursion of~(4) for $n\le3\ell(\ell+1)$.
(4) Induction on $n$, the base case being $n\le\ell-1$, where
$\Mn$ is irreducible (item~(1) with $r=n$): for $n\ge\ell$ the exact sequence of (2)
together with (3) at $(s,t)=(n,0)$ gives the two displayed factors and the twisted factors
of $\Mn[-\ell-1]$; each step lowers $n$ by $\ell+1$, so the chain
$n,\,n-(\ell+1),\dots$ contains exactly $\lfloor (n+1)/(\ell+1)\rfloor$ values $\ge\ell$,
each emitting two raw factors, and terminates in a value of $[-1,\ell-1]$ emitting one or
none --- whence the two counts of (4) --- and the dimension count is the alternating sum of the
same exact sequence. The central character of every emitted pair satisfies
$r+2j\equiv n\pmod{\ell-1}$: for $(r,0)$ because $r\equiv n$, and for $(\ell-1-r,n)$
because $(\ell-1-r)+2n\equiv n-r+n-n\equiv0$ is $n\equiv r$ again.
\end{proof}

\noindent\emph{Implemented} (\texttt{HeckeExt/GloverFactors}, slot-4 WP-7): the recursion
of~(4), with the convention of the preamble pinned to the pipeline's own $\rho_n$
(\texttt{ComputeSymmetricPowerMatrixFM} branch \texttt{[C1]},
$\mathrm{entry}(t,t)=a^{\,n-t}d^{\,t}$), checked against the Brauer character of $\rho_n$ on
the $\ell$-regular classes of $\SL_2(\Fl)$, computed independently from those matrices, for
every $n\le3\ell$ at $\ell\in\{5,7\}$ (\texttt{HeckePredictionsValidation}). Identity~6 of
\S\ref{sec:identities} is \emph{not} thereby asserted --- that needs the atlas (E6) --- but
the factor list it quantifies over is computed, and its $\det$-twists have one meaning.

\subsection{Identities the computation must satisfy}\label{sec:identities}
All \STD\ unless marked; each becomes an assertion in the validation that will gate this
feature.
\begin{enumerate}
\item \HA{cl:dimH1}; $\dim X_2=\dim X_5=\lfloor n/\ell\rfloor+1$; exactness of
  \HA{eq:les} with the dimension bookkeeping of \HA{cl:les};
  $\dim X_1=\#\{a(\ell+1)+b(\ell^2-\ell)=n\}$.
\item $[T_q,T_{q'}]=0$, $[T_q,\iota]=0$, equivariance of every map in \HA{eq:les};
  the recursion of \HA{thm:relations} at $q^2$.
\item Boundary eigenvalues $\chi_a$ (\HA{cl:chi}, \VERIFY{} --- the $H^0$ characters and the
  $n<\ell$ Eisenstein eigenvalue are asserted, and \HA{cl:chi}(2)'s general-$n$ $H^1$ multiset
  recorded, $330/330$ cells, by \texttt{HeckeBoundaryOperatorValidation} since slot 1; the
  tag stays until the claim is proved); $1+q^{n+1}$ on the reduced
  Harder class; $\tau(q)\bmod\ell$ at $n=10$; the known cusp forms of \HA{sec:example}.
\item $r^1=0$ on the Eisenstein class iff $\ell\mid\num(B_k/k)$ ($n<\ell-1$;
  Claim~\ref{cl:harder}) --- asserted in the form ``$r^1$ vanishes on the joint eigenspace
  $E$'' on the Kummer list and ``$\rank(r^1E)=1$'' off it; verified $37/37$
  (\texttt{HeckeBoundaryParabolicValidation}, 2026-09-16).
\item Poincar\'e duality (\HA{cl:pd}) once $\Dn$ runs (E7, scheduled --- v1.8.2: \emph{accepted}
  at the slot-7 close, \HA{cl:pd}(1)--(2) tested on all $54$ cells of the E2 case list at
  $\ell\ge5$, failing on none; the adjointness of
  \HA{cl:pd}(2) is the acceptance gate for $X_3$ at $\ell\ge5$,
  \texttt{HeckeDualityValidation}): equal characteristic polynomials
  between $X_3(\Mn)$ and $H^1(\Dn)$, and the Loewy layers of each equal to the socle layers of
  the other in reverse order (Prediction~\ref{pr:dual}; \emph{not} reversed Loewy series).
\item Compatibility along $L\cdot$ and along the Ash--Stevens filtration \Mnote{5}: the
  semisimplification of $H^1(\Ga;\Mn)$ is the sum over the Glover factors
  $\Sym^{n_i}\otimes\det^{j_i}$ ($0\le n_i\le\ell-1$, $n_i+2j_i\equiv n\bmod(\ell-1)$) of
  $\theta^{j_i}$-twists of weight-$(n_i+2)$ systems, plus the $H^0$-corrections.
  Claim~\ref{cl:glover} supplies the factor multiset itself \emph{and} the convention that
  fixes the individual $j_i$; without that convention this item does not have a
  determinate meaning (v1.5.0). \textbf{And this is the item on which the convention gates
  E7} (v1.6.0): it is here that the $\det$-twists meet $\Dn$, so the two readings of
  $\Sym^n$ differ by $\theta^{-n}$ in exactly this comparison and nowhere a validation can
  see them (Claim~\ref{cl:glover}'s second preamble paragraph). E7 starts against the
  standard-module reading or not at all.

  \emph{Implemented} (\texttt{HeckeExt/GloverFactors} for the factor multiset, slot-4 WP-7;
  bound in the console as \texttt{glover\_factors(n, ell)}, slot-5 WP-10, which reports the
  same convention). The \emph{identity} is not thereby asserted --- that needs the atlas.
  \emph{Status} (v1.8.0; the slot-6 WP-9 result, recorded in v1.7.2's changelog and now stated
  here). The identity \textbf{is asserted}, as item~6 of \texttt{HeckeIdentityCertificate}, on
  $108$ cells, $\ell\in\{5,7,11\}$ times $n\in\{0,2,\dots,60\}\cup\{80,100,120,160,200\}$,
  against the in-process $\Sigma(\ell)$ (\texttt{HeckeExtEndToEndValidation}, arms~[B]
  and~[C]). It holds under the standard-module convention, and the
  contragredient reading fails on all $63$ cells where the two differ. It is \textbf{not
  evaluated} at $\ell=691$. $\Sigma(691)$ is slot-7 WP-5's, and it is not built at this revision
  (v1.8.2: built by slot-7 WP-5 and published by hash and build command only; the identity is
  still not evaluated at $\ell=691$).
  Even once built, it is a stored atlas, which neither the record nor golden reads (slot-7 owner
  ruling~3b~(B); \S\ref{sec:schema}, ``The stage-2 keys''~(b)). At $\ell=691$ the record carries
  \texttt{atlas\_unavailable}, and the atlas is joined to the records only outside them, by the
  campaign's analysis tool.

  \emph{v1.12.0-pre.1: the $L\cdot$ half extended by its long exact sequence} (E1 identity
  class~5 extended; slot-11 WP-5 and WP-6; the objects, their names and indexing are
  \S\ref{sec:thetaobj}, the mathematics is [C]'s section on the Dickson invariant and $\theta$).
  Through v1.11.1 the $L\cdot$ half was the one identity
  $A^{(1)}_n(T_p)\,L_\ast=p\,L_\ast\,A^{(1)}_{n-\ell-1}(T_p)$ (\texttt{HeckeIdentityCertificate}
  item~5, slot-5 WP-8). From v1.12.0-pre.1 it is the following list, written with
  $\theta=\theta_n:H^q(\Ga;\Mn\otimes\det)\to H^q(\Ga;M_{n+\ell+1})$ (source-indexed,
  \S\ref{sec:thetaobj}~(a)), $Q_n$ the quotient local system, $\ell\ge5$, $p$ a prime $\ne\ell$:
  \begin{enumerate}[label=(6\alph*)]
  \item \emph{fibre}: $\rho_{n+\ell+1}(A)\,L=\det(A)\,L\,\rho_n(A)$ for every $A\in M_2(\Fl)$,
    singular included, and hence $\rho_{n+\ell+1}(A)(L\cdot\Sym^n)\subseteq L\cdot\Sym^n$, so
    $\rho_{Q_n}(A)$ is defined for every $A\in M_2(\Fl)$ (the well-definedness certificate of
    \texttt{FibreRepresentation::QuotientL});
  \item \emph{exactness}: the six-term sequence
    $0\to H^0(\Mn)\to H^0(M_{n+\ell+1})\to H^0(Q_n)\xrightarrow{\delta}H^1(\Mn)\xrightarrow{\theta}
    H^1(M_{n+\ell+1})\to H^1(Q_n)\to0$ is exact at every spot, as \emph{subspaces} (a basis of the
    image and of the kernel at each spot, the inclusion both ways), with $\delta$ computed by the
    snake construction on cochains and never inferred from ranks; the three $H^2$ terms are $0$
    (\HA{lem:h2torsion});
  \item \emph{dimension}: the alternating sum of the six dimensions is $0$, and $\rank\delta$
    equals the closed form in $(n,\ell)$ of \HA{thm:thetales};
  \item \emph{Dickson ring}: $H^0(\Ga;\Mn)=\Fl[L,D]_n$, the span of the explicit invariants
    $L^aD^b$, $a(\ell+1)+b(\ell^2-\ell)=n$, equal to the computed $H^0$ in dimension and span
    (\HA{prop:dicksonH0}, with $D$ as \HA{def:dicksoninv} defines it; consistent with
    \HA{rmk:dickson}'s semigroup and with the dimension formula of item~1);
  \item \emph{Hecke}: for $p\ne\ell$, $A(T_p)_{n+\ell+1}\,\theta=p\,\theta\,A(T_p)_n$ in degrees
    $0$ and~$1$, entrywise, with the untwisted identity asserted to \emph{fail} wherever
    $\theta\,A(T_p)_n\ne0$ and $p\not\equiv1\pmod\ell$ (the exponent pinned, as in item~5); $\pi$ intertwines $A(T_p)_{n+\ell+1}$ with $A(T_p)_{Q_n}$
    untwisted, and $\delta$ intertwines $A(T_p)_{Q_n}$ with $p\,A(T_p)_n$; and over $\Fl$
    $A(U_\ell)_{n+\ell+1}\,\theta=0$, already on cochains (Claim~\ref{cl:Tlnil}(1),
    \HA{prop:thetatwist}), and $\delta\,A(U_\ell)_{Q_n}=0$;
  \item \emph{$H^0$ eigenvalues} (\HA{prop:heckeH0}): $A(T_p)$ acts on $L^aD^b$ by
    $(1+p)\,p^a$ for $p\ne\ell$; $A(U_\ell)$ on $H^0(\Ga;\Mn)$, $n\ge1$, has rank $\le1$ and
    $A(U_\ell)^2=0$, kills every $L^aD^b$ with $a\ge1$, and is \emph{zero for $n<\ell^3-\ell$} but
    \emph{not} in general: at $n=\ell^3-\ell$ it is non-zero, the first instance being $(120,5)$,
    $U_5D^6=-L^{20}$ (the only one with $\ell\in\{5,7,11,13\}$, $n\le200$); at $n=0$,
    $U_\ell=1$. What ``every Dickson class is non-ordinary'' needs is $U_\ell^2=0$
    (\HA{cor:thetanonordinary}), which holds at every $n\ge1$. \emph{v1.12.0-pre.1 correction}:
    the slot-11 plan (and this item's first draft) said ``$U_\ell=0$ on all of $H^0$''; [C]
    refutes it (\HA{rmk:heckeH0plan}), and the assertion is the one above, with $(120,5)$ an arm.
    On $H^0(Q_n)$ (dimension $[(\ell-1)\mid n+2]$, spanned by $\mathbf 1_n$):
    $T_p\mathbf 1_n=(p+1)\mathbf 1_n$ and $U_\ell\mathbf 1_n=0$ (\HA{lem:Qnfunctions});
  \item \emph{kernel}: $\ker(\theta|_{H^1})=\im\delta$ as subspaces, with its decomposition by
    systems predicted in writing before the run (slot-11 rule~20);
  \item \emph{negative seams}: replacing $L$ by $x^\ell y$ (not $\det$-equivariant) fails (6a)'s
    certificate, and a perturbed $\delta$ fails (6b).
  \end{enumerate}
  Every item is a statement of [C] or follows from one by the twist argument of
  Claim~\ref{cl:Tlnil}(1) (the Hecke operators of $\Mn\otimes\det$ are $p\,A(T_p)_n$ and
  $\ell\,A(U_\ell)_n=0$, and every map of the sequence commutes with the operators of its own
  terms); its STATUS is [C]'s. This list names only what the code asserts; where the two differ,
  [C] wins and the package reports, and a clause [C] does not state is run as a recorded
  observation, not as a claim (slot-11 rule~21). (6a)--(6d) and~(6h) are
  slot-11 WP-5's (\texttt{HeckeDicksonThetaValidation}, \texttt{standard}, and its \texttt{slow}
  coverage companion), (6e)--(6g) WP-6's (arms of the same name or a second \texttt{standard}
  name). \textbf{They are validation arms, not certificate items}: the item codes $1$--$7$ of
  \texttt{HeckeIdentityCertificate} are its tamper vocabulary and stay as they are (item~5 keeps
  its one identity), no record carries any of (6a)--(6h), and no golden line is added (slot-11
  decision~6~(A)). At $\ell\in\{2,3\}$ none of them is asserted (\HA{lem:h2torsion}'s hypothesis
  $6\in R^\times$ fails, and the sequence need not stop at $H^1(Q_n)$).
  \emph{v1.12.0, as built}: (6a)--(6d) and~(6h) are \texttt{HeckeDicksonThetaValidation}
  (\texttt{standard}, $108$ cells, $979$ checks) and \texttt{HeckeDicksonThetaCoverageSlowValidation}
  (\texttt{slow}, $404$ cells, $3590$ checks); (6e)--(6g) are \texttt{HeckeDicksonThetaHeckeValidation}
  (\texttt{standard}, $5482$ checks) and \texttt{HeckeDicksonThetaHeckeCoverageSlowValidation}
  (\texttt{slow}, $3836$ checks), with (6e) also at cochain level ($n\le20$, $U_\ell$ included). Every item
  held on every cell. In (6e) the ``untwisted identity fails'' check adds nothing to the twisted one: for
  $p\not\equiv1\pmod\ell$, $p\,X=X$ iff $X=0$, so the twisted identity alone pins the exponent (slot-11
  WP-14, m-5).
\item Sum over $\m$ of the Loewy layer dimensions equals $\dim X_i$; the Jordan partition of
  $n_j$ on $V$ refines the Loewy length ($\nu\ge$ largest part).
\end{enumerate}

\subsection{The Dickson--\texorpdfstring{$\theta$}{theta} objects}\label{sec:thetaobj}
This subsection is the specification slot-11 WP-5 and WP-6 build against, written before either
commits a claim (slot-11 plan \S3, rule~21). \emph{Division of labour}: [C]'s section on the
Dickson invariant and $\theta$ (\HA{sec:dicksontheta}; slot-11 WP-2, [C] v1.11.0-pre.1) owns
the mathematics --- every statement, its hypotheses, its proof or
its reference, the complement basis of $Q_n$ and the closed form of the dimensions; this subsection
owns the names, the indexing, what the code computes, where each statement is asserted, the
hand-off to the census (\S\ref{sec:census}) and the refusals. Where the two differ, [C] wins and
the package that finds the difference reports it (slot-11 rule~21; WP-2's author repairs).
\begin{enumerate}[label=(\alph*)]
\item \textbf{Indexing and the map.} Throughout, $\ell\ge5$, $\Mn=\Sym^n(\Fl^2)$ in the
  standard-module convention of \HA{def:rhon}, and $L=x^\ell y-xy^\ell$. The map is
  \emph{source-indexed}, as in [C] and the slot-11 plan:
  \[
    \theta_n:=L_\ast:\ H^q(\Ga;\Mn\otimes\det)\longrightarrow H^q(\Ga;M_{n+\ell+1}),\qquad q\in\{0,1\},
  \]
  induced by $f\mapsto Lf$, which is $\De$-equivariant by the fibre identity
  $\rho_{n+\ell+1}(A)\,L=\det(A)\,L\,\rho_n(A)$ for every $A\in M_2(\Fl)$ (the identity~$(\ast)$
  of \texttt{HeckeExt/HeckeDicksonMultiplication.h}). On $\Ga$ the character $\det$ is trivial, so
  $H^q(\Ga;\Mn\otimes\det)=H^q(\Ga;\Mn)$ as groups and the pipeline's own objects at $n$ serve;
  the twist is visible only to the Hecke operators, which on $\Mn\otimes\det$ are $p\,A(T_p)_n$ and
  $\ell\,A(U_\ell)_n=0$ over $\Fl$. \texttt{HeckeDicksonMultiplication} and the console verb
  \texttt{L\_map(H\_lower, H)} are \emph{target-indexed} with $m=n-\ell-1$; they compute the same
  map, $\theta_m$. The census (\S\ref{sec:census}) is target-indexed too: at a cell $(n,\ell)$ the
  $\theta$-image is $\im\theta_{n-\ell-1}\subseteq H^1(\Ga;\Mn)$, which is $0$ for $n<\ell+1$ (in
  particular at every published $\ell=691$ cell, $n\le200$). \emph{Not to be confused} with the
  atlas label's \texttt{theta\_twist} (\S\ref{sec:schema}), which records a system's twist relative
  to its atlas base system, a labelling convention; the census reports both and asserts no
  relation between them.
\item \textbf{The quotient local system.} $Q_n:=M_{n+\ell+1}/L\cdot\Mn$, of dimension $\ell+1$ for
  every $n\ge0$. Its fibre action $\rho_{Q_n}(A)$ is defined for \emph{every}
  $A\in M_2(\Fl)$, singular included (by $(\ast)$, $L\cdot\Mn$ is stable under every $A$), written
  on the monomial complement basis $\mathcal B_n=(\bar u_0,\dots,\bar u_{\ell-1},\bar u_{n+\ell+1})$
  that \HA{def:Qn} fixes, with its class map --- the code implements that basis and never chooses
  one (the normal form $(\bar u_0,\bar u_{n+2},\dots,\bar u_{n+\ell+1})$ is \emph{not} it). It enters the pipeline as a third \texttt{FibreRepresentation} enumerator,
  \texttt{QuotientL} (\texttt{LocalSystemUtils}), appended, with its relation list for
  \texttt{LocalSystemSL2Z} and its arm in \texttt{HeckeOperatorOnK} (slot-11 WP-5 only,
  append-only): \texttt{SymN} and \texttt{ContragredientN} stay byte-identical
  (\texttt{bench\_identity.sh} $45/0$, the golden gate). Well-definedness is certified entrywise
  ((6a) of \S\ref{sec:identities}) on seeded invertible \emph{and singular} $A$ and on every
  relation matrix, because WP-6 applies $U_\ell$, whose representatives are singular mod~$\ell$, on
  all six terms.
\item \textbf{The six terms and five maps; code names.} For $\ell\ge5$ (\HA{lem:h2torsion}: the
  $H^2$ terms vanish at $6\in R^\times$) the long exact sequence of
  $0\to\Mn\otimes\det\xrightarrow{L\cdot}M_{n+\ell+1}\to Q_n\to0$ is
  \[
    0\to H^0_n\xrightarrow{\ \theta^0\ }H^0_{\uparrow}\xrightarrow{\ \pi^0\ }H^0_{Q}
    \xrightarrow{\ \delta\ }H^1_n\xrightarrow{\ \theta^1\ }H^1_{\uparrow}
    \xrightarrow{\ \pi^1\ }H^1_{Q}\to0,
  \]
  $H^q_n=H^q(\Ga;\Mn)$, $H^q_\uparrow=H^q(\Ga;M_{n+\ell+1})$, $H^q_Q=H^q(\Ga;Q_n)$. The code names
  are fixed here: terms \texttt{H0\_n}, \texttt{H0\_up}, \texttt{H0\_Q}, \texttt{H1\_n},
  \texttt{H1\_up}, \texttt{H1\_Q}; maps \texttt{theta0}, \texttt{pi0}, \texttt{delta},
  \texttt{theta1}, \texttt{pi1}. Each map is a matrix over $\Fl$ \emph{in the stored bases} of the
  pipeline's \texttt{EquivariantCohomologyOnK} objects of its two terms (column $j$: the image of
  the source's $j$-th stored basis vector, in the target's stored coordinates), the convention of
  \texttt{DicksonInducedMaps}. $\delta$ is computed by the snake construction on cochains (lift
  to $M_{n+\ell+1}$-cochains, apply the coboundary, divide by $L$), never inferred from ranks.
  Home: \texttt{HeckeExt/HeckeDicksonTheta.\{h,cpp\}} (new; WP-5, then WP-6), \texttt{namespace
  hecke\_ext}, under the layer's one-way dependency.
\item \textbf{The Dickson ring.} $D:=\sum_{i=0}^{\ell}x^{i(\ell-1)}y^{(\ell-i)(\ell-1)}$ as
  \HA{def:dicksoninv} defines it ($=(x^{\ell^2}y-xy^{\ell^2})/L$, \HA{lem:dicksontwist}(2)), of
  degree $\ell^2-\ell$, fixed by every $A$ invertible mod~$\ell$ but \emph{not} by singular ones
  (\HA{lem:dicksontwist}(3)) --- which is why $U_\ell$ need not vanish on $H^0$ ((6f));
  $H^0(\Ga;\Mn)=\Fl[L,D]_n$ with basis $\{L^aD^b:a(\ell+1)+b(\ell^2-\ell)=n\}$
  (\HA{prop:dicksonH0}), built as explicit invariants and compared with the computed $H^0$ in
  dimension and span ((6d)).
\item \textbf{What is certified, what is asserted.} (6a)--(6h) of \S\ref{sec:identities}, with the
  packages and names stated there. The $\theta$-cycle, Ash--Stevens and companion-form instances
  (slot-11 WP-6 scope~3) are \emph{the list \HA{rmk:thetacheckable} states as checkable on the
  model at level~1 and $n\le200$} (\HA{sec:thetaliterature}); each is predicted in writing, with its source, in the package's report skeleton and
  committed before it is run (slot-11 rule~20), and each is compared against an independent object
  --- the label of $\Sigma(\ell)$ (\S\ref{sec:atlas}) and the $T_p$-eigenvalues for $p\le13$.
  Agreement and disagreement are both results; a disagreement is a finding, reported, never
  re-fitted.
\item \textbf{The hand-off to the census.} For each published cell $(n,\ell)$ of
  \texttt{ext\_first} with $\ell\le13$ and $n\ge\ell+1$ and each maximal ideal $\m$ of its record,
  slot-11 WP-6 computes $d_\theta(\m):=\dim_{\Fl}\big(E_\m\cap\im\theta_{n-\ell-1}\big)$, $E_\m\subseteq
  H^1(\Ga;\Mn)$ the generalised eigenspace of~$\m$. Since $\im\theta$ is stable under every $T_p$
  and killed by $U_\ell$ ((6e)), it is a Hecke submodule, so $\im\theta=\bigoplus_\m(\im\theta)_\m$ and
  $\sum_\m d_\theta(\m)=\rank\theta^1$ --- asserted per cell. The run identifies $E_\m$ with the
  record's ideal by the matching rule of \S\ref{sec:census}~(e) (one-to-one; a mismatch is a
  finding, never a guess). Output: one line per ideal,
  \begin{center}\footnotesize
  \texttt{THETA\_IMAGE n=<n> ell=<ell> ideal=<id> witness=<c0,c1,...> dim\_E=<d> theta\_image\_dim=<t> rank\_theta\_cell=<r>}
  \end{center}
  (\texttt{witness} the record's \texttt{witness\_poly}, low$\to$high), printed by WP-6's
  validation only under the read-only interlock \texttt{HECKE\_THETA\_IMAGE\_ALL=1} (not an
  emitter; the \texttt{HECKE\_INTEGRAL\_CAMPAIGN} pattern), and collected by the census tool into
  \texttt{theta\_image.tsv} with the format tag \texttt{hecke\_census\_theta/v1} (\S\ref{sec:census}~(h)).
  Cells with $n<\ell+1$, and every $\ell=691$ cell, are filled by the tool with $d_\theta=0$ and the
  source \texttt{n\_lt\_ell\_plus\_1}; they are not run.
\item \textbf{What these objects are not.} No member of \texttt{hecke\_ext/v3},
  \texttt{hecke\_report/v2} or \texttt{hecke\_ring\_report/v1} carries them: the record's per-ideal
  \texttt{canonical} slot, which \ref{D6} reserves for $L\cdot H^1(M_{n-\ell-1})$ among others,
  keeps emitting \texttt{\{"status": "unavailable"\}}. No settings key is added (eleven stay
  eleven): $\theta$ is a derived object with a validation, not a setting. No golden line is added
  (slot-11 decision~6~(A)); no \texttt{HeckeIdentityCertificate} item code is added.
\item \textbf{Refusals} (\S\ref{sec:ring}~(b), the v1.12.0-pre.1 paragraph). Every entry point of
  \texttt{HeckeDicksonTheta}, the \texttt{QuotientL} relation list and \texttt{HeckeOperatorOnK}'s
  \texttt{QuotientL} arm are $\Fl$-only in slot~11: under a non-default
  \texttt{coefficient\_ring}, or handed a composite object, each raises
  \texttt{hecke::NotImplementedInThisBuild} naming itself (console exit~$5$). $\ell<5$ is refused as
  a mathematical precondition, by \texttt{ENGINE\_CHECK\_MSG} naming \HA{lem:h2torsion} (a
  \texttt{std::logic\_error}, console exit~$3$), not as a build capability. Odd $n$ is not
  refused: every term vanishes on this model ($-I$), and the sequence is the zero sequence.
  $\theta$ over $\Z/\ell^m$ is slot~12's, after [C]'s Hida section (slot-11 plan \S3.15).
\end{enumerate}
\emph{v1.12.0, as built} (slot-11 WP-5, WP-6). (a)--(e) and (g) as specified; the code names of~(c)
are \texttt{hecke\_ext::DicksonThetaCell}'s. (f): the \texttt{THETA\_IMAGE} lines cover $384$ cells and
$1934$ ideals, \emph{every} block matched its record ideal one-to-one by \texttt{witness\_poly} and exact
restricted \texttt{t\_p}, and $\sum_\m d_\theta(\m)=\rank\theta^1$ held on every cell; the table is
committed as \texttt{docs/implementation/slot11/WP-6\_theta\_image.tsv} (\texttt{hecke\_census\_theta/v1})
and copied into the census edition as \texttt{theta\_image.tsv}. The match uses the block's own
\texttt{witness\_poly} in place of the driver's local-algebra witness, which is exact under the shipped
\texttt{consume\_block\_locality = true} (slot-11 WP-14). (h) \emph{differs}: $\ell\in\{2,3\}$ is refused
by \texttt{hecke::NotImplementedInThisBuild} (console exit~$5$), not by an \texttt{ENGINE\_CHECK\_MSG} naming
\HA{lem:h2torsion} (exit~$3$); the guard's message states the reason (the sequence is not available on the
model, \HA{rmk:thetasmallell}). The ring refusals as built are the first item of \S\ref{sec:ring}'s
``Published v1.12.0'' paragraph.
\emph{v1.12.1 (the slot-11 WP-15 referee's m-1, the ruling on F-2): the specification stands, and the
build now follows it.} $\ell<5$ is a \emph{mathematical precondition}, not a build capability: on this
model the equivariant cochain functor is not exact at $\ell\in\{2,3\}$ and the model does not compute the
group long exact sequence there (\HA{rmk:thetasmallell}), which is definitional unavailability --- the
class of a field accessor on a ring handle, exit~$3$ --- whereas exit~$5$ promises an implementation that
can never come. Since slot-11 WP-17 \texttt{DicksonThetaCell} refuses $\ell<5$ by
\texttt{ENGINE\_CHECK\_MSG} naming \HA{lem:h2torsion} (a \texttt{std::logic\_error}, console exit~$3$), and
\texttt{HeckeDicksonThetaValidation}'s refusal arm asserts that kind and that name. The ``differs''
sentence above is the v1.12.0 record of the build before that change.

%======================================================================
\section{Scalar extension to \texorpdfstring{$\Flr$}{F(l\textasciicircum r)}}\label{sec:field}
%======================================================================

\subsection{Why, and when}
The residue field of $\T_\m$ is $\F=\Flr$, $r=f_\m$; the eigenvalues $a_q(\m)$ live there, and
``the'' simple $\T_\m$-module is $\F$. Over $\Fl$ every object is a Galois \emph{orbit}
(irreducible factors, blocks of dimension $r\cdot(\text{length})$). Extension is needed
\emph{only} when some $r_\m>1$, and then only to (i) report eigenvalues as elements, (ii)
compare characters across different $n$ in the atlas, (iii) present the structure maps
\ref{D4} as $\F$-matrices. Nothing in \S\ref{sec:alg} needs it:

\begin{claim}[Extension commutes with everything]\label{cl:split}\PROVED\
Let $V=H_\m$ over $\Fl$ with $\T_\m/\m=\F$. Then $\T_\m\otimes_{\Fl}\F\cong\prod_{\sigma\in\mathrm{Gal}(\F/\Fl)}\T_{\m,\sigma}$
with $r$ local factors permuted simply transitively by $\Frob$, each with residue field $\F$,
and $V\otimes_{\Fl}\F=\bigoplus_\sigma V_\sigma$ accordingly. For every $\sigma$:
$\dim_\F V_\sigma=\dim_{\Fl}V/r$, the Loewy and socle layer dimensions of $V_\sigma$ over $\F$
are those of $V$ over $\Fl$ divided by $r$, the Jordan type of $t_j-a_{p_j}$ on $V_\sigma$ is the
partition whose $s$-th nullity is $\dim_{\Fl}\ker\pi_j(t_j)^s/r$ --- i.e.\ the S10
generalized-Jordan partition of $t_j$ w.r.t.\ $\pi_j$ (which is in units of $\deg\pi_j$) with
every multiplicity divided by $r/\deg\pi_j$; the two coincide only when $\deg\pi_j=r$ ---, the
embedding dimension and
Loewy length of $\T_{\m,\sigma}$ are those of $\T_\m$, and $\T_{\m,\sigma}$ is Gorenstein iff
$\T_\m$ is. In particular \emph{no extension class is created or destroyed by extending
scalars}, and the record of \S\ref{sec:classdata} computed over $\Fl$ is the record of every
$V_\sigma$.
\end{claim}
\begin{proof}
$\F\otimes_{\Fl}\F\cong\F^r$ via $a\otimes b\mapsto(\sigma(a)b)_\sigma$; tensoring the local
ring $\T_\m$ (an $\F$-vector space through any coefficient-field lift, which exists since
$\F/\Fl$ is separable) gives the product decomposition, and $\m\otimes\F$ is the product of
the maximal ideals. All listed invariants are dimensions of $\F$-subquotients defined by
$\T_\m$-linear operations, which commute with base change; the Galois group permutes the
factors, so all factors have equal invariants, and the $\Fl$-dimension of $V$ is the sum over
$\sigma$ of $\dim_\F V_\sigma$. Jordan types: on $V_\sigma$ only the factor
$t_j-a^{\sigma}_{p_j}$ of $\pi_j(t_j)=\prod_\tau(t_j-a^{\tau}_{p_j})$ is nilpotent, the others
are invertible, so $\ker(t_j-a_{p_j})^s|_{V_\sigma}=\ker\pi_j(t_j)^s|_{V_\sigma}$, and the
$\Fl$-nullity of $\pi_j(t_j)^s$ on $V$ is $r$ times this common value.
\end{proof}

\begin{decisionbox}
\textbf{Compute over $\Fl$ from start to finish; extend scalars only in the report layer.}
All arithmetic in the certified pipeline stays \texttt{uint32}/\texttt{uint64} mod $\ell$.
This is the route recommended by the capability inventory \S4: \emph{no new element type in the
engine}; $\F$ is represented as the $\Fl$-algebra $\Fl[X]/(\pi)$ through the existing
\texttt{FpPolyT} (\texttt{powmod}, \texttt{frobenius\_powmod}), on the per-maximal-ideal blocks only
($\dim\lesssim200$), and never on a result path. Two tiers: (i) the \emph{block's own field}
$\Fl[X]/(\pi_w)$, in which the character is already available as the polynomials $P_j$ of
Claim~\ref{cl:resfield} with zero new arithmetic --- always emitted, as
\texttt{coefficient\_field}; (ii) the re-embedding into the canonical field of \S\ref{sec:canon}
for cross-run comparison, the only place needing polynomials over $\F$ --- \emph{since v1.4.0
the default report form} (\texttt{split\_field\_report = canonical}, owner decision OQ10; through
v1.3.3 tier (i) was the default), \emph{conditional on E6's performance gate}: at the atlas cells
$\ell\in\{5,7,11,13\}$ (every even $n\le\ell-1$) and at $(10,691)$, $(160,691)$, $(198,691)$ the
canonical-field time must satisfy
$\texttt{canonical\_field\_s}\le0.10\times\texttt{total\_s}$ on every cell, with the numerator,
the denominator and the protocol pinned by [X] \S6.2~E6~(d), which is the one governing statement
of the gate (v1.7.1: that item is written into [X] at v1.10, the slot-6 close; until then its text
is \texttt{docs/implementation/e6\_gate\_brief.md} \S6.1, which is binding,
\texttt{docs/implementation/hecke\_extensions\_slot6\_plan.md} \S4.3; v1.7.2: [X] v1.10 carries it,
and the interim pointer is superseded) (owner decisions \texttt{[OD]} E6-1 and E6-2 of 2026-09-22,
\texttt{docs/implementation/e6\_gate\_brief.md}); $\texttt{own\_field\_s}$ is measured and
reported beside it and is not a bound. (Through v1.6.2 the gate read
$\texttt{canonical\_field\_s}\le0.10\times\texttt{ext\_s}$ \emph{and}
$\le2\times\texttt{own\_field\_s}$; since tier~(i) is by this document's own words ``zero new
arithmetic'', the second clause compared a certified field computation with a rendering and was
not meetable, and the first had a denominator of microseconds on 20 of the 23 cells --- both
measured on all 23 cells in \texttt{docs/implementation/slot5/WP-7\_report.md} \S4. E6-1 changed
the denominator to \texttt{total\_s}, the wall of \texttt{HeckeExtDriver::run()}, and deleted the
second clause; the constant $0.10$ is unchanged.) If it fails, the default flips to
\texttt{own\_field} for the first campaign (\emph{superseded} by rule~Z below; v1.8.1), the measurement goes to the owner, and the $\Flr$
element type is the reopening question ([X] \S0.3, \S2.2 Stage~6). \emph{The cell list and the
failure rule} (v1.8.0; slot-7 owner ruling~4 of 2026-09-23, slot-7 plan \S10; this document only
points to them). The cells are the $23$ above plus, as \textbf{gating} cells, the first
campaign's highest-$r$ cells, those with \texttt{canonical\_field\_r\_max} $\ge10$ in WP-4's
$r_{\max}$ scan (\texttt{docs/implementation/slot7/WP-4\_rmax\_ell691.tsv}, read by
\texttt{e6\_measure.sh --cells-from}; option C1). (v1.8.1: the threshold $10$ is WP-4's reading
of the ruling's ``highest-$r$ cells'', the default of \texttt{e6\_measure.sh --r-min}; [X] v1.11
pins it. With it the list is $23+19-1=41$ cells, $(198,691)$ being both one of the $23$ and at
$r_{\max}=10$.) $(210,691)$ and $(256,134217757)$ are
\textbf{reported only}, since neither is a campaign cell ($n=210>200$; $\ell=134217757$ is not a
campaign prime; v1.8.1). The failure rule is option Z: a miss in the verdict window goes to the
owner before the campaign's publication run and never flips the report form by itself. This
replaces the automatic ``flips'' of the previous sentence. The owner then chooses among three
options: flip the campaign; hold \texttt{canonical} with (d) recorded not met; or delay the
campaign for a stronger lever. A miss on a reported cell flips nothing. The list and the rule are written into [X] \S6.2~E6~(d) at [X] v1.11. Until then,
the slot-7 plan's \S10 and \S2.4 row T3 are binding (v1.8.2: superseded --- [X] v1.11 carries
both, and pins the threshold $10$). \emph{Status}: W1 measured the worst gating
cell, $(170,691)$, at $0.0104$ over the $41$ gating cells, so the bound holds and the default
stands (Algorithm~\ref{alg:split}'s cost note); E6 is declared at the slot-7 close (v1.8.2). The defining polynomial is
the setting \texttt{field\_polynomial = least} (Definition~\ref{def:canon}); \texttt{conway} is
reserved, parsed and rejected (v1.7.2: parsed and validated as a setting, its fingerprint
moving; its consumer \texttt{FieldPolynomialFromSettings} raises
\texttt{hecke::NotImplementedInThisBuild}; v1.8.2: every production entry point --- the
field's constructor, \texttt{--bench ext}, the certificate, the atlas and the record ---
routes through it since slot-7 WP-4). Working in $\F$ from the start would multiply
every cost by $r^2$--$r^3$, would need a new engine width (the \texttt{ModCtx}/kernel/tape layers
are hard-wired to integer residues, inventory \S4.1), and would certify nothing that
Claim~\ref{cl:split} does not already give.
\end{decisionbox}
\noindent\emph{Where the two tiers sit in the record} (v1.8.0-pre.1): tier~(i) is the per-ideal
\texttt{coefficient\_field}, always; tier~(ii) is the per-ideal \texttt{split}, the
\texttt{hecke\_canonical\_split/v1} fragment of Algorithm~\ref{alg:split} verbatim, present only
under \texttt{split\_field\_report = canonical}. Under \texttt{own\_field} the record carries
tier~(i) alone. The one statement of the layout, with its presence rule, is the paragraph ``The
stage-2 keys'' of \S\ref{sec:schema}; it holds whichever report form the gate above selects.

\noindent\emph{The coefficient ring is a different axis} (v1.9.0-pre.1). This section extends the
\emph{residue field} of $\T_\m$ over a run whose coefficients are $\Fl$; it says nothing about
coefficients in $\Z/\ell^k$. The setting \texttt{coefficient\_ring} (E8), its grammar, its refusal
rule and the objects a $\Z/\ell^k$ run computes are specified in \S\ref{sec:ring}. Under a
non-default ring nothing in this section runs in slot~8: the canonical split, the atlas and the
extension driver are field notions, and each refuses by the rule of \S\ref{sec:ring}~(b).

\subsection{Canonical, deterministic choice of \texorpdfstring{$\F$}{F}}\label{sec:canon}
\begin{definition}[Project-canonical defining polynomial]\label{def:canon}
$C_{\ell,r}\in\Fl[Y]$ is the \texttt{CanonicalLess}-least monic irreducible polynomial of
degree $r$ (order: degree, then lexicographic on the low$\to$high coefficient vector),
found by enumerating candidates in that order and applying \texttt{CertifyIrreducible};
$\F:=\Fl[Y]/(C_{\ell,r})$.
\end{definition}
\begin{claim}\label{cl:canon}\PROVED\
$C_{\ell,r}$ is a deterministic function of $(\ell,r)$ alone, reproducible by any
implementation of the stated order, and every step of its construction is certified. It is
\emph{not} in general the Conway polynomial (which imposes compatibility across subfields
and a different ordering, and is tabulated only for small $\ell,r$ \cite{lubeck}); when the
Conway polynomial is available the record may name it as well, and the record always stores
$C_{\ell,r}$ itself, so the identification is part of the data.
\end{claim}

\begin{algorithm}[Splitting a block, report layer only]\label{alg:split}
Input: a certified block with witness polynomial $\pi_w$ ($\deg\pi_w=r$) and the polynomials
$P_j$ of Claim~\ref{cl:resfield}.
\begin{enumerate}
\item Find a root $\rho\in\F$ of $\pi_w$: $\pi_w$ splits completely in $\F$; use
  equal-degree splitting over $\F$ with a \emph{fixed} candidate sequence
  ($\gcd(\pi_w(X),(X+s)^{(\ell^r-1)/2}-1)$, $s=0,1,Y,\dots$), and certify $\pi_w(\rho)=0$ in $\F$
  --- the repository's find-then-certify pattern. The $r$ roots are $\rho^{\ell^i}$; take
  $\rho_{\mathrm{can}}:=$ the \texttt{CanonicalLess}-least of them (as coefficient vectors).
\item Split eigenvalues: $a_q(\m)=P_q(\rho_{\mathrm{can}})\in\F$ for $q\in S$; the orbit is
  $\{(P_q(\rho^{\ell^i}))_q\}_i$, all reported.
\item $\Frob$ acts on $\F$ by $Y\mapsto Y^\ell\bmod C_{\ell,r}$ and on the split data by
  $\rho\mapsto\rho^\ell$; it permutes the $r$ conjugate blocks of Claim~\ref{cl:split}
  transitively. The record reports both forms: the $\Fl$-form (factors, orbit length $r$) and
  the split form (one canonical tuple plus the Frobenius orbit).
\item Structure maps in $\F$-coordinates: choose the $\F$-basis of each Loewy layer given by
  the $\Fl$-basis reduced modulo the action of $\rho_{\mathrm{can}}$ (the layer is an
  $\Fl[w]/(\pi_w)$-module); the $\F$-matrix of $\theta(n_a)$ is then read off. (This step is
  optional; ranks are already known by Claim~\ref{cl:frank}.)
\end{enumerate}
\end{algorithm}
Only step 1 requires arithmetic in $\F$ (univariate polynomials over $\F$); it is new code
confined to the reporting layer, $O(r^3\log\ell)$ $\F$-operations per block $=O(r^5\log\ell)$
$\Fl$-multiplications with schoolbook arithmetic in $\F$; \emph{not} negligible at high $r$, see
[X] \S6.2~E6~(d). (v1.8.0; slot-7 WP-4, commits \texttt{fa829c7}, \texttt{cab4e38}. Through
v1.8.0-pre.1 this sentence read ``$O(r^3\log\ell)$ per block --- negligible'', with no unit. The
unit is $\F$-operations, and ``negligible'' was \textbf{refuted} before the slot-7 fix on the two
high-$r$ cells measured, at $r=14$ and $r=19$ (v1.8.1; through v1.8.0 this read ``at $r\ge11$'',
but no cell at $r\in\{11,12,13\}$ was measured before the fix). The canonical-field time exceeded
the $0.10\times\texttt{total\_s}$ bound of \S\ref{sec:field}
on both high-$r$ cells measured: $0.26$ at $(210,691)$, $r=14$, and $1.12$ at $(256,134217757)$,
$r=19$ ([X] v1.10 \S6.2~E6~(d)). \emph{Neither is a campaign cell} ($n=210>200$, and
$\ell=134217757$ is not a campaign prime), so under slot-7 owner ruling~4 both are
\textbf{reported only} and are not among the gating cells below (\S\ref{sec:field}; v1.8.1). The
first campaign's own $\ell=691$, $n\le200$ cells reach $r=15$
(\texttt{docs/implementation/slot7/WP-4\_rmax\_ell691.tsv}). WP-4's operation counter,
\texttt{HeckeCanonicalFieldValidation} arm~H \emph{as of commit \texttt{fa829c7}}, which counted
the unmodified route, found
the shipped route at its own count, a fitted exponent of $4.8$--$5.6$ in $r$ in
$\Fl$-multiplications (over $r=10..20$ at $\ell\in\{5,691,134217757\}$), with no defect above it.
The fix keeps the candidate sequence and every
value: it evaluates $u\mapsto u^\ell$ in $\F[X]/(f)$ as a precomputed $\Fl$-linear map and
multiplies in a flat quotient ring, so the bound above is unchanged and its constant drops, at
$r=19$ from $144.1$M (arm~H at \texttt{fa829c7}) to $12.1$M (arm~H at this revision)
$\Fl$-multiplications at $\ell=691$. (v1.8.1: at this revision arm~H counts the fixed route and
fits $4.89$, $4.23$ and $3.87$ at $\ell=5$, $691$, $134217757$ over $r=10..20$; it counts the
pre-fix route only for $r\le10$, as its byte-identity oracle, and fits nothing to it.) After the
fix, the slot-7
verdict window W1 measured the worst gating cell, $(170,691)$ at $r=14$, at ratio $0.0104$ over
the $41$ gating cells of slot-7 owner ruling~4: \textbf{E6~(d) MET}. The two reported cells
measured $0.0062$ at $(210,691)$ and $0.0151$ at $(256,134217757)$ in the same window, under the
bound but not gating (v1.8.1).)

%======================================================================
\section{Output schema}\label{sec:schema}
%======================================================================
Two records exist (v1.4.0). The \emph{per-degree} record \texttt{hecke\_report/v2}
(\texttt{Hecke/HeckeReport}, since E0b, slot 2 --- the golden-visible writer, about the bulk
complex) carries level (c) with, per block, \texttt{local}, \texttt{witness\_poly},
\texttt{coefficient\_field} $=\{$\texttt{w\_hat\_exponent}, \texttt{P}$\}$, the corrected
multiplicity note, \texttt{level\_c.locality} $=\{$\texttt{mode}, \texttt{budget},
\texttt{source}, \texttt{complete}, \texttt{splits}$\}$ and the four degree fields below at
record level; it carries no resumable state --- the two scalars \texttt{complete} /
\texttt{splits} say whether to look for one, and the state is R3's sidecar material (slot-2
WP-2). The \emph{extension} record \texttt{hecke\_ext/v2} (E2; \texttt{v1} until v1.6.0 of this
document --- the bump is the record-level \texttt{u\_ell} and the per-ideal
\texttt{groups} list below, made before the first emission) is one JSON per $(n,\ell)$
cell, extending that schema and \texttt{Hecke/README.md} (conventions kept: polynomial
coefficient lists low$\to$high, matrices as arrays of rows, entries in $[0,\ell)$,
\texttt{residue\_degree} versus \texttt{residue\_degree\_lower\_bound} never conflated;
hand-rolled writer):
{\footnotesize
\begin{verbatim}
{ "schema": "hecke_ext/v5", "ell": 691, "n_sym": 10, "weight": 12, "sign": "+|-|all",
  "commit": "<hash>", "model": "sl2z_7vertex", "iota_sign": "eisenstein_in_H1_plus",
  "settings": { "locality": {"value": "budget=64", "source": "default"},
                "u_ell":    {"value": "...", "source": "default|spec|console"},
                ...,        (every REGISTERED key, in registry order; see below)
                "presentation": {"value": "off", "source": "default"},
                "predictions":  {"value": "annotate", "source": "default"},
                "fingerprint": "0000000000000000" },
  "coefficient_ring": "F_ell", "resumable_items": 0,
  "predictions_mode": "annotate",
  "predictions": {"kummer": [{"k": 12, "applicable": true, "divides": true,
                              "k_reduced": 12, "b_over_k_mod_ell": 0, "reason": "direct"}],
                  "glover_factors": {"schema": "hecke_glover_factors/v1", "n": 10,
                                     "ell": 691, "dimension": 11, "length": 1,
                                     "factors": [{"sym": 10, "det": 0, "multiplicity": 1}]}},
  "anemic_pool": [2,3,5,7,11,13,17,19,23,29,31],
  "s_sep": {"status": "atlas_unavailable"},
  "u_ell": {"included": true, "reason": "within_budget", "predicted_mb": 104.06,
            "budget_mb": 2048, "model_version": "u_ell_rss/v2-2026-09-17"},
  "convention_check": {"tau_mod_ell_ok": true, "eis_eigenvalue_ok": true},
  "complexes": {
    "bulk":     {"q_max": 2, "truncated": false, "q_max_source": "derived", "degrees": [0,1,2]},
    "boundary": {"q_max": 1, "truncated": false, "q_max_source": "derived", "degrees": [0,1]},
    "pair":     {"q_max": 2, "truncated": false, "q_max_source": "derived", "degrees": [0,1,2]} },
  "groups": { "H0": {"dim":0}, "H0_bd": {"dim":1, "status":"unavailable"|"computed"},
              "H1c": {"dim":3, "status":"unavailable"}, "H1": {"dim":3}, "H1_bd": {"dim":1, ...},
              "H2c": {"dim":0, "status":"unavailable"},
              "dim_M_k": 2, "dim_S_k": 1, "dim_dickson": 0, "dim_H1_shriek": 2 },
  "les": { "exact": true, "ranks": {"r0": 0, "r1": 1}, "hecke_equivariant": null },
  "operators_computed": [2,3,5,7,691],
  "ideals": [
    { "id": 0, "sign": "all", "residue_degree": 1,
      "residue_degree_certified": true, "local": true,
      "groups": [{"group": "H1", "own": true, "match": true, "from": 0,
                  "witness": {"pi_w": [24,1], "N": 2, "N_used": 2}, "failure": null}],
      "canonical": {"H1": {"H1_shriek": "codim 1 (ker r^1; rank r^1 = 1)",
                           "eisenstein": "socle E, 2-dim on the unsplit H1; a line per sign",
                           "dickson": "n/a", "L_image": "n/a"}},
      "residue_witness": [[2,1]], "witness_poly": [24,1],
      "coefficient_field": {"w_hat_exponent": 1,
                            "P": {"2": [667], "3": [252], "5": [684], "7": [531], "691": [1]}},
      "split": {"schema": "hecke_canonical_split/v1", "ok": true, "ell": 691, "r": 1,
                "field_polynomial": [0,1], "witness_poly": [24,1], "candidates_used": ..,
                "rho_canonical": [667], "conjugates": [[667]], "primes": [2,3,5,7,691],
                "eigenvalues": [[667],[252],[684],[531],[1]],
                "frobenius_orbit": [[[667],[252],[684],[531],[1]]]},
      "label": {"status": "atlas_unavailable"},
      "char_Fl": {"2": {"pi": [24,1]}, "3": {"pi": [439,1]}},
      "dims": {"H1": 3, "H1c": null, ...},
      "algebra": {"dim": 2, "embdim": 1, "nilpotency": 2, "gorenstein": true,
                  "hilbert": [1,1], "basis": [[..],[..]]},
      "generators_certified": [2],
      "certificate": "cyclic"|"commutant"|"sturm(B')"|"dickson"|"boundary",
      "certificate_routes": {"a": "cyclic|not_cyclic",
                             "b": "excluded_by_socle_bound|done|over_budget",
                             "c": "found|none_at_cap|not_applicable|off",
                             "d0": "passed|failed:<class>|not_applicable:u_ell|not_run",
                             "d1": "passed|not_injective|failed:<p>|not_run:budget|not_applicable:u_ell|not_run"},
      "boundary_reps": [],     (R(ell) iff d1 is passed or failed:<p>; v4)
      "boundary_data": {"rank_r1": D, "r1": [[..]], "t": {"<p>": [[..]], ..}},
                               (present iff d1 = passed; p in boundary_reps; v5)
      "final": true|false|"heuristic", "final_reason": null|"<reason>",
      "sturm": {"B_prime": 2, "cap": 24, "closure_dim": 2, "dim_M_k": 2,
                "mode": "nonfinal", "outcome": "found", "sturm_primes": [2],
                "proof_range": true, "proof_range_source": "cell_local",
                "resumable": {"kind": "sturm_bound", "complete": true, "budget": {"cap": 24},
                              "reached": {"B_reached": 2}, "budget_history": [{"cap": 24}]}},
      "tangent": {"2": "active", "3": "active", "5": "active", "7": "active", "691": "active"},
      "modules": { "H1": {"loewy": [2,1], "socle": [2,1], "cyclic": false, "free": false,
                           "min_generators": 2,
                           "jordan": {"2": [2,1], "3": [2,1], "691": [2,1]},
                           "structure_maps": [{"piece": [0,1], "rank": 1,
                                               "active": [2,3,5,7,691],
                                               "matrices": {"2": [[23,0]], "3": [[120,0]],
                                                            "691": [[221,0]]}}]},
                   "H1c": null, ... },
      "predicted": ["eisenstein_cuspidal:k=12"], "confirmed": [true],
      "unpredicted": [] } ],
  "identities": {"passed": 7, "failed": 0},
  "certificates_ran": ["locality","exhaustive","cyclic|commutant|sturm","module"],
  "hashes": {"A1_T2": "fnv64", ...} }
\end{verbatim}}
\paragraph{\texttt{final} is tri-state (v1.4.0; [X] \S1.3.1 R3.2; owner decisions OQ8 /
OQ15, slot-1 confirmations~6 and~8).} Per ideal,
$\texttt{final}\in\{\texttt{true},\texttt{false},\texttt{"heuristic"}\}$. \texttt{true} iff
\texttt{local} is true, every applicable certificate verifies with $0$, and the generation
certificate is \texttt{cyclic} or \texttt{commutant}, or \texttt{sturm(B\textquotesingle)} with
$B'$ existing on a non-Dickson $\m$ of $X_4$ (Claim~\ref{cl:sturmmodl}), \textbf{or (v1.13.0,
\texttt{hecke\_ext/v4}) \texttt{dickson} / \texttt{boundary}}: Certificate~\ref{cert:dickson} (d0) passed on an
$X_1$ block, or Certificate~\ref{cert:boundary} (d1) passed on an $X_4$ block --- each a proof of the same
predicate as (a)--(c), so \texttt{true} in every mode, the default \texttt{closure} included; a failed (d0)
/ (d1) leaves the reason (a)--(c) produced. \texttt{"heuristic"}
only under the user-selected modes \texttt{final\_mode = sturm\_heuristic} (the heuristic range:
Dickson-type $\m$, $X_2,X_3,X_5,X_6$) or \texttt{sturm\_unavailable = heuristic} ($B'$ missing or
cap hit): the ideal's only generation evidence is a heuristic sweep, recorded with the primes
actually swept (\texttt{sturm\_primes}) and the mode (\texttt{sturm.mode}) --- never in the
default configuration. \texttt{false} otherwise, with a reason from the fixed list
\texttt{\{not\_local, sturm\_unavailable, sturm\_heuristic, closure\_unreachable,
certificate\_failed:<code or certificate name>, budget\_exhausted:<kind>, interrupted\}}
(\textbf{v1.5.0}, corrected from \texttt{certificate\_failed:<code>}: the payload is
\emph{either} the non-zero clause index a \texttt{verify()} returned \emph{or} one of the
closed list of certificate names \texttt{\{cyclic, commutant, sturm, locality, module,
identity, dickson, boundary\}} (the last two v1.13.0; the shipped driver lets a failed (d0) / (d1) fall
through to the existing reason, so it emits neither). The name form is not a slip: a certificate \emph{route} may report ``the
sweep ran and its answer was no'' --- \texttt{certificate\_failed:sturm}, emitted by
Certificate~\ref{cert:sturm}'s route since slot~2 and pinned by
\texttt{HeckeExtSturmSweepValidation} --- with no \texttt{verify()} having failed and so
no clause index to give. [C] \HA{sec:recordschema} item~6(d) carries the same rule; slot-4
WP-1 \S6 item~1):
\texttt{budget\_exhausted:locality} is the budget-mode locality survivor
(Algorithm~\ref{alg:stage2}(2)), \texttt{budget\_exhausted:commutant} the budgeted structured
system of Certificate~\ref{cert:comm}, \texttt{closure\_unreachable} the proof that (b) fails
(the socle bound, or a computed $\dim C>\dim A_S$) with (a) failing and (c) off its proof range,
\texttt{interrupted} a strict loop stopped at a step boundary; \texttt{not\_local} is reachable
only through strict mode's interruption path, since a block proved non-local splits. A
\texttt{budget\_exhausted} reason always comes with \texttt{resumable.complete = false};
\texttt{closure\_unreachable} and \texttt{sturm\_heuristic} are not budget outcomes and carry no
state. In the default mode (\texttt{final\_mode = closure},
\texttt{sturm\_unavailable = nonfinal}) an ideal that cannot reach closure is \emph{emitted} with
\texttt{final = false}, its reason and its resumable object, listed in the campaign's
\texttt{analysis.txt}, and never withheld. \texttt{final} never reads \texttt{resumable}; no
field is branched on by the pipeline that wrote it ([C] rule A6). Through v1.3.3 this paragraph
read ``a record is final at $\m$ iff \ldots; otherwise it is marked \texttt{"final": false} with
the reason'' --- a two-valued field with no fixed reason list.

\paragraph{Settings, resumable state, coefficient ring (v1.4.0).} \texttt{settings} carries
every mathematical setting in force with its \texttt{source} $\in\{$\texttt{default},
\texttt{spec}, \texttt{console}$\}$ --- \texttt{final\_mode}, \texttt{sturm\_unavailable},
\texttt{sturm}, \texttt{commutant\_budget}, \texttt{locality}, \texttt{field\_polynomial},
\texttt{split\_field\_report},
\texttt{coefficient\_ring}, \texttt{predictions}, \texttt{presentation}, \texttt{u\_ell},
\texttt{integral\_fallback}, and \texttt{q\_max} where settable ([X] \S1.5; owner decision R1:
these are the \emph{second} family of switches, each of which may change a published value,
as no \texttt{ENGINE\_*} variable may) --- and the \texttt{fingerprint} of the non-default
entries ($0$ at the defaults), which enters every extension hash input and the campaign job id
(the \texttt{\_s<hex>} suffix), so that only the default configuration is golden and a golden
mismatch can never be explained by a setting. (v1.8.0-pre.1: the list just given is [X]
\S1.5's inventory of the family; the block carries the \emph{registered} members of it, in the
shape stated in item~(f) of ``The stage-2 keys'' below, which corrects the example through
v1.7.2.) Every budgeted item writes at its own site
\texttt{resumable} $=\{$\texttt{kind} $\in\{$\texttt{sturm\_bound}, \texttt{locality},
\texttt{commutant}, \texttt{u\_ell}, \texttt{fill}$\}$ (v1.9.0-pre.1: and
\texttt{ladic\_level}, \S\ref{sec:ring}~(d), for the day a record carries it), \texttt{complete}, \texttt{budget},
\texttt{reached}, \texttt{budget\_history}, \texttt{state} $\mid$ \texttt{state\_ref}
$=\{$\texttt{path}, \texttt{fnv64}$\}\}$ --- inline if $\le64$\,KiB, else a sidecar named by
the record layer and hashed by the one FNV-1a-64 of \texttt{Fnv.h} --- and the record carries
\texttt{resumable\_items}, the count of incomplete items ([X] \S1.3.1 R3.1); a resume with a
higher budget is a new recorded setting, a resume under a changed non-budget setting is refused.
\textbf{Two rules on the sidecar} (v1.5.0; slot-4 WP-1 \S6 items 2 and 3; [C]
\HA{sec:recordschema} item~6(b)): \texttt{state\_ref.path} is \emph{relative to the
directory holding the record}, so that the pair of files moves together, and \emph{the
name must be injective} --- a record may carry several oversize states of the \emph{same}
kind, one \texttt{commutant} per ideal, so only the first of a kind may take the bare
\texttt{<record-stem>.state.<kind>.json} and every further one is disambiguated by an
index. A writer that has no stem, and so cannot name a file, inlines the state and says so
rather than emitting a path that names nothing.
\textbf{\texttt{phase\_reached}} (v1.5.0; slot-3 WP-3 finding~6): where a route reports the
phase it reached --- strict locality's \texttt{linear} / \texttt{berlekamp} /
\texttt{residue\_search} --- the record carries the \emph{per-block} value, under
\texttt{phase\_reached} inside that block, \emph{and} the run's high-water mark over every
block including the parents a split has since removed, under the separate record-level key
\texttt{phase\_reached\_high\_water}. They are different facts and neither substitutes for
the other; both are emitted only where a route that has phases ran.
\textbf{\texttt{groups} holds groups and four scalars} (v1.5.0, stated because a consumer
cannot tell them apart from the shape alone; slot-4 WP-1 \S6 item~5): every key of
\texttt{groups} other than the closed list \texttt{\{dim\_M\_k, dim\_S\_k, dim\_dickson,
dim\_H1\_shriek\}} is a cohomology group and carries \texttt{dim} and, where it is not
computed, \texttt{status}; those four are integers \textbf{or \texttt{null}}. A later
schema revision may move them to a sibling object; until it does, the closed list is the
rule.
\textbf{\texttt{null}, never dropped} (v1.6.0; slot-5 WP-1 finding~2; [C]
\HA{sec:recordschema} item~6(c)). The list is a constant of the schema, but
\texttt{dim\_H1\_shriek} is $\dim\ker r^{1}$, which does not exist until the boundary
complex does. A record that has not computed it emits it as \texttt{null}; dropping the
key would make the closed list a function of the cell, which is exactly the property that
made it a rule, and emitting $0$ would say something false.

\textbf{\texttt{canonical} is keyed by the ambient group, at \emph{ideal} level}
(v1.6.0; slot-5 WP-1 finding~1). Through v1.5.2 the example above nested
\texttt{"canonical"} \emph{inside} \texttt{modules.H1}, while the shipped writer emitted it
beside \texttt{modules}; the example is corrected to the writer's placement, and the object
is a map keyed by the ambient group (\texttt{"H1"}, \texttt{"H1\_plus"}, \ldots), which
carries the same information, reads like \texttt{groups}, and keeps \texttt{modules} a
verbatim per-group \texttt{ModuleStructure} serialisation with no string splicing. A group
whose canonical positions were not computed is absent from the map, not \texttt{null}
inside it --- \texttt{canonical} is keyed by what was computed, not by a closed list.

\textbf{\texttt{les.hecke\_equivariant} is the \emph{action}'s claim} (v1.6.0; slot-5 WP-1
finding~4; [C] \HA{sec:recordschema} item~6(g)). The cone knows nothing about $T_q$, so a
record carrying the cone but no cone action has not evaluated the flag. It is therefore
tri-state: \texttt{true} asserted, \texttt{false} evaluated and failed, \texttt{null} not
evaluated. Through v1.5.2 the example showed \texttt{true} without saying who asserts it,
and a cone-only record emitted \texttt{false}, which under [C] rule~A6 reads ``evaluated
and it did not hold'' --- the one thing it did not mean. The same rule governs
\texttt{exact} at a spot the cone did not assert.

\textbf{A \emph{resumed} run writes the settings block a \emph{fresh} run at the same
effective values would write} (v1.5.0; slot-4 WP-2 finding~4; [C] \HA{sec:recordschema}
item~6(a)). The rule is not implied by the paragraph above, and it decides a fingerprint:
an implementation that installs the budget it was resumed \emph{to} makes
\texttt{resume(2048)} on a refused $U_\ell$ emit
\texttt{"u\_ell": "budget\_mb=2048"} where a fresh run at the default emits
\texttt{"u\_ell": "auto"}, so the two records differ in their settings text, their
fingerprint and every hash downstream of it although they computed the same thing. The
binding rule is the equality: a resume to a value \emph{equal to the default} installs the
default's canonical text and any other value installs its own, whereupon ``a resume to the
default budget is byte-identical to a fresh default run, fingerprint included'' is a
statement the implementation asserts (\texttt{HeckeExtRamanujanValidation} arm~4, slot-4
WP-2).

\textbf{\texttt{u\_ell} is carried at \emph{record} level, with a single
\texttt{resumable}} (v1.6.0, \texttt{hecke\_ext/v2}; the repair v1.5.0 booked, approved as
owner ruling \texttt{[OD]} slot-4 close-out 5 of 2026-09-22 and made by slot-5 WP-1).
The admission decision of \S\ref{sec:alg} Stage~1 is a property of $(n,\ell)$ and of the
budget, not of a maximal ideal, and the record now says so: \textbf{one} \texttt{u\_ell}
object, carrying its verdict and --- when $U_\ell$ is refused --- \textbf{one}
\texttt{resumable}, beside \texttt{resumable\_items}, which is accordingly a count of
\emph{decisions}. No \texttt{ideals[]} entry carries a \texttt{u\_ell} object. A cell that
formed no verdict emits the placeholder \texttt{"u\_ell": \{"status": "unavailable"\}} ([C]
\HA{sec:recordschema} item~6(c)); an \texttt{included} verdict carrying a
\texttt{resumable}, a \texttt{resumable} with no verdict, and a reason outside the
enumerated three are each refused by the writer's own \texttt{verify()}.
\emph{What v1 did, and why this is a version bump}: through \texttt{hecke\_ext/v1} the
object and its \texttt{resumable} were attached to \emph{each} \texttt{ideals[]} entry, so
a four-ideal cell emitted one decision four times and \texttt{resumable\_items} counted it
four times --- a count of \emph{entries}, which [C] \HA{sec:recordschema} item~6(b) still
permits a schema that replicates a run-level verdict, provided it says so. The repair
changes the \emph{value} of an existing key rather than adding one, so it is a schema
version bump and not an amendment in place, and it was made \textbf{before} the first
emission of an extension record as a golden object, which is what made it free. A consumer
distinguishes the two shapes by the \texttt{schema} string alone.

\textbf{An \texttt{ideals[]} entry is per \emph{group}, and since
\texttt{hecke\_ext/v2} it reports the matching} (v1.5.0; v1.6.0 for the addition; slot-4
WP-2 finding~1, taken by slot-5 WP-1). Stage~2 of \S\ref{sec:alg} matches the per-group
decompositions by the cross-block evaluation test ([XA M5], \texttt{SameMaximalIdeal}), so
a maximal ideal is in principle an object spanning $X_1,\dots,X_6$; the writer still emits
\textbf{one entry per (group, block)}, each with a single \texttt{modules} entry --- that
shape is unchanged, which is why the addition needed no other change --- and an entry now
additionally carries a per-ideal \texttt{groups} list. Each element names the ambient group,
whether it is the block's \emph{own} group, the verdict of \texttt{SameMaximalIdeal}, the
ideal the test was run \emph{from}, and the \textbf{witness} it used --- the reduced $\pi_w$
low$\to$high, the exponent $N$, and the $N'$ actually used on \emph{this} block ([XA M5]:
nilpotency is bounded on the block the test evaluates \emph{on}, not on the one it comes
from) --- with the failure code on a non-match. The entry for the block's own group is the
identity match, which is what makes two such lists comparable. Through
\texttt{hecke\_ext/v1} the match count was a driver diagnostic that no key reported; [C]
\HA{sec:recordschema} item~6(i) now makes ``a matching a record performed is a matching it
reports, with its witness'' a rule, on the ground that a count cannot be re-checked by a
reader who has only the record.
\texttt{coefficient\_ring} $\in\{$\texttt{F\_ell}, \texttt{Z/ell\textasciicircum k}$\}$ names
the ring of the run (\texttt{F\_ell} today; E8, owner decision OQ18; v1.9.0-pre.1: the setting of
that name, its grammar, its refusal rule and the E8 objects are specified in \S\ref{sec:ring};
the record's top-level key keeps this meaning and reads \texttt{F\_ell} in every record of
slot~8, since every record writer refuses a non-default ring). \texttt{P} is an object
keyed by the prime, one reduced polynomial per generator --- on a single-operator presentation
it has one entry ($538/538$ blocks of the golden grid) --- so that
Certificate~\ref{sec:localcert}'s clause ``for every $j$'' is checkable from the record.

\paragraph{The stage-2 keys (v1.8.0-pre.1; slot-7 WP-12a, spec first).} What follows
specifies every key that slot-7 WP-3 (stage~2 of the staged re-baseline, [X] \S8) and slot-7
WP-2 (\texttt{presentation}) will emit, \emph{before} either emits it; a key not specified here
is not emitted ([X] \S8, slot-7 plan WP-3 ``content check''). Each key has its type, its
presence rule and its default. All of it rides on [C] \HA{sec:recordschema} item~6: rule~(c)
(a required key not computed is \texttt{null}; an unavailable object is a placeholder, never an
omission), rule~(d) (reasons are named from a closed list), and item~6(f), which leaves the
inventory to this section.
\begin{enumerate}[label=(\alph*)]
\item \textbf{\texttt{sturm.proof\_range}} (boolean or \texttt{null}) and
  \textbf{\texttt{sturm.proof\_range\_source}} (string), inside every \texttt{sturm} object,
  so present exactly when \texttt{sturm} is. \texttt{proof\_range} says whether the ideal lies
  in the range on which Claim~\ref{cl:sturmmodl} lets Sturm certify, i.e.\ whether the ideal is
  \emph{not} Dickson-type; it records that test alone --- the claim's other hypotheses (the
  group, the existence of $B'$) are not part of this field. The source is one of the closed
  list:
  \begin{itemize}
  \item \texttt{cell\_local} --- \textbf{the default, at every $\ell$} (slot-7 owner ruling~3a
    of 2026-09-23, $\ell\ge691$ included). Decided \emph{at the cell} by the method of
    Claim~\ref{cl:sturmmodl}: the ideal's anemic system over the pool \texttt{anemic\_pool}
    below, computed at the cell, is compared with the systems of the same cell's $H^0$ ideals
    over the same pool ($\mathrm{tors}\,H^1(\Ga;\Mn(\Z))\otimes\Fl\cong H^0$);
    \texttt{proof\_range = true} iff it equals none of them. No atlas is read.
  \item \texttt{caller\_input} --- the value was supplied by the caller
    (\texttt{ExtDriverOptions::proof\_range}) and is not verified. Never the default; a record
    carrying it says so, and a validation arm that forces a branch uses it.
  \item \texttt{method\_undefined} --- the method cannot be applied at this cell; the one case
    known is a missing $B'(k,\ell)$ (Claim~\ref{cl:bprime}). Then \texttt{proof\_range} is
    \texttt{null}, and the ideal's \texttt{final} is governed by \texttt{sturm\_unavailable}
    as before.
  \end{itemize}
  \texttt{dickson\_count} is the value every \texttt{hecke\_ext/v2} record emitted through
  slot~6 --- a caller input under another name (\texttt{proof\_range} was the unverified
  default \texttt{true}, the M1 qualification of [X] v1.10 \S0) --- and is \textbf{retired}: no
  record after stage~2 item~3 emits it, and a reader treats it as \texttt{caller\_input}.
  \emph{How the decision reaches \texttt{final}}: only through the Sturm route's outcome, which
  the consumer obtains by re-running that route with the decided range --- \texttt{final} never
  reads \texttt{proof\_range} ([C] rule A6; slot-6 WP-9 finding~4: flipping the flag alone and
  re-evaluating \texttt{final} reports no change, which is wrong).
  \textbf{\texttt{anemic\_pool}} (record level; list of primes, ascending; always present from
  the same item on): the pool $T=\{q\le31\}\setminus\{\ell\}$ over which the anemic systems of
  this item and of~(b) are compared. The operators of the pool are computed at the cell for
  that purpose; \texttt{operators\_computed} keeps its meaning (the operators the decomposition
  ran on).
\item \textbf{The atlas label and \texttt{atlas\_unavailable}} (slot-7 owner ruling~3b~(B)).
  $\Sigma(\ell)$ is built in-process at $\ell\in\{5,7,11,13\}$ (the range
  \texttt{HeckeAtlasValidation} pins) and at no other $\ell$; the record never reads a stored
  \texttt{hecke\_atlas/v1}, so golden stays self-contained, and making an atlas available at
  another $\ell$ later is a new owner-gated event, not a correction. Two keys:
  \begin{itemize}
  \item \textbf{\texttt{label}}, per ideal, always present. On a local ideal at an atlas
    $\ell$: \texttt{\{"status": "computed", "key": <the CanonicalLess-least Frobenius
    conjugate of the ideal's anemic system over \texttt{anemic\_pool}, one element of
    $\F_{\ell^r}$ per prime, low$\to$high>, "r": <r>, "system": <its id in $\Sigma(\ell)$>,
    "type": "eisenstein"|"cuspidal", "chi": [a,b]|null, "base": <id>,
    "theta\_twist": <j>\}}, the fields being those of the atlas's system record ([WP6-8] \S1:
    \texttt{type} is \texttt{eisenstein} iff the system equals some $\chi_{a,b}$ on $T$,
    \texttt{chi} the least such $(a,b)$, and the system is $\theta^{j}\otimes$ system
    \texttt{base}). \texttt{key} is the label proper --- a function of the system alone, which
    is what identifies systems across $n$; \texttt{system} and \texttt{base} are ids in the
    in-process atlas at the same settings, informative only. At every other $\ell$:
    \texttt{\{"status": "atlas\_unavailable"\}} and nothing else. On a non-local ideal, which
    has no character: \texttt{null}.
  \item \textbf{\texttt{s\_sep}}, record level, always present: at an atlas $\ell$,
    \texttt{\{"status": "computed", "primes": [...], "added": [...]\}} --- $S_{\mathrm{sep}}(\ell)$
    ascending, and $S_{\mathrm{sep}}(\ell)\setminus S_0$ ascending, the operators
    Algorithm~\ref{alg:stage2}(3) added (empty when $S_{\mathrm{sep}}(\ell)\subseteq S_0$, the
    vacuous case); the added operators appear in \texttt{operators\_computed}. At every other
    $\ell$: \texttt{\{"status": "atlas\_unavailable"\}}.
  \end{itemize}
\item \textbf{The split's location, stated once} (tier~(ii) of \S\ref{sec:field}; [OD] OQ10,
  E6-3). The per-ideal key \textbf{\texttt{split}} is the \texttt{hecke\_canonical\_split/v1}
  fragment of Algorithm~\ref{alg:split} \emph{verbatim}, as its one renderer writes it
  (\texttt{schema}, \texttt{ok}, \texttt{ell}, \texttt{r}, \texttt{field\_polynomial}
  $=C_{\ell,r}$, \texttt{witness\_poly}, \texttt{candidates\_used}, \texttt{rho\_canonical},
  \texttt{conjugates}, \texttt{primes}, \texttt{eigenvalues}, \texttt{frobenius\_orbit};
  elements of $\F$ as coefficient vectors of length $r$, low$\to$high). It is present
  \textbf{iff} the ideal is \texttt{local} \textbf{and} \texttt{split\_field\_report =
  canonical}; its \texttt{primes} are the keys of \texttt{coefficient\_field.P}, ascending, and
  it is computed through one \texttt{CanonicalField} per $(\ell,r)$ under the
  \texttt{field\_polynomial} in force. \texttt{ok = false} is a certificate failure of the
  split and is refused by the record's \texttt{verify()}. Under \texttt{split\_field\_report =
  own\_field} the record carries \texttt{coefficient\_field} (tier~(i)) only and no
  \texttt{split} key --- the key list is then a function of a recorded setting, as for
  \texttt{presentation}, not of the cell. \texttt{candidates\_used} makes the fragment depend
  on the root finder's candidate sequence as well as on $\rho_{\mathrm{can}}$, so that sequence
  is part of the record (slot-7 plan \S2.3 item~2). \emph{Retired}: the example through v1.7.2
  showed \texttt{field} $=\{$\texttt{C\_ell\_r}, \texttt{conway\_available},
  \texttt{rho\_can}$\}$ and \texttt{char\_split}; no writer emitted either, and \texttt{split}
  replaces both (\texttt{conway\_available} returns only with the Conway tables that
  Claim~\ref{cl:canon} permits the record to name).
\item \textbf{\texttt{predictions\_mode} and the prediction objects} (Stage~6; [X] \S1.5,
  OQ14). \textbf{\texttt{predictions\_mode}} (record level, string, always present) is the value
  of the \texttt{predictions} setting in force, \texttt{annotate} (the default) or
  \texttt{assert}; it is emitted beside \texttt{settings} because golden sees
  \texttt{settings} only as its fingerprint (item~(f)). \textbf{\texttt{predictions}} (record
  level, always present) holds the cell's prediction \emph{inputs}, attached before the
  pipeline runs: \texttt{kummer}, a list, ascending in $k$, of Kummer verdicts
  \texttt{\{k, applicable, divides, k\_reduced, b\_over\_k\_mod\_ell, reason\}} ---
  \texttt{kummer\_verdict} of \texttt{HeckePredictions}, \texttt{k\_reduced} and
  \texttt{b\_over\_k\_mod\_ell} \texttt{null} unless \texttt{reason} is \texttt{direct} or
  \texttt{kummer\_congruence} --- holding at least the weight $k=n+2$
  (Claim~\ref{cl:harder}); and \texttt{glover\_factors}, the
  \texttt{hecke\_glover\_factors/v1} object of $\Sym^n$ under the standard-module convention of
  Claim~\ref{cl:glover}. Per ideal, always present: \textbf{\texttt{predicted}}, a list of
  prediction ids \texttt{"<name>:<argument>"} with \texttt{<name>} from the closed list
  \texttt{\{eisenstein\_cuspidal\}} at this version (Prediction~\ref{pr:eiscusp};
  \texttt{<argument>} is \texttt{k=<w>}, the weight whose verdict predicts it); a new name is an
  edit of this list before it is emitted. \textbf{\texttt{confirmed}}, index-aligned with
  \texttt{predicted}, each \texttt{true}, \texttt{false} (evaluated and refuted) or
  \texttt{null} (not evaluated). \textbf{\texttt{unpredicted}}, the ambient groups $X_i$ of this
  entry with $\m X_{i,\m}\ne0$ (a non-split piece, \S\ref{sec:classdata}) that no entry of
  \texttt{predicted} covers. Under \texttt{annotate} a refuted prediction is recorded and the
  run goes on; under \texttt{assert} it fails the run, and \texttt{assert} is used by
  validations, never by a campaign ([X] \S1.5, OQ14).
\item \textbf{\texttt{presentation}} (per ideal; slot-7 owner confirmation~15, audit
  finding~m-5). Present \textbf{only} under \texttt{presentation = on} (default \texttt{off},
  so absent from the default record and from golden); then on every ideal entry, the per-sign
  entries carrying the per-sign presentation. Its value is
  \texttt{\{"generator\_basis": "block", "relation\_basis": "algebra\_rref", "fragment":
  <hecke\_presentation/v1>\}}: the fragment is \texttt{HeckePresentation}'s text verbatim.
  \emph{Encoding}: a generator index $j$ is an index into the \textbf{block basis} --- the
  stored basis of $V=X_{i,\m}$ in which the driver computed $A_S$ and the entry's
  \texttt{modules}, the basis the Nakayama lift is taken in ([X] \S2.2 Stage~5) --- and
  \emph{not} an $H^q$-column; the object names that basis in \texttt{generator\_basis}. A
  relation is a $t$-tuple over the RREF $\Fl$-basis $b_0,\dots$ of $A_S$ (its entry
  $i\cdot\dim A_S+j$ is the coefficient of $b_j$ in $a_i$), named by \texttt{relation\_basis}.
  On a non-local entry the fragment is the unit's \texttt{available = false} rendering. The
  presentation is a derived object of stored bases and is \textbf{not} a hash input. Turning it
  on therefore leaves the record's \texttt{hashes} and the golden \texttt{ext\_ideals} and
  \texttt{ext\_certify} hashes unchanged. It does \emph{not} leave the record \emph{text}
  unchanged. Beyond the added objects, \texttt{presentation = on} is a non-default setting, so it
  moves the \texttt{settings} block and its fingerprint (item~(f)), and with them every hash of the
  record text. Remove the objects and restore the setting, and the text is the default text again
  ([X] \S5; \texttt{HeckePresentationValidation} arm~(k4) asserts all three statements, and
  arm~(k3) that the fingerprint moves). (v1.8.1: precisely, (k4) asserts that the three hashes are
  equal, that the \texttt{settings} member of the record text differs, and that the stripped text
  with the default \texttt{settings} member put back \emph{is} the default text; (k3) asserts that
  the \texttt{presentation = on} settings have a non-zero fingerprint. That every hash taken over
  the text moves --- the golden \texttt{ext\_json\_canonical} / \texttt{ext\_json\_full} texts
  among them, although golden runs at the defaults only --- follows from the second and is not
  separately asserted.) (v1.8.0; slot-7 WP-2 finding. Through v1.8.0-pre.1 this
  read ``leaves every hash of the record unchanged'', which is false of any hash taken over the
  text.)
\item \textbf{The \texttt{settings} shape} (audit finding~m-12; reconciled to the writer,
  \texttt{MathSettings::to\_json}). \texttt{settings} maps each key to an object
  \texttt{\{"value": "<canonical text>", "source": "<source>"\}} and ends with
  \texttt{"fingerprint": "<16 hex>"}: one entry per \emph{registered} key --- exactly
  \texttt{MathSettings::keys()}, in registry order, which is append-only --- and then the
  fingerprint, rendered as sixteen lower-case hex digits (\texttt{"0000000000000000"} at the
  defaults: the ``$0$'' of the paragraph above). A member of [X] \S1.5's inventory that is not
  registered is absent until its consumer registers it (\texttt{coefficient\_ring} is today a
  top-level key, not a settings entry; v1.9.0-pre.1: once slot-8 WP-3 registers it, it is also
  the \emph{last} settings entry, appended at the tail of the registry, and the top-level key
  stays, \S\ref{sec:ring}~(a)). Through v1.7.2 the example showed
  \texttt{\{key: value, \ldots, "source": \{\ldots\}, "fingerprint": "0"\}}, which no writer
  produced. \emph{Golden}: after slot-7 stage~1~(1) the golden harness reduces the
  \emph{top-level} \texttt{settings} object of its \texttt{ext\_json\_*} texts to
  \texttt{\{"fingerprint": \ldots\}} by a path-scoped redaction, so golden sees only the
  fingerprint and a key registered at its default moves no golden line; the record itself still
  carries the full block.
\item \textbf{\texttt{source} is binding and closed}: $\texttt{source}\in\{\texttt{default},
  \texttt{spec}, \texttt{console}\}$, the three values of [C] \HA{sec:recordschema} item~6(a)
  and [X] \S1.5; a fourth value is an edit of this document and of [X], never a package
  decision. \texttt{default}: the key was not set. \texttt{spec}: set by a spec --- a campaign
  spec, \emph{and the golden harness}, which stamps \texttt{spec} on every key it installs
  (slot-7 owner confirmation~13). \texttt{console}: set in a console session.
\item \textbf{The schema string.} Items (a)--(d) add keys to the default record
  (\texttt{anemic\_pool}, \texttt{s\_sep}, \texttt{label}, \texttt{split},
  \texttt{predictions\_mode}, \texttt{predictions}, \texttt{predicted}, \texttt{confirmed},
  \texttt{unpredicted}), so under slot-7 owner confirmation~10 the string becomes
  \texttt{hecke\_ext/v3}, as the example shows, in the last commit of stage~2 (stage~3); a
  record emitted before that commit is \texttt{v2} and carries none of them except the two
  \texttt{sturm} fields, which v2 already emitted. Every new key is covered by the record's
  \texttt{verify()}.
\item \textbf{The finality certificates: \texttt{hecke\_ext/v4} (v1.13.0; slot-12 S12-FIN,
  owner-approved 2026-09-26).} Every ideal carries \texttt{certificate\_routes.d0} and \texttt{.d1}
  (Certificates~\ref{cert:dickson}, \ref{cert:boundary}; closed value lists as in the example) and the list
  \texttt{boundary\_reps} (empty unless (d1) reached its representatives), after
  \texttt{generators\_certified}; \texttt{certificate} may be \texttt{"dickson"} or \texttt{"boundary"}.
  The routes run on the unsplit record only (a per-sign ideal carries \texttt{not\_run}, as for (b) and
  (c)), at $\ell\ge5$, $n$ even, and only where (a)--(c) did not pass. \texttt{verify()} refuses (code~30)
  a value outside the lists, a route on the wrong group or sign, \texttt{passed} disagreeing with the
  verdict's input, \texttt{not\_applicable:u\_ell} disagreeing with $\ell\in$ \texttt{operator\_primes},
  \texttt{boundary\_reps} $\ne R(\ell)$ where (d1) reached them, a certificate name the routes do not
  carry; and (code~31) a (d0) that does not re-derive from the stored $t_p$. The string becomes
  \texttt{hecke\_ext/v4}; a \texttt{v3} record carries none of these keys, and its \texttt{final} is the
  (a)--(c) verdict. The published \texttt{v3} campaign sets (\texttt{ext\_first}, \texttt{ext\_u691}) are
  not re-emitted.
\item \textbf{(d1)'s data: \texttt{hecke\_ext/v5} (v1.14.0-pre.1; slot-13 WP-7, owner-approved 2026-09-26,
  slot-13 Q3; slot-12 FIN audit G2).} An ideal with \texttt{certificate\_routes.d1 = passed} carries, after
  \texttt{boundary\_reps}, the object \texttt{boundary\_data} $=\{$\texttt{"rank\_r1"}: $D$, \texttt{"r1"}:
  $r^1|_V$ ($h^1(\partial)\times D$), \texttt{"t"}: $\{$\texttt{"<p>"}: $t_p|_V$ ($D\times D$)$\}$, $p\in$
  \texttt{boundary\_reps} in class order$\}$ --- the block basis of the stored \texttt{t\_p}, matrices as arrays
  of rows, entries in $[0,\ell)$; every other ideal carries no such key. \texttt{verify()} code~32 re-derives
  (H1) and (H2) from it (Certificate~\ref{cert:boundary}). The string becomes \texttt{hecke\_ext/v5}; a
  \texttt{v4} record carries no \texttt{boundary\_data} and its (d1) is re-derivable in structure only. No
  verdict moves: \texttt{boundary\_data} records what the route decided on, never an input to
  \texttt{final}.
\end{enumerate}

\paragraph{The degree fields are per complex (v1.4.0; [C] \HA{sec:recordschema} item~5,
slot-2 owner confirmation~1).} \texttt{q\_max} (the top degree of \emph{that} complex),
\texttt{truncated}, \texttt{q\_max\_source} $\in\{$\texttt{"derived"}, \texttt{"set"}$\}$ (set
only where \texttt{truncated} is true, E9) and \texttt{degrees} $=[0,\dots,q_{\max}]$ are
carried by every complex a record is about, \textbf{never as a maximum over a record}: the
per-degree \texttt{hecke\_report/v2}, about the bulk, at record level (bulk $2$ /
\texttt{false} / \texttt{"derived"} / $[0,1,2]$ on the model); the extension record inside each
complex's object (boundary $1$ / \texttt{false} / \texttt{"derived"} / $[0,1]$; the pair as the
bulk). The campaign's job-level \texttt{degrees} is the \emph{reported selection}, a different
field that keeps its meaning and is part of \texttt{Job::id()}. The writer derives the fields
from the complex's own degree count and refuses a record that cannot name its complex.

\noindent The example's values are those of \S\ref{sec:tests} at $(10,691)$, \emph{corrected}
(v1.4.0): through v1.3.3 this paragraph read ``at a Kummer prime $r^1=0$ on the Eisenstein
class, so $H^1_!=H^1$ (\texttt{"all"})'', which the computation refutes --- $r^1$ vanishes on
the socle $E$ of the block but $\rank r^1=1$, and $H^1_!=\ker r^1$ has codimension~$1$
($\dim H^1_!=2$, hence $\dim X_3=3$ and $X_6=0$ by \HA{eq:dimbook}; \S\ref{sec:tests}); the
Eisenstein line is $\soc(H^1_+)=\m H^1_+$, the bottom Loewy layer, \emph{per sign}, and
undefined on the unsplit $H^1$ (\ref{D6}); and $T_{691}$ \emph{is} active ($c_{691}=542\ne0$).
The fields \texttt{local}, \texttt{sturm}, \texttt{coefficient\_field} (the $P_j$: the
character in $\Fl[X]/(\pi_w)$, each $P_j$ the reduced representative $X\bmod\pi_w$ of degree
$<r$ --- a constant at $r=1$, never the unreduced $X$ --- is not reproducible from the $\pi_q$
alone), the structure-map \texttt{matrices} (rows: bottom-layer basis, columns: top-layer
basis), \texttt{settings} with its fingerprint, \texttt{coefficient\_ring}, the per-complex
degree fields, \texttt{commit}, \texttt{model} and \texttt{iota\_sign} are what make a record
reproducible by a third party.
\texttt{iota\_sign} \emph{names the normalisation} of the involution relative to the geometric
one --- \texttt{"eisenstein\_in\_H1\_plus"} means $\iota^{\ast}=-\iota^{\ast}_{\mathrm{geo}}$,
[C] v1.2.0 \HA{def:iota}(iv), owner decision OQ3 of 2026-09-16 (\ref{add:iota}) --- and is
true by definition; a third party reproduces every $\pm$ field of the record from that one
formula. It is not a free parameter and is never branched on.

\subsection{The coefficient ring and the E8 objects}\label{sec:ring}
This subsection is the specification slot-8 WP-3 and WP-4 build against: the key WP-3 registers,
the rule by which everything that does not honour it refuses, and the objects a $\Z/\ell^k$ run
computes, with their golden names. It is \textbf{spec first}, written before any of it is coded
(slot-8 plan WP-2), under the owner's rulings of 2026-09-24 on the slot-8 decisions~2 (E8-P1
\ldots~P5 adopted), 3 (item~4 split; 4a in slot~8), 4 (refusal first, over the closure) and 5
(the seven E8 points), each at recommendation~(A); nothing here is provisional.
\emph{Division of labour}: this document owns the key --- its grammar, default, canonical text
and refusal rule; [C] owns the mathematics of integral coefficients (\HA{prop:coneintegral},
\HA{cl:bockstein}, \HA{cl:amalgam}, \HA{lem:h2torsion}) and points here for the key;
\texttt{src/ModArithmetic/ENGINE\_API.md} \S13 owns the arithmetic and the elimination
(\texttt{ModCtx::prime\_power}, the Howell form, \texttt{kernel\_module}, \texttt{image\_module},
\texttt{solve\_multi\_rhs\_module}, \texttt{length}, \texttt{invariant\_factors}); where [X]
v1.11's E8 item-3 text and \S13 differ, \S13 wins (slot-8 plan \S2.3 item~11). [L] is the E8
investigation, \texttt{docs/implementation/e8\_ladic\_investigation.md}, whose proposals E8-P1
\ldots~P5 the owner adopted. [IB] is the integral-line design brief,
\texttt{docs/implementation/integral\_line\_design\_brief.md}, whose open decisions OD-$n$ this
document cites (v1.9.1: the abbreviation was used from v1.9.0-pre.1 on without being defined). \emph{No record carries any of this in slot~8}: the objects of~(d) are
computed by \texttt{HeckeReductionValidation} and by the golden harness's E8 points only, and are
console-unreachable (rule R19 refuses \texttt{golden()} at a non-zero fingerprint, and the
pipeline verbs refuse by~(b)). Campaign rows and record fields for $\Z/\ell^k$ come with C5, in
slot~9.

\paragraph{Pre-release v1.10.0-pre.1 (slot-9 WP-2, spec first).} Slot~9 lifts four refusals of
this subsection and adds one schema, each specified here before it is coded (slot-9 plan WP-2):
\texttt{EquivariantCohomologyOnK} at $\ell\in\{2,3\}$ (slot-9 WP-3; (b3), (c1$'$)); the boundary
and the cone over $\Z/\ell^k$, E8 item~4b (slot-9 WP-4; (f)); the campaign's job loop with the
job-level schema \texttt{hecke\_ring\_report/v1} (slot-9 WP-8, C5c; (g)); and the console's
ring-aware verbs (slot-9 WP-12, C5a; (h)). The lift is listed member by member in~(b1). Every
text below marked ``v1.10.0-pre.1'' is binding from this pre-release; slot-9 WP-15 publishes
v1.10.0, recording what was built at the same sites. Written under the owner's rulings of
2026-09-24 on the slot-9 decisions 2, 3 (with~(iii)), 4 and~6, each at recommendation~(A);
nothing is provisional. \emph{Record carriage in slot~9}: the objects of~(d) are carried by
\texttt{hecke\_ring\_report/v1} (the campaign, under a non-default ring only) and by the console
handles of~(h); the objects of~(f) by the golden harness's E8 points (slot-9 WP-4's event R1,
\emph{a/s}) and by the validations only --- no record and no console verb carries them in slot~9.

\paragraph{Status at v1.10.0-pre.2 (slot-9 WP-15, phase~1).} What the pre-release specified is
\textbf{built as specified}, with the as-built facts added at their sites, each marked
``v1.10.0-pre.2''. \textbf{$\ell\in\{2,3\}$} (decision~4): slot-9 WP-3 (\texttt{97d0a8d}, merged
at \texttt{b53836f}), (b1), (b3), (c1$'$); \texttt{HeckeBredonTreeRingValidation}
(\texttt{standard}, $1394$ checks). \textbf{E8 item~4b} (decision~3 with (iii)): slot-9 WP-4
(\texttt{d0aca96}, merged at \texttt{8c9faf4}; review repairs \texttt{2933b65}, \texttt{7ddb555}),
(f); the event R1 \textbf{landed} at \texttt{762dbc3}: the $126$ entries of~(f) at the seven
points, predicate-gated ($126$ added, $91$ tier~A and $35$ tier~B, $0$ changed, $0$ removed; tier~A
$11\,575\to11\,666$, tier~B $5866\to5901$, entries $17\,441\to17\,567$), both determinism digests
unchanged by construction and checked before the dump
(\texttt{docs/matrix-engine/rebaseline-2026-09-24-e8-4b-addonly.md}); measured by
\texttt{HeckeReductionValidation} ($69\,122$ checks at the base, the 4b arms among them),
\texttt{HeckeReductionCoverageSlowValidation} and \texttt{GOLDEN\_STRICT=1
GoldenRegressionValidation}. \textbf{The console} (decision~2): C5a, slot-9 WP-12
(\texttt{b3cd0d3}, merged at \texttt{4bf2289}), (a) and (h); C5b, slot-9 WP-13 (\texttt{80eb184},
merged at \texttt{53066cd}), which binds the extension layer and refuses a non-default ring
throughout; their follow-up \texttt{f28a0fa}. The campaign (C5c, slot-9 WP-8) and the sites
slot-9 WP-6, WP-7 and WP-9 touch were still being built; the next paragraph records them.

\paragraph{Status at v1.10.0 (slot-9 WP-15, phase~2).} Everything the pre-releases specified is
\textbf{built as specified}, with four recorded deviations of~(g) and one amendment of~(g11); the
as-built facts of this phase are marked ``v1.10.0'' at their sites. Base \texttt{9b004f2}, every
slot-9 code package merged. \textbf{The campaign} (decision~6): slot-9 WP-8, C5c
(\texttt{ad65bd9}, with the (g11) amendment \texttt{f4fc3d2}, merged at \texttt{8853544}); (g)
as built, with (b1)'s and (b4)'s campaign rows; \texttt{CampaignRingValidation}
(\texttt{standard}, $693$ checks); zero movement at the defaults, measured against the base
binary. The merge of WP-8 with WP-9 produced one defect, repaired at \texttt{158c0be} as the lead
ruled: at $\ell\in\{2,3\}$ the ring job had named \texttt{SmallEllPolicy::Bredon}, which WP-9's
new scope guard forbids outside \texttt{EquivariantCohomologyOnK}; the job loop now builds its
relation list from the raw \texttt{ModCtx::prime\_power} context (the same value), and the Bredon
lift stays with \texttt{EquivariantCohomologyOnK} ((b1), the scoping row). \textbf{The carries of
slot-9 WP-9} (\texttt{fb8a63a}, merged at \texttt{e078600}), three of them closing open review
findings of this subsection: review finding m-1 by \texttt{SmallEllPolicyScopeValidation}
((b1)); m-2 by \texttt{HeckeBredonU64FlatRouteValidation} ((c1$'$)); the named guard of
\texttt{RunPipelineModEll} ((b1), ``refuses structurally''); and
\texttt{EquivariantCohomologyOnK::get\_prime()} refusing a composite object, with its one composite
caller re-pointed to \texttt{ring\_context()} ((b1), the constant-fibre constructor).
\texttt{IsotropyModelDeltaSquaredValidation} (review finding m-4) belongs to line~2, not to this
subsection. \textbf{The item-4b checks trimmed} (slot-9 WP-4t, \texttt{45acc4f}, merged at
\texttt{7e379bd}; owner-approved, test code only): the $(60,5)$ cell of the 4b arms moved to
\texttt{HeckeReductionCoverageSlowValidation}, and rows implied by kept rows were dropped
((b1), the 4b row); \texttt{HeckeReductionValidation} is $68\,970$ checks at this base (the
$69\,122$ above was the pre-release's base), its 4b arms $306$. \textbf{Slot-9 WP-6 and WP-7}
(the elimination driver under the caller's pass options and the effective-policy cache key; the
arena first-fit in $O(\log L)$ and the narrower $\tau$ lock) are engine and scheduling changes that
move no value, record, golden line or digest; no site of this document names them, and [P] v3.7
records the two that reach the console. \textbf{E8's declaration, the next paragraph, is unchanged.} Measured at
this base: \texttt{GOLDEN\_STRICT=1 GoldenRegressionValidation} tier~A $11\,666$ failing~$0$,
tier~B $5901$ mismatching~$0$ (re-run at this base by this revision), and the three
determinism digests unchanged (slot-9 WP-7's gates on this tree).

\textbf{Phase E8 is complete as specified} (slot-9 decision~3~(iii)~(A); [X] E8): \textbf{items
1--4b as specified; 4b over $\Z/\ell^k$ by route~1; \HA{obl:coneexact}(vii)'s composite half
open} (the two cross-routes and \texttt{EquivariantCohomologyOnK}'s relative sub-complex
constructor refuse a composite ring, with arms). Two further limits are recorded, not open items
of E8's specification: \HA{obl:coneexact}(vi) over the ring asserts the commutator and does not
compute the $\nu=2$ recursion ((f)); and at $\ell\in\{2,3\}$ only the lengths and invariant
factors of the Bredon object exist ((b3), (c1$'$)). [C] v1.9.0 states the same status
(\HA{sec:notimplemented}). (v1.11.0-pre.2: \textbf{superseded} by the amended declaration of the
``Status at v1.11.0-pre.2'' paragraph; this paragraph is v1.10.0's text, kept.)

\paragraph{Pre-release v1.11.0-pre.1 (slot-10 WP-2, spec first).} Slot~10 lifts three refusals of
this subsection, amends the console's list, and declines one lift by name, each specified here
before it is coded (slot-10 plan WP-2), under the owner's rulings of 2026-09-25 on the slot-10
decisions 3, 4, 6 and~8, each at recommendation~(A); nothing is provisional. Every text below
marked ``v1.11.0-pre.1'' is binding from this pre-release; slot-10 WP-15 publishes v1.11.0,
recording what was built at the same sites. The sites, and the package each binds:
\begin{itemize}
\item \textbf{\HA{obl:coneexact}(vii) over $\Z/\ell^k$} (decision~3~(A); slot-10 WP-3): (b1)'s
  ``the lift (slot~10)'', the three members that move; (f)'s (vii) bullet, the route over the ring,
  the regime rule and the acceptance statement. \emph{The declaration on acceptance}, which slot-10
  WP-15 writes into v1.11.0 in place of the paragraph above, and not before: ``\textbf{Phase E8 is
  complete}: items 1--4b as specified; \HA{obl:coneexact}(vii) over the ring \textbf{discharged at
  every E8 point} ($\ell^k<2^{32}$): route~1 in every regime, the two \texttt{uint32\_t}
  cross-routes below $2^{32}$, which follow the \texttt{uint32\_t} regime rule above it.'' If
  slot-10 WP-3 stops at the strict-pair route (its recorded risk; slot-10 plan WP-3), the
  declaration names route~2 only and (vii)'s route-3 half stays open, as the owner rules.
\item \textbf{The Bredon object at $\ell\in\{2,3\}$} (decision~4~(A)): (b3)'s ``the Hecke action on
  the Bredon object is declined'', and one sentence in~(c1$'$). No code package lifts anything
  there; the refusals of \texttt{HeckeOperatorOnK} and \texttt{HeckeReduction} at $2$ and $3$ stand,
  with their arms, and their message keeps pointing to~(b3).
\item \textbf{Line~2} (slot-10 WP-10, WP-13): (b1)'s row ``outside the closure'', the one
  specifying sentence those packages need from this document --- that the free model's objects
  are not members of~(b)'s closure and carry no refusal, and how their read-only use of
  \texttt{HeckeOperatorOnK} relates to~(b). Their mathematics is [C]'s (\HA{sec:integral},
  \HA{prop:strictlift}, \HA{lem:bulkperiodic}) and their design [IB]'s; neither is restated here.
\item \textbf{The console} (decision~6~(ii)~(A); slot-10 WP-12): (h)'s ``v1.11.0-pre.1'' block, the
  five structural members and the ring-mismatch rule. The other two items of decision~6 --- the
  member-getter position and the campaign verbs' exit~$5$ --- are console behaviour, not ring
  semantics, and are specified in [P] v3.8-pre.1 (\S6, \S8), which this document points to.
\item \textbf{The survey} (decision~8; slot-10 WP-4): (g)'s one sentence, no schema change.
\end{itemize}
\emph{Record carriage in slot~10} is slot~9's, unchanged: the cross-routes of (vii) are validated
equal to route~1, whose objects are the emitted values ((f), the event R1), so no record, golden
line or console verb carries a cross-route object; line~2 emits nothing (\texttt{[IB]} OD-9~(a)).
\textbf{Two pointers into [C]}, recorded and not patched (slot-10 plan \S4.3; slot~9's M-4
pattern): \HA{obl:coneexact}(vii)'s ``\textbf{its composite half is open}'' and
\HA{sec:notimplemented}'s ``(vii)'s composite half is open'' are \textbf{stale from slot-10 WP-3's
merge} until [C]'s next edit (slot-10 WP-15, [C] v1.10.0, decision~12~(A)), which records the
discharge; in that window this subsection and the tree are right. [C] v1.10.0 also carries
decision~4's remark beside \HA{prop:strictlift}(4) (below, (b3)); until then (b3) is its only
statement. (v1.11.0-pre.2: \textbf{both pointers closed} --- [C] v1.10.0 restates
\HA{obl:coneexact}(vii) and \HA{sec:notimplemented} over the ring, the window running from
\texttt{c9684e9} to that version's commit, and carries the remark as \HA{rmk:bredonhecke}.)

\paragraph{Status at v1.11.0-pre.2 (slot-10 WP-15, phase~1).} What the pre-release specified is
\textbf{built as specified}, with one recorded deviation in~(f) (route~3's transport) and one
amendment of the declaration (below); the as-built facts are marked ``v1.11.0-pre.2'' at their
sites. Base \texttt{f0640d0}. \textbf{(vii) over the ring} (decision~3~(A)): slot-10 WP-3
(\texttt{a4abc39}, merged at \texttt{c9684e9}); (b1), (f); \texttt{HeckeConeCrossRouteRingValidation}
(\texttt{standard}, $386$ checks, new); three rows of \texttt{HeckeReductionValidation}'s arms~(X)
and~(X4b) changed in place, its count unchanged ($68\,970$), and
\texttt{HeckeReductionCoverageSlowValidation} unchanged ($3731$); the $\Fl$ cone byte-identical
(\texttt{cone\_output\_diff.sh}, $0$ hunks). \textbf{Line~2} (outside the closure): slot-10 WP-10
(\texttt{ee5bdd4}, merged at \texttt{cb1c3c4}; \texttt{HeckeStrictLiftValidation}, $1824$ checks),
WP-13 (\texttt{99c6284}, merged at \texttt{3e371a2}; \texttt{HeckePeriodicityClassValidation},
$2482$) and WP-11 (\texttt{94ea619}, merged at \texttt{8f82016}, an \emph{a/s} package;
\texttt{HeckeIntegralActionValidation}, $4601$); their mathematics is [C] v1.10.0's
(\HA{rmk:strictliftbuilt}, \HA{rmk:integralevidence}). \textbf{The Bredon object at
$\ell\in\{2,3\}$} (decision~4~(A)): nothing lifted, (b3); [C] v1.10.0 carries the reason as
\HA{rmk:bredonhecke}. \textbf{The console} (decision~6~(ii)): slot-10 WP-12 (\texttt{3c96d63},
merged at \texttt{6726075}); (h). \textbf{The survey} (decision~8): \emph{v1.11.0}: \textbf{not run in slot~10} --- slot-10 WP-4,
the $\ell$-adic survey under $\Z/\ell^2$, is optional and belongs to the owner's window W2, which had
not run it when this version was written; it is deferred, and (g)'s sentence says the same. Nothing
of~(g) moves by it. \emph{v1.11.1}: \textbf{run in W2} after the close window (owner's consent,
2026-09-25) and published as \texttt{docs/matrix-engine/results/ring\_survey\_l2/} (owner-signed):
$\Z/\ell^2$, $\ell\in\{2,3,5,7,11,13,691\}$, even $n\le64$, $p\in\{2,3\}$ --- $462$ jobs, $0$
failed ($132$ \texttt{ok\_partial}, every one at $\ell\in\{2,3\}$, the Bredon object with (b3)'s
refusals recorded). At $\ell\ge5$, $n>0$: $43$ cells with Dickson classes, of which exactly three
carry $\ell^2$-torsion, $(30,5)$, $(60,5)$, $(56,7)$ --- the loci $n=\ell m$ of
\HA{rmk:bocksteinexp}(2) inside the grid, and the probe's cells of its clause~(3); the
bookkeeping identity of (d) and the universal-coefficient shape $H^0=\{2\}^{h}\{1\}^{b}$,
$H^1=\{2\}^{2\dim M_{n+2}-1+h}\{1\}^{b}$, $H^2=0$ hold on all $160$ cells; the $K$-mode at $K=2$ is
uncertified exactly at the three $\ell^2$-torsion cells; the O4 rows of [IB] hold; every job object
at the five $\Z/\ell^2$ E8 points equals its golden entry ($164$ objects). A second optional
publication ran in the same window: slot-10 WP-5's \texttt{results/ext\_u691/} (the $45$
\texttt{u\_ell\_over\_budget} cells of \texttt{ext\_first} with $U_{691}$), an $\Fl$ campaign outside
this section. Nothing of~(g) moves by either.
\textbf{Lane~B} (\emph{v1.11.0}; base of this version \texttt{47e3ff9}, every slot-10 code
package merged). \emph{Slot-10 WP-7} (\texttt{2f423af}, merged at \texttt{4560624};
\texttt{docs/implementation/slot10/WP-7\_report.md}) profiled the cone and the sign split at
$(160,691)$ and shipped two levers, both scheduling only: \textbf{the $\iota$ lever} ---
\texttt{HeckeInvolution::compute}'s fifteen independent checks of \HA{obl:iota}(iii) and the
per-degree splits of the $\pm$ decomposition run as host tasks in the serial order, the first failure
in task order, so the aborting check and its message are the serial loop's (the sign split
$1.80\to0.79$\,s at $(160,691)$, indicative) --- and \textbf{L4}, route~2's block rows of
\texttt{HeckeOperatorOnK} on the consumer pool, with a Fenwick tree and an \texttt{Executor} per row,
every rank, pattern hash, \texttt{ShapeKey} and recorded program the serial loop's. \textbf{The
cone-certificate lever} (\texttt{ConeCertificate::verify}'s twelve span tests as host tasks) is
\textbf{recorded for slot~11, not shipped}. Evidence that no value moves: \texttt{cone\_output\_diff.sh}
against the WP-7 base, $0$ lines changed and $0$ new; \texttt{ext\_compare\_mini.sh}, $84/84$
\texttt{ext/} records byte-identical with \texttt{commit} redacted; \texttt{bench\_identity.sh}, $45$
cases, $0$ differ; the golden file and the four digests unchanged; $0$ ThreadSanitizer reports.
Which site of this subsection names them: (b1)'s honouring list (L4 is inside
\texttt{HeckeOperatorOnK}, an honouring member; the $\iota$ lever is inside \texttt{HeckeInvolution},
a refusing member whose guard and arm are untouched). \emph{Slot-10 WP-8} (\texttt{00eb1a3},
\texttt{aeb8853}, merged at \texttt{7551dcc}; \texttt{docs/implementation/slot10/WP-8\_report.md}):
\textbf{the ninth bit} of \texttt{pass\_options\_fingerprint()} (bit~$256$, set iff the effective
device policy is \texttt{gpu}, so \texttt{hybrid} and \texttt{gpu} no longer share a program-cache
slot; \texttt{ENGINE\_API.md} \S10.6; \texttt{EngineFingerprintDeviceBitValidation}, \texttt{fast},
$8$ checks) --- a cache key, never a value; and \textbf{validation carries}, none a change of
specification: \texttt{HeckeFreeModelValidation}'s (T6) refuses on the 7-vertex complex's own number
of degrees rather than on the caller's label, and arm~(R) no longer re-runs two registered names
in process ($803\to801$ checks); \texttt{HeckePeriodicityClassValidation}'s ball arm is a genuinely
oriented disc, so the fill reaches a second round ($2482\to2483$ checks, with the review's m-9);
\texttt{HeckeConeCrossRouteRingValidation} asserts route~3's non-vacuity (review m-12, (f) below);
and \texttt{slot10\_close\_window.sh}. No site of this subsection names WP-8 except (f)'s check
count. Slot-10 WP-6 (the thread pool's cache-line placement,
merged at \texttt{2e5f073}) and WP-9 (\texttt{ENGINE\_API.md} \S13.12, arbitrary $k$ designed and
not built, merged at \texttt{4664b40}) move no value and no site of this subsection names them; WP-9
\emph{proposes} for a later version of this document the grammar amendment $k\le63\to k\le126$ of
(a), with the cell refusal $\ell^k\ge2^{64}\to\ell^k\ge2^{127}$, which v1.11.0 does not adopt ---
and it confirms that (a) already places the word-size refusal at the cell, by~(b3), not in
\texttt{MathSettings::set()}, which enforces the grammar cap only (its finding F-3, a precision for
the plan, not an error here). Measured at this base by the packages' own gates (each on its own
tree, the merges verified by step~0): \texttt{GOLDEN\_STRICT=1 GoldenRegressionValidation} tier~A
$11\,666$ failing~$0$, tier~B $5901$ mismatching~$0$; the four determinism digests unchanged.

\textbf{The declaration, as amended} (v1.11.0-pre.2; it replaces the ``Phase E8 is complete as
specified'' paragraph below, which stays as v1.10.0's text, and the declaration quoted in
v1.11.0-pre.1's paragraph above, which the slot's independent review found to claim more than was
built --- \texttt{docs/implementation/slot10/WP-14\_review.md} finding M-1): ``\textbf{Phase E8 is
complete}: items 1--4b as specified; \HA{obl:coneexact}(vii) over the ring \textbf{discharged in its
first two sentences at every E8 point} ($\ell^k<2^{32}$) --- route~1 in every regime, the two
\texttt{uint32\_t} cross-routes below $2^{32}$, which follow the \texttt{uint32\_t} regime rule
above it; \textbf{its third sentence}, the Poincar\'e-duality adjointness gate of
\HA{prop:coneroutes}(3), \textbf{recorded as unavailable over the ring}, by the clause's own rule,
since the E7 pairing and the contragredient refuse a composite ring.'' The third sentence was never
an E8 item ((f)'s ``Not lifted''), so E8's completeness does not wait on it. Two limits stay
recorded, not open items of E8's specification: \HA{obl:coneexact}(vi) over the ring asserts the
commutator and does not compute the $\nu=2$ recursion ((f)); and at $\ell\in\{2,3\}$ only the
lengths and invariant factors of the Bredon object exist ((b3), (c1$'$)). [C] v1.10.0 states the
same status (\HA{obl:coneexact}(vii), \HA{sec:notimplemented}).

\paragraph{Pre-release v1.12.0-pre.1 (slot-11 WP-4, spec first).} \textbf{Slot~11 lifts no
refusal of this subsection}: everything it builds on the pipeline is $\Fl$-only. What follows is
specified here, before it is coded, so that (b1)'s ``no row may be added by a code package'' holds
(slot-11 plan \S3 ``WP-4''; owner rulings of 2026-09-25, every decision at recommendation,
decision~3 at~(A$'$)); nothing is provisional. Slot-11 WP-15 publishes v1.12.0 with what was built
at the same sites.
\begin{itemize}
\item \textbf{New members of~(b1)'s closure} --- each \emph{refuses, by a named guard with an arm,
  setting and input}, raising \texttt{hecke::NotImplementedInThisBuild} naming itself (console
  exit~$5$): \texttt{HeckeExt/HeckeDicksonTheta}, every entry point (slot-11 WP-5, WP-6;
  \S\ref{sec:thetaobj}~(h)); \texttt{HeckeExt/HeckeCensusSplit} (WP-7; \S\ref{sec:census}~(j));
  \texttt{HeckeOperatorOnK} constructed with \texttt{FibreRepresentation::QuotientL} (whatever the
  other enumerators do over the ring, \texttt{QuotientL} refuses: $\theta$ over $\Z/\ell^m$ is
  slot~12's); and \texttt{LocalSystemSL2Z}'s \texttt{QuotientL} relation list with the
  $\rho_{Q_n}$ fibre-matrix builder of \texttt{LocalSystemUtils} (a composite \texttt{ModCtx} or a
  non-default setting). \texttt{SymN} and \texttt{ContragredientN} keep their rows of~(b1)
  unchanged. At $\ell\in\{2,3\}$ the ring refusal (the default \texttt{SmallEllPolicy}) comes first;
  under the default ring the $\theta$ entry points refuse $\ell<5$ as a precondition
  (\S\ref{sec:thetaobj}~(h)). \emph{The arms} live in the packages' own validations
  (\texttt{HeckeDicksonThetaValidation}, \texttt{HeckeCensusSplitValidation}), not in arm~(X) of
  \texttt{HeckeReductionValidation}, whose file no slot-11 package edits (slot-11 plan \S3.0,
  shared-file rule); the as-built list is WP-15's to copy into v1.12.0, and a member a package
  finds at its base that this list does not name keeps refusing and is reported.
\item \textbf{The console} (slot-11 WP-16, \emph{a/s}): \texttt{dickson\_les} and \texttt{theta}
  refuse a non-default ring at dispatch and again in the lazy compute, and a ring handle through
  \texttt{bind::RequireFieldHandle} --- exit~$5$, a surface of~(b4) like every C5b verb ([P]
  v3.9-pre.1 \S3.6, \S8).
\item \textbf{Outside the closure}, as (b1)'s row for line~2 says: slot-11 WP-11's
  \texttt{Integral/Model/StrictLift} for $\mathrm{diag}(1,p^2)$ and $\mathrm{diag}(p,p)$ (explicit
  \texttt{ModCtx}, no setting read); and slot-11 WP-18's $128$-bit element type (next item).
\item \textbf{L1-X under decision~3~(A$'$).} Slot~11 builds \texttt{ENGINE\_API.md} \S13.12's
  \emph{step~1 only} (slot-11 WP-18, after WP-12's F-2): the engine's $128$-bit composite element
  type (\texttt{ModCtx} over \texttt{u128}, the \texttt{U128\_Montgomery} regime, \texttt{u256}
  products, the primitives, Hensel inversion), \texttt{DenseBuf} and the scalar kernels, the
  \S13.6 entry points, \texttt{ref::howell} on naive arithmetic, the block model's composite
  surface --- \emph{and its certificate\,/\,oracle rule}: the \texttt{u128} route equal to the
  \texttt{u64} route entrywise on \texttt{HeckeReductionValidation}'s elimination matrices, and
  arithmetic arms at $\ell^k\ge2^{64}$ against \texttt{RefLinAlg}'s naive arithmetic
  (\texttt{EngineU128CompositeValidation}). It is \textbf{value-neutral and has no consumer}: no
  member of~(b1)'s closure, no setting, no campaign cell, no console verb and no golden point
  reaches it. Consequently \textbf{(a)'s grammar $k\le63$ and (b3)'s cell refusal at
  $\ell^k\ge2^{64}$ do not move in slot~11}, and every ring surface keeps refusing $\ell^k\ge2^{64}$
  exactly as at v1.11.1. The grammar amendment $k\le126$ with the cell refusal at
  $\ell^k\ge2^{127}$ (slot-10 WP-9's proposal, recorded in the v1.11.0 status above) is
  \textbf{slot~12's}, with L1-X's steps~2--4 (the pipeline route, the grammar, campaign and
  console, the validation and review), beside [C]'s Hida chapter, where the per-cell certified
  precision bound is written; a later pre-release of this document specifies it, not this one.
  \emph{F-1} (a Montgomery-64 \texttt{prep}\,/\,\texttt{mulp} pair in the composite row updates of
  the \texttt{u64} Howell\,/\,module route; canonical in and out) and \emph{F-2}
  (\texttt{ExactKthRoot}'s trial division replaced by the ring carried as $(\ell,k)$, or a
  deterministic Miller--Rabin; the blocker \S13.12 names in \texttt{Hecke/CoefficientRing.cpp})
  are value-preserving: no value, member, refusal kind or exit code of this subsection moves, and
  a changed refusal \emph{message} is reported by the package.
\item \textbf{Record carriage in slot~11}: none. $\theta$'s objects are carried by validations
  only (\S\ref{sec:thetaobj}~(g)); the census by its own artefact \texttt{hecke\_census/v1}
  (\S\ref{sec:census}), which reads $\Fl$ records only. No record, golden line or digest moves.
\end{itemize}

\paragraph{Pre-release v1.12.0-pre.2 (slot-11 WP-18c, spec first; the (C2) re-ruling).} The
owner's clarification of 2026-09-25 (latest) puts \emph{arbitrary $k$} into slot~11, by the (C2)
$\ell$-adic digit induction of \texttt{ENGINE\_API.md} \S13.11 (slot-11 WP-18b, built, value-neutral)
and a pipeline route on it (slot-11 WP-18c). Pre.1's sentence that (a)'s grammar $k\le63$ and (b3)'s
cell refusal at $\ell^k\ge2^{64}$ ``do not move in slot~11'' is \textbf{superseded} by the items
below; the rest of pre.1 stands. Every text marked ``v1.12.0-pre.2'' is binding from this
pre-release; slot-11 WP-15 publishes v1.12.0 with what was built. The [C] side is
\HA{obl:backend}(iv)'s digit-plane sentence ([C] v1.11.0-pre.2): an element of $\Z/\ell^{K}$ held
as $L=\lceil K/m\rceil$ fixed-width words, $L$ fixed per cell, is not the arbitrary-precision
arithmetic that clause excludes.
\begin{itemize}
\item \textbf{(a) The grammar} (v1.12.0-pre.2). $\texttt{Z/ell\textasciicircum}k$ with
  $2\le k\le4096$; everything else of~(a) is unchanged (spelling, reserved values, default,
  canonical text, registration; eleven keys stay eleven). The cap is a \emph{parse} cap, stated so
  that the refusal is reviewable; raising it is an owner ruling, not a code change. Whether a value
  is executable at a cell is decided at the cell, by the next item.
\item \textbf{The per-cell rule} (v1.12.0-pre.2; it refines (b3)'s word-size refusal). At a cell
  $(n,\ell)$ under $\texttt{Z/ell\textasciicircum}k$: (i)~$\ell^k<2^{64}$ --- the shipped native
  route, \emph{unchanged}, with every one of its refusals (in particular the console's
  \texttt{uint32\_t} limit); (ii)~$\ell^k\ge2^{64}$ and $\ell\ge5$ --- the \textbf{digit route}
  below; (iii)~$\ell^k\ge2^{64}$ and $\ell\in\{2,3\}$ --- \textbf{refused by name} (exit~$5$): the
  Bredon object is the flat-coordinate route, $C^q_{\mathrm{eq}}$ is not free there and the
  precision-split lemma (L1) of \S13.11.3 does not apply; line~2 serves those primes;
  (iv)~at a digit cell, a working precision with $K\cdot\lceil\log_2\ell\rceil>8192$ bits, or a
  predicted working set above the slot's memory cap ($40$\,GB, rule~19 of the slot-11 plan), is
  refused by name, exit~$5$, the message carrying the predicted figure.
\item \textbf{The digit route} (\texttt{Hecke/LadicCell}; one new consumer, instantiating no
  existing class template for a new element type). For a digit cell with $n>0$ even (the
  $K$-mode's applicability range, (d); elsewhere the route refuses by name) and requested output
  precision $k$ it computes, over $\Z/\ell^{K}$ in digit planes: $C^q_{\mathrm{eq}}$ (the kernel
  module of the relation system, asserted free), $X_q$, $Z^q$, the (c1)--(c5) cyclic decomposition
  of $H^q$ and the invariant factors of $d^q$, all by the \S13.11 entry points, which return the
  \S13.5--\S13.6 canonical forms; $\rho_n(g)$ from the exact integral matrices by the division-free
  Pascal route of \S13.10 F4 (the $\mathrm{diag}(\ell,1)$ split at non-units); $\Theta_g$ for $T_p$,
  $r=1$, $p\ne\ell$, and for $U_\ell$ ($p=\ell$), from the shipped integral coset and $\tau$ tables,
  \emph{both routes built and asserted equal}, with assertions (i)--(iii) of the ring route;
  $A^{(q)}(T_p)$ on the cyclic summands; and $\mathrm{charpoly}_{A(T_p)}^{1,\mathrm{free}}$ (the
  division-free characteristic polynomial of the free block, as in (d)). \emph{Honoured} at a digit
  cell: the lengths and invariant factors of $H^q(\Z/\ell^{k})$, the $K$-mode ($r_K$ at the
  working level, \texttt{ladic\_type\_certified}, $\texttt{ladic\_type}^q$), the free characteristic
  polynomial and the free block of $A^{(1)}(T_p)$ at precision $k$, and the ordinary idempotent
  (below). \emph{Refused by name at a digit cell}, exit~$5$ (members of (b1)'s closure that the
  route does not reach in slot~11): the Bockstein and the reduction maps of (d) other than
  truncation, the boundary and the cone ((f)), the pairing and the contragredient, the extension
  driver, $T_{p^r}$ with $r\ge2$, the operator family and the decompositions, and every stored
  basis as a console \texttt{Matrix}.
\item \textbf{The per-cell certified precision bound} (v1.12.0-pre.2; the owner's request). For a
  requested precision $k$: $e_{\max}$ is read from the \emph{certified $K$-mode at word size}
  (\HA{prop:certifiede}(3)): one shipped native run at
  $K_0=\max(2,\lfloor1+\log_\ell(n/(\ell+1))\rfloor+1)$ ([L] \S2.5's observed bound, used only as the
  starting level), escalated once to the largest word level $K_w=\max\{k':\ell^{k'}<2^{64}\}$ if $K_0$
  does not certify; if $K_w$ does not certify either, the cell is \textbf{refused, resumable}:
  ``precision insufficient'', with the R3 state $\{\texttt{kind}=\texttt{ladic\_precision},\
  \texttt{budget}=\{K_w\},\ \texttt{reached}=\{r_{K_w}(d^0),r_{K_w}(d^1)\},\
  \texttt{complete}=\texttt{false}\}$ --- never a guessed value. $e_{\max}$ is the largest torsion
  exponent of $\texttt{ladic\_type}^1$ (asserted: $\texttt{ladic\_type}^2$ has no torsion at
  $\ell\ge5$). The working precision is $K=k+e_{\max}$; the digit run \emph{re-evaluates} the
  stopping rule at $K$ (a second, independent certificate) and checks its invariant factors of
  $H^q(\Z/\ell^{K})$ against the universal-coefficient reading of the certified type. The lengths
  and invariant factors of $H^q(\Z/\ell^{k})$ are that reading at $k$. A digit record states
  \texttt{precision\_requested} $=k$, \texttt{precision\_working} $=K$, \texttt{e\_max},
  \texttt{word\_level} (the level that certified), \texttt{precision\_certified},
  \texttt{digit\_width} and \texttt{digits\_L}; the free characteristic polynomial and every Hecke
  datum on $\Lambda^1\otimes\Z/\ell^{k}$ are certified to precision $k$ exactly when
  \texttt{precision\_certified} is true.
\item \textbf{The ordinary idempotent} (v1.12.0-pre.2; a certified datum, not a choice of $m$).
  At $p=\ell$ the route computes, on the free block $U$ ($r\times r$) of $A^{(1)}(U_\ell)$ over
  $\Z/\ell^{k}$, $e=-p_0^{-1}W\,s(W)$ with $W=U^{rk}$ and $\chi_W(x)=x^{a}p(x)$, $p(x)=p_0+x\,s(x)$,
  $p_0$ a unit (Fitting: $W$ vanishes on the nilpotent summand and is invertible on the ordinary
  one, so $W p(W)=0$); $e$ is a polynomial in $U_\ell$, hence the ordinary projector
  $\lim U_\ell^{N!}$. It is certified by $e^2=e$, $eA(T_p)=A(T_p)e$ for the other primes of the
  cell, $(U(1-e))^{r}\equiv0\pmod\ell$, and $\rank_{\Fl}(e\bmod\ell)=\rank_{\Fl}(U\bmod\ell)^{r}$. The
  depth $m$ a Hida computation needs is supplied per question by [C]'s later Hida section; the
  route's bound $K=m+e_{\max}$ is its input, and no cap on $m$ beyond (a) and the per-cell rule is
  imposed.
\item \textbf{The width switch} (v1.12.0-pre.2). \texttt{ENGINE\_LADIC\_WIDTH} $\in\{32,64,128\}$
  (default $64$) selects the digit word; it is an \emph{engine-family} switch --- it cannot change a
  value (the width oracle matrix of \S13.11.7 and the route's cross-width arm prove it), enters no
  fingerprint, job id or record hash, and appears in a digit record only as \texttt{digit\_width}
  and \texttt{digits\_L}.
\item \textbf{(g) Digit jobs} (v1.12.0-pre.2). The campaign's ring job loop accepts digit cells: the
  spec's refusal of $\ell^k\ge2^{64}$ moves to the per-cell rule (a grid cell of case (iii) or (iv)
  is refused at spec validation, naming the cell). A digit job writes \texttt{hecke\_ring\_report/v1}
  with the members (g1)--(g16) as a native job does, and in addition \texttt{"route":"digit"} and the
  object \texttt{"precision"} holding the bound's keys; its \texttt{modulus} is the text
  \texttt{"<ell>\textasciicircum<k>"}; a residue is written as its signed decimal (the
  representative of least absolute value) when that fits a signed $64$-bit integer, else as the
  string \texttt{"ell-adic:[}$a_0$\texttt{,}$a_1$\texttt{,\ldots]"} of its base-$\ell$ digits, low
  to high; the Bockstein and the reduction are recorded refusals ((g14)); at $p=\ell$ the record
  carries \texttt{"hida"}. The precision keys appear on digit jobs only; no native job's record, id,
  file name or summary line moves.
\item \textbf{(h) Digit handles} (v1.12.0-pre.2). Under $\texttt{Z/ell\textasciicircum}k$ with
  $\ell^k\ge2^{64}$, $\ell\ge5$: \texttt{symmetric\_power}, \texttt{cohomology} and
  \texttt{hecke(H, p)} ($r=1$) return digit handles. On a digit \texttt{Cohomology} handle the
  members \texttt{computed}, \texttt{q\_max}, \texttt{truncated}, \texttt{q\_max\_source},
  \texttt{degrees}, \texttt{coefficient\_ring}, \texttt{lengths}, \texttt{invariant\_factors},
  \texttt{ladic\_type\_certified}, \texttt{r\_K} and \texttt{ladic\_type} honour, and the new member
  \texttt{precision} (a record of the bound's keys) is honoured on digit handles only; the rest of
  (h)'s ring members refuse, exit~$5$. \emph{v1.12.0} (slot-11 WP-14 pass~2, M-1, built): on a digit
  \texttt{Cohomology} handle the field members (\texttt{dims}, \texttt{C}, \texttt{Z}, \texttt{B},
  \texttt{H}, \texttt{P}) exit~$3$, as on every ring handle; every other member not on the digit list
  exits~$5$. On a digit \texttt{HeckeOperator} handle
  \texttt{coefficient\_ring}, \texttt{p}, \texttt{r} and \texttt{charpoly\_free} honour (its
  coefficients as Strings in (g)'s residue form); \texttt{A} and every other member refuse, exit~$5$.
  At $\ell\in\{2,3\}$ beyond $2^{64}$: the per-cell refusal, exit~$5$.
\end{itemize}

\paragraph{Published v1.12.0 (slot-11 WP-15): the two pre-release paragraphs as built.} Sources: the
reports of slot-11 WP-5, WP-6, WP-7, WP-8, WP-12, WP-18a, WP-18b and WP-18c and the slot's independent
review, passes~1 and~2 (\texttt{docs/implementation/slot11/}). Nothing specified above is withdrawn
except where this paragraph says so; where the build differs from the text, the difference is stated
here and is a carry, not a silent change.
\begin{itemize}
\item \textbf{The $\theta$ and census members of~(b1)'s closure} (pre.1's first item) \emph{refuse}, as
  specified, but not every one by a guard naming itself: \texttt{HeckeExt/HeckeDicksonTheta}
  (\texttt{DicksonThetaCell}) refuses a non-default ring by its own named guard; \texttt{HeckeExt/HeckeCensusSplit}
  refuses through \texttt{HeckeExtDriver}'s guard (the message names the driver);
  \texttt{HeckeOperatorOnK} constructed with \texttt{QuotientL} over a ring refuses through its
  non-\texttt{SymN} guard, whose message says ``contragredient''; the \texttt{QuotientL} relation list and
  fibre builder carry no ring guard of their own and are reached only through those two refusing entry
  points. No validation arm drives a ring into \texttt{DicksonThetaCell} or the split driver (the
  $\ell\in\{2,3\}$ arm of \texttt{HeckeDicksonThetaValidation} exists). \emph{Carry to slot~12}: a guard
  naming itself at each of the four sites and one arm each, with $\theta$ over $\Z/\ell^m$.
\item \textbf{The console} (pre.1's second item): \textbf{not built}. Slot-11 WP-16 gave way under the
  owner's re-ruling of decision~3 (slot-11 plan \S10); \texttt{dickson\_les} and \texttt{theta} do not
  exist, and [P] v3.9 records them as specified, not built.
\item \textbf{L1-X} (pre.1's fourth item): superseded by pre.2 and by the owner's later rulings (slot-11
  plan \S10: arbitrary $k$ by the (C2) digit induction in slot~11). The native \texttt{u128} layer of
  \texttt{ENGINE\_API.md} \S13.12 is built (slot-11 WP-18a) as the \emph{third digit width and an
  oracle}, not as a pipeline route: \texttt{EngineU128CompositeValidation} ($26\,587$ checks). F-1 and F-2
  are built and value-preserving (slot-11 WP-12; \texttt{ModCtxMontgomeryPairValidation},
  \texttt{CoefficientRingKthRootValidation}); no value, member, refusal kind or exit code moved. One
  pre-existing undefined shift in \texttt{ExactKthRoot} at $k<2$ (reached by no pipeline caller) was made
  defined with the observed result.
\item \textbf{The (C2) engine} (\texttt{ENGINE\_API.md} \S13.11, slot-11 WP-18b):
  \texttt{engine::ladic}, header-only, with its width oracle matrix, \texttt{EngineLadicInductionValidation}
  ($1888$ checks). Domain $2\le\ell<2^{31}$, $K\ge2$, no engine cap on $K$.
\item \textbf{The digit route} (pre.2, slot-11 WP-18c, \texttt{Hecke/LadicCell}), as specified, with:
  (i)~the caps of the per-cell rule --- $8192$ bits and $40$\,GB --- checked at the \emph{largest reachable}
  working precision $K=k+K_w-1$ (since $e_{\max}<K_w$) at spec validation and at the console's dispatch,
  so no job is refused inside its run (slot-11 WP-14 pass~2, m-3; at $\ell=691$: $\texttt{Z/ell\textasciicircum}814$
  admitted, $815$ refused); (ii)~$\ell\ge2^{31}$ with $\ell^k\ge2^{64}$ refused by name, exit~$5$ (the
  engine's domain); (iii)~$n=0$ or $n$ odd refused by name, exit~$5$ (the $K$-mode's range);
  (iv)~$\texttt{ladic\_type}^2$ asserted to have free rank~$0$ and no torsion, and every $X_q$ residual-checked.
  \texttt{HeckeLadicCellValidation} ($891$ checks) carries the oracles of \texttt{ENGINE\_API.md}
  \S13.11.11, the M-0 charpolys at $(10,691)$ beyond $2^{127}$, the bound at $(20,5)$, $(30,5)$, $(60,5)$
  ($e_{\max}=1,2,2$), a cross-$K$ and cross-route comparison of $\chi(T_2)$ at those three cells, and the
  forced-uncertified refusal.
\emph{v1.18.0} (slot-18 WP-19b, \texttt{hecke::ladic\_cell::DigitCellSession}): the route is cut at the Hecke step --- steps~1--2 (the cell over $\Z/\ell^{K}$, prime-independent) run once per cell, the Hecke step per prime $p\ne\ell$ on the shared cell; \texttt{RunCell} \emph{is} the session run on the request's primes (one code path), and the console keeps a bounded session cache keyed by the cell. Value-free: \texttt{HeckeLadicSplitValidation} ($56$ checks) compares the shared-cell records with fresh \texttt{RunCell} records bit for bit.
\item \textbf{The ordinary idempotent} (pre.2): as specified. Certified at $(10,691)$ with $U_{691}$
  (fully ordinary on the free block, $e=I$) and at $(30,5)$ with $U_5$ ($r=5$, $\chi_W=x^{4}p(x)$,
  $\rank_{\F_5}e=1$: a genuine Fitting split). A digit job's record carries
  \texttt{"hida":\{"precision", "rank", "rank\_e\_mod\_ell", "ordinary\_rank\_mod\_ell", "certified", "e"\}},
  \texttt{"e"} being the idempotent's $r\times r$ entries, row-major, as residue Strings in~(g)'s form
  (\texttt{Campaign/CampaignRing.cpp}). \emph{v1.12.1} (the slot-11 WP-15 referee's M-2): v1.12.0 listed
  five keys and omitted \texttt{"e"}; this paragraph is the record's only inventory of the object, and
  (b) makes that inventory binding for a non-default-ring record, so the sixth key is named here.
\item \textbf{(g) as built}: as specified; a digit job books the Hecke time into \texttt{t\_cohomology}
  (a timing field, no value). \textbf{(h) as built}: as specified, with the field members exiting~$3$
  (the v1.12.0 sentence in~(h) above); a digit \texttt{HeckeOperator}'s \texttt{charpoly\_free} is the Record
  \texttt{\{precision, coeffs\}} with the coefficients as Strings in~(g)'s residue form; and
  \texttt{hecke(H, p)} on a digit handle recomputes the whole cell once per operator --- deterministic,
  so no value moves; a cost, recorded here (slot-11 WP-14 pass~2, m-2).
\item \textbf{Owner item, recorded, not decided: the console's \texttt{uint32\_t} gap.} By rule~(i) the
  console keeps the native route's \texttt{uint32\_t} limit below $2^{64}$, so it refuses
  $\texttt{Z/ell\textasciicircum}k$ at $\ell=5$ for $14\le k\le27$ and at $\ell=691$ for $4\le k\le6$ while
  it serves $k\ge28$ and $k\ge7$ there by the digit route; the campaign serves every $k$. Closing it
  (the console's \texttt{uint64\_t} instantiation, or the digit route below $2^{64}$ at the console) is the
  owner's to rule (slot-11 WP-14 pass~2, m-4; [P] v3.9).
\end{itemize}

\paragraph{The Hecke algebra over $\Z/\ell^{K}$, read mod~$\ell$ (v1.15.1-pre.2; slot-14 WP-6; certified verdict and digit route v1.15.1-pre.5, slot-15 WP-5).}
\emph{Definitional; computed objects, no certificate beyond the checks named and the certified verdict of (ii), which is [C]'s statement, not
one of this subsection.}
\begin{enumerate}
\item[(i)] \textbf{Objects.} For $\ell\ge5$, $n\ge1$ even, $K\ge2$ (since v1.15.1-pre.5 any $K$: $\ell^{K}\ge2^{64}$ by the digit route of (v)), let
  $M_K=H^1(\Ga;\Sym^n\Z/\ell^{K})$ by the module route of this subsection and, per $\Fl$-rational joint system $\m$ of $(T_p)_{p\in S}$, $S=\{2,3,5,7,\ell\}$: the joint idempotent
  $e_\m\in\Z/\ell^{K}[T_p]$ (per-$p$ idempotent polynomials, Newton lift $e\leftarrow3e^2-2e^3$; $e^2=e$ and $eT_p=T_pe$ are checked), the block
  $B_K=e_\m M_K\cong\bigoplus_s\Z/\ell^{b_s}$ by \texttt{cyclic\_decomposition}, $A_K$ the subalgebra of $\End(B_K)$ generated by the $T_p$ (Howell
  forms), its length, the Hilbert function of $A_K\otimes\Fl$ at $\m$, $A_1$ the same on $B_K/\ell$ (the census algebra $A$ of \S\ref{sec:census}), and
  $\varepsilon_K:=K\cdot\dim_{\Fl}(A_K/\ell)-\mathrm{length}(A_K)$ ($=\varepsilon^{\mathrm{sat}}_K$ of \HA{prop:satdefect} at the certified step, where
  $\dim=\lceil f/2\rceil$; v1.15.1-pre.6, Fable audit M2 --- the text before read $K\cdot\mathrm{rank}_{\Z_\ell}(B_K)$, which the code never computed). At $K_{\mathrm{cls}}=\max(2,e_{\max}+1)$ ($e_{\max}$ from \HA{prop:torsionprepass})
  the summands of order $<\ell^{K_{\mathrm{cls}}}$ are the torsion of $\Lambda_\m$ and the rest its free part of rank $f$; on a torsion block the free part is read on
  $\ell^{t}$-scaled generators, $t$ the largest torsion exponent.
\item[(ii)] \textbf{The \texttt{stable} verdict is certified} (v1.15.1-pre.5; owner ruling Q4, 2026-09-30). Until v1.15.1-pre.4 it rested on
  \texttt{docs/implementation/slot14/WP-M1\_loewy\_note.md}, Proposition~1, and \texttt{stable}, \texttt{torsion\_free\_stable} and \texttt{torsion\_stable}
  were \emph{readings, not certificates}. [C] v1.15.0 proves the stopping rule --- \HA{prop:saturation}(3) (P15), \HA{prop:satdefect} (P17),
  \HA{lem:heckerankone}, \HA{lem:freepart} --- and the route's verdict is \HA{cor:stablecert}, in the form of the WP-4 referee's repair B1: the block
  $e_\m M_K$ is the block of a maximal ideal $\m'$ of $\mathbb{T}'=\Z_\ell[T_p:p\in S]$ (the joint $\Fl$-rational system of $S$), which may be a \emph{sum} of
  blocks of the full algebra $\mathbb{T}^{\mathbb L}$; $\mathrm{rank}\,e\mathbb{T}^{\mathbb L}=\lceil f/2\rceil$ for every idempotent $e$ ($f=\mathrm{rank}(e\mathbb L')$,
  at level~$1$ at most one Eisenstein system in any block), so $\dim\mathsf A'_K/\ell\le\mathrm{rank}\,\mathbb{T}'_{\m'}\le\lceil f/2\rceil$, and
  equality certifies $\mathsf A'_K/\ell\cong\mathbb{T}'_{\m'}/\ell$ --- Hilbert function and Loewy length included --- at this and every larger $K$.
  The Hilbert function \emph{printed} is the route's reading of $\mathsf A'_K/\ell$ (its sum and $h_0=1$ are checked); the certificate is of the
  isomorphism, which carries it (v1.15.1-pre.6, audit M3).
  \emph{What it certifies:} $\mathbb{T}'_{\m'}/\ell$, the algebra generated by the $T_p$, $p\in S$; \emph{not} $\mathbb{T}^{\mathbb L}_\m/\ell$, which needs the
  separate generation certificate of \HA{cor:stablecert}(2) that this route does not run. The full algebra of a torsion block with $f>0$ stays
  uncertified (\HA{cor:stablecert}(3)); its free part (\texttt{torsion\_free\_stable}, \HA{lem:freepart}(2)) and a pure-torsion block
  (\texttt{torsion\_stable}, \HA{lem:freepart}(3)) are certified, as is \texttt{reduced\_rank1} (the case $K=1$ of \HA{cor:stablecert}(1)).
  \emph{The certificate} (\texttt{CertifyStable}) re-derives, in exact $\Z/\ell^{K}$ arithmetic and from the module runs themselves, (1) the block:
  $e_\m$ a polynomial in the $T_p$ (so $e_\m\in\mathbb{T}'$), $e_\m^2=e_\m\ne0$, the $T_p$ commuting on the block, $T_p-a_p$ nilpotent on it mod~$\ell$
  --- which make $e_\m$ the idempotent of $\m'$ --- and, as a consistency check only (v1.15.1-pre.6, audit M1), $\mathrm{rank}_{\Fl}(e_\m\bmod\ell)$ the joint
  generalised eigenspace's dimension; the primes of $S$ are checked prime (audit M4); (2) $f$ and the torsion exponents from the \emph{Smith} invariant
  factors of $e_\m M_{K_{\mathrm{cls}}}$ (a second entry point beside the cyclic decomposition), and on the whole module
  $M_{K_{\mathrm{cls}}}$ the summands of order $\ell^{K_{\mathrm{cls}}}$ number $\mathrm{rank}_{\Z_\ell}H^1(\Ga;\Sym^n\Z_\ell)=2d_k-1$ while the others are
  exactly the pre-pass torsion of \HA{prop:torsionprepass} (so $K_{\mathrm{cls}}>e_{\max}$ and no torsion summand is read as free); (3) $\lceil f/2\rceil$; (4) $\mathsf A'_K$ by a second
  construction, the $\Z$-span of the monomials $\prod_pT_p^{i_p}$, $0\le i_p<D$ (Cayley--Hamilton on the $D$-generated block), with its length and
  $\dim\mathsf A'_K/\ell$ (on $\ell^{t}$-scaled generators for a free part); (5) the criterion $\dim\mathsf A'_K/\ell=\lceil f/2\rceil$ ($f=0$ and
  $K\ge t$ for \texttt{torsion\_stable}); (6) $\varepsilon_K=\mathrm{level}\cdot\lceil f/2\rceil-\mathrm{length}(\mathsf A'_K)\ge0$ and, along the recorded
  levels, $\varepsilon^{\mathrm{sat}}$ non-decreasing (\HA{prop:satdefect}(2)). Each is compared with the route's reading; \emph{any} mismatch throws
  (\texttt{ENGINE\_CHECK}) --- the verdict is never silently downgraded. Every \texttt{LOEWY\_RING} line gains \emph{one} appended field,
  \texttt{stable\_cert=certified|reading|none} (\texttt{reading}: the certificate not run, \texttt{LoewyOptions::certify = false}; \texttt{none}: no stable
  verdict); the old text, \texttt{verdict=stable} included, is byte-identical.
\item[(iii)] \textbf{Refusals.} $\ell<5$, $n=0$ or odd, $K<2$ by \texttt{ENGINE\_CHECK}; beyond $2^{64}$ the digit route's per-cell rule
  (\S\ref{sec:ring}, v1.12.0-pre.2: $K\lceil\log_2\ell\rceil\le8192$, the predicted working set, $\ell<2^{31}$) by \texttt{hecke::NotImplementedInThisBuild};
  the Hida module runner (an oracle) beyond $2^{64}$ by name; a non-local algebra by name.
\item[(iv)] \textbf{Surface.} \texttt{HeckeRingLoewyValidation} (\texttt{standard}, $181$ checks since v1.15.1-pre.6 ($169$ at pre.5), $123$ before; the Python drivers of
  \texttt{wpm1/} as oracle at $(58,7)$; arms (C1) the certified verdict on every stable verdict met, (C2) negative controls --- a perturbed claim (a wrong $n$, $K_{\mathrm{cls}}\le e_{\max}$ included) and a
  perturbed $\mathsf A_K$ are refused, (C3) the digit route);
  \texttt{HECKE\_RING\_LOEWY\_CELLS="n:ell[:KMAX],\dots"} (read-only), one \texttt{LOEWY\_RING} line per block, \texttt{LOEWY\_RING\_SUMMARY} (since
  v1.15.1-pre.5 with an appended \texttt{stable\_cert} tally), C-F9 failures \emph{reported} as \texttt{LOEWY\_RING\_CF9\_FINDING}, never gated. No record
  key, no golden object, no setting; \texttt{coefficient\_ring} is not read (the route takes $K$ itself). Evidence: the counts quoted at
  Prediction~\ref{pr:latticesplit}.
\item[(v)] \textbf{The digit route} (v1.15.1-pre.5; the carry of [X] v1.19 (6)). At $\ell^{K}\ge2^{64}$ the module run is the (C2) digit route of this
  subsection (\texttt{hecke::ladic\_cell::RunCell} at exactly the working precision $K$, its level driver and canonicalisation, its UCT check against the
  pre-pass type \HA{prop:torsionprepass}), and the algebra reading runs on \texttt{engine::ladic}'s ring, Howell form, length, invariant factors and
  cyclic decomposition --- the \S13.5--\S13.6 objects in digit arithmetic, no new elimination. Every reading is basis-free, so the two routes give the
  same \texttt{LOEWY\_RING} lines wherever both apply ((C3): $(58,7)$, $(30,5)$, $(34,5)$ with the digit route forced below $2^{64}$); beyond $2^{64}$,
  at $\ell=5$, $K=28$ ($5^{28}>2^{64}$, the NP5R family), the digit module's readings at $(30,5)$ and $(34,5)$ are those predicted from the
  \texttt{uint64\_t} route at $K=26,27$ by \HA{prop:saturation}(3)--(5) and \HA{prop:satdefect}(2) (the plateau $\varepsilon^{\mathrm{sat}}_{26}=
  \varepsilon^{\mathrm{sat}}_{27}=\delta^{\mathrm{sat}}$, then $\mathrm{length}(\mathsf A_{28})=28\,r-\delta^{\mathrm{sat}}$), and \texttt{CertifyStable} certifies them.
\end{enumerate}

\paragraph{Status at v1.9.0 (slot-8 WP-13).} Everything this subsection specifies is
\textbf{built as specified}; the text below it is the pre-release's, with the as-built facts
added at their sites, each marked ``v1.9.0''. \textbf{E8 item~2} (the composite
\texttt{ModCtx}, the Howell and Smith elimination, the module entry points of \S13) is slot-8 WP-1
(commit \texttt{9e02013}); \textbf{E8 item~3} (the pipeline over $\Z/\ell^k$, the $\rho_n$ repair
E8-P5, the key and its refusals) is slot-8 WP-3 (\texttt{9071ab0}), with the repairs of the
independent review (slot-8 WP-12, pass~2: findings M-1, M-2, m-6, m-7) as slot-8 WP-3r
(\texttt{28fdb73}); \textbf{E8 item~4a} (the reductions, the Bockstein, the identity of~(d), the
controls, the \ref{D6} arm) is slot-8 WP-4 (\texttt{d96d89d}), \textbf{accepted}, and its golden
emission R1 \textbf{landed} in the merge \texttt{53a1909}: the $124$ entries of~(e) at the seven
points, predicate-gated ($124$ added, $0$ changed, $0$ removed; tier~A $11\,472\to11\,575$, tier~B
$5845\to5866$), both determinism digests unchanged by construction
(\texttt{docs/matrix-engine/rebaseline-2026-09-24-e8-addonly.md}). \textbf{E8 item~4b} (the
boundary and cone halves, \HA{obl:coneexact}(iv)'s non-field half) is slot~9's (slot-8 ruling~3).
What measures it: \texttt{HeckeReductionValidation} (\texttt{standard}; WP-1's elimination arms,
WP-3's arms (K), (B), (P), (C), (T), (U), (X), WP-3r's arm (M) and WP-4's (R0)--(R5), (B1)--(B3),
\ref{D6}), \texttt{HeckeReductionCoverageSlowValidation} (\texttt{slow}; (R0)--(R4) and
(B1)--(B3) at $\Z/\ell^2$ on the E2 case list at $\ell\ge5$, $52$ cells, two of them behind the
read-only \texttt{HECKE\_REDUCTION\_COVERAGE\_LARGE=1}) and
\texttt{GOLDEN\_STRICT=1 GoldenRegressionValidation} (tier~A $11\,575$ failing~$0$, tier~B $5866$
mismatching~$0$). \textbf{Every predicted value of~(e) held}; the one literal that did not ---
footnote~b's ``$51\ne52$'' --- is in another model's coordinates, and is restated below.

\paragraph{(a) The key \texttt{coefficient\_ring}.}
\begin{itemize}
\item \textbf{Grammar.} $\texttt{coefficient\_ring}\in\{\texttt{F\_ell}\}\cup\{\texttt{Z/ell\textasciicircum}k : 2\le k\le 63\}$,
  spelled exactly so: the literal \texttt{F\_ell}, or the literal \texttt{Z/ell\textasciicircum}
  followed by $k$ in decimal without sign, leading zero or surrounding space. The token
  \texttt{ell} is literal: the setting is run-level and the prime is the cell's, so at a cell
  $(n,\ell)$ the value \texttt{Z/ell\textasciicircum2} means $\Z/\ell^2$ ($\Z/25$ at $\ell=5$,
  $\Z/691^2$ at $\ell=691$). The field has exactly one spelling:
  \texttt{Z/ell\textasciicircum1} and \texttt{Z/ell\textasciicircum0} are \emph{invalid}, as
  \texttt{prime\_power(ell, 1)} is rejected (\S13.1). The cap $63$ is the largest $k$ with
  $2^k<2^{64}$; whether $\ell^k$ fits a machine word is decided at the cell, by~(b3). \emph{v1.12.0-pre.2}: the cap is $4096$ (the pre-release paragraph
  ``v1.12.0-pre.2'' above); the cell decides, by that paragraph's per-cell rule.
\item \textbf{Reserved values.} \texttt{Z/<m>} with a literal integer $m\ge2$ (a general $\Z/m$; [X] \S1.5) is reserved because
  \S13 does not cover it (it needs gcd pivoting, not valuation pivoting). A reserved value is
  \emph{legal} --- a build-capability statement, not a user error --- and
  \texttt{MathSettings::set()} refuses it itself, by raising
  \texttt{hecke::NotImplementedInThisBuild} naming the key and the value (console exit~5): no
  consumer anywhere honours it. \emph{v1.23.0-pre.1} (D6, audit M2): \texttt{Z} (integral
  coefficients; line~2, [IB] OD-10~(a)) is \emph{no longer} reserved in this sense. By \HA{def:ringsemantics}(3) and
  \HA{sec:recordschema} item~7 it is stored, with a non-zero fingerprint, and honoured by the
  integral descent only (the verb \texttt{descent}, the campaign's descent cells, the record
  \texttt{integral\_descent\_record/v1}); every other consumer refuses it by name with
  \texttt{NotImplementedInThisBuild} (exit~5). The text before v1.23.0-pre.1 --- \texttt{Z} refused inside
  \texttt{set()}, where it was never in force --- is superseded by those two statements of~[C]. Every other text is
  invalid and raises \texttt{std::invalid\_argument}, as for every key.
\item \textbf{Default and canonical text.} The default is \texttt{F\_ell}. The canonical text of a
  value is its one spelling above; it is what \texttt{text(key)} returns, what the settings block
  carries under \texttt{value}, and what enters the fingerprint, as
  \texttt{coefficient\_ring=Z/ell\textasciicircum<k>}.
\item \textbf{Registration.} Slot-8 WP-3 registers the key by the four-edit rule of
  \texttt{Hecke/HeckeSettings}, \textbf{appended at the tail} of the registry (after
  \texttt{predictions}; the registry is append-only), making eleven keys, with
  $\texttt{source}\in\{\texttt{default},\texttt{spec},\texttt{console}\}$ as for every key
  (item~(g) of ``The stage-2 keys'' above). The record's \emph{top-level}
  \texttt{coefficient\_ring} key stays: it names the ring the record was computed over, and reads
  \texttt{F\_ell} in every record of slot~8. \emph{v1.9.0, as built} (slot-8 WP-3): registered
  exactly so, the eleventh key, with the three arms of \texttt{set()} (valid, reserved, invalid);
  the fingerprint is $0$ at the defaults (\texttt{HeckeReductionValidation} arm~(K),
  \texttt{ConsoleSettingsValidation}; the standing measure is
  \texttt{GOLDEN\_STRICT=1 GoldenRegressionValidation} at $0$ / $0$). A fresh golden emission is
  byte-identical to the reference, and the $22$ console transcripts re-emit unchanged; those two
  are one-off gate runs of slot-8 WP-3's report \S4~(c), not validations (v1.9.1: through v1.9.0
  they were attributed to the two arms). At the console the reserved values surface with
  exit~$5$ on the \texttt{set} statement (slot-8 WP-10's binding seam; \texttt{ConsoleErrorValidation}
  arm~(K), \texttt{HeckeReductionValidation} arm~(X)); \texttt{--set coefficient\_ring=Z} on the
  command line is exit~$1$, as every \texttt{--set} refusal is ([P] v3.6.1 \S6). Since v1.9.1
  (slot-8 WP-13, referee finding M-1) that code is pinned by one check in arm~(K) of
  \texttt{ConsoleErrorValidation}. Whether it should be exit~$5$ is an owner question recorded
  for C5. \emph{v1.10.0-pre.1: ruled} (slot-9 decision~2~(ii)~(A), 2026-09-24; it supersedes the
  lead default S8-D4). A \texttt{--set} argument whose value is \emph{legal but unexecutable} in
  this build --- the reserved \texttt{Z/<m>} of this key (and, before v1.23.0-pre.1, \texttt{Z}; see (a)), which
  \texttt{MathSettings::set()} refuses with \texttt{hecke::NotImplementedInThisBuild} --- is
  \textbf{exit~$5$} from slot-9 WP-12's commit, as on the \texttt{set} statement; arm~(K) of
  \texttt{ConsoleErrorValidation} changes its pinned~$1$ to~$5$ in that same commit, and the
  report names the change. An \emph{invalid} value (\texttt{Z/ell\textasciicircum1}, any other
  text) and an undeclared key stay exit~$1$. Until that commit, exit~$1$ as above. The same commit
  makes a \texttt{set}-statement refusal carry the statement's \texttt{file:line:col} (the
  second authorised \texttt{Console/Core/} edit of that decision, [P] v3.7).
  \emph{v1.10.0-pre.2, as built} (slot-9 WP-12, \texttt{b3cd0d3}): exactly so. The CLI maps a
  \texttt{--set} refusal of kind \texttt{NotImplementedInThisBuild} to exit~$5$, and an invalid
  value or an undeclared key to exit~$1$ (\texttt{Core/ConsoleCli}); \texttt{Evaluator::exec\_set}
  back-fills the statement's position into the refusal (\texttt{set coefficient\_ring = "Z"} now
  prints \texttt{file:1:1: not implemented in this build: \ldots}). Arm~(K) of
  \texttt{ConsoleErrorValidation} changed its pin $1\to5$ in the same commit and pins the three
  codes $5$ / $1$ / $0$ and the position. These are the only two \texttt{Console/Core/} edits of
  slot~9.
\item \textbf{Fingerprint and golden cost.} At the default the key contributes nothing to the
  fingerprint, which stays $0$; a non-default value makes it non-zero and so makes a new campaign
  job id, as for every key (the paragraph ``Settings, resumable state, coefficient ring'' above).
  \textbf{Registering the key at its default, and extending its grammar later to admit
  \texttt{Z}, cost no golden line}: since slot-7 stage~1~(1) golden sees the record's settings
  block only as its fingerprint (item~(f) of ``The stage-2 keys''), and a key registered at its
  default is proven key-additive ([X] v1.11 \S1.5). That proof covers a key at its default and
  nothing else: \textbf{a record field added for a ring costs golden lines} --- emitted at the
  defaults it moves the $62$ \texttt{ext\_json\_*} lines and the ext digest, as \texttt{s\_sep}
  did in slot~7 --- \textbf{unless it is emitted only under a non-default ring}, whose
  fingerprint-$\ne0$ texts are not golden. (Line~2's record keys \texttt{q\_max},
  \texttt{truncated}, \texttt{q\_max\_source} and \texttt{object} are therefore to be emitted only
  under \texttt{coefficient\_ring = Z}; v1.9.0: the owner confirmed this, slot-8 ruling~6d, at
  OD-10~(a) of \texttt{docs/implementation/integral\_line\_design\_brief.md}.) Registration still moves every campaign per-job
  \texttt{settings} block, never an id, a file name or a header.
\item \textbf{Where the ring is read.} The ring is chosen \emph{inside} the pipeline (\S13.8~(1)):
  an object that honours it reads \texttt{coefficient\_ring} \textbf{once, at construction}, from
  the \texttt{MathSettings} in force on the constructing thread, builds \texttt{ModCtx(ell)} for
  \texttt{F\_ell} and \texttt{ModCtx::prime\_power(ell, k)} otherwise, and keeps that context for
  its lifetime; it never re-reads the setting, and an object never changes ring. The composite
  route is selected by \texttt{is\_field()} alone, and the ring travels in the context, never as a
  bare modulus (\S13.1, \S13.7).
\end{itemize}

\paragraph{(b) The refusal rule (normative; slot-8 decision~4).}
\textbf{Every construction site and every field-only consumer that does not honour a non-default
\texttt{coefficient\_ring} raises \texttt{hecke::NotImplementedInThisBuild} naming itself
(console exit~5); none computes over $\Fl$ while the setting names $\Z/\ell^k$, and none reports a
dimension where a length is meant.} Precisely:
\begin{enumerate}[label=(b\arabic*)]
\item \emph{The closure.} The guarded set is the closure of the pipeline under ``constructs or
  consumes'': every construction site of \texttt{EquivariantCohomologyOnK} or
  \texttt{HeckeOperatorOnK} outside \texttt{*Validation.*}, and every consumer downstream of one
  whose algebra is a field notion (an RREF-defined basis, a dimension, a polynomial over $\Fl$, a
  primary decomposition, a pairing's perfectness, a residue field). The set is \emph{stated as this
  rule}, not as a list: slot-8 WP-3 derives its members by grep at its base and reports them, an
  independent reviewer re-derives them (slot-8 WP-12), and WP-13 copies the list into v1.9.0.

  \emph{v1.9.0: the list.} Derived by grep at \texttt{3ab190f} (slot-8 WP-3, report \S3),
  re-derived independently at \texttt{06bebdc} (slot-8 WP-12, pass~2, which found one member
  missing, finding m-5, because it was merged after WP-3's base), and re-derived at this revision's
  base \texttt{5c70e95} (slot-8 WP-13, which found one more construction site that the earlier
  greps did not list, the campaign's job loop; the campaign surface already covered it). Each arm
  named below is a member of arm~(X) of \texttt{HeckeReductionValidation} ($36$ members, plus the
  console's exit-code checks); ``setting and input'' means the member refuses both a non-default
  setting and a composite input object, with an arm for each.
  \begin{itemize}
  \item \textbf{Honours}: \texttt{EquivariantCohomologyOnK} and \texttt{HeckeOperatorOnK} (arms
    (C), (T) of \texttt{HeckeReductionValidation} at nine cell-levels, the U64\_Wide cell
    $(10,4294967291)$ over $\Z/\ell^2$ among them); \texttt{Hecke/HeckeReduction} (slot-8 WP-4:
    \texttt{red\_H\textasciicircum q}, the Bockstein, their certificates; it consumes honouring
    objects only and checks their rings); \texttt{HeckeReductionValidation},
    \texttt{HeckeReductionCoverageSlowValidation}; the golden harness inside \texttt{RunE8Point}
    only.
  \item \textbf{Refuses, by a named guard with an arm}: \texttt{HeckeModuleDecomposition} (and
    with it \texttt{FpPoly}'s consumers); \texttt{HeckeOperatorFamily} ($T_{p^r}$, $T(p,p)$,
    $U_\ell$; setting and input); \texttt{ContragredientCell} and
    \texttt{PoincareDualityPairing} (E7); \texttt{HeckeCrossChecks} (setting and input);
    \texttt{BarCocycleH1}; \texttt{HeckeBoundaryOperator}; \texttt{SL2ZBoundaryPair} (setting and
    input); \texttt{HeckeInvolution} (setting and input); \texttt{EquivariantCochainComplexMorphismOnK};
    \texttt{HeckeExt/HeckeBoundary} (setting and input); \texttt{HeckeConeStrictPairRoute} and
    \texttt{HeckeConeRelativeRoute}; and, inside the honouring class,
    \texttt{EquivariantCohomologyOnK}'s relative sub-complex constructor (the cone-base
    constructor is guarded too, without an arm of its own: its one caller is the mapping cone, a
    surface of~(b4)).
  \item \textbf{Refuses, by a named guard, unreachable} (no arm):
    \texttt{cochain\_commutes} (\texttt{HeckeBoundaryDiagnostics}); its
    \texttt{HeckeBoundaryOperator} and \texttt{HeckeOperatorOnK} arguments cannot be built over
    $\Z/\ell^k$ or under a non-default setting.
  \item \textbf{Refuses structurally}: every site below builds its local system with the field
    generator, so a prime modulus reaches \texttt{EquivariantCohomologyOnK}, whose own guard
    refuses it by name under a non-default ring. The golden harness outside \texttt{RunE8Point};
    \texttt{LocalSystemSL2Z}'s demonstration functions and \texttt{LocalSystemPresets};
    \texttt{Integral/IntegralCharpolyCrt} (\texttt{RunPipelineModEll}, slot-8 WP-7, reached only
    from \texttt{HeckeIntegralCharpolyValidation}), because line~2 is CRT over $\Fl$ runs by
    definition (WP-12 finding m-5; it carries no guard naming itself --- \emph{v1.10.0}: since
    slot-9 WP-9, \texttt{fb8a63a}, it does: \texttt{RunPipelineModEll} opens with a call of
    \texttt{hecke::RequireFieldCoefficientRing} naming itself, so a caller
    under a composite ring gets (b2)'s named refusal; the refusal was already total); and, behind
    their surfaces of~(b4), the campaign's job loop (\texttt{Campaign/CampaignDriver}; it honours
    from slot-9 WP-8, below), the \texttt{--bench} runners and the console bindings (in part, from
    slot-9 WP-12, below).
  \item \textbf{Unreachable} under a non-default ring, because every input is built by a refusing
    member (\texttt{HeckeModuleDecomposition}, \texttt{HeckeInvolution}, \texttt{HeckeExtDriver},
    or \texttt{HeckeCrossChecks}' own local system): \texttt{HeckeSturmSweep},
    \texttt{HeckeSignSplit}, \texttt{HeckeExtComplexRecords}, \texttt{HeckeAtlas},
    \texttt{HeckeTangentPair}, \texttt{HeckeLocalAlgebra}, \texttt{InducedLocalSystem}; and
    \texttt{kernel.cu}'s legacy benchmark, which runs at the process defaults: the argument-less
    run installs no setting, and no command-line switch sets \texttt{coefficient\_ring}. (v1.9.1:
    through v1.9.0 this read ``no command-line switch sets a mathematical key outside
    \texttt{--console} and \texttt{--campaign}'', which is false. \texttt{--bench ext --u-ell}
    installs \texttt{u\_ell}, source \texttt{spec}, through \texttt{ScopedUEll}
    (\texttt{BenchDriver.cpp:1273--1276}). It sets no ring, and \texttt{--bench} refuses a
    non-default ring at \texttt{RunMode}'s entry in any case.)
  \item \textbf{The seven surfaces} of~(b4), each with its arm.
  \end{itemize}

  \begingroup\sloppy
  \emph{v1.10.0-pre.1: the lift} (slot~9). The rule above is unchanged; slot~9 moves the members
  below, and no other. Each row names the member, its status through slot~8, its status from the
  named slot-9 package's commit, and the arm that tests it; ``arm~(X)'' is
  \texttt{HeckeReductionValidation}'s. A member that moves to ``honours'' is tested both ways: an
  honouring arm over $\Z/\ell^k$, \emph{and} the refusal arms of its field-only downstream
  consumers kept. \textbf{No row may be added by a code package}: a member found at a package's
  base that this list does not name --- the grep includes the template spellings
  \texttt{EquivariantCohomologyOnKT<} and \texttt{HeckeOperatorOnKT<} (slot-8 WP-13, finding~1) ---
  keeps refusing, and the package reports it for slot-9 WP-15 to record.
  \begin{itemize}
  \item \textbf{Honours at $\ell\in\{2,3\}$} (new row; slot-9 WP-3, decision~4):
    \texttt{EquivariantCohomologyOnK}'s ordinary constructor and \texttt{compute\_cohomology()},
    through the flat-coordinate route of~(c1$'$), with the ring's relation list
    (\texttt{GenerateSL2ZRelationsRing}) and $\rho_n$ over $\Z/\ell^k$ (E8-P5's division-free
    Pascal route, validated entrywise against the division-free oracle on every $g$ the tested
    cells touch). It reports \texttt{lengths\_H\textasciicircum q} and
    \texttt{invariant\_factors\_H\textasciicircum q} only. Test:
    \texttt{HeckeBredonTreeRingValidation} (\texttt{standard}, new) against the frozen tree model
    over the same ring (O7, \HA{prop:bredon}(2)). \emph{The lift is scoped} (as built, slot-9
    WP-3, \texttt{b53836f}): \texttt{CoefficientRingContext} honours $\ell\in\{2,3\}$ \emph{only}
    under the policy argument \texttt{SmallEllPolicy::Bredon}, which \texttt{EquivariantCohomologyOnK}
    alone passes; under the default policy it keeps refusing at $2$ and $3$ by name, and that
    refusal is its own row of~(X). Every other member that honours at $\ell\ge5$ therefore keeps
    refusing $\ell\in\{2,3\}$ through the default policy, with an arm; inside
    \texttt{EquivariantCohomologyOnK} the four stored-basis accessors and the six $K$-mode
    accessors refuse there, and so do the console's pipeline verbs (exit~$5$). There,
    \textbf{refuse by name, each with an arm}: \texttt{HeckeOperatorOnK} (setting and input:
    a composite object at $\ell\in\{2,3\}$ can now be built, so the input arm is reachable) and
    \texttt{HeckeReduction} (its four entry points), by \texttt{hecke::NotImplementedInThisBuild} (a
    build capability; M1 of line~2 is slot~10's); and the $K$-mode accessors of
    \texttt{EquivariantCohomologyOnK} (\texttt{ladic\_type}, \texttt{ladic\_type\_certified},
    \texttt{ladic\_e\_max}, \texttt{ladic\_level\_state}, \texttt{r\_K}, and the sixth WP-3
    guards), each by its named refusal. The $K$-mode is \emph{undefined} there, not merely unbuilt
    (\HA{prop:certifiede}(5): the model complex over $\Z/\ell^e$ at $\ell\in\{2,3\}$ carries no
    integral lattice to certify). This pre-release first specified a named \texttt{std::logic\_error}
    for it; WP-3 built the default-policy \texttt{NotImplementedInThisBuild} (exit~$5$), which also
    satisfies~(b2), and WP-15 records which kind stands. The flipped arms are named in~(b3).
    \emph{v1.10.0-pre.2: \texttt{hecke::NotImplementedInThisBuild} stands} (exit~$5$ at the
    console). Both readings are conformant; the built one is kept because the refusal is made
    where every other member refuses $\ell\in\{2,3\}$ --- by the default policy of
    \texttt{CoefficientRingContext}, or by the accessor's own guard naming it --- and because a
    caller cannot tell ``undefined here'' from ``not in this build'' by anything it could do
    differently; the text of the refusal names \HA{prop:certifiede}(5). The slot's review
    (WP-14, pass~1; v1.10.1: \texttt{docs/implementation/slot9/WP-14\_review.md}) recommended recording it so. As built (\texttt{97d0a8d}): the ten refusing
    accessors are the four stored-basis ones (\texttt{module\_degree}, \texttt{C\_eq\_rank},
    \texttt{H\_module\_basis}, \texttt{H\_class\_coordinates}) and the six $K$-mode ones
    (\texttt{r\_K}, \texttt{ladic\_applicable}, \texttt{ladic\_type\_certified},
    \texttt{ladic\_type}, \texttt{ladic\_e\_max}, \texttt{ladic\_level\_state}); the honouring ones
    are \texttt{num\_degrees}, \texttt{H\_lengths}, \texttt{H\_invariant\_factors},
    \texttt{ring\_context}, \texttt{flat\_module\_degree} and \texttt{is\_flat\_route}; the field
    accessors keep their named \texttt{std::logic\_error}. \texttt{HeckeOperatorOnK} refuses in
    \texttt{init\_composite()}; \texttt{HeckeReduction} refuses through the stored-basis accessors,
    and since slot-9 WP-4 its new cone and boundary entry points carry an explicit guard. Arm~(X)
    of \texttt{HeckeReductionValidation} went from $36$ to $66$ members (two changed rows, $30$
    new), and WP-4's (X4b) adds $21$ rows. \emph{Scoping}: \texttt{SmallEllPolicy::Bredon} is a
    public enumerator; that no caller but \texttt{EquivariantCohomologyOnK} passes it is held by
    that class's comment and by arm~(X)'s rows for the callers that exist, not by a guard script
    (review finding m-1, open). \emph{v1.10.0: closed} (slot-9 WP-9, \texttt{fb8a63a}):
    \texttt{SmallEllPolicyScopeValidation} (\texttt{fast}) scans the tree and allows the token
    \texttt{SmallEllPolicy::Bredon} only in \texttt{EquivariantCohomologyOnK.cpp},
    \texttt{Hecke/CoefficientRing.\{h,cpp\}}, validation files and (v1.10.1) the guard's own
    translation unit \texttt{Registry/Slot9/WP9.cpp}, and asserts that the one caller's
    site still exists. It caught the one violation the slot produced: WP-8's ring job named the
    policy at $\ell\in\{2,3\}$, a merge interaction of WP-8 and WP-9, repaired by \texttt{158c0be}
    as the lead ruled --- the job loop builds its relation list from the raw
    \texttt{ModCtx::prime\_power} context (the same value), and the Bredon lift stays with
    \texttt{EquivariantCohomologyOnK} alone, whose constructor re-derives the context under the
    policy. The allow-list did not change.
  \item \textbf{Moves from ``refuses'' to ``honours'' at $\ell\ge5$} (slot-9 WP-4, E8 item~4b,
    decision~3~(A) with~(iii)~(A); objects in~(f)): \texttt{SL2ZBoundaryPair} (the boundary
    complex, an \texttt{EquivariantCohomologyOnK} on $K_{\mathrm{bd}}$, and the restriction);
    \texttt{EquivariantCochainComplexMorphismOnK}'s \emph{constant-fibre} constructor (the
    restriction $r=j^{\ast}$ with fibre $I_{d_V}$, both complexes over the same ring; a ring
    mismatch between domain and codomain is a wrong input and raises a named
    \texttt{std::logic\_error}; \emph{v1.10.0}: since slot-9 WP-9, \texttt{fb8a63a}, its composite
    branch reads the modulus from \texttt{ring\_context()}, because
    \texttt{EquivariantCohomologyOnK::get\_prime()} now refuses a composite object by
    \texttt{ENGINE\_CHECK} --- on one it held $\ell^k$, not a prime --- the value is unchanged),
    its induced maps given in the (c5) coordinates of~(c);
    \texttt{EquivariantMappingCone} (the dense module route of~(f), selected by
    \texttt{is\_field()}; the tape route stays the $\Fl$ path); \texttt{HeckeConeAction},
    \textbf{route~1 only}; \texttt{HeckeBoundaryOperator} ($T_p^{\partial}$ on the cyclic summands
    of $H^q(\partial)$, by \texttt{solve\_multi\_rhs\_module}); \texttt{Hecke/HeckeReduction}, which
    grows the reductions and the Bockstein on the boundary and the cone (its $\ell\ge5$ scope
    unchanged); the golden harness inside \texttt{RunE8Point}, which grows the objects of~(f)
    under the (a/s) event R1 only. Tests: arms of \texttt{HeckeReductionValidation} at the four
    $\ell\ge5$ E8 cells at $k\in\{2,3\}$ and at the U64\_Wide cell $(10,4294967291)$ over
    $\Z/\ell^2$; the E2 case list at $\ell\ge5$ in \texttt{HeckeReductionCoverageSlowValidation}.
    \emph{v1.10.0: the tiers as trimmed} (slot-9 WP-4t, \texttt{45acc4f}, owner-approved; test
    code only): the $(60,5)$ cell of the item-4b arms, unchanged in content (both levels, both
    primes, (4b-B/C/R/H/K), the commutators, $K=3\to2$, the Bockstein), runs in
    \texttt{HeckeReductionCoverageSlowValidation} (\texttt{slow}, $99$ checks on its own
    \texttt{E8\_4B\_POINT} line), so \texttt{--all} gates the ring Hecke action on the cone and
    boundary at $(6,5)$, $(16,5)$, $(30,5)$ and $(10,691)$ and the two word-regime cells, and
    $(60,5)$'s values stay gated in \texttt{--all} by R1; and rows implied by a kept row were
    dropped (the \texttt{coneexact\_iv\_certified} and \texttt{routes\_bit\_identical} flag reads,
    the image-length vector and $\ell(H^1_c)$, implied by \texttt{ConeModuleCertificate::verify()}
    $=0$ with $H^0_c=0$ and $H^2=0$, and the boundary's equal lengths, implied by its equal
    invariant factors). The item-4b arms went from $458$ to $306$ checks,
    \texttt{HeckeReductionValidation} from $69\,122$ to $68\,970$, and the coverage list's 4b part
    from $1482$ to $1232$ (\texttt{HeckeReductionCoverageSlowValidation} $3882\to3731$).
    Each of these refuses at $\ell\in\{2,3\}$ and wherever $\ell^k\ge2^{64}$ by the same named
    error, with an arm. \emph{v1.10.0-pre.2, as built} (slot-9 WP-4, \texttt{d0aca96}; the closure
    re-derived by grep at \texttt{b53836f}, the template spellings included: guard sites in $28$
    files, construction sites in $13$ non-validation TUs, the same set as this list). Every
    member above honours, and \textbf{only} when its inputs are over the ring in force; every
    other combination keeps slot~8's named refusal, and every field accessor of an honouring
    object throws a named \texttt{std::logic\_error}. \textbf{The width}: the route itself ---
    \texttt{ModuleConeT}, \texttt{ModuleConeActionT}, \texttt{HeckeBoundaryOperatorT} --- is
    templated on the ring context's element type and runs in \textbf{every} word regime (the
    \texttt{U64\_Montgomery} cell $(10,134217757)$ and the \texttt{U64\_Wide} cell
    $(10,4294967291)$ over $\Z/\ell^2$ run it end to end); the four \emph{classes}
    \texttt{EquivariantMappingCone}, \texttt{SL2ZBoundaryPair}, \texttt{HeckeConeAction} and
    \texttt{EquivariantCochainComplexMorphismOnK} are \texttt{uint32\_t}-instantiated, as they
    always were, so they honour $\ell^k<2^{32}$ and refuse above it by the ring context's named
    refusal ``in the \texttt{uint32\_t} instantiation'' --- a state no caller reaches, because a
    \texttt{uint32\_t} bulk object cannot exist there either. At $\ell=5$ the cone's reductions to
    $\Fl$ run (the $\Fl$ cone class being \texttt{uint32\_t}); at the two word-regime cells they do
    not, and the dense route, the boundary, both Hecke actions and the boundary reductions do.
    \texttt{SL2ZBoundaryPair} builds the circle model over the ring with a new generator,
    \texttt{GenerateSL2ZBoundaryRelationsRing}, under the default \texttt{SmallEllPolicy}. Tests:
    \texttt{HeckeReductionValidation} arms (4b-B), (4b-C), (4b-H), (4b-R), (4b-K) and (X4b).
  \item \textbf{Keeps refusing, with an arm} (decision~3~(iii)~(A), slot-9 plan M-3):
    \texttt{HeckeConeStrictPairRoute}, \texttt{HeckeConeRelativeRoute}, and
    \texttt{EquivariantCohomologyOnK}'s relative sub-complex constructor. So
    \HA{obl:coneexact}(vii) over $\Z/\ell^k$ is \textbf{route~1 only}: the agreement of the three
    routes is asserted over $\Fl$ as before, and \textbf{(vii)'s composite half is open}, recorded
    here and in [C] by slot-9 WP-15 (v1.11.0-pre.2: closed for the first two sentences of (vii) by
    slot-10 WP-3 --- ``the lift (slot~10)'' below and the declaration of the status paragraph). Also kept: \texttt{EquivariantCochainComplexMorphismOnK}'s
    \texttt{FiberBlockTable} constructor (its callers are the involution and the E7 objects, field
    notions); \texttt{HeckeExt/HeckeBoundary} (setting and input: it consumes
    \texttt{SL2ZBoundaryPair}, which now honours, and keeps its own guard); every other
    member of the list above. \texttt{EquivariantCohomologyOnK}'s protected (prime, rep\_dim)
    constructor keeps its guard: no class derives from \texttt{EquivariantCohomologyOnKT} at the
    slot-9 base (grep), so it is unreachable, and the cone's module route does not construct through
    it --- a package that must, stops and reports. \emph{v1.10.0-pre.2, as built} (slot-9 WP-4):
    all kept; the module route did not need the protected constructor. The two cross-routes
    now also refuse a composite \emph{normative action} handed to them (their guard gained the
    input argument), with an arm each in (X4b), beside the pre-existing (X) rows.
  \item \textbf{Moves from ``unreachable'' to ``refuses, with an arm''}: \texttt{cochain\_commutes}
    (\texttt{HeckeBoundaryDiagnostics}). Once \texttt{HeckeBoundaryOperator} and
    \texttt{HeckeOperatorOnK} can be built over $\Z/\ell^k$, its named guard is reachable, and
    slot-9 WP-4 gives it an arm in~(X). \emph{v1.10.0-pre.2}: done, a row of (X4b).
  \item \textbf{Moves from ``refuses structurally, behind its surface'' to ``honours''}: the
    campaign's job loop (\texttt{Campaign/CampaignDriver}; slot-9 WP-8, C5c, decision~6), under the
    schema of~(g); its ext pass keeps refusing, at spec validation (b4). \emph{v1.10.0, as built}
    (slot-9 WP-8, \texttt{ad65bd9}; merge fix \texttt{158c0be}): the job loop honours
    $\Z/\ell^k$ at every $\ell$ with $\ell^k<2^{64}$ (\texttt{uint32\_t} below $2^{32}$,
    \texttt{uint64\_t} above). At $\ell\in\{2,3\}$ it carries the Bredon object's lengths and
    invariant factors, built by \texttt{EquivariantCohomologyOnK} (the job loop hands it a relation
    list over the raw \texttt{ModCtx::prime\_power} context and names no policy), and records
    \texttt{hecke} (\texttt{HeckeOperatorOnK}), \texttt{bockstein} and \texttt{reduction}
    (\texttt{HeckeReduction}) as refused; at $r\ge2$ it records \texttt{hecke} refused at its own
    dispatch. Tested by \texttt{CampaignRingValidation} arms (A), (T), (B), (R), (S).
  \item \textbf{Moves in part}: the console bindings (slot-9 WP-12, C5a, decision~2): exactly the
    verbs and members of~(h) honour; every other pipeline verb keeps exit~$5$, through a lazy handle
    too (\texttt{ConsoleErrorValidation} arm~(K), extended). \emph{v1.10.0-pre.2, as built}
    (slot-9 WP-12, WP-13, \texttt{f28a0fa}): exactly (h)'s list honours at $\ell\ge5$; every other
    member refuses through one mechanism (\texttt{bind::GateTypeOnRing}), so a member added later
    refuses until (h) lists it; every C5b verb that runs the pipeline refuses a non-default ring at
    dispatch and again in its lazy compute, and every C5b verb handed a ring handle refuses it at
    dispatch (\texttt{bind::RequireFieldHandle}), exit~$5$. Arm~(K): $61$ single-access scripts
    under $\Z/\ell^2$ at $(10,691)$ --- $25$ honour, $8$ exit~$3$, $28$ exit~$5$ --- and five rows at
    the uncertified $(30,5)$; $210$ checks at the base.
  \end{itemize}
  \endgroup

  \begingroup\sloppy
  \emph{v1.11.0-pre.1: the lift (slot~10).} The rule above is unchanged, and so is slot~9's rule
  for a lift: each row names the member, its status through slot~9, its status from the named
  slot-10 package's \emph{lift commit} (the last commit of that package's branch, made after this
  pre-release merges; slot-10 plan \S4.2), and the arm that tests it. A member that moves is tested
  both ways; \textbf{no row may be added by a code package} --- a member found at a package's base
  that this list does not name keeps its status (the grep includes the template spellings
  \texttt{EquivariantCohomologyOnKT<}, \texttt{HeckeOperatorOnKT<}, \texttt{ModuleConeT<} and \texttt{ModuleConeActionT<}),
  and the package reports it for slot-10 WP-15 to record.
  \begin{itemize}
  \item \textbf{Moves from ``keeps refusing'' to ``honours'' at $\ell\ge5$, $\ell^k<2^{32}$}
    (slot-10 WP-3; decision~3~(A)): \texttt{HeckeConeStrictPairRoute},
    \texttt{HeckeConeRelativeRoute}, and \texttt{EquivariantCohomologyOnK}'s \emph{relative
    sub-complex constructor} (the refusal flipped is its \texttt{ring\_k\_ >= 2} guard). Their
    route over the ring, and the acceptance, are~(f)'s (vii) bullet. \emph{Regime rule} (slot-10
    plan audit M-3): both cross-routes consume the \texttt{uint32\_t}-instantiated
    \texttt{HeckeConeAction} and are \textbf{not templated}; they honour exactly where that class
    honours, $\ell^k<2^{32}$, and above it refuse by the ring context's named refusal ``in the
    \texttt{uint32\_t} instantiation'', as every \texttt{uint32\_t} class of the slot-9 row above
    does --- with \textbf{one arm} at a 4b word-regime cell (\texttt{U64\_Montgomery}
    $(10,134217757)$ or \texttt{U64\_Wide} $(10,4294967291)$ over $\Z/\ell^2$), where route~1
    (\texttt{ModuleConeActionT}) runs and routes~2 and~3 refuse. The relative constructor is a
    member of the \texttt{EquivariantCohomologyOnKT} template and follows its object's element
    type; it honours in both instantiations at $\ell\ge5$, $\ell^k<2^{64}$, but only its
    \texttt{uint32\_t} instantiation has a consumer (route~3). \emph{Kept}: each of the three
    refuses at $\ell\in\{2,3\}$ by name (the relative constructor included: its $\ell\ge5$ module
    route does not apply there, and no flat-coordinate relative route is specified), with an arm;
    the \textbf{cone-base constructor} of \texttt{EquivariantCohomologyOnK} keeps its guard (the
    module route of~(f) does not derive from it), and so do the protected (prime, rep\_dim)
    constructor, \texttt{EquivariantCochainComplexMorphismOnK}'s \texttt{FiberBlockTable}
    constructor and \texttt{HeckeExt/HeckeBoundary}; \HA{prop:coneroutes}(3)'s adjointness over the
    ring is \textbf{not} lifted (it needs the E7 pairing, which refuses a composite ring, (b1)'s
    list). The two cross-routes' \emph{input} guard (a composite normative action under a
    non-default setting at $\ell\in\{2,3\}$ or above $2^{32}$) keeps its (X4b) rows; their
    \emph{honouring} rows change from ``refuses'' to ``honours'' and are reported as \emph{changed}
    rows, with the reason. Tests: \texttt{HeckeConeCrossRouteRingValidation} (new, slot-10 WP-3; its
    arms are~(f)'s), and arms~(X)/(X4b) of \texttt{HeckeReductionValidation}, whose check count
    moves only by the changed rows.
    \emph{v1.11.0-pre.2, as built} (slot-10 WP-3). One guard, \texttt{CrossRouteRingGuard}, in both
    cross-routes: they honour iff the normative action is composite ($k\ge2$), its $k$ equals the
    $k$ in force, and \texttt{CoefficientRingContext<uint32\_t>} admits $\ell^k$ (the regime
    refusal); every other combination is slot~8's \texttt{RequireFieldCoefficientRing} refusal, as
    before. The relative sub-complex constructor honours at $\ell\ge5$ and refuses at
    $\ell\in\{2,3\}$ by name (``\ldots\ the cone and its route~3 are not honoured at ell in
    \{2, 3\}''); the cone-base guard is untouched. Additive API, all within the three members: a
    static \texttt{RequireRingRegime(ell, k)} on each route; a \emph{validation seam} constructor of
    the strict route taking an external strict table, non-owning and re-certified in
    \texttt{compute()} (so it cannot make an uncertified table look agreed), used only by the
    negative control; the relative route's \texttt{*\_module} accessors, its field accessors now
    throwing a named \texttt{std::logic\_error} on a composite route. No member outside the list was
    found at the base, so no row was reported for this version to record. Arms (X4b) and (X) of
    \texttt{HeckeReductionValidation} (in \texttt{HeckeReductionConeValidation.cpp} and
    \texttt{HeckeReductionPipelineValidation.cpp}): three rows changed in place from ``refuses'' to
    ``honours'', each with its reason in the code, the count unchanged; the input-guard rows stay.
    \emph{v1.11.0} (slot-10 WP-7): \textbf{none} in the cone classes
    (\texttt{EquivariantMappingCone*}, \texttt{HeckeConeAction}) or in \texttt{HeckeSignSplit} ---
    the cone-certificate lever was recorded for slot~11, not shipped. Two levers landed elsewhere in
    the closure, both value-preserving: L4 in \texttt{HeckeOperatorOnK}'s route~2 (an honouring
    member; rows run concurrently, every key and program the serial loop's, and the E8 golden
    points, computed through it over the ring, unchanged under \texttt{GOLDEN\_STRICT=1}), and the
    $\iota$ lever in \texttt{HeckeInvolution} (a refusing member; only the order of independent
    checks on its field path, its guard and arm untouched). Evidence: the status paragraph's lane~B.
  \item \textbf{Keeps refusing at $\ell\in\{2,3\}$, now for a stated reason} (decision~4~(A)):
    \texttt{HeckeOperatorOnK} (setting and input) and \texttt{HeckeReduction} (its entry points),
    by \texttt{hecke::NotImplementedInThisBuild}. The reason is~(b3)'s ``v1.11.0-pre.1'' block;
    the rows, arms and messages are unchanged.
  \item \textbf{Outside the closure: line~2's free model} (slot-10 WP-10, WP-13; slot-9 WP-10,
    WP-11). \texttt{Integral/Model/IsotropyModel}, \texttt{Integral/Model/IntegralGroups} and
    slot~10's \texttt{Integral/Model/StrictLift} ($\tau_{\mathrm{Iso}}$, $\widehat\Theta_g$) and
    \texttt{Integral/Model/PeriodicityClass} ($e$, $E_q$) --- and slot-10 WP-11's
    \texttt{Integral/Model/IntegralHecke}, if it lands --- are \textbf{not members} of~(b1)'s
    closure: none constructs \texttt{EquivariantCohomologyOnK} or \texttt{HeckeOperatorOnK} outside
    \texttt{*Validation.*}, none consumes one, and each takes its ring as an explicit
    \texttt{engine::ModCtx} argument and reads no mathematical setting (no settings key, no record
    key, no golden object). \texttt{coefficient\_ring} therefore neither selects nor refuses their
    ring, at any $\ell$, $2$ and $3$ included; what they compute over $\Z/\ell^k$ at $\ell\in\{2,3\}$
    is $H^{\bullet}(\Ga;V_R)$ (\HA{prop:isomodel}(3)), not the Bredon object of~(b3), and they do
    not refuse there. \emph{The read-only use of an honouring member} (slot-10 WP-10 scope~4): the
    strict cochain map $\Theta^{s}_g$ is \texttt{HeckeOperatorOnK} built, inside
    \texttt{HeckeStrictLiftValidation}, on the \texttt{BoundaryStrict} $\tau$ table ($\tau^{s}$,
    integral, solve regime \texttt{"Z"}), over $\Fl$ at every $\ell$ and over $\Z/\ell^k$ at
    $\ell\ge5$; that is \texttt{HeckeOperatorOnK} honouring as listed, on another admissible table.
    \textbf{No row moves for it}: if the composite constructor refuses that table, the refusal
    stands, slot-10 WP-10 does not edit the class, and its $\Z/\ell^k$ arm drops to~$\Fl$ and is
    reported. Over $\Z/\ell^k$ at $\ell\in\{2,3\}$ no $\Theta^{s}_g$ exists in this build (decision~4,
    (b3)); the strict compatibility of \HA{prop:strictlift}(3) there follows from the exact
    $\Z$-certificate $\varepsilon\circ\tau_{\mathrm{Iso}}=\tau^{s}\circ\varepsilon$ over $\Z$, the coset transport
    \HA{eq:cosettransport} being functorial in the ring, and is not separately checked.
    \emph{v1.11.0-pre.2, as built} (slot-10 WP-10, WP-13, WP-11). As specified:
    \texttt{Integral/Model/StrictLift}, \texttt{PeriodicityClass} and \texttt{IntegralHecke} take
    an explicit \texttt{engine::ModCtx} (\texttt{uint32\_t} or \texttt{uint64\_t}), read no setting,
    add no key, record or golden object, and construct pipeline objects only inside
    \texttt{*Validation.*}; \texttt{IsotropyModel} and \texttt{IntegralGroups} are unedited. WP-10's
    step~1 found that \texttt{HeckeOperatorOnK}'s composite constructor \textbf{accepts} the
    \texttt{BoundaryStrict} table at $k\ge2$ (its three refusals --- the Bredon primes, the
    contragredient, a $\tau$ built by the mod-$\ell$ fallback --- do not apply: the table's solve
    regime is \texttt{"Z"}), so the $\Z/\ell^k$ arm ran at $\ell\in\{5,7\}$, $k\in\{2,3\}$, and no
    row moved. The sentence above is \textbf{completed} (the slot's review,
    \texttt{docs/implementation/slot10/WP-14\_review.md}, finding m-1): the identity over $\Z/\ell^k$
    at $2$ and $3$ follows from the proof of \HA{prop:strictlift}(3), a formal computation valid over
    every commutative ring, whose \emph{three} inputs are certified over $\Z$ or hold by
    construction --- the strict identity $\varepsilon\circ\tau_{\mathrm{Iso}}=\tau^{s}\circ\varepsilon$,
    the twisted equivariance (T1) of $\tau_{\mathrm{Iso}}$ (which makes the coset transport well
    typed, \HA{lem:cosettransport}(1)), and the $\Ga$-equivariance of $\varepsilon$; the certificate
    alone is one of the three ([C] v1.10.0 \HA{rmk:strictliftbuilt}(3)).
  \item \textbf{Moves in part: the console} (slot-10 WP-12; decision~6~(ii)~(A)): the five
    structural members of~(h)'s ``v1.11.0-pre.1'' block move from ``refuses'' (through
    \texttt{bind::GateTypeOnRing}) to ``honours'' on a ring handle at $\ell\ge5$; every other
    unlisted member keeps exit~$5$; arm~(K) of \texttt{ConsoleErrorValidation}.
    \emph{v1.11.0-pre.2, as built} (slot-10 WP-12): as listed --- (h)'s ``v1.11.0-pre.2'' block.
  \end{itemize}
  \endgroup
\item \emph{Each member is exactly one of three.} It \textbf{honours} the ring, with a test over
  $\Z/\ell^k$; or it \textbf{refuses}, by one constructor-entry check raising
  \texttt{hecke::NotImplementedInThisBuild} whose text names the class and the canonical value
  (for example \texttt{"<Class>: coefficient\_ring=Z/ell\textasciicircum<k> is not honoured in this
  build"}), with a test arm; or it is \textbf{unreachable} under a non-default ring, with the
  reason. None is silent. A refusing member refuses when the setting in force is non-default
  \emph{or} when an input object it is handed was constructed over a composite ring: by~(a), an
  object's ring is its own, whatever is in force when it is consumed.
\item \emph{What honours it in slot~8.} \texttt{EquivariantCohomologyOnK},
  \texttt{HeckeOperatorOnK}, \texttt{HeckeReductionValidation}, the golden harness \textbf{inside}
  its E8 points only (\texttt{RunE8Point}, which installs the ring around the E8 computations and
  nowhere else, (e)), and optionally one \texttt{--bench} switch of WP-3's choosing. An honouring
  object refuses, by the same named error, at a cell where it cannot honour the value: at
  $\ell\in\{2,3\}$ (there $6\notin R^\times$, and \HA{prop:coneintegral}~(2)'s isomorphism $C\otimes R\cong C_R$ is not available,
  so a $\Z/\ell^k$ run is not the reduction of one integral complex; the finer model is line~2's), and wherever $\ell^k\ge2^{64}$ (\S13's
  scope; the $\ell$-adic induction E8-P6 is deferred) --- never by an \texttt{ENGINE\_CHECK} abort.
  Over $\Z/\ell^k$ no honouring object reports a \texttt{dim}: \texttt{length} (\S13.6) replaces
  it, and a dimension accessor on a composite object throws a named \texttt{std::logic\_error}.
  \emph{v1.9.0, as built.} WP-3 chose \textbf{no} \texttt{--bench} switch, so \texttt{--bench}
  refuses every mode (arm~(X)). The honouring object's refusals are arms of~(X): $(10,3)$ under
  $\Z/\ell^2$, and, through \texttt{CoefficientRingContext}, $\ell=2$, a \texttt{u32} overflow
  and $\ell^k\ge2^{64}$. The range $\ell^k\in[2^{63},2^{64})$ (U64\_Wide) is
  \emph{honoured and exercised} end to end by the cell $(10,4294967291)$ over $\Z/\ell^2$ (slot-8
  WP-3r, on review finding M-2); no narrowing of this item is needed. The refusal at
  $\ell\in\{2,3\}$ has its mathematics in [C] v1.8.0: \HA{def:ringsemantics}(2) (the object at
  $\ell\in\{2,3\}$ is the Bredon one, not a reduction of $C(\Z)$) and \HA{prop:certifiede}(5); it
  was WP-2's choice, put to the owner in WP-2's report, and not objected to (v1.10.0-pre.1:
  \textbf{superseded} by slot-9 decision~4~(A), which supersedes the lead default S8-D2; see the
  next paragraph). The converse caller
  error --- a local system over $\ell^j$, $j\ge2$, handed to \texttt{EquivariantCohomologyOnK}
  while \texttt{F\_ell} is in force --- raises a named \texttt{std::logic\_error}, not
  \texttt{NotImplementedInThisBuild}: it is a wrong input, not a build capability (review finding
  m-6; a member of arm~(X)).

  \begingroup\sloppy\emph{v1.10.0-pre.1: $\ell\in\{2,3\}$, and what honours in slot~9} (decision~4~(A); slot-9 plan
  audit B-1, M-4). From slot-9 WP-3's commit, \texttt{EquivariantCohomologyOnK} \textbf{honours}
  $\Z/\ell^k$ at $\ell\in\{2,3\}$ (with $\ell^k<2^{64}$ and the element-type rule unchanged). The
  object it computes there is the one [C] names: the Bredon cohomology over $\Z/\ell^k$
  (\HA{def:ringsemantics}(2)), which is the tree complex \HA{eq:treemodel} over the ring in
  force (\HA{prop:bredon}(2)--(3)) --- not a reduction of $C(\Z)$ and not group cohomology
  (\HA{prop:ringinforce}(2)--(3)). It reports the \textbf{lengths and invariant factors} of
  $H^q$ for $q=0,1,2$, by the flat-coordinate route of~(c1$'$), and nothing else: \textbf{the
  certified $K$-mode refuses there} (\HA{prop:certifiede}(5); by its named refusal,
  (b1)), and \texttt{HeckeOperatorOnK}, \texttt{HeckeReduction}, the boundary and cone classes of~(f)
  and every field-only consumer \textbf{keep refusing} at $2$ and $3$ with
  \texttt{NotImplementedInThisBuild}, each with an arm. Its oracle is \textbf{O7}, the tree model
  over the same ring ([IB] \S6.2 M0 item~3), frozen from the probe's \texttt{model\_R} in
  \texttt{HeckeBredonTreeRingValidation}; an in-process tree model over $\Z/\ell^k$ reported beside
  it shares the Howell code with the pipeline and is a consistency check, not an oracle ([IB] K6).
  \textbf{The flipped arms} (M-4). Two as-built refusal rows of~(X) name $\ell\in\{2,3\}$, and
  both \emph{flip under \texttt{EquivariantCohomologyOnK}'s policy} (\texttt{SmallEllPolicy::Bredon})
  in slot-9 WP-3's commit (\texttt{b53836f}), each reported there as a \emph{changed} row:
  ``\texttt{CoefficientRingContext} at $\ell=2$'' (the one the slot-9 plan names; under the default
  policy the refusal stays, as its own row) and ``\texttt{EquivariantCohomologyOnK} at $(10,3)$''
  under $\Z/\ell^2$ (the paragraph above's ``$(10,3)$''), which now honours and asserts O7. The $\ell=3$ refusal they carried moves to new rows for
  \texttt{HeckeOperatorOnK} and \texttt{HeckeReduction} at $(10,3)$, and for the $K$-mode
  accessors there. The \texttt{u32}-overflow and $\ell^k\ge2^{64}$ rows keep refusing. \textbf{A pointer into [C]}: [C]
  \HA{def:ringsemantics}(2)'s clause ``the implementation of this build refuses it there'' is
  superseded by this item from v1.10.0-pre.1 --- in the tree, from slot-9 WP-3's merge --- until
  [C]'s next edit (slot-9 WP-15), which repairs it under slot-9 decision~10 either way; the window
  is recorded, not closed by a second [C] editor (slot-9 plan \S4.3). \emph{v1.10.0-pre.2}: the
  window is closed --- [C] v1.9.0 (slot-9 WP-15) replaces the clause by what this build computes
  at $\ell\in\{2,3\}$ and what it refuses. \emph{What honours in slot~9},
  beyond slot~8's list above: \texttt{EquivariantCohomologyOnK} at $2$ and $3$ (WP-3); the boundary
  and cone members of~(f) at $\ell\ge5$ (WP-4); the campaign's job loop (WP-8, (g)); the console
  verbs of~(h) (WP-12) --- the full list is~(b1)'s ``the lift''. Still \textbf{no} \texttt{--bench}
  switch.\par\endgroup

  \begingroup\sloppy\emph{v1.11.0-pre.1: the Hecke action on the Bredon object at $\ell\in\{2,3\}$ is
  declined} (slot-10 decision~4~(A); slot-10 plan \S2.3 item~2, audit B-1, M-2, M-6). At
  $\ell\in\{2,3\}$ the object \texttt{EquivariantCohomologyOnK} computes over $\Z/\ell^k$ is the
  Bredon object $C^{\mathrm{eq}}(\Z/\ell^k)$ (\HA{def:ringinforce}(3), \HA{prop:bredon}(2)--(3)), and
  two different Hecke objects on it could be lifted. \textbf{Neither is, in slot~10}, for two
  different reasons, stated separately:
  \begin{enumerate}[label=(\alph*)]
  \item \emph{The cochain map $\Theta_g$ on $C^{\mathrm{eq}}(\Z/\ell^k)$} --- the coset transport
    \HA{eq:cosettransport} of an admissible $\tau$ (the shipped canonical one, or the strict
    $\tau^{s}$), which \texttt{HeckeOperatorOnK} would assemble. \textbf{Not needed}, and so neither
    built nor specified. Line~2's Hecke action (slot-10 WP-10, WP-11) reads $T_p$ through
    $\widehat\Theta_g$ on the free model $\widehat C=C_{\mathrm{Iso}}(R)$, whose cohomology is
    $H^{\bullet}(\Ga;V_R)$ over every ring (\HA{prop:isomodel}(3)), and never through the Bredon object;
    the strict compatibility $\varepsilon^{\ast}\circ\Theta^{s}_g=\widehat\Theta_g\circ\varepsilon^{\ast}$
    of \HA{prop:strictlift}(3) is checked in slot~10 only on rings where a $\Theta^{s}_g$ is shipped
    ($\Fl$ at every $\ell$; $\Z/\ell^k$ at $\ell\ge5$), and over $\Z/\ell^k$ at $\ell\in\{2,3\}$ it
    follows from the exact $\Z$-certificate of $\varepsilon\circ\tau_{\mathrm{Iso}}=\tau^{s}\circ
    \varepsilon$, \HA{eq:cosettransport} being functorial in the ring ((b1), ``outside the
    closure'').
  \item \emph{The operator on cohomology $A^{(q)}(T_p)$ of the Bredon object.} \textbf{Declined}: it
    is \textbf{not shown to be independent of the choice of $\tau$}. Homotopy uniqueness of
    admissible $\tau$ needs $6\in R^{\times}$ (\HA{lem:tauexists}(2)), which fails at $\ell\in\{2,3\}$,
    and \HA{prop:strictlift}(4) compares $\Theta^{s}_g$ with the shipped $\Theta_g$ on cohomology only
    at $6\in R^{\times}$ --- at $\ell\in\{2,3\}$ ``they may differ, and no statement of [that] chapter
    reads them there''. This is the \textbf{absence of a proof of independence, not a proof of
    dependence}: [C] exhibits no two admissible $\tau$ with different induced operators on the Bredon
    object, and this document claims none. Lifting it would ship the canonical $\tau$'s operator
    with no statement of what it is an invariant of, so no consumer may read it and none is built.
  \end{enumerate}
  \emph{Consequences, binding from this pre-release.} \texttt{HeckeOperatorOnK}'s refusal at
  $\ell\in\{2,3\}$ (in \texttt{init\_composite()}, by \texttt{NotImplementedInThisBuild}, citing this
  item) and \texttt{HeckeReduction}'s stand, with their arms; so do the refusals of the boundary and
  cone classes of~(f) and of the cross-routes there ((b1)). The campaign's job loop keeps recording
  \texttt{hecke} (with \texttt{bockstein} and \texttt{reduction}) as refused at $\ell\in\{2,3\}$
  ((g)), and the console lifts nothing there ((h)). [X]~\S8's re-cut row ``\texttt{HeckeOperatorOnK}
  over $\Z/\ell^K$ at $\ell\in\{2,3\}$ (M1's need)'' is withdrawn (slot-10 WP-17, [X] v1.14). \emph{If
  a line-2 package exhibits two admissible $\tau$ whose Bredon-object operators differ} at $2$ or
  $3$, that is a finding for [C], recorded by slot-10 WP-15 in place of the ``not shown'' reading ---
  not assumed here. [C] v1.10.0 (slot-10 WP-15, decision~12~(A)) records the same reason in a remark
  beside \HA{prop:strictlift}(4), in this form: ``not shown independent of $\tau$; [C] compares
  nothing there''.
  \emph{v1.11.0-pre.2, as built.} Nothing was lifted: \texttt{HeckeOperatorOnK}'s refusal in
  \texttt{init\_composite()} and \texttt{HeckeReduction}'s stand with their arms and messages, and no
  code package of slot~10 edited either class. \textbf{No line-2 package exhibited two admissible
  $\tau$ with different operators} at $2$ or $3$ --- none computed an operator on the Bredon object
  over $\Z/\ell^k$, $k\ge2$, at all --- so the ``not shown'' reading stands. Item~(a)'s last sentence is
  completed by (b1)'s ``outside the closure'' as built (the functoriality argument's three inputs,
  review finding m-1). [C] v1.10.0 records the reason as \HA{rmk:bredonhecke}, in the form quoted,
  with one precision this item left implicit: at $k=1$ the shipped $\Fl$ pipeline computes (a) and
  (b) for the \emph{canonical} $\tau$ at $\ell\in\{2,3\}$ --- a definite object, the canonical
  datum's operator, which no statement identifies with anything $\tau$-independent there; the
  declined lift is the one over $\Z/\ell^k$, $k\ge2$. (\emph{v1.11.0}, the phase-1 referee's M-1 and
  m-5: ``has always computed'' is replaced by ``computes'' --- the canonical $\tau$ is
  \HA{alg:tau}'s since [C] v1.1.0 (\HA{rmk:iotatausym}), and its operator (b) did not move in that
  re-baseline and equals the legacy $\tau$'s at every $\ell\in\{2,3\}$ golden point,
  \HA{rmk:tauexistsscope}(1), an observation, not a theorem; [C] v1.10.1 qualifies ``compares
  nothing there'' to the operators it names, $\Theta^{s}_g$, the canonical $\Theta_g$ and
  $\Theta_{\mathrm{Iso},g}$, and points to that observation.)\par\endgroup
\item \emph{The seven user-facing surfaces} keep named, user-facing refusals of their own, each
  with a test arm (the \texttt{conway} precedent, slot-7 WP-4): \textbf{\texttt{HeckeExtDriver}}
  (at construction); \textbf{the campaign} (at spec validation, \texttt{Campaign/CampaignSpec}, for
  the job loop and the ext pass alike, before any job runs); \textbf{\texttt{--atlas}}
  (\texttt{HeckeAtlasCli}); \textbf{\texttt{--bench}}, every mode other than the ones WP-3
  supports; \textbf{the console's pipeline verbs} (one dispatch check in
  \texttt{Console/Bindings/}); \textbf{\texttt{HeckeReport}} (\texttt{hecke\_report/v2} never
  moves); and \textbf{the cone} (\texttt{EquivariantMappingCone}, \texttt{HeckeConeAction}).
  Every message carries the tag \texttt{not\_implemented\_in\_this\_build:}.
  \emph{v1.9.0, as built} (arm~(X), one member each): \texttt{HeckeExtDriver} at construction;
  the campaign at spec validation (\texttt{CampaignSpec}), before any job, for
  \texttt{Z/ell\textasciicircum2}, \texttt{Z/ell\textasciicircum3} and \texttt{Z} (refused by every consumer but the integral descent, from v1.23.0-pre.1: (a));
  \texttt{--atlas} with exit~$5$; \texttt{--bench} in every mode, the guard sitting at
  \texttt{bench::RunMode}'s entry before the mode dispatch (the arm drives \texttt{cohom}); the console at the dispatch of \texttt{symmetric\_power},
  \texttt{cohomology}, \texttt{hecke} and \texttt{family}, exit~$5$ --- also through a lazy handle,
  since slot-8 WP-10's \texttt{bind::LazyCompute} ([P] v3.6 \S8); \texttt{HeckeReport} at
  \texttt{StreamHeckeReportJson}'s entry, the emitted \texttt{hecke\_report/v2} text being
  byte-identical; and \texttt{EquivariantMappingCone} and \texttt{HeckeConeAction}.

  \begingroup\sloppy\emph{v1.10.0-pre.1: the surfaces in slot~9.} Four keep their refusal unchanged:
  \texttt{HeckeExtDriver}, \texttt{--atlas}, \texttt{--bench} (every mode; no switch is added)
  and \texttt{HeckeReport} (\texttt{hecke\_report/v2} never moves). Three change.
  \textbf{The campaign} (slot-9 WP-8): spec validation \emph{admits} a composite ring for the job
  loop, which runs under~(g). It \emph{still refuses at spec validation, by name and before any
  job file exists}, a spec with \texttt{ext = true} under a non-default ring, and a spec whose
  grid holds a cell with $\ell^k\ge2^{64}$ (naming the cell). So no refusal ever follows a job
  that has run. The refusal of \texttt{Z} inside \texttt{MathSettings::set()} is superseded (v1.23.0-pre.1): see (a), \HA{def:ringsemantics}(3).
  \textbf{The console} (slot-9 WP-12): the verbs and members of~(h) honour, and every other
  pipeline verb keeps exit~$5$, through a lazy handle too. \textbf{The cone}
  (slot-9 WP-4): \texttt{EquivariantMappingCone} and \texttt{HeckeConeAction} (route~1)
  \emph{honour} at $\ell\ge5$. So the cone is no longer a refusing surface there; it refuses at
  $\ell\in\{2,3\}$ and wherever $\ell^k\ge2^{64}$, and its two cross-routes refuse (b1).
  \emph{v1.10.0-pre.2, as built.} The console (slot-9 WP-12, WP-13) and the cone (slot-9 WP-4)
  are as stated, with the class width of (b1) (the cone classes refuse from $\ell^k\ge2^{32}$).
  \emph{v1.10.0, as built} (slot-9 WP-8): \texttt{CampaignSpec} admits $\Z/\ell^k$ for the job
  loop, and refuses, by name and before any job file or output directory exists
  (\texttt{--campaign} exits~$1$), a spec with \texttt{ext = true} under a non-default ring
  (``\ldots\ is not honoured in this build for the ext pass \ldots'') and a spec whose grid holds a
  cell with $\ell^k\ge2^{64}$ (naming \texttt{n\_sym}, $\ell$ and the section's line; $\ell=4294967311$
  at $k=2$ and $3^{41}$ are refused, $2^{63}$ is admitted). Arms: \texttt{CampaignRingValidation}~(X);
  \texttt{HeckeReductionValidation} arm~(X), whose \texttt{surface campaign (spec validation)} row
  now carries \texttt{ext = true}, the surface that still refuses.\par\endgroup
\end{enumerate}
Everything \S\ref{sec:field}, \S\ref{sec:alg} and the extension record compute is a field notion,
and is refused under a non-default ring in slot~8 (v1.10.0-pre.1: and in slot~9; the
\texttt{hecke\_ext} record stays $\Fl$-only).

\paragraph{(c) The cyclic-summand basis rule ([X] v1.11 E8 finding M3), stated once.}
$H^q=Z^q_{\mathrm{eq}}/B^q_{\mathrm{eq}}$ over $R=\Z/\ell^k$ is \texttt{invariant\_factors(ctx, Z,
B)} (\S13.6), and its cyclic summands get a basis by the following rule, a deterministic function
of the input matrices $(G,R_l)=(Z,B)$, as the D5b bases of \S8 are of theirs. It is the one
definition tier~B reads.
\begin{enumerate}[label=(c\arabic*)]
\item \emph{Presentation.} $\Gamma=[\gamma_1,\dots,\gamma_s]$ is the Howell generator matrix of
  $N=\operatorname{colspan}[R_l\mid G]$ (\S13.6 step~1), its columns in Howell-row order. The
  relation matrix $\mathrm{Rel}$ ($s\times t$; rows = generators, columns = relations) is
  $[\,K\mid X\,]$: first the rows of \texttt{kernel\_module}$(\Gamma)$, in Howell order, as
  columns; then the coordinates \texttt{solve\_multi\_rhs\_module}$(\Gamma,R_l)$ of the columns of
  $R_l$, in their order. So $M\cong R^s/\operatorname{colspan}\mathrm{Rel}$.
  \emph{The coordinates of $(G,R_l)$} (v1.9.0; slot-8 WP-3's choice, which the pre-release left
  coordinate-free; review finding m-9). On the pipeline route $(G,R_l)$ is read in
  \textbf{$C^q_{\mathrm{eq}}$ coordinates}. $C^q_{\mathrm{eq}}$ is
  \texttt{kernel\_module}$(A_{C^q})$, the kernel of the equivariance constraints on the flat
  cochains, with its Howell rows as the basis: that basis is unique (\S13.5), and every pivot is
  asserted to be a unit, so $C^q_{\mathrm{eq}}$ is free (\HA{prop:coneintegral}(1)--(2) at
  $\ell\ge5$; v1.9.1: through v1.9.0 this cited~(2) alone).
  $X_q$ is the matrix of $d^q$ from the basis of $C^q_{\mathrm{eq}}$ to that of
  $C^{q+1}_{\mathrm{eq}}$, by \texttt{solve\_multi\_rhs\_module}. Then $G$ is the Howell basis of
  $\ker X_q$ (the cocycles), as columns, and $R_l=X_{q-1}$ (the coboundaries), both in
  $C^q_{\mathrm{eq}}$ coordinates. So $(G,R_l)$, and with it the basis of~(c4), is a function of
  the cell and the ring alone. The flat cochain of $b_i$ is the basis of $C^q_{\mathrm{eq}}$
  applied to $b_i$'s coordinates. Tier~B (\texttt{red\_H\textasciicircum q}) reads this basis.
  \emph{v1.10.0-pre.1 (slot-9 plan audit B-1): this item --- the unit-pivot assertion, the
  freeness of $C^q_{\mathrm{eq}}$, and the $C^q_{\mathrm{eq}}$ coordinates of $(G,R_l)$ --- holds
  at $\ell\ge5$ only}, where \HA{prop:coneintegral}(1)--(2) and \HA{prop:ringinforce}(1) make
  $C^q_{\mathrm{eq}}$ free. At $\ell\in\{2,3\}$ it is not free (below), the unit-pivot assertion
  would fail, and coordinates in a non-free module are not unique; there the route is~(c1$'$).
\item[(c1$'$)] \begingroup\sloppy\emph{The flat-coordinate route at $\ell\in\{2,3\}$} (v1.10.0-pre.1; decision~4).
  At $\ell\in\{2,3\}$, $C^q_{\mathrm{eq}}(\Z/\ell^k)$ is in general \textbf{not free}: at odd $n$
  the central relation gives $2f=0$ on every cochain value (\HA{prop:ringinforce}(3)'s
  computation read over $\Z/2^k$), and at even $n$, $\ell=2$, already $n=2$ gives
  $(V_R)^{\langle S\rangle}=R(x^2+y^2)\oplus\{b\,xy:2b=0\}$ (\HA{prop:ringinforce}(2)'s
  computation read over $\Z/2^k$). So \texttt{EquivariantCohomologyOnK} computes there, in
  \textbf{flat cochain coordinates}, with the flat coboundary $d^q=D_q\otimes I$ and the
  equivariance constraint matrix $A_{C^q}$ of the pipeline's own relation list over the ring
  (unchanged):
  \[
    Z^q=\texttt{kernel\_module}\!\begin{bmatrix}A_{C^q}\\ d^q\end{bmatrix},\qquad
    B^q=d^{q-1}\,\Gamma_{q-1},\qquad
    H^q\cong\texttt{invariant\_factors}(\mathrm{ctx},Z^q,B^q),
  \]
  where $\Gamma_{q-1}$ is the matrix whose columns are the Howell generators of
  $C^{q-1}_{\mathrm{eq}}=\texttt{kernel\_module}(A_{C^{q-1}})$ (with $B^0=0$), and $Z^q$ is taken
  as the matrix whose columns are its Howell generators. Both are generator matrices of submodules of
  the flat cochains, $B^q\subseteq Z^q$, and \S13.6's \texttt{invariant\_factors} takes column
  spans, so no basis of $C^q_{\mathrm{eq}}$ is needed. It yields
  \texttt{invariant\_factors\_H\textasciicircum q} and
  $\texttt{lengths\_H\textasciicircum q}=\texttt{length}(\operatorname{colspan}[B^q\mid Z^q])-
  \texttt{length}(\operatorname{colspan}B^q)=\sum_ie_i$, asserted equal, for $q=0,1,2$, and
  \textbf{nothing else}. There is no stored basis of $C^q_{\mathrm{eq}}$ or of $H^q$, so (c2)--(c5)
  are \textbf{not applied} there, and no map on $H^q$ (no \texttt{red\_H\textasciicircum q}, no
  $A^{(q)}(T_p)$, no Bockstein) is defined (v1.11.0-pre.1: and slot~10 defines none --- the Hecke action on the Bredon object is declined, for the two reasons of~(b3)'s ``v1.11.0-pre.1'' item). There is also no $X_q$, so no $r_K$, no
  \texttt{ladic\_type}, no free rank and no \texttt{charpoly}: those are the $K$-mode, which
  \HA{prop:certifiede}(5) refuses there. The relation list over the ring is the pipeline's,
  unchanged, and its $2$- and $3$-torsion is what the tree model predicts (\HA{prop:bredon}(2):
  $H^0=V_R^{\Gamma}$, $H^1=V_R^{-I}/(V_R^S+V_R^U)$, $H^q=0$ for $q\ge2$; the pipeline's
  $H^2=0$ is checked by \texttt{HeckeBredonTreeRingValidation}, arms~(T) and~(A) --- v1.10.0-pre.2:
  through v1.10.0-pre.1 this read ``is asserted'', but the route itself asserts nothing about
  $H^2$, the slot's review, WP-14, finding m-3; v1.10.1: \texttt{docs/implementation/slot9/WP-14\_review.md}). \emph{As built} (slot-9 WP-3,
  \texttt{compute\_cohomology\_flat\_module()}): exactly this route, with three certificates
  re-derived from the dense $A_{C^q}$, $d^q$ and the two Howell lengths, independently of the
  invariant-factor result --- $A_{C^q}B^q=0$ and $d^qB^q=0$ entrywise ($B^q\subseteq Z^q$);
  $\sum_ie_i=\texttt{length}(Z^q)-\texttt{length}(B^q)$; and
  $\texttt{length}(C^q_{\mathrm{eq}})=\texttt{length}(Z^q)+\texttt{length}(B^{q+1})$. Measured on
  $96$ cells ($\ell\in\{2,3\}$, $k\in\{2,3,4\}$, even $n\le20$, odd $n\le9$) against O7, with O7
  controls at $k=1$ (the shipped $\Fl$ route, $32$ cells) and at $\ell=5$ (the $\ell\ge5$ module
  route, $5$ cells). The \texttt{uint64\_t} instantiation of the route is compiled and reached
  by no honouring cell (review finding m-2, open). \emph{v1.10.0: closed} (slot-9 WP-9,
  \texttt{fb8a63a}): \texttt{HeckeBredonU64FlatRouteValidation} (\texttt{fast}) runs one
  honouring \texttt{uint64\_t} cell, $\ell=2$, $n=4$, $k=33$, against the frozen O7 row at
  $(\ell,k,n)=(2,4,4)$ (lengths $(2,6,0)$, $H^0=[2]$, $H^1=[4,2]$, $H^2=[\,]$) read through the
  route's $K$-scaling: an exponent equal to $K$ is a free summand $\Z/\ell^K$ and moves with $K$,
  one below $K$ is $\ell$-power torsion and does not. It predicts and finds lengths $(2,35,0)$,
  $H^1=[33,2]$ at $K=33$. \textbf{The $\ell\ge5$ route is byte-identical}: the branch is selected
  by $\ell$ at the ring context, and no line of the $\ell\ge5$ module route moves (slot-9 plan
  M-2; its oracle is \texttt{GOLDEN\_STRICT=1} over the $124$ entries of~(e)).\par\endgroup
\item \emph{Smith elimination} of $\mathrm{Rel}$ by \S13.6 step~2's pivot rule on the trailing
  submatrix (smallest valuation anywhere, then leftmost column, then topmost row), moved to the
  diagonal by one row swap and one column swap. The pivot is made exactly $\ell^\sigma$ by scaling
  its \emph{column} by \texttt{inv(unit\_part)} of the pivot (a relation operation). Then the
  pivot's \emph{row} is cleared first (column operations) and its \emph{column} second (row
  operations), each by the \texttt{divide\_exact} quotient; there is no other normalisation. The
  elimination stops when every trailing entry is $0$.
\item \emph{Mirroring.} Column operations on $\mathrm{Rel}$ act on relations and leave the
  generators alone. Every row operation is mirrored on $\Gamma$ by the inverse column operation:
  $\mathrm{row}_i\leftarrow\mathrm{row}_i-q\,\mathrm{row}_j$ becomes
  $\gamma_j\leftarrow\gamma_j+q\,\gamma_i$, and a row swap swaps the two columns of $\Gamma$. (If
  $U\,\mathrm{Rel}\,V=D$, the transformed generators are the columns of $\Gamma U^{-1}$.)
\item \emph{Read-off.} Diagonal position $i$ with entry $\ell^\sigma$, $1\le\sigma<k$, gives a
  summand $\Z/\ell^\sigma$ generated by the transformed $\gamma_i$; a generator whose diagonal
  entry is $0$ or missing gives $\Z/\ell^k$; a unit diagonal entry gives nothing (its generator is
  $0$ in $M$). The summands are then ordered by exponent, \emph{descending}, ties by final
  generator index, ascending (a stable sort), so that the basis is indexed $b_1,b_2,\dots$ with
  $b_i$ generating the $\Z/\ell^{e_i}$ summand and $e_1\ge e_2\ge\cdots$. The identity
  $\sum_ie_i=\texttt{length}(N)-\texttt{length}(\operatorname{colspan}R_l)$ is asserted (\S13.6
  step~4).
\item \emph{Coordinates and maps.} A class $x\in M$ has coordinates $(c_i)$ with
  $x=\sum_ic_ib_i$, the $i$-th \emph{reduced} into $[0,\ell^{e_i})$; they are unique. A map
  between such modules (the Hecke matrix $A^{(q)}(T_p)$ on the cyclic summands, WP-3; a reduction
  map, (d)) is the matrix whose column $j$ holds the reduced coordinates of the image of $b_j$,
  entry $(i,j)$ in $[0,\ell^{e_i})$ of the target. Such a matrix is basis-dependent, and its
  characteristic polynomial is not an invariant of a non-free module ([L] \S2.6), which is why (d)
  defines \texttt{charpoly} on the free quotient only.
\end{enumerate}
The rule fixes the two choices \S13.6 leaves open --- the order of the relation columns and the
side on which the pivot is scaled --- and nothing else. It is not used at $k=1$: the $\Fl$ path
keeps the D5b bases. \emph{v1.9.0, how the rule is tested} (slot-8 WP-3r, on review finding M-1,
which found that no arm tested the basis rule directly). Arm~(M) of \texttt{HeckeReductionValidation}
runs over $\Z/4$, $\Z/25$, $\Z/125$ and $\Z/691^2$ (\texttt{u32}), and over $\Z/134217757^2$
(Montgomery) and $\Z/4294967291^2$ (U64\_Wide). Its cases are fixed-seed known-answer and random
$(G,R_l)$. It asserts: the exponents against \texttt{invariant\_factors} and the naive oracle; the
coordinates of the basis are the identity; each $b_i$ has order exactly $\ell^{e_i}$; the basis
generates; and the coordinates of~(c5) are unique and linear. It also rejects three corrupted
decompositions, the ones a wrong mirror on $U$, a wrong mirror on $\Gamma$ or a wrong read-off
position would produce. Each of those three mutations of the code fails the validation hundreds of
times. On every pipeline cell, the coordinates of the cells' own basis are asserted to be the
identity, for $q=0,1,2$ (arm~(C)).

\paragraph{(d) The E8 objects: names and meaning.}
They are for the golden harness and the validations, not for a record (decision~4). Throughout,
$C$ is the integral equivariant cochain complex of the model; for $\ell\ge5$,
$C_{\Z/\ell^k}=C\otimes\Z/\ell^k$ (\HA{prop:coneintegral}~(2)). On the \emph{applicability range}
$n>0$ even, $\ell\ge5$ (\HA{cl:bockstein}~(1), \HA{cl:amalgam}~(3), \HA{lem:h2torsion}~(1)),
$H^1(C)_{(\ell)}\cong\Z_{(\ell)}^{r}\oplus\bigoplus_{i=1}^{t}\Z/\ell^{e_i}$ and
$H^0(C)_{(\ell)}=H^2(C)_{(\ell)}=0$, with $r=\dim M_{n+2}+\dim S_{n+2}$ and $t=\dim H^0(\Fl)$ the
number of Dickson classes; by universal coefficients
$H^0(\Z/\ell^k)\cong\bigoplus_i\Z/\ell^{\min(e_i,k)}$,
$H^1(\Z/\ell^k)\cong(\Z/\ell^k)^r\oplus\bigoplus_i\Z/\ell^{\min(e_i,k)}$ and $H^2(\Z/\ell^k)=0$
([L] \S2.1, \S2.4). The level of a run is $K:=k$, the exponent of the ring in force; slot~8 runs
no other $K$. The $K$-mode objects (\texttt{r\_K}, \texttt{ladic\_type*}, the free-quotient
\texttt{charpoly}) and the assertion of the Bockstein meaning are defined on the applicability
range only, and are absent elsewhere; every E8 point lies in it.
\begin{itemize}
\item \textbf{\texttt{lengths\_H\textasciicircum q}} ($q=0,1,2$; tier~A): the composition length
  of $H^q$ over $\Z/\ell^k$, $\sum_ie_i$ over its invariant factors, so that
  $|H^q|=\ell^{\text{length}}$. Over $\Z/\ell^k$ it replaces \texttt{dim\_H\textasciicircum q}.
  \emph{v1.10.0-pre.1}: this object and the next are the \emph{only} objects of this list defined
  at $\ell\in\{2,3\}$. There they are the Bredon object's (\HA{def:ringsemantics}(2)), computed
  by~(c1$'$); no universal-coefficient statement of this paragraph's preamble applies to them, and
  every other object below is undefined at $\ell\in\{2,3\}$. No golden point lies at
  $\ell\in\{2,3\}$ in slot~9 (decision~4~(A); (A$'$) was not ruled).
\item \textbf{\texttt{invariant\_factors\_H\textasciicircum q}} ($q=0,1,2$; tier~A): the list
  $e_1\ge e_2\ge\cdots\ge1$ with $H^q\cong\bigoplus_i\Z/\ell^{e_i}$, from
  \texttt{invariant\_factors(ctx, Z, B)} (\S13.6); $e_i=k$ is a summand that is free over
  $\Z/\ell^k$. The list is empty for the zero module.
\item \textbf{\texttt{r\_K(d\textasciicircum0)}}, \textbf{\texttt{r\_K(d\textasciicircum1)}}
  (tier~A) and \textbf{\texttt{ladic\_type\_certified}} (tier~A, boolean). $r_K(d)$ is the number
  of non-zero diagonal entries, i.e.\ entries $\ell^s$ with $s<K$, in the Smith form of $d$ over
  $\Z/\ell^K$ (by (c2)'s pivot rule). The $K$-mode is \textbf{certified} iff
  $r_K(d^0)=\dim C^0_{\mathrm{eq}}$ and $r_K(d^1)=\dim C^2_{\mathrm{eq}}$ (E8-P2; [L] \S2.4). Then
  every elementary divisor of $d^0$ and $d^1$ over $\Z_{(\ell)}$ is $<\ell^K$, and the $e_i$ are
  read off exactly. Both counts are emitted whether or not the mode certifies. An uncertified
  level is a loud outcome, never an error. Its state is R3-shaped (kind \texttt{ladic\_level},
  budget $\{K\}$, reached $\{r_K(d^0),r_K(d^1)\}$, \texttt{complete = false}), and its continuation
  is a \emph{fresh} run at a larger $k$, not a \texttt{resume}: \texttt{coefficient\_ring} is a
  non-budget setting, and a resume under a changed non-budget setting is refused
  (\S\ref{sec:schema}). \emph{v1.9.0, as built} (slot-8 WP-3, in \texttt{EquivariantCohomologyOnK}):
  $r_K(d^q)$ is the number of invariant factors of $X_q$ (the matrix of~(c1)), and the
  \texttt{ladic\_level} state is carried by the object. Three things are asserted on every cell
  by arm~(C) of \texttt{HeckeReductionValidation}: \texttt{ladic\_type\textasciicircum q} is checked
  against \texttt{invariant\_factors\_H\textasciicircum q} by universal coefficients; the cells
  whose torsion exponents are all $<2$ certify at $K=2$; and the two controls are uncertified at $K=2$, each with a
  deficit of exactly one, and certified at $K=3$. The deficit at $K=2$ is $r_2(d^0)=\dim
  C^0_{\mathrm{eq}}-1$, with $r_2(d^1)=\dim C^2_{\mathrm{eq}}$.
\item \textbf{\texttt{ladic\_type\textasciicircum q}} ($q=0,1,2$; tier~A; present \textbf{only}
  where \texttt{ladic\_type\_certified = true}): the $\Z_{(\ell)}$-type
  $(r_q;e^q_1\ge e^q_2\ge\cdots)$ of $H^q(C)_{(\ell)}$, its free rank and torsion exponents. On the
  applicability range it is $(0;)$ at $q=0,2$ and $(r;e_1,\dots,e_t)$ at $q=1$. Write
  $e_{\max}:=\max_ie_i$, and $e_{\max}:=0$ when $t=0$.
\item \textbf{\texttt{charpoly\_A(T\_p)\textasciicircum1\_free\_p<p>}} (tier~A; E8-P4, [L] F-C;
  slot-8 plan \S2.3 item~13), \textbf{with its precondition}. It is defined \emph{only} at a level
  $K$ where \texttt{ladic\_type\_certified = true}, as the characteristic polynomial of
  $A^{(1)}(T_p)$ on the free quotient
  \[
    H^1(\Z/\ell^K)\,/\,H^1(\Z/\ell^K)[\ell^{e_{\max}}]\;\cong\;(H^1(C)/\mathrm{tors})\otimes\Z/\ell^{K-e_{\max}},
  \]
  a free $\Z/\ell^{K-e_{\max}}$-module of rank $r$, stable under $T_p$ because the
  $\ell^{e_{\max}}$-torsion is characteristic ([L] \S2.6). \textbf{The precision $K-e_{\max}$ is
  part of the object}: the object is the exponent $K-e_{\max}$ together with the monic polynomial
  of degree $r$ over $\Z/\ell^{K-e_{\max}}$, coefficients low$\to$high in $[0,\ell^{K-e_{\max}})$.
  It is basis-independent, hence tier~A. At an uncertified level it has no definition ([L] \S2.4:
  a summand $\Z/\ell^e$ with $e\ge K$ reads as free), so it is \textbf{absent}, and its accessor
  throws a named \texttt{std::logic\_error} and never returns a value. It exists for $q=1$ only: on
  the applicability range $H^0$ and $H^2$ have free rank $0$. \emph{v1.9.0, as built} (slot-8
  WP-3, \texttt{free\_charpoly\_H1}). It is computed by Berkowitz over $\Z/\ell^K$ on the free
  block, then reduced to $\Z/\ell^{K-e_{\max}}$. At weight $12$ it is \emph{asserted}, not
  printed (slot-8 WP-3r, review finding m-7): at every certified torsion-free cell with $n=10$ and
  every prime of the cell, the value must be exactly $(x-(1+p^{11}))(x-\tau(p))^2$ over
  $\Z/\ell^{K-e_{\max}}$. That is $(10,691)$ over $\Z/691^2$ at $p\in\{2,3,5\}$, and
  $(10,134217757)$ and $(10,4294967291)$ over $\Z/\ell^2$ at $p=2$ (arm~(T)). Wherever $H^1$ is
  free, its reduction mod~$\ell$ equals the $\Fl$ characteristic polynomial. Elsewhere it divides
  that polynomial.
\item \textbf{\texttt{bockstein\_rank\textasciicircum0}} (tier~A): the rank over $\Fl$ of the
  mod-$\ell$ Bockstein $\beta^0:H^0(\Fl)\to H^1(\Fl)$, the connecting map of
  $0\to\Z/\ell\to\Z/\ell^2\to\Z/\ell\to0$ (\HA{cl:bockstein}), computed from the stored bases (E8
  item~4a). \textbf{\texttt{higher\_torsion\_count}} (tier~A) is
  $\dim H^0(\Fl)-\texttt{bockstein\_rank\textasciicircum0}$. Hence the \textbf{bookkeeping identity}
  \[
    \texttt{bockstein\_rank\textasciicircum0}+\texttt{higher\_torsion\_count}=\dim H^0(\Fl),
  \]
  true by definition. The content is its \emph{meaning} ([L] \S2.4): $\beta^0$ is the injective
  integral Bockstein $H^0(\Fl)\cong H^1(C)[\ell]$ followed by reduction mod~$\ell$, and a summand
  $\Z/\ell^e$ lies in its kernel iff $e\ge2$, so
  \[
    \texttt{bockstein\_rank\textasciicircum0}=\#\{i:e_i=1\},\qquad
    \texttt{higher\_torsion\_count}=\#\{i:e_i\ge2\},
  \]
  and the second equality is \textbf{asserted wherever the $K$-mode certifies}, against the $e_i$
  of \texttt{ladic\_type\textasciicircum1}. It is \emph{not} the equality
  $\texttt{bockstein\_rank\textasciicircum0}=\dim H^0(\Fl)$: that one means ``no
  $\ell^2$-torsion'', and it fails by design at the two positive controls $(30,5)$ and $(60,5)$.
  Both objects are mod-$\ell$ objects, and so are the same at every $k$ of a cell.
  \emph{v1.9.0.} The meaning is now a statement of [C]: \HA{rmk:bocksteinexp}(1), proved in [C] since
  v1.8.0 (slot-8 WP-9; finding F-A of [L], owner ruling~10), is exactly
  $\rank\beta^0=\#\{i:e_i=1\}$. Its clause~(2) explains why the $K$-mode needs its certificate.
  A class of $H^0(M_m)$ produces torsion of exponent $\ge j+1$ at $n=\ell^jm$, so no fixed $k$
  reads every exponent correctly (\HA{prop:certifiede}). \emph{As built} (slot-8 WP-4,
  \texttt{Hecke/HeckeReduction}), $\beta^0$ is computed from the stored bases. A D5b class is
  lifted into $C^0_{\mathrm{eq}}(\Z/\ell^2)$ through the free Howell basis; $X_0c=\ell w$ gives
  $\beta^0=[w]$. Arm~(B3) asserts the identity, and asserts the two counts wherever the $K$-mode
  certifies; arm~(B2) asserts the long exact sequence's length bookkeeping; arm~(B1) asserts
  $\beta\circ A(T_p)=A(T_p)\circ\beta$. At every E8 point \texttt{higher\_torsion\_count} equals
  its pre-stated value in~(e).
\item \textbf{\texttt{red\_H\textasciicircum q}} ($q=0,1,2$; \textbf{tier~B}; [X] v1.11 E8 finding
  F3): the reduction map $H^q(\Z/\ell^k)\to H^q(\Fl)$ induced by $\Z/\ell^k\to\Fl$, as the matrix
  of (c5) from the cyclic-summand basis of (c) to the D5b basis of the $\Fl$ run at the same
  $(n,\ell)$, with entries in $[0,\ell)$. Reductions compare as maps, never as bases. On
  cohomology the reduction is injective after $\otimes\Fl$, with a Tor cokernel (universal
  coefficients), so ``the $\Z/\ell^k$ object reduces entrywise to the $\Fl$ object'' is false even
  up to a change of basis. The certificates are the map equalities
  $A(T_p)\circ\mathrm{red}=\mathrm{red}\circ A(T_p)$ and $\beta\circ A(T_p)=A(T_p)\circ\beta$
  (item~4a). Cochain-level matrices reduce entrywise and are not objects. \emph{v1.9.0, as built}
  (slot-8 WP-4, \texttt{Hecke/HeckeReduction}). The map is
  $\mathrm{red}=P_q(b_j\bmod\ell)$, with $P_q$ the $\Fl$ run's projector to the D5b basis. Five
  arms of \texttt{HeckeReductionValidation} assert it at every E8 cell-level and every prime:
  (R1) it is the class map; (R2) the map equality above, for $q=0,1,2$; (R3) $\Theta$ reduces
  entrywise; (R4) at $K=2$ it is exact against $\beta$; (R5) functoriality from $K=3$ to $K=2$.
  At the E8 cells its Tor cokernel is non-zero exactly where a summand $\Z/\ell^e$ of $H^0$ has
  $e<K$, and it is $0$ on $H^1$ and $H^2$.
\end{itemize}
The Hecke matrix $A^{(q)}(T_p)$ over $\Z/\ell^k$ on the cyclic summands (c5) is computed and
validated, but is \textbf{not} an E8 golden object in slot~8. Nor is \ref{D6}'s refinement at
$(10,691)$, which is a validation arm and a report line (E8 item~4a). At $p=\ell$, $T_p$ is
$U_\ell$ (the harness builds its cosets with $p=\ell$). Over $\Z/\ell^k$ that is item~3's newest
route (E8-P5's $\operatorname{diag}(\ell,1)$ split), validated in slot~8 and not frozen.
\emph{v1.9.0: what the \ref{D6} arm measured} (slot-8 WP-4, \texttt{HeckeReductionValidation}).
At $(10,691)$ over $\Z/691^2$ the free characteristic polynomial of $T_2$ is
$[252219,379705,475480,1]=(x-2049)(x+24)^2$, which is asserted. So the pipeline's normalisation is
the classical one: $a_E=1+2^{11}$ and $a_\Delta=\tau(2)$, with $a_E\equiv a_\Delta\pmod{691}$ and
$a_E\not\equiv a_\Delta\pmod{691^2}$. $\ker(A-a_E)$ contains a free Eisenstein eigenline of
rank~$1$, and $\ker(A-a_\Delta)$ a free cuspidal eigenplane of rank~$2$; the arm asserts this
$1+2$ split. The free eigenline and the free eigenplane together span a submodule of index~$691$
in $H^1$: the determinant of their three generators is $415291=601\cdot691$, of $691$-adic
valuation~$1$. That index is recorded, not asserted. It is the Ramanujan congruence seen
integrally. The $U_5$ route is exercised by arms (P) and (T) at
every $\ell=5$ cell: $115$ distinct $g$ per cell against the oracle, the split route among them.

\paragraph{(e) The E8 golden points (slot-8 decision~5).}
Every E8 object is emitted at a \textbf{new point} named
\[
  \texttt{E8 ring=Z/<ell>\textasciicircum<k> n=<n> ell=<ell>}
\]
with $\ell$, $k$ and $n$ in decimal and the ring written with the cell's prime, for example
\texttt{E8 ring=Z/691\textasciicircum2 n=10 ell=691} (the \texttt{S6 sl2z\_restrict} naming
precedent). This holds at $(10,691)$ and $(16,5)$ too, which are also grid points under their
\texttt{n=\ldots\ ell=\ldots} names. No existing point gains or loses an entry, and no E8 marker is
emitted anywhere else. The E8 point list is its own ($(6,5)$, $(30,5)$ and $(60,5)$ are not grid
points and do not join the grid); it is looped from the full golden run only, and the ring is
installed (\texttt{ScopedMathSettings}, source \texttt{spec}) around the E8 computations only. The
seven points and their primes: $(10,691)$ over $\Z/691^2$ with $p\in\{2,3,5\}$; $(6,5)$ and
$(16,5)$ over $\Z/25$, and $(30,5)$ and $(60,5)$ over $\Z/25$ and over $\Z/125$, each with
$p\in\{2,3\}$ (at $\ell=5$, $p=5$ is $U_5$, not frozen, (d)).

\emph{Predicted values} ([L] \S2.5's exact probe; the invariant factors follow from it by
universal coefficients; slot-8 plan \S2.3 item~13 and decision~2). A prediction that fails is a
finding for the owner, and is never re-predicted.
\begin{center}\small
\begin{tabular}{@{}lcccccccc@{}}
\toprule
cell & $K$ & cert. & $r$ & $e$ ($e_{\max}$) & \texttt{charpoly} over & $H^0$ & $H^1$ & $\beta$-rank / h.t.c. \\
\midrule
$(10,691)$ & 2 & yes & 3 & --- (0) & $\Z/691^2$ & $()$ & $(2,2,2)$ & $0$ / $0$ \\
$(16,5)$ & 2 & yes & 3 & --- (0) & $\Z/25$ & $()$ & $(2,2,2)$ & $0$ / $0$ \\
$(6,5)$ & 2 & yes & 1 & $(1)$ (1) & $\Z/5$ & $(1)$ & $(2,1)$ & $1$ / $0$ \\
$(30,5)$ & 2 & no$^{a}$ & 5 & --- & absent & $(2)$ & $(2^6)$ & $0$ / $1$ \\
$(30,5)$ & 3 & yes & 5 & $(2)$ (2) & $\Z/5$ & $(2)$ & $(3^5,2)$ & $0$ / $1$ \\
$(60,5)$ & 2 & no$^{b}$ & 9 & --- & absent & $(2,1)$ & $(2^{10},1)$ & $1$ / $1$ \\
$(60,5)$ & 3 & yes & 9 & $(2,1)$ (2) & $\Z/5$ & $(2,1)$ & $(3^9,2,1)$ & $1$ / $1$ \\
\bottomrule
\end{tabular}
\end{center}
\noindent$^{a}$~$r_2(d^0)=\dim C^0_{\mathrm{eq}}-1$: the $\Z/25$ summand reads as free.
$^{b}$~$r_2(d^0)=\dim C^0_{\mathrm{eq}}-1$, as at~$a$: a deficit of \textbf{exactly one}, the
$\Z/25$ summand reading as free (v1.9.0, restated model-independently; slot-8 WP-3's finding,
confirmed by review finding m-8). Through v1.9.0-pre.1 this footnote read
``$r_2(d^0)=51\ne52=\dim C^0_{\mathrm{eq}}$ ([L] \S2.5)''. Those figures are [L]'s, in the
coordinates of the amalgam model $V^S\oplus V^U$. They are not the shipped seven-vertex model's,
where $\dim C^0_{\mathrm{eq}}=235$ and $r_2(d^0)=234$ (and at~$a$, $119$ and $118$). The
invariant content, a deficit of one, is model-free. It is what the arm asserts
(\texttt{HeckeReductionValidation}, arm~(C), which asserts $\operatorname{rank}-1$, not a literal),
and it is what golden freezes (\texttt{r\_K(d\textasciicircum0)} at the two uncertified points). In both, $r_2(d^1)=\dim
C^2_{\mathrm{eq}}$. The $H^0$ and $H^1$ columns are
\texttt{invariant\_factors\_H\textasciicircum q} at $\Z/\ell^K$, exponents with multiplicity;
$H^2=0$ throughout. $r=\dim M_{n+2}+\dim S_{n+2}$; $e$ and $e_{\max}$ are those of
\texttt{ladic\_type\textasciicircum1}, shown only where certified. At an uncertified level the
$\ell^2$-torsion reads as free, so \texttt{invariant\_factors} is correct for $\Z/\ell^K$ but does
not determine the $\Z_{(\ell)}$-type, exactly as the certificate says.

\emph{The object list at each point.} Tier~A: \texttt{lengths\_H\textasciicircum q},
\texttt{invariant\_factors\_H\textasciicircum q}, \texttt{r\_K(d\textasciicircum0)},
\texttt{r\_K(d\textasciicircum1)}, \texttt{ladic\_type\_certified},
\texttt{bockstein\_rank\textasciicircum0} and \texttt{higher\_torsion\_count}, and, at a certified
point only, \texttt{ladic\_type\textasciicircum q} and
\texttt{charpoly\_A(T\_p)\textasciicircum1\_free\_p<p>} for each of the point's primes. Tier~B:
\texttt{red\_H\textasciicircum q}. That list, with the prime restriction and the \texttt{charpoly}
precondition, is the generating rule of WP-4's pre-stated set. The set is WP-4's to generate, and
the event is the lead's to land after its gate (slot-8 plan \S4.4).

\emph{v1.9.0: the landed values} (R1, merge \texttt{53a1909}; \S3 of the re-baseline note named
in the status paragraph above). Every value in the table
above is the frozen one, line for line. The generated set was $124$ entries ($103$ tier~A, $21$
tier~B): $2$ uncertified points $\times14$, plus $5$ certified points $\times17$, plus $11$
\texttt{charpoly} objects. \texttt{golden\_ref\_diff.py --expect} held it as a set equality, and
\texttt{golden\_explain.py} found no text or decomposition moved. The emission is byte-identical
across three runs, one of them under \texttt{ENGINE\_THREADS=1}. In the shipped model the counts
are as follows (point order as in the table).
\begin{itemize}
\item $\dim C^0_{\mathrm{eq}}=41,65,27,119,119,235,235$ and
  $r_K(d^0)=41,65,27,118,119,234,235$.
\item $r_K(d^1)=\dim C^2_{\mathrm{eq}}=55,85,35,155,155,305,305$.
\item $\texttt{lengths\_H\textasciicircum1}=6,6,3,12,17,21,30$ and
  $\texttt{lengths\_H\textasciicircum0}=0,0,1,2,2,3,3$.
\end{itemize}
The free characteristic polynomials, coefficients low$\to$high, are listed below. At $(10,691)$
over $\Z/691^2$ they are $[252219,379705,475480,1]$ ($p=2$), $[323249,57149,299829,1]$ ($p=3$)
and $[82572,23641,342757,1]$ ($p=5$). At $(16,5)$ over $\Z/25$ they are $[18,21,8,1]$ and
$[16,4,4,1]$. At $(6,5)$ over $\Z/5$ they are $[1,1]$ and $[2,1]$. At $(30,5)$, $k=3$, over
$\Z/5$ they are $[1,1,3,3,1,1]$ and $[2,1,4,2,2,1]$. At $(60,5)$, $k=3$, over $\Z/5$ they are
$[2,1,3,4,2,1,3,4,2,1]$ and $[1,1,1,1,1,1,1,1,1,1]$. Where two are given, the first is for $p=2$
and the second for $p=3$. Measured by
\texttt{GOLDEN\_STRICT=1 GoldenRegressionValidation} (tier~A $11\,575$ / $0$, tier~B $5866$ /
$0$). At $(6,5)$, and at $(30,5)$ and $(60,5)$ with $k=3$, the precision $K-e_{\max}$ is~$1$. So
those six polynomials carry nothing beyond their reduction mod~$\ell$, as~(d) says.

\begingroup\sloppy % v1.10.0-pre.1: long \texttt names below
\paragraph{(f) E8 item~4b: the boundary and the cone over $\Z/\ell^k$ (v1.10.0-pre.1).}
Slot-9 decision~3~(A), with its sub-item (iii)~(A); built by slot-9 WP-4. \emph{Scope}: $\ell\ge5$
and $\ell^k<2^{64}$, as for the bulk. At $\ell\in\{2,3\}$ every member below refuses by name
((b1), (b3)). \emph{Notation}: $C^{\bullet}_{\partial}$ is the circle model of the boundary
(\HA{def:pair}, \HA{prop:circlemodel}), the equivariant complex of $K_{\mathrm{bd}}$ in degrees
$0,1$, built by \texttt{SL2ZBoundaryPair}; $r=j^{\ast}$ is the restriction; and $\mathrm{Cone}(r)$
is \HA{eq:conediff}. \textbf{Everything about the cone is indexed by the degree of $H^{\ast}_c$},
as \texttt{EquivariantMappingCone} and the $\Fl$ golden objects \texttt{dim\_Hc\textasciicircum q}
are. The stored cochain group in degree $q$ is
$Cc^q:=\mathrm{Cone}(r)^{q-1}=C^q_{\mathrm{eq}}(K)\oplus C^{q-1}_{\mathrm{eq}}(K_{\mathrm{bd}})$, and
$H^q_c=H^q(Cc^{\bullet})=H^{q-1}(\mathrm{Cone}(r))$ for $q\in\{0,1,2\}$.
\begin{itemize}
\item \emph{The mathematics used.} (i)~The terms of $C^{\bullet}_{\partial}$ and of
  $\mathrm{Cone}(r)$ over $\Z$ are free (\HA{prop:coneintegral}(1)). The reduction
  $\mathrm{Cone}(r)_{\Z}\otimes\Z/\ell^k\to\mathrm{Cone}(r)_{\Z/\ell^k}$ is an isomorphism at
  $\ell\ge5$ (\HA{prop:coneintegral}(2)); so is the boundary's, whose stabilisers are $\{\pm I\}$
  (the same averaging argument, needing only $2\in R^{\times}$). So both $\Z/\ell^k$ complexes are
  reductions of integral complexes of free modules, and universal coefficients apply to them.
  (ii)~The long exact sequence \HA{eq:coneles} is exact over every commutative ring
  (\HA{lem:conerelative}(2)). $H^0_c=0$ and $H^q_c=0$ for $q\ge3$ hold over every ring
  (\HA{lem:conerelative}(4)), and $H^2(C(\Z/\ell^k))=0$ holds at $\ell\ge5$ (\HA{cl:amalgam}). (iii)~The
  reductions and the Bocksteins are equivariant on the nose (\HA{prop:coneintegral}(3)).
\item \emph{The route.} A \textbf{dense module route}, selected by \texttt{is\_field()} alone
  (\texttt{ENGINE\_API.md} \S13.6--\S13.7: there is no block or tape Howell form). The tape route
  of the cone stays the $\Fl$ path, byte-identical (\texttt{cone\_output\_diff.sh}). The route
  inherits the ring context's element type and Montgomery regime and uses \S13.6's exact
  \texttt{solve\_multi\_rhs\_module}; no new regime and no new bound.
\item \emph{Coordinates, and the basis rule.} (c1)--(c5) apply \textbf{verbatim} to both complexes,
  with $C^q_{\mathrm{eq}}$ read as follows.
  \begin{itemize}
  \item \emph{Boundary}: $C^q_{\partial,\mathrm{eq}}=\texttt{kernel\_module}(A_{C^q_{\partial}})$.
    It is free at $\ell\ge5$, and its unit pivots are asserted.
  \item \emph{Cone}: $Cc^q_{\mathrm{eq}}=C^q_{\mathrm{eq}}(K)\oplus C^{q-1}_{\mathrm{eq}}(K_{\mathrm{bd}})$
    in the class's flat layout, bulk rows first. Its Howell basis is the concatenation of the two
    Howell bases, bulk first. Two row modules on disjoint column blocks have, as their sum, the
    concatenation of their reduced Howell forms, and \S13.5's form is unique. This is asserted
    once per degree.
  \item $X_q$ is the complex's own differential in those coordinates. For the cone that is
    $D(c,c_{\partial})=(-dc,\,d_{\partial}c_{\partial}+rc)$ of \HA{eq:conediff}, the class's
    \texttt{D\_block}.
  \end{itemize}
  So each of $H^q(\partial)$ and $H^q_c$ over $\Z/\ell^k$ has the cyclic-summand basis of~(c4).
  The three map families $\pi^q_{\ast}$, $r^q$, $\iota^q_{\ast}$ of \HA{eq:coneles}, $T_p^{\partial}$
  and $A_c^{(q)}(T_p)$ are (c5) matrices between those bases.
\item \emph{Certificates} (validation and object level; none of these is a golden object).
  \begin{itemize}
  \item \HA{obl:coneexact}(i)--(iii) over the ring, as over $\Fl$.
  \item \textbf{(iv)'s non-field half}. At every spot of \HA{eq:coneles}, the composite of
    consecutive maps is zero. Each Howell generator of $\ker$ of the outgoing map is exhibited as
    the image of a vector under the incoming map, modulo the target's relations, by
    \texttt{solve\_multi\_rhs\_module}. The solve's residual is re-derived: the image minus the
    generator lies in the colspan of the relations. At every bulk spot,
    $\operatorname{im}\pi^q_{\ast}=\ker r^q$ is checked with \emph{both} inclusions by solves.
    A failure aborts. The E5 certificate's item numbering is kept: the new item is appended,
    never renumbered. \emph{v1.10.0-pre.2, as built} (slot-9 WP-4, \texttt{ConeModuleCertificate}
    in \texttt{EquivariantMappingConeModule.cpp}): \textbf{item~14}. Every composite is solved
    into the target's relations; the Howell generators of $\ker$ of the outgoing map are the
    first coordinates of \texttt{kernel\_module}$([\mathrm{out}\mid D_N])$, each exhibited as
    $\mathrm{in}(w_1)+D_Mw_2$ by \texttt{solve\_multi\_rhs\_module}, and the witnesses are stored and
    re-derived by \texttt{verify()}. \emph{Why the check is complete} (review finding m-W4-2):
    \texttt{verify()} re-derives each stored witness, not that the witnesses generate the kernel;
    completeness comes from item~11's length identity at every term together with the
    composite-zero check --- $\operatorname{im}\mathrm{in}\subseteq\ker\mathrm{out}$ and
    $\ell(\operatorname{im}\mathrm{in})=\ell(M)-\ell(\operatorname{im}\mathrm{out})
    =\ell(\ker\mathrm{out})$ force equality in a finite module --- so exactness is certified even if
    a witness list were short, and the witnesses are a second discharge. Two negative controls in
    arm (4b-C): a tampered kernel generator fails item~14 (since the review repair
    \texttt{2933b65} it tampers the generator $x$, compared exactly, not a witness, so it cannot
    false-alarm), and a wrong stored image length fails item~11.
  \item (v) over the ring: every certificate item that states a dimension reads
    \texttt{length} (\S13.6), and every rank of a map reads the length of its image. So the first
    identity of \HA{eq:conedims} reads
    $\ell(H^1_c)=\ell(H^0(\partial))-\ell(H^0)+\ell(\ker r^1)$, which holds unconditionally, and the
    second reads $\ell(H^2_c)=\ell(H^1(\partial))-\ell(\operatorname{im}r^1)$, which holds under
    $H^2=0$ and is therefore asserted at $\ell\ge5$.
  \item (vi): \textbf{the commutator is asserted; the $\nu=2$ recursion is not computed over the
    ring} (v1.10.0-pre.2, narrowed to what is built; through v1.10.0-pre.1 this item read
    ``(vi) is asserted where two primes are computed ($6\in R^{\times}$)'', which nothing
    implemented --- the slot's review, WP-14, finding M-W4-1; v1.10.1: \texttt{docs/implementation/slot9/WP-14\_review.md}). As built (review repair
    \texttt{2933b65}): at every level of every 4b cell and for every pair of computed primes,
    $[A_c(T_p),A_c(T_{p'})]=0$ for $q\le2$ and $[A^{\partial}(T_p),A^{\partial}(T_{p'})]=0$ for
    $q\le1$, over $\Z/\ell^k$ in (c5) form (products over the ring, row $i$ compared modulo
    $\ell^{e_i}$ of the target); arm (4b-H) of \texttt{HeckeReductionValidation}. [C]
    \HA{obl:coneexact}(vi) is therefore \emph{partial} over the ring, as [C] v1.9.0 records; the
    recursion over the ring is a later slot's.
  \item (viii)'s two arms stay $\Fl$-only in slot~9.
  \item The map equalities, as~(d) states them for the bulk:
    $A^{\partial}(T_p)\circ\mathrm{red}=\mathrm{red}\circ A^{\partial}(T_p)$ and
    $A_c(T_p)\circ\mathrm{red}=\mathrm{red}\circ A_c(T_p)$; every map of \HA{eq:coneles} commutes
    with $\mathrm{red}$; and $\beta\circ A=A\circ\beta$ on both complexes.
  \item The cone's mod-$\ell$ Bockstein $\beta^1_c\colon H^1_c(\Fl)\to H^2_c(\Fl)$ is computed and
    validated like the boundary's, but it is \textbf{not} an object in slot~9.
  \end{itemize}
  \emph{Pointer into [C]}: \HA{obl:coneexact}(iv)'s ``until it exists the non-field case is
  \textbf{recorded} as unchecked'' is stale from slot-9 WP-4's merge until [C]'s next edit
  (slot-9 WP-15). \emph{v1.10.0-pre.2}: closed --- [C] v1.9.0 restates the clause and records the
  non-field half as checked over $\Z/\ell^k$, $\ell\ge5$.
\item \emph{\HA{obl:coneexact}(vii) over the ring} (decision~3~(iii)~(A), slot-9 plan M-3).
  \textbf{Route~1 only} (\texttt{HeckeConeAction}) honours. The strict-pair route, the relative
  route and \texttt{EquivariantCohomologyOnK}'s relative sub-complex constructor refuse, each with
  an arm ((b1)). The entrywise agreement of the three routes stays an $\Fl$ statement, and
  \textbf{(vii)'s composite half is open}. E8 is declared ``complete as specified: 4b by route~1;
  (vii) over the ring open''. \emph{v1.10.0-pre.2, as built and declared} (status paragraph
  above): items 1--4b as specified; 4b over $\Z/\ell^k$ by route~1; (vii)'s composite half open,
  the strict-pair route, the relative route and the relative sub-complex constructor refusing
  with arms ((b1)). (v1.11.0-pre.2: superseded by the next item and the amended declaration of the
  status paragraph.)

  \emph{v1.11.0-pre.1: (vii) over the ring, specified} (slot-10 decision~3~(A); built by slot-10
  WP-3; the three members are (b1)'s ``the lift (slot~10)''). \textbf{Scope}: $\ell\ge5$ and
  $\ell^k<2^{32}$. Route~1 is unchanged and runs every regime; the two cross-routes run where the
  \texttt{uint32\_t} classes they consume run, and refuse above by the regime rule of~(b1).
  \begin{itemize}
  \item \emph{Route~2, the strict-pair route, over the ring.} $\tau^{s}$ (the
    \texttt{BoundaryStrict} table) is integral, with solve regime \texttt{"Z"}, so its $\Theta^{s}_g$
    over $\Z/\ell^k$ is the reduction of the integral tables, assembled by the honouring
    \texttt{HeckeOperatorOnK} over the ring; the boundary operator is the honouring
    \texttt{HeckeBoundaryOperator} of~(f); the strict datum's pair homotopy $\eta^{s}$ is asserted
    $0$, as over $\Fl$. The strict action $A^{s,(q)}_c$ is assembled by route~1's module machinery
    (\texttt{ModuleConeActionT}) from the strict data, on \texttt{ModuleConeT}'s stored generators:
    the \textbf{same} cyclic-summand basis (c4) of $H^q_c$ as route~1's, so the comparison is between
    (c5) matrices in one basis. The record's fields keep their $\Fl$ meaning
    (\texttt{six\_invertible} and \texttt{asserted} true at $\ell\ge5$; \texttt{solve\_regime}
    \texttt{"Z"}; \texttt{eta\_zero}).
  \item \emph{Route~3, the relative route, over the ring.} The relative sub-complex
    $C^{\bullet}_{\mathrm{eq}}(K,K_{\mathrm{bd}})=\ker r$ is built by
    \texttt{EquivariantCohomologyOnK}'s relative constructor on the \textbf{module route of
    the bulk}, with (c1)--(c5) verbatim and $C^q_{\mathrm{eq}}$ read as
    \texttt{kernel\_module} of the equivariance constraints stacked with the vanishing constraints
    on the simplices of $K_{\mathrm{bd}}$; its Howell pivots are asserted to be units (it is expected
    free at $\ell\ge5$: a simplex's stabiliser is the same in $K$ as in $K_{\mathrm{bd}}$, so the
    restriction of equivariant cochains is onto, and the kernel of an onto map of free modules over
    the local ring $\Z/\ell^k$ is free). A non-unit pivot raises a named \texttt{std::logic\_error} and
    is a finding for slot-10 WP-15, never a silent fallback. The flat-coordinate route of~(c1$'$) is
    \textbf{not} used at $\ell\ge5$. Then $A'^{(q)}$ is the (c5) matrix of route~3's cochain map on
    the cyclic-summand basis of $H^q(K,K_{\mathrm{bd}})$, the correspondence
    $M^q\colon H^q(K,K_{\mathrm{bd}})\to H^q_c$ is a (c5) matrix, certified an isomorphism over the
    ring (equal invariant factors, and every basis element of $H^q_c$ exhibited as an image by
    \texttt{solve\_multi\_rhs\_module} modulo the target's relations), and the transported action
    $M\,A'\,M^{-1}$ is computed \textbf{column by column} --- for each $b_j$ of $H^q_c$ a preimage under
    $M$ by the solve, then $A'$, then $M$, reduced in (c5) form. No inverse matrix is formed: a (c5)
    matrix between non-free modules has none in general.
  \item \emph{Acceptance} (\HA{obl:coneexact}(vii) as written, over the ring). At each of the seven
    E8 points of~(e), in every computed degree, for $T_p$ at the point's primes ($p\in\{2,3\}$, and
    $p=5$ at $(10,691)$) and for $U_\ell$ at the cells where route~1's arms compute it, the
    $A^{(q)}_c$ of route~2 and of route~3 each \textbf{equal route~1's entrywise} in (c5) form (row
    $i$ compared as reduced residues in $[0,\ell^{e_i})$). Both agreements \textbf{abort} on failure
    at $\ell\ge5$ ($6\in R^{\times}$). A negative control: one perturbed entry of $\tau^{s}$ is
    detected (by the table's own certificate or by the disagreement; the arm says which). One arm at
    a 4b word-regime cell above $2^{32}$: route~1 runs, routes~2 and~3 refuse by the regime refusal.
    One arm each at $\ell\in\{2,3\}$: the three members refuse. Measured by
    \texttt{HeckeConeCrossRouteRingValidation} (\texttt{standard}; a point whose wall would push the
    name over its tier runs under a separate \texttt{slow} name, not a re-tier). The $\Fl$ cone is
    byte-identical (\texttt{cone\_output\_diff.sh} against the slot-9 close, $0$ hunks), and
    \texttt{HeckeReductionValidation} and \texttt{HeckeReductionCoverageSlowValidation} keep their
    check counts except for (b1)'s changed rows.
  \item \emph{Not lifted}: \HA{prop:coneroutes}(3)'s adjointness over the ring (it needs the E7
    pairing, which refuses a composite ring); (vi)'s $\nu=2$ recursion; (viii)'s two arms, which stay
    $\Fl$-only. \emph{No object}: the cross-routes are validated equal to route~1, whose objects are
    the emitted ones, so no golden line, record member or console verb is added, and golden does not
    move.
  \item \emph{The declaration}, on acceptance, is the one quoted in the ``Pre-release
    v1.11.0-pre.1'' paragraph above; slot-10 WP-15 writes it here and in the status paragraph, and
    [C] v1.10.0 records (vii) over the ring as discharged at $\ell^k<2^{32}$, every E8 point, with the
    \texttt{uint32\_t} regime rule above.
    \emph{v1.11.0-pre.2}: written, \textbf{amended} as the status paragraph says --- the first two
    sentences of (vii) discharged, the third (the adjointness gate this bullet lists as ``not
    lifted'') recorded unavailable over the ring (review finding M-1); [C] v1.10.0 records the same.
  \item \emph{v1.11.0-pre.2, as built} (slot-10 WP-3; \texttt{docs/implementation/slot10/WP-3\_report.md}).
    \textbf{Route~2} as specified: $\tau^{s}$ and $\eta^{s}$ integral exactly as over $\Fl$,
    $\eta^{s}=0$ asserted; $\Theta^{s}_g$ from \texttt{HeckeOperatorOnK}'s module route on the same
    composite bulk cohomology; $\Phi^{s}$ from route~1's \texttt{ModuleConeActionT} on the same cone
    and boundary operator; so route~1's item~2 asserts $r\,\Theta^{s}=\Theta^{\partial}r$ on the nose,
    re-derived over the ring and recorded as \texttt{cochain\_commutes}. \textbf{Route~3} with
    \textbf{one recorded deviation}: the transport is not computed column by column by solves, but by
    an \emph{explicit inverse on classes} --- with the cone differential
    $D_c(c,c_{\partial})=(-dc,\,d_{\partial}c_{\partial}+rc)$ and the equivariant section
    $\mathrm{ext}_0$ of $r$ ($r\,\mathrm{ext}_0=I$, checked), $M\colon[z]\mapsto[(z,0)]$ and
    $N\colon[(c,c_{\partial})]\mapsto[c+d\,\mathrm{ext}_0c_{\partial}]$ are homomorphisms, mutually
    inverse on classes, asserted $MN=I$ and $NM=I$ in (c5) form, and $A'$ is transported as
    $M\,A'\,N$; the equal invariant factors of $H^q(K,K_{\mathrm{bd}})$ and $H^q_c$ are then a
    consistency check. Still no inverse \emph{matrix} is formed. The review verified that $N$ is well
    defined on classes and that (c5) products of homomorphism matrices are well defined (finding
    m-11 asks for the argument in the header note). $E=\mathrm{ext}_0h$ and
    $\Theta'=\Theta+dE+Ed$, with R3.1--R3.6 as over $\Fl$; the relative constructor's pivots were
    units at every point. \textbf{Acceptance, as measured}: at the seven points ($(10,691)$ with
    $T_2$, $T_3$, $T_5$; $(6,5)$, $(16,5)$ at $\Z/\ell^2$; $(30,5)$, $(60,5)$ at $\Z/\ell^2$ and
    $\Z/\ell^3$, with $T_2$, $T_3$), routes 1, 2 and~3 are entrywise equal in every degree
    $0\ldots2$, with non-vacuity rows ($\tau^{s}\ne\tau$, $A_c\ne0$, $A_c(T_2)\ne A_c(T_3)$ at every
    level); $H_c$'s invariant factors equal $H^q(\ker r)$'s, e.g.\ $()\,|\,(2,2,2)\,|\,()$ at
    $(10,691)$. No $U_\ell$ arm: route~1's arms compute none at these cells. \textbf{Negative
    control}: a strict table perturbed after construction is refused by its own certificate
    (code~30); a table perturbed \emph{past} its certificate and pushed through by hand aborts at
    route~1's item~4 ($D_c\Phi\ne\Phi D_c$), not at the final comparison; nothing perturbs route~3,
    which has no table of its own (review finding m-12 records a stronger seam for a later slot).
    (\emph{v1.11.0}: slot-10 WP-8 asserts route~3's non-vacuity, which was printed only ---
    $E\ne0$, $\Theta'\ne\Theta$ and $r\Theta\ne\Theta^{\partial}r$ on the shipped $\tau$ --- at every
    point, level and prime, $3\times15$ checks more, so $431$ by the source; the slot close
    re-derives the count from the run. The negative seam on route~3 itself stays for slot~11.)
    \textbf{Regime arm} at $(10,134217757)$ over $\Z/\ell^2$: route~1
    (\texttt{ModuleConeActionT<uint64\_t>}) runs, both \texttt{RequireRingRegime} refuse by the
    named \texttt{uint32\_t} refusal. \texttt{HeckeConeCrossRouteRingValidation}, $386$ checks,
    $0$ failing, \texttt{standard} (no \texttt{slow} sibling was needed). The $\Fl$ cone is
    byte-identical by \texttt{cone\_output\_diff.sh} with $0$ hunks against the \texttt{bd6d1bb}
    reference binary (the plan's \texttt{2ac9ee8} reference did not exist when WP-3 ran; the slot
    close re-runs it, review finding m-13).
  \end{itemize}
\end{itemize}
\emph{The objects}, all at the seven E8 points of~(e) and nowhere else. Tier~A:
\begin{itemize}
\item \textbf{\texttt{bdry\_lengths\_H\textasciicircum q}}, $q=0,1$: the composition length of
  $H^q(C^{\bullet}_{\partial}\otimes\Z/\ell^k)$, as \texttt{lengths\_H\textasciicircum q}.
\item \textbf{\texttt{bdry\_invariant\_factors\_H\textasciicircum q}}, $q=0,1$: its list
  $e_1\ge e_2\ge\cdots\ge1$, as \texttt{invariant\_factors\_H\textasciicircum q}.
  \emph{What is known.} Over $\Q$, $\rho_n(T)$ is unipotent with one Jordan block. So
  $H^0(C_{\partial,\Z})\cong\Z$ and $H^1(C_{\partial,\Z})_{(\ell)}\cong\Z_{(\ell)}\oplus\bigoplus_t\Z/\ell^{a_t}$
  (even $n$; the $\ell^{a_t}$ are the non-unit elementary divisors of $\rho_n(T)-1$ over
  $\Z_{(\ell)}$). Universal coefficients give \textbf{the same list in both degrees}:
  $(k,\min(a_t,k)\ldots)$ sorted. So $\texttt{bdry\_lengths\_H\textasciicircum0}=
  \texttt{bdry\_lengths\_H\textasciicircum1}$, which is asserted.
\item \textbf{\texttt{bdry\_bockstein\_rank\textasciicircum0}}: the rank over $\Fl$ of the mod-$\ell$
  Bockstein $\beta^0_{\partial}\colon H^0(\partial;\Fl)\to H^1(\partial;\Fl)$, the connecting map of
  $0\to\Fl\to\Z/\ell^2\to\Fl\to0$ on $C^{\bullet}_{\partial}$, computed from the stored bases.
  \emph{Why degree~$0$.} $C^{\bullet}_{\partial}$ lives in degrees $0,1$, so $\beta^1_{\partial}$
  lands in $H^2(\partial)=0$, and degree~$0$ is the only degree in which the object is non-trivial.
  It is \textbf{defined at every point}. It is a mod-$\ell$ object, the same at every $k$ of a
  cell, and is computed from the cell's $\Z/\ell^2$ boundary object, as
  \texttt{bockstein\_rank\textasciicircum0} is.
  \emph{Meaning.} It equals $\#\{t:a_t=1\}$, by the argument of \HA{rmk:bocksteinexp}(1) applied
  to $C_{\partial,\Z}$:
  \begin{itemize}
  \item the classes of $H^0(C_{\partial,\Z})/\ell$ lift to $\Z/\ell^2$-cocycles, so
    $\beta^0_{\partial}$ kills them;
  \item on $H^1(C_{\partial,\Z})[\ell]$, $\beta^0_{\partial}$ is $G[\ell]\to G/\ell G$.
  \end{itemize}
  Since $H^2(C_{\partial,\Z})=0$, no certified level is needed: at every $k\ge2$,
  $\texttt{bdry\_bockstein\_rank\textasciicircum0}$ is the number of entries equal to~$1$ in
  \texttt{bdry\_invariant\_factors\_H\textasciicircum1}. This is \textbf{asserted at every point}.
  There is no boundary analogue of \texttt{higher\_torsion\_count}: $H^0(C_{\partial,\Z})$ is not
  zero, so $\dim H^0(\partial;\Fl)-\rank\beta^0_{\partial}$ counts a free class as well, and
  would misname it.
\item \textbf{\texttt{cone\_lengths\_Hc\textasciicircum q}}, $q=0,1,2$: the composition length of
  $H^q_c$ over $\Z/\ell^k$. This is the slot-9 plan's \texttt{cone\_lengths\_H\textasciicircum q},
  spelled \texttt{Hc} after the $\Fl$ object \texttt{dim\_Hc\textasciicircum q}, whose degree
  convention it shares. \texttt{cone\_lengths\_Hc\textasciicircum0} is $0$, asserted.
\item \textbf{\texttt{cone\_invariant\_factors\_Hc\textasciicircum q}}, $q=0,1,2$: its list, as
  above. This is the plan's \texttt{cone\_invariant\_factors\_H\textasciicircum q}.
\item \textbf{\texttt{coneexact\_iv\_certified}} (boolean): \texttt{true} iff the non-field half
  of \HA{obl:coneexact}(iv) above ran at every spot all three of whose terms are reported, and
  passed. On the model complexes every spot is reported. A failure aborts, so on an emitted point
  the value is \texttt{true}. \texttt{false} is the pre-slot-9 state, ``recorded as unchecked'',
  and no slot-9 build emits it at an E8 point.
\item \textbf{\texttt{cone\_les\_im\_lengths}}: the plan's ``cone LES length identity'', as an
  object with content. It is the list of the lengths of the images of the eight map instances of
  \HA{eq:coneles}, in [C]'s order
  $\pi^0_{\ast},r^0,\iota^0_{\ast},\pi^1_{\ast},r^1,\iota^1_{\ast},\pi^2_{\ast},r^2$
  (\HA{obl:coneexact}, preamble). Here $\pi^q_{\ast}\colon H^q_c\to H^q$, and
  $\iota^q_{\ast}\colon H^q(\partial)\to H^{q+1}_c$. It is emitted only after \textbf{the length
  identity} is asserted at each of the seven terms
  $H^0,H^0(\partial),H^1_c,H^1,H^1(\partial),H^2_c,H^2$ of the sequence: the length of the term
  equals the image length of the incoming map plus that of the outgoing map. This is
  \HA{eq:coneles}'s exactness in lengths, with $H^0_c=0$ and $H^2(\partial)=0$ at its ends.
\end{itemize}
Tier~B (basis-dependent; (c)'s bases):
\begin{itemize}
\item \textbf{\texttt{bdry\_red\_H\textasciicircum q}}, $q=0,1$: the reduction
  $H^q(\partial;\Z/\ell^k)\to H^q(\partial;\Fl)$. It is the (c5) matrix from the cyclic-summand
  basis to the D5b basis of the $\Fl$ boundary complex, the one \texttt{SL2ZBoundaryPair} builds
  at the defaults. As in~(d), $\mathrm{red}=P^{\partial}_q(b_j\bmod\ell)$.
\item \textbf{\texttt{cone\_red\_Hc\textasciicircum q}}, $q=0,1,2$: the same map for $H^q_c$, into
  the $\Fl$ cone's stored basis \texttt{Hc\textasciicircum q}, through its projector
  \texttt{Pc\textasciicircum q}.
\end{itemize}
Reductions compare as maps, never as bases, exactly as in~(d). The Hecke matrices
$A^{\partial}(T_p)$ and $A_c(T_p)$ over $\Z/\ell^k$ are computed and validated, but are
\textbf{not} objects in slot~9, as $A^{(q)}(T_p)$ over the ring is not.

\emph{Hash recipes}, as~(e):
\begin{itemize}
\item a scalar stores its value, and the boolean stores $0$ or $1$;
\item a list hashes its length and then its entries. \texttt{rows} is the length and
  \texttt{cols} is~$1$. \texttt{cone\_les\_im\_lengths} has $8$ rows;
\item a reduction matrix hashes \texttt{rows}, then \texttt{cols}, then its entries row-major.
  $0\times0$ is legal, for example \texttt{cone\_red\_Hc\textasciicircum0}.
\end{itemize}
\emph{The generating rule of slot-9 WP-4's pre-stated set} (the event R1, (a/s); the lead lands it,
after slot-9 WP-14 pass~1). At each of the seven points of~(e) the set is:
\begin{itemize}
\item tier~A: \texttt{bdry\_lengths\_H\textasciicircum\{0,1\}},
  \texttt{bdry\_invariant\_factors\_H\textasciicircum\{0,1\}},
  \texttt{bdry\_bockstein\_rank\textasciicircum0}, \texttt{cone\_lengths\_Hc\textasciicircum\{0,1,2\}},
  \texttt{cone\_invariant\_factors\_Hc\textasciicircum\{0,1,2\}}, \texttt{coneexact\_iv\_certified}
  and \texttt{cone\_les\_im\_lengths}, which is $13$;
\item tier~B: \texttt{bdry\_red\_H\textasciicircum\{0,1\}} and
  \texttt{cone\_red\_Hc\textasciicircum\{0,1,2\}}, which is $5$.
\end{itemize}
No object depends on a prime or on the certified level. Every object is defined at every point, so
\textbf{the ``absent where undefined'' rule removes nothing at these seven points}. The set is
$7\times18=\mathbf{126}$ entries, $91$ tier~A and $35$ tier~B.
\begin{itemize}
\item They are emitted at each point after its existing E8 objects, in the order above; no other
  point gains or loses an entry.
\item Predicted counts after R1: tier~A $11\,575+91=11\,666$, tier~B $5866+35=5901$, entries
  $17\,441+126=17\,567$.
\item Both determinism digests are unchanged by construction, because the E8 points are looped from
  \texttt{ComputeAll()} only.
\end{itemize}

\emph{Predicted values}. A prediction that fails is a finding for the owner, and is never
re-predicted. The boundary is predicted exactly, from the elementary divisors of $\rho_n(T)-1$
over $\Z_{(\ell)}$, computed independently of the pipeline by exact rational elimination (script and
output in slot-9 WP-2's report). Write $B:=\texttt{bdry\_lengths\_H\textasciicircum0}$,
$L_q:=\texttt{lengths\_H\textasciicircum q}$ (the landed values of~(e)), and
$c_2:=\texttt{cone\_lengths\_Hc\textasciicircum2}$. Exactness and $H^2=0$ then determine the
whole of \texttt{cone\_les\_im\_lengths} from $c_2$:
\[
  (0,\;L_0,\;B-L_0,\;L_1-B+c_2,\;B-c_2,\;c_2,\;0,\;0),\qquad
  \texttt{cone\_lengths\_Hc\textasciicircum1}=L_1-L_0+c_2 .
\]
Here $r^0$ is injective (\HA{lem:conerelative}(4)), $\pi^2_{\ast}$ lands in $H^2=0$, and
$\iota^1_{\ast}$ is therefore onto $H^2_c$.
\begin{center}\small
\begin{tabular}{@{}lccccp{0.45\textwidth}@{}}
\toprule
cell & $K$ & $H^0(\partial)=H^1(\partial)$ & $B$ & $\beta_{\partial}$ & cone \\
\midrule
$(10,691)$ & 2 & $(2)$ & 2 & 0 & $H^1_c=(2,2,2)$, $H^2_c=()$; im-lengths $(0,0,2,4,2,0,0,0)$ \\
$(16,5)$ & 2 & $(2,1,1,1)$ & 5 & 3 & $H^1_c=(2,2,2,m)$, $H^2_c=(m)$, $m\in\{1,2\}$ \\
$(6,5)$ & 2 & $(2,1)$ & 3 & 1 & $\ell(H^1_c)=2+c_2$ \\
$(30,5)$ & 2 & $(2,2,1^5)$ & 9 & 5 & $\ell(H^1_c)=10+c_2$ \\
$(30,5)$ & 3 & $(3,2,1^5)$ & 10 & 5 & $\ell(H^1_c)=15+c_2$ \\
$(60,5)$ & 2 & $(2^3,1^{10})$ & 16 & 10 & $\ell(H^1_c)=18+c_2$ \\
$(60,5)$ & 3 & $(3,2^2,1^{10})$ & 17 & 10 & $\ell(H^1_c)=27+c_2$ \\
\bottomrule
\end{tabular}
\end{center}
\noindent The third column is \texttt{bdry\_invariant\_factors\_H\textasciicircum q}, the same list at $q=0,1$. Exponents are listed with multiplicity. $\beta_{\partial}$ is
\texttt{bdry\_bockstein\_rank\textasciicircum0}. $H^0_c=()$ and \texttt{coneexact\_iv\_certified}
$=1$ at every point.
\begin{itemize}
\item \emph{The two complete cone rows.} They use the $\Fl$ golden values of the ext cone at the same
  cell: $\dim H^2_c(\Fl)=0$ at $(10,691)$, and $\dim H^1_c(\F_5)=4$, $\dim H^2_c(\F_5)=1$ at
  $(16,5)$.
\item \emph{How they follow.} $H^2_c$ is the top cohomology of the cone, so
  $H^2_c(\Z/\ell^k)\otimes\Fl=H^2_c(\Fl)$, and it vanishes by Nakayama when the latter does. The
  $\Q$-ranks $\operatorname{rk}H^1_c=3$ and $\operatorname{rk}H^2_c=0$ (weight $12$ and $18$, one
  Eisenstein class) then fix the rest by universal coefficients.
\item \emph{What is left open.} $m=\min(a,2)$ with $5^a=\lvert H^2_c(\Z)_{(5)}\rvert$ is not
  predicted. At $(16,5)$ the im-lengths are $(0,0,5,1+m,5-m,m,0,0)$. So $\ker r^1$ over $\Z/25$
  has length $1+m\ge2$, although $r^1\otimes\F_5$ is injective: $H^1(\partial;\Z/25)$ has
  $5$-torsion summands, and elements of order~$5$ can map into them.
\item $c_2$ at the other four rows is not predicted.
\end{itemize}
\emph{v1.10.0-pre.2: the landed values} (slot-9 WP-4; R1 at \texttt{762dbc3}). \textbf{Every
prediction above held}: the boundary lists exactly, row by row; the complete $(10,691)$ row; and
the pattern $(0,L_0,B-L_0,L_1-B+c_2,B-c_2,c_2,0,0)$ at the other rows. The values left open are
\textbf{$m=1$ at $(16,5)$} --- $H^1_c=(2,2,2,1)$, $H^2_c=(1)$, im-lengths $(0,0,5,2,4,1,0,0)$, so
$\lvert H^2_c(\Z)_{(5)}\rvert=5$: $H^3_c=0$ (\HA{lem:conerelative}(4)) makes $H^2_c(\Z/25)=H^2_c(\Z)
\otimes\Z/25$, one cyclic summand $\Z/5^a$ by $\dim H^2_c(\F_5)=1$, and $m=\min(a,2)=1$ forces
$a=1$ --- and $c_2=0$ at $(6,5)$, $c_2=1$ at $(30,5)$ at both $K$, $c_2=2$ at $(60,5)$ at both
$K$, with $H^2_c=(1)$ and $(1,1)$ there at $K=2$ and at $K=3$ alike, so those summands are
$\Z/5$ exactly. The cone Bockstein $\beta^1_c$ has rank $0$ at $(6,5)$ and $(10,691)$, $1$ at
$(16,5)$ and $(30,5)$ and $2$ at $(60,5)$, consistent with
$\ell(H^2_c(\Z/\ell^2))=2\dim H^2_c(\Fl)-\operatorname{rk}\beta^1_c$, which the validation asserts
at its \texttt{uint32\_t} cells (v1.10.1: not at the two word-regime cells, which have no $\Fl$
cone class to compare with).
Arm (4b-C) pins $m=1$ at $(16,5)$ since \texttt{7ddb555} (review finding m-W4-4). The landed set is
exactly the generating rule's: $126$ entries, $91$ tier~A and $35$ tier~B, emitted after each
point's existing E8 objects; three emissions byte-identical, one of them under
\texttt{ENGINE\_THREADS=1}; \texttt{golden\_ref\_diff.py --expect} holds it as a set equality and
\texttt{golden\_explain.py} exits~$0$ with no flag. \emph{Four of the eighteen objects per point
are structural constants} --- \texttt{cone\_lengths\_Hc\textasciicircum0} $=0$,
\texttt{cone\_invariant\_factors\_Hc\textasciicircum0} $=()$, \texttt{cone\_red\_Hc\textasciicircum0}
$0\times0$ and \texttt{coneexact\_iv\_certified} $=1$ --- $28$ entries that cannot move without a
certificate failing first; they are frozen as specified (review finding m-W4-4). Measured by
\texttt{GOLDEN\_STRICT=1 GoldenRegressionValidation}: tier~A $11\,666$ failing~$0$, tier~B $5901$
mismatching~$0$.

\paragraph{(g) The job-level schema \texttt{hecke\_ring\_report/v1} (v1.10.0-pre.1).}
Slot-9 decision~6~(A); built by slot-9 WP-8 (C5c). Documented for users in \texttt{campaign.md}
\S15, which slot-9 WP-8 writes. \emph{v1.11.0-pre.1} (slot-10 decision~8~(A)): \textbf{no schema
change in slot~10} --- slot-10 WP-4's $\ell$-adic survey under $\Z/\ell^2$ writes
\texttt{hecke\_ring\_report/v1} exactly as frozen in v1.10.0, (g1)--(g16) with the four recorded
deviations and the (g11) amendment, and a finding that would need a member is reported for a later
version, not added. \emph{v1.11.0-pre.2}: no code package of slot~10 touched the schema or
\texttt{Campaign/}. \emph{v1.11.0}: slot-10 WP-4's survey was \textbf{not run in slot~10} (the
optional W2; deferred), so nothing here moves. \emph{v1.11.1}: the survey \textbf{ran in W2}
(\texttt{docs/matrix-engine/results/ring\_survey\_l2/}, $462$ records): every record is
\texttt{hecke\_ring\_report/v1} exactly as frozen, (g1)--(g16) with the recorded deviations, and no
finding of its analysis needs a member --- so nothing here moves.

\begingroup\sloppy\emph{v1.10.0, as built} (slot-9 WP-8, \texttt{ad65bd9}, with the (g11)
amendment \texttt{f4fc3d2}, merged at \texttt{8853544}; \texttt{Campaign/CampaignRing.\{h,cpp\}};
\texttt{campaign.md} \S15). \texttt{RunOneJob} dispatches to the ring job whenever
\texttt{coefficient\_ring} is not \texttt{F\_ell}; nothing of the field path runs. The members
(g1)--(g16) are emitted in that order, with \texttt{certify} after (g5); the file is assembled in
memory and written once. At $\ell\ge5$ the job computes the cohomology by the module route; at
$\ell\in\{2,3\}$ by the flat-coordinate route of~(c1$'$), through
\texttt{EquivariantCohomologyOnK}'s own lift: since the merge fix \texttt{158c0be} the job loop
builds its relation list there from the raw \texttt{ModCtx::prime\_power} context (the value
\texttt{CoefficientRingContext} would return) and does \textbf{not} name
\texttt{SmallEllPolicy::Bredon}, which stays with \texttt{EquivariantCohomologyOnK} alone ((b1),
the scoping row); at every other $\ell$ the context keeps its refusal rule. The element width is
\texttt{uint32\_t} for $\ell^k<2^{32}$ and \texttt{uint64\_t} otherwise. \textbf{Four deviations
from the text below, recorded as built} (slot-9 WP-8 report, findings 1, 2, 4 and~6):
(i)~at $\ell\ge5$ \emph{off the $K$-mode's range} ($n=0$ or $n$ odd) \texttt{kmode} is
\texttt{"not\_applicable"}, as at $\ell\in\{2,3\}$ ((g7) named $2$ and $3$ only), and
\texttt{ladic}, $r_K(d^q)$ and \texttt{ladic\_type*} are absent, as~(d) says, while
\texttt{bockstein} and \texttt{red\_H\textasciicircum q} are present (they are mod-$\ell$ objects,
defined there); (ii)~the uncertified \texttt{ladic.state} of~(g9) also carries
\texttt{need: \{r\_K0, r\_K1\}}, the deficit (the console's \texttt{ladic\_level} state carries the
same); (iii)~\texttt{hecke} is computed iff $\ell\ge5$, $r=1$ and $1\in$ \texttt{point.degrees}
((g11)); otherwise, at $\ell\ge5$, it is absent and \textbf{not} refused --- at $\ell\in\{2,3\}$ it
is refused whatever the degree selection ((g14)); (iv)~each refusal's \texttt{who} names the member
the job loop called (\texttt{HeckeOperatorOnK}, \texttt{HeckeReduction (BocksteinCounts)},
\texttt{HeckeReduction (ReduceHToField)}, \texttt{campaign job loop}) and its \texttt{message} is
that guard's own text, obtained by calling the real guard (at $\ell\in\{2,3\}$ the Bockstein and
reduction messages name the Bredon object's accessor reached, \texttt{module\_degree} or
\texttt{H\_module\_basis}); a guard that no longer refuses \textbf{fails the job}, so a later lift
needs this paragraph amended first. Three further readings, not deviations: \texttt{timings\_s}
keeps \texttt{campaign\_job/v1}'s nine keys (\texttt{cosets}, \texttt{tau}, \texttt{theta} the
$r=1$ Hecke part; \texttt{family} and \texttt{external} $0$; \texttt{report} the objects' own
time, including the $\Fl$ and $\Z/\ell^2$ comparison objects); \texttt{certify.mode} is the spec's
certify text verbatim; in \texttt{summary.tsv} the \texttt{status} column carries
\texttt{ok\_partial} and a degree outside \texttt{point.degrees} holds \texttt{-}, while
\texttt{index.json}'s per-job \texttt{status} stays \texttt{ok} ((g) adds exactly one top-level key
there). \textbf{(g11) amended} (lead decision of 2026-09-24, before this version froze the
schema): the free \texttt{charpoly}'s coefficients are decimal strings (below). \textbf{Acceptance}
is measured by \texttt{CampaignRingValidation} (\texttt{standard}, $693$ checks): arm~(A), two ring
campaigns at \texttt{jobs = 2} over the seven E8 points of~(e) with each point's primes
($17$ jobs), every carried object \textbf{equal to the frozen golden entry} (tier, rows, columns,
hash) and the object set equal to the point's 4a golden set for the job's $p$, the (e) table's
certified flag, $H^0$ / $H^1$ invariant factors, Bockstein rank and higher-torsion count literally,
through \texttt{summary.tsv}; arm~(T), the weight-$12$ free \texttt{charpoly} over $\Z/\ell^2$ at
$(10,691)$ and at the \texttt{uint64\_t} cell $(10,134217757)$, in string form; arm~(B), four cells
at $\ell\in\{2,3\}$ against the frozen O7 rows, \texttt{not\_applicable}, the three refusals and
\texttt{ok\_partial}; arm~(R), $r=2$; arms~(S) and~(X), the spec surface of~(b4).
\textbf{The fingerprint-$0$ byte identity} is measured against the base binary: on
\texttt{ext\_compare\_mini.spec} ($84$ ext records raw-identical; the $84$ job files,
\texttt{index.json} and \texttt{summary.tsv} identical once the timing and RSS numbers are blanked),
on a reduced field sweep ($104$ files) and a resident \texttt{level\_a = always},
\texttt{certify = full} spec, by \texttt{--dry-run} of the seven shipped specs (byte-identical, so
\texttt{sweep\_full}'s job ids are unchanged), and by the existing campaign validations with their
id assertions unchanged. The job loop is clean under \texttt{CheckThread} at
\texttt{jobs = 2}.\par\endgroup
\begin{itemize}
\item \emph{When.} The campaign's job loop writes this schema \textbf{only} under a non-default
  \texttt{coefficient\_ring}, and then instead of \texttt{campaign\_job/v1}.
\item \emph{A fingerprint-$0$ job} writes \texttt{campaign\_job/v1} with its
  \texttt{hecke\_report/v2} reports \textbf{byte for byte as before}. The same holds for its
  \texttt{summary.tsv} and \texttt{index.json}. No default id, file name, header or shape
  signature moves.
\item \emph{Golden cost: none.} The schema exists only at a fingerprint $\ne0$, and those texts are
  not golden ((a), ``Fingerprint and golden cost'').
\item \emph{The ext pass} keeps refusing a non-default ring by name, at spec validation, (b4).
\end{itemize}
\emph{The job file.} It is \texttt{jobs/<id>.json}, where the id is the existing
\texttt{Job::id()} with its settings suffix \texttt{\_s<16 hex>}, because the fingerprint is
non-zero; there is no new id rule. It is one JSON object whose members come in this order:
\begin{enumerate}[label=(g\arabic*)]
\item \texttt{"schema": "hecke\_ring\_report/v1"}.
\item \texttt{"job\_id"}: the id above.
\item \texttt{"point"}, \texttt{"model\_descriptor"}, \texttt{"shape\_signature"}: as in
  \texttt{campaign\_job/v1} (the same writers). \texttt{point.degrees} is the job's
  \emph{reported selection}, which is \texttt{Job::id()}'s, not a complex's degree fields.
\item \texttt{"coefficient\_ring"}: the canonical text \texttt{"Z/ell\textasciicircum<k>"}.
  \texttt{"modulus"}: $\ell^k$ as a \emph{decimal string}, since it may exceed $2^{53}$.
  \texttt{"ladic\_level"}: $K=k$, an integer.
\item \texttt{"settings"}: the \texttt{MathSettings::to\_json} block of
  \texttt{campaign\_job/v1}, with every registered key, its value and \texttt{source}, and then
  \texttt{"fingerprint"}, which is never \texttt{"0000000000000000"} here.
\item \texttt{"degree\_fields"}: the bulk complex's per-complex fields \texttt{q\_max},
  \texttt{truncated}, \texttt{q\_max\_source} and \texttt{degrees}
  ($2$ / \texttt{false} / \texttt{"derived"} / $[0,1,2]$ on the model).
\item \texttt{"kmode"}: \texttt{"certified"}, \texttt{"uncertified"} or
  \texttt{"not\_applicable"}. The last is used at $\ell\in\{2,3\}$
  (\HA{prop:certifiede}(5)) and is \textbf{not} a refusal. \emph{v1.10.0, as built}: also at
  $\ell\ge5$ off the $K$-mode's range, $n=0$ or $n$ odd (deviation~(i) above).
\item \texttt{"cohomology"}: one object per $q$ in \texttt{point.degrees}, ascending, each
  \texttt{\{"q", "length", "invariant\_factors"\}}. These are (d)'s
  \texttt{lengths\_H\textasciicircum q} and \texttt{invariant\_factors\_H\textasciicircum q}; the
  list is descending, and \texttt{[]} is the zero module. \textbf{No member anywhere in the file
  is a dimension.}
\item \texttt{"ladic"} (\textbf{present only at $\ell\ge5$}): \texttt{\{"r\_K": [$r_K(d^0)$,
  $r_K(d^1)$], "certified": <bool>, \ldots\}}. It carries \texttt{"type"} (a list of
  \texttt{\{"q", "free\_rank", "torsion"\}} for $q=0,1,2$) \emph{only} when certified. It carries
  \texttt{"state"}, the R3 state of kind \texttt{ladic\_level} of~(d) (budget $\{K\}$, reached
  $\{r_K(d^0),r_K(d^1)\}$, \texttt{complete = false}), \emph{only} when uncertified. The state's
  continuation is a fresh job at a larger $k$, never \texttt{resume}. \emph{v1.10.0, as built}:
  the state also carries \texttt{need: \{r\_K0, r\_K1\}} (deviation~(ii) above).
\item \texttt{"bockstein"} (\textbf{present only at $\ell\ge5$}):
  \texttt{\{"bockstein\_rank\_0", "higher\_torsion\_count", "dim\_H0\_F\_ell"\}}, (d)'s objects
  with the bookkeeping identity asserted. They are computed from the cell's $\Z/\ell^2$ object and
  the $\Fl$ object, and each of those is built with its own ring installed for its construction
  only (the \texttt{RunE8Point} precedent). \texttt{dim\_H0\_F\_ell} is a dimension of the
  \emph{field} run, named as one.
\item \texttt{"hecke"} (\textbf{present only at $\ell\ge5$, $r=1$, $1\in$
  \texttt{point.degrees}}): \texttt{\{"p", "r": 1, "charpoly\_free"\}}.
  \texttt{"charpoly\_free"} is \texttt{\{"precision": $K-e_{\max}$, "coeffs": [\ldots]\}}, low to
  high, each coefficient a \emph{decimal string} in $[0,\ell^{K-e_{\max}})$ (it may exceed
  $2^{53}$, as (g4)'s \texttt{modulus} may; v1.10.0, lead decision of 2026-09-24 --- through
  v1.10.0-pre.2 this item wrote JSON numbers), and is \emph{present only when certified}. At
  $p=\ell$ the operator is $U_\ell$. \emph{v1.10.0, as built} (slot-9 WP-8, \texttt{f4fc3d2}): no
  other member of the file can exceed $2^{53}$ --- under $k\ge2$, $\ell^k<2^{64}$ forces
  $\ell<2^{32}$; every count is bounded by the complex; hashes and \texttt{modulus} are strings ---
  and \texttt{CampaignRingValidation} arm~(T) asserts the string form at $(10,691)$ and at the
  \texttt{uint64\_t} cell $(10,134217757)$ (for example
  \texttt{"18014406292930825"}).
\item \texttt{"objects"}: a list of \texttt{\{"name", "tier", "rows", "cols", "fnv1a64"\}}, one
  for every (d) object present, under its golden name of~(e) and its hash recipe of~(e), in (e)'s
  order. \texttt{red\_H\textasciicircum q} (tier~B) is included at $\ell\ge5$, and
  \texttt{"fnv1a64"} is sixteen lower-case hex digits. Wherever a job's cell, $k$ and $p$ coincide
  with an E8 point of~(e), each entry equals that point's golden entry. This is the comparison
  slot-9 WP-8's acceptance~(a) makes.
\item \texttt{"refusals"}: a list of \texttt{\{"part", "who", "message"\}}, one per part of the job
  that was refused by name. \texttt{part} is one of \texttt{"hecke"}, \texttt{"bockstein"},
  \texttt{"reduction"}. \texttt{message} carries the tag
  \texttt{not\_implemented\_in\_this\_build:}. A refused part is \textbf{never silently skipped}.
  \begin{itemize}
  \item At $\ell\in\{2,3\}$, \texttt{hecke}, \texttt{bockstein} and \texttt{reduction} are refused
    (\texttt{HeckeOperatorOnK}, \texttt{HeckeReduction}).
  \item At $r\ge2$, \texttt{hecke} is refused at the job loop's own dispatch: $T_{p^r}$ over the
    ring is not validated in this build.
  \end{itemize}
  The list is empty when nothing was refused.
\item \texttt{"status"}: \texttt{"ok"} (no refusal), \texttt{"ok\_partial"} (the cohomology
  computed, and at least one part refused and listed) or \texttt{"failed"} (an error, written as
  \texttt{campaign\_job/v1} writes a failed job).
\item \texttt{"timings\_s"}, \texttt{"peak\_rss\_mb"}, \texttt{"peak\_rss\_scope"}: as in
  \texttt{campaign\_job/v1}.
\item \texttt{"campaign\_job\_complete"}: the same completion trailer, \texttt{true}, or
  \texttt{false} on failure. So \texttt{resume} works unchanged.
\end{enumerate}
The \texttt{certify} knob of the spec is accepted and recorded as \texttt{"certify": \{"mode":
<spec value>, "applicable": false\}} after (g5). Every external cross-check is a field route, and
\texttt{HeckeCrossChecks} refuses a ring. The in-process certificates always run, and a failure
among them fails the job.

\emph{\texttt{summary.tsv} and \texttt{index.json} of a ring campaign.}
\begin{itemize}
\item \emph{\texttt{summary.tsv}.} The 21 columns keep their names and order. The ten field-valued
  ones hold \texttt{-}: \texttt{dim\_H0}, \texttt{dim\_H1}, \texttt{dim\_H2},
  \texttt{charpoly\_A0}, \texttt{charpoly\_A1}, \texttt{charpoly\_A2}, \texttt{trace\_A1},
  \texttt{rank\_A1}, \texttt{theta\_flat\_hash\_q1} and \texttt{external}. Fifteen columns are
  appended: \texttt{coefficient\_ring ladic\_level kmode length\_H0 length\_H1 length\_H2
  invariant\_factors\_H0 invariant\_factors\_H1 invariant\_factors\_H2 r\_K0 r\_K1
  bockstein\_rank0 higher\_torsion\_count charpoly\_free\_A1 refusals}.
\item \emph{How the new columns are written.} A list is semicolon-separated (the \texttt{charpoly}
  convention: coefficients low to high, then the precision after a \texttt{@}). An empty list is
  \texttt{()}. An absent or undefined value is \texttt{-}. \texttt{refusals} holds the refused
  parts, semicolon-separated, or \texttt{-}.
\item \emph{\texttt{index.json}} gains one top-level key, \texttt{"coefficient\_ring"}, only in a
  ring campaign.
\item A campaign with a fingerprint-$0$ spec writes neither change.
\end{itemize}

\paragraph{(h) The console under a non-default ring (C5a; v1.10.0-pre.1).}
Slot-9 decision~2~(A); built by slot-9 WP-12. [P] v3.7 (slot-9 WP-15) records the binding surface.
This paragraph fixes \textbf{which verbs and members honour $\Z/\ell^k$}, and that
\textbf{every other pipeline verb keeps exit~$5$} by~(b), through a lazy handle too.
\begin{itemize}
\item \emph{Adding to the list.} A member not listed here refuses with exit~$5$, naming itself. It
  may honour only after this list is amended, as a pre-release amendment by the lead or at
  slot-9 WP-15.
\item \emph{Exit codes.} A named \texttt{std::logic\_error} surfaces as exit~$3$, and
  \texttt{hecke::NotImplementedInThisBuild} as exit~$5$.
\item \emph{$\ell\in\{2,3\}$.} \textbf{In slot~9 the console lifts nothing at $2$ and $3$}
  (slot-9 plan WP-3, acceptance~(b)): under a composite ring, every console pipeline verb at
  $\ell\in\{2,3\}$ keeps exit~$5$. The Bredon object at $2$ and $3$ is reachable from the
  validations and from the campaign's job loop, (g).
\end{itemize}
At $\ell\ge5$:
\begin{itemize}
\item \texttt{symmetric\_power(Model, n, ell)} \textbf{honours}. It builds the ring's relation
  list, \texttt{GenerateSL2ZRelationsRing}, over $\Z/\ell^k$.
\item \texttt{cohomology(L)} \textbf{honours} on a ring \texttt{LocalSystem}, and returns a
  \texttt{Cohomology} handle over the ring (the E8 type; [P] \S3.6: ``\texttt{coefficient\_ring}
  becomes a member of \texttt{Cohomology}''). Its honouring members:
  \begin{itemize}
  \item \texttt{.coefficient\_ring}, the canonical text; also \texttt{F\_ell} on a field handle,
    and on every \texttt{HeckeOperator};
  \item \texttt{.computed}, and the per-complex degree members \texttt{.q\_max},
    \texttt{.truncated}, \texttt{.q\_max\_source} and \texttt{.degrees};
  \item \texttt{.lengths}, a list over $q=0,1,2$, and \texttt{.invariant\_factors(q)};
  \item the $K$-mode: \texttt{.ladic\_type\_certified}, \texttt{.r\_K(q)} for $q\in\{0,1\}$, and
    \texttt{.ladic\_type(q)}, a Record \texttt{\{free\_rank, torsion\}}. At an uncertified level
    \texttt{.ladic\_type(q)} raises its named \texttt{std::logic\_error}, which is exit~$3$;
  \item \texttt{state(H)} and \texttt{save\_state(H)}, which return the \texttt{ladic\_level}
    state. \texttt{resume(H, budget)} on it is refused, by \S\ref{sec:schema}'s rule for a changed
    non-budget setting, with that rule's exit code: the continuation is a fresh run at a larger
    $k$;
  \item \texttt{.bockstein\_rank} and \texttt{.higher\_torsion\_count}, (d)'s objects, computed as
    in~(g10);
  \item \texttt{.red(q)}, the tier-B matrix \texttt{red\_H\textasciicircum q}. Its $\Fl$ comparison
    object is built with \texttt{F\_ell} installed for that construction only.
  \end{itemize}
  The field members \texttt{.dims}, \texttt{.Z(q)}, \texttt{.B(q)}, \texttt{.H(q)},
  \texttt{.C(q)}, \texttt{.P(q)}, \texttt{.A\_C(q)} and \texttt{.A\_Z(q)} raise the object's named
  \texttt{std::logic\_error}, which is exit~$3$. There is no \texttt{dim} over $\Z/\ell^k$
  ((b3)).
\item \texttt{hecke(H, p)} with $r=1$ \textbf{honours} on a ring \texttt{Cohomology}, $p=\ell$
  included (then it is $U_\ell$). Its honouring members:
  \begin{itemize}
  \item \texttt{.coefficient\_ring};
  \item \texttt{.A(q)}, the (c5) matrix on the cyclic summands;
  \item \texttt{.charpoly\_free}, the Record \texttt{\{precision, coeffs\}} of~(d), defined only at
    a certified level; otherwise it raises the named \texttt{std::logic\_error}, which is exit~$3$.
  \end{itemize}
  Every other member exits~$5$.
\item \texttt{iota(H)} \textbf{honours} on a ring \texttt{Cohomology} (v1.20.0-pre.6, an amendment of this list by slot-20 WP-LV3(d); Remark~\ref{rmk:lviotak}), at $\ell\ge5$, when the ring in force
  equals the handle's (a mismatch refuses at dispatch, exit~$5$, as \texttt{hecke(H, p)}; a digit handle refuses when forced). Its honouring members: \texttt{n}, \texttt{ell}, \texttt{status},
  \texttt{coefficient\_ring}, \texttt{certificate}, \texttt{json} (schema \texttt{hecke\_involution\_ring/v1}), and the module members \texttt{lengths\_plus}, \texttt{lengths\_minus},
  \texttt{invariant\_factors\_plus(q)}, \texttt{invariant\_factors\_minus(q)} (also on $\Fl$ handles, as $k=1$). The field members \texttt{dims}$\ast$ raise a named \texttt{std::logic\_error}
  (exit~$3$); the field methods \texttt{matrix}, \texttt{plus}, \texttt{minus}, \texttt{sign\_split}, \texttt{eisenstein\_line} exit~$5$. Nothing changes on an $\Fl$ handle. Script
  \texttt{049\_lv3\_iota\_ring.oc}.
\end{itemize}
Everything else keeps exit~$5$ under a non-default ring:
\begin{itemize}
\item \texttt{hecke(H, p, r)} with $r\ge2$, refused at the verb's dispatch;
\item \texttt{family}, \texttt{decompose}, \texttt{report} and \texttt{crosschecks};
\item every extension-layer verb of C5b (slot-9 WP-13), including the reserved
  \texttt{boundary\_cohomology(H)} and \texttt{compact\_support(H)}. \textbf{The objects of~(f) are
  console-unreachable in slot~9}, as E8's were in slot~8;
\item the reserved \texttt{reduce(H, k)} and \texttt{bockstein(H)}, which stay reserved and unbound.
\end{itemize}
\texttt{golden()} still refuses at a non-zero fingerprint (R19). The harness verbs lift the
session's settings to the baseline as before.
\emph{Tests}: arm~(K) of \texttt{ConsoleErrorValidation}, extended: each honoured verb and member
succeeds, and every other one exits~$5$ (or~$3$, where named above). The first-freeze scripts
\texttt{023\_}\ldots\ assert the values against \texttt{HeckeReductionValidation}'s printed ones
at the E8 points. [P]'s ``no edit to \texttt{Console/Core/}'' holds, except the two named
exceptions of~(a).

\emph{v1.10.0-pre.2, as built} (slot-9 WP-12, \texttt{b3cd0d3}; follow-up \texttt{f28a0fa}).
The list above honours, exactly, at $\ell\ge5$; nothing is lifted at $\ell\in\{2,3\}$. The ring is
read once, at the verb's dispatch, and kept in the handle; a lazy handle installs exactly that ring
around its construction, so a handle's ring never depends on who forces it. The comparison
objects of \texttt{.red} and of the Bockstein members (the $\Fl$ object and, at $k\ne2$, the
$\Z/\ell^2$ one) are each built with their own ring installed for their construction only, the
\texttt{RunE8Point} precedent. \texttt{.ladic\_type(q)} and \texttt{.charpoly\_free} are Records
(\texttt{\{free\_rank, torsion\}}, \texttt{\{precision, coeffs\}}); the \texttt{ladic\_level} state is
the handle's R3 state (\texttt{budget\_setting = coefficient\_ring}, \texttt{resumable} exactly
when uncertified); \texttt{save\_state(H)} forces a declared ring handle first, \texttt{state(H)}
reports once it has computed, and \texttt{resume} is refused --- by the core's ``is complete'' at
a certified level, by the type's non-budget-rule text (exit~$3$) at an uncertified one. The
\texttt{hash()} of a List of Ints follows the golden list recipe (length, then entries), so
\texttt{hash(H.invariant\_factors(q))} equals \texttt{invariant\_factors\_H\textasciicircum q}
byte for byte (the empty list hashed identically under both recipes, so no earlier value moved).
Scripts \texttt{023\_e7\_duality\_10\_691}, \texttt{024\_e8\_ring\_points} (the seven E8 points
against the literals \texttt{HeckeReductionValidation} prints, and by golden hash where a console
value has the golden recipe) and \texttt{025\_e8\_console\_surface}, each a first freeze; the
$22$ earlier transcripts are byte-unchanged. \textbf{Four as-built facts this paragraph did not
state} (WP-12 report \S4):
\begin{itemize}
\item \emph{A ring mismatch is exit~$5$.} A verb handed a handle built over a different ring
  refuses at dispatch with \texttt{hecke::NotImplementedInThisBuild}, naming both rings --- the
  pipeline's own classification (\texttt{HeckeOperatorOnK}); likewise a C5b verb handed a ring
  handle after the setting has returned to \texttt{F\_ell} (\texttt{f28a0fa}).
\item \emph{\texttt{.A(q)} over the ring is a List of Int rows, not a \texttt{Matrix}}: row $i$
  lives modulo $\ell^{e_i}$, and the \texttt{Matrix} type's field algorithms would be meaningless
  on it.
\item \emph{Unlisted structural members exit~$5$ on a ring handle}, as ``Adding to the list''
  says: \texttt{LocalSystem.n}, \texttt{.p}, \texttt{.d\_V}, \texttt{.relations};
  \texttt{HeckeOperator.p}, \texttt{.r}; \texttt{Cohomology.ell}, \texttt{.d\_V},
  \texttt{.dim\_C}. Listing any of them (with its code) is an owner item for a later amendment of
  this list; it is not made here.
\item \emph{A member getter's error carries no source position.} The gated types convert their
  members' exceptions into console errors in the bindings, so the exit codes are right ($5$ and
  $3$), but the core's \texttt{Evaluator::object\_member} does not wrap a getter's exception as a
  method call's is wrapped, and a getter has no position; the type-level \texttt{resume} refusal
  lacks one for the same reason. The fix is one \texttt{try}/\texttt{catch} in
  \texttt{Console/Core/}, the analogue of~(a)'s \texttt{exec\_set} back-fill, and is recorded for
  the next authorised Core batch ([P] v3.7).
\end{itemize}
The E7 verbs \texttt{dual} and \texttt{pairing} and every C5b verb (slot-9 WP-13) refuse a
non-default ring with exit~$5$, as listed above.\par

\emph{v1.11.0-pre.1: the structural members, and the ring-mismatch rule} (slot-10
decision~6~(ii)~(A); built by slot-10 WP-12, from its lift commit). The list above is amended, as
``Adding to the list'' requires, by \textbf{five members that honour on a ring handle} at
$\ell\ge5$:
\begin{itemize}
\item \texttt{LocalSystem.n}: $n$, an Int, as on a field handle.
\item \texttt{LocalSystem.d\_V}: $n+1$, an Int --- over $\Z/\ell^k$ the rank of the free module
  $\Sym^n(R^2)$, not a dimension over a field; the number is the same.
\item \texttt{LocalSystem.coefficient\_ring}: \textbf{a new member of the type}, present on every
  handle, as \texttt{Cohomology}'s and \texttt{HeckeOperator}'s are: the canonical text
  \texttt{"Z/ell\textasciicircum<k>"} on a ring handle, \texttt{"F\_ell"} on a field handle. It is
  the one member slot~10 adds to a type; \texttt{docs/console/console.md} is regenerated by the lead
  after slot-10 WP-12 merges.
\item \texttt{HeckeOperator.p}: the Hecke prime, an Int ($p=\ell$ is $U_\ell$).
\item \texttt{HeckeOperator.r}: an Int, always $1$ on a ring handle (\texttt{hecke(H, p, r)} with
  $r\ge2$ is refused at dispatch, above).
\end{itemize}
Each is a read of the handle's stored data: it forces no computation, so on a declared lazy handle
it returns without computing, as on a field handle. \textbf{Every other unlisted member keeps
exit~$5$} through \texttt{bind::GateTypeOnRing}, in particular \texttt{LocalSystem.p} (it would read
$\ell$ while the ring is $\Z/\ell^k$; \texttt{.coefficient\_ring} carries both) and
\texttt{.relations}, \texttt{HeckeOperator}'s other unlisted members, and \texttt{Cohomology.ell},
\texttt{.d\_V}, \texttt{.dim\_C}. At $\ell\in\{2,3\}$ no ring handle can be built (the verbs refuse
at dispatch, above), so nothing changes there. \textbf{The ring-mismatch rule}, recorded as built by
slot-9 WP-12 (its finding~4, ``four as-built facts'' above) and \textbf{binding from this
pre-release}: a verb handed a handle built over a ring other than the one in force refuses
\emph{at dispatch} with \texttt{hecke::NotImplementedInThisBuild} naming both rings, exit~$5$ --- the
pipeline's own classification (\texttt{HeckeOperatorOnK}) --- and so does a C5b verb handed a ring
handle while \texttt{F\_ell} is in force; a \emph{member} read on an existing handle is not a
dispatch and follows the handle's own ring, whatever is in force ((b2): an object's ring is its
own). \emph{Tests}: arm~(K) of \texttt{ConsoleErrorValidation} --- its four rows \texttt{L.n},
\texttt{L.d\_V}, \texttt{T.p}, \texttt{T.r} change from exit~$5$ to ``honours'', each with its
value checked, and are reported as changed rows; \texttt{L.coefficient\_ring} is a new row, on a
ring handle and on a field handle; the other unlisted members (\texttt{L.p}, \texttt{L.relations},
\texttt{H.ell}, \texttt{H.d\_V}, \texttt{H.dim\_C(1)}, \ldots) keep their exit-$5$ rows; new scripts from \texttt{026\_}, first freezes. \emph{Positions}:
a member getter's error also carries its \texttt{file:line:col} from slot-10 WP-12's commit (the one
named \texttt{Console/Core/} exception of slot-10 decision~6~(i)); that is console behaviour, and
[P] v3.8-pre.1 \S2~C7, \S6 and \S8 specify it. No exit code of this paragraph changes by it.

\emph{v1.11.0-pre.2, as built} (slot-10 WP-12, \texttt{3c96d63}, merged at \texttt{6726075};
\texttt{docs/implementation/slot10/WP-12\_report.md}). \textbf{The five members} honour exactly as
listed: \texttt{LocalSystem} gains \texttt{coefficient\_ring} (\texttt{"F\_ell"} on a field or
preset handle), and \texttt{bind::GateTypeOnRing}'s rules admit \texttt{LocalSystem.n}, \texttt{.d\_V},
\texttt{.coefficient\_ring} and \texttt{HeckeOperator.p}, \texttt{.r}; each is a read of the payload,
and a declared lazy handle stays \texttt{declared} after the reads (pinned). Every other unlisted
member keeps exit~$5$. Arm~(K) of \texttt{ConsoleErrorValidation}: \texttt{L.n}, \texttt{L.d\_V},
\texttt{T.p}, \texttt{T.r} changed from exit~$5$ to~$0$, each with its value; \texttt{L.coefficient\_ring}
new; a new arm~(W) ($24$ scripts) pins the member positions, the handle's own ring, and \textbf{the
mismatch rule in both directions} ($\Z/\ell^3$ in force against a $\Z/\ell^2$ handle, $\Fl$ against a
ring handle, a ring against an $\Fl$ handle) --- the rule held at every handle-taking verb at the base,
so no code changed for it; $218\to294$ checks, tier \texttt{fast}. Script \texttt{026\_c5\_leftovers\_ring\_members}
is a first freeze; no existing transcript moved; \texttt{docs/console/console.md} regenerated (one
row, \texttt{LocalSystem.coefficient\_ring}). \textbf{Positions and one exit code}: the named
\texttt{Console/Core/} exception adds a member access's \texttt{file:line:col} to a getter's error
([P] v3.8 \S2~C7). Every gated getter of this paragraph's types already converted
\texttt{NotImplementedInThisBuild} to exit~$5$, so no code of this paragraph moved; but an
\emph{unwrapped} getter raising \texttt{hecke::NotImplementedInThisBuild} was exit~$3$ at the base
(it escaped as a plain \texttt{std::exception}) and is \textbf{exit~$5$} now, the kind classifier
deciding as for a method --- no shipped getter was found to raise it unwrapped, an audit WP-12
records as not exhaustive over the extension layer's types ([P] v3.8 \S8).\par
\endgroup

\subsection{The congruence-census catalogue \texttt{hecke\_census/v1}}\label{sec:census}
This subsection is the specification slot-11 WP-7 (the tool, the computed split, the catalogue),
WP-6 (the $\theta$-image column, \S\ref{sec:thetaobj}~(f)), WP-8 (the second edition) and WP-19
(the examples document [EX], which selects from~(g)) build against, written before any of them
commits a result. The catalogue classifies every maximal ideal of the published extension records
with the owner's column list (slot-11 plan \S3 ``WP-7'' scope~3). It is a \textbf{derived
artefact with its own format tag}, \texttt{hecke\_census/v1}: \emph{not} a record, not a campaign
output and not golden. \textbf{No record schema changes}: \texttt{hecke\_ext/v3},
\texttt{hecke\_report/v2} and \texttt{hecke\_ring\_report/v1} are untouched, the published
directories \texttt{ext\_first/}, \texttt{ext\_u691/} and \texttt{ring\_survey\_l2/} stay
byte-unchanged, and no digest or golden line moves. [CEN] below is the day-zero census note,
committed verbatim at \texttt{docs/implementation/slot11/census/} (\texttt{congruence\_census.md},
\texttt{census.py}), which the tool must reproduce.
\begin{enumerate}[label=(\alph*)]
\item \textbf{Inputs and the supersession rule.} Under \texttt{docs/matrix-engine/results/},
  the records \texttt{ext/*.json} of \texttt{ext\_first/} ($505$ cells, $\ell\in\{5,7,11,13,691\}$,
  even $n\le200$) and of \texttt{ext\_u691/} ($45$ cells), read-only. \emph{Supersession}: a cell $(n,\ell)$
  of \texttt{ext\_first} is \emph{superseded} iff \texttt{ext\_u691} holds a record of the same
  $(n,\ell)$ and \texttt{ext\_first}'s record has \texttt{u\_ell.included = false}; a superseded
  record contributes no row and is listed, with its hash, in the provenance. The tool asserts that
  the superseded set is exactly the $45$ cells $\ell=691$, $n=112,\dots,200$ of [CEN]; any other
  outcome stops the tool. The second edition adds WP-8's new directory by the same rule (a cell in
  two directories is a finding unless a supersession rule of that edition's \texttt{README.md} names
  it). Auxiliary inputs: \texttt{ring\_survey\_l2/} for~(d5); \S\ref{sec:ring}~(e)'s table for the
  E8 points' exponents; the $\theta$-image and split lines of~(e).
\item \textbf{Rows.} One row per entry of a record's \texttt{ideals} list after supersession ---
  [CEN]'s unit, $2569$ rows in the first edition. A row's key is
  $(\texttt{source\_set},n,\ell,\texttt{ideal\_id})$. Its \emph{own group} is the group whose
  \texttt{modules} entry the record carries (\texttt{H1}, else \texttt{H0}), as [CEN]'s
  \texttt{census.py} reads it; its \emph{groups} are the record's per-ideal \texttt{groups} list.
\item \textbf{Columns}, in this order; the TSV header names them exactly so. ``Record'' means the
  row's ideal entry, copied verbatim (lists as compact JSON); every cell holds a value or one of the
  explicit tokens \texttt{not\_computed}, \texttt{not\_surveyed}, \texttt{not\_evaluated},
  \texttt{not\_applicable}, \texttt{undecided} --- never an empty field.
  \begingroup\footnotesize
  \begin{longtable}{@{}>{\raggedright\arraybackslash}p{3.7cm}>{\raggedright\arraybackslash}p{3.3cm}>{\raggedright\arraybackslash}p{7.2cm}@{}}
  \toprule
  column(s) & source & rule \\
  \midrule
  \endhead
  \texttt{edition}, \texttt{source\_set}, \texttt{record\_path}, \texttt{record\_sha256} & the file & path relative to \texttt{docs/matrix-engine/results/}; sha256 of the record's bytes \\
  \texttt{n}, \texttt{ell}, \texttt{weight}, \texttt{ideal\_id} & record & \texttt{weight} $=n+2$ \\
  \texttt{own\_group}, \texttt{groups} & record & (b) \\
  \texttt{witness\_poly} & record & low$\to$high \\
  \texttt{boundary\_split}, \texttt{rank\_r1}, \texttt{dim\_E} & computed, (d1), (e) & \texttt{eisenstein}, \texttt{cuspidal}, \texttt{mixed}, \texttt{undecided}, \texttt{not\_applicable} \\
  \texttt{sign}, \texttt{dim\_plus}, \texttt{dim\_minus} & record; computed, (d2), (e) & the owner's $\pm$ column \\
  \texttt{residue\_degree}, \texttt{residue\_degree\_certified} & record & a \texttt{residue\_degree\_lower\_bound} is written \texttt{>=<f>}, never as a degree (\S\ref{sec:schema}) \\
  \texttt{length\_A}, \texttt{reduced}, \texttt{loewy\_A}, \texttt{embdim}, \texttt{gorenstein}, \texttt{hilbert} & record \texttt{algebra.*} & (d3) \\
  \texttt{module\_D}, \texttt{module\_length}, \texttt{module\_loewy}, \texttt{module\_loewy\_length}, \texttt{min\_generators}, \texttt{cyclic}, \texttt{free} & record \texttt{modules.<own group>.*} & (d3) \\
  \texttt{label\_status}, \texttt{label\_type}, \texttt{label\_system}, \texttt{label\_theta\_twist} & record \texttt{label} & verbatim; \texttt{label\_type} is the atlas's (reducible\,/\,irreducible $\bar\rho$), \emph{not} (d1) \\
  \texttt{predicted}, \texttt{confirmed}, \texttt{unpredicted} & record & verbatim; only the Kummer predictor exists ([CEN]), so \texttt{unpredicted} is not informative at $\ell\le13$ --- stated in the edition's \texttt{README.md}, not per row \\
  \texttt{kummer\_tag} & record \texttt{predicted}, cell \texttt{predictions.kummer} & \texttt{eisenstein\_cuspidal:k=<k>} with its \texttt{confirmed} flag, else \texttt{none}; the two $\ell=691$ occurrences cross-referenced (tag \texttt{pair\_691}, (g)) \\
  \texttt{theta\_image\_dim}, \texttt{theta\_image} & computed, (d4) & \texttt{none}, \texttt{partial}, \texttt{all}, \texttt{not\_computed} \\
  \texttt{zl\_depth}, \texttt{zl\_depth\_source} & (d5) & \\
  \texttt{glover\_multiplicity}, \texttt{glover\_main}, \texttt{glover\_correction}, \texttt{glover} & computed, (d6), (e) & \texttt{equal}, \texttt{differs\_by:<d>}, \texttt{class\_level:<ids>}, \texttt{not\_evaluated}, \texttt{not\_applicable} \\
  \texttt{final}, \texttt{final\_reason}, \texttt{certificate}, \texttt{certificate\_routes} & record & verbatim; (d7) \\
  \texttt{sturm\_B\_prime}, \texttt{sturm\_cap}, \texttt{sturm\_proof\_range}, \texttt{sturm\_outcome} & record \texttt{sturm.*} & \texttt{not\_applicable} when the record has no \texttt{sturm} object \\
  \texttt{certified} & derived, (d7) & \texttt{true}, \texttt{heuristic} \\
  \texttt{extremal} & derived, (g) & a JSON list of tags, possibly empty \\
  \bottomrule
  \end{longtable}
  \endgroup
\item \textbf{The rules.}
  \begin{enumerate}[label=(d\arabic*)]
  \item \emph{The boundary split} (the one computed classification; slot-11 WP-7,
    \texttt{HeckeExt/HeckeCensusSplit}). For a row whose own group is \texttt{H1}: $E_\m\subseteq
    H^1(\Ga;\Mn)$ its generalised eigenspace over $\Fl$ on the unsplit $H^1$ (every published record
    has \texttt{sign = "all"}), $\dim E_\m=\texttt{dim\_E}$, and
    $\texttt{rank\_r1}:=\rank_{\Fl}\big(r^1|_{E_\m}\big)$ with $r^1:H^1(\Ga;\Mn)\to H^1(\partial;\Mn)$
    the restriction \texttt{HeckeBoundary} builds (its joint-eigenspace use at
    \texttt{HeckeBoundary.h:293} is the precedent). Verdict: \texttt{eisenstein} iff
    $\texttt{rank\_r1}=\dim E_\m$; \texttt{cuspidal} iff $\texttt{rank\_r1}=0$; \texttt{mixed} iff
    $0<\texttt{rank\_r1}<\dim E_\m$. A row whose own group is \texttt{H0}: \texttt{not\_applicable}
    (its $H^1$ part, if any, is its own row). \emph{Where computed}: every cell with $\ell\le13$, and
    the two $691$-pattern cells $(10,691)$, $(198,691)$; at every other $\ell=691$ cell the verdict
    is \texttt{undecided}, and the row keeps its \texttt{kummer\_tag}. \emph{What the verdict
    means}: it is a mod-$\ell$ statement about $\ker r^1=H^1_!(\Fl)$, not about reductions of
    characteristic-zero forms --- \S\ref{sec:tests} records that for $\ell\le n+1$ ``$H^1_!$ is the
    cuspidal part'' fails in both directions, so \texttt{cuspidal} means ``inside $\ker r^1$'' and
    \texttt{eisenstein} ``injects into the boundary''. \emph{Expected at the benchmark}: at
    $(10,691)$ the Eisenstein ideal is \texttt{mixed}, $\dim E_\m=3$, $\texttt{rank\_r1}=1$
    (Claim~\ref{cl:ramanujan}; \S\ref{sec:tests}, the ten Kummer cells: $\rank r^1=1$ on the
    block). The owner's question --- an Eisenstein-labelled ideal that is \texttt{mixed} --- is the
    conjunction $\texttt{label\_type}=\texttt{eisenstein}$ (at $\ell\le13$) or a confirmed
    \texttt{kummer\_tag} (at $\ell=691$) with $\texttt{boundary\_split}=\texttt{mixed}$.
    \emph{Invariants} (per cell, asserted by the split driver): $\sum\dim E_\m=\dim H^1$ over the
    \texttt{H1} rows; $\sum\texttt{rank\_r1}=\rank r^1=\dim H^1-\dim H^1_!$.
  \item \emph{$\pm$.} \texttt{sign} is the record's, verbatim (\texttt{all} in every published
    record, so it carries no information there). The owner's $\pm$ content is
    $\texttt{dim\_plus}:=\dim(E_\m\cap H^1_+)$ and $\texttt{dim\_minus}:=\dim(E_\m\cap H^1_-)$ for the
    production involution $\iota^\ast$ (\texttt{HeckeSignSplit}; $\iota^\ast$ commutes with the
    Hecke family, so $E_\m$ is $\iota^\ast$-stable), computed by the split driver with the verdict
    of~(d1) and where it is; asserted $\texttt{dim\_plus}+\texttt{dim\_minus}=\dim E_\m$ and
    $\sum\texttt{dim\_plus}=\dim H^1_+$ of the record's \texttt{groups}. \texttt{not\_computed}
    where (d1) is \texttt{undecided}; \texttt{not\_applicable} on \texttt{H0} rows.
  \item \emph{Algebraic columns} (the keys [CEN]'s \texttt{census.py} reads).
    $\texttt{length\_A}:=\texttt{algebra.dim}/\texttt{residue\_degree}$, the length of the local
    algebra $A$ over itself; $\texttt{loewy\_A}:=\texttt{algebra.nilpotency}$, its Loewy length;
    $\texttt{reduced}:=(\texttt{loewy\_A}=1)$, equivalently $A$ a field;
    \texttt{embdim}, \texttt{gorenstein}, \texttt{hilbert} verbatim. On the module:
    $\texttt{module\_D}:=\texttt{D}$, $\texttt{module\_length}:=\texttt{D}/f$ (the semisimple
    multiplicity of the ideal's system in its own group; it may exceed \texttt{length\_A}, the
    module being neither cyclic nor free in general), $\texttt{module\_loewy}:=\texttt{loewy}$,
    $\texttt{module\_loewy\_length}$ its number of layers. ``Length'' in the owner's list is
    \texttt{length\_A}; \texttt{dim\_E} is the $\Fl$-dimension of the ideal's block, never
    called a length.
  \item \emph{$\theta$-image.} $\texttt{theta\_image\_dim}:=d_\theta(\m)$ of
    \S\ref{sec:thetaobj}~(f) (cell target-indexed: the image of $\theta_{n-\ell-1}$);
    \texttt{theta\_image} is \texttt{none} iff $d_\theta=0$, \texttt{all} iff $d_\theta=\dim E_\m$,
    \texttt{partial} otherwise. $d_\theta=0$ by definition at $n<\ell+1$ (every $\ell=691$ row), with
    source \texttt{n\_lt\_ell\_plus\_1}. \texttt{not\_applicable} on \texttt{H0} rows (the
    $H^0$-level image of $\theta^0$ is the Dickson ring's $L\cdot$, (6d), not a census column).
    \texttt{not\_computed} where no \texttt{THETA\_IMAGE} line exists (e.g.\ a second-edition cell
    WP-6's producer has not run on).
  \item \emph{$\Z/\ell^k$ depth} --- the exponents of the $\ell$-primary torsion of
    $H^1(\Ga;\Sym^n\Z^2)$ localised at~$\m$. By \HA{cl:bockstein}(2)--(3) (Hecke-equivariant
    Bockstein, $\ell\ge5$, $n>0$) that torsion is non-zero at $\m$ iff $H^0(\Ga;\Mn)_\m\ne0$,
    i.e.\ iff the row's own group is \texttt{H0} or its \texttt{groups} list carries an \texttt{H0}
    entry with \texttt{matched = true}. Values: \texttt{0} (source \texttt{bockstein}) for every
    other row --- at every $\ell\ge5$, $n>0$, surveyed or not; \texttt{not\_applicable} at $n=0$
    ($H^0(\Ga;\Z)=\Z$ is not torsion). For a row meeting $H^0$ at a cell of
    \texttt{ring\_survey\_l2} ($\Z/\ell^2$, even $n\le64$): the $H^0$ invariant factors give each
    cyclic summand's exponent as \texttt{1} (factor~$1$) or \texttt{>=2} (factor~$2$; the
    \texttt{higher\_torsion\_count} summands), refined to the exact exponent at the E8 points over
    $\Z/5^3$ by \S\ref{sec:ring}~(e)'s table ($(30,5)$: $(2)$; $(60,5)$: $(2,1)$), which the tool
    carries as a frozen table citing that site and never recomputes (source \texttt{e8\_point}; at a
    ring level $k$ an exponent is exact below~$k$ and written \texttt{>=k} at~$k$). The multiset is
    \emph{attributed to the row} iff the cell has exactly one \texttt{H0}-own row (the row and every
    row matched to it carry it) or all the cell's exponents are equal (each $H^0$-meeting row
    carries the common value); otherwise every $H^0$-meeting row reads
    \texttt{cell\_level:\{<multiset>\}} (source \texttt{ring\_survey\_l2}). At an unsurveyed cell a
    row meeting $H^0$ reads \texttt{>=1} with source \texttt{not\_surveyed}. Expected in the first
    edition (from \texttt{ring\_survey\_l2}'s analysis): $\ell^2$-torsion exactly at $(30,5)$,
    $(60,5)$, $(56,7)$.
  \item \emph{Glover reconciliation} (Claim~\ref{cl:glover}; identity~6 of
    \S\ref{sec:identities}). For an \texttt{H1} row at $\ell\in\{5,7,11,13\}$ (the in-process
    atlas range): the ideal's residue class $c$ (one Galois orbit of anemic systems over the
    canonical split field, as \texttt{HeckeIdentityCertificate} item~6 defines it);
    $\texttt{glover\_multiplicity}:=\dim_{\Fl}(E_\m)/r$, the semisimple multiplicity of $c$ in $H^1$;
    $\texttt{glover\_main}:=$ the coefficient of $c$ in $\sum_i\theta^{j_i}[H^1(M_{n_i})]$ over the
    Glover factors $\Sym^{n_i}\otimes\det^{j_i}$ of \texttt{glover\_factors(n, ell)} (the record's
    \texttt{hecke\_glover\_factors/v1}, re-derived by the certificate's own recursion, standard-module
    convention); $\texttt{glover\_correction}:=$ the coefficient of $c$ in the $H^0$ (and $H^2$)
    corrections of identity~6. \texttt{glover} is \texttt{equal} iff
    $\texttt{glover\_multiplicity}=\texttt{glover\_main}$, else \texttt{differs\_by:<d>} with
    $d=\texttt{glover\_multiplicity}-\texttt{glover\_main}$; identity~6 then requires
    $d=\texttt{glover\_correction}$, and a row where it does not is a \textbf{finding} (tag
    \texttt{identity\_6\_mismatch}), never re-fitted. If $c$ is shared by several rows of the cell
    (ideals that differ only outside the anemic family), each reads \texttt{class\_level:<ids>} with
    the class's numbers. \texttt{not\_evaluated} at $\ell=691$ (identity~6 is not evaluated there,
    \S\ref{sec:identities} item~6); \texttt{not\_applicable} on \texttt{H0} rows. Computed by the
    split driver through \texttt{HeckeIdentityCertificate}'s item-6 evaluation, never by a second
    implementation. This column is where [CEN]'s inference ``the $\ell=5$ multiplicity-$4$ ideals
    at $n=34\dots44$ are Glover multiplicities'' is confirmed or refuted.
  \item \emph{Certification.} Every column read from the record's \texttt{algebra} or
    \texttt{modules} is certified iff $\texttt{final}=\texttt{true}$; the row's \texttt{certified}
    is \texttt{true} then and \texttt{heuristic} otherwise (every non-final published ideal has
    \texttt{final = false}, \texttt{final\_reason = "sturm\_heuristic"}), and the \texttt{sturm\_*}
    columns carry its proof range against the Sturm bound. A statement of the catalogue about a
    heuristic row (``reduced'', ``Loewy length~$2$'') is marked heuristic wherever it is made.
  \end{enumerate}
\item \textbf{The computed columns' line formats and the matching rule.} The split driver runs
  each cell's $\Fl$ pipeline and the shipped ext driver's decomposition (seconds per cell at
  $\ell\le13$), and \texttt{HeckeCensusSplitValidation} (\texttt{standard}, $\le60$\,s) asserts the
  invariants of (d1), (d2) and identity~6 on a grid, with $(10,691)$'s \texttt{mixed} among them.
  Under the read-only interlock \texttt{HECKE\_CENSUS\_SPLIT\_ALL=1} (not an emitter) it runs every
  census cell of~(d1) and prints, per \texttt{H1} row,
  \begin{center}\footnotesize
  \texttt{CENSUS\_SPLIT n=<n> ell=<ell> ideal=<id> witness=<c0,...> dim\_E=<d> rank\_r1=<r> verdict=<v> dim\_plus=<a> dim\_minus=<b> glover\_mult=<m> glover\_main=<g> glover\_corr=<c> glover=<token>}
  \end{center}
  and the $\theta$ producer prints the \texttt{THETA\_IMAGE} lines of \S\ref{sec:thetaobj}~(f).
  \emph{Matching}: a run's block is the record's ideal iff its \texttt{witness\_poly} and its
  restricted \texttt{t\_p} matrices equal the record's (the pipeline and driver are deterministic,
  so equality is exact); the matching must be one-to-one on every cell, and a block or ideal left
  unmatched is a \textbf{finding} that stops the cell, never a nearest guess.
\item \textbf{The category table.} The tool reproduces [CEN]'s category table (ideals, reduced,
  non-reduced, maximal nilpotency, maximal \texttt{length\_A}, maximal \texttt{embdim},
  non-Gorenstein, $f>1$, maximal reduced module multiplicity, not final, per~$\ell$) \emph{exactly}
  from the published records --- the acceptance test of the tool --- and adds the
  \texttt{boundary\_split} and \texttt{theta\_image} counts per~$\ell$ below it.
\item \textbf{Extremal selection} (the list [EX] selects its examples from; a row carries every tag
  it meets; a heuristic row's tag carries the suffix \texttt{\_heuristic}):
  \texttt{len\_ge\_3} ($\texttt{length\_A}\ge3$); \texttt{embdim\_ge\_2};
  \texttt{non\_gorenstein}; \texttt{loewy\_ge\_3} ($\texttt{loewy\_A}\ge3$ or
  $\texttt{module\_loewy\_length}\ge3$ --- none in the first edition by [CEN], kept for the second);
  \texttt{family:<key>} for a non-reduced row whose key
  $(\ell,\texttt{length\_A},\texttt{embdim},\texttt{gorenstein},\texttt{hilbert})$ occurs at $\ge3$
  values of~$n$ (the recurring families); \texttt{reduced\_high\_mult} (\texttt{reduced} and
  $\texttt{module\_length}\ge4$); \texttt{mixed} (every \texttt{mixed} row); \texttt{pair\_691}
  (a non-reduced $\ell=691$ row with a confirmed \texttt{kummer\_tag}: $(10,691)$ and $(198,691)$ in
  the first edition, the pattern $A\cong\F[\varepsilon]/\varepsilon^2$, $\dim E_\m=3$, module Loewy
  layers $[2,1]$); \texttt{theta\_partial} ($0<d_\theta<\dim E_\m$); \texttt{zl\_depth\_ge\_2};
  \texttt{identity\_6\_mismatch}. The selection is the union; [EX] chooses among it and says why.
\item \textbf{The \texttt{results/} layout}, published owner-signed as the directory
  \texttt{congruence\_census/} of \texttt{docs/matrix-engine/results/}:
  \begin{center}\footnotesize
  \begin{tabular}{@{}>{\raggedright\arraybackslash}p{4.3cm}>{\raggedright\arraybackslash}p{9.6cm}@{}}
  \texttt{README.md} & the index of editions (the one file a later edition edits) \\
  \texttt{edition1/README.md} & the format tag, the edition, the inputs with each directory's \texttt{index.json} sha256, (a)'s rule, the tool's version and command line, the binary's commit, the owner's sign-off line \\
  \texttt{edition1/census.tsv} & every row, the columns of~(c) in order; first line \texttt{\# schema=hecke\_census/v1 edition=1 \dots} \\
  \texttt{edition1/extremal.tsv} & the rows with a non-empty \texttt{extremal}, same columns \\
  \texttt{edition1/category\_table.txt} & (f) \\
  \texttt{edition1/inputs.sha256} & one line per record read (\texttt{<sha256>  <relpath>}), superseded ones marked \\
  \texttt{edition1/split\_lines.txt} & the \texttt{CENSUS\_SPLIT} lines, verbatim \\
  \texttt{edition1/theta\_image.tsv} & \texttt{hecke\_census\_theta/v1}: the \texttt{THETA\_IMAGE} lines as a table \\
  \texttt{edition1/findings.txt} & the answers of slot-11 WP-7 scope~5, each with its rows \\
  \texttt{edition2/} & the same files over edition~1's inputs and WP-8's directory \\
  \end{tabular}
  \end{center}
  A later edition never rewrites an earlier one's files. The tool is \texttt{congruence\_census.py}
  under \texttt{docs/matrix-engine/tools/} (WP-7; WP-8 appends its second-edition mode), read-only
  on every input; it writes only into the edition directory it is given.
\item \textbf{Non-default settings.} A record computed under a non-default mathematical setting
  (WP-8's \texttt{ell = 5} cells under a raised \texttt{sturm}) is \emph{not golden-comparable}:
  its \texttt{settings.fingerprint} is non-zero, and the edition's \texttt{README.md} says so and
  names the fingerprint. The catalogue carries the record's fingerprint in the provenance, not as
  a column.
\item \textbf{Refusal.} The split driver and the $\theta$ producer are $\Fl$-only
  (\S\ref{sec:ring}~(b), the v1.12.0-pre.1 paragraph); the census tool reads only
  $\Fl$ records and stops on a record whose \texttt{coefficient\_ring} is not \texttt{F\_ell}.
\end{enumerate}
\emph{v1.12.0, as built} (slot-11 WP-7, WP-8). The tool, \texttt{HeckeCensusSplit} and the first edition
are built to (a)--(k), with these differences, each recorded in the edition's \texttt{README.md}:
\begin{itemize}
\item \textbf{(e)} \texttt{CENSUS\_SPLIT} carries one field after the ones listed, \texttt{t\_p=<fnv64>}:
  the FNV-1a-64 of the canonical text \texttt{"<p>:[[a,b],\ldots];\ldots"} of the restricted operators,
  ascending in~$p$. The tool recomputes it from the published record, which makes~(e)'s matching exact
  without printing matrices; its negative controls (one changed entry, one changed prime) move it.
\item \textbf{(f)} The acceptance is \emph{cell-exact} (the $55$ cells of [CEN]'s table compared with
  whitespace stripped, $55/55$, and its two summary lines by value), not byte-exact: [CEN]'s padding is
  hand-made. Two negative controls stop the tool.
\item \textbf{(c), (d)} \texttt{kummer\_tag} is written \texttt{eisenstein\_cuspidal:k=<k>:confirmed} (or
  \texttt{:unconfirmed}); at the $\ell=691$ cells other than $(10,691)$ and $(198,691)$ the computed
  columns read \texttt{not\_computed} and the split \texttt{undecided}. (d6)'s Glover classes run over the
  cell's anemic primes $\{2,3,5,7\}\smallsetminus\{\ell\}$ ($S_{\mathrm{sep}}(\ell)$ lies inside it at every
  $\ell\le13$); a class shared by several rows is checked against $\sum\dim E_\m=\text{mult}\cdot f$ by the
  review on all $298$ classes of edition~1, but not asserted in the code (a carry, slot-11 WP-14 pass~2,
  m-6). \texttt{HeckeCensusSplitValidation} is \texttt{standard} ($117$ checks); under
  \texttt{HECKE\_CENSUS\_SPLIT\_ALL=1} it runs the edition-1 grid ($406$ cells, $5643$ checks).
\item \textbf{(i)} Edition~1 is published: \texttt{docs/matrix-engine/results/congruence\_census/edition1/},
  $2569$ rows, owner-signed on 2026-09-25 (commit \texttt{bb7c860}; the edition's \texttt{README.md} line
  still reads ``pending'' and is the slot close's to update). \textbf{Edition~2 is not built in slot~11}:
  the owner moved it to slot~12 (slot-11 plan \S10, ruling~(i)). In its place slot-11 WP-8 published the
  W2 campaigns \texttt{ext\_census\_l7l11} ($112$ cells, $\ell\in\{7,11\}$, $n\le300$) and
  \texttt{ext\_sturm\_l5} ($101$ cells, $\ell=5$ under a raised \texttt{sturm}, fingerprint
  \texttt{29cc925d211962b3}, \emph{not golden-comparable}, (j)) and a \emph{data set}, not an edition:
  \texttt{congruence\_census/w2\_dataset/}, format tag \texttt{hecke\_census\_dataset/v1} (the tool's
  dataset mode: the record-derived columns only; the split, $\pm$, $\theta$-image and Glover columns read
  \texttt{not\_computed}\,/\,\texttt{not\_evaluated}), owner-signed on 2026-09-25 (commit \texttt{63f2191}).
  \texttt{HECKE\_CENSUS\_SPLIT\_ALL} hard-codes the edition-1 grid; widening it is edition~2's code change.
\end{itemize}

\noindent\emph{A prediction read off edition~2} (v1.14.0-pre.2): on every non-reduced $H^1$ row the local algebra
has $\m_A^2=0$ --- Prediction~\ref{pr:latticesplit} (\CONJ), tested on the lattice mod~$\ell^2$ by slot-13 WP-9
(v1.14.0-pre.3: on non-Dickson rows only; on a non-reduced Dickson-type block it is untested; v1.15.0: a third route and
slot-14 WP-6's $348$ cells, no non-reduced Dickson-type block among them; v1.15.1-pre.7: \textbf{refuted as stated} at $\ell=37$
beyond the census window, on Kummer--Frobenius ideals, and restated with their exclusion --- every edition-2 row lies below every
Kummer--Frobenius onset; v1.17.0-pre.2: the v1.16.0 form refuted on cohomology at $(608,547)$, beyond every edition, and restated again on Kummer-free classes;
v1.17.0-pre.4, owner ruling on the referee F1's B-1: that form contradicted on the forms side at irreducible congruence classes, the regular $\ell=139$ among them,
and (R) restated on congruence-free classes; the rows of editions~2 and~3 unaffected).

\noindent\emph{Edition~2 reproduced on a second machine} (v1.15.0, provenance): the first cluster run (job \texttt{34550470}, commit
\texttt{07f413f}; \texttt{docs/matrix-engine/results/noctua2\_run1/cpu2\_arms.txt}, published 2026-09-28) recomputed edition~2's
\texttt{CENSUS\_SPLIT} lines (the union of its two split cases is exactly edition~2's $2568$-line \texttt{split\_lines.txt}), the
$\theta$-image table of slot-11 WP-6 ($1934$ rows) and the forms census of slot-12 WP-12 ($655$ rows), each byte-identical to the
tracked reference.

\noindent\emph{The depth window beyond $n=64$} (v1.15.0, pointer): (d5)'s $\Z/\ell^{k}$ depth is surveyed at even $n\le64$
(\texttt{ring\_survey\_l2}). The first cluster run extended the survey over $\Z/\ell^2$ to even $66\le n\le300$ at $\ell\in\{7,11,13\}$ and
$66\le n\le200$ at $\ell\in\{5,691\}$ (\texttt{docs/matrix-engine/results/noctua2\_run1/depth\_window\_l2.tsv}, published 2026-09-28;
\HA{rmk:bocksteinexp}(3); [EX] \texttt{nc:depth}). No edition carries it: a \texttt{zl\_depth} column for $n>64$ would be a later edition's
(census edition~3, deferred to slot~15). (v1.16.1) Edition~3 (below) carries a \texttt{depth\_window} column from the
$\Z/\ell^2$ window at $\ell\in\{5,7,11,13,691\}$ only; at its new primes the column reads \texttt{not\_surveyed}.

\noindent\emph{Census edition~3} (v1.16.1; slot-15 WP-15, owner request of 2026-10-01; a pointer and its evidence, no statement of
this document changes). \texttt{docs/matrix-engine/results/congruence\_census/edition3/}, format tag
\texttt{hecke\_census\_dataset/v2} (edition~1's $58$ columns unchanged, then $13$ trailing columns), is published \emph{beside}
edition~2 (a sibling; editions~1 and~2 are untouched and still reproduced byte for byte, \texttt{EDITION2\_NATIVE rows=3285
identical=3285}). It has $80\,293$ rows: edition~2's $3285$ and $77\,008$ at the $21$ primes $\ell=17\ldots103$, in $4891$ cells, the
new ones all \texttt{hecke\_ext/v5} records under \texttt{u\_ell = budget\_mb=16384} (settings fingerprint \texttt{43364ef4020ac327},
so \emph{not golden-comparable}); every row is \texttt{final = true}. \textbf{Coverage, stated honestly}: the C1 grid is every even
$n\le700$ at each new prime ($7350$ cells); edition~3 is \textbf{complete at $\ell=61\ldots103$} ($350$ cells each, the C1b block,
$62\,113$ rows) and \textbf{partial at $17\ldots59$} ($n\le120$ only at $17\ldots31$; the heavy end of C1a and the local smoke at
$37\ldots59$): $3064$ of the $7350$ C1 cells are missing (\texttt{missing.txt}); they are to be added, add-only, as edition~3.1. A
statement ``at $\ell$'' about edition~3 means ``on the rows present''. Columns not computed on every row: the split columns ($604$
local-smoke cells only), \texttt{hidden\_congruence} ($3035$ new rows joined to a weight-shift block, no \texttt{true} among them),
\texttt{saturated}\,/\,\texttt{mult\_one} ($321$ new rows; arm~(T) was not run at $n>500$ and its lines were dropped on $230$ cells at
$61\ldots103$, $n\ge302$, where its default \texttt{u\_ell} leaves $U_\ell$ out --- a slot-16 repair, then edition~3.1).
\emph{What edition~3 adds as evidence} (computed, on the rows present):
\begin{itemize}
\item Prediction~\ref{pr:latticesplit}~(R): \texttt{loewy\_A}$\;\ge3$ occurs on exactly $18$ rows, all at $\ell=37$ and all
  Kummer--Frobenius, so (R) has $0$ violations; at $\ell=61\ldots103$, where the coverage is complete, all $2165$ non-reduced rows
  have Loewy length $\le2$, consistent with~(R). The C-F9 lattice test held on all $155$ evaluable cheap new-prime rows
  (\texttt{cf9.tsv}).
\item Prediction~\ref{pr:kfonset}: the onset law predicts no Loewy-$4$ cell at $\ell=61\ldots103$ inside $n\le700$, and none occurs;
  the law's positive half is untested there.
\item Kummer agreement (\texttt{kummer\_exact\_agree}): true on $55\,015$ rows of edition~3, false on none; on the complete C1b block
  $40\,713/40\,713$.
\item The k2 law of slot-14 WP-C14 I-8 (a twisted Eisenstein ideal $\operatorname{eis}(i,j)$, $0<i<j$, outside the degenerate family,
  is non-reduced $[3,1]$ iff $k\ge\ell(j-i+1)+i(\ell+1)$): $174/174$ on the unpredicted twisted-Eisenstein ideals at $61\ldots103$,
  no miss, none early (slot-15 WP-C18~\S3).
\item The rationality clause of slot-15 WP-12's Proposition~D on the dihedral $\Q(\sqrt{-\ell})$ blocks: $3$, $2$, $1$ rational
  systems per cell at $\ell=71$, $79$, $83$ ($55$, $45$, $39$ cells), and at $\ell=103$ one block of residue degree~$2$ with a
  character of order~$5$ ($28$ cells), exactly as predicted; none at $\ell\equiv1\pmod4$ or at~$67$.
\item F9: the $\operatorname{eis}(0,k-1)$ ideal is non-reduced in every $v_\ell(B_k/k)\ge1$ cell present, the depth-$2$ cells
  $(282,37)$ and $(226,103)$ included.
\end{itemize}
Report: \texttt{docs/implementation/slot15/WP-15\_report.md}; the C1b analysis \texttt{WP-C18\_noctua2\_rb\_report.md}; the
edition's \texttt{README.md} and \texttt{findings.txt} (E3-1 \ldots\ E3-8).

\noindent\emph{Census edition~3.1} (v1.17.0-pre.1; slot-16 WP-4 and WP-5, published at the slot-16 close by the owner's default ruling S16-N4; a
pointer and its evidence, no statement of this document changes). \texttt{docs/matrix-engine/results/congruence\_census/edition3/3.1/} is an
\textbf{add-only} extension of edition~3: the same $80\,293$ rows and the same $70$ non-edition columns, byte for byte (\texttt{addonly.txt}:
\texttt{ADDONLY rows=80293 \dots{} violations=0}), with \texttt{edition} $=3.1$; \texttt{saturated} and \texttt{mult\_one} filled on the $2544$
residue-degree-$1$ non-reduced rows that edition~3's arm~(T) had dropped (all $2544$ read \texttt{saturated = true}; with edition~3's, $2865$ of $2865$
new-prime rows are saturated); and three trailing columns from slot-16 WP-5's \texttt{HeckeSaturationDefectValidation} (\texttt{sat\_defect\_dim},
\texttt{sat\_defect\_loewy}, \texttt{sat\_length}: the $\dim I$, the Loewy length of $I$ and $\mathrm{length}(\T^{\mathrm{sat}}/\T)$ of
\HA{rmk:sattor}, read on a torsion-free block at a stable~$K$), \textbf{partial}: computed on the cells counted by \texttt{sat\_cells\_done} in the
edition's \texttt{FINGERPRINT.txt} ($897$ of $1459$ at its publication, 2026-10-02; $560$ when it was staged, 2026-10-01), \texttt{not\_computed:no\_log} on the other rows and
\texttt{not\_computed:torsion\_block} on Dickson-locus blocks. Edition~3's $3064$ missing C1 cells are still missing; the cluster returns become
edition~3.2. A statement ``at $\ell$'' about edition~3.1 means ``on the rows present'', as for edition~3. Report:
\texttt{docs/implementation/slot16/WP-4\_report.md}; the edition's \texttt{README.md}.

\noindent\emph{Census edition~3.2} (v1.18.0-pre.1; a pointer and its evidence only, no statement of this document changes; slot-17 WP-9,
\texttt{docs/implementation/slot17/WP-9\_report.md}). \texttt{docs/matrix-engine/results/congruence\_census/edition3/3.2/} is \textbf{staged, not published}
(owner-gated; the slot-18 plan's ruling~O7: it is completed after the cluster returns and published add-only then). It has edition~3.1's $80\,293$ rows and
$74$ columns; the only changes are the three saturation-defect columns filled on $176$ further cells (the cells $(n,\ell)$ with $546\le n\le608$ of
\texttt{satdef\_cells.tsv}), so \texttt{sat\_cells\_done=1073 of 1459}
in its \texttt{FINGERPRINT.txt}; the add-only proof against~3.1 reads \texttt{violations=0} (\texttt{addonly.txt}; $4982$ fills per column, one reason refined from
\texttt{not\_computed:no\_log} to \texttt{not\_computed:torsion\_block}). The remaining $386$ cells ($608\le n\le700$, the edition's \texttt{README.md}) read
\texttt{not\_computed:no\_log}. (v1.18.0-pre.2, the referee F1's m-E6: pre.1 read ``$n\ge608$'' for the remaining cells; the ranges are in $n$ and share
$n=608$, a cell being a pair $(n,\ell)$: the $7$ cells $(608,\ell)$ filled in~3.2 are among the $176$.) Until it is published this document cites edition~3.1.

%======================================================================
\section{First test cases and expected outcomes}\label{sec:tests}
%======================================================================

\subsection{Ramanujan: \texorpdfstring{$(n,\ell)=(10,691)$}{(n,l)=(10,691)}}
Campaign data: $\dim H^1=3$, one block $(3{:}1)$ for each of $p=2,3,5$; oracle O6:
$\min(T_2)=(X+24)^e$, $e\ge2$ \HA{sec:example}.

\begin{claim}[The local algebra at the Ramanujan ideal]\label{cl:ramanujan}\PROVED\ (given \HA{cl:bockstein}, Claim~\ref{cl:sturm} and Eichler--Shimura over $\Z_{691}$)
Let $\m$ be the unique maximal ideal of $\T_\infty$ on $H^1(\Ga;M_{10}(\F_{691}))$. Then
$\T_\m\cong\F_{691}[\varepsilon]/(\varepsilon^2)$ with
\[
T_q\;=\;\sigma_{11}(q)+c_q\,\varepsilon,\qquad
c_q\equiv\frac{\tau(q)-\sigma_{11}(q)}{691}\pmod{691},
\]
for every prime $q$, $q=\ell=691$ included ($U_{691}=T_{691}$, \HA{cor:ulrecursion}); concretely
$c_2=-3$, $c_3=-256$, $c_5\equiv-174$, $c_7\equiv-137$, and $c_{691}\equiv542$ from
$\tau(691)=-2747313442193908$, $\sigma_{11}(691)\equiv1$. Hence the \emph{matrix identities} on
all of $H^1$
\begin{gather*}
3\,(T_3-252)=256\,(T_2+24),\qquad 3\,(T_5-4830)=174\,(T_2+24),\\
3\,(T_7+16744)=137\,(T_2+24),\qquad 3\,(U_{691}-1)=-542\,(T_2+24),
\end{gather*}
every $T_q-\tau(q)$ has Jordan type $(2,1)$, embedding dimension $1$, $\T_\m$ Gorenstein,
every $T_q$ with $c_q\ne0$ is an active direction, and $H^1_\m\cong\T_\m\oplus\F_{691}$
(Loewy layers $(2,1)$, socle of dimension $2$, not cyclic; cyclic on each $\iota$-sign once
\ref{add:iota} exists: $H^1_+\cong\T_\m$ free of rank one, $H^1_-\cong\F_{691}$).
\end{claim}
\begin{proof}
$T_\Z(M_{12})\subseteq\Z\times\Z$, $T_m\mapsto(\sigma_{11}(m),\tau(m))$, is the subring generated
by $(1,1)$ and these pairs; $\sigma_{11}(m)\equiv\tau(m)\pmod{691}$ for all $m$ (Ramanujan),
and $\gcd\bigl((\tau(2)-\sigma_{11}(2))/691,(\tau(3)-\sigma_{11}(3))/691\bigr)=\gcd(3,256)=1$, so
$T_\Z(M_{12})=\{(a,b):a\equiv b\ (691)\}=:R$, a free $\Z$-module on $(1,1)$ and
$\varepsilon:=(0,691)$ with $\varepsilon^2=691\varepsilon$. $R\otimes\F_{691}=\F_{691}[\varepsilon]/(\varepsilon^2)$
is local. Since $10<690$, $H^1(\Ga;M_{10}(\Z))$ has no $691$-torsion (\HA{cl:bockstein}),
the Hecke algebra on $H^1(\Ga;M_{10}(\Z_{691}))$ is $R\otimes\Z_{691}$ (both torsion-free,
equal after $\otimes\Q$), and $\T_\m$ is its reduction, a quotient of
$\F_{691}[\varepsilon]/(\varepsilon^2)$; O6 (non-scalar $T_2$) excludes the quotient $\F_{691}$.
The coefficients: $(a,b)=a(1,1)+\frac{b-a}{691}\varepsilon$. On the $3$-dimensional $H^1_\m$,
$\varepsilon^2=0$ and $\varepsilon\ne0$ force $\rank\varepsilon=1$, whence the Jordan type,
the Loewy layers and $H^1_\m\cong\T_\m\oplus\F$; the sign split is Eichler--Shimura
($H^1_+\cong M_{12}$, $H^1_-\cong S_{12}$ \Mnote{5}).
\end{proof}
\textbf{Verified.} On the existing $A^{(1)}(T_q)$ (campaign run at $(10,691)$, $p\in\{2,3,5,7\}$,
2026-09-14) the stored matrices are $T_q=a_q I+c'_qE_{31}$ with
$(a_q,c'_q)=(667,23),(252,120),(684,643),(531,129)$, and all three anemic identities hold
exactly in $M_3(\F_{691})$; moreover $c'_q/c'_2\equiv c_q/c_2$ for $q=3,5,7$ (ratios
$546,58,276$), which is stronger. \textbf{That one line is basis-dependent, and the rest of
this paragraph is not} (v1.5.0; slot-4 WP-2 finding~6; [XA~m19] draws the distinction and
the code honours it): $a_q$, the ratios, the Jordan type, $\rank\varepsilon$ and the
$U_{691}$ identity below are invariants of the module, whereas the individual $c'_q$ --- and
with them the shape $a_qI+c'_qE_{31}$ --- are read off \emph{the stored $H^1$ basis}, a
golden \emph{tier-B} object. A re-derivation on another admissible basis reproduces every
other number here and need not reproduce the $c'_q$; a claim that it must would be a claim
about the basis choice and not about $\T_\m$. Rerun for this revision with $p=691$ (\texttt{--campaign},
\texttt{certify = none}): $A^{(1)}(U_{691})=I+221\,E_{31}$, and
$221\equiv542\cdot23\cdot(-3)^{-1}\pmod{691}$ exactly as predicted, so
$3(U_{691}-1)=-542(T_2+24)$ holds too --- the strongest single test in this document, since it
exercises $U_\ell$ and the value $\tau(691)$. The Jordan type $(2,1)$ and $\rank\varepsilon=1$
are read off directly. This is the first acceptance test of the programme, passed with no new
code.

\subsection{Two further campaign rows at \texorpdfstring{$\ell=691$}{l=691}}
\begin{center}\small
\begin{tabular}{@{}>{\raggedright\arraybackslash}p{1.6cm}>{\raggedright\arraybackslash}p{5.2cm}>{\raggedright\arraybackslash}p{8.6cm}@{}}
\toprule
$(n,\ell)$ & campaign data ($p=2,3,5$; \texttt{results\_table.tsv}) & expected \\
\midrule
$(198,691)$ & $\dim H^1=33$; blocks $2{:}1,3{:}1,4{:}2,4{:}2,20{:}10$; the $3{:}1$ block is the
Eisenstein class merged with one $\F_{691}$-rational cuspidal system ($(691,200)$ is an irregular
pair \cite{washington}) & \VERIFY\ As Claim~\ref{cl:ramanujan}, $\T_\m\cong\F[\varepsilon]/(\varepsilon^2)$,
\emph{provided} the congruence has index one: $691\,\|\,\num(B_{200}/200)$ holds (recomputed for
this revision: $|\num(B_{200}/200)|\equiv217665\pmod{691^2}$; $B_{200}<0$, and the signed numerator is
$\equiv259816$, v1.12.0). \textbf{Uniqueness: closed at v1.12.0.} Exactly one $\overline{\Q}_{691}$-eigenform
of weight~$200$ is congruent to $E_{200}$, with $\F_{691}$-rational reduction: at $(198,691)$ $H^0=0$, so
the block is torsion-free (\HA{cl:bockstein}), and $\dim E_\m=3$ with $\pm$ content $2/1$ (the census,
\S\ref{sec:census}) is one Eisenstein class plus two per congruent eigenform (Eichler--Shimura), which
leaves room for exactly one; independently, on $S_{200}$ (dimension~$16$) the characteristic polynomial of
$T_p$ modulo $691$ factors with degrees $1,1,2,2,10$ for $p=2,3,5,7$ --- the census's five blocks
$2{:}1,3{:}1,4{:}2,4{:}2,20{:}10$ --- and the joint generalised eigenspace of $1+p^{199}$ is
one-dimensional (slot-11 WP-19's reading of the record, recomputed from $q$-expansions by its referee,
WP-20, \texttt{docs/implementation/slot11/WP-20\_referee.md}; the examples document [EX] \S5.1). Through
v1.11.1 this read: the uniqueness ``is \emph{not checked} and the possible ambiguity is documented here
(v1.4.0; owner decision OQ13: uniqueness does not concern this stage --- cohomology over $\Z$ and $\Q$
with functoriality is the long-term route; an LMFDB lookup \ldots\ is a to-do \ldots; through v1.3.3
`still to be checked')''. Loewy $(2,1)$, $r^1=0$ on the
Eisenstein class in the sense of Claim~\ref{cl:harder} --- on the socle $E$, with
$\rank r^1=1$; the other blocks semisimple with $\T_\m$ a field; the $20{:}10$ block one
Galois orbit of residue degree $10$ with multiplicity $2$, cyclic per sign. \\
$(160,691)$ & $T_5$ alone has a $4{:}1$ block where $T_2$ and $T_3$ each have
$2{:}1,2{:}1,2{:}1,4{:}2$ & \textbf{Verified 2026-09-16} (v1.3.3, \texttt{HeckeLocalityProbe},
recorded, not asserted): the \texttt{ComputeMultiOp} run on $\{T_2,T_3,T_5\}$ gives seven blocks
$2{:}1,1{:}1,2{:}1,2{:}1,4{:}2,6{:}3,10{:}5$ on $\dim H^1=27$, all seven certified local
(\texttt{certify() = 0}); $T_2$'s $4{:}2$ is a residue-degree-$2$ system which certifies local
with $f=2$. This is the value \texttt{HeckeExtSeparationValidation} (E2) asserts.
\textbf{The second alternative, not the first} (v1.7.0; measured by that validation's
\texttt{SEPARATION\_SINGLE} line, slot-5 WP-8): $T_2$ alone, $T_3$ alone and $\{T_2,T_3\}$ each
give exactly the family's seven blocks, and $T_5$ alone gives the \emph{same} seven-block
partition, differing only in reporting $f=1$ on the four-dimensional block --- $T_5$'s $4{:}1$
\emph{is} the family's $4{:}2$, seen with a coarser invariant. Through v1.6.2 this row said the
first alternative holds ($T_5$'s $4{:}1$ two $\F_{691}$-rational systems merged at $a_5$ and
split by the refinement); that cannot be right, since a split would give eight blocks and every
family gives seven (the family's decomposition refines each single operator's; the validation
compares $(\dim{:}f)$ lists and block counts, not subspaces --- v1.7.1). So the block partition
at $(160,691)$, with its residue degrees, is already $T_2$'s and $T_3$'s, and only $T_5$ alone
under-reports a residue degree (Claim~\ref{cl:bk}; v1.7.0's ``no single $T_q$ certifies the
residue degrees that $\{T_2,T_3\}$ does'' is refuted by the same line). Every $\T_\m$ here is a field except where the residue degree exceeds $1$
--- the row now \emph{is} evidence for the refinement logic; Claim~\ref{cl:bk}'s isomorphism
stays recalled. \\
\bottomrule
\end{tabular}
\end{center}

\subsection{Dickson and boundary phenomena at small \texorpdfstring{$\ell$}{l}}
\begin{itemize}
\item $\ell\in\{5,7,11\}$, $n=\ell+1$: first Dickson class $L$, character $\chi_1$
  (\HA{cl:chi}); $\dim H^0=1$, $\dim H^1=\dim M_k+\dim S_k+1$ (\HA{cl:dimH1}; the
  known loci $(6,5),(12,5),(8,7)$ of \texttt{CLAUDE.md}); check its gluing in the cone once
  \ref{add:cone} exists. \VERIFY\ \emph{Recorded 2026-09-16}
  (\texttt{HeckeBoundaryParabolicValidation}, slot-2 WP-7): $\dim H^1_!=\dim H^1-\rank r^1$ with
  $\rank r^1\le\lfloor n/\ell\rfloor+1$, and the slot-2 plan's expectation ``$\dim H^1_!\ge
  2\dim S_k$ at a Dickson locus'' --- never a claim of this document --- fails on $3$ of $6$
  cells: $(6,5),(12,5),(8,7)$ give $0=2\dim S_k$, $(16,7)$ gives $1<2$, $(20,5),(24,5)$ give
  $0<2$; for $\ell\le n+1$ the classical ``$H^1_!$ is the cuspidal part, of dimension
  $2\dim S_k$'' fails in both directions ($25$ cells), with $H^1_!=0$ at
  $(10,5),(14,5),(16,5),(14,7)$ where $\rank r^1=3$; $\dim H^1_!$ is even on all $39$ cells with
  $\ell>n$.
\item $\ell=5$, $n\in\{4,8\}$: $(\ell-1)\mid n$, $\chi_0\equiv\chi_1$, the boundary characters
  collapse (Prediction~\ref{pr:cone}). $\ell=5$, $n=6$: $E_8\equiv1\pmod5$, Eisenstein character
  $1+q^{-1}$, and simultaneously the Dickson locus. \VERIFY
\item $\ell=5$, $n=34$ and $\ell=7$, $n=58$ ($k=12+\ell^2-1$): multiplicity $2$ for $\De$'s
  full system, since $\tau(5)=4830\equiv0\ (5)$ and $\tau(7)=-16744\equiv0\ (7)$
  (Claim~\ref{cl:Tl}). Measured at $(34,5)$ (inventory, Stage 9): $T_2$ and $T_5$ are each
  semisimple on $H^1$, hence so is $\F_5[T_2,T_5]$ (commuting semisimple elements over a perfect
  field generate a reduced algebra); the block is split unless some other $T_q$ acts
  non-semisimply --- expected \emph{split} \CONJ, to be settled by Stage 4.
\item $\ell=5$, $n=58$ and $\ell=7$, $n=82$ ($k=12\ell$): the $U_\ell$ self-extension of
  Claim~\ref{cl:Tl}: $T_\ell$ with a Jordan block of size $2$ at $\De$'s ideal, the anemic
  operators inactive there. \VERIFY
\item The Sturm counterexamples of Claim~\ref{cl:sturm}(b) on the pipeline's own $H^1$ (rerun
  for this revision, $p\le13$): at $(n,\ell)=(22,5)$ and $(26,7)$, $\dim H^1=5$, the $\Fl$-algebra
  generated by $A^{(1)}(T_2),A^{(1)}(T_3)$ has dimension $2$, the one generated by all
  $A^{(1)}(T_q)$, $q\le13$, has dimension $3=\dim M_k$: Certificate~\ref{cert:sturm} with
  $B(k)=3$ would have certified falsely, with $B'=5$, resp.\ $7$, it is sound. At $(32,131)$
  the algebra generated by $T_2$ alone is already everything ($B'=2$) although the $\Z$-span of
  $T_1,T_2,T_3$ is not: the certificate tests \emph{algebra} membership, so that point breaks
  the $\Z$-module statement only. \VERIFY\ (the $B'$ side is computed; the $H^1$ side measured)
\item Small-$\ell$ Kummer cases with $n<\ell-1$ (Claim~\ref{cl:harder}):
  $(n,\ell)$ one of $(14,3617)$, $(16,43867)$, $(18,283)$, $(18,617)$, $(20,131)$, $(20,593)$,
  $(22,103)$, $(22,2294797)$ (all $k\le24$) and $(24,657931)$: exactly one non-split
  Eisenstein--cuspidal $\m$ in $X_4^+$; all other $\m$ cyclic per sign. \VERIFY\ (needs E3).
  \textbf{Refuted and corrected, 2026-09-16.} Through v1.3.3 this item continued ``$r^1=0$ on
  it, $X_3$ picking up the boundary class'', and \S\ref{sec:schema}'s example drew $H^1_!=H^1$
  from it. \texttt{HeckeBoundaryParabolicValidation} (slot-2 WP-7), on all ten Kummer cells:
  $r^1$ vanishes on the socle $E$ of the block ($\rank(r^1E)=0$ --- Claim~\ref{cl:harder}'s
  Kummer direction, $10/10$; the converse $\rank(r^1E)=1$ on $27$ non-Kummer cells with
  $n<\ell-1$) but \emph{not} on the block: $\rank r^1=1$ and
  $\dim H^1_!=\dim H^1-1=2\dim S_k$ on every cell ($(10,691)$: $3/2/1$; $(22,103)$: $5/4/1$;
  the other eight $3/2/1$), because $r^1$ factors through the cosocle of the
  $\m_{\mathrm{Eis}}$-block and the class restricting non-trivially is its \emph{top}, not its
  socle. Consequently, by \HA{eq:dimbook}, $X_6=H^2_c=0$ and $\dim X_3=\dim H^1$ at a Kummer
  cell --- not $X_6\cong H^1(\partial)$ --- while $X_3$ does pick up the boundary class
  ($X_2\hookrightarrow X_3$ regardless). At $(10,5)$ and $(10,7)$, outside $n<\ell-1$, the
  criterion is not asserted: $E$ is a line with $\rank(r^1E)=1$, and $H^1_!=0$ resp.\
  $\dim H^1_!=1$ (recorded).
\item $\Dn$ versus $\Mn$ for one $n\ge\ell$ (e.g.\ $(n,\ell)=(12,5)$): the Poincar\'e-duality
  reversal of Loewy series (Prediction~\ref{pr:dual}), once $\Dn$ is a local system in the
  code. \VERIFY
\end{itemize}

%======================================================================
%======================================================================
\section{The \texorpdfstring{$L$}{L}-value interface, and its Eisenstein instance}\label{sec:lvalue}
%======================================================================
(New at v1.19.0-pre.1, hecke-extensions slot~19, WP-2; owner ruling O6 of 2026-10-03 and the owner's interface requirement of the same day,
\texttt{docs/implementation/slot18/lead\_carry.md}~C9. Placed after \S\ref{sec:tests} and before the appendices, so that no section number moves. Rule~21:
the code, slot-19 WP-LV, merges only after this text is at pre.1; its normalisations are fixed here \emph{before} any claim.) This section fixes (i) an abstract
interface for computing $L$-values, or their $\ell$-adic valuations, from the cohomology the project computes, general enough to carry later Birch's
$\GL_2$ modular symbols (linear functionals on $H^{1}$ from cycles, twisted values $L(f,\chi,1)$) and Hida's $\GL_2\times\GL_2$ pairing (a bilinear form on $H^{1}$,
$L(1,\mathrm{Ad}f)$, Hida's congruence criterion), and its own first instance, a Hecke character with its congruence module (v1.19.0-pre.2); (ii) its first implemented instance, the Eisenstein\,/\,Bernoulli one (since v1.19.0-pre.6 the \textsf{linear} and \textsf{bilinear} instances are implemented too, Remark~\ref{rmk:lvshapes}), with its exact index
statements; (iii) the normalisation table. \textbf{Not used}: Harder's route to rationality through the boundary spectral sequence (owner, 2026-10-03). Throughout
$\ell\ge5$ and $v=v_\ell$; [U] is \texttt{docs/latex/hida\_hecke\_algebra.tex}, cited as \Ucite{label}.

\subsection{The interface}\label{sec:lviface}

\begin{definition}[An $L$-value route]\label{def:lvroute}
An \emph{$L$-value route} is a tuple $\mathcal R=(D,\kappa,\mathcal N,\mathrm{ev},\mathcal I)$:
\begin{enumerate}
\item \textbf{Domain} $D$: a named model of [C] and a coefficient ring $\mathsf R\in\{\Z_\ell,\Z/\ell^{m}\}$ --- $H^{1}(\Ga;\Mn\otimes\mathsf R)$ at level one, $H^{1}_c$ (the cone of
  \HA{sec:cone}), or the character model $H^{1}(\Ga_0(N);V_{\mathsf R}(\chi))$ of \HA{def:nebmodel} --- with its Hecke operators, $\iota^{\ast}$ (\HA{def:iota}; ``$+$'' is the
  Eisenstein side) and a named Hecke-stable subquotient $H'$ (a sign, an $\m$-part, the cuspidal part) on which the route lives.
\item \textbf{Kind} $\kappa\in\{\textsf{linear},\textsf{bilinear},\textsf{congruence}\}$ (the third added at v1.19.0-pre.2, referee F1 M-2):
  \textsf{linear} --- a $\mathsf R$-linear functional $\lambda\colon H'\to\mathsf R$ (Birch's modular symbols, (B) of Remark~\ref{rmk:lvshapes});
  \textsf{bilinear} --- a $\mathsf R$-bilinear pairing $\langle\,,\rangle\colon H'\times H''\to\mathsf R$ with a stated adjunction of the Hecke operators
  ($\langle T_px,y\rangle=\langle x,T_p^{\ast}y\rangle$, $T_p^{\ast}$ named) (Hida's pairing, (H));
  \textsf{congruence} --- a \emph{Hecke character} $e=(e_p)_p$, i.e.\ a system of Hecke eigenvalues given by its values $e_p\in\mathcal O$ on the $T_p$ (a ring homomorphism
  from the abstract Hecke algebra $\mathsf R[T_p:p]$ to $\mathcal O$, $\mathcal O$ the coefficient ring $\mathsf R$ or a named finite extension of it; $e$ need \emph{not} occur in $H'$),
  together with the Hecke algebra $\T(H')$ (the $\mathsf R$-subalgebra of $\mathrm{End}(H')$ generated by the $T_p$), the ideal $I_e\subset\T(H')_{\mathcal O}:=\T(H')\otimes_{\mathsf R}\mathcal O$
  generated by the $T_p-e_p$, and the \emph{congruence module}
  \[
    Q_e\;:=\;\T(H')_{\mathcal O}\big/\bigl(I_e+\mathrm{Ann}_{\T(H')_{\mathcal O}}(I_e)\bigr).
  \]
  When $e$ does not occur in $H'\otimes\Q_\ell$ and $\T(H')$ is torsion-free, $\mathrm{Ann}(I_e)=0$ and $Q_e=\T(H')_{\mathcal O}/I_e$ (the Eisenstein instance (E)); when
  $e=\lambda_f$ is the system of an eigenform $f$ occurring in $H'$, $I_e=\ker\lambda_f$ and $Q_e\cong\mathcal O/\lambda_f(\mathrm{Ann}(\ker\lambda_f))$ is Hida's congruence module of $f$.
\item \textbf{Normalisation record} $\mathcal N$: the $L$-function; the point $s_0$; the twist character $\chi$ (conductor, values); the sign at infinity
  $\varepsilon\in\{+,-\}$ (the $\iota$-eigenvalue); the Euler factors removed, each as an exact rational at $s_0$; the period convention (none; $\Omega_f^{\pm}$; the
  Petersson norm); every normalising constant (Gauss sums, powers of $2\pi i$, factorials) as an exact rational or an algebraic number with its $\ell$-valuation;
  and the \emph{ambiguity} allowed --- an index statement is an equality of ideals of $\Z_\ell$ (of $\ell$-valuations), never of $\ell$-adic units unless the
  record says so.
\item \textbf{Evaluation} $\mathrm{ev}$: $\lambda(c)$, or $\langle c,c'\rangle$, for named classes, a value in $\mathsf R$; for the \textsf{congruence} kind, the exact values
  $e_p$ for $p$ in a named finite set $S$ (with a certificate that the $T_p$, $p\in S$, generate $\T(H')$, \S\ref{sec:certgen}) and the finite module $Q_e$ (its
  elementary divisors).
\item \textbf{Index statement} $\mathcal I$: an assertion $\mathrm{length}_{\Z_\ell}(Q)=v(L^{\mathrm{alg}})$, $Q$ a named finite module computed from $D$ (a congruence
  module, the cokernel of $\lambda$ on a lattice, the discriminant of $\langle\,,\rangle$, the module $Q_e$ of a \textsf{congruence} route), $L^{\mathrm{alg}}\in\Q$ (or in a number field, with a place over $\ell$)
  given exactly by $\mathcal N$; over $\mathsf R=\Z/\ell^{m}$ it reads $\mathrm{length}(Q\otimes\Z/\ell^{m})=\min(m,v(L^{\mathrm{alg}}))$ when $Q$ is cyclic. Every
  index statement carries its own status tag.
\end{enumerate}
An instance fixes all five before anything is computed (rule~20); the code carries $\mathcal N$ as data and prints it with every value.
\end{definition}

\begin{remark}[The three shapes]\label{rmk:lvshapes}
\SPEC{} (E) \emph{Eisenstein\,/\,Bernoulli}, \textsf{congruence} (ring form; its module form is a \textsf{linear} reading on $H^{+}$), \S\ref{sec:lveis}: implemented by WP-LV. (B) \emph{Birch / Manin--Stevens modular symbols}, linear: $D=H^{1}_c$ (the
modular symbols of weight $k$ are $H^{1}_c(\Ga;\Mn)$, \cite{ash_stevens}), $\lambda_{\chi,j}(\varphi)=\sum_{a\bmod m}\bar\chi(a)\,\langle\varphi(\{a/m\}-\{\infty\}),\,
P_{j,a}\rangle$ with $P_{j,a}:=\rho_n^{\vee}(\operatorname{adj}B_{a,m})\,u^{\vee}_{n-j}\in D_n$ (the $u_{n-j}$-coordinate of $\rho_n(B_{a,m})\Phi(\{a/m,\infty\})$, \HA{prop:pairing}(3); R19-8~\S6; F1-delta m-E2) (at pre.1 written as a polynomial $(mX-aY)^{j}$, of the wrong degree; F1 m-E5 written $(mX-aY)^{j}$, of the wrong degree; F1 m-E5), giving the critical values $L(f,\chi,j+1)$, $0\le j\le k-2$, divided by
$\Omega_f^{\pm}$, in the normalisation $L(f,\chi,j+1)/\Omega_f^{\pm}=\chi(-1)\,(-2\pi i)^{j+1}\tau(\chi)\,\lambda_{\chi,j}/(j!\,\binom{n}{j}\,m^{j+1})$ \cite{mtt} (\S8; v1.19.0-pre.4, slot-19 R19-8 \S1.4 as confirmed by isign; v1.19.0-pre.5: the formulas (8.3), (8.6), p.~10, read on the page by R19-8, the constant derived there from them). Here $\pm$ is the \emph{geometric} sign $\chi(-1)(-1)^{j}$, which in the convention of [C] is $-\chi(-1)(-1)^{j}$; the earlier constant was off by $\binom{n}{j}\tau(\bar\chi)/(\chi(-1)\tau(\chi))$ (R19-8 tested only quadratic and trivial $\chi$; for non-real $\chi$ the factor is derived, not tested; v1.21.0-pre.3: for every primitive $\chi$ of conductor $25$, $49$, $121$ at $k=12$ and $25$, $49$ at $k=16$,
non-real ones included, \CMPD{} by slot-21 P4, Remark~\ref{rmk:lvwild}(1); for non-real $\chi$ of any other conductor (a prime, or not a prime power) still derived, not tested). The functional $\lambda_{\chi,j}$ lives on $H^{1}_c$ and descends to $H^{1}_\m$ only at a residually non-Eisenstein $\m$.
\emph{Prototype of~(B)} (v1.19.0-pre.5; slot-19 R19-8, \texttt{docs/implementation/slot19/R19-8\_report.md} with its errata, Python, in-house, predictions committed first; not
code of this project): on [C]'s model a class $(c,c_\partial)$ of $H^{1}_c=H^{0}(\mathrm{Cone})$ has the modular symbol $\Phi(\{0,\infty\})=(\rho_n(S)-1)\bigl(c(v_0v_4)+c(v_4v_6)-c_\partial(v_0)\bigr)$,
well defined, equivariant, zero on coboundaries (\PROVED{} there), and $H^{1}_c(\text{model})\to\{\text{modular symbols}\}$ is an isomorphism at $8$ cells over $\F_p$ (\CMPD); the
normalisation above matches the numerically computed $L(f,\chi,j+1)$ at $945$ critical values (the weights $k\le26$ with $\dim S_k=1$, $9$ twists, all $j$) to $58$ digits
(\CMPD; trivial and quadratic $\chi$ only); leg~(c) in modular-symbol form holds at the $8$ Eisenstein cells with $\dim S_k=1$, each to depth exactly $v(B_k/k)=1$ (\CMPD); at
$(258,37)$ the free sign of $H^{1}_c$ is the geometric-even one (Claim~\ref{cl:258lattice}) and the critical values of the four eigenforms agree modulo~$\pi$ and not
modulo~$\pi^{2}$ (\CMPD). The instance (\texttt{BirchModularSymbolRoute}; gates G1--G5; index statements (L1) \CMPD{} at depth~$1$, \CONJ{} at depth $\ge2$ (v1.20.0-pre.6: \CMPD{} at depth~$2$ at $(282,37)$, $(226,103)$, slot-20 WP-LV3,
Claim~\ref{cl:lvtriangle}), (L2) \PROVED{}
from (L1), (L3) \CMPD{} at $(258,37)$) was specified in R19-8~\S6 for slot~20 and is \emph{implemented} in slot~19 (v1.19.0-pre.6, WP-LIN): on the C++ cone's $H^{1}_c$
(the field cone over $\F_\ell$, the module cone over $\Z/\ell^{m}$) $\Psi(c,c_\partial):=\Phi(\{0,\infty\})$ is an isomorphism $H^{1}_c(\mathsf R)\cong W(\mathsf R)=\ker(1+\rho_n(S))\cap\ker(1+\rho_n(U)+\rho_n(U)^{2})$
--- lengths and invariant factors, including the torsion summand at $(72,37)$ over $\Z/37^{3}$ --- at $7$ cells over $\F_\ell$ and $\Z/\ell^{3}$ and at $(258,37)$ over $\F_{37}$ (and $\Z/37^{2}$,
env arm) (\CMPD); [C]'s $T_p$ on $H^{1}_c$ \emph{is} the symbol $T_p$, $\Psi\circ T^{c}_p=T^{\mathrm{sym}}_p\circ\Psi$ ($p=2,3$, over $\Z/\ell^{2}$; no adjoint, no scalar), and
$\Psi\circ\iota^{\ast}_{[C]}=-\rho_n(\iota)\circ\Psi$ read through $H^{1}_c\to H^{1}$ --- a geometric-even symbol is [C]-minus (\CMPD); \emph{torsion} (v1.21.0-pre.2; slot-21 F4 on [U] part~4,
S2, S3): $W(\Z/\ell^{R})$ need not be free --- $H^{2}_c(\Ga;\Mn\otimes\Z_\ell)[\ell^{\infty}]\neq0$ iff $(\Dn\otimes\F_\ell)^{\Ga}\neq0$, e.g.\ at $(72,37)$ --- so slot-21 P8 reads the $\m$-blocks, at
$\m_{\mathrm{Eis}}$ on \emph{both} signs, as the image of $W(\Z/\ell^{R+d})$ in $W(\Z/\ell^{R})$, which kills torsion of exponent $\le d$ (a reading convention of the code, \SPEC; no statement here changes); the sign rule (the functional vanishes on
the sign opposite to $\chi(-1)(-1)^{j}$, geometric) holds also for three non-real $\chi$ (\CMPD); R19-8's frozen numbers (the eigenlines $r^{\pm}$ at $k=12,\dots,26$, $108$ twisted rows)
are reproduced by the C++ symbols modulo $\ell^{2}$ (\CMPD); leg~(c), Remark~\ref{rmk:lvlegc}. (H) \emph{Hida's
pairing}, bilinear: $D=H^{1}_!\otimes\Z_\ell$ (or $\Z/\ell^{m}$) for $n<\ell$ (so that $\Mn$ is self-dual, \HA{prop:pairing}), $\langle x,y\rangle=(x\cup\iota^{\ast}y)[\Ga\backslash\mathbb H]$
with $T_p^{\ast}=T_p$; for a newform $f$, $\langle\delta_f^{+},\delta_f^{-}\rangle$ for integral generators of the $f$-eigenlines of the two signs is, up to a unit, the
congruence number of $f$, and equals $L(1,\mathrm{Ad}f)/(\Omega_f^{+}\Omega_f^{-})$ times explicit factors (Hida's congruence criterion, \cite{hida_cong}). The factors
(v1.19.0-pre.5; slot-19 R19-9, \texttt{docs/implementation/slot19/R19-9\_report.md}~\S2): at level one, with $r_f^{\pm}=\Omega^{\pm}\delta^{\pm}$,
\[
  L(k,\Sym^{2}f)\big/(\Omega^{+}\Omega^{-})\;=\;-(-1)^{k/2}\,i\,\frac{(2\pi)^{k+1}}{(k-1)!}\,\langle\delta^{+},\delta^{-}\rangle_{!},
\]
$L(k,\Sym^{2}f)=L(1,\mathrm{Ad}f)$ in the unitary normalisation, \emph{no Euler factor removed}: \STD{} there (\cite[Thm.~5.1, pp.~247--248, and (3.2), p.~237]{hida_cong} at
$N=1$, read on the page by R19-9) and the identity $\langle f,f\rangle=-2(2i)^{1-k}\Omega^{+}\Omega^{-}\langle\delta^{+},\delta^{-}\rangle_{!}$ (\CMPD{} at the $10$ level-one eigenforms of
weights $12$--$28$, to $2.4\cdot10^{-15}$; the Haberland\,/\,Kohnen--Zagier form of it is \STD, recalled, not read, so the constant is \STD{} (Hida, read; Haberland, recalled) and \CMPD{} (10 forms), not \PROVED; F1-delta M-1). For $\ell\ge k$ the constant is an $\ell$-unit.
\emph{Prototype of~(H)} (v1.19.0-pre.5; R19-9, Python, in-house, predictions committed first; not code of this project): the pairing $\langle\Phi,c\rangle:=\langle\Phi(\{0,\infty\}),y_c\rangle$,
$y_c$ the potential of $c\in H^{1}(\Ga;D_n)$, is the Poincar\'e pairing $H^{1}_c(\Ga;V_n)\times H^{1}(\Ga;D_n)\to\Z/\ell^{m}$ for every even $n$ and $\ell\ge5$ (\PROVED{} there); on [C]'s model
it equals $-1$ times the cone-level cup form of \texttt{Hecke/HeckePairing.h} (\CMPD, $9$ cells, $n\ge\ell$ among them), with $T_p^{\ast}=T_p$ ($p=2,3$, and $p=5$: slot-19 WP-BIL, \CMPD{} at the same cells, predicted from $\Ga\alpha\Ga=\Ga\,\mathrm{adj}(\alpha)\Ga$ at level one) and $\iota$ anti-isometric
(\CMPD, $27$ cells and $(258,37)$). Its value on integral eigenlines is the length of the congruence module \emph{of the lattice it is computed on} (a lemma, \PROVED{} there): on
Hida's $H^{1}_{!}$ it sees exactly the cusp--cusp congruences ($144169$ at $k=24$; $131$, $139$ at $k=28$) and none of the Eisenstein ones, on the mixed lattice
$H^{1}_c(V)\times H^{1}(D)$ it adds $v(B_k/k)=1$ at the $8$ Eisenstein cells on the sign carrying the boundary class (\CMPD). At $(258,37)$ ($n\ge\ell$; modulo $37^{10}$) the mixed
value of each of the four eigenforms is $v_\pi=3$ on the free sign and $2$ on the $\m$-sign, and $3$ on both through $H^{1}_{!}(D)$ (index $37$ in $H^{1}(D)_\m$ per sign) (\CMPD):
$v_{37}=3/2$ per eigenform, not the pairwise depth $1/2$, and $\sum_fv_{37}(\eta_f)=6=v(\mathrm{disc}\,P)$ --- the pairing reads the discriminant, not $\tilde\T/\T_\m$. The mixed
reading at depth~$2$ (e.g.\ $(282,37)$) was \CONJ{} through v1.20.0-pre.5 (v1.20.0-pre.1: Claim~\ref{cl:lvtriangle}(H); v1.19.0-pre.6: still; the C++ $\m$-block reader was specific to R19-1's quartic at $(258,37)$, a generic $\m_{\mathrm{Eis}}$-block of
rank $>2$ needs the eigensystem split over a general $\mathcal O$); \textbf{v1.20.0-pre.6}: \CMPD{} at $(282,37)$, $(226,103)$ --- $v=2$ on the $p_0$ sign, $0$ on the other --- by slot-20 WP-LV3's
generic reader (Check~\ref{chk:lvtriangle}(a), (b)). The instance (\texttt{HidaPairingRoute}; gates G1--G5; index statements (H1)--(H3)) was specified in R19-9~\S6 for
slot~20 and is \emph{implemented} in slot~19 (v1.19.0-pre.6, WP-BIL): $\langle\Phi,c\rangle=\langle\Phi(\{0,\infty\}),y_c\rangle$ over $\Z/\ell^{m}$ on three named lattices --- the mixed
perfect form $H^{1}_c(\Ga;V_n)\times H^{1}(\Ga;D_n)$ (every even $n$, $\ell\ge5$), Hida's self-dual form on $H^{1}_{!}$ ($n+1<\ell$, $G_{!}(P,Q)=\langle P,\theta\iota_{!}Q\rangle=-\tfrac16B(P,Q)$,
radical $\Z_\ell p_0$), and $H^{1}_c(V_n)\times H^{1}_{!}(D_n)$ at $(258,37)$. Its gates (\texttt{HeckeLValueValidation}, \CMPD): on the code's \texttt{PoincareDualityPairing} the cone-level
Gram matrix $G^{1}$ equals $-\langle\Psi(c,c_\partial),y(v_4v_6)\rangle$ at R19-9's $9$ cells ($n\ge\ell$ among them) and at $(258,37)$ (env arm), and the route reproduces it;
$A_D(T_p)=A_c(T_p)^{\mathsf T}$ for $p=2,3,5$ and $\iota$ anti-isometric at the $27$ cells of weights $12$--$28$ and at $(258,37)$ modulo $37^{10}$; R19-9's valuations at all $27$ cells and at
$(258,37)$ (mixed $3/2$, through $H^{1}_{!}(D)$ $3/3$, $\sum_fv_{37}=6$); every value equals the length of the congruence module of its lattice read from $A_c(T_2)$ alone (the
lemma's right-hand side), and the Eisenstein values equal $v(B_k/k)$ of the \textsf{congruence} kind's Bernoulli routine. The change-of-lattice constant for $n\ge\ell$ (between
$\theta(H^{1}(V))$ and $H^{1}(D)$ on the $f$-line) was not computed in slot~19; the normalisation record printed it as open (v1.20.0-pre.1: defined in Definition~\ref{def:lvlattice}; v1.20.0-pre.6:
its value is Proposition~\ref{prop:lvlattice}, \PROVED, with the values of Remark~\ref{rmk:lvlatticevalues}, and the record computes it, slot-20 WP-LV3).
(B) and (H) were \emph{declared} in the code (throwing \texttt{NotImplementedInThisBuild}) through v1.19.0-pre.5; since v1.19.0-pre.6 both are implemented (slot-19 WP-LIN, WP-BIL), on
[C]'s seven-vertex model at level one only (other models and levels refuse by name). The implementations reproduce the research prototypes and are gated; no statement of this
document outside this remark, Remark~\ref{rmk:lvlegc} and \S\ref{sec:lvnorm} rows~(B), (H) rests on them. The three kinds are distinct and each has its instance: (B) \textsf{linear}, (H) \textsf{bilinear}, (E) \textsf{congruence}; (H)'s index
statement ``discriminant $=$ congruence number'' equates a \textsf{bilinear} evaluation with the module $Q_{\lambda_f}$ of the \textsf{congruence} kind read on $\T(H')$
(Definition~\ref{def:lvroute}(2), $e=\lambda_f$), so a route may carry two kinds that its index statement compares (v1.19.0-pre.2). At the block of
\S\ref{sec:block258} the congruence modules of the block are computed already, by other means: Hida's congruence module of each eigenform $f_i$,
$\mathcal O_{K_i}/(P'(\rho_i))\cong\mathcal O_{K_i}/\pi^{3}$ (order $37^{3}$: three root distances $1/2$, so $v_\pi(P'(\rho_i))=3$; it includes the Galois conjugate), and the
inter-orbit congruence module $\tilde\T/\T_\m$ (order $37^{2}$) of Claim~\ref{cl:258alg}(4) (v1.19.0-pre.2, F1 m-E4). The first is what (H)'s prototype reads on the free sign
($v_\pi=3$, R19-9; v1.19.0-pre.5).
(M) (v1.20.0-pre.1; owner ruling O5 of 2026-10-04) \emph{The Mazur--Tate--Teitelbaum measure of an ordinary symbol} is not a fourth shape but a family of \textsf{linear} evaluations of~(B),
one per ball of $\Z_\ell$: \S\ref{sec:lvmeasure} (Definitions~\ref{def:lvmeasure}, \ref{def:lvmtt}, Proposition~\ref{prop:lvmtt}, Claim~\ref{cl:lvmeascong}).
\emph{Coefficients} (v1.21.0-pre.1; slot-21 P3, \texttt{docs/implementation/slot21/P3\_spec.md}): (M) is defined for an ordinary eigen-system with values in a local monogenic $\Z_\ell$-order
$\mathcal O=\Z_\ell[x]/(F)$ --- unramified, $W(\F_q)$, or totally ramified --- by $\mathcal O$-linearity of Definition~\ref{def:lvmtt} in the symbol ($P_\theta=\sum_jx^{j}P_j$, $P_j$ with
$\Z_\ell$-coordinates), and Proposition~\ref{prop:lvmtt} holds verbatim over $\mathcal O$ (\PROVED: its proof is an identity of $c$-tables, valid over any $\Z_\ell$-algebra); \emph{implemented}
(slot-21 P3, \texttt{HeckeLValueCoefficientRingValidation}, Remark~\ref{rmk:lvcoeffring}); refused by name: residual multiplicity $\ge2$ of the eigenvalue on a sign, and non-ordinary.
\end{remark}

\begin{remark}[Leg~(c) in modular-symbol form (v1.19.0-pre.6)]\label{rmk:lvlegc}
(Slot-19 WP-LIN on R19-8~\S\S3, 6; predictions \texttt{docs/implementation/slot19/wplin/predictions.md}, committed before any run.) Let $p_0:=u_0-u_n\in W$ (the image of
$H^{0}(\Ga_\infty;V_n)=\langle u_0\rangle$, the symbol of the boundary class; $T_pp_0=(1+p^{k-1})p_0$), and
$\mathrm{Bern}_k:=\sum_{i\ \mathrm{even},\,2\le i\le k-2}\frac{B_iB_{k-i}}{i!\,(k-i)!}X^{i-1}$, (the interior of the Eisenstein period polynomial; shape recalled, not read; R19-8~\S2.2), written in the coordinates $u_t$ and made primitive. At the $8$ Eisenstein cells with
$\dim S_k=1$, $(k,\ell)=(12,691)$, $(16,3617)$, $(18,43867)$, $(20,283)$, $(20,617)$, $(22,131)$, $(22,593)$, $(26,657931)$, on $H^{1}_c(\Ga;V_n\otimes\Z/\ell^{3})$ of the C++ module cone
(\CMPD{} at all $8$; the controls $(10,37)$, $(22,37)$ with $v(B_k/k)=0$ give depth~$0$):
\begin{enumerate}
\item the $\m_{\mathrm{Eis}}$-block of the geometric-even part has rank~$2$ ($p_0$ and the cusp line);
\item \textbf{(L1)} the cusp eigenline $r^{+}$ satisfies $r^{+}\equiv u\cdot p_0\pmod\ell$, $u$ a unit, to depth exactly $v(B_k/k)=1$ (the route's index statement);
\item $r^{-}\equiv s\cdot\mathrm{Bern}_k\pmod\ell$ to the same depth, with $\mathrm{Bern}_k\in W(\F_\ell)\smallsetminus W(\Z/\ell^{2})$ and $T_p\mathrm{Bern}_k\equiv(1+p^{k-1})\mathrm{Bern}_k\pmod\ell$, $p\le7$;
\item \textbf{(L2)} every twisted value ($\chi\neq\mathbf1$, the $8$ quadratic characters of R19-8) of $r^{+}$ is $\equiv0\pmod\ell$; on $r^{-}$ the number of right-sign twisted values
  that are $\ell$-units is R19-8's count.
\end{enumerate}
(L2) follows from (L1) and $\lambda_{\chi,j}(p_0)=0$ for $\chi\neq\mathbf1$ (\PROVED{} in R19-8~\S1.4). Item~2 is the $H^{1}_c$ form of $\Delta\equiv E_{12}$: the integral cusp vector is congruent
to the boundary symbol. At depth $\ge2$ ($v(B_k/k)\ge2$, e.g.\ $(282,37)$) (L1) and item~3 are \CONJ{} (not computed: none of the six weights with $\dim S_k=1$ has
$v(B_k/k)\ge2$ at any $\ell$, and the gated reading is on a block with one cusp line). \emph{Evidence} (v1.20.0-pre.1; slot-20 R20-1, \texttt{docs/implementation/slot20/R20-1\_measures.md}~\S4.2,
Python, one route, predictions committed first): at $(282,37)$ and $(226,103)$ ($d=2$; on the $\m_{\mathrm{Eis}}$-block separated by $T_3$) $r^{+}\equiv u\cdot p_0$ and
$r^{-}\equiv s\cdot P^{\mathrm{Eis}}$ (Claim~\ref{cl:lvmeascong}(1),(2)) hold to depth exactly~$2$, and at $(30,37)$ to depth exactly~$1$. Both items \textbf{stay \CONJ} at depth $\ge2$ (lead decision,
slot-20 WP-2) until the C++ gate of Check~\ref{chk:lvtriangle}(b) reads them; Claim~\ref{cl:lvtriangle}. \textbf{v1.20.0-pre.6}: the C++ gate read (L1) --- $r^{+}\equiv u\cdot p_0$ to depth
\emph{exactly}~$2$ at $(282,37)$ and $(226,103)$ (slot-20 WP-LV3, \texttt{BirchModularSymbolRoute::index\_statement()} on the group-cocycle $W$ over $\Z/\ell^{3}$ and, in the heavy arm, on the C++ cone
over $\Z/37^{3}$, $\Z/103^{3}$; the $\m_{\mathrm{Eis}}$-block cut jointly by $T_2,T_3,T_5$), so \textbf{(L1) is \CMPD{} at depth~$2$} at these two cells (R20-1's Python and WP-LV3's C++, two routes).
Item~3 was not gated by WP-LV3 and stays \CONJ{} at depth $\ge2$ as an index statement of this route (R20-1's reading, one route; the congruence itself is \CMPD{} as a run, Claim~\ref{cl:lvmeascong}(2)).
\textbf{v1.22.0-pre.2} (slot-22 P5$''$, the Fable referee \texttt{docs/implementation/slot22/P5pp\_referee.md}~\S\S1.3--1.4): (L1) at depth \emph{exactly}~$1$ at $(66,37)$ and $(102,37)$ ($d=1$, $n>\ell$) --- $r^{+}\equiv p_0\pmod{37}$,
$r^{+}\not\equiv p_0\pmod{37^{2}}$, $r^{+}$ rebuilt over $\Z/37^{3}$ on the group-cocycle $W$ from the period relations and Merel's $T_2$, $T_3$ with the $a_2$, $a_3$ of the verb \texttt{lfunction}
(\texttt{docs/implementation/slot22/p5pp/referee/ref\_direct.py}, outputs \texttt{referee/out/direct\_66\_37.txt}, \texttt{direct\_102\_37.txt}): computed, \textbf{the referee's Python, one route} --- not yet \CMPD{}
in this document's two-route sense; \texttt{BirchModularSymbolRoute::index\_statement()} has not read these cells (no arm of the gate runs them; a follow-up).
\end{remark}

\subsection{The Eisenstein / Bernoulli instance}\label{sec:lveis}

\begin{definition}[The instance (E)]\label{def:lveis}
Let $k\ge4$ be even, $n=k-2$, $a:=k\bmod(\ell-1)$ with $a\not\equiv0$. $D$: $H:=H^{1}(\Ga;\Mn\otimes\Z_\ell)$ modulo torsion, and $H_m:=H^{1}(\Ga;\Mn\otimes\Z/\ell^{m})$; $H'$:
the minus sign $H^{-}$ (cuspidal: the Eisenstein class lies in $H^{+}$) localised at the Eisenstein maximal ideal $\m_{\mathrm{Eis}}=(\ell,\,T_p-1-p^{k-1})$. $\kappa$:
\textsf{congruence} (v1.19.0-pre.2; at pre.1 mislabelled \textsf{linear}, referee F1 M-2) --- the Eisenstein character $e\colon T_p\mapsto e_p:=1+p^{k-1}$ (values in
$\mathcal O=\Z_\ell$; it does not occur in $H^{-}$) against the Hecke algebra $\T^{-}_k=\T(H')$ of $H^{-}_{\m_{\mathrm{Eis}}}$, $I_{\mathrm{Eis}}:=I_e=(T_p-1-p^{k-1})_p$, and
$Q_e=\T^{-}_k/I_{\mathrm{Eis}}\T^{-}_k$ ($\mathrm{Ann}(I_{\mathrm{Eis}})=0$: $\T^{-}_k$ is torsion-free and $I_{\mathrm{Eis}}\otimes\Q_\ell=\T^{-}_k\otimes\Q_\ell$). This is the ring form,
Claim~\ref{cl:lveis}(1); the module form, Claim~\ref{cl:lveis}(2), reads the same valuation through a \textsf{linear} functional on the \emph{other} sign, $r^{1}|_{H^{+}}$ with its
$\omega_\partial$-coordinate, and is \STD. $\mathcal N$: $L=\zeta$ at $s_0=1-k$, $\zeta(1-k)=-B_k/k$ (the constant term of $E_k$; not $L(E_k,s)=\zeta(s)\zeta(s-k+1)$ at $s=k$, which is a pole; F1 m-E1); the
Euler factor at $\ell$ removed, $(1-\ell^{k-1})$ (a unit, so it changes no valuation); no period; $\chi=\mathbf 1$; sign $-$; the ambiguity: $\Z_\ell^{\times}$.
$\mathcal I$: Claim~\ref{cl:lveis}.
\end{definition}

\begin{claim}[The Eisenstein index statements]\label{cl:lveis}
In the setting of Definition~\ref{def:lveis}:
\begin{enumerate}
\item \textbf{(Ring.)} \PROVED{} from \Ucite{pr:u-U1} (itself \PROVED{} from \STD{} (Emerton, Thm.~5, verified in [U])) $\T^{-}_k/I_{\mathrm{Eis}}\T^{-}_k\cong\Z_\ell/\ell^{v(B_k/k)}$, for
  every $k\ge4$ in a class $a\not\equiv0$.
\item \textbf{(Module.)} \STD{} (Harder's denominator of the Eisenstein class, \cite{harder3}, \VERIFYloc; Claim~\ref{cl:harder} is its mod-$\ell$
  reading) for $n<\ell-1$: the primitive integral Eisenstein class $E\in H^{+}$ restricts to $r^{1}(E)=\ell^{v(B_k/k)}\cdot u\cdot\omega_\partial$, $u\in\Z_\ell^{\times}$,
  $\omega_\partial$ a generator of the free part of $H^{1}(\Ga_\infty;\Mn\otimes\Z_\ell)$; equivalently $[H^{+}:H^{+}_{!}\oplus\Z_\ell E]=\ell^{v(B_k/k)}$. \CONJ{} for $n\ge\ell-1$.
\item \textbf{(Finite precision.)} \PROVED{} from (1) and \HA{cor:stablecert}. Let $\mathsf A_m:=\T^{-}_{k,\m_{\mathrm{Eis}}}/(\T^{-}_{k,\m_{\mathrm{Eis}}}\cap\ell^{m}\mathrm{End})$ be the Hecke
  algebra of the $\m_{\mathrm{Eis}}$-block of $H^{-}_m$ (\S\ref{sec:ring}). If $\T^{-}_{k,\m_{\mathrm{Eis}}}$ is \emph{saturated} in $\mathrm{End}(H^{-}_{\m_{\mathrm{Eis}}})$ --- the
  \texttt{stable} verdict of \HA{cor:stablecert}(1) \emph{at $K=1$} together with the generation certificate (2), equivalently \HA{prop:cohforms}(5)'s condition; automatic on a
  rank-$1$ block --- then $\mathsf A_m=\T^{-}_{k,\m_{\mathrm{Eis}}}/\ell^{m}$, and for a set $S$ of $T_p$ generating $\mathsf A_m$ (\S\ref{sec:certgen})
  \[
    \mathsf A_m\big/\bigl(T_p-1-p^{k-1}:p\in S\bigr)\;\cong\;\Z/\ell^{\min(m,\,v(B_k/k))}.
  \]
  Without saturation (a \texttt{stable} verdict first at some $K>1$; \HA{cor:stablecert} allows $K\le c^{\mathrm{sat}}+1$), with $\ell^{c^{\mathrm{sat}}}\T^{\mathrm{sat}}\subset\T^{-}_{k,\m_{\mathrm{Eis}}}$,
  only $\min(m-c^{\mathrm{sat}},v(B_k/k))\le\mathrm{length}\le\min(m,v(B_k/k))$ (v1.19.0-pre.2, referee F1 M-1: a \texttt{stable} verdict at $K>1$ does not give
  $\mathsf A_m=\T/\ell^{m}$ --- e.g.\ $\T=\Z_\ell[\ell e]\subset\Z_\ell\times\Z_\ell$, $e=(1,0)$, is \texttt{stable} at $K=2$ but $|\mathsf A_2|=\ell^{3}\neq|\T/\ell^{2}\T|=\ell^{4}$).
  Independently of any certificate, $\min_{p\neq\ell}v\bigl(c_{k,p}(1+p^{k-1})\bigr)\ge v(B_k/k)$ for the characteristic polynomials $c_{k,p}$ of $T_p$ on
  $H^{-}_{m,\m_{\mathrm{Eis}}}$ (read modulo $\ell^{m}$; \Ucite{thm:u-sandwich}, over $p\neq\ell$: at $p=\ell$ a non-ordinary form of the residual class contributes
  $v(a_\ell-1-\ell^{k-1})=0$ on the full $H^{-}$).
\end{enumerate}
\end{claim}

\begin{proof}
(1) The Eisenstein system is ordinary ($T_\ell\equiv1+\ell^{k-1}\equiv1$ modulo $\m_{\mathrm{Eis}}$), so $\T^{-}_{k,\m_{\mathrm{Eis}}}$ is the $\m_{\mathrm{Eis}}$-factor of the ordinary cuspidal algebra
$\mathbf h_k$ of [U] (\Ucite{lem:u-levelone}(1): the algebras generated by the $a_p(f)$ and by the unit roots $\alpha_f$ agree, $\alpha\in\Z_\ell[T_\ell]$ by Hensel; and
$H^{-}\otimes\bar\Q_\ell$ carries each cuspidal eigensystem once, \HA{prop:cohforms}). On the $\m_{\mathrm{Eis}}$-factor, $I_{\mathrm{Eis}}$ (with $T_\ell-1-\ell^{k-1}$) and [U]'s $I^{0}$ (with $U-1$) generate the same ideal: $T_\ell-1-\ell^{k-1}=(\alpha-1)(1-\ell^{k-1}/\alpha)$,
$\alpha$ the unit root, the second factor a unit (v1.19.0-pre.2, F1 m-E2). By (U1-$\Lambda$) $\mathbf h(a)/I^{0}(a)\cong\Lambda/(G_a)$; tensoring with $\Lambda/P_{k}$
(right exact; \Ucite{thm:u-control}) gives $\mathbf h_k/I^{0}\mathbf h_k\cong\Z_\ell/G_a(u^{k}-1)$, and $v(G_a(u^{k}-1))=v(B_k/k)$ (\Ucite{pr:u-U1}). The quotient is
supported at $\m_{\mathrm{Eis}}$ only. (3) Saturation gives $\T\cap\ell\,\mathrm{End}=\ell\T$, hence by induction $\T\cap\ell^{j}\mathrm{End}=\ell^{j}\T$ for all $j$, so $\mathsf A_m=\T^{-}_{k,\m_{\mathrm{Eis}}}/\ell^{m}$; the ideal generated by the $T_p-e_p$, $p\in S$, is the image of $I_{\mathrm{Eis}}$
(the $T_p$, $p\in S$, generate, so $I_{\mathrm{Eis}}$ is generated by the $T_p-e_p$, $p\in S$, as an ideal); reduce (1) modulo $\ell^{m}$. Without saturation, $\ell^{c^{\mathrm{sat}}}\T^{\mathrm{sat}}\subset\T$ gives $\ell^{m}\T\subset\T\cap\ell^{m}\mathrm{End}\subset\ell^{m-c^{\mathrm{sat}}}\T$, so
$\mathsf A_m/I$ is a quotient of $\T/(\ell^{m},I_{\mathrm{Eis}})$ and surjects onto $\T/(\ell^{m-c^{\mathrm{sat}}},I_{\mathrm{Eis}})$: the two-sided bound. The inequality is
\Ucite{thm:u-sandwich}'s upper bound (over $p\neq\ell$), $L_k\le v(c_{k,p}(1+p^{k-1}))$, with $L_k=v(B_k/k)$.
\end{proof}

\noindent\emph{A reading of Claim~\ref{cl:lveis}(2) on the ordinary block} (v1.21.0-pre.1; slot-21 F4, \texttt{docs/implementation/slot21/P8\_F4\_referee\_report.md}, M3 and \S4 item~3; no tag moves).
For $n<\ell-1$, (2) is equivalent to: the ordinary $\m_{\mathrm{Eis}}$-block of $H^{+}$ is free of rank one over $\T_{k,\m_{\mathrm{Eis}}}$ --- this is \Ucite{pr:u-msym}~(U17)(a), which [U] proves
from (2) on the ordinary block ($\ell^{j_k}=[H^{+}:H^{+}_{!}\oplus\Z_\ell E]$ there, $H^{1}(\partial;\Mn\otimes\Z_\ell)$ being torsion-free for $n<\ell$); for $n\ge\ell-1$ both stay \CONJ.

\begin{claim}[The twisted instance through $\Ga_1(\ell)$]\label{cl:lvtwist}
\PROVED{} from \Ucite{pr:u-U1}, \Ucite{lem:u-twistvalues} and \STD{} (Hida's control at the nebentypus points: for $k\ge2$ \emph{not} in the class $a$,
$\mathbf h(a)/P_{k}\cong\mathbf h^{\mathrm{ord}}_k(\Ga_0(\ell),\omega^{a-k})$ \cite{hida_ens} (pp.~249--250: Thm.~3.1 and Cor.~3.2 (3.3b), $h_0(N,a;\mathcal O_K)/P_kh_0(N,a;\mathcal O_K)\cong
h^{0}_k(\Ga_1(N_0p),a;\mathcal O_K)$ for \emph{every} $k\ge2$, $P_k=(1+X)-(1+p)^{k}$, with p.~249's $S_k(\Ga_1(N_0p),a;A)=S_k(\Phi,\omega^{a-k};A)$, $\Phi=\Ga_0(p)\cap\Ga_1(N_0)$,
$\omega$ a character of $\Phi$ via $\bigl(\begin{smallmatrix}a&b\\c&d\end{smallmatrix}\bigr)\mapsto\omega(d)$; here $N_0=1$, $p=\ell$; read on the page by referee F1, slot~19).
Convention check: \HA{def:nebmodel} twists by $\chi(a_\gamma)=\bar\chi(d_\gamma)$, but \HA{prop:nebcorr}(3) ($T(p,p)=\chi(p)p^{n}$) and \HA{thm:nebrelations}(2)
($T_{p^{2}}=T_p^{2}-\chi(p)p^{n+1}$) pin the model's eigensystems to those of $S_k(\Ga_0(\ell),\chi)$ in the standard ($d$-)convention, which is Hida's. Let $\chi=\omega^{j}$
with $j$ even (so $\chi(-1)=1=(-1)^{k}$), $j\not\equiv0\pmod{\ell-1}$, and $a:=k+j\not\equiv0$. On the cuspidal cohomology $H^{1}_{!}$ (both copies of each eigensystem; Eichler--Shimura) of
the character model $H^{1}(\Ga_0(\ell);V_{\Z_\ell}(\chi))$ (\HA{def:nebmodel}), localised at the maximal ideal of the Eisenstein system $p\mapsto1+\chi(p)p^{k-1}$ (that of $E_k(\mathbf1,\chi)$), with Hecke algebra $\T^{-}$:
\[
  \T^{-}/I_{\mathrm{Eis}}\T^{-}\;\cong\;\Z_\ell\big/\tfrac{B_{k,\chi}}{k}\Z_\ell.
\]
\end{claim}

\begin{proof}
As Claim~\ref{cl:lveis}(1), with the nebentypus specialisation in place of \Ucite{thm:u-control}: $\epsilon_p(x_k)=1+\omega(p)^{a-1}\langle p\rangle^{k-1}=1+\omega^{a-k}(p)p^{k-1}$
(\Ucite{def:u-eischar}), and $G_a(x_k)=-B_{k,\omega^{a-k}}/k$ (\Ucite{lem:u-twistvalues}), $\omega^{a-k}=\chi$. The identification of $\mathbf h^{\mathrm{ord}}_k(\Ga_0(\ell),\chi)$
with the Hecke algebra of the minus sign of the character model is Eichler--Shimura for $\Ga_1(\ell)$ (\STD, \cite{shimura}~Ch.~8).
\end{proof}

\begin{remark}[Checkable form and gate cells for WP-LV]\label{rmk:lvgates}
\textbf{Untwisted.} At $(\ell,k)\in\{(37,32),(691,12)\}$ ($v(B_k/k)=1$), at the regular controls $(37,k)$ with $v(B_k/k)=0$, and at $(59,44)$, $(67,58)$
($v=1$): Claim~\ref{cl:lveis}(3) at $m=1,2,3$ --- $\Z/\ell^{\min(m,1)}$ or $0$ --- and the boundary reading of~(2), all with $n<\ell-1$. The $\m_{\mathrm{Eis}}$-blocks
there have rank $1$ (one cusp form congruent to $E_k$), so the saturation hypothesis of Claim~\ref{cl:lveis}(3) is met at all four untwisted cells (a rank-$1$
algebra $\Z_\ell$ is saturated; v1.19.0-pre.2, F1 M-1). A deeper rung, $v(B_k/k)=2$, is $(282,37)$ ($k=284$) and $(226,103)$ ($k=228$),
both with $n\ge\ell-1$: there (1) is \PROVED, (3) is \PROVED{} under its saturation hypothesis, and (2) is \CONJ. At pre.2 the rank of the $\m_{\mathrm{Eis}}$-block at these two
cells was not established; since [U] 0.3.0-pre.3 it is: (U13), \Ucite{pr:u-lambda}, is \PROVED{} at $37$ and $103$ (from the \CMPD{} reading $d_{\mathrm{Eis}}=1$ at
one weight, slot-19 R19-2, and Hida's control; slot-19 referee F4, m8), so the cuspidal component through $\m_{\mathrm{Eis}}$ is $\Z_\ell[[X]]$, the block has rank~$1$ at every
weight of the two classes, and the saturation hypothesis of (3) is met there (v1.19.0-pre.3). Neither cell is a gate cell of WP-LV; both are cells of WP-LV2's env arm
(\texttt{HECKE\_LVALUE\_HEAVY=1}, read-only; v1.19.0-pre.5, Status below);
\texttt{HeckeLValueValidation} prints the rank and the saturation verdict of every $\m_{\mathrm{Eis}}$-block it gates, falling back to the two-sided bound of (3) where
saturation fails. \textbf{Twisted.} Claim~\ref{cl:lvtwist} at $\ell\le37$ only (cost), on \HA{sec:nebentypus}'s
model (as built, v1.19.0-pre.5: at $37$ in the default arm, at $59$, $67$ in the env arm). \textbf{Bernoulli numbers.} Exact rationals from WP-NF's \texttt{BigInt} (\S\ref{sec:lvnf}), and independently modulo $\ell^{m}$ (the Kummer\,/\,Voronoi route);
the two must agree. The $\lambda$-readings and the first depth-$2$ rungs that [U] derives from two Bernoulli numbers per class (\Ucite{prop:u-eistrunc}: $284$ at $37$,
$228$ at $103$, and at $691$ the rungs $k=342252$ (class $12$) and $k=342440$ (class $200$), where $v(B_k/k)=2$ is now \emph{computed}, with $v=1$ at $k\mp690$ ---
referee F1, power-sum route, \texttt{docs/implementation/slot19/wp3/f1\_lambda.py} and \texttt{.out}) are the cross-check of the Bernoulli code. As built
(v1.19.0-pre.3): the exact route is WP-LV's own tangent-number recurrence on a minimal unsigned big integer (\texttt{HeckeExt/LValue/Bernoulli}; WP-NF's
\texttt{BigInt} may replace it, value-free), the modulo-$\ell^{m}$ routes are Voronoi's congruence and slot-15's Kummer--Mahler route
(\texttt{HeckeExt/HeckeHidaComponents}); all three agree, and the depth-$2$ rungs $284$, $228$, $342252$, $342440$ read $v=2$. [U] 0.3.0-pre.3 adds the depth-$3$ rungs
$37580$, $914250$, $49928412$, $159113510$, \PROVED{} from five exact Bernoulli numbers per class (\Ucite{prop:u-eistrunc}); they are not gated here.
\textbf{Status} (v1.19.0-pre.3). \CMPD{} (slot-19 WP-LV, \texttt{HeckeLValueValidation}, $311$ checks; reproduced on the forms side by the slot-19 referee F4, m18): at
the four untwisted cells, $(n,\ell)=(30,37)$, $(10,691)$, $(42,59)$, $(56,67)$, the $\m_{\mathrm{Eis}}$-block is rank~$1$ and saturated, $\#Q_m=\ell^{\min(m,1)}$ for
$m=1,2,3$, the module form's index is $\ell^{1}=\ell^{v(B_k/k)}$ ($n<\ell-1$ at all four), and $v_\ell(c_{k,p}(1+p^{k-1}))=1$ at $p=2,3,5,7$; at the regular controls
$(10,37)$, $(22,37)$ the module $Q_m$ is $0$. The saturated rank-$1$ path is the only one exercised: the two-sided-bound fallback of Claim~\ref{cl:lveis}(3) has no
test cell on cohomology.
\textbf{Status} (v1.19.0-pre.5; slot-19 WP-LV2, \texttt{HeckeLValueValidation} extended, $537$ checks by default, $657$ with \texttt{HECKE\_LVALUE\_HEAVY=1};
\texttt{docs/implementation/slot19/WP-LV2\_predictions.md}, predictions committed before the runs). \emph{Twisted} (\CMPD): Claim~\ref{cl:lvtwist} on the character model
$H^{1}(\Ga_0(\ell);V_{\Z/\ell^{K}}(\omega^{j}))$ (slot-18 WP-6's model route; $\iota^{\ast}$ on the induced model new) at $(\ell;k,j)=(37;4,28)$, $(37;6,26)$ (default arm), $(59;4,40)$,
$(67;4,54)$ (env arm): $v(B_{k,\chi}/k)=1$ by two routes (the definition with exact $B_i$ and Teichm\"uller $\omega$, and Kummer's congruence), the $\m_{\mathrm{Eis}}$-block of
$H^{-}$ of rank~$1$ and saturated, $\#Q_m=\ell^{\min(m,1)}$ at $m=1,2,3$, and $\rank H^{+}=\rank H^{-}+2$; the controls $(37;4,8)$ and $(37;4,34)$ (the latter with one torsion summand,
read through the free part) have $v=0$ and $Q=0$; the conjugate twist and the untwisted character give $Q=0$. The twisted depth-$2$ rung $(37;62,6)$ has $v(B_{62,\omega^{6}}/62)=2$
by both Bernoulli routes; its module was \emph{not} run (above $7.3$\,GiB, stopped by a $6$\,GiB watchdog; v1.20.0-pre.1: slot-20 WP-LV3 runs it under a cap taken from a measured growth curve,
Check~\ref{chk:lvtriangle}(e); v1.20.0-pre.6: \textbf{run} at $n=60$, $16.43$\,GB, $7.3$\,h --- the $\m_{\mathrm{Eis}}$-block of $H^{-}$ of rank~$1$, saturated, $\#Q_m=37^{\min(m,2)}$, $m=1,2,3$, \CMPD). \emph{Deep rungs} (\CMPD): at $(282,37)$ and $(226,103)$,
$H^{1}(\Ga;\Mn\otimes\Z/\ell^{4})$ is free (no $\ell$-torsion; the prediction of torsion was wrong), of rank $47$ resp.\ $39$ ($24+23$, $20+19$), the $\m_{\mathrm{Eis}}$-block of $H^{-}$
has rank~$1$ and is saturated, $\#Q_m=\ell^{\min(m,2)}$ at $m=1,2,3$, and $v_\ell(c_{k,p}(1+p^{k-1}))=2$ at $p=2,3,5,7$; the module form's index reads $c=2=v(B_k/k)$, printed and not
gated --- (2) stays \CONJ{} for $n\ge\ell-1$. Slot-19 R19-7 read the same two cells independently (console over $\Z/\ell^{3}$ and an outside reading of the matrices, the joint
Eisenstein block, $f=1$): $\#Q_2=\ell^{2}$, $v(a_p-1-p^{k-1})=2$ at $p=2,3,5,7,\ell$ (\CMPD). \emph{Fallback} of~(3): the two-sided bound is tested on planted algebras
($\Z_\ell[\ell e]$, $\Z_\ell[\ell+\ell^{2}e]$, with $c^{\mathrm{sat}}$), still on no cohomology cell. A read-only cell hook (\texttt{HECKE\_LVALUE\_CELLS}, level one or $\chi=\omega^{j}$)
runs further cells and gates only the index statement.
\end{remark}

\subsection{The normalisation table}\label{sec:lvnorm}
\begin{center}\footnotesize
\begin{longtable}{@{}>{\raggedright\arraybackslash}p{2.2cm}>{\raggedright\arraybackslash}p{3.2cm}>{\raggedright\arraybackslash}p{3.3cm}>{\raggedright\arraybackslash}p{3.0cm}>{\raggedright\arraybackslash}p{3.6cm}@{}}
\toprule
route & $L$ at $s_0$, twist & Euler factor, period, sign & index statement & status \\
\midrule
\endhead
(E) level $1$ & $\zeta(1-k)=-B_k/k$; $\chi=\mathbf1$ & $(1-\ell^{k-1})$ removed (unit); none; $-$ & $\T^{-}/I_{\mathrm{Eis}}\cong\Z_\ell/(B_k/k)$ & \PROVED{} from \STD{} (Claim~\ref{cl:lveis}(1)); module form \STD{} for $n<\ell-1$ \\
(E$_\chi$) $\Ga_1(\ell)$ & $L(1-k,\chi)=-B_{k,\chi}/k$; $\chi=\omega^{j}$ & none needed ($\chi(\ell)=0$); none; $-$ & $\T^{-}/I_{\mathrm{Eis}}\cong\Z_\ell/(B_{k,\chi}/k)$ & \PROVED{} given Hida's control at nebentypus points (Claim~\ref{cl:lvtwist}) \\
(KL) $\ell$-adic & $L_\ell(1-k,\omega^{a})=G_a(u^{k}-1)$ & $(1-\omega^{a-k}(\ell)\ell^{k-1})$; --- ; --- & --- (an interpolation identity) & \PROVED{} from \STD{} (\Ucite{lem:u-twistvalues}) \\
(a) Kummer\,/\,Coates--Wiles & $\varphi_k(\eta_i):=D^{k}\log F_{\eta_i}(T)|_{T=0}$, $D=(1+T)\tfrac{d}{dT}$, $F$ the Coleman series of the norm-coherent system
  $\eta_i^{(n)}=\prod_c\bigl((\zeta_{\ell^n}^{c}-1)/(\zeta_{\ell^n}-1)\bigr)^{f_c}$: Kummer's logarithmic derivative in Washington's normalisation
  (\cite[\S13.7, pp.~299--300]{washington1st}, defined there modulo~$\ell$), $\Z_\ell$-valued as the Coates--Wiles homomorphism $\delta_k$ on norm-coherent units (there,
  Lemma~13.52; Coates--Wiles' own $\varphi_k$, \cite[p.~230, (8)]{coates_wiles}, is modulo~$p$); $\varphi_k(\eta_i)=S_{i,k}B_k/k$, $S_{i,k}=\sum_cf_c(c^{k}-1)\equiv-1\pmod\ell$
  for $k\equiv i\not\equiv0\pmod{\ell-1}$; twist $\omega^{i}$ (the radical character of $\eta_i$)
  & $(1-\ell^{k-1})$: the restriction of the measure to $\Z_\ell^{\times}$, $(1-\ell^{k-1})\varphi_k(\eta_i)=-S_{i,k}L_\ell(1-k,\omega^{i})$ ($\varphi_k$ itself has none);
  auxiliary factor $S_{i,k}$ (the $(c^{k}-1)$-smoothing); no period; even
  & $v(\varphi_k(\eta_i))=v(B_k/k)$; modulo $\ell^{m}$ a function of the level-$m$ unit $\eta_i^{(m)}$ (the layer-$1$ unit gives it modulo $\ell$ only, $1\le k\le\ell-2$:
  Washington p.~300's argument); in the pole class $k\equiv0\pmod{\ell-1}$ the correspondence $k\leftrightarrow i$ fails
  & identity \PROVED{} (elementary; Kummer's computation, \cite[pp.~64--65]{washington1st}, and \S13.7, pp.~299--300, both read on the page); values \CMPD{} at $37$
  (all even $i$, levels $\le4$) and $691$ ($i=12,200$), four routes, slot-19 R19-4 (v1.19.0-pre.3).
  \emph{Level $m$, Bernoulli-free} (v1.19.0-pre.5; slot-19 R19-7): on the $\omega^{w}$-part, $\nu(a+\ell^{m}\Z_\ell)=\ell^{-m}\sum_{j=1}^{m}\mathrm{Tr}_j(\zeta_{\ell^{j}}^{-a}\lambda_j)$,
  $\lambda_m=\log_\ell(1-\zeta_{\ell^{m}})$, $\lambda_j=\mathrm{Tr}_{m/j}\lambda_m$, and $\Psi_k:=\sum_aa^{k}\nu_a\equiv(1-\ell^{k-1})B_k/k\pmod{\ell^{m}}$ (\PROVED{} there); \CMPD{} at
  $m=2$ ($37$, $59$, $67$, $103$, $691$; $5769/5769$ rows), $m=3$ ($37$, $103$), $m=4$ ($37$): the depth-$m$ rung is the unique zero of $\Psi$ modulo $\ell^{m}$ --- $284$, $914$,
  $3292$, $228$, $342252$, $342440$ ($m=2$), $37580$, $914250$ ($m=3$), $1072544$ ($m=4$, predicted from [U]'s $\alpha$ before the run); the level-$(m-1)$ datum is right
  modulo $\ell^{m-1}$ only. That the level-$2$ zero is the unramified $\Z/\ell^{2}$ Kummer lift: \PROVED{} there from \STD{} (\cite[Thm.~13.54, p.~311; Prop.~13.51, p.~309;
  Lemma~13.52, p.~310]{washington1st}, read on the page); at level $m\ge3$ \SPEC \\
(b) Leopoldt & $L_\ell(1,\omega^{j})$ $=-\tfrac{\tau(\omega^{j})}{\ell}\sum_{a}\omega^{-j}(a)\,\cdot$ $\log_\ell(1-\zeta^{a})$ $=-\sum_c\omega^{j}(c)x_c$,
  $\log_\ell(1-\zeta)=\sum_cx_c\zeta^{c}$ (Iwasawa's $\log$); $j$ even, $j\not\equiv0$
  & none ($\omega^{j}(\ell)=0$); no period; even
  & $L_\ell(1,\omega^{j})=G_j(0)$; $\equiv G_j(x_k)\pmod{\ell^{1+v(k)}}$, $k\equiv j$; $v\ge1\iff\ell\mid B_j/j$ (equality with $v(B_j/j)$ is not general; it holds,
  $v=1$, at $(37,32)$, $(691,12)$, $(691,200)$)
  & formula \STD{} (\cite[Thm.~5.18, p.~63; Prop.~8.12, p.~156]{washington1st}, read on the page); reduction \PROVED; values \CMPD{} modulo $\ell^{5}$ at $37$ (all even $j$;
  $L_{37}(1,\omega^{32})\equiv47910634$) and $691$ ($j=2,4,12,100,200$), three routes, slot-19 R19-4 (v1.19.0-pre.3) \\
(B) Birch & $L(f,\chi,j+1)$, $0\le j\le k-2$ & none; $\Omega_f^{\pm}$; $\chi(-1)(-1)^{j}$ (geometric; $-\chi(-1)(-1)^{j}$ in [C]) & $\lambda_{\chi,j}$ integral, $v$ of its value & \SPEC{} (Remark~\ref{rmk:lvshapes}); research prototype \CMPD{} (slot-19 R19-8), instance spec R19-8~\S6 (v1.19.0-pre.5);
  \emph{implemented} (v1.19.0-pre.6, slot-19 WP-LIN, \texttt{BirchModularSymbolRoute}): gates G1--G5 \CMPD, (L1) \CMPD{} at depth~$1$ ($8$ cells), \CONJ{} at depth $\ge2$ (v1.20.0-pre.6: \CMPD{} at depth~$2$ at $(282,37)$, $(226,103)$, slot-20 WP-LV3), (L2) \PROVED{} from (L1)
  and \CMPD{} (Remark~\ref{rmk:lvlegc}) \\
(H) Hida & $L(1,\mathrm{Ad}f)=L(k,\Sym^{2}f)$ & none removed (\cite[Thm.~5.1, p.~247]{hida_cong}, $N=1$, read; v1.19.0-pre.5); $\Omega_f^{+}\Omega_f^{-}$, $\delta^{\pm}$ primitive in a named lattice; $+\times-$; constant $-(-1)^{k/2}i(2\pi)^{k+1}/(k-1)!$ & $v\langle\delta_f^{+},\delta_f^{-}\rangle=$ length of the congruence module of that lattice & \SPEC{} (Remark~\ref{rmk:lvshapes});
  the constant \STD{} (Hida, read; Haberland, recalled) and \CMPD{} (10 forms); research prototype \CMPD{} (slot-19 R19-9), instance spec R19-9~\S6 (v1.19.0-pre.5); \emph{implemented} (v1.19.0-pre.6, slot-19 WP-BIL,
  \texttt{HidaPairingRoute}): index statements (H1) $H^{1}_{!}$, (H2) mixed, (H3) $(258,37)$ \CMPD{} against the congruence module of the lattice; (H2) at depth~$2$ ($(282,37)$) \CONJ{} (v1.20.0-pre.6: \CMPD{} at $(282,37)$, $(226,103)$, slot-20 WP-LV3);
  the change-of-lattice constant for $n\ge\ell$ open (v1.20.0-pre.1: Definition~\ref{def:lvlattice}, value open until slot-20 R20-3; v1.20.0-pre.6: Proposition~\ref{prop:lvlattice} \PROVED, values
  Remark~\ref{rmk:lvlatticevalues} \CMPD, computed by the record); (H2) at depth~$2$: Claim~\ref{cl:lvtriangle} \\
(M) MTT measure (v1.20.0-pre.1) & $L_\ell(f,\alpha,\cdot)$: $\int_{\Z_\ell^{\times}}\bar\chi(x)x^{j}\,d\mu_f$, $0\le j\le k-2$, $\chi=\mathbf1$ or of $\ell$-power conductor; branch series $\mathcal L_{f,a}(T)$;
  values in $\mathcal O$, the Hecke order of the eigen-system (v1.21.0-pre.1, slot-21 P3: $\Z_\ell$, $W(\F_q)$ or a totally ramified $\Z_\ell[x]/(F)$; Remark~\ref{rmk:lvcoeffring})
  & $(1-\ell^{j}/\alpha)(1-\beta/\ell^{j+1})$ ($\chi=\mathbf1$), $\alpha^{-m}$ (conductor $\ell^{m}$); $\Omega_f^{\pm}$ of $H^{1}_c$ (on the boundary sign $\ell^{d}$ times the cuspidal one); geometric $\chi(-1)(-1)^{j}$
  & congruence of measures with $-s\,G_{a+1}G_{k-1-a}$ modulo $\ell^{d}$ off the two pole branches, with $-s\,(G_{a+1}G_{k-1-a}+\mathrm{polar}_a)$ on them (v1.21.0-pre.3); $\mu^{+}=d$
  & measure, moments and the algebraic interpolation \PROVED{} (Proposition~\ref{prop:lvmtt}, slot-20 R20-1); the analytic constant (B)'s status ($\tau(\chi)$ for $\ell$-power conductor derived, not
  tested; \CMPD{} since v1.21.0-pre.3, slot-21 P4, Remark~\ref{rmk:lvwild}); the congruence \CMPD{} at $(30,37)$, $(282,37)$, $(226,103)$, the pole branches an exception (Claim~\ref{cl:lvmeascong}(4)),
  explained by an explicit polar term since v1.21.0-pre.3 (slot-21 P2, Proposition~\ref{prop:lveismeas}); literature (Mazur, Greenberg--Vatsal) held since 2026-10-05, not yet read for this claim (v1.21.0-pre.1);
  planned code slot-20 WP-MEAS (\texttt{measure()}, \texttt{HeckeLValueMeasureValidation}, Check~\ref{chk:lvmeas}); over $\mathcal O$: slot-21 P3 (Check~\ref{chk:lvcoeffring}); the two-variable
  family by interpolation: slot-21 P9, Remark~\ref{rmk:lvtwovar} (\CONJ{} as a $\Lambda$-adic object, \Ucite{pr:u-msym} (U16)) \\
\bottomrule
\end{longtable}
\end{center}

\subsection{The number-field layer}\label{sec:lvnf}
\SPEC{} at v1.19.0-pre.1; \textbf{built} at v1.19.0-pre.3 (slot-19 WP-NF, the seventh static target \texttt{NumberField}; owner's no-CAS rule of 2026-10-03). The objects the $L$-value legs (a), (b) and the Galois side of \S\ref{sec:block258} need, in in-house
exact arithmetic: $\Z[\zeta_{\ell^{k}}]$ and, modulo $\ell^{M}$, its decomposition by the idempotents $e_{\omega^{i}}$ of the group ring of $(\Z/\ell^{k})^{\times}$
(\emph{equivariant before global}: the $\omega^{12}$-part at $691^{2}$ is one module of rank $691$ over $\Z/691^{2}$, not a field of degree $476\,790$); the cyclotomic
units $u_c=(1-\zeta^{c})/(1-\zeta)$ and $\eta_i:=\prod_{c=2}^{\ell-1}u_c^{f_c}$, $f_c\equiv c^{-i}\pmod\ell$ lifted to $(-\ell/2,\ell/2)$ (slot-18 R18-2~\S1; $\eta_i$ spans the
$\omega^{i}$-part of $E/E^{\ell}$); $v_\pi$; the local test \PROVED{} there: for a unit $\eta$ with $v:=v_\pi(\eta^{\ell-1}-1)$, $K(\eta^{1/\ell})/K$ is unramified at
$\lambda=(1-\zeta)$ iff $v\ge\ell$ and $\lambda$ splits iff $v\ge\ell+1$; its layer-$2$ analogue over $\Q(\zeta_{\ell^{2}})$; $\Q(\alpha)$ arithmetic and $\ell$-adic
factorisation (Hensel, Newton polygon). \emph{Gate} (\texttt{NumberFieldValidation}, \texttt{fast}, $117$ checks; \textbf{held}, v1.19.0-pre.3): R18-2's exact $\eta_{32}$, $\eta_6$ at $37$ (FNV-1a-64 of the
$36$ decimal coefficients \texttt{0xcba3d1938a721f7b}, \texttt{0xef292f5a4e321b01}; $v_\pi(\eta^{36}-1)=38$, $6$), on the equivariant layer-$2$ route of slot-19 R19-3
$\kappa=284$ the only unramified lift at $37$ ($u=y_0=13$), $\kappa=914$ at $59$ ($y_0=31$), $\kappa=3292$ at $67$ ($y_0=59$), and the layer-$1$ test at $691$
($v=692$ at $i=12,200$). The $\Z/691^{2}$ layer itself ($\kappa_0=342252$, $342440$; [L37], \texttt{sec:fieldsribet}) is slot-19 R19-3's Python computation and is
not a gate of the library. The library realises no statement of [C]; of this document it realises the objects above and exact Bernoulli numbers
(\texttt{BigRational}); the Bernoulli numbers gated in Remark~\ref{rmk:lvgates} come from WP-LV's own exact route.
\emph{Built at v1.21.0-pre.1} (slot-21 T1, T2; STL only, the include guard unchanged; no golden line, digest or record key). \textbf{T1}: an in-house LLL (Cohen, GTM~138,
Alg.~2.6.3 rational and 2.6.7 integral; Definition~2.6.1's reducedness; the substitute for Lenstra--Lenstra--Lov\'asz, not held) and POLRED (Alg.~4.4.11\,/\,4.4.12), \texttt{NumberField/\{Lll,Polred\}},
gate \texttt{NumberFieldLllValidation} (\texttt{fast}, $674$ checks: known lattices ($\Z^{n}$, $E_8$, $A_2$, integer relations), random lattices against an own Gram--Schmidt test and a brute-force
$\lambda_1$, POLRED of $\Q(\zeta_7)^{+}$, $\Q(\zeta_{37})^{+}$, $\Q(\zeta_\ell)$ for $\ell=7,11,37$, negative controls). \textbf{T2}: multiprecision reals \texttt{NumberField/MpReal} and the
degree-$\ell$ Ribet field \texttt{NumberField/RibetField}: $P=\mathrm{charpoly}(\theta)$, $\theta=\mathrm{Tr}_{H/F}(\eta^{1/\ell})$, by two numerical precisions and exactly by CRT at split primes (all
equal), the Frobenius root counts against the Kummer prediction, at $(37,32)$ the maximal order ($\mathrm{disc}\,\mathcal O_F=+37^{35}$ exactly) and T1's POLRED; at $(691,12)$, $(691,200)$ the
unreduced polynomial (degree $691$, $24\,995$ resp.\ $14\,669$ digits, both routes equal) in the read-only env arm \texttt{NUMBERFIELD\_RIBET\_CELLS}; gate \texttt{NumberFieldRibetFieldValidation}
(\texttt{standard}, $29$ checks; \texttt{docs/implementation/slot21/T2\_report.md}). Status of T2's claims after the slot-21 referee F4 (\texttt{docs/implementation/slot21/rc/F4\_referee\_report.md}~\S4):
$\mathrm{Gal}(H/\Q)\cong\mathrm{AGL}(1,\ell)$ and $\mathrm{disc}\,F=(-1)^{(\ell-1)/2}\ell^{\ell-2}$ with signature $(1,(\ell-1)/2)$ \PROVED{} (the derivation is in the non-binding
\texttt{docs/latex/ribet\_cells.tex}, \texttt{rc:sec:methfield}); the complement formula for $\theta_j$, the Frobenius counts, the CRT route, the maximal order at $37$ and the POLRED output \CMPD. Nothing
of [C] or of this document's earlier statements rests on T1 or T2.

\noindent\emph{Code map of \S\ref{sec:lvalue} (\textbf{built}, slot-19 WP-LV and WP-NF; planned at v1.19.0-pre.1, contract first; v1.19.0-pre.3).} WP-LV: \texttt{HeckeExt/LValue/} --- \texttt{LValueRoute} (Definition~\ref{def:lvroute}: domain, kind $\in\{$\texttt{linear}, \texttt{bilinear}, \texttt{congruence}$\}$,
\texttt{NormalisationRecord}, \texttt{evaluate}, \texttt{index\_statement}), \texttt{EisensteinBernoulliRoute} (a \texttt{congruence} route; Definitions~\ref{def:lveis}, Claims~\ref{cl:lveis},
\ref{cl:lvtwist}), \texttt{BirchModularSymbolRoute} (\texttt{linear}) and \texttt{HidaPairingRoute} (\texttt{bilinear}) (declared, throwing \texttt{NotImplementedInThisBuild}, through v1.19.0-pre.5);
\texttt{HeckeLValueValidation} (\texttt{standard}, $311$ checks, $\approx5$\,s; planned at $\le4$\,GB) on Remark~\ref{rmk:lvgates}'s untwisted cells. At v1.19.0-pre.3 not built: the twisted
route of Claim~\ref{cl:lvtwist}. \textbf{Built at v1.19.0-pre.5} (slot-19 WP-LV2, same files): the twisted route on the character model over $\Z/\ell^{K}$ ($\iota^{\ast}$ on the induced
model, $T_p$\,/\,$U_\ell$ by \HA{prop:nebcorr}), twisted Bernoulli numbers by two routes, the free part $(H/\mathrm{tors})/\ell^{K}$ at finite precision (torsion-aware), $c^{\mathrm{sat}}$ and the
two-sided bound of Claim~\ref{cl:lveis}(3); the character model is accepted at level~$\ell$, the Moerdijk--Pronk and Shapiro models refused; \texttt{HeckeLValueValidation} now
$537$ checks, $\approx14$\,s, $0.17$\,GB; env arm \texttt{HECKE\_LVALUE\_HEAVY=1} ($657$ checks, $\approx419$\,s, $0.98$\,GB) and the read-only cell hook
\texttt{HECKE\_LVALUE\_CELLS}. \textbf{Built at v1.19.0-pre.6} (Remark~\ref{rmk:lvshapes}, Remark~\ref{rmk:lvlegc}): slot-19 WP-BIL, the \texttt{bilinear} instance
(\texttt{HeckeExt/LValue/HidaPairing.\{h,cpp\}}, \texttt{HidaPairingRoute.cpp}; arms of \texttt{HeckeLValueValidation} (G1--G5, negative controls), env arm
\texttt{HECKE\_LVALUE\_BILINEAR\_HEAVY=1} (the G1 gate at $(258,37)$)), and slot-19 WP-LIN, the \texttt{linear} instance (\texttt{HeckeExt/LValue/BirchModularSymbol.\{h,cpp\}}, frozen
R19-8 data \texttt{BirchFrozenR198.inc}; arms (BS-$\ast$) of \texttt{HeckeLValueValidation}; \texttt{HECKE\_LVALUE\_HEAVY=1} adds $(258,37)$ over $\Z/37^{2}$ and G3 at $(72,37)$); both on
[C]'s seven-vertex model at level one, other models refused by name; check counts and timings in \texttt{docs/validation-catalogue.md}. WP-NF: the seventh target \texttt{NumberField/} (STL only, include guard
\texttt{check\_numberfield\_includes.sh}), \texttt{NumberFieldValidation} (\texttt{fast}, $117$ checks). Both value-free with respect to golden, the digests and the record schema.
\textbf{Planned at v1.20.0-pre.1} (slot~20, contract first, rule~21). WP-MEAS (\S\ref{sec:lvmeasure}): \texttt{HeckeExt/LValue/ModularSymbolMeasure.\{h,cpp\}} (the measure, the branch elements,
$\mu$\,/\,$\lambda$, the Eisenstein symbol, the two Kubota--Leopoldt group-ring routes), \texttt{LValueRoute::measure} with its default refusal and \texttt{BirchModularSymbolRoute::measure},
the frozen table \texttt{MeasureFrozenR201.inc}, \texttt{HeckeLValueMeasureValidation} (\texttt{standard}, Check~\ref{chk:lvmeas}); no addition to
\texttt{NumberField/CwLayer} was needed (\texttt{cw\_lambda} + \texttt{cw\_measure} return $\nu$; v1.20.0-pre.2, WP-MEAS's report). WP-LV3 (\S\ref{sec:lvdepth2}): the generic $\m_{\mathrm{Eis}}$-block reader replacing WP-BIL's quartic-specific one, the cross-kind gate,
$c^{\varepsilon}_f$ in \texttt{NormalisationRecord}, $\iota$ over $\Z/\ell^{k}$ in \texttt{Hecke/HeckeInvolution} and the console, the twisted rung behind \texttt{HECKE\_LVALUE\_HEAVY}
(Check~\ref{chk:lvtriangle}); no new validation name. \textbf{Built at v1.20.0-pre.6} (slot-20 WP-LV3, merged \texttt{c1fda051}; WP-MEAS merged before it, Check~\ref{chk:lvmeas}):
\texttt{HeckeExt/LValue/MBlockReader.\{h,cpp\}} (the generic $\m$-block reader: the Newton polygon of the block's characteristic polynomial through \texttt{NumberField}, the roots refined digit by digit in
tame, totally ramified orders $\Z_\ell[\pi]/(\pi^{e}-\ell u)$, $e\mid\ell-1$, $f=1$; an unramified residue extension, $e\nmid\ell-1$, a non-integral refinement slope and a multiple root refused by name),
the route change in \texttt{HidaPairingRoute.cpp} (WP-BIL's quartic-specific reader kept as the cross-check arm), \texttt{HeckeLValueLV3Arms.cpp}, the joint $T_2,T_3,T_5$ cut in
\texttt{BirchModularSymbol.cpp}, $c^{\varepsilon}_f$ in \texttt{NormalisationRecord}, \texttt{Hecke/HeckeInvolution}'s \texttt{BuildRingInvolution} (schema \texttt{hecke\_involution\_ring/v1}) and the
console's \texttt{iota} under $\Z/\ell^{k}$, script \texttt{049\_lv3\_iota\_ring.oc}; \texttt{HeckeLValueValidation} $1471$ checks by default ($\approx31$\,s, $0.58$\,GB), $1614$ with
\texttt{HECKE\_LVALUE\_HEAVY=1} ($729$\,s, $1.55$\,GB), the twisted rung behind \texttt{HECKE\_LVALUE\_HEAVY=twist62}; no new validation name, no golden line or pinned digest moved.
\textbf{Built at v1.20.0-pre.5} (slot-20 stretch S1, S2): \texttt{HeckeExt/LValue/LFunctionRecord.\{h,cpp\}} (the record
\texttt{lfunction/v1}, \S\ref{sec:lvlfunction}), the console verb \texttt{lfunction} (\texttt{read\_only}), arm (M7) and the env arm (M691) of \texttt{HeckeLValueMeasureValidation} with
\texttt{MeasureFrozenS2.inc}; no new validation name, no golden line, digest or existing record key. Both on [C]'s seven-vertex model at level one, the model taken as a parameter; neither moves a golden line, a pinned digest or a record key.

\subsection{Measures: the measure-valued \textsf{linear} evaluation}\label{sec:lvmeasure}
(New at v1.20.0-pre.1, hecke-extensions slot~20, WP-2; owner ruling O5 of 2026-10-04: \emph{no fourth kind}. Source: slot-20 R20-1,
\texttt{docs/implementation/slot20/R20-1\_measures.md}, Python, in-house, predictions committed before the gated runs (\texttt{ffb34f70}); the code is slot-20 WP-MEAS,
\texttt{docs/implementation/slot20/WP-MEAS\_spec.md}, which merges only after this text is at pre.1 (rule~21). Placed at the end of \S\ref{sec:lvalue}, so that no
statement number moves.) Throughout: the domain of Remark~\ref{rmk:lvshapes}~(B) --- [C]'s seven-vertex model at level one, $H^{1}_c(\Ga;\Mn\otimes\mathsf R)\cong W(\mathsf R)$
by $\Psi$ ($W=\ker(1+\rho_n(S))\cap\ker(1+\rho_n(U)+\rho_n(U)^{2})$, \CMPD, slot-19 WP-LIN), $\mathsf R=\Z_\ell$ or $\Z/\ell^{R}$, $\ell\ge5$, $n=k-2$ even, $u_t=x^{n-t}y^{t}$,
$\rho_n(A)P(x,y)=P(ax+cy,bx+dy)$ on all of $\mathrm{Mat}_2(\Z)$, $\Phi(D)$ the modular symbol of a class with $\Phi(\{0,\infty\})=\Psi(x)$, and $\Phi(\{a/M,\infty\})$ computed along
Manin's continued-fraction path. The geometric sign $\varepsilon$ is the eigenvalue of $\rho_n(\mathrm{diag}(-1,1))$ ($u_t\mapsto(-1)^{t}u_t$); geometric-even is [C]-minus
(Remark~\ref{rmk:lvshapes}~(B)). $u:=1+\ell$, $\omega$ the Teichm\"uller character, $x=\omega(x)\langle x\rangle$, $\langle x\rangle=u^{\ell_u(x)}$ for $x\in\Z_\ell^{\times}$.

\begin{definition}[The measure-valued evaluation of a \textsf{linear} route]\label{def:lvmeasure}
\SPEC. A \emph{measure-valued evaluation} of a \textsf{linear} route $\mathcal R=(D,\textsf{linear},\mathcal N,\mathrm{ev},\mathcal I)$ (Definition~\ref{def:lvroute}) is a family of values
$\mu_x(a+\ell^{M}\Z_\ell)\in\mathsf R$, for $x\in H'$, $M\ge0$, $a\in\Z/\ell^{M}$, such that (i) for each ball $x\mapsto\mu_x(a+\ell^{M}\Z_\ell)$ is an $\mathsf R$-linear functional on
$H'$ --- a \textsf{linear} evaluation in the sense of Definition~\ref{def:lvroute}(2), so Definition~\ref{def:lvroute} is unchanged; (ii) it is a \emph{distribution},
$\mu_x(a+\ell^{M}\Z_\ell)=\sum_{b=0}^{\ell-1}\mu_x(a+b\ell^{M}+\ell^{M+1}\Z_\ell)$; (iii) it is bounded (values in $\mathsf R$). Its normalisation record states the
interpolation formula of its moments; its index statements are \emph{congruences of measures}: equalities of the branch elements in $\mathsf R[\Ga_u/\Ga_u^{\ell^{M-1}}]$
($\Ga_u=1+\ell\Z_\ell$) for every level $M$ (equivalently in $\mathsf R[\![\Z_\ell^{\times}]\!]$) modulo a named ideal. The code carries it as the accessor \texttt{measure()} of
\textsf{linear} routes; a \textsf{congruence} or \textsf{bilinear} route refuses it (\texttt{std::logic\_error}), a \textsf{linear} route that does not implement it throws
\texttt{NotImplementedInThisBuild}. Slot-19 R19-7's level-$m$ Coates--Wiles measure (\S\ref{sec:lvnorm} row~(a)) has the same shape on the layer, not on cohomology.
\end{definition}

\begin{definition}[The Mazur--Tate--Teitelbaum measure of an ordinary eigensymbol (M)]\label{def:lvmtt}
Let $x\in H^{1}_c(\Ga;\Mn\otimes\mathsf R)$ be a $T_\ell$-eigenclass, $T_\ell x=a_\ell x$, with $a_\ell\in\mathsf R^{\times}$ (\emph{ordinary}), $P:=\Psi(x)$, $\alpha$ the unit root
of $X^{2}-a_\ell X+\ell^{k-1}$ (unique, $\alpha\equiv a_\ell\pmod{\ell^{k-1}}$), $\beta:=\ell^{k-1}/\alpha$. Put $\gamma_{a,m}:=\begin{pmatrix}\ell^{m}&-a\\0&1\end{pmatrix}$ and
\[
  c(a,m):=\bigl\langle\rho_n(\gamma_{a,m})\,\Phi(\{a/\ell^{m},\infty\}),\,u_n^{\vee}\bigr\rangle\quad(m\ge0),\qquad c(a,-1):=P_n,
\]
($c(a,m)=-\sum_iP(q_i,q'_i)$ over the bottom rows $(q_i,q'_i)$ of Manin's matrices of $a/\ell^{m}$, with $P(c,d):=\sum_tP_tc^{n-t}d^{t}$). The measure (M) is
\[
  \mu_x(a+\ell^{m}\Z_\ell)\;:=\;\alpha^{-m}\,c(a,m)\;-\;\alpha^{-m-1}\,\ell^{k-2}\,c(a,m-1)\qquad(m\ge0).
\]
The factor $\ell^{k-2}$ on the stabilisation term is part of the definition: without it the family is not a distribution (a negative control,
Check~\ref{chk:lvmeas}(M6); the formula of the slot-20 plan, \texttt{00-slot20-plan.md} v1.0 \S2 WP-MEAS, omits it --- corrected by R20-1).
\end{definition}

\begin{proposition}[(M) is a measure, its moments, its interpolation formula]\label{prop:lvmtt}
\PROVED{} (elementary; slot-20 R20-1~\S\S2.1--2.7, re-derived for this text). In the setting of Definition~\ref{def:lvmtt}:
\begin{enumerate}
\item \textbf{(The $\ell$-Hecke relation at the $\ell$-adic cusps.)} $\sum_{b=0}^{\ell-1}c(a+b\ell^{m},m+1)=a_\ell\,c(a,m)-\ell^{k-2}c(a,m-1)$ for all $m\ge0$ and $a$.
\item \textbf{(Measure.)} $\mu_x$ depends only on $a\bmod\ell^{m}$, is a distribution and is bounded: it is the $u_n^{\vee}$-coordinate of the $\Mn(\mathsf R)$-valued measure
  $D(a,m):=\alpha^{-m}\rho_n(\gamma_{a,m})\Phi_\alpha(\{a/\ell^{m},\infty\})$, where $\Phi_\alpha:=\Phi-\alpha^{-1}V\Phi$, $V\Phi(D):=\rho_n(\mathrm{diag}(1,\ell))\Phi(\ell D)$, is the
  $\ell$-stabilisation: $U_\ell\Phi_\alpha=\alpha\Phi_\alpha$ with $U_\ell\Phi(D):=\sum_b\rho_n\bigl(\begin{smallmatrix}\ell&-b\\0&1\end{smallmatrix}\bigr)\Phi((D+b)/\ell)$, from $U+V=T_\ell$ and
  $UV=\ell^{k-1}$ on level-one symbols.
\item \textbf{(Moments.)} For compact open $U\subset\Z_\ell$ and $0\le j\le n$: $\langle D(U),u_{n-j}^{\vee}\rangle=\binom nj\int_U(-x)^{j}\,d\mu_x$; the level-$M$ Riemann sum is
  correct modulo $\ell^{M}$.
\item \textbf{(Interpolation, algebraic form.)} For $\chi$ primitive of conductor $\ell^{m}$, $m\ge1$, and for $\chi=\mathbf 1$:
  \[
    (-1)^{j}\binom nj\int_{\Z_\ell^{\times}}\bar\chi(x)\,x^{j}\,d\mu_x\;=\;E_{\chi,j}\cdot\lambda_{\chi,j}(\Phi),\qquad
    E_{\mathbf 1,j}=\Bigl(1-\frac{\ell^{j}}{\alpha}\Bigr)\Bigl(1-\frac{\ell^{n-j}}{\alpha}\Bigr)=\Bigl(1-\frac{\ell^{j}}{\alpha}\Bigr)\Bigl(1-\frac{\beta}{\ell^{j+1}}\Bigr),\quad E_{\chi,j}=\alpha^{-m},
  \]
  $\lambda_{\chi,j}$ the functional of Remark~\ref{rmk:lvshapes}~(B) at modulus $\ell^{m}$ (at modulus $1$ for $\chi=\mathbf 1$: $\lambda_{\mathbf 1,j}(\Phi)=P_{n-j}$).
\item \textbf{(Parity, branches.)} $\mu_x(-U)=\varepsilon\,\mu_x(U)$. The \emph{branch series} $\mathcal L_{x,a}(T):=\int\omega^{a}(x)(1+T)^{\ell_u(x)}\,d\mu_x$ satisfy
  $\mathcal L_{x,a}(u^{j}-1)=\int\omega^{a-j}(x)x^{j}\,d\mu_x$, vanish for $(-1)^{a}\ne\varepsilon$, and level $M$ determines $\mathcal L_{x,a}$ modulo
  $\omega_{M-1}:=(1+T)^{\ell^{M-1}}-1$, by the branch element $g_i=\sum_{\delta\in\mu_{\ell-1}}\omega^{a}(\delta)\mu_x(\delta u^{i}+\ell^{M}\Z_\ell)$; so $\mu=0$ and $\lambda<\ell^{M-1}$
  are read at level~$M$, the $T^{t}$-coefficients for $t<\ell$ modulo $\ell^{M-1}$.
\item \textbf{(The boundary symbol.)} Write $\mu^{(\alpha)}_{Q}$ and $D^{(\alpha)}_{Q}(a,m):=\alpha^{-m}\mathrm{vec}_Q(a,m)-\alpha^{-m-1}\ell^{k-2}\mathrm{vec}_Q(a,m-1)$ for the formula of
  Definition~\ref{def:lvmtt} applied to an arbitrary symbol $Q\in W$ with a fixed unit $\alpha$ (linear in $Q$; for $Q=\Psi(x)$ and $\alpha$ the unit root of $x$ it is $D(a,m)$). For
  $p_0=u_0-u_n$ (Remark~\ref{rmk:lvlegc}) and $m\ge1$, $\mathrm{vec}_{p_0}(a,m)=\ell^{mn}p_0$, so $D^{(\alpha)}_{p_0}(a,m)=\alpha^{-m-1}(\alpha-1)\ell^{mn}p_0$ (and $0$ for $\alpha=1$, the
  trivial zero of the Eisenstein measure). Hence if $\Psi(x)\equiv u\cdot p_0\pmod{\ell^{e}}$ then, by linearity, $\mu_x(U)\equiv0\pmod{\ell^{\min(e,n)}}$ for every compact open
  $U\subset\Z_\ell^{\times}$ (v1.20.0-pre.2, F1 M-1: the object of the identity pinned; with (2)'s $D(a,m)$ of the eigenclass $p_0$ itself, unit root $1$, it would be $0$).
\end{enumerate}
\emph{Analytic form} (status of (B)'s constant). With Remark~\ref{rmk:lvshapes}~(B)'s normalisation, (4) reads
\[
  \int_{\Z_\ell^{\times}}\bar\chi(x)\,x^{j}\,d\mu_f\;=\;(-1)^{j}\,E_{\chi,j}\,\frac{j!\,M_\chi^{\,j+1}}{\chi(-1)(-2\pi i)^{j+1}\tau(\chi)}\cdot\frac{L(f,\chi,j+1)}{\Omega_f^{\pm}},
\]
$M_\chi$ the conductor, $\pm=\chi(-1)(-1)^{j}$ geometric; the $\binom nj$ cancels, so no unit hypothesis on it is needed. The constant has (B)'s status: \CMPD{} for trivial and quadratic $\chi$
of conductor prime to~$\ell$; for $\chi$ of $\ell$-power conductor the factor $\tau(\chi)$ is \CMPD{} since v1.21.0-pre.3 (slot-21 P4, Remark~\ref{rmk:lvwild}(1): every primitive $\chi$ of conductor
$25$, $49$, $121$ at $k=12$ and $25$, $49$ at $k=16$, and $125$, $169$ at $k=12$ in an env arm, against complex $L(f,\chi,j+1)$; the quotient is the conductor-$1$ $\Omega_f^{\pm}$ to $\le10^{-59}$; through
v1.21.0-pre.2 it read ``\emph{derived, not tested} (no complex $L(f,\omega^{i},s)$ was computed)''). The shape of
$E_{\mathbf 1,j}$ is MTT's $e_p$ (\cite{mtt} \S11, recalled, \VERIFYloc; R19-8 read (8.3), (8.6) on the page for the non-$\ell$ part).
\end{proposition}

\begin{proof}[Proof sketch (R20-1~\S2)]
(1) $T_\ell\Phi(D)=\sum_b\rho_n(\begin{smallmatrix}\ell&-b\\0&1\end{smallmatrix})\Phi((D+b)/\ell)+\rho_n(\mathrm{diag}(1,\ell))\Phi(\ell D)$ at $D=\{a/\ell^{m},\infty\}$, with
$\gamma_{a,m}\begin{pmatrix}\ell&-b\\0&1\end{pmatrix}=\gamma_{a+b\ell^{m},m+1}$, $\gamma_{a,m}\mathrm{diag}(1,\ell)=\ell\gamma_{a,m-1}$, $\rho_n(\ell I)=\ell^{n}$; at $m=0$ the last term has
$u_n$-coordinate $\ell^{n}P_n$. (2) $D(a,m)=\alpha^{-m}\mathrm{vec}(a,m)-\alpha^{-m-1}\ell^{k-2}\mathrm{vec}(a,m-1)$ by (1)'s matrix identity; $\gamma_{a+\ell^{m},m}=\gamma_{a,m}T^{-1}$ and
$T$-equivariance give the dependence on $a\bmod\ell^{m}$; additivity is $U_\ell\Phi_\alpha=\alpha\Phi_\alpha$; boundedness because $\alpha$ is a unit. (3)
$\gamma_{a,m}=T^{-a}\mathrm{diag}(\ell^{m},1)$ gives $D_j(a,m)\equiv\binom nj(-a)^{j}D_0(a,m)\pmod{\ell^{m}}$; sum over a partition and let $M\to\infty$. (4) For $\chi$ of conductor
$\ell^{m}$ the stabilisation term depends on $a\bmod\ell^{m-1}$ and is killed by $\sum\bar\chi$; for $\chi=\mathbf 1$, $\sum_{a\in(\Z/\ell)^{\times}}D(a,1)=D(0,0)-D(0,1)$ with
$\gamma_{0,1}=\mathrm{diag}(\ell,1)$. (5) $\Phi(\{-r,\infty\})=\varepsilon\rho_n(\iota)\Phi(\{r,\infty\})$ and $\iota\gamma_{a,m}=\gamma_{-a,m}\iota$; the rest is the binomial expansion of
$\omega_r$ ($v(\binom{\ell^{r}}{t})=r-v(t)$). (6) The first column of $\gamma_{a,m}g$ is $(0,\ell^{m})^{\top}$ for every Manin matrix $g$, and $p_0$ is the symbol of
$\Phi(\{r,s\})=F(s)-F(r)$, $F(\gamma\infty)=\rho_n(\gamma)u_0$.
\end{proof}

\begin{claim}[The Eisenstein congruence of measures]\label{cl:lvmeascong}
\CMPD{} at $(n,\ell)=(30,37)$, $(282,37)$, $(226,103)$, $d:=v(B_k/k)=1,2,2$ (slot-20 R20-1~\S\S3--4, \texttt{r201/run\_cells.py}, \texttt{out/*\_full*.out}; each value compared with an
independent object: exact twisted Bernoulli numbers at every critical moment, and the Kubota--Leopoldt group ring built from the regularised Bernoulli measure $E_{1,c}$). The
reduction to the Eisenstein symbol --- $\mu_{r^{-}}\equiv s\,\mu^{(1)}_{P^{\mathrm{Eis}}}\pmod{\ell^{d}}$ on every ball --- is \PROVED{} given its three inputs (R20-1 Prop.~3.1: $r^{-}\equiv sP^{\mathrm{Eis}}$,
$\alpha\equiv1$ and $\ell^{k-2}\equiv0$ modulo $\ell^{d}$, then linearity); its Kubota--Leopoldt form (3) is \CMPD{} (the twisted values of $P^{\mathrm{Eis}}$, \STD{} recalled, then the runs). The expected
literature is Mazur and Greenberg--Vatsal (v1.21.0-pre.1: both held since the library inventory of 2026-10-05, \texttt{docs/implementation/slot21/library\_inventory.md}; neither yet read for this claim; through v1.20.0 this sentence read ``literature pending: neither is in the library''). Let $f$ be the ordinary eigenform with
$a_p\equiv1+p^{k-1}\pmod\ell$ found by a separating operator ($T_2+3T_3+7T_5$ at $(30,37)$, $T_3$ at the two depth-$2$ cells), $r^{\pm}$ its integral eigensymbols of the
two geometric signs.
\begin{enumerate}
\item \textbf{(The Eisenstein symbol.)} $P^{\mathrm{Eis}}_{n-j}:=\binom nj\frac{B_{j+1}}{j+1}\frac{B_{n-j+1}}{n-j+1}$ for odd $j$, $0$ otherwise ($=n!\,\mathrm{Bern}_k$ of
  Remark~\ref{rmk:lvlegc}, up to the primitive scaling), is integral and primitive, and $P^{\mathrm{Eis}}\in W(\Z/\ell^{d})\smallsetminus W(\Z/\ell^{d+1})$ (\CMPD; that it lies in
  $W(\Z/\ell^{d})$ follows from Zagier's Laurent polynomial $p_k^{-}\in\hat W_k^{-}$ --- read on the page since v1.21.0-pre.4 (F4m): Zagier, Invent.\ Math.\ 104 (1991), Prop.~(i), eq.~(11), p.~453,
  read on the page; through pre.3 ``Zagier's Eisenstein period function, \STD{} recalled, not read'' (the name is the repository's, not Zagier's; Proposition~\ref{prop:lveismeas}(1)): its terms outside $\Mn$ are $\frac{B_k}{k(k-1)}(X^{-1}+X^{n+1})$, of valuation $d$ when
  $\ell\nmid k-1$).
\item \textbf{(Minus side, values.)} On the geometric-odd sign, $r^{-}\equiv s\cdot P^{\mathrm{Eis}}\pmod{\ell^{d}}$, $s$ a unit, to depth exactly~$d$.
\item \textbf{(Minus side, measures.)} On every odd branch $a\not\equiv-1,\,k-1\pmod{\ell-1}$,
  \[
    \mathcal L^{-}_{f,a}(T)\;\equiv\;-s\;G_{a+1}\bigl(u(1+T)-1\bigr)\;G_{k-1-a}\bigl(u^{k-1}(1+T)^{-1}-1\bigr)\pmod{\ell^{d}},
  \]
  $\mathcal L^{-}_{f,a}$ the branch series of $\mu_{r^{-}}$ (Proposition~\ref{prop:lvmtt}(5)), $G_b$ of \Ucite{lem:u-twistvalues}: as whole elements of $\Z/\ell^{d}[\Ga_u/\Ga_u^{\ell^{M-1}}]$ at
  level $M=3$ ($37$) and $M=2$ ($103$) ($16/16$, $16/16$, $49/49$ branches), and at every critical moment $T=u^{j}-1$, $0\le j\le n$, against exact twisted Bernoulli numbers
  ($0$ failures in $496$, $4528$, $11123$). The group-ring form is stronger than the moments: at $n=30$ the $31$ critical values fix a series only modulo $(\ell,X^{31})$
  (\Ucite{thm:u-interp}). Negative control: the wrong twist $G_{a+3}$ matches $0/78$ branches. \emph{A second derivation} (v1.21.0-pre.3, slot-21 P2; no tag moves): by the reduction above,
  $\mathcal L^{-}_{f,a}\equiv s\,\mathcal L^{\mathrm{poly}}_{a}\pmod{\ell^{d}}$, $\mathcal L^{\mathrm{poly}}$ the branch series of $\mu^{(1)}_{P^{\mathrm{Eis}}}$; P2 computes
  $\mathcal L^{\mathrm{poly}}_a\equiv\mathcal L^{\mathrm{Eis}}_a\pmod{\ell^{d}}$ off the pole branches (to depth exactly $d$; Proposition~\ref{prop:lveismeas}(5), \CMPD) and
  $\mathcal L^{\mathrm{Eis}}_a=-G_{a+1}(u(1+T)-1)\,G_{k-1-a}(u^{k-1}(1+T)^{-1}-1)$ exactly in $\Lambda\otimes\Q$ (Proposition~\ref{prop:lveismeas}(3), \PROVED{} given \STD), which gives (3) at every cell
  where (2) and that comparison hold.
\item \textbf{(The pole branches: an explicit exception.)} On the two branches $a\equiv-1$ and $a\equiv k-1\pmod{\ell-1}$ one factor of (3) is $G_0$, the pole of $\zeta_\ell$, and (3)
  is \textbf{not} claimed there. \CMPD: (3)'s right-hand side, read at the critical moments, fails modulo $\ell^{d}$ at every point of both branches ($62/62$, $566/566$, $454/454$);
  modulo $\ell^{d-1}$ (at $d=2$) it fails at exactly $14$ of $566$ points at $(282,37)$ and at exactly $4$ of $454$ at $(226,103)$, namely $(a,j)=(23,21)$, $(23,124)$, $(101,102)$,
  $(101,205)$ --- the points where the argument of the $G_0$-factor ($x_{k-1-j}$ on the branch $k-1$, $x_{j+1}$ on the branch $-1$) has valuation $\ge2$. At $(282,37)$ the $14$
  points are $(a,j)=(31;\,24,61,98,135,172,209,246)$ and $(35;\,36,73,110,147,184,221,258)$ --- the rule's prediction ($j\equiv k-1\pmod{37}$ on the branch $31$, $j\equiv-1\pmod{37}$
  on the branch $35$), committed by slot-20 WP-MEAS before its first run (\texttt{374c2394}) and located by its C++ gate (\CMPD, Check~\ref{chk:lvmeas}(M4)(d)); so $18/18$ exceptions are
  located at both cells by one rule. On both pole branches $\mu=\lambda=0$ at all three cells (\CMPD). \emph{Explained} (v1.21.0-pre.3, slot-21 P2, Proposition~\ref{prop:lveismeas}(4)--(6)):
  the exact Eisenstein measure has, on the two pole branches, a polar term $\mathrm{Res}/(u(1+T)-1)$, resp.\ $\mathrm{Res}/(u^{k-1}(1+T)^{-1}-1)$, with
  $\mathrm{Res}=(1-\ell)(\log_\ell u/\ell)(1-\ell^{k-1})B_k/k$, $v(\mathrm{Res})=d$ (the constant \CMPD), and on every odd branch, pole branches included,
  $\mathcal L^{-}_{f,a}\equiv s\,(\mathcal L^{\mathrm{Eis}}_a-\mathrm{polar}_a)\pmod{\ell^{d}}$; at the critical points
  $\mathcal L^{-}_{f,a}(u^{j}-1)\equiv-s\,\bigl(F_1F_2+\mathrm{Res}/(u^{r}-1)\bigr)\pmod{\ell^{d}}$, $F_1F_2$ the product of (3)'s two $G$-factors at $T=u^{j}-1$ (R20-1's naive product of twisted Bernoulli numbers), $r=j+1$ on the branch $-1$ and $r=k-1-j$ on the branch $k-1$
  (\CMPD{} with $r^{-}$ at the three cells: $0/62$, $0/566$, $0/454$ critical points, Check~\ref{chk:lveismeas}(E8)). As $v(u^{r}-1)=1+v(r)$, the polar value has valuation exactly $d-1-v(r)$
  (\PROVED{} from $v(\mathrm{Res})=d$), so the naive product agrees modulo $\ell^{d-1-v(r)}$ and no further: the exceptions modulo $\ell^{d-1}$ are exactly the points with $\ell\mid r$ --- the rule above,
  and the $14+4$ located points, now \emph{derived} (gate: $\mathrm{Res}\cdot(1+\ell)$ in place of $\mathrm{Res}$ fails exactly there). Greenberg--Vatsal and Mazur (held, library inventory 2026-10-05,
  not yet read for it) are literature context, not a dependency; R20-L is not needed. Through v1.21.0-pre.2 this sentence read: ``\emph{Not explained here}: the regularised Eisenstein measure on
  the pole branches (R20-1~\S4.4 reads a polar term of valuation $d-1-v(r)$ [\dots]) is \textbf{deferred} [\dots]; no statement of this document rests on it.'' The congruence~(3) is still not claimed on
  the pole branches; the polar-term form is Proposition~\ref{prop:lveismeas}(6). At $(10,691)$ the polar-term form is checked on $\Delta$'s own eigensymbol at level $2$, at all $3795$ critical points and as group-ring elements on every odd branch, pole branches included (\CMPD, Check~\ref{chk:lveismeas}(E8L2); Proposition~\ref{prop:lveismeas}(6)).
\item \textbf{(Minus side, invariants.)} $\mu=0$ and $\lambda=0$ on every odd branch at the three cells (\CMPD; off the pole branches this follows from (3), both $G$-factors being
  units --- $32$ is the only irregular index at $37$, $24$ at $103$). The general form is the level-one Greenberg--Vatsal formula $\lambda^{-}_a=\lambda(G_{a+1})+\lambda(G_{k-1-a})$ off the pole branches,
  the cell's own irregular index always on a pole branch (Remark~\ref{rmk:lviwasawa}(7), v1.21.0-pre.7, slot-21 P5; \CMPD{} at nine cells, among them $(10,691)$ with $\lambda^{-}=1$ at $199$, $501$).
\end{enumerate}
The plus side is Remark~\ref{rmk:lvmeasplus}.
\end{claim}

\begin{remark}[The plus side: the boundary, $\mu^{+}=d$, and an observation]\label{rmk:lvmeasplus}
(Slot-20 R20-1~\S\S2.7, 4.2, 4.5.) (i) On the geometric-even sign the $H^{1}_c$-primitive eigensymbol satisfies $r^{+}\equiv u\cdot p_0\pmod{\ell^{d}}$, $u$ a unit, to depth exactly
$d$, at the three cells of Claim~\ref{cl:lvmeascong} --- (L1) of Remark~\ref{rmk:lvlegc}, computed at depth~$2$ at $(282,37)$, $(226,103)$ (R20-1, one route, predictions committed
first); its tag stays \CONJ{} at depth $\ge2$ (Remark~\ref{rmk:lvlegc}). (ii) Given (i), $\mu_{r^{+}}(U)\equiv0\pmod{\ell^{d}}$ on every $U\subset\Z_\ell^{\times}$ (Proposition~\ref{prop:lvmtt}(6),
linearity in the symbol, $n\ge d$; \PROVED{} from (i)), and the minimum valuation on the level-$M$ unit balls is exactly $d$ (\CMPD{} at the three cells): $\mu^{+}=d$, i.e.\ the $H^{1}_c$-period on the
boundary sign is $\ell^{d}$ times the cuspidal one --- the measure form of slot-19 WP-BIL's finding that the mixed lattice adds $v(B_k/k)$ on that sign
(Remark~\ref{rmk:lvshapes}~(H)). (iii) \emph{Observation, not a claim.} The reduced plus measure $\ell^{-d}\mu_{r^{+}}$ (from $w=(r^{+}-p_0)/\ell^{d}$) has $\mu=0$ and $\lambda=1$ on the
branch pair $\{4,26\}$ at $(30,37)$ and $(282,37)$ and on $\{46,62,78\}$ at $(226,103)$, $\lambda=0$ elsewhere; the pairs are $a\leftrightarrow n-a\pmod{\ell-1}$ (the functional
equation; $62$ self-paired). At the regular controls $(10,37)$ ($\Delta$) and $(16,37)$, $\mu=0$ throughout and $\lambda=1$ on the branch pairs $\{21,25\}$, $\{21,31\}$ (minus) and
$\{8,22,26,30\}$ at $(16,37)$ (plus). No interpretation is proposed; the plus-side Eisenstein ``measure'' involves the trivial zero $L(\chi,j-n)=0$, so a reading needs the
derivative or boundary term (Greenberg--Vatsal's $\lambda$-formula for the other sign is the expected place; held, not yet read; v1.21.0-pre.1). \emph{Partly placed}
(v1.21.0-pre.7, slot-21 P5, Remark~\ref{rmk:lviwasawa}): $62$ at $(226,103)$ and $8$, $26$ at $(16,37)$ are even self-dual branches, where $\lambda$ is odd by the functional equation (\PROVED,
Remark~\ref{rmk:lviwasawa}(3)); the pairs $\{4,26\}$ and $\{46,78\}$ are $\{-k,n+k\}\pmod{\ell-1}$, a pattern seen at all nine Eisenstein cells (through pre.8 ``recorded, not derived''), since v1.21.0-pre.9 identified with the
McCallum--Sharifi forced zero (\CMPD{} at $32$ cells; $\lambda^{+}\ge1$ there through v1.22.0-pre.1 a theorem \emph{conditional on Sharifi's conjecture}; since v1.22.0-pre.2 \emph{unconditional at the computed cells} (slot-22 P5$''$; a per-pair computation, \PROVED{} modulo \STD{} inputs, given \CMPD{} L3 at the pair and \CMPD{} (H1) at the cell; the uniform statement open), Remark~\ref{rmk:lviwasawa}(9)); the controls' other pairs stay unexplained (Remark~\ref{rmk:lviwasawa}(8)).
\end{remark}

\begin{checkitem}[The gate of WP-MEAS: \texttt{HeckeLValueMeasureValidation}]\label{chk:lvmeas}
\SPEC{} (the checks); the expected values are R20-1's (\CMPD{} or \PROVED{} as tagged above), frozen in an in-source table \texttt{MeasureFrozenR201.inc}, never this code's output.
\emph{Registration}: \texttt{Registry/Slot20/WPMEAS.cpp}, \texttt{cmake/slot20/WPMEAS.cmake}; tier \texttt{standard}; quiet budget (not a gate: no timing is a verdict) $\le15$\,s and $\le1$\,GB (estimate $5$--$8$\,s,
$\le0.2$\,GB); first run $3.5$\,s, $37$\,MB, $178$ checks, $0$ failing (WP-MEAS). \emph{Cells}: $(10,37)$ ($\Delta$, $d=0$, level $2$), $(16,37)$ ($d=0$, level $3$), $(30,37)$ ($d=1$, level $3$), $(282,37)$ ($d=2$, level $3$), $(226,103)$ ($d=2$,
level $2$), all modulo $\ell^{4}$; $(10,5)$ the non-ordinary refusal ($\tau(5)\equiv0\pmod5$). Arms:
\begin{enumerate}
\item[(M0)] \textbf{The route} on the C++ module cone at $(30,37)$ over $\Z/37^{4}$: \texttt{measure()} of the cone classes mapping to $r^{\pm}$ equals the symbol-side measure (FNV equal);
  linearity on two eigenclasses of different sign; the \textsf{bilinear} and \textsf{congruence} routes refuse \texttt{measure()}; $(10,5)$ refuses ``not ordinary''; a non-eigenclass
  refuses ``not a $T_\ell$ eigenclass''. \textbf{This arm stays in the \texttt{standard} tier} (lead decision, slot-20 WP-2: $\approx3$\,s of the $15$\,s budget).
\item[(M1)] Proposition~\ref{prop:lvmtt}(1)--(5) at every cell and both signs: three-term and distribution failures $0$ at $m=0,\dots,M-1$; the wrong-parity branches vanish; the moment
  identity (4) for $\chi=\omega^{-i}$, $i\in\{0,1,2,3,(\ell-1)/2,\ell-2\}$, every $j$, modulo $\ell^{\min(M,R)}$, against slot-19 WP-LIN's $\lambda_{\chi,j}$ (\texttt{EvAllJ}); the Riemann-sum
  moment equals the group-ring branch value at three $j$ per cell.
\item[(M2)] R20-1's frozen values modulo $\ell^{4}$ at level $2$ ($r^{-}$ scaled to first unit coordinate $1$, $r^{+}$ to $r^{+}_0=1$; FNV-1a-64 of $\mu(a+\ell^{2}\Z_\ell)$, units $a$
  ascending, $8$ LE bytes each):
  \begin{center}\scriptsize
  \begin{tabular}{@{}lrrrrrlrl@{}}
  \toprule
  cell & $a_2$ & $a_3$ & $a_\ell=\alpha$ & $s$ & $\mu^{-}(2+\ell^{2})$ & FNV $\mu^{-}$ & $\mu^{+}(2+\ell^{2})$ & FNV $\mu^{+}$\\
  \midrule
  $(10,37)$ & 1874137 & 252 & 1454464 & --- & 272431 & \texttt{0x721edc172f8e24bd} & 847858 & \texttt{0xcc0c00522e7e0429}\\
  $(16,37)$ & 1873633 & 1869877 & 1454991 & --- & 57498 & \texttt{0xed8c569e28ac74be} & 698836 & \texttt{0x57a02c809eb73309}\\
  $(30,37)$ & 1622769 & 150214 & 270360 & 922897 & 1022643 & \texttt{0xf6fef89b62e18907} & 78662 & \texttt{0x6ddd3cd8d0e6a6ad}\\
  $(282,37)$ & 1079424 & 1744174 & 1770118 & 575834 & 1326561 & \texttt{0xc1d253ded2d790cd} & 1815294 & \texttt{0xa812126c962c7079}\\
  $(226,103)$ & 15697592 & 81971822 & 50764066 & 83864971 & 44186485 & \texttt{0x5fb753f5ed92f150} & 95024813 & \texttt{0xde51edced2aacde1}\\
  \bottomrule
  \end{tabular}
  \end{center}
  and $\mu^{-}(\Z_\ell^{\times})=0$.
\item[(M3)] Claim~\ref{cl:lvmeascong}(1),(2) and Remark~\ref{rmk:lvmeasplus}(i),(ii) at the three Eisenstein cells: $v(a_p-1-p^{k-1})=d$ for $p\in\{2,3,5,7,\ell\}$; (L1) to depth exactly $d$;
  $\mu^{+}=d$; $P^{\mathrm{Eis}}$ integral, primitive, in $W(\Z/\ell^{d})\smallsetminus W(\Z/\ell^{d+1})$; $r^{-}\equiv sP^{\mathrm{Eis}}$ to depth exactly $d$.
\item[(M4)] Kubota--Leopoldt: (a) the two group-ring routes for $G_b$ --- from $E_{1,c}$ and from \texttt{NumberField}'s Coates--Wiles layer (\S\ref{sec:lvnf}; $G_b=-(\ell-1)\nu_b$ in
  \texttt{CwLayer}'s index convention) --- agree for every even $b\not\equiv0$ at $37$ (level $3$) and $103$ (level $2$), and both reproduce $-(1-\ell^{k-1})B_k/k$ at $\ge4$ weights (a
  new prediction of WP-MEAS, \PROVED{} in principle, R19-7~\S\S1.1--1.2; held at its first run: $17/17$ and $50/50$ classes, $68/68$ and $200/200$ values); (b) Claim~\ref{cl:lvmeascong}(3) at the critical moments; (c) (3) as group-ring
  elements, with $G_b$ from both routes; (d) Claim~\ref{cl:lvmeascong}(4): the located failure sets modulo $\ell^{d-1}$ gated as computed (the $14$ at $(282,37)$ predicted and located).
\item[(M5)] The invariants of Claim~\ref{cl:lvmeascong}(5) and Remark~\ref{rmk:lvmeasplus}(iii) per branch and sign, as tables (recorded values; gating them gates the code against R20-1, not
  the observation).
\item[(M6)] Negative controls, each gated as \emph{failing} its positive check: no $\ell^{k-2}$ (distribution fails at $m=1$ at every cell); the wrong root $\alpha(1+\ell)$; the wrong twist
  $G_{a+3}$; the wrong sign $+s$; at $(282,37)$ and $(226,103)$ the separator $T_2+3T_3+7T_5$ must be \emph{refused} with multiplicity~$2$.
\item[(M7)] (v1.20.0-pre.5, slot-20 stretch~S1.) The record \texttt{lfunction/v1} of \S\ref{sec:lvlfunction} at the five cells, both signs, every branch of the right parity: values, truncation
  against level $M-1$ (Lemma~\ref{lem:lvlfunprec}), $(\mu,\lambda)$ from the coefficients, determinism ($42$ checks; the default arm becomes $220$ checks, $\approx4.5$\,s, $37$\,MB). Output
  lines \texttt{MEAS\_LFUN}. Its checks are not extended to the heavy cells (the level-$3$ cell at $103$ has $G=10\,609$); under
  \texttt{HECKE\_LVALUE\_MEASURE\_HEAVY=1} the $42$ default checks run unchanged.
\item[(M691)] (v1.20.0-pre.5, slot-20 stretch~S2; env-armed, read-only.) \texttt{HECKE\_LVALUE\_MEASURE\_691=1} prints \texttt{EFFECTIVE HeckeLValueMeasureValidation arm=ELL691 cells=1} and runs
  Claim~\ref{cl:lvmeas691}: the two Kubota--Leopoldt routes at $691$, level $2$ ($344$ classes); the cell $(10,691)$ at $R=2$ (the in-TU twisted Bernoulli table needs $\ell^{R+4}<2^{62}$), level $2$,
  through (M1), (M3)--(M7) exactly as at the other Eisenstein cells ($d=1$); $a_p\equiv\tau(p)\pmod{691^{2}}$; the oracle's $s$, $\mu^{\pm}(2+\ell^{2})$ and FNVs (\texttt{MeasureFrozenS2.inc});
  $v=0$ at the $22$ pole points. $43$ checks ($263$ with the default arm), $0$ failing at the first run; $\approx34$\,s of the $38.5$\,s and $66$\,MB (beyond the \texttt{standard} budget of $15$\,s, hence
  env-armed). Output line \texttt{MEAS\_691}.
\end{enumerate}
Output lines \texttt{MEAS\_CELL}, \texttt{MEAS\_ID}, \texttt{MEAS\_EIS}, \texttt{MEAS\_KL}, \texttt{MEAS\_INV}, \texttt{MEAS\_NEG}, and \texttt{MEAS\_ARMS checks=<c> failing=<f>}. A read-only env arm
\texttt{HECKE\_LVALUE\_MEASURE\_HEAVY=1} prints \texttt{EFFECTIVE HeckeLValueMeasureValidation arm=HEAVY cells=3} and adds $(226,103)$ at level~$3$, $(282,37)$ at $R=5$ and every
character in (M1) ($\le60$\,s, $\le2$\,GB). \emph{Exit criteria of WP-MEAS}: WP-MEAS's predictions (Q1--Q9 of its spec, \S5) committed before its first run; \texttt{failing=0}; the
roster $181=66/98/17$ (\texttt{--all} runs $164$), \texttt{EXPECT\_NAMES} $181$, the [CLI] tables regenerated in the same commit and again after the merge, the catalogue entry; no golden
line, digest, record key or schema key moves; a disagreement with R20-1 is reported as a finding against this section, never re-pinned.
\end{checkitem}

\subsection{The depth-2 triangle, the change of lattice, and \texorpdfstring{$\iota$}{iota} over \texorpdfstring{$\Z/\ell^{k}$}{Z/l\^{}k}}\label{sec:lvdepth2}
(New at v1.20.0-pre.1, hecke-extensions slot~20, WP-2; the contract text of slot-20 WP-LV3, which merges only after it (rule~21), and of slot-20 R20-3's constant. Placed
at the end of \S\ref{sec:lvalue}.)

\begin{claim}[The three kinds at depth~2]\label{cl:lvtriangle}
At the two depth-$2$ rungs $(n,\ell)=(282,37)$ ($k=284$) and $(226,103)$ ($k=228$), $d=v(B_k/k)=2$, $n\ge\ell-1$, on the $\m_{\mathrm{Eis}}$-block of $H^{1}$, cut by the \emph{joint} ideal
$\m_{\mathrm{Eis}}=(\ell,\,T_p-1-p^{k-1}:p=2,3,5)$ (v1.20.0-pre.6, slot-20 WP-LV3: $(\ell,T_2-1-2^{k-1})$ alone is too large at $(282,37)$ --- a cusp form there has $a_2\equiv1+2^{k-1}$, and the
$T_2$-block has rank~$3$ on the $p_0$ sign):
\begin{enumerate}
\item[(C)] \textsf{congruence}: $\#Q_m=\ell^{\min(m,2)}$, $m=1,2,3$ --- \CMPD{} (slot-19 WP-LV2 and, independently, R19-7; Remark~\ref{rmk:lvgates});
\item[(L)] \textsf{linear}: (L1) to depth exactly $2$, $r^{+}\equiv u\cdot p_0\pmod{\ell^{2}}$ and not modulo $\ell^{3}$ --- \CMPD{} (v1.20.0-pre.6; \CONJ{} through pre.5): slot-20 R20-1 (Python,
  predictions committed first; Remark~\ref{rmk:lvmeasplus}(i), Claim~\ref{cl:lvmeascong}(2)) and slot-20 WP-LV3 (C++, on $W$ over $\Z/\ell^{3}$ and on the C++ cone over $\Z/\ell^{3}$); and
  Remark~\ref{rmk:lvlegc}'s item~3 to depth exactly $2$ --- \CONJ{} as an index statement of this route (R20-1
  only, one route, not gated by WP-LV3; the congruence itself is \CMPD{} as a run, Claim~\ref{cl:lvmeascong}(2));
\item[(H)] \textsf{bilinear}: (H2) on the mixed lattice $H^{1}_c(\Ga;\Mn)\times H^{1}(\Ga;\Dn)$, the value on integral eigenlines has $v=2$ on the sign carrying the boundary class
  $p_0$ and $v=0$ on the other --- \CMPD{} (v1.20.0-pre.6, slot-20 WP-LV3's generic reader, every value equal to the congruence module of its lattice read from $A_c$ alone, also for the boundary
  system; \CONJ{} through pre.5: slot-19 R19-9~\S8's prediction, which it confirms).
\end{enumerate}
\textbf{The triangle} (\CMPD{} since v1.20.0-pre.6; \CONJ{} through pre.5): the three depths are equal, to $d$ --- WP-LV3's cross-kind gate, congruence $=$ bilinear $=$ linear $=v(B_k/k)$,
holds at $(282,37)$, $(226,103)$ ($2$) and at the depth-$1$ controls $(30,37)$, $(10,691)$ ($1$), and in the heavy arm against [C]'s $H^{-}$ ($\#Q_3$ of the \textsf{congruence} route) and (L1) on the C++
cone (Check~\ref{chk:lvtriangle}, outcome). As a statement for every depth-$2$ rung it stays \CONJ. At depth~$1$ the three kinds agree at the cells where each was computed (\CMPD; Remarks~\ref{rmk:lvgates},
\ref{rmk:lvlegc}, \ref{rmk:lvshapes}~(H)). Gated by Check~\ref{chk:lvtriangle}; a disagreement there is a finding against this claim and against Remark~\ref{rmk:lvshapes}, never a
re-pin (the first suspect is the lattice: mixed versus $H^{1}_{!}$; the second the change of lattice, Definition~\ref{def:lvlattice}).
\end{claim}

\begin{definition}[The change-of-lattice constant for $n\ge\ell$]\label{def:lvlattice}
\SPEC. Let $\theta:=\Phi_n\colon\Mn\to\Dn$ be the comparison map of \HA{lem:symdual} (Gram matrix $(-1)^{t}t!\,(n-t)!$; the map of Remark~\ref{rmk:lvshapes}~(H), Hida's form
$G_{!}(P,Q)=\langle P,\theta\iota_{!}Q\rangle$, whose code \texttt{HidaPairing::theta} is $\Phi_n/n!$, the same up to the unit $n!$ when $n<\ell$): an isomorphism after $\otimes\Q$, and over
$\Z_\ell$ iff $n<\ell$ (\HA{lem:symdual}: $\Phi_n\otimes\Z_\ell$ is an isomorphism iff $n!\in\Z_\ell^{\times}$). Let $f$ be a Hecke eigensystem with values in $\mathcal O$ (a finite extension of $\Z_\ell$, uniformiser $\pi$), $\varepsilon$ a sign, and for
a $\Z_\ell$-lattice $L$ in a Hecke module write $L[f]^{\varepsilon}:=L_{\mathcal O}^{\varepsilon}\cap(\text{the }f\text{-eigenline of }L^{\varepsilon}\otimes\mathrm{Frac}\,\mathcal O)$ (a free
$\mathcal O$-module of rank~$1$ where $f$ occurs with multiplicity one in that sign). The \emph{change-of-lattice constant} is
\[
  c^{\varepsilon}_f\;:=\;\mathrm{length}_{\mathcal O}\Bigl(H^{1}(\Ga;\Dn)/\mathrm{tors}\,[f]^{\varepsilon}\Big/\theta_\ast\bigl(H^{1}(\Ga;\Mn)/\mathrm{tors}\,[f]^{\varepsilon}\bigr)\Bigr),
\]
defined where $\theta_\ast$ carries the $f$-line of $H^{1}(\Ga;\Mn\otimes\Q_\ell)$ onto that of $H^{1}(\Ga;\Dn\otimes\Q_\ell)$. \emph{This holds} (v1.20.0-pre.2, F1 m-E5; \PROVED{} from [C]):
by \HA{lem:symdual}, $\Phi_n\circ\rho_n(g)=\rho_n^{\vee}(g)\circ\Phi_n$ for every $g\in\mathrm{Mat}_2(\Z)$, untwisted, so $\Phi_n$ commutes on cochains with $T_p$ formed from the same cosets,
i.e.\ with $A(\Dn)(T_p)$ of \HA{sec:symconv}'s code map (\texttt{HeckeOperatorOnK}'s \texttt{FibreRepresentation}; \texttt{HeckeDualityValidation} asserts the intertwining
untwisted), and with $\iota$ (\HA{sec:involution}); $\theta_\ast\otimes\Q_\ell$ is then a Hecke- and $\iota$-equivariant isomorphism and carries each eigenline of each sign onto the
corresponding one. It is no longer an input (Proposition~\ref{prop:lvlattice}(1)); R20-3 checked that its prototype's $T_p$ on $H^{1}(\Ga;\Dn)$ is that normalisation (\texttt{EQUIV} $0$ failures at four cells; a $\det$-twist would break it). $c^{\varepsilon}_f=0$ for $n<\ell$ (\PROVED: $\theta$ is then an isomorphism over $\Z_\ell$). It converts a pairing value on $H^{1}_c(\Mn)\times H^{1}(\Mn)$ (through $\theta$)
into one on the mixed lattice of (H2), and is what slot-19 WP-BIL printed as ``open''. \emph{Status}: at v1.20.0-pre.1 the value was \textbf{open}; since v1.20.0-pre.6 it is given by Proposition~\ref{prop:lvlattice}(2) and Remark~\ref{rmk:lvlatticevalues} (slot-20 R20-3,
theory and two prototypes, predictions first, \texttt{92336edc}), and slot-20 WP-LV3 computes it in \texttt{NormalisationRecord} for $n\ge\ell$ with $\theta=\Phi_n$ (the record's row ``change of
lattice $c_f$ (\texttt{def:lvlattice}; $\theta=\Phi_n$, $n\ge\ell$)'', printing $e\,v_\ell(n!)$ and $d_f$ separately; no longer ``open'').
\end{definition}

\begin{proposition}[The change-of-lattice constant]\label{prop:lvlattice}
\PROVED{} (v1.20.0-pre.6; slot-20 R20-3, \texttt{docs/implementation/slot20/R20-3\_lattice\_constant.md}~\S2). Take $\theta=\Phi_n$ of \HA{lem:symdual} and let
$V'_n:=\bigoplus_t\binom nt\Z_\ell u_t=n!\,\Phi_n^{-1}(\Dn)\subset\Mn$ (a $\De$-stable lattice, since $\Dn$ is $\De$-stable under $\rho_n^{\vee}$
and $\Phi_n$ is $\De$-equivariant and injective; $\ell^{m}\Mn\subset V'_n\subset\Mn$ with $m=\max_t v_\ell\binom nt$, the maximal number of base-$\ell$
carries in $t+(n-t)$). It is \emph{not} the $\Z_\ell$-span of the $n$-th powers $(ax+by)^{n}$ for $n\ge\ell+1$: modulo $\ell$ those span at most
$\ell+1$ dimensions, since $a^{n-1}b\equiv a^{n-\ell}b^{\ell}$ on $\F_\ell^{2}$ (v1.20.0-pre.7, F1$'$ M-E1: through pre.6 the definition read
``$\sum_{a,b}\Z_\ell(ax+by)^n=\bigoplus_t\binom nt\Z_\ell u_t$'', false for $n\ge\ell+1$; every value and the code concern the diagonal lattice, so nothing computed moves).
\begin{enumerate}
\item $\theta_\ast$ commutes with every $T_p$ and with $\iota^\ast$, so the input of Definition~\ref{def:lvlattice} holds for every $f$ of multiplicity one in its sign
  (\HA{lem:symdual}: $\Phi_n\circ\rho_n(g)=\rho_n^{\vee}(g)\circ\Phi_n$ for all $g\in\mathrm{Mat}_2(\Z)$, with $\rho_n^{\vee}(g)=\rho_n(\mathrm{adj}\,g)^{\mathsf T}$; double-coset operators are natural in
  $\De$-module maps).
\item $\Phi_n^{-1}(\Dn)=(n!)^{-1}V'_n$ (as $t!\,(n-t)!=n!/\binom nt$), hence
  \[ c^{\varepsilon}_f=e\,v_\ell(n!)-d^{\varepsilon}_f,\qquad d^{\varepsilon}_f:=\mathrm{length}_{\mathcal O}\Bigl(H^{1}(\Ga;\Mn)/\mathrm{tors}\,[f]^{\varepsilon}\Big/\mathrm{im}\,H^{1}(\Ga;V'_n)/\mathrm{tors}\,[f]^{\varepsilon}\Bigr),\]
  $e=v_\pi(\ell)$, and $0\le d^{\varepsilon}_f\le e\,m$; in particular $c_f^{\varepsilon}=e\,v_\ell(n!)$ for every $f$ when all base-$\ell$ digits of $n$ below the leading one equal $\ell-1$ ($V'_n=\Mn$), and
  $0\le d_f^{\varepsilon}\le e$ for $\ell\le n<\ell^{2}$.
\item $d^{\varepsilon}_f=0$ whenever $H^{1}(\Ga;\Mn/V'_n)^{\varepsilon}_{\m_f}=0$ (the sequence $H^{1}(V'_n)\to H^{1}(\Mn)\to H^{1}(\Mn/V'_n)\to H^{2}(V'_n)=0$ ($\ell\ge5$: $\SL_2(\Z)=\Z/4\ast_{\Z/2}\Z/6$ and the orders $4$, $6$ are units, Mayer--Vietoris), localised at $\m_f$ and taken on the
  $\varepsilon$-part: a surjection of free modules of equal rank).
\end{enumerate}
\end{proposition}

\begin{remark}[Values]\label{rmk:lvlatticevalues}
\CMPD{} (v1.20.0-pre.6; slot-20 R20-3, two routes: the transport by $\Phi_n$ on $H^{1}_c$ (\texttt{r203/lattice.py}), and the definition on $H^{1}(\Ga;\Mn)$ with the group-cocycle $T_p$
(\texttt{r203/direct.py}); $49/49$ eigenlines agree at four cells; and slot-20 WP-LV3's C++ \texttt{NormalisationRecord}, which reads $d_f$ on the $\Dn$ side through the perfect pairing with
$H^{1}_c(\Dn)$, a third route). On the $\m_{\mathrm{Eis}}$-blocks $d_f=0$ on both signs, for the cusp form and for $E_k$: $c^{\pm}_f=7$ at $(282,37)$, $2$ at $(226,103)$; at $(258,37)$ $V'_n=\Mn$ and
$c_f=6e$ for every $f$ (R19-1's quartic, $e=2$: $v_\pi=12$, $v_{37}=6$, all four forms, both signs); $0$ for $n<\ell$ ($(30,37)$, $(10,691)$). Generic eigenlines can carry $d_f=1$: at $(40,37)$ five
($\lambda\equiv3,7,29$ on $+$, $3,29$ on $-$, with $\lambda$ the eigenvalue of $T_2+5T_3+11T_5+13T_7$; WP-LV3's negative control \texttt{LV3\_LATTICE\_NEG}). For $\ell\le n<\ell^{2}$,
$\Mn/V'_n\otimes\F_\ell\cong\bigoplus_j\Sym^{a_j}\otimes{\det}^{b_j}$ ($a_j=\ell-3-n_0+n_1-2j$, $b_j=n_0+1+j$, $n=n_1\ell+n_0$; the module structure recalled from Glover's filtration, \STD, the graded form
\PROVED, the dimension counts \CMPD), and $d_f^{\varepsilon}=1$ exactly on the $\omega^{b_j}$-twists of the weight-$(a_j+2)$ systems, the sign changed by $(-1)^{b_j}$ (\CONJ{} in general; \CMPD{} line by
line at $(40,37)$, $(282,37)$, $(226,103)$, with $\sum_m(D_m^{+}+D_m^{-})=\dim H^{1}(\Mn/V'_n)$ at all five cells $n\ge\ell$). Through $\theta_{R199}=\Phi_n/n!$ (slot-19 R19-9; the code's
\texttt{HidaPairing::theta}, not integral for $n\ge\ell$) the same index is $-d_f$. Whole-lattice indices (WP-LV3, \texttt{LV3\_THETA\_INDEX},
$v_\ell[H^{1}(\Dn)/\mathrm{tors}:\Phi_{n,\ast}H^{1}(\Mn)/\mathrm{tors}]$\,/\,$v_\ell[H^{1}(\Mn)/\mathrm{tors}:\Theta H^{1}(\Dn)/\mathrm{tors}]$, $\Theta=(\Phi_n/n!)^{-1}$): $(10,5)$ $4/2$, $(10,7)$ $2/1$,
$(12,5)$ $2/0$, $(24,7)$ $8/1$, $(40,37)$ $2/5$, $(72,37)$ $11/0$, $(30,37)$ $0/0$, $(258,37)$ $258/0$.
\end{remark}

\begin{remark}[$\iota$ over $\Z/\ell^{k}$]\label{rmk:lviotak}
Nothing new is claimed: [C] already gives $\iota^{\ast}$ over every coefficient ring (\HA{def:iota}), its commutation with every $T_p$ (\HA{prop:iotacommutes}) and the eigenspace
decomposition whenever $2$ is a unit (\HA{prop:iotasplit}) --- so over $\Z/\ell^{k}$, $\ell$ odd, $H^{q}=H^{q,+}\oplus H^{q,-}$, Hecke-stably and compatibly with the reductions
$\Z/\ell^{k}\to\Z/\ell^{k'}$ (\PROVED{} from [C]). What was missing is the \emph{code}: in the slot-19 build the console refuses \texttt{iota} under
\texttt{coefficient\_ring = "Z/ell\^{}k"} (exit~$5$; R19-7~\S6.3 read the sign split off the matrices). Slot-20 WP-LV3(d) makes \texttt{Hecke/HeckeInvolution} and the console honour it,
gated in Check~\ref{chk:lvtriangle}(d). \emph{Built} (v1.20.0-pre.6): \texttt{hecke::BuildRingInvolution} applies the field object's cochain formula ($\iota^{\ast}=-\iota^{\ast}_{\mathrm{geo}}$,
$\rho_n(\iota)=\mathrm{diag}((-1)^{n-j})$) on the module basis of every degree, with the always-on checks $\iota^{2}=\mathrm{id}$ and $\mathrm{length}\,H^{+}+\mathrm{length}\,H^{-}=\mathrm{length}\,H$
($\ell\ge5$; $\ell\in\{2,3\}$ refused by name); the record is \texttt{hecke\_involution\_ring/v1} (a new schema, not golden); the console list is \S\ref{sec:ring}~(h).
\end{remark}

\begin{checkitem}[The gates of WP-LV3]\label{chk:lvtriangle}
\SPEC{} (the checks; expected values from the sources named, committed as WP-LV3's predictions before its first run). All arms are inside \texttt{HeckeLValueValidation} (no new name;
its default wall stays $\le$ its slot-19 figure $+10$\,s, heavier cells behind the read-only \texttt{HECKE\_LVALUE\_HEAVY=1}).
\begin{enumerate}
\item[(a)] \textbf{The generic $\m_{\mathrm{Eis}}$-block reader.} Given a Hecke block over $\Z/\ell^{m}$ and the Hensel factor $P$ of $\mathrm{charpoly}(T_{p^\ast})$ at $\m_{\mathrm{Eis}}$, it
  returns the eigensystem over $\mathcal O=\Z_\ell[x]/(P)$ (a product of local orders in general; $\Q(\alpha)$ and $v_\pi$ through \texttt{NumberField}, \S\ref{sec:lvnf}) and the
  per-eigenform lines. \emph{As built} (v1.20.0-pre.6): the block is that of $t=T_2-\lambda_0$, at $\m_{\mathrm{Eis}}$ the joint block of $T_2,T_3,T_5$ (Claim~\ref{cl:lvtriangle}); the roots of
  its characteristic polynomial $P\equiv x^{r}\pmod\ell$ are read from the Newton polygon of $P$ and its residual polynomials, each in a \emph{tame, totally ramified} order
  $\mathcal O=\Z_\ell[\pi]/(\pi^{e}-\ell u)$, $e\mid\ell-1$, residue degree $f=1$, refined digit by digit; an unramified residue extension ($f>1$), $e\nmid\ell-1$, a non-integral refinement slope
  and a multiple root to the precision are \emph{refused by name}, never read. No global $\Q(\alpha)$ is formed at the gate cells. At $(258,37)$ $\m=(37,T_2-29)$, WP-BIL's and R19-1's convention. First gate: slot-19 WP-BIL's values at $(258,37)$ (Remark~\ref{rmk:lvshapes}~(H): mixed $3/2$, through $H^{1}_{!}(D)$ $3/3$, $\sum_fv_{37}=6$) reproduced bit for
  bit; the quartic-specific path stays as a cross-check arm until they match. The reader takes the model as a parameter (no hard-wired seven-vertex model, no trivial centre).
\item[(b)] \textbf{The triangle} (Claim~\ref{cl:lvtriangle}) at $(282,37)$ and $(226,103)$: $\#Q_2=\ell^{2}$; (H2) $v=(2,0)$ on (the $p_0$ sign, the other); (L1) depth exactly $2$; and
  their agreement --- the \emph{cross-kind gate} --- at these two cells and at two depth-$1$ controls, $(10,691)$ (a gate cell of all three kinds) and $(30,37)$ (expected depth $1$
  in all three readings). A disagreement is a finding for Claim~\ref{cl:lvtriangle}.
\item[(c)] \textbf{The constant} $c^{\varepsilon}_f$ of Definition~\ref{def:lvlattice} in \texttt{NormalisationRecord} for $n\ge\ell$, against R20-3's predicted values.
\item[(d)] \textbf{$\iota$ over $\Z/\ell^{k}$} (Remark~\ref{rmk:lviotak}): \texttt{Hecke/HeckeInvolution} and the console accept $\mathsf R=\Z/\ell^{k}$ (exit~$5$ removed); a new console script
  (the next number, frozen \texttt{.expected}: an addition, not a re-baseline) reads $H^{\pm}$ at $(282,37)$ over $\Z/37^{k}$; expected: free of ranks $24$ ($+$, the Eisenstein side) and
  $23$ ($-$) (slot-19 WP-LV2), the $T_p$-blocks equal to R19-7~\S5's matrix-read split. \emph{Amended} (v1.20.0-pre.6, lead decision on WP-LV3's report): the script
  \texttt{049\_lv3\_iota\_ring.oc} reads $(30,37)$ and the torsion cell $(38,37)$ over $\Z/37^{3}$ (the $\Z/37^{k}$ cohomology at $n=282$ costs $\approx25$\,s per run and every numbered script runs
  several times inside \texttt{ConsoleDeterminismValidation}'s $45$\,s threshold); the $(282,37)$ reading runs in the heavy arm (\texttt{LV3\_IOTA\_RING\_282}). \texttt{ConsoleDeterminismValidation}'s pinned C4~(a) digest does not move; the unpinned
  all-scripts digest moves, and the report says so.
\item[(e)] \textbf{The twisted depth-$2$ rung} $(\ell;k,j)=(37;62,6)$ of Claim~\ref{cl:lvtwist} ($v(B_{62,\omega^{6}}/62)=2$ by two routes, Remark~\ref{rmk:lvgates}): first the character model
  at $n=20$, $40$ (RSS and wall measured), then the curve's prediction for $n=60$, then the run behind \texttt{HECKE\_LVALUE\_HEAVY} under a cap $\le24$\,GB; expected: rank~$1$,
  saturated, $\#Q_m=\ell^{\min(m,2)}$. If the curve predicts more than $24$\,GB the cell is not run and is recorded with its curve (or handed to a cluster kit).
\end{enumerate}
\emph{Exit criteria of WP-LV3}: (a) held before (b) is trusted; (b)--(e) held or each failure reported as a finding the same turn; zero golden movement, the four digests unchanged
(the console addition excepted as stated in (d)).

\emph{Outcome} (v1.20.0-pre.6; \CMPD, slot-20 WP-LV3, \texttt{docs/\allowbreak implementation/\allowbreak slot20/\allowbreak WP-LV3\_report.md}; \texttt{Heck\allowbreak eLValue\allowbreak Validation} $1471$ checks by default, $1614$ heavy, $0$ failing; zero golden
movement, the digests unchanged, the console C4~(a) digest unchanged and the unpinned all-scripts digest moved, as~(d) said). (a) At $(258,37)$ the generic reader equals WP-BIL's bit for bit (cell
fields, block ranks, the characteristic polynomial modulo $37^{8}$, the Gram pivots; R19-1's factors $P_1$, $P_2$ \emph{recovered}, two roots each; per factor (mixed, P9, $v_\pi P'$, content) geo-even
$(3,3,3,0)$, geo-odd $(2,3,3,1)$). (b) The cross-kind gate holds at all four cells: congruence $=$ bilinear ($p_0$ sign) $=$ linear (L1) $=v(B_k/k)$, $2$ at $(282,37)$ and $(226,103)$, $1$ at
$(30,37)$ and $(10,691)$, the bilinear value $0$ on the other sign; in the heavy arm $\#Q_3$ on [C]'s $H^{-}$ and (L1) on the C++ cone over $\Z/\ell^{3}$ give $2$ at both depth-$2$ cells. No finding.
(c) Remark~\ref{rmk:lvlatticevalues}'s values, all equal to R20-3's predictions. (d) $\iota$ over $\Z/37^{3}$ at $(22,37)$, $(30,37)$, $(38,37)$: \texttt{uint32} $=$ \texttt{uint64}; at $k=1$ the field
object's dimensions; at $(38,37)$ the torsion on $+$; heavy: at $(282,37)$ $H^{1}_{+}=(\Z/37^{3})^{24}$, $H^{1}_{-}=(\Z/37^{3})^{23}$. (e) The twisted rung $(37;62,6)$ at $n=60$ (character model, $K=4$):
predicted from the growth curve ($n=20$: $2.10$\,GB; $n=40$: $7.48$\,GB) at $17.5$\,GB, $\approx7.3$\,h, under a $22$\,GB cap; measured $16.43$\,GB, $7$\,h\,$18$\,min; rank~$386$ ($194+192$),
$f^{-}=1$, saturated ($c^{\mathrm{sat}}=0$), $\#Q_m=37^{\min(m,2)}$ at $m=1,2,3$, $v(B_{62,\omega^{6}}/62)=2$, $\rank H^{+}=\rank H^{-}+2$, the negative controls $Q=0$ --- the first twisted rung at
depth~$2$, \CMPD{} for Claim~\ref{cl:lvtwist} (all $35$ checks of the arm).
\end{checkitem}

\subsection{The truncated one-variable \texorpdfstring{$\ell$}{l}-adic \texorpdfstring{$L$}{L}-function as a record}\label{sec:lvlfunction}
(New at v1.20.0-pre.5, hecke-extensions slot~20, stretch~S1, the owner-Q3 roadmap's P1; placed at the end of \S\ref{sec:lvalue}, so that no statement number moves. \emph{No new
mathematics}: this subsection writes the branch series of Proposition~\ref{prop:lvmtt}(5) as a record with a certified precision per coefficient. Predictions
\texttt{docs/implementation/slot20/s1s2/predictions.md}, committed before the first run (\texttt{c361d65c}).) Notation of \S\ref{sec:lvmeasure}; $\mu=\mu_x$ the measure of
Definition~\ref{def:lvmtt} of an ordinary eigensymbol, read at level $M\ge2$ modulo $\ell^{R}$; $G:=\ell^{M-1}$; $\omega_r(T):=(1+T)^{\ell^{r}}-1$.

\begin{definition}[The truncated $\ell$-adic $L$-function of a branch]\label{def:lvlfun}
\SPEC. For a branch $a\in\Z/(\ell-1)$ the \emph{truncated $\ell$-adic $L$-function} is the image of $\mathcal L_{x,a}(T)=\int\omega^{a}(x)(1+T)^{\ell_u(x)}\,d\mu_x$ in
$\Lambda/(\ell^{R},\omega_{M-1}(T))$, $\Lambda=\Z_\ell[\![T]\!]$, written in the basis $T^{t}$, $0\le t<G$:
\[
  \mathcal L^{(M,R)}_{x,a}(T)=\sum_{t<G}c_tT^{t},\qquad c_t=\sum_{i<G}g_i\binom it\bmod\ell^{R},
\]
$g$ the branch element of Proposition~\ref{prop:lvmtt}(5) at level $M$ ($(1+T)^{i}\leftrightarrow[u^{i}]$; $\sum_tc_tT^{t}$ is the unique representative of degree $<G$, since $\omega_{M-1}$ is
monic of degree $G$). Its \emph{values} are $\mathcal L^{(M,R)}_{x,a}(u^{j}-1)=\int\omega^{a-j}(x)x^{j}\,d\mu_x$ modulo $\ell^{\min(M,R)}$ ($0\le j\le n$: the critical moments of
Proposition~\ref{prop:lvmtt}(4)); its $\mu$ and $\lambda$ are those of Proposition~\ref{prop:lvmtt}(5) (read on $t<\min(G,\ell)$), and $\lambda_{\bmod\ell}$ is the first $t<G$ with $c_t$ a unit.
\end{definition}

\begin{lemma}[The certified precision of the coefficients]\label{lem:lvlfunprec}
\PROVED{} (elementary). Let $F=\sum_tF_tT^{t}\in\Lambda$ and $F=Q\,\omega_r+\mathrm{Rem}$ with $\deg\mathrm{Rem}<\ell^{r}$ (Weierstrass division by the distinguished polynomial $\omega_r$; $Q\in\Lambda$).
Then $\mathrm{Rem}_0=F_0$ and, for $1\le t<\ell^{r}$, $\mathrm{Rem}_t\equiv F_t\pmod{\ell^{\,r-\lfloor\log_\ell t\rfloor}}$. Hence the coefficients $c_t$ of Definition~\ref{def:lvlfun} agree with the
$T^{t}$-coefficients of the true series $\mathcal L_{x,a}$ modulo $\ell^{\mathrm{prec}(t)}$,
\[
  \mathrm{prec}(0)=R,\qquad \mathrm{prec}(t)=\min\bigl(R,\;M-1-\lfloor\log_\ell t\rfloor\bigr)\quad(1\le t<G):
\]
modulo $\ell^{\min(R,M-1)}$ for $t<\ell$ (Proposition~\ref{prop:lvmtt}(5)) and modulo $\ell$ for every $t<G$.
\end{lemma}

\begin{proof}
$\mathrm{Rem}_t=F_t-\sum_{1\le s\le t}Q_{t-s}\binom{\ell^{r}}{s}$ for $t<\ell^{r}$ (the coefficient $\omega_{r,s}=\binom{\ell^{r}}s$, $1\le s<\ell^{r}$), and
$v\bigl(\binom{\ell^{r}}s\bigr)=r-v(s)\ge r-\lfloor\log_\ell s\rfloor\ge r-\lfloor\log_\ell t\rfloor$. Apply it to $F=\mathcal L_{x,a}$, $r=M-1$: the level-$M$ branch element is the image of $\mathcal L_{x,a}$
modulo $\omega_{M-1}$ (Proposition~\ref{prop:lvmtt}(5)), exact modulo $\ell^{R}$.
\end{proof}

\begin{remark}[The record \texttt{lfunction/v1} and the verb \texttt{lfunction}]\label{rmk:lvlfunschema}
\SPEC{} (slot-20 S1; code \texttt{HeckeExt/LValue/LFunctionRecord.\{h,cpp\}}, the console verb in \texttt{HeckeExt/RegisterHeckeExtensionVerbs.cpp}, [P]~v3.11.1~\S3.6). One record per
(form, sign, truncation), one entry per branch; canonical JSON, keys in this order, no whitespace; identified by the FNV-1a-64 (\texttt{Fnv.h}) of that text. It is a new schema: no key
of \texttt{hecke\_ext/v5} or of any golden object moves, and the record is not golden.
\begin{center}\scriptsize
\begin{tabular}{@{}lp{0.72\textwidth}@{}}
\toprule
key & content \\
\midrule
\texttt{schema} & \texttt{"lfunction/v1"} \\
\texttt{n}, \texttt{k}, \texttt{ell}, \texttt{R}, \texttt{M}, \texttt{G}, \texttt{u} & the cell $(n,\ell)$, $k=n+2$, values modulo $\ell^{R}$, level $M$, $G=\ell^{M-1}$, $u=1+\ell$ \\
\texttt{sign}, \texttt{geometric} & \texttt{"minus"} ($r^{-}$, odd coordinates) or \texttt{"plus"} ($r^{+}$); the geometric sign \texttt{"odd"}\,/\,\texttt{"even"} (\S\ref{sec:lvmeasure}) \\
\texttt{form} & \texttt{separator} (the operator $\sum c_pT_p$), \texttt{target\_root} (its root modulo $\ell$), \texttt{eisenstein\_congruent}, \texttt{a\_p} ($p\in\{2,3,5,7,\ell\}$ and the separator's primes,
  modulo $\ell^{R}$), \texttt{a\_ell}, \texttt{alpha} (the unit root), \texttt{symbol} (the eigensymbol's normalisation: $r^{-}$ first unit coordinate $1$, $r^{+}_0=1$) \\
\texttt{truncation} & the ring $\Z/\ell^{R}[T]/(\omega_{M-1}(T))$, $\omega$, the basis, the precision rule of Lemma~\ref{lem:lvlfunprec} (text) \\
\texttt{normalisation} & \texttt{ModularSymbolMeasure::normalisation()} (Proposition~\ref{prop:lvmtt}(4) and its analytic form; the ambiguity $\mathcal O^{\times}$ of the eigensymbol's scaling) \\
\texttt{branches[]} & per branch: \texttt{a}, \texttt{vanishes} (the wrong parity), \texttt{mu}, \texttt{lambda}, \texttt{readable}, \texttt{lambda\_mod\_ell}, \texttt{coefficients} ($c_0,\dots,c_{G-1}$ modulo
  $\ell^{R}$), \texttt{precision} ($\mathrm{prec}(0),\dots,\mathrm{prec}(G-1)$), \texttt{values\_modulus} ($\ell^{\min(M,R)}$), \texttt{critical\_values} ($\mathcal L^{(M,R)}_{x,a}(u^{j}-1)$, $0\le j\le n$) \\
\bottomrule
\end{tabular}
\end{center}
The console verb \texttt{lfunction(n, ell, a, M = 2, R = 4, sign = "minus", sep = nil)} returns one branch of the record as a \texttt{Record} (the keys above, flattened, plus \texttt{hash} and an
\texttt{LFUNCTION} line). The eigensymbol is picked as in Check~\ref{chk:lvmeas}: the root of $\mathrm{charpoly}(\mathrm{sep}\mid W^{-})$ congruent to the Eisenstein value if it occurs, else the
first root modulo $\ell$, simple modulo $\ell$ (else refused); the measure is the one \texttt{BirchModularSymbolRoute::measure()} returns on the class with this symbol (Check~\ref{chk:lvmeas}(M0)).
Refusals (exit~$3$): $\ell<5$ or not prime, $n$ odd, $M<2$, $\ell^{R}\ge2^{62}$, $\ell^{M}>4\cdot10^{6}$ balls (the read-only budget), a non-simple root, no unit $r^{+}_0$, $a_\ell\equiv0$. The verb is
tagged \texttt{read\_only} (pure in-memory arithmetic; \texttt{Console/Bindings/ReadOnlyTags.cpp} and \texttt{Console/serve/read\_only\_verbs.tsv}).
\end{remark}

\noindent\emph{The gate} is Check~\ref{chk:lvmeas}(M7) (v1.20.0-pre.5): at the five default cells and both signs, on every branch of the record's parity, (1) $\sum_tc_t(u^{j}-1)^{t}$ equals the
group-ring branch value at every $0\le j\le n$ (exact modulo $\ell^{R}$); (2) at the three level-$3$ cells, the level-$3$ element folded modulo $\Ga_u^{\ell}$ equals the level-$2$ element
built from the level-$2$ balls, and the level-$2$ coefficients agree with the level-$3$ ones modulo $\ell^{\mathrm{prec}_2(t)}$ for $t<\ell$ (Lemma~\ref{lem:lvlfunprec} tested against a deeper
truncation; the bound is attained --- the coefficients differ modulo $\ell^{\mathrm{prec}_2(t)+1}$ --- at $628$--$631$ of the $648$ coefficients $1\le t<\ell$ of each sign on which $\mu=0$, so it is sharp there; on the plus side of the Eisenstein cells, where $\mu^{+}=d$, at none); (3) $(\mu,\lambda)$ recomputed from the coefficients equal Proposition~\ref{prop:lvmtt}(5)'s,
and $\lambda_{\bmod\ell}=\lambda$ where $\mu=0$; (4) the canonical text is deterministic. $42$ checks, $0$ failing at the first run (\CMPD).

\begin{claim}[Ramanujan's congruence as a congruence of measures: $\Delta$ at $691$]\label{cl:lvmeas691}
\CMPD{} at $(n,\ell)=(10,691)$, $k=12$, $d=v(B_{12}/12)=1$ (slot-20 stretch~S2; predictions \texttt{docs/implementation/slot20/s1s2/predictions.md} committed first (\texttt{c361d65c});
R20-1's Python oracle at this cell committed before the C++ run (\texttt{s1s2/oracle691.md}, \texttt{da0687d7}); the C++ arm of Check~\ref{chk:lvmeas} (\texttt{HECKE\_LVALUE\_MEASURE\_691=1}) at
$R=2$, level $M=2$). $\Delta$ is ordinary at $691$: $\tau(691)\equiv1\pmod{691}$, $v(\tau(691)-1)=1$ (the $q$-expansion, exact). On the seven-vertex model, with the separator $T_2+3T_3+7T_5$:
\begin{enumerate}
\item $W^{-}$ has rank $1$; the Eisenstein root occurs and is simple; the eigensymbol $r^{-}$ has $a_p\equiv\tau(p)\pmod{691^{2}}$ for $p\in\{2,3,5,7,691\}$ (Ramanujan on the symbol side), and
  $v(a_p-1-p^{11})=1$ at all five; $\alpha\equiv\tau(691)\pmod{691^{2}}$; (L1) $r^{+}\equiv u\,p_0$ and $r^{-}\equiv s\,P^{\mathrm{Eis}}$ to depth exactly~$1$.
\item \textbf{(The congruence of measures.)} Claim~\ref{cl:lvmeascong}(3) holds at this cell: on all $343$ off-pole odd branches,
  $\mathcal L^{-}_{\Delta,a}(T)\equiv-s\,G_{a+1}(u(1+T)-1)\,G_{11-a}(u^{11}(1+T)^{-1}-1)\pmod{691}$ as elements of $\Z/691[\Ga_u/\Ga_u^{691}]$, with $G_b$ from \emph{both} Kubota--Leopoldt routes
  (which agree for all $344$ even classes and give $-(1-691^{k-1})B_k/k$ at $1376/1376$ weights), and at the $3773$ critical moments against exact twisted Bernoulli numbers ($0$ failures).
  Negative controls: the wrong twist $G_{a+3}$ matches $0/342$, the wrong sign $0/343$.
\item \textbf{(The pole branches $a=11$, $689$.)} One factor is $G_0$ and the other is $G_{12}$, the $\Delta$-class. The naive product fails modulo $691$ at all $22$ critical points, and the
  difference is a unit at all $22$ --- the valuation $d-1-v(r)=0$ of R20-1~\S4.4's polar term (Claim~\ref{cl:lvmeascong}(4)), predicted before the run. $\mu=\lambda=0$ on both.
\item \textbf{(Invariants.)} Minus side: $\mu=0$ on every odd branch, $\lambda=1$ exactly on $a=199$ and $a=501$, $\lambda=0$ elsewhere --- predicted from (2): these are the two off-pole branches carrying
  the other irregular class, $G_{200}\cdot G_{502}$ and $G_{502}\cdot G_{200}$, with $\lambda(G_{200})=1$ (\Ucite{thm:u-691classes}); $199+501\equiv n\pmod{690}$. Plus side: $\mu^{+}=1$ on every even
  branch (Remark~\ref{rmk:lvmeasplus}(ii)); the reduced plus measure has $\lambda=1$ exactly on $\{22,350,678\}$, a functional-equation orbit $a\leftrightarrow n-a$ ($350$ self-paired) --- recorded,
  not predicted (Remark~\ref{rmk:lvmeasplus}(iii)).
\end{enumerate}
The expected literature is the same as for Claim~\ref{cl:lvmeascong} (Mazur; Greenberg--Vatsal; held, not yet read; v1.21.0-pre.1). Every value above equals the Python oracle's where the oracle computed it.
\end{claim}

\subsection{Coefficient rings, the two-variable family, the Eisenstein measure, wild characters and Iwasawa invariants}\label{sec:lvslot21}
(New at v1.21.0-pre.1, hecke-extensions slot~21, the $\ell$-adic track P3 and P9; placed at the end of \S\ref{sec:lvalue}, so that no statement number moves. \emph{No statement of
\S\S\ref{sec:lvmeasure}--\ref{sec:lvlfunction} changes}: P3 is a change of coefficients of Definition~\ref{def:lvmtt}, P9 assembles the measures of several weights. Specs, predictions
committed before the code: \texttt{docs/implementation/slot21/P3\_spec.md} (\texttt{bfd63a34}), \texttt{P9\_spec.md} (\texttt{7ff1cb29}); neither is refereed.)

\begin{remark}[The measure over the Hecke order of the eigen-system]\label{rmk:lvcoeffring}
(Slot-21 P3.) Let $\chi=\mathrm{charpoly}(T_2\mid W^{-})$ and $F$ its Hensel factor over $\Z/\ell^{R}$ with $F\equiv\bar g^{e}\pmod\ell$, $\bar g$ irreducible, $\gcd(\bar g,\bar\chi/\bar g^{e})=1$;
$\mathcal O=(\Z/\ell^{R})[x]/(F)$ --- for $e=1$ the unramified $W(\F_{\ell^{f}})/\ell^{R}$, $f=\deg\bar g$; for $f=1$, $e\ge2$ and $F(X+\lambda)$ Eisenstein the totally ramified $\Z_\ell[x]/(F)$. The
eigen-symbol $P_\theta$ of $\theta=x$ on $W^{\pm}\otimes\mathcal O$ is found by Gauss--Jordan with unit pivots (applicable iff $T_2-\bar\theta$ has rank $r-1$ residually: \emph{residual
multiplicity one}, else refused by name), normalised at its first unit coordinate; $a_\ell\in\mathcal O^{\times}$ is required (else refused, \emph{not ordinary}); $\alpha\in\mathcal O^{\times}$ the unit root;
then $\mu_{P_\theta}=\sum_jx^{j}\mu^{(\alpha)}_{P_j}$ by Remark~\ref{rmk:lvshapes}~(M). Proposition~\ref{prop:lvmtt}(1)--(6) hold over $\mathcal O$ (\PROVED; the proof is coefficient-free).
\emph{Computed} (\CMPD: a Python oracle sharing no $\mathcal O$-arithmetic with the C++, its values frozen before the C++ ran; $q$-expansion characteristic polynomials of $T_2$, $T_3$ on $S_k$
as a third, symbol-free route) at $(22,23)$, $(22,29)$, $(26,31)$, $(26,37)$, $(28,29)$ ($W(\F_{\ell^{2}})$), $(34,37)$ ($W(\F_{37^{3}})$), $(26,131)$, $(26,139)$ (totally ramified, $e=2$; $S_{28}$, whose Hecke
field has discriminant divisible by $131$ and $139$), with the degree-one controls $(10,37)$, $(16,37)$ bit-identical to Check~\ref{chk:lvmeas}(M2)'s FNVs: the identities of
Proposition~\ref{prop:lvmtt} with $0$ failures; Galois equivariance $\mu_{\sigma\theta}=\sigma\circ\mu_\theta$ (the eigen-symbol of $\sigma\theta$ recomputed, not conjugated); at the ramified cells a
second construction through WP-LV3's $\m$-block reader in $\Z_\ell[\pi]/(\pi^{e}-\ell u)$, equal to the first through $x\mapsto\lambda_0+\rho$ for both roots $\rho$; $(30,29)$ refused, not ordinary.
\emph{Hida control modulo $\m$} (\CMPD{} at one residually non-split component; the oracle found it before the C++): at $\ell=29$, $\bar g=X^{2}+22X+15$, the weights $k=24,52,80,108$ of one class
modulo $\ell-1$ (ordinary rank~$2$, one Galois orbit at each), the level-$2$ unit-ball measures of $k=52,80,108$ are proportional modulo $\m$ to that of $k=24$ with scalars $19,4,21$ (geometric odd)
and $1,1,1$ (geometric even) in $\F_{29}^{\times}$, $0/812$ balls mismatching each; the residual eigen-system is the same at all four weights. This is the first numerical instance in this project of
the $\Lambda$-adic coherence of measures at a residually non-split component; its natural statement is \Ucite{pr:u-msym}~(U16), which stays \CONJ.
\emph{Scope} (finding F-P3-1, P3's survey, \CMPD{} as a census, one route): at level one, $\ell\le59$ and $n\le130$--$150$, no ordinary eigen-system has a repeated residual $T_2$-eigenvalue,
so the ordinary ramified cells come from residually distinct, characteristic-zero-ramified Hecke fields ($131$, $139$ at $k=28$; $144169$ at $k=24$, the env arm); WP-LV3's ramified
$(258,37)$ is \emph{not} ordinary ($T_{37}$ nilpotent modulo $37$ on $W^{-}$; $k=260\equiv8\pmod{36}$), so it carries no measure.
\end{remark}

\begin{checkitem}[The gate of P3: \texttt{HeckeLValueCoefficientRingValidation}]\label{chk:lvcoeffring}
\SPEC{} (the checks); expected values P3's Python oracle's, frozen in \texttt{CoefficientRingFrozenP3.inc} before the C++. Tier \texttt{standard}; arms (R0) setup and the $q$-expansion
charpolys, (R1) eigen-symbols, (R2) Proposition~\ref{prop:lvmtt} over $\mathcal O$, (R3) frozen measures, (R4) Galois, (R5) the ramified route~B, (R6) Hida modulo $\m$, (R7) degree-one controls and base
change into $(\Z/37^{4})[x]/(x^{2}-2)$, (R8) negative controls and the refusals by name. First run: $603$ checks, $0$ failing, $1.2$\,s, $24$\,MB; env arm \texttt{HECKE\_LVALUE\_COEFFRING\_HEAVY=1}
(adds $(22,144169)$ at its ramified prime and $(34,37)$ at level~$3$): $674$ checks, $0$ failing, $138$\,s, $74$\,MB. No golden line, record key or digest. \emph{P2/P4} (v1.21.0-pre.3; through
pre.2 ``pending''): the Eisenstein measure exact in $\Lambda$ (P2) is Proposition~\ref{prop:lveismeas}; wild characters on the $L$-side (P4) are Remark~\ref{rmk:lvwild}, built over
$\Z_\ell[\zeta_{\ell^{m-1}}]$ only --- the coefficient ring $\mathcal O\otimes\Z_\ell[\zeta_{\ell^{m-1}}]$ of a non-trivial Hecke order $\mathcal O$ is \emph{not} built (P4 adds no tensor-product
constructor; \texttt{P4\_spec.md}~\S6).
\end{checkitem}

\begin{remark}[The two-variable family by interpolation (slot-21 P9)]\label{rmk:lvtwovar}
(Slot-21 P9, \texttt{HeckeLambdaTwoVarValidation}, \texttt{standard}; inputs P7's $\Lambda$-adic $q$-expansions and P8's normalised masses, \Ucite{pr:u-msym}, \Ucite{prop:u-msymnorm}.) For a class
$(\ell,a)$ and an étale cuspidal component, the normalised unit-ball masses $\nu_B(x_k)$ of $N=5$ consecutive weights are interpolated in $X$ (\Ucite{thm:u-interp}), and the branch
elements of Proposition~\ref{prop:lvmtt}(5) and the coefficients of Definition~\ref{def:lvlfun} are formed in $\Lambda_{K,m}(x_c)$: a two-variable $L_a(X,T)$ modulo $(\ell^{K},(X-x_c)^{m},\omega_{M-1}(T))$.
It is the specialisation family of one $\mathbf h(a)_\m$-valued measure \emph{only if} \Ucite{pr:u-msym}~(U16) holds; nothing here proves (U16), and the gates resting on it are findings. Assembly
and the agreement with P1's own pipeline (\S\ref{sec:lvlfunction}) at the computed weights are \PROVED{} statements checked by the code; at the cells $(11,2)$, $(13,0)$, $(23,2)$, $(37,32)$ the
first run held every gate and fired no finding ($71$ checks, $41$ finding gates; \CMPD{} as a run, one C++ route, against Python predictions committed first for 11a):
\begin{enumerate}
\item \textbf{(Numerical control.)} The interpolants predict every branch coordinate at three further weights modulo $\ell^{5}$, and P1's \texttt{lfunction/v1} at the further weights equals
  the predicted $L_a$ up to P1's normaliser modulo $\ell^{5}$ at $(13,0)$ and $(37,32)$ --- evidence for (U16)(b), \CONJ.
\item \textbf{(The $\Lambda$-adic Eisenstein congruence of measures.)} At $(37,32)$, on the component through $\m_{\mathrm{Eis}}$, the cross products of $L_a(X,T)$ with
  $G_{a+1}(u(1+T)-1)\,G_{k-1-a}(u^{k-1}(1+T)^{-1}-1)$ (Claim~\ref{cl:lvmeascong}(3) at each weight) vanish at $X=\alpha_{32}$ --- the zero of $G_{32}$, $\alpha_{32}\equiv281570\pmod{37^{4}}$
  (\Ucite{pr:u-lambda}, (U13)), recomputed from exact Bernoulli numbers --- modulo $37^{3}$ on all $592$ off-pole odd slots (the full Kubota--Leopoldt precision at level $M_E=3$), against
  $37^{1}$ at every integer weight with $L_k=1$ and $37^{2}$ at $k=284$; at $\alpha_{32}+37$ only modulo $37$. In the env arm \texttt{HECKE\_LAMBDA\_TWOVAR\_ALL=1}: $21\,904$ slots at level $3$,
  the same $37^{3}$, and at $(103,24)$ modulo $103^{2}$ at $X=\alpha_{24}\equiv2266\pmod{103^{2}}$ on $5047$ slots. The expected statement, a congruence of $\Lambda$-adic measures modulo
  $(X-\alpha_a)$ (Mazur--Kitagawa, recalled; the interpolation of Claim~\ref{cl:lvmeascong} in $k$), is \CONJ; the run is evidence for it and for (U16).
\item \textbf{(Weight two.)} At $(11,2)$ (the $\Delta$ family) the values predicted at $x_2$ from $k=12,\dots,52$ equal 11a's own normalised MTT masses (Manin symbols on $\Gamma_0(11)$, no
  level-one routine) on every unit ball to level~$3$, both signs, modulo $11^{5}$; the trivial zero $\nu_{\Z_{11}^{\times}}(x_2)\equiv0$ and Greenberg--Stevens' derivative identity hold modulo $11^{5}$
  (\CONJ: they rest on (U16) at $k=2$).
\item \textbf{(The $\mathcal L$-invariant, three ways.)} $\mathcal L_{\mathrm{GS}}:=-2\,d\alpha/dk(2)$ from the level-one family alone $\equiv\mathcal L_{\mathrm{Tate}}(11a)\pmod{11^{8}}$
  ($\mathcal L_{\mathrm{Tate}}\equiv14882352232\pmod{11^{10}}$, committed before the C++), and $\mathcal L_{\mathrm{MTT}}$ from 11a's plus measure $\equiv\mathcal L_{\mathrm{Tate}}\pmod{11^{5}}$. The theorem
  $\mathcal L_p(E)=-2a_p'(2)$ and the MTT conjecture for 11a are Greenberg--Stevens' (\STD, recalled: their pp.~422--423 not read on the page; \texttt{docs/implementation/slot21/gs\_reading\_notes.md});
  the run reproduces them to the stated precision (\CMPD{} as a run).
\end{enumerate}
\emph{P1's two limits} (a remark for \S\ref{sec:lvlfunction}, found by P9's oracle arm; no statement of \S\ref{sec:lvlfunction} changes): P1's pipeline (i) refuses a cell where $W(\Z/\ell^{R})$ is not free
($H^{2}_c$ with $\ell$-torsion; at $(11,2)$ every weight $k\ge22$), whereas P8 reads through $W(\Z/\ell^{R+2})$; and (ii) at a class through $\m_{\mathrm{Eis}}$ takes the Eisenstein root for
every separator, so it cannot select another component of that class. Both are recorded per side by P9 (\texttt{P1\_COVERAGE}), never a failure. \emph{P2/P4} (v1.21.0-pre.3; through pre.2
``pending''): P2 gives the Eisenstein side of item~2 exactly at each \emph{single} weight, pole branches included (Proposition~\ref{prop:lveismeas}); the comparison of two $\Lambda$-adic objects
--- $L_a(X,T)$ against an Eisenstein measure interpolated in $X$ (P2's \texttt{EisDistribution} as P9's Eisenstein side, \texttt{P2\_spec.md}~\S6) --- is \emph{not} built, so item~2 is unchanged.
P4's wild characters (Remark~\ref{rmk:lvwild}) are at single weights as well; none is evaluated on $L_a(X,T)$.
\end{remark}

\paragraph{The Eisenstein measure exactly, and wild characters (v1.21.0-pre.3).} (New at v1.21.0-pre.3, hecke-extensions slot~21, the $\ell$-adic track P2 and P4; appended to
\S\ref{sec:lvslot21}, so that no statement number moves. \emph{No statement of \S\S\ref{sec:lvmeasure}--\ref{sec:lvlfunction} changes}; Claim~\ref{cl:lvmeascong}(4) gains its explanation as added
text. Specs with predictions committed before the gated runs: \texttt{docs/implementation/slot21/P2\_spec.md} (v0.1 \texttt{2c82a469}, v0.2 with the first runs; code \texttt{74c7e86c}) and
\texttt{P4\_spec.md} (v0.1 \texttt{8f9d4cf4}; code \texttt{f2c07b20}; v0.2 \texttt{b7b0e2c3}); Python oracles \texttt{p2/run\_p2.py} (output of record \texttt{p2/out\_run\_p2.txt}) and
\texttt{p4/cal\_*.py} (\texttt{p4/out\_cal.txt}), in-house, no CAS; neither spec is refereed. Tags as in the specs: their COMP is \CMPD, their STD-recalled is \STD{} recalled.)
Notation of \S\ref{sec:lvmeasure}; in addition $(\ell-1)\nmid k$ (so the two pole branches $a\equiv-1$, $a\equiv k-1$ differ and $d:=v(B_k/k)\ge0$), $E_r(x+\ell^{m}\Z_\ell):=\ell^{m(r-1)}B_r(x/\ell^{m})/r$
the Bernoulli distributions (Washington \S12.1, \STD{} recalled), $\bar B_1(t):=\{t\}-\tfrac12$, $\iota(x):=x^{-1}$, $x_r:=u^{r}-1$.

\begin{proposition}[The Eisenstein measure (slot-21 P2)]\label{prop:lveismeas}
Tags per item. Let $\ell\ge5$, $n=k-2$ even, $(\ell-1)\nmid k$.
\begin{enumerate}
\item \textbf{(The symbol side is a Dedekind--Rademacher sum.)} Let $\tilde P:=n!\,p_k^{-}$, Zagier's Laurent polynomial $p_k^{-}(X)=\sum_{-1\le m\le k-1,\ m\text{ odd}}\beta_m\beta_{k-2-m}X^{m}$,
  $\beta_m=B_{m+1}/(m+1)!$ (read on the page: Zagier, Invent.\ Math.\ 104 (1991) 449--465, Prop.~(i), eq.~(11), p.~453, $p_k^{-}\in\hat W_k^{-}$; read on the page for P2, \texttt{P2\_spec.md}~\S1.2, and again by the referee F4m;
  the phrase ``Eisenstein period function'' is the repository's, not Zagier's, which is all that \texttt{missing\_literature.md} LIT-05 leaves open; the same polynomial is Kohnen--Zagier \S1.3,
  p.~211; Claim~\ref{cl:lvmeascong}(1) carries the tag since v1.21.0-pre.4). Its interior is $P^{\mathrm{Eis}}$, its outer coefficients are
  $B_k/(k(k-1))$ at $u_{-1}$, $u_{n+1}$, and $\tilde P$ is annihilated by $1+S$ and $1+U+U^{2}$ as a Laurent polynomial in the convention of \S\ref{sec:lvmeasure}, the interior alone is not (\CMPD).
  Feed $\tilde P$ into the Manin-path formula of Definition~\ref{def:lvmtt} and drop the one singular monomial ($u_{-1}=x^{n+1}/y$ at the first Manin matrix, bottom row $(1,0)$: ``the boundary''). Then,
  \emph{as rational numbers}, for every $m\ge1$ and $a\in[0,\ell^{m})$,
  \[
    \tilde c(a,m)=-D(a,m)+\frac{B_k}{k}\cdot\frac{a}{\ell^{m}},\qquad D(a,m):=\sum_{y\bmod\ell^{m}}\bar B_1\Bigl(\frac{ay}{\ell^{m}}\Bigr)E_{k-1}(y+\ell^{m}\Z_\ell).
  \]
  \CMPD{} (exact, at every residue; Check~\ref{chk:lveismeas}(E1)). The term $(B_k/k)a/\ell^{m}$ is the period around $\infty$ (not $T$-invariant). That the Eisenstein modular symbol has
  Dedekind--Rademacher form is classical (Stevens, \emph{Arithmetic on modular curves}, Ch.~2; \STD{} recalled, not read; LIT-07 not held); no proof of the identity is claimed here.
\item \textbf{(The measure.)} $\mu^{\mathrm{Eis}}(a+\ell^{m}\Z_\ell):=-D(a,m)+\ell^{k-2}D(a,m-1)$ (the formula (M) with $\alpha=1$, the unit root of $X^{2}-(1+\ell^{k-1})X+\ell^{k-1}$) is a
  distribution over $\Q$ and odd (\CMPD, exact); on $\Z_\ell^{\times}$ it is a convolution of two Bernoulli distributions,
  \[
    \mu^{\mathrm{Eis}}|_{\Z_\ell^{\times}}=-E_1|_{\Z_\ell^{\times}}\star\iota_\ast E_{k-1}|_{\Z_\ell^{\times}},\qquad\text{at level }M:\ \mu^{\mathrm{Eis}}(a)=-\sum_{x\in(\Z/\ell^{M})^{\times}}E_1(x)\,E_{k-1}(xa^{-1})
  \]
  (\CMPD). Hence, for $\chi$ of conductor $\ell$ and every level $M$, $\sum_a\chi(a)\mu^{\mathrm{Eis}}(a+\ell^{M}\Z_\ell)=-B_{1,\chi}B_{k-1,\bar\chi}/(k-1)$ exactly, no Riemann sum (\PROVED{} from the
  convolution form and the distribution property of the generalised Bernoulli numbers, \STD; \CMPD{} at every odd $\chi=\omega^{i}$, (E4)).
\item \textbf{(The branch series.)} The branch map is a ring homomorphism, so $\mathrm{branch}_a(\mu^{\mathrm{Eis}})=-\mathrm{branch}_a(E_1)\cdot\mathrm{branch}_a(\iota_\ast E_{k-1})$, and on every odd branch
  \[
    \mathcal L^{\mathrm{Eis}}_a(T)=-G_{a+1}\bigl(u(1+T)-1\bigr)\cdot G_{k-1-a}\bigl(u^{k-1}(1+T)^{-1}-1\bigr)\qquad\text{in }\Lambda\otimes\Q,
  \]
  exactly, not modulo $\ell^{d}$ (\PROVED{} given Washington Thm.~12.2, p.~239, read on the page by the referee F4m, v1.21.0-pre.4: $\int\chi\omega^{-1}(a)\langle a\rangle^{s}dE_{1,c}$
  $=-(1-\chi(c)\langle c\rangle^{s+1})\,L_\ell(-s,\chi)$; through pre.3 ``\STD{} recalled, not re-read for P2'', with the leg mis-recalled as ``$x\,\zeta_\ell=E_1$, $dE_{k,c}=k\,x^{k-1}dE_{1,c}$'' (F4m, M1):
  in this text's normalisation $E_r=\ell^{m(r-1)}B_r/r$ the correct forms are $E_1=-x\,\zeta_\ell$ on $\Z_\ell^{\times}$, $\zeta_\ell$ the pseudo-measure with $\int\omega^{b}(x)\langle x\rangle^{s}d\zeta_\ell=G_b(u^{s}-1)$,
  and $dE_r=x^{r-1}dE_1$ (Washington's $E_k=i^{k-1}B_k(\{a/i\})$, p.~232, carries no $1/k$, whence his $k\,x^{k-1}$); $G_b$ of
  \Ucite{lem:u-twistvalues}, read as in Claim~\ref{cl:lvmeascong}(3)). Hence each factor carries the sign $-1$: $\mathrm{branch}_a(E_1)(T)=-G_{a+1}(u(1+T)-1)$ and
  $\mathrm{branch}_a(\iota_\ast E_{k-1})(T)=-G_{k-1-a}(u^{k-1}(1+T)^{-1}-1)$ (\CMPD{} by F4m at every critical value of the regular odd branches at $(14,7)$, $(30,37)$, $(66,37)$, $(10,691)$,
  $(282,37)$; \texttt{P2\_spec.md} through v0.2 and [U] 0.5.0-pre.5 stated them without the sign; the product is unchanged). \CMPD{} at every regular odd branch against both Kubota--Leopoldt routes of Check~\ref{chk:lvmeas}(M4)(a), which are
  independent of this construction; the even branches vanish.
\item \textbf{(The pole branches: an explicit polar term.)} On $a\equiv-1$ and $a\equiv k-1$ one factor is $G_0$:
  \begin{gather*}
    \mathcal L^{\mathrm{Eis}}_{-1}(T)=\frac{\mathrm{Res}}{u(1+T)-1}+(\Lambda),\qquad\mathcal L^{\mathrm{Eis}}_{k-1}(T)=\frac{\mathrm{Res}}{u^{k-1}(1+T)^{-1}-1}+(\Lambda),\\
    \mathrm{Res}=(1-\ell)\,\frac{\log_\ell u}{\ell}\,(1-\ell^{k-1})\,\frac{B_k}{k},
  \end{gather*}
  the same residue on both, $\mathrm{Res}=-\rho_\infty\,G_k(x_k)$ with $\rho_\infty:=(1-\ell)\log_\ell(u)/\ell$ ($\zeta_\ell$'s residue in the variable $u^{s}-1$, \STD{} recalled) and
  $G_k(x_k)=-(1-\ell^{k-1})B_k/k$ (\Ucite{lem:u-twistvalues}); so $v(\mathrm{Res})=d$. The shape is \PROVED{} for the branch $-1$ up to the constant (elementary, from $E_1$'s ball values);
  the constant was \CMPD{} (calibrated at $(12,5)$, level~$4$, and $(14,7)$, level~$3$; then predicted at the other cells) and is \PROVED{} given (3) and $\zeta_\ell$'s residue (F4m,
  v1.21.0-pre.4): $T\,G_0(T)\in\Lambda$ with $(TG_0)(0)=\rho_\infty$, so $G_0(T)=\rho_\infty/T+(\Lambda)$; with $X:=u(1+T)-1$ and $Y:=u^{k-1}(1+T)^{-1}-1$ one has $Y-x_k=-u^{k-1}(1+T)^{-1}X$,
  so $G_k(Y)\equiv G_k(x_k)\pmod{X\Lambda}$ and $-G_0(X)\,G_k(Y)=-\rho_\infty G_k(x_k)/X+(\Lambda)$; on the branch $k-1$ the same with the two factors exchanged (the pole at $Y=0$, where $X=x_k$). At level $M$ the polar part is
  $(1-\ell)(1-\ell^{k-1})(B_k/k)\ell^{-M}N_u$, $N_u=\sum_{i<\ell^{M-1}}u^{i}[\gamma^{i}]$ (branch $-1$; on the branch $k-1$, $\sum_iu^{-(k-1)i}[\gamma^{i}]$ with an extra $1/(k-1)$); \CMPD: the branch
  element has valuation exactly $d-M$ when $d<M$, and minus the polar part it is integral.
\item \textbf{(WP-MEAS's polynomial symbol.)} Let $\mu^{\mathrm{poly}}:=\mu^{(1)}_{P^{\mathrm{Eis}}}$ (Proposition~\ref{prop:lvmtt}(6)'s notation) with branch series $\mathcal L^{\mathrm{poly}}_a$. At every cell with
  $d\ge1$ of Check~\ref{chk:lveismeas}: $v(\mathcal L^{\mathrm{poly}}_a-\mathcal L^{\mathrm{Eis}}_a)=d$ exactly off the pole branches, and
  $v\bigl(\mathcal L^{\mathrm{poly}}_a-(\mathcal L^{\mathrm{Eis}}_a-\mathrm{polar}_a)\bigr)=d$ exactly on them (\CMPD, (E7)).
\item \textbf{(The cusp form's pole branches.)} Given Claim~\ref{cl:lvmeascong}(2) and its reduction (R20-1 Prop.~3.1, \PROVED{} given its inputs), (5) gives, on every odd branch, pole branches
  included,
  \begin{gather*}
    \mathcal L^{-}_{f,a}(T)\equiv s\,\bigl(\mathcal L^{\mathrm{Eis}}_a(T)-\mathrm{polar}_a(T)\bigr)\pmod{\ell^{d}},\\
    \mathcal L^{-}_{f,a}(u^{j}-1)\equiv-s\,\Bigl(F_1F_2+\frac{\mathrm{Res}}{u^{r}-1}\Bigr)\pmod{\ell^{d}},
  \end{gather*}
  $r=j+1$ (branch $-1$), $r=k-1-j$ (branch $k-1$), $F_1F_2$ the product of the two $G$-factors at $T=u^{j}-1$. \CMPD{} with the eigensymbol $r^{-}$ at $(30,37)$ (level $3$), $(282,37)$ (level $3$),
  $(226,103)$ (level $2$), $R=4$: $0/62$, $0/566$, $0/454$ critical points (E8); with $\mu^{\mathrm{poly}}$ in place of $f$ (the Python oracle) also at $(66,37)$, $(42,59)$, $(56,67)$, $(22,103)$ and
  $(10,691)$ ($134$, $86$, $114$, $46$, $22$ points, $0$ failures; at $(66,37)$, $k=68>\ell$, the two points with $r=37$ have a polar value of valuation $-1$ and the
  identity still holds modulo $37$: F4m). The polar value has valuation exactly $d-1-v(r)$ (\PROVED{} from $v(\mathrm{Res})=d$ and $v(x_r)=1+v(r)$), so the naive product
  $-s\,F_1F_2$ agrees modulo $\ell^{d-1-v(r)}$ and no further: Claim~\ref{cl:lvmeascong}(4)'s exceptions modulo $\ell^{d-1}$ are the points with $\ell\mid r$, the $14+4$ located points of
  $(282,37)$, $(226,103)$ (\CMPD: the located sets of (E8) equal the claim's; the control $\mathrm{Res}\cdot(1+\ell)$ fails exactly at them), and at a depth-one cell with $r<\ell$ throughout ---
  $(30,37)$, $(10,691)$ --- the naive product fails at every pole point by a unit (Claim~\ref{cl:lvmeas691}(3): predicted there by R20-1, now derived). At $(10,691)$ the congruence is moreover \emph{checked on $\Delta$'s own eigensymbol at level $2$} ($R=2$): $\mathcal L^{-}_{\Delta}\equiv s(\mathcal L^{\mathrm{Eis}}-\mathrm{polar})\pmod{691}$
  holds at all $3795$ critical points of the $345$ odd branches ($3773$ regular, $22$ pole; $0$ failing, located set $\emptyset$; negative controls: no polar term $22/22$ failing, $-\mathrm{Res}$ $22/22$,
  $\mathrm{Res}\cdot(1+\ell)$ $0/22$ --- invisible at critical points), and as elements of $\Z/691[C_{691}]$ on every odd branch, pole branches included: $g^{\Delta}_a\equiv s(g^{\mathrm{Eis}}_a-\mathrm{polar}_a)$ coefficientwise,
  $0/343$ regular and $0/2$ pole branches ($s=286647$ as in S2; at level $2$ $\min v(g^{\mathrm{Eis}})=-1$ on both pole branches, $11$ and $689$, the polar term making the difference integral; controls: no polar term $2/2$, $-\mathrm{Res}$ $2/2$,
  $\mathrm{Res}\cdot(1+\ell)$ $2/2$ --- here the control fails, though it passes at every critical point --- and $s+1$ $345/345$; the C++ fingerprint of $g^{\Delta}$ on the pole branches, $11$: $(443652,194679)$, $689$: $(443652,11103)$, equals the oracle's)
  (\CMPD, \texttt{HeckeLValueEisensteinMeasureValidation} arms E8L2 and E9L2; independent Python oracle \texttt{docs/implementation/slot21/p2/p2x\_oracle.py}, output \texttt{out\_p2x\_oracle.txt}; predictions committed first, \texttt{P2\_spec.md} v0.4 \S10).
  At $d=1$ the critical values modulo $\ell$ are level-independent ($u\equiv1$); the level-$2$ content is the group-ring identity.
\end{enumerate}
\emph{Not claimed} (\texttt{P2\_spec.md}~\S1.5): a proof of (1); the plus side ($p_0$, the trivial zero; Remark~\ref{rmk:lvmeasplus}(iii) stays an observation); $(\ell-1)\mid k$ (one pole branch, a double
pole). Characters of $\ell$-power order: Remark~\ref{rmk:lvwild}(6), (7).
\end{proposition}

\begin{checkitem}[The gate of P2: \texttt{HeckeLValueEisensteinMeasureValidation}]\label{chk:lveismeas}
\SPEC{} (the checks); the expected values are P2's predictions and its Python oracle's, committed first, never this code's output. \emph{Registration}: \texttt{Registry/Slot21/P2\_eismeas.cpp},
\texttt{cmake/slot21/P2\_eismeas.cmake}; sources \texttt{HeckeExt/LValue/EisensteinMeasure.\{h,cpp\}}, \texttt{HeckeLValue\allowbreak Eisenstein\allowbreak Measure\allowbreak Validation.\{h,cpp\}}; WP-MEAS's public API only, no
existing $L$-value file edited. Tier \texttt{standard}, read-only. Exact rational identities modulo $q_1=2^{61}-1$ and $q_2=2^{62}-57$ (one prime at $\ell\ge100$); $\ell$-adic statements in a scaled
representation $\ell^{E}x\bmod\ell^{R+E}$, $E=2M+1+v(k-1)$, $\ell^{R+E}<2^{62}$. \emph{Cells}: Eisenstein $(12,5)$ level $4$, $(14,7)$ level $3$, $(16,11)$, $(20,13)$, $(10,37)$ ($d=0$),
$(30,37)$, $(66,37)$, $(42,59)$, $(56,67)$, $(22,103)$ ($d=1$), $(282,37)$, $(226,103)$ ($d=2$) at level $2$, $(10,691)$ at level $1$ (E1--E7, E9) and, for $\Delta$'s own eigensymbol, at level $2$ (E8L2, E9L2); cusp forms (WP-MEAS's) $(30,37)$, $(282,37)$ at level $3$,
$(226,103)$ at level $2$, $R=4$. Arms: (E0) Zagier's relations at $6$ random points modulo $q_1$, the interior alone fails; (E1) item (1) at every residue $a<\ell^{m}$, $m\le M$; (E2) distribution and parity,
negatives: the cusp term kept, $\mu^{\mathrm{poly}}$ over $\Q$; (E3) the convolution form; (E4) the twisted masses; (E5) item (3) by both Kubota--Leopoldt routes ($E_{1,c}$ only at $\ell>103$), even
branches $0$; (E6) item (4); (E7) item (5); (E8) item (6) and the located sets; (E9) negatives --- no polar term, $-\mathrm{Res}$ (fail at every point), $\mathrm{Res}\cdot(1+\ell)$ (fails exactly where
$\ell\mid r$), the wrong twist $G_{a+3}$ (at level $\ge2$ only: at level~$1$ a comparison is one residue and coincides with probability $1/\ell$; first run $3/342$ at $(10,691)$). Lines
E8L2 ($7$ checks: the critical-point identity at $3795$ points, the group-ring identity on the $345$ odd branches, the C++ fingerprints on the pole branches) and E9L2 ($7$ checks: the four negative controls at critical points and in the group ring). Lines
\texttt{EISMEAS}, \texttt{EISMEAS\_POLE}, \texttt{P2\_ARMS checks=<c> failing=<f>}; \texttt{EFFECTIVE \dots{} arm=DEFAULT cells=13}. \emph{First run} (\CMPD): $225$ checks, $0$ failing,
$5.3$\,s, $20$\,MB (E1 $0/2812$ at $37$, $0/10712$ at $103$, $0/691$ at $691$; E6 $d-M$ exactly at $d<M$, $0$ at $d=M=2$; E7 $d$, $d$ at the eight cells with $d\ge1$; E8 $0/62$, $0/566$, $0/454$; no polar term
and $-\mathrm{Res}$ fail at every point, $\mathrm{Res}\cdot(1+\ell)$ at $0$, $14$, $4$). E5 is vacuous at $\ell=5$ (both odd branches are pole branches). Read-only env arm
\texttt{HECKE\_LVALUE\_EISMEAS\_HEAVY=1} (\texttt{arm=HEAVY cells=15}): adds $(30,37)$ at level $3$ (E1 $0/104118$, E3 $0/49284$, E6 $(-2,0)$) and $(20,131)$ at level $2$; $259$ checks, $0$ failing,
$11.9$\,s, $20$\,MB. \emph{With E8L2 / E9L2} (P2x, \CMPD): default arm $239$ checks, $0$ failing, $6.2$\,s; heavy arm $273$ checks, $12.6$\,s (E8L2 $0.8$\,s on $16$ threads, $3.9$\,s single; the Python oracle $78$\,s, $0.22$\,GB). \texttt{GOLDEN\_STRICT=1}: A $11827/0$, B $6231/0$. No golden line, record key or digest.
\end{checkitem}

\begin{remark}[Wild characters on the $L$-side (slot-21 P4)]\label{rmk:lvwild}
Let $\ell\ge5$, $m\ge2$, $F=\ell^{m}$, $\zeta$ of order $\ell^{m-1}$, $\mathcal O_m=\Z_\ell[\zeta]$, $e=\varphi(\ell^{m-1})$, $\pi=1-\zeta$; for $a\in(\Z/F)^{\times}$, $\langle a\rangle=u^{L(a)}$,
$\psi_t(a):=\zeta^{tL(a)}$ ($\ell\nmid t$; conductor $\ell^{m}$), $\chi_{i,t}:=\omega^{i}\psi_t$ --- every primitive character of conductor $\ell^{m}$, once $\zeta$ is fixed. Tags per item; everything
is at single weights, over $\mathcal O_m$ (not over $\mathcal O\otimes\mathcal O_m$, Check~\ref{chk:lvcoeffring}).
\begin{enumerate}
\item \textbf{(The analytic constant.)} With Remark~\ref{rmk:lvshapes}~(B)'s normalisation the quotient $\Omega_{\chi,j}:=L(f,\chi,j+1)\,j!\binom nj M^{j+1}/(\chi(-1)(-2\pi i)^{j+1}\tau(\chi)\lambda_{\chi,j})$
  is, for every primitive $\chi$ of conductor $M=25$, $49$, $121$ at $k=12$ ($\Delta$) and $M=25$, $49$ at $k=16$, and every $j$ with $\lambda_{\chi,j}\ne0$, one constant per sign, equal to the
  conductor-$1$ $\Omega_f^{\pm}$ to relative $10^{-59}$--$10^{-67}$ ($16$, $36$, $100$, $16$, $36$ characters; $176$, $396$, $1100$, $240$, $540$ values; env arm: $M=125$, $169$ at $k=12$) --- \CMPD{}
  ($L$ by the two-sided incomplete-Gamma series, $256$-bit \texttt{nf::MpReal}; $P^{\pm}$ and $\lambda_{\chi,j}$ exactly over $\Q$); $\tau(\bar\chi)$ in place of $\tau(\chi)$ fails (relative deviation
  $\approx2$). No ordinarity is needed: the run is at primes where $f$ is not ordinary as well. The root number of $f_\chi\in S_k(M^{2},\chi^{2})$, in
  $\Lambda(f_\chi,s)=w\,i^{k}N^{k/2-s}\Lambda(f_{\bar\chi},k-s)$, $N=M^{2}$, is $w=\tau(\chi)^{2}/M$ (\PROVED, elementary, F4m, v1.21.0-pre.4: $f_\chi=\tau(\bar\chi)^{-1}\sum_{v\in(\Z/M)^{\times}}\bar\chi(v)\,f|_k\bigl(\begin{smallmatrix}1&v/M\\0&1\end{smallmatrix}\bigr)$
  and $\bigl(\begin{smallmatrix}1&v/M\\0&1\end{smallmatrix}\bigr)W_{M^{2}}=\gamma_v\bigl(\begin{smallmatrix}M&u\\0&M\end{smallmatrix}\bigr)$ with $\gamma_v=\bigl(\begin{smallmatrix}v&-x\\M&-u\end{smallmatrix}\bigr)\in\mathrm{SL}_2(\Z)$,
  $u\equiv-v^{-1}\pmod M$, $W_{M^{2}}=\bigl(\begin{smallmatrix}0&-1\\M^{2}&0\end{smallmatrix}\bigr)$; as $f|_k\gamma_v=f$ and $v\mapsto u$ permutes $(\Z/M)^{\times}$ with $\bar\chi(v)=\chi(-1)\chi(u)$,
  $f_\chi|_kW_{M^{2}}=\chi(-1)\,\tau(\chi)\tau(\bar\chi)^{-1}f_{\bar\chi}=(\tau(\chi)^{2}/M)\,f_{\bar\chi}$ by $\tau(\bar\chi)=\chi(-1)\overline{\tau(\chi)}$ and $|\tau(\chi)|^{2}=M$, and the Mellin
  transform turns $f_\chi|_kW_{M^{2}}=w\,f_{\bar\chi}$ into the displayed functional equation; \CMPD{} to $10^{-59}$; P4's v0.1 predicted $\chi(-1)\tau(\chi)^{2}/M$ from a sign-blind
  calibration, and the C++ gate refuted it --- \texttt{P4\_spec.md}~\S8, deviation~1; the constant of Proposition~\ref{prop:lvmtt} does not use $w$).
\item \textbf{(Proposition~\ref{prop:lvmtt}(4) at wild $\chi$.)} $(-1)^{j}\binom nj\int\bar\chi(x)x^{j}d\mu_f=\alpha^{-m}\lambda_{\chi,j}$ in $\mathcal O_m/\ell^{\min(M,R)}$ (level-$M$ Riemann sum) --- a
  check of the code path of a \PROVED{} statement: \CMPD{} at $(10,11)$, $(10,13)$, $(10,37)$ and $(30,37)$ ($R=4$, $M=3$, $m=2$), both signs, every $i$, $t\in\{1,2\}$, every $j$ ($0/2232$ per sign at
  $(30,37)$; env arm: $(10,11)$ at $m=3$, $e=110$); $\alpha^{-(m-1)}$ in place of $\alpha^{-m}$ matches nowhere.
\item \textbf{(A functional-equation symmetry.)} $\lambda_{\chi,j}(P)=-\varepsilon\,\lambda_{\bar\chi,n-j}(P)$ modulo $\ell^{R}$, $\varepsilon$ the geometric parity of $P$ (\CMPD{} at the cells of (2) and, in the
  calibration, at $(14,7)$, $(16,7)$, $(20,11)$). \PROVED{} (F4m, v1.21.0-pre.4; through pre.3 ``looks provable from $\rho_n$ of $\bigl(\begin{smallmatrix}0&-1\\F&0\end{smallmatrix}\bigr)$; no proof is given''): for
  $u\equiv-a^{-1}\pmod F$ and $\gamma=\bigl(\begin{smallmatrix}a&-x\\F&-u\end{smallmatrix}\bigr)\in\mathrm{SL}_2(\Z)$ one has $\{a/F,\infty\}=\gamma\{\infty,u/F\}$, so
  $\Phi(\{a/F,\infty\})=-\rho_n(\gamma)\Phi(\{u/F,\infty\})$, and $\gamma_{a,m}\gamma=S\,\gamma_{u,m}$ with $S=\bigl(\begin{smallmatrix}0&-1\\1&0\end{smallmatrix}\bigr)$, while
  $\langle\rho_n(S)v,u^{\vee}_{n-j}\rangle=(-1)^{j}\langle v,u^{\vee}_{j}\rangle$ ($u_t\mapsto(-1)^{t}u_{n-t}$); summing against $\bar\chi(a)=\chi(-1)\chi(u)$ gives
  $\lambda_{\chi,j}(P)=-\chi(-1)(-1)^{j}\lambda_{\bar\chi,n-j}(P)$ for \emph{every} $P\in W$, which is $-\varepsilon\lambda_{\bar\chi,n-j}(P)$ wherever either side is non-zero
  ($\chi(-1)(-1)^{j}=\varepsilon$ there); checked by the referee on a generic $P\in W$: $0/176$ at $(10,5)$, $0/1100$ at $(10,11)$, $m=2$. Consequence (\CMPD): the control
  ``$\chi$ for $\bar\chi$'' coincides with the true value exactly at $j=n/2$, for the $\chi$ with $\chi(-1)(-1)^{n/2}=\varepsilon$, and nowhere else.
\item \textbf{(Wild values read $\lambda$.)} $\int\omega^{b-j}\psi_tx^{j}d\mu=\mathcal L_b(\zeta^{t}u^{j}-1)=\sum_ig_i\zeta^{ti}u^{ij}$ (Proposition~\ref{prop:lvmtt}(5) at the character $[u]\mapsto\zeta^{t}u^{j}$ of
  $\Ga_u$; \CMPD{} modulo $\ell^{\min(M,R)}$), and $v_\pi(\mathcal L_b(\zeta^{t}-1))=e\mu_b+\lambda_b$ whenever $\lambda_b<e$ (\PROVED: Weierstrass preparation and $v_\pi(\zeta^{t}u^{j}-1)=1$). \CMPD{} at
  every readable branch: $\lambda=1$ at branch $0$ of $r^{+}$ at $(10,11)$ and at the branches $21$, $25$ of $r^{-}$ at $(10,37)$; at $(30,37)$ $\mu=1$ on every even branch of $r^{+}$
  ($v_\pi=36$ or $37$) --- Remark~\ref{rmk:lvmeasplus}(ii)'s $\mu^{+}=d$ seen at wild points.
\item \textbf{(Kubota--Leopoldt at a wild point.)} With $G_b$ as WP-MEAS's group-ring element (index $i\leftrightarrow[u^{i}]$), $G_b(\zeta^{t}u^{r}-1)=-B_{r,\omega^{b-r}\psi_t}/r$, $\psi_t(u)=\zeta^{t}$ (not
  $\zeta^{-t}$), no Euler factor (\PROVED{} from \Ucite{lem:u-twistvalues}'s reading of Washington Thm.~7.10 and Thm.~5.11, F4m, v1.21.0-pre.4: $G_b(X)=f(u(1+X)^{-1}-1,\omega^{b})$, and at
  $1+X=\zeta^{t}u^{r}$ the argument is $\zeta^{-t}u^{1-r}-1=\zeta_\psi(1+q_0)^{s}-1$ with $s=1-r$ and $\zeta_\psi=\psi_t(u)^{-1}$, so $G_b(\zeta^{t}u^{r}-1)=L_\ell(1-r,\omega^{b}\psi_t)=-B_{r,\omega^{b-r}\psi_t}/r$,
  the Euler factor $1-\omega^{b-r}\psi_t(\ell)\ell^{r-1}$ being $1$ as $\ell$ divides the conductor; \CMPD, by both routes, against the definition of $B_{r,\chi}$, and independently by the referee at
  $(\ell,m)=(5,2)$, $(7,2)$ with $\zeta^{-t}$ failing at every value; the wild extension of \Ucite{lem:u-twistvalues}, recorded there as a remark). On the pole class $b=0$,
  $v_\pi(B_{r,\omega^{-r}\psi_t}/r)=-1$ exactly: at a wild point the pole of $\zeta_\ell$ is a simple $\pi$-adic denominator (\PROVED{} from $G_0(T)=\rho_\infty/T+(\Lambda)$,
  Proposition~\ref{prop:lveismeas}(4), $v(\rho_\infty)=0$ and $v_\pi(\zeta^{t}u^{r}-1)=1$; \CMPD).
\item \textbf{(The Eisenstein measure at wild $\chi$.)} Proposition~\ref{prop:lveismeas}(2)'s identity $\sum_a\chi(a)\mu^{\mathrm{Eis}}(a+\ell^{M}\Z_\ell)=-B_{1,\chi}B_{k-1,\bar\chi}/(k-1)$ holds at wild odd $\chi$
  (the argument does not use the conductor; \CMPD{} at $(12,5)$, $m=2,3$, $(14,7)$, $(16,11)$, $(30,37)$), the even masses vanish, and the polar element of Proposition~\ref{prop:lveismeas}(4) evaluated
  at $\zeta^{t}$ has $v_\pi=ed-1$ exactly, the value minus it being integral (\CMPD; $35$ at $(30,37)$).
\item \textbf{(The pole branches at wild points.)} $\mathcal L^{-}_{f,a}(\zeta^{t}u^{j}-1)\equiv-s\,\bigl(F_1F_2+\mathrm{Res}/(\zeta^{\pm t}u^{r}-1)\bigr)\pmod{\ell^{d}\mathcal O_m}$,
  $F_1=B_{j+1,\omega^{a-j}\psi_t}/(j+1)$, $F_2=B_{n-j+1,\omega^{j-a}\psi_t^{-1}}/(n-j+1)$, $+t$ and $r=j+1$ on the branch $-1$, $-t$ and $r=k-1-j$ on the branch $k-1$ (\CMPD: $(30,37)$ $0/62$ at $t=1$
  and, env arm, $t=2$; $(282,37)$ $0/566$). As $v_\pi(\zeta^{\pm t}u^{r}-1)=1$ for every $r$, the polar value has $v_\pi=ed-1$, i.e.\ $\ell$-adic valuation $d-1/e$, \emph{independent of $r$}: it is visible modulo $\ell^{d}$ at every critical
  point (dropping it, or $-\mathrm{Res}$, fails at $62/62$, $566/566$; as at the tame points) and invisible modulo $\ell^{d-1}$ at every point, so Claim~\ref{cl:lvmeascong}(4)'s located set
  $\{\ell\mid r\}$ has no analogue at wild points (\texttt{P4\_spec.md}~\S0.9); and $\mathrm{Res}\cdot(1+\ell)$ changes the value only by
  $\pi^{e(d+1)-1}$ times a unit, so that control has no power at wild points ($0$ failures, recorded, not a gate on failure). The pole-branch identity at wild points (W6) is \CMPD{} also at
  $(226,103)$ ($d=2$, $\ell=103$, $M=2$, $m=2$; v1.21.0-pre.7, slot-21 P4x): $0/454$, the polar term visible at every point ($v_\pi=ed-1=203$; dropping it, or $-\mathrm{Res}$, fails at $454/454$),
  $\mathrm{Res}\cdot(1+\ell)$ invisible ($0$); the scaled comparison needs $\mathcal O/103^{10}$ ($>2^{62}$), computed in $128$-bit arithmetic (\texttt{OArith128} through \texttt{ModCtx<u128>}, $P=12$) and
  cross-checked against the $64$-bit route at $(30,37)$, $(282,37)$ (every input and every point's polar / product element equal, the verdict counts identical); independent exact-integer Python
  oracle \texttt{docs/implementation/slot21/p4/p4x\_pole\_oracle.py} ($P=13$: $0/454$, $454$, $454$, $0$, its $s=83864971$ equal to the C++ line's); predictions committed first
  (\texttt{P4\_spec.md} v0.4~\S10, \texttt{487057e6}; code \texttt{ab5a1ee1}). Through v1.21.0-pre.6 this item read ``Not covered: $(226,103)$ (the scaled representation exceeds $2^{62}$)''.
  Not covered: wild characters of non-trivial tame level, $\ell=2,3$, $(\ell-1)\mid k$.
\end{enumerate}
\end{remark}

\begin{checkitem}[The gate of P4: \texttt{HeckeLValueWildCharacterValidation}]\label{chk:lvwild}
\SPEC{} (the checks); expected values P4's predictions and calibration, committed first. \emph{Registration}: \texttt{Registry/\allowbreak Slot21/\allowbreak P4\_wild.cpp}, \texttt{cmake/\allowbreak slot21/\allowbreak P4\_wild.cmake}; sources
\texttt{HeckeExt/\allowbreak LValue/\allowbreak WildCharacter.\{h,cpp\}}, \texttt{HeckeLValue\allowbreak Wild\allowbreak Character\allowbreak Validation.\{h,cpp\}}; the public APIs of WP-MEAS, P2 and \texttt{NumberField} (\texttt{CycloModRing}, \texttt{MpReal})
only. Tier \texttt{standard}, read-only. Arms: (W0) the characters at $(\ell,m)=(5,2)$, $(5,3)$, $(7,2)$, $(11,2)$, $(37,2)$, P4's $\mathcal O$-arithmetic against \texttt{nf::CycloModRing}, negative $t=\ell$ not primitive;
(W1) Remark~\ref{rmk:lvwild}(5), $r\le10$, $t\in\{1,2\}$, both routes, at $(\ell,m,M,R)=(5,2,2,6)$, $(5,3,3,6)$, $(7,2,3,5)$, $(11,2,2,5)$, $(37,2,2,4)$, negative $\zeta^{-t}$; (W2) items (2), (3);
(W3) item (4); (W4) item (6); (W5) item (1), negative controls $\tau(\bar\chi)$ and conductor $\ell M$; (W6) item (7); (W7) the non-ordinary refusal of \texttt{UnitRoot} at $(10,5)$, $(10,7)$. Lines
\texttt{WILD\_W0}, \texttt{WILD\_W1}, \texttt{WILD\_CUSP}, \texttt{WILD\_EIS}, \texttt{WILD\_COMPLEX}, \texttt{WILD\_POLE}, \texttt{WILD\_W7}, \texttt{P4\_ARMS checks=<c> failing=<f>};
\texttt{EFFECTIVE \dots{} arm=DEFAULT cells=27}. \emph{First run} (\CMPD): $201$ checks, $0$ failing, $4.0$\,s, $24$\,MB. Read-only env arm \texttt{HECKE\_LVALUE\_WILD\_HEAVY=1}
(\texttt{arm=HEAVY cells=31}: W5 at $125$, $169$; W2/W3 $(10,11)$ at $m=3$; W6 $(30,37)$ at $t=2$): $238$ checks, $0$ failing, $7.4$\,s, $25$\,MB. \texttt{GOLDEN\_STRICT=1}: A $11827/0$, B $6231/0$.
No golden line, record key or digest. \emph{One prediction was wrong} and is recorded as such: the root number (item (1)); the calibration had compared $\min(|x-\mathrm{pred}|,|x+\mathrm{pred}|)$,
which cannot see the sign it was meant to calibrate (\texttt{P4\_spec.md}~\S8, the methodological lesson). \emph{With P4x} (v1.21.0-pre.7, \CMPD): W6 takes a $128$-bit route at a cell exactly
when its demand exceeds the $64$-bit representation ($(226,103)$, line \texttt{WILD\_POLE cell=(226,103) route=u128 P=12 Pneed=10 \dots{} points=454 fail=0}), and a new route-agreement check
(\texttt{WILD\_ROUTE}) forces it at $(30,37)$, $(282,37)$ against the $64$-bit route ($0$ mismatches); default arm \texttt{arm=DEFAULT cells=28}, $222$ checks, $0$ failing, $5.3$\,s, $0.09$\,GB; heavy arm
\texttt{arm=HEAVY cells=32}, $259$ checks, $0$ failing, $8.9$\,s. No new env hook or registered name.
\end{checkitem}

\begin{remark}[Iwasawa invariants of the measure branches, and the tame functional equation (slot-21 P5)]\label{rmk:lviwasawa}
(New at v1.21.0-pre.7. Source: \texttt{docs/implementation/slot21/P5\_spec.md} v0.2 --- predictions v0.1 committed before any run, \texttt{0dd6c269}; code and first runs \texttt{58c24210}; gate
Check~\ref{chk:lviwasawa}. Tags per item as in the spec; every run \CMPD; nothing upgraded. Refereed at v1.21.0-pre.8 by F4 (\texttt{rc/F4\_referee\_p5\_report.md}): the pole-branch formula of (7)
added, five wording repairs marked ``F4''. Item (9) extended at v1.21.0-pre.9 by P5$'$ (the McCallum--Sharifi identification; not refereed), and at v1.22.0-pre.2 by P5$''$ (the modular-symbol proof at the computed cells; Fable-refereed).) Notation of Proposition~\ref{prop:lvmtt}: $\ell\ge11$, $n$ even, $k=n+2$, $f$ ordinary of level one with unit
root $\alpha$, $r^{\pm}$ its eigensymbols, $\varepsilon=\pm1$ their geometric parity (the eigenvalue of $u_t\mapsto(-1)^{t}u_t$: $r^{-}$ has odd coordinates, so $\varepsilon(r^{-})=-1$; through pre.7 ``$\mp1$'', F4); $\mathcal L_b(T)\in\Z_\ell[\![T]\!]$ the branch $b\in\Z/(\ell-1)$ of $\mu_{r^{\pm}}$ (non-zero only for $(-1)^{b}=\varepsilon$), $(\mu_b,\lambda_b)$ its Iwasawa
invariants; $G_b$ as in \Ucite{lem:u-twistvalues}; $\sigma(T):=u^{n}(1+T)^{-1}-1$.
\begin{enumerate}
\item \textbf{(The tame functional equation.)} For every $P\in W$ (not only eigenvectors), every Dirichlet character $\chi$ of conductor $F\in\{1,\ell\}$ and every $0\le j\le n$,
  \[\lambda_{\chi,j}(P)=-\chi(-1)(-1)^{j}\,\lambda_{\bar\chi,n-j}(P)\quad\text{exactly}.\]
  \PROVED: Remark~\ref{rmk:lvwild}(3)'s proof uses nothing about $F$ but $\gcd(a,F)=1$, and holds verbatim at $F=\ell$; at $F=1$ ($a=0$, $\gamma_{0,0}=1$) the identity is $P+\rho_n(S)P=0$ read on the
  coordinate $u_{n-j}$, i.e.\ $P_{n-j}=-(-1)^{j}P_j$. At an eigensymbol of parity $\varepsilon$ the values vanish unless $\chi(-1)(-1)^{j}=\varepsilon$, so the non-zero ones carry the sign $-\varepsilon$.
  \CMPD{} at every cell of Check~\ref{chk:lviwasawa} on $r^{\pm}$ and on generic (non-eigen) $P\in W$, every Teichm\"uller power $\omega^{i}$ and conductor~$1$ ($31413$ values at $(282,37)$, $70143$ at
  $(226,103)$; $0$ failures); the flipped sign matches nowhere; ``$\chi$ for $\bar\chi$'' matches exactly at $j=n/2$ on $r^{-}$ (as in Remark~\ref{rmk:lvwild}(3)).
\item \textbf{(The branch form.)} $\mathcal L_b(T)=-(-1)^{b}\,\mathcal L_{n-b}\bigl(\sigma(T)\bigr)$. \PROVED{} given Proposition~\ref{prop:lvmtt}(4) at every conductor $\ell^{m}$ and (1) at every
  $F=\ell^{m}$ (Remark~\ref{rmk:lvwild}(3)): $E_{\chi,j}$ and $\binom nj$ are symmetric under $j\leftrightarrow n-j$, so the difference of the two sides vanishes at the points $\zeta u^{j}-1$, $\zeta$ of $\ell$-power
  order, $\zeta^{-1}u^{n-j}-1=\sigma(\zeta u^{j}-1)$, infinitely many in the open disc; Weierstrass preparation. At level $M$ it holds in $\Z/\ell^{R}[\Ga_u/\Ga_u^{\ell^{M-1}}]$ modulo
  $\ell^{\min(M,R)}$ (\CMPD, every branch of the measure's parity, both signs; the flipped sign fails at every branch with $g_b\not\equiv0$).
\item \textbf{(Consequences for $(\mu,\lambda)$.)} \PROVED{} from (2), $\sigma$ being an automorphism of $\Z_\ell[\![T]\!]$ ($\sigma(T)\in(\ell,T)$ and $\sigma\circ\sigma=\mathrm{id}$; F4) which preserves $\mu$
  (it is an automorphism of $\F_\ell[\![T]\!]$) and $\lambda$ (modulo $\ell$, $\bar\sigma(T)=(1+T)^{-1}-1=-T+T^{2}-\dots$ is $T$ times a unit, so $\mathrm{ord}_T$ is preserved):
  \emph{(FE-inv)} $(\mu_b,\lambda_b)=(\mu_{n-b},\lambda_{n-b})$ for every branch; \emph{(FE-parity)} at a self-dual branch ($2b\equiv n\bmod\ell-1$; never a pole branch of the Eisenstein measure, since $a\equiv-1$ or $k-1$ with $2a\equiv k-2$ forces $(\ell-1)\mid k$; F4) $\sigma$ is an involution of the open disc (over $\C_\ell$) with the single fixed point
  $T_0=u^{n/2}-1$ ($-u^{n/2}-1$ is a unit), $\sigma'(T_0)=-1$, the zeros of $\ell^{-\mu}\mathcal L_b$ (counted by Weierstrass preparation) off $T_0$ pair up, and the order $r$ at $T_0$ satisfies $(-1)^{r}=-(-1)^{b}$: \textbf{$\lambda_b$ is odd at an even self-dual branch
  (plus measure) and even at an odd one (minus measure)} --- an $\ell$-adic \emph{forced central zero} at $j=n/2$ on the even self-dual branches. \CMPD: (FE-inv) at every branch of every cell; at every
  self-dual branch of every cell $\lambda=1$ at the even ones ($(10,\ell)$ at $\ell=11,19,23,31,43,47,59,67,103,131$; $(16,37)$ at $8$ and $26$; $(42,59)$, $(56,67)$, $(22,103)$, $(20,131)$,
  $(226,103)$, $(10,691)$ at $350$) and $\lambda=0$ at the odd ones.
\item \textbf{(The near-trivial zero.)} $\mathcal L_0(0)=E_{\mathbf 1,0}\,\lambda_{\mathbf 1,0}$ with $E_{\mathbf 1,0}=(1-1/\alpha)(1-\ell^{n}/\alpha)$; so if $\alpha\equiv1\pmod\ell$ then on the plus measure
  $\mu_0\ge1$ or $\lambda_0\ge1$, and likewise at $b\equiv n$ by (FE-inv) (\PROVED). $\alpha\equiv\tau(\ell)\equiv1\pmod\ell$ for $\Delta$ at $\ell=11$, $23$ (the values \STD, recalled; the code reads $a_\ell$ from the
  symbol), and at every Eisenstein-congruent cell, where $\mu^{+}=d\ge1$ anyway. \CMPD: $\lambda=1$ at $(10,11)$ branch $0$ (also its even self-dual branch) and at $(10,23)$ branches $0$, $10$.
\item \textbf{(A ball-free reading of $(\mu,\lambda)$: the values oracle.)} For one branch the critical values $V_j=\mathcal L_b(u^{j}-1)$ come from the eigensymbol by Proposition~\ref{prop:lvmtt}(4) alone (no
  ball table, no group ring); their Newton divided differences $\delta_t$ at nodes $j_0<\dots<j_L$ with $\ell\nmid\binom n{j_s}$ satisfy $v(\delta_t-c_t)\ge\mu_b+1$, $c_t$ the $T^{t}$-coefficient (\PROVED:
  $v(T_j-T_{j'})=1+v(j-j')$ and the complete homogeneous polynomials $h_m$ have $v\ge m$), so $v(\delta_t)\ge\mu_b+1$ for $t<\lambda_b$ and $v(\delta_{\lambda_b})=\mu_b$. \CMPD: consistent with the group-ring
  invariants at every readable branch, integral everywhere, and the only reading where the group-ring invariants are not readable ($\mu\ge M-1$: the depth-$2$ cells' plus side, $\mu^{+}=2$ at $M=2$).
\item \textbf{(The Kubota--Leopoldt invariants.)} $\lambda(G_b)=0$ iff $\ell\nmid B_b/b$ (\PROVED{} \emph{given} $\mu(G_b)=0$: $c_0\equiv G_b(u^{b}-1)=-(1-\ell^{b-1})B_b/b\pmod\ell$, so a unit $c_0$ gives
  $(\mu,\lambda)=(0,0)$ unconditionally, while $\ell\mid c_0$ gives $\lambda\ge1$ only once $\mu=0$ is known; F4); $\mu(G_b)=0$ for every even $b$
  (Ferrero--Washington, \STD, recalled, \VERIFYloc); $\lambda(G_b)=1$ at every irregular pair $(\ell,b)$ with $\ell<2^{31}$ (read on the page: \cite{hho}, abstract item~(3), p.~1, ``the Iwasawa invariants satisfy
  $\lambda_p=\nu_p=i_p$'', text layer read for P5 and again by F4, who also read \S7 of the manuscript, where $\lambda_p$, $\nu_p$ are the Iwasawa invariants of the $p$-class groups of $\Q(\zeta_{p^{n}})$ and the
  test is four incongruences per irregular index; the step $\lambda_p=\sum_{b\text{ even}}\lambda(G_b)$ is \STD, recalled (Iwasawa's theorem for the minus part through the analytic class number formula, the plus
  part trivial by Kummer--Vandiver, verified in \cite{hho} item~(2); locator unverified), and with $\lambda(G_b)\ge1$ at each irregular index --- Ferrero--Washington again --- it forces $\lambda(G_b)=1$). \CMPD{} by two routes --- the
  Kubota--Leopoldt group ring at level~$2$, and exact Bernoulli numbers alone (divided differences of $G_b(u^{r}-1)=-(1-\ell^{r-1})B_r/r$ at $r=b+s(\ell-1)$, $s=0,\dots,3$), equal at every class ($344/344$ at
  $691$): at the $16$ primes $11,\dots,47,59,67,103,131,691$, $\mu(G_b)=0$ at every even $b$, and $\lambda(G_b)=1$ exactly at $(37,32)$, $(59,44)$, $(67,58)$, $(103,24)$, $(131,22)$, $(691,12)$,
  $(691,200)$, $0$ elsewhere. \emph{No $\lambda(G_b)\ge2$ occurs.}
\item \textbf{(Greenberg--Vatsal at level one.)} At a cell with $d=v(B_k/k)\ge1$, on every odd \emph{regular} branch $a$ ($a+1\not\equiv0$, $k-1-a\not\equiv0\bmod\ell-1$):
  \[\mu^{-}_a=0,\qquad\lambda^{-}_a=\lambda(G_{a+1})+\lambda(G_{k-1-a}),\]
  \PROVED{} given Claim~\ref{cl:lvmeascong}(3) (or Proposition~\ref{prop:lveismeas}(3), (5) with Claim~\ref{cl:lvmeascong}(2)) and $\mu(G)=0$: modulo $\ell$ the two substitutions are $T\mapsto T$ and
  $T\mapsto(1+T)^{-1}-1$, both $\lambda$-preserving --- the level-one analogue of Greenberg--Vatsal's formula (no Euler factor at level one; Greenberg--Vatsal held, not read for it). As $\ell\mid B_k/k$,
  $k\equiv b_0$ for an irregular index $b_0$, and $a+1\equiv b_0\iff k-1-a\equiv0$: \textbf{the cell's own irregular index always lands on a pole branch} (\PROVED, elementary); only a second irregular index at the
  same $\ell$ contributes. \CMPD{} at all nine cells with $d\ge1$ ($(30,37)$, $(66,37)$, $(42,59)$, $(56,67)$, $(22,103)$, $(20,131)$; $(282,37)$, $(226,103)$; $(10,691)$): $(\mu^{-},\lambda^{-})=(0,0)$ on every regular
  odd branch, except $(10,691)$, $\lambda^{-}=1$ at $a=199$ and $a=501$ ($\lambda(G_{200})=1$), as predicted. \emph{The pole branches} (through pre.7 a record; \PROVED{} at v1.21.0-pre.8, F4, given
  Proposition~\ref{prop:lveismeas}(4), (6) --- the same inputs as the regular branches plus the polar term --- and $\mu(G)=0$): on both pole branches
  \[(\mu^{-},\lambda^{-})=(0,\ \lambda(G_{b_0})-1),\qquad b_0\equiv k\ \text{the cell's own irregular index},\]
  i.e.\ the formula above holds on every odd branch with $\lambda(G_0):=-1$ for the pole. Proof on the branch $-1$ (the branch $k-1$ is its (FE-inv) partner): with $X=u(1+T)-1$, $Y=u^{k-1}(1+T)^{-1}-1$,
  $x_k=u^{k-1}-1$ and $G_0(X)=\rho_\infty/X+H(X)$, $H\in\Lambda$, Proposition~\ref{prop:lveismeas}(4) gives
  $\mathcal L^{\mathrm{Eis}}_{-1}-\mathrm{polar}_{-1}=-H(X)\,G_k(Y)-\rho_\infty\bigl(G_k(Y)-G_k(x_k)\bigr)/X$, both terms in $\Lambda$ ($Y-x_k\in X\Lambda$). Modulo $\ell$: $X\equiv T$, $x_k\equiv0$,
  $Y\equiv\bar\sigma_1(T):=(1+T)^{-1}-1=-T+\dots$, and $\bar G_k=\bar G_{b_0}$ has $\mathrm{ord}_T=\lambda(G_{b_0})\ge1$ with $\bar G_k(0)=0$; so the second term has $\mathrm{ord}_T=\lambda(G_{b_0})-1$ exactly (its leading
  coefficient $\pm\bar\rho_\infty\bar c_{\lambda}$, $v(\rho_\infty)=0$) and the first has $\mathrm{ord}_T\ge\lambda(G_{b_0})$; by Proposition~\ref{prop:lveismeas}(6) $\mathcal L^{-}_{f,-1}\equiv s(\mathcal L^{\mathrm{Eis}}_{-1}-\mathrm{polar}_{-1})\pmod{\ell^{d}}$
  with $s$ a unit, hence $\mu^{-}_{-1}=0$ and $\lambda^{-}_{-1}=\lambda(G_{b_0})-1$. \CMPD{} at the nine cells: $(0,0)$ on both pole branches, as $\lambda(G_{b_0})=1$ there by (6).
  (So the own irregular index \emph{is} seen by the minus measure, on the pole branches, as $\lambda(G_{b_0})-1$; a cell with $\lambda(G_{b_0})\ge2$, none known, would show $\lambda^{-}\ge1$ there.)
  $\mu^{+}=d$ on every plus branch (Remark~\ref{rmk:lvmeasplus}(ii); at $d=2$ through (5)).
\item \textbf{(The tables, a record.)} $\mu=0$ on every branch of every control ($d=0$) cell: $(10,\ell)$ for the $15$ primes $11\le\ell\le131$ of (6) and $(16,37)$ --- evidence, at finite level, for Greenberg's
  conjecture ($\mu=0$ for residually irreducible $\bar\rho_f$; \STD, recalled), predicted, not proved. The branches with $\lambda=1$ ($\lambda=0$ elsewhere; $\mu=0$ except $\mu^{+}=d$ on the Eisenstein plus side;
  $^{*}$ = even self-dual, forced by (3)):
  \begin{center}\footnotesize
  \begin{tabular}{@{}lll@{}}
  \toprule
  cell & minus & plus\\
  \midrule
  $(10,11)$ & --- & $0^{*}$ (also (4))\\
  $(10,13)$, $(10,17)$, $(10,29)$, $(10,41)$ & --- & ---\\
  $(10,19)$, $(10,31)$, $(10,43)$, $(10,59)$, $(10,103)$ & --- & $14^{*}$, $20^{*}$, $26^{*}$, $34^{*}$, $56^{*}$ resp.\\
  $(10,23)$ & --- & $0$, $10$ ((4)), $16^{*}$\\
  $(10,37)$ & $21$, $25$ & ---\\
  $(10,47)$ & $13$, $43$ & $16$, $40$, $28^{*}$\\
  $(10,67)$ & $31$, $45$ & $38^{*}$\\
  $(10,131)$ & $63$, $77$ & $70^{*}$\\
  $(16,37)$ & $21$, $31$ & $22$, $30$, $8^{*}$, $26^{*}$\\
  $(30,37)$, $(66,37)$, $(282,37)$ & --- & $4$, $26$\\
  $(42,59)$ & --- & $14$, $28$, $50^{*}$\\
  $(56,67)$ & --- & $8$, $48$, $28^{*}$\\
  $(22,103)$, $(226,103)$ & --- & $46$, $78$, $62^{*}$\\
  $(20,131)$ & --- & $42$, $108$, $22$, $128$, $10^{*}$\\
  $(10,691)$ & $199$, $501$ ((7)) & $22$, $678$, $350^{*}$\\
  \bottomrule
  \end{tabular}
  \end{center}
  Every unstarred entry comes in an (FE-inv) pair $\{b,n-b\}$. No $\lambda\ge2$ and no $\mu>0$ occur (other than $\mu^{+}=d$). The tables equal the frozen records of WP-MEAS (\texttt{MeasureFrozenR201.inc}:
  $(10,37)$, $(16,37)$, $(30,37)$) and S2 (\texttt{MeasureFrozenS2.inc}: $(282,37)$, $(226,103)$, $(10,691)$), both signs, no branch skipped. \emph{Unexplained} (recorded, not derived): the $\lambda=1$ pairs of
  the controls away from the self-dual and near-trivial-zero branches --- $(10,37)$ minus $\{21,25\}$, $(10,47)$ minus $\{13,43\}$ and plus $\{16,40\}$, $(10,67)$ minus $\{31,45\}$, $(10,131)$ minus $\{63,77\}$,
  $(16,37)$ minus $\{21,31\}$ and plus $\{22,30\}$. Each such pair is one congruence $c_0\equiv0\pmod\ell$ on a branch (all of its critical values divisible by $\ell$); a branch with unit-sized
  random coefficients has $\lambda\ge1$ with probability $1/\ell$, and the pairs then come in (FE-inv) orbits, so over the $32$ sign-cells of the controls ($(\ell-1)/4$ orbits each) about $\sum(\ell-1)/(4\ell)\approx7.7$ orbits are
  expected by chance; $7$ are seen ($47$ is a regular prime, so no Eisenstein mechanism is available there). The pairs are therefore consistent with chance and carry no claim (F4).
\item \textbf{(The plus side: the McCallum--Sharifi forced zero; through pre.8 ``an observation, recorded, not derived''.)} On the plus side of \emph{every} Eisenstein cell (nine cells, $d=1,2$) $\lambda^{+}=1$ at the (FE-inv) pair $\{-k,\,n+k\}\pmod{\ell-1}$, $k\equiv b_0$
  the cell's own irregular index ($4\equiv-32$, $26\equiv30+32$ at $(30,37)$; $678\equiv-12$, $22=10+12$ at $(10,691)$, while the second index $200$ gives nothing on the plus side); at $(20,131)$ a further
  pair $\{22,128\}=\{k,n-k\}$. This places Remark~\ref{rmk:lvmeasplus}(iii)'s pairs. Not a candidate (F4, pre.8; through pre.7 ``a derivation through the boundary part $r^{+}\equiv u\cdot p_0$ is the natural
  candidate''): $p_0$ contributes nothing to these branches --- $\lambda_{\chi,j}(p_0)=0$ for $\chi\ne\mathbf1$ (Remark~\ref{rmk:lvlegc}~(L2)), so $\mu_{p_0}$ lives on the branches $0$ and $n$ alone, and it is
  $\alpha^{-m-1}(\alpha-1)\ell^{mn}p_0$ on a level-$m$ ball (Proposition~\ref{prop:lvmtt}(6)), $\equiv0\pmod{\ell^{d}}$. The pattern sits in the reduced measure $\ell^{-d}\mu_{r^{+}}\bmod\ell$, i.e.\ in the
  measure of $w=(r^{+}-u\,p_0)/\ell^{d}\bmod\ell$ (Remark~\ref{rmk:lvmeasplus}(iii)), a mod-$\ell$ symbol on which $T_q-(1+q^{k-1})$ acts by $u\,(a_q-1-q^{k-1})\ell^{-d}\,p_0$ --- the next $\ell$-adic digit of
  $a_q$ (the constants of [RC]'s Frobenius shadow). Through pre.8 this item ended: ``A derivation would need the twisted values of $w$; not attempted.''
  \emph{Identified} (v1.21.0-pre.9, slot-21 P5$'$, \texttt{docs/implementation/slot21/P5prime\_note.md}, commit \texttt{08d83ab3}; scripts \texttt{p5p/}, outputs \texttt{p5p/out/}).
  Put $h(b):=\ell^{-d}\mathcal L^{+}_b(0)\bmod\ell$ on the even branches. Given $\mu^{+}_b=d$, $\lambda^{+}_b\ge1\iff h(b)=0$ (\PROVED: every critical point $u^{j}-1$ lies in $\ell\Z_\ell$, so
  $\mathcal L^{+}_b(u^{j}-1)\equiv\mathcal L^{+}_b(0)\pmod{\ell^{d+1}}$, and Weierstrass preparation; $\mu^{+}_b=d$ per branch is read at $M=3$ only, so at the $M=2$ cells a zero of $h$ reads
  ``$\lambda^{+}_b\ge1$ or $\mu^{+}_b\ge d+1$''). Let $x=(x_i)_{i\ \mathrm{odd}}$ be the null vector over $\F_\ell$ of McCallum--Sharifi's relations for the pairing values
  $e_{i,r}=\langle\eta_i,\eta_{r-i}\rangle_r$, $r$ the cell's own irregular index (their relation from $\{\rho_a,\rho_{a-1}\}=0$, $4\le a\le\ell-1$ even, and skew-symmetry $x_i+x_{r-i}=0$;
  \cite{mccallum_sharifi} \S5, read on the page of the authors' web PDF, pp.~12--14, not held in the project library). Then (\CMPD, \texttt{p5p\_sharifi.py}, \texttt{p5p\_report.py}; $32$ cells,
  $28$ irregular pairs $(\ell,r)$ with $37\le\ell\le691$, every even branch, $4961$ branch values, $0$ mismatches, the nullity $1$ at every irregular $(\ell,r)$ and $0$ at the regular $r\pm2$ tried --- not a general fact: the P5$'$ audit finds nullity $10$ at the regular $(257,162)$, where $2$ has order $16$; cf.\ McCallum--Sharifi's
  regular $(73,38)$; audit v1.21.0-pre.10):
  \[h(b)=c\cdot x_{b+1}\quad\text{for every even }b,\ \text{one }c\in\F_\ell^{\times}\text{ per cell}\]
  --- consistent with Sharifi's conjecture, stated in \cite{sharifi_reciprocity}, Introduction, (1.1) (read on the page of the arXiv version by the P5$'$ audit): $\varphi_f((\alpha_t,\alpha_{k-t}))=c_{p,k}\cdot L_p(f,\omega^{t-1},1)$
  with $c_{p,k}\in(\Z/p)^{\times}$ independent of the odd $t$ and of $f$, $f$ of weight $2$, level $p$, character $\omega^{k-2}$ (announced in \cite{sharifi_eisenstein}, Introduction, p.~4: $(\eta_i,\eta_{k-i})$ against $L_p(f,\omega^{i-1},1)$ modulo
  $\mathfrak p^{2}$), seen at level one in this normalisation. The index $t\leftrightarrow\omega^{t-1}$ is Sharifi's; that the branch $b$ (at $j=0$) of the level-one weight-$k$ measure carries the twist $\omega^{t-1}$, $t=b+1$, of the weight-$2$ form
  is read off the data, not derived. McCallum--Sharifi show that the cup-product pairing itself has $e_{\ell-r,r}=0$ (\cite{mccallum_sharifi}, the remark after Cor.~5.2, read on the page:
  $\eta_{\ell-r}$ is the Kummer generator of Ribet's unramified extension, so it pairs trivially; their \S5 assumes Vandiver's conjecture at $\ell$, verified for $\ell<2^{31}$ \cite{hho}), and by skew-symmetry $e_{2r-1,r}=0$; the computed
  vector has $x_{\ell-r}=x_{2r-1}=0$ as well (they observe it for $p<10^{4}$; \CMPD{} here); under $i=b+1$ these are exactly $b\equiv-k$ and $b\equiv n+k$.
  \textbf{Theorem, conditional on (S)} ``$h(b)=c\cdot e_{b+1,r}$ with $c\in\F_\ell^{\times}$ at a $d=1$ cell'' (Sharifi's conjecture in this level-one normalisation; not \PROVED{} as a uniform statement; at the computed cells see \emph{The modular-symbol proof} below, v1.22.0-pre.2):
  \emph{at every level-one Eisenstein cell with $d=1$, $h(b)=0$ at $b\equiv-k$ and $b\equiv n+k\pmod{\ell-1}$, hence $\lambda^{+}_b\ge1$ there wherever $\mu^{+}_b=d$} (in general ``$\lambda^{+}_b\ge1$ or
  $\mu^{+}_b\ge d+1$''; $\mu^{+}_b=d$ is read at $M=3$, $\ell\le157$) --- from (S), $e_{\ell-r,r}=e_{2r-1,r}=0$ and the reduction above: the pair $\{-k,n+k\}$ is the McCallum--Sharifi forced zero.
  \emph{Status of (S)} (P5$'$ audit, \texttt{docs/implementation/slot21/rc/P5prime\_audit\_report.md}): Fukaya--Kato \cite{fukaya_kato} (Introduction read on the page of the authors' 2012 preprint, not held) prove Sharifi's
  conjectures in the $\Lambda$-adic form --- $\varpi\circ\Upsilon=1$ and $\Upsilon\circ\varpi\equiv1$ modulo torsion (Thm.~0.15), $\Upsilon(C)\equiv L^{-}$ modulo torsion (Thm.~0.20), componentwise (Rmk.~0.17(3)) --- under
  (i) ``$\xi$ has no multiple zero'' or (ii); (i) holds on the relevant component for $\ell<2^{31}$ by \cite{hho} ($\lambda(G_b)=1$ at each irregular index, with the \STD{} steps of (7)). \emph{Not checked}: the descent from
  their $\Lambda$-adic statement to the finite-level (1.1) modulo $\mathfrak p^{2}$, and from the weight-$2$ level-$\ell$ form to level one, weight $k$, in this normalisation (Hida's control theorem, recalled). The theorem therefore
  stays conditional \emph{as a uniform statement}; at the computed cells it no longer needs (S) (below). \emph{No novelty is claimed}: given (1.1), the forced zero is an immediate consequence of McCallum--Sharifi's remark; the audit did not find it written out in the four texts read
  (\cite{mccallum_sharifi} \S5, \cite{sharifi_eisenstein} and \cite{sharifi_reciprocity} Introductions, \cite{fukaya_kato} Introduction) and did not search later work (e.g.\ Wake--Wang-Erickson). A second irregular index $c\not\equiv r$ forces nothing on the plus side of the $r$-cell (its generator $\eta_{\ell-c}$ kills the projection to a different eigenspace,
  the pairing $\langle\ ,\ \rangle_c$): the comparison with the McCallum--Sharifi vector of every other irregular index fails (\CMPD{} at $691$ both ways --- $(198,691)$ sees $\{490,398\}=\{-200,n+200\}$, not $\{678,210\}$ ---, $617$, $467$, $379$, $353$, $157$).
  The further zeros ($\{22,128\}$ at $(20,131)$, $\{12,54\}$ at $(66,101)$, $\{196,438\}$ at $(18,617)$, $\{118,196\}$ at $(82,233)$, \dots) are zeros of the McCallum--Sharifi vector itself, properties of
  $(\ell,r)$ constant along the Hida family (\CMPD), not forced by the relations' visible symmetries (a pointer, not checked: by \cite{sharifi_reciprocity}, Introduction, under mild assumptions $(\alpha_t,\alpha_{k-t})$ vanishes
  only if the Selmer group over $\Q$ of $T_f(t)$ is non-zero). Audit (v1.21.0-pre.10): the identity re-derived independently (own code from the relations as printed, the measure read through the verb) at $(30,37)$,
  $(282,37)$ ($d=2$) and at three irregular pairs outside the $28$: $(257,164)$, $(293,156)$, $(353,300)$ (cells $(162,257)$, $(154,293)$, $(298,353)$): every even branch, one scalar, $0$ mismatches, the zeros exactly $\{-k,n+k\}$ (\CMPD). \emph{The modular-symbol proof} (v1.22.0-pre.2; slot-22 P5$''$, \texttt{docs/implementation/slot22/P5pp\_note.md} v0.2, scripts \texttt{p5pp/}; Fable referee \texttt{P5pp\_referee.md}: ACCEPT WITH REPAIRS,
  $0$\,B\,/\,$2$\,M\,/\,$7$\,m, the repairs R1--R8 applied in v0.2). It uses neither Sharifi's conjecture nor Fukaya--Kato. Write $r^{+}=u\,p_0+\ell^{d}w$ --- the hypothesis (H1), i.e.\ (L1) of
  Remark~\ref{rmk:lvlegc} at depth $d$. Then $h(b)$ is the $\omega^{b}$-Fourier coefficient of the Ash--Stevens image $\Psi=A_{\bar w}$, $A_P(c:d)=P(c,d)$, on the Manin symbols $[a:1]=\{\infty\to a/\ell\}$ of
  Fukaya--Kato--Sharifi's module $M'_1=H_1(X_1(\ell),C',\Z)^{+}$ (T.~Fukaya, K.~Kato, R.~Sharifi, \emph{Modular symbols in Iwasawa theory}, Iwasawa Theory 2012, Springer 2014, \S\S2.1.2--2.1.8, read on the authors'
  web copy; only the soundness of its presentation of $M'_1$ and $\langle a\rangle^{-1}[u:v]=[au:av]$ are used). $\Psi$ is even, of character $\omega^{2-r}$, kills the relations of $M'_1$, and is Eisenstein for Merel's sum on these labels with
  eigenvalue $q^{2-r}+q$, because the defect $(T_q-1-q^{k-1})\bar w=\bar c_q\,\bar p_0$ lies on the boundary symbol, whose image vanishes on $M'_1$ (\PROVED{} there, from Proposition~\ref{prop:lvmtt}(6) and the \STD{} inputs below). At each
  computed pair the functionals with these properties form one line, McCallum--Sharifi's $\mathrm{span}(x)$ (``L3'': P5$''$'s script at $33$ irregular pairs --- P5$'$'s $28$, $(257,164)$, $(293,156)$, $(353,300)$ and the second indices $(617,174)$, $(617,338)$ ---
  re-run with independent code by the referee at three of them and three regular controls; a computation per pair). Hence $h(b)=y_{b+1}=c\cdot x_{b+1}$. The two convention-dependent steps (the eigenvalue $q^{2-r}+q$ and the character
  $\omega^{2-r}$, derived from Atkin--Li and Fukaya--Kato--Sharifi's normalisation) were verified directly by the referee on the object: an operator identity on all of $W^{+}$ for $q=2,3$, the character $2-r$ alone, the eigenvalue unique among five candidates,
  $h(b)=y_{b+1}$ with scalar $1$, at $(30,37)$, $(66,37)$, $(102,37)$. (H1) holds: at the $28$ cells with $d=1$, $n<\ell$ it is \PROVED{} from P5$'$'s own computed input $\ell\mid\mathcal L^{+}_b(0)$ (every even $b$), through the cusps of $X_1(\ell)$
  over~$0$ and the injectivity of $P\mapsto A_P$ for $n\le\ell$; at $(282,37)$, $(226,103)$ it is \CMPD{} at depth~$2$ (Remark~\ref{rmk:lvlegc}); at $(66,37)$, $(102,37)$ ($n>\ell$) it is computed at depth exactly~$1$ by the
  referee's Python, one route (Remark~\ref{rmk:lvlegc}; the C++ gate has not read it). \textbf{Status (v1.22.0-pre.2), per pair:} at the $30$ cells with the project's own readings of (H1), and at $(66,37)$, $(102,37)$ once (H1) is gated (so far the referee's Python, one route, the C++ gate has not read it), the theorem above --- $h(b)=0$ at $b\equiv-k,\,n+k$, hence $\lambda^{+}_b\ge1$ there wherever $\mu^{+}_b=d$, now at $d=1$ and $d=2$ alike --- and (S), $h=c\cdot e_{b+1,r}$ with $c\in\F_\ell^{\times}$ (additionally given $h\not\equiv0$, computed, and
  $e\neq0$, \cite{sharifi_eisenstein} Thm.~1.2), are \textbf{\PROVED{} modulo \STD{} inputs, given \CMPD{} L3 at the pair and \CMPD{} (H1) at the cell}. The \STD{} inputs (recalled, not read): Manin's presentation of $M_1(\Gamma_1(\ell))$;
  Merel's sum as $T_q$ on $W$ and on $M_1(\Gamma_1(\ell))$ (Merel, LNM 1585, 1994, \S2; the path condition is \HA{lem:merelpath}); the cusps of $X_1(\ell)$ over~$0$ and the diamond action on them. \emph{Nothing uniform in $(\ell,r)$ is
  claimed}: L3 is a computation per pair, and the statement for every irregular pair stays open (and conditional on (S) as above). (S) is P5$'$'s level-one statement $h=c\cdot e$, of the shape of Sharifi's (1.1) and
  \emph{consistent with Sharifi's conjecture (1.1)}; its identification with (1.1) --- the index map $t=b+1$, the descent from weight~$2$, level~$\ell$ --- is untouched, and (1.1) itself is not claimed here.
  The $d=2$ match at $(282,37)$, $(226,103)$ (through pre.1 ``a record'') is thereby explained. No novelty is claimed.
  Conversely the minus side sees the second index through (7)'s Greenberg--Vatsal formula: $\lambda^{-}=1$ at
  $\{107,109\}$ of $(60,157)$ and at $\{47,61\}$ of $(108,157)$ (\CMPD, $M=3$), as at $(10,691)$. \emph{Still open}: $\lambda^{+}=1$ exactly at the forced branches (\CMPD{} where $M=3$ reads it, $\ell\le157$;
  at $M=2$ only ``$\ge1$''); the non-forced zeros of $x$; (S) as a statement uniform in $(\ell,r)$ (at the computed cells \PROVED{} modulo \STD{} inputs since pre.2, above), and the index map $i=b+1$ against Sharifi's $t$; the $d=2$ cells other than $(282,37)$,
  $(226,103)$ (Sharifi's statement is modulo $\mathfrak p^{2}$; at these two the digit $\ell^{-2}\mathcal L^{+}_b(0)\bmod\ell$ matches $x$, through pre.1 a record, since pre.2 explained by the modular-symbol proof); a \emph{uniform}
  modular-symbol proof (the Eisenstein part of $M'_1\otimes\F_\ell$ one-dimensional at every irregular pair, or a symbolic derivation of McCallum--Sharifi's relations from the Manin and $T_2$-Eisenstein relations;
  through pre.1 ``not attempted'' for any pair, since pre.2 done pair by pair at the $33$ computed pairs; the literature's map between Manin symbols and cup products is
  Sharifi's $\varpi$, \cite{sharifi_reciprocity} \S5, with Fukaya--Kato's $\Upsilon$ the other way; not compared).
\end{enumerate}
\emph{Not covered}: P5~(a) (the $\mathcal L$-invariant of 11a) is P9's V12 (Remark~\ref{rmk:lvtwovar}); characters of $\ell$-power conductor enter only through (2)'s proof; $\lambda\ge2$ is nowhere seen.
\end{remark}

\begin{checkitem}[The gate of P5: \texttt{HeckeLValueIwasawaInvariantsValidation}]\label{chk:lviwasawa}
\SPEC{} (the checks); expected values P5's predictions, committed first (\texttt{P5\_spec.md} v0.1, \texttt{0dd6c269}). \emph{Registration}: \texttt{Registry/\allowbreak Slot21/\allowbreak P5\_iwasawa.cpp},
\texttt{cmake/\allowbreak slot21/\allowbreak P5\_iwasawa.cmake}; source \texttt{HeckeExt/\allowbreak LValue/\allowbreak HeckeLValue\allowbreak Iwasawa\allowbreak Invariants\allowbreak Validation.\{h,cpp\}}; the public APIs of WP-MEAS, WP-LIN
and WP-LV only, no existing file's code edited. Tier \texttt{standard}, read-only. Arms: (I0) Remark~\ref{rmk:lviwasawa}(6) by both routes at $15$ primes ($691$ in the heavy arm); (I1) item (1), negatives: the flipped sign, and
``$\chi$ for $\bar\chi$'' counted only on equal-valuation pairs of power $\ell^{R-v}\ge10^{5}$; (I2) item (2) in the group ring, negative: the flipped branch sign; (I3) the tables of item (8) by the group-ring invariants
and the values oracle (5), (FE-inv), (FE-parity), item (4), $\mu=0$ at the controls; (I4) item (7) with $\lambda(G)$ from (I0)'s Bernoulli-only route, $\mu^{+}=d$; (I5) the tables against
\texttt{MeasureFrozenR201.inc}, \texttt{MeasureFrozenS2.inc}. Lines \texttt{IWA\_KL}, \texttt{IWA\_FE}, \texttt{IWA\_GRFE}, \texttt{IWA\_TABLE}, \texttt{IWA\_GV}, \texttt{IWA\_FROZEN},
\texttt{P5\_ARMS checks=<c> failing=<f>}; \texttt{EFFECTIVE \dots{} arm=DEFAULT cells=22}. \emph{First run} (\CMPD): $580$ checks, $0$ failing (I0 $60$, I1 $132$, I2 $88$, I3 $270$, I4 $24$, I5 $6$),
$21.6$\,s, $0.10$\,GB. Read-only env arm \texttt{HECKE\_LVALUE\_IWASAWA\_HEAVY=1} (\texttt{arm=HEAVY cells=25}: the Kubota--Leopoldt table at $691$, the cells $(10,691)$, $(282,37)$, $(226,103)$):
$658$ checks, $0$ failing, $72$\,s, $0.10$\,GB. \texttt{GOLDEN\_STRICT=1}: A $11827/0$, B $6231/0$. No golden line, record key or digest; \texttt{--list} $192$. Two negative controls were repaired before the first green run,
no predicted value changed (\texttt{P5\_spec.md}~\S8: the ``$\chi$ for $\bar\chi$'' control had no power at high valuation --- a match at power $\ell^{2}$ is not evidence; the group-ring negative is vacuous where
$g_b\equiv0\pmod{\ell^{\min(M,R)}}$ and is then not counted).
\end{checkitem}

\appendix
\section{Index of claims with STATUS}\label{app:claims}
%======================================================================
\noindent The seven claims of \S\ref{sec:imported} that v1.2.2 indexed here --- the amalgam
presentation, the dimension of $H^1$, the torsion origin of the Dickson classes, the boundary
groups, the six-term sequence, the boundary characters and Poincar\'e duality --- moved into
[C] with the scope cut and are no longer statements of this document. Their status is [C]'s
own STATUS index, where four of them are also \emph{corrected}
(\S\ref{sec:imported}); \HA{cl:pd} in particular is carried there as \textsc{open}, not as
\STD.

\begin{center}\footnotesize
\begin{longtable}{@{}l>{\raggedright\arraybackslash}p{3.6cm}>{\raggedright\arraybackslash}p{10.6cm}@{}}
\toprule
ref & status & content \\
\midrule
\endhead
\ref{cl:harder} & \STD\ / verified ($10$ Kummer $+$ $27$ converse cells, 2026-09-16) & Harder denominator / Ribet; $r^1=0$ iff $\ell\mid\num(B_k/k)$, read mod $\ell$ as ``$r^1=0$ on the joint eigenspace $E$'', which is two-dimensional at a Kummer cell \\
\ref{cl:local} & \PROVED & $\Ext^1$ vanishes across maximal ideals \\
\ref{cl:gen} & \PROVED & generation $\Leftrightarrow$ separation + tangent-space spanning \\
\ref{cl:sigma} & \STD & finiteness of $\Sigma(\ell)$ (Serre, Jochnowitz, Ash--Stevens) \\
\ref{cl:sep} & \PROVED & existence of $S_{\mathrm{sep}}(\ell)$ \\
\ref{cl:Tl} & \STD & $T_\ell$ indispensable; $\theta^{\ell-1}$ and $E_{\ell-1}$ doublings \\
\ref{cl:sturm} & \STD\ / \CONJ\ (a) / \PROVED\ (b) & cuspidal Sturm bound $\lfloor k/12\rfloor$; the $\Z$-module statement for $M_k$ with $B(k)$ is \emph{false} ($\ell=5,7,131,\dots$) \\
\ref{cl:bprime} & \PROVED & the computed bound $B'(k,\ell)$ certifies generation over $\Z_\ell$ \\
\ref{cl:sturmmodl} & \PROVED\ (given a recorded $B'$) & Sturm certifies generation on non-Dickson $\m$ of $X_4$ only \\
\ref{cl:mkfaithful} & \PROVED\ (v1.15.1-pre.4, WP-19; referee pending) & $\T/\ell$ of $M_k$ has dimension $\dim M_k$ on $M_k(\Fl)$, $\ell\ge5$; $B'(k,\ell)$ always exists \\
\ref{cl:formsF9} & \PROVED\ (v1.15.1-pre.4, WP-19; referee pending) & F9 $=$ Eichler--Shimura per block, exact up to $\dim H^0(\Ga;\Mn)_\lambda$ (the torsion guard) \\
\ref{cl:bk} & \STD\ (recalled) / verified ($(160,691)$, 2026-09-16) & Bella\"iche--Khare: $\T_{\bar\rho}\cong\Fl[[x,y]]$; the campaign example is resolved --- seven local blocks (v1.3.3); $T_5$'s $4{:}1$ is the family's $4{:}2$ block, not split by the refinement, and the block partition is already $T_2$'s (v1.7.0, \texttt{HeckeExtSeparationValidation}) \\
\ref{cl:Tlnil} & (1) proved in [C] (v1.12.0-pre.1; through v1.11.1 \STD); (2) \STD & $U_\ell\circ\theta=0$ over $\Fl$, on cochains (\HA{prop:thetatwist}; checked by (6e) of \S\ref{sec:identities}, slot-11 WP-6, v1.12.0); $T_\ell$'s Jordan type \\
\S\ref{sec:thetaobj} & specification, no claim (v1.12.0-pre.1; built, v1.12.0; (h)'s refusal as specified in the build, v1.12.1) & the Dickson--$\theta$ objects: names, indexing, the six terms and five maps, the census hand-off, refusals; the mathematics is [C]'s \\
\S\ref{sec:census} & specification, no claim (v1.12.0-pre.1; edition~1 built, v1.12.0) & the congruence-census catalogue \texttt{hecke\_census/v1}: rows, columns, the boundary-split, $\pm$, $\theta$-image, Glover and $\Z/\ell^k$-depth rules, extremal criteria, layout \\
\ref{cl:glover} & \STD\ (1) / \PROVED\ (2)--(4) (v1.5.0; (3) \PROVED\ since v1.21.0-pre.5) & the Glover factors of $\Sym^n$ by the $\theta$-recursion, and the standard-module convention that fixes the $\det$-twists \\
\ref{cl:gap} & \PROVED & per-operator refinement can miss a split; locality criterion; enforced in the loop since E0b (slot 2) \\
\ref{rmk:coef} & \STD & the coefficient field $K\cong\Flr$ exists and is unique (Hensel) \\
\ref{cl:resfield} & \PROVED & $K=\Fl[\hat w]$, $\hat w=w^{\ell^N}$; the $P_j$, $n_j=t_j-P_j(\hat w)$ \\
\ref{cl:sharpbound} & \PROVED\ (v1.5.0) & $\dim_{\Fl}\bar A=\operatorname{lcm}_p\deg\pi_p$ on a \emph{local} block, so budget mode's \texttt{residue\_degree\_lower\_bound} is sharp there \\
\ref{cl:ext1} & \PROVED & $\Ext^1_{\T_\m}(\F,\F)=(\m/\m^2)^\vee$ \\
\ref{cl:yoneda} & \PROVED & Yoneda class $=$ structure map, complete invariant of two-step modules \\
\S\ref{sec:classdata} & \STD & wildness for embedding dimension $\ge2$, Loewy length $\ge3$ \\
\ref{pr:eiscusp} & \STD\ (reworded, v1.12.0; ``cuspidal'' read in characteristic zero, v1.12.1) & Eisenstein--cuspidal congruences (Ribet, Mazur), for $E_k$'s own class ($U_\ell=1$), $k\not\equiv0\pmod{\ell-1}$; the literal pre-v1.12.0 ``iff'' over every twist fails at $(10,7)$ \\
\ref{pr:cuspcusp} & \CONJ & cuspidal--cuspidal extensions, $T_\ell$ usually active \\
\ref{pr:cone} & \CONJ & boundary/cone gluings \\
\ref{pr:dual} & \STD; computed in E7 (accepted, slot 7; v1.8.2) & duality: Loewy layers $\leftrightarrow$ socle layers in reverse order \\
\ref{pr:galois} & \CONJ & Galois-cohomological interpretation of active directions \\
\ref{pr:conc} & \STD & concentration at few $\m$ for $n\gg\ell$ \\
\ref{pr:latticesplit} & \CONJ\ as restated a fifth time (v1.25.0-pre.1, owner ruling of 2026-11-06: irreducible classes \emph{of one weight}, replacing ``not dihedral''); the v1.22.0-pre.1 form \textbf{refuted 2026-11-06} on cohomology at the non-ordinary class of $f_{16}$ at $\ell=59$ (four census-4.0 rows, $n=884,942,944,1000$, $(T_2-a_2)^2\ne0$; \texttt{slot25/R59\_verification.md}); restated a fourth time (v1.22.0-pre.1, owner ruling of 2026-10-09 on census finding F33-2: dihedral classes excluded); the v1.17.0-pre.4 form \textbf{refuted 2026-10-09} on cohomology at the dihedral classes at $\ell=23$, $31$ ($161$ census-3.3 rows, $A\cong\Fl[x]/x^3$; the slot-22 referee \texttt{F33-2\_referee.md}); restated a third time at v1.17.0-pre.4 (owner ruling of 2026-10-02 on the referee F1's B-1); the v1.17.0-pre.2 Kummer-free form \textbf{contradicted 2026-10-02} on the forms side at irreducible congruence classes, $(2210/2278,67)$, $(3798/3930,131)$, $(4030/4170,139)$ ($139$ regular; $\Fl[x]/x^4$; congruences confirmed over~$\Z$; six further pairs the same; slot-17 R17-2~\S7); the v1.16.0 form \textbf{refuted 2026-10-02} on cohomology at $(608,547)$ ($g=2$ first window, $A\cong\F_{547}[x]/x^3$, saturated; slot-17 R17-2) and contradicted on the forms side at the swapped classes at $37$, $103$, $691$; the v1.14.0 statement \textbf{refuted 2026-10-01} at $\ell=37$ (18 certified C1a records, the forms side, the $\Z/37^2$ lattice at $(654,37)$) & (R): $\m_A^2=0$ on every $\Fl$-rational non-Dickson $H^1$ block that is congruence-free (for every even $w\le\ell+1$ and every twist the generalised eigenspace of the system in $M_w(\Fl)$ has dimension $\le1$; on reducible classes this is Kummer-free) and, when $\bar\rho_\m$ is irreducible, of one weight (exactly one pair $(w,j)$ of that test with a non-zero eigenspace; this excludes the dihedral, the non-ordinary and the companion classes); nothing asserted on the other classes; no counterexample in census~4.0 ($672$ rows of Loewy length $\ge3$, none in scope; $5805$ irreducible rows of one weight with $\dim V\ge6$, all with $\m_A^2=0$). History: (v1.22.0-pre.1) the same with ``not dihedral'' ($\bar\rho_\m$ irreducible with $\bar\rho_\m\cong\bar\rho_\m\otimes\omega^{(\ell-1)/2}$) for ``of one weight''; no counterexample among the $100\,025$ non-Dickson, non-dihedral residue-degree-$1$ rows of census~3.3; (v1.17.0-pre.4) the same on congruence-free classes; (v1.17.0-pre.2) the same on Kummer-free classes (residual irreducible, or reducible with both seeds $\ell$-regular), ``in particular at every regular~$\ell$''; (v1.15.1-pre.7, released v1.16.0) $\m_A^2=0$ off the Kummer--Frobenius locus; (v1.14.0-pre.2; held on $654/654$ census rows, slot-13 WP-9) C-F9: $\m_A^2=0$ on every $\Fl$-rational $H^1$ block at $\ell\ge5$ (the lattice nearly split); not explained by the Serre-weight count; evidence non-Dickson only, zero evidence on a non-reduced Dickson-type $\m$ (v1.14.0-pre.3); a third route and $452/452$ blocks on $348$ cells, the $206$ Dickson-type blocks read all reduced (v1.15.0, slot-14 WP-M1, WP-6); $1233/1233$ non-reduced blocks beyond the census window on the first cluster, \texttt{noctua2\_section\_m/} (v1.15.1-pre.3, P12) \\
\ref{pr:kfonset} & \CONJ\ (v1.15.1-pre.7; restated v1.17.0-pre.1; hypothesis $g_{\ell+1-w}=1$ added v1.17.0-pre.3) & on a Kummer--Frobenius $\m$ with $1\le t\le J(w)-1$ and $g_{\ell+1-w}=1$, $A\cong\Fl[x]/x^4$ in the window $\ell w+t(\ell+1)\le k<w+t(\ell+1)+(\ell^2-1)$ (``from $k$ on'' withdrawn: forms side $(1,2,2,1)$ at $k=1412$, $1450$); $24/24$ onset hits at $\ell=37$, $0$ false positives; (v1.17.0-pre.2) the lower end held off $37$: $(736,67)$ on cohomology, $(804,67)$, $(1002,59)$, $(1062,59)$, $(3534,101)$ on the forms side; the first Loewy-$4$ cell at $103$ is the swapped $(2574,103)$, outside the prediction; (v1.17.0-pre.3) nothing asserted at $g\ge2$: the pre.2 form failed at $(547,486)$ ($(608,547)$ has Loewy length~$3$ below the onset, $\dim V=12$ in the window);
(v1.18.0-pre.1, generalised, \CONJ) (A) as above; (B) the swapped family at $g=1$, the same window with $w=w_K$ (forms side: onsets at $37$, $103$, $691$; last cell
and first cell past the window at $(1400,37)$, $(1438,37)$, $(10632,103)$ / $(1436,37)$, $(1474,37)$, $(10734,103)$); (C) $g=2$, $A\cong\Fl[x]/x^3$ in the first
window, both families ($(608,547)$ on cohomology, six forms-side cells; v1.18.0-pre.2: three more, the $a=1$ window sampled at cells $1$--$4$, $31$, $62$ of $62$;
(B) five more at $37$, $103$, the referee F1); nothing asserted at the $g=2$ meeting (there Prediction~\ref{pr:depth}'s datum $\F_{547}[x]/x^6$), at $g\ge3$, or past the windows \\
\ref{cl:seedtwist} & \PROVED\ (v1.18.0-pre.1; slot-17 R17-1 Lemma~3.1) & two classes $1\oplus\omega^{w-1}$, $1\oplus\omega^{w'-1}$ agree up to twist iff $w'\in\{w,\ell+1-w\}$; the Kummer classes of a pair are the two Eisenstein ends, the standard and the swapped family, each with one presentation $\omega^t(1\oplus\omega^{w-1})$, $1\le t\le J(w)-1$ \\
\ref{cl:selfmirror} & \PROVED\ (given the class number formula, \STD, locator not verified; v1.18.0-pre.1; R17-1 Lemma~3.2) & $\ell\equiv3\ (4)$, $\ell>3$: $\ell\nmid B_{(\ell+1)/2}$ \\
\ref{cl:dimformula} & \PROVED\ (given \STD\ inputs; Edixhoven Thm.~2.6: attribution verified in a secondary source, locator pending (v1.18.0-pre.2, R18-4); second locator Berger--Li--Zhu Thm.~4.5 for $w\le\ell$; v1.18.0-pre.1; R17-1 Theorem~4.1) & $\dim H^1(\Ga;\Mn)_\m=m_1(1+2g_{w_1})+m_2(1+2g_{w_2})$ on every reducible non-Dickson $\m$, $\ell\ge5$; cuspidal: torsion-free, $\rank\T_\m=\frac12\dim$; census check $21\,690/21\,690$ \\
\ref{rmk:dimwindows} & \PROVED\ (from Claim~\ref{cl:dimformula}) and computed (v1.18.0-pre.1) & the first and second windows on the multiplicity side, $1492/1492$ standard and $1906/1906$ swapped classes of $13$ pairs; the exact event sequence checked by \texttt{wp3/window\_check\_f1.py} (v1.18.0-pre.2) \\
\ref{pr:depth} & \CONJ\ (v1.18.0-pre.1; the $\ell=37$ report's \texttt{conj:depth}) & $\T_\m/\ell\cong\Fl[x]/x^{m r_0}$ for $m\le2$ on Kummer classes ($r_0=1+g$) and irreducible congruence classes ($r_0=2$); $73$ curvilinear blocks on the forms side, $26$ on cohomology (v1.18.0-pre.2; pre.1: $65$); fails at $m=3$; $(m,r_0)=(2,3)$ held at $(34460,547)$, $\F_{547}[x]/x^6$ (v1.18.0-pre.2, the referee F1; pre.1: untested), so the hypothesis is $m\le2$, not $m r_0\le4$ \\
\ref{cl:258alg} & \CMPD\ (1),(3)--(5), slot-19 R19-1, two routes; \PROVED\ (2) from \HA{prop:cohforms} (v1.19.0-pre.1) & $(258,37)$: $\T_\m=\Z_{37}[T_2-29]$ the full Hecke algebra, $=\T^{\mathrm f}_\m$; $\tilde\T/\T_\m\cong\F_{37}[\varepsilon]/\varepsilon^{2}$; depth $1/2$; the over-orders $\T_\m\subsetneq R=(\m:\m)\subsetneq\tilde\T$, all Gorenstein (Bass) \\
\ref{cl:258lattice} & \PROVED\ given \CMPD\ inputs and \STD\ (Bass, \S6, locator pending, not in the library; the Gorenstein input proved inline at pre.2) (v1.19.0-pre.1); the sign \CMPD\ (WP-LC, pre.3) & $\Lambda_\m\cong\T_\m\oplus\m$; $\Lambda_\m^{+}\cong\T_\m$, $\Lambda_\m^{-}\cong\m$ \\
\ref{chk:latcmp} & (a) \PROVED; (b)--(e) \PROVED\ given the \CMPD\ inputs (v1.19.0-pre.1); run by slot-19 WP-LC: held at $M=1,\dots,4$, \CMPD\ (pre.3) & $H^{1}(\Ga;M_{258}\otimes\Z/37^{M})_\m$, $M\le4$: rank, $\mu=3$, Smith valuations of $t$ and $U_{37}$, per sign, $T_p=Q_p(t)$ \\
\ref{rmk:latcmpseven} & \CMPD\ (slot-19 WP-LC; F4 m5--m7) (v1.19.0-pre.3) & $\ell=7$: per-sign freeness over $\T_\m$ fails on $5$ of $6$ signs at $(48,7)$, $(90,7)$; $(48,7)$ plus $\cong\Z_7[t/7]$, told from $\mathcal O_{\max}$ by $U_7$ \\
\ref{def:lvroute}, \ref{rmk:lvshapes} & \SPEC\ (v1.19.0-pre.1; kind \textsf{congruence} added at pre.2, F1 M-2) & the $L$-value interface: domain, kind (\textsf{linear}, \textsf{bilinear}, \textsf{congruence}), normalisation record, evaluation, index statement; Birch (linear) and Hida (bilinear) shapes declared at pre.1, implemented at pre.6 (below); Eisenstein = congruence; v1.19.0-pre.5: research prototypes of (B), (H) \CMPD{} (slot-19 R19-8, R19-9), (H)'s constant from Hida Thm.~5.1 (read), slot-20 specs; v1.19.0-pre.6: (B), (H) implemented (slot-19 WP-LIN, WP-BIL), gates \CMPD; (H2) at depth~$2$ \CONJ \\
\ref{rmk:lvlegc} & \CMPD\ at $8$ cells, (L2) \PROVED\ from (L1); \CONJ\ at depth $\ge2$ (v1.19.0-pre.6, WP-LIN) & leg~(c) on $H^{1}_c$: $r^{+}\equiv u\,p_0$, $r^{-}\equiv s\,\mathrm{Bern}_k$ modulo~$\ell$ to depth $v(B_k/k)=1$; twisted values of $r^{+}$ $\equiv0$; v1.20.0-pre.1: computed to depth~$2$ at $(282,37)$, $(226,103)$ (slot-20 R20-1), still \CONJ{} there; v1.20.0-pre.6: (L1) \CMPD{} at depth~$2$ there (WP-LV3, C++), item~3 still \CONJ; v1.22.0-pre.2: (L1) at depth exactly~$1$ at $(66,37)$, $(102,37)$ computed (the P5$''$ referee's Python, one route) \\
\ref{def:lvmeasure} & \SPEC\ (v1.20.0-pre.1, owner O5) & a measure is a family of \textsf{linear} evaluations, one per ball; index statements are congruences of measures; the accessor \texttt{measure()} \\
\ref{def:lvmtt}, \ref{prop:lvmtt} & definition; \PROVED\ (v1.20.0-pre.1, slot-20 R20-1); the analytic constant has (B)'s status; $\tau(\chi)$ at $\ell$-power conductor \CMPD\ (v1.21.0-pre.3, slot-21 P4) & the $\ell$-stabilised MTT measure $\alpha^{-m}c(a,m)-\alpha^{-m-1}\ell^{k-2}c(a,m-1)$: a measure, its moments, $E_{\mathbf1,j}=(1-\ell^{j}/\alpha)(1-\ell^{n-j}/\alpha)$, $E_{\chi,j}=\alpha^{-m}$ \\
\ref{def:lvlfun}, \ref{lem:lvlfunprec} & definition; \PROVED\ (v1.20.0-pre.5, slot-20 S1; elementary) & the branch series modulo $(\ell^{R},\omega_{M-1})$ as $T$-coefficients; $\mathrm{prec}(0)=R$, $\mathrm{prec}(t)=\min(R,M-1-\lfloor\log_\ell t\rfloor)$; the record \texttt{lfunction/v1} \\
\ref{cl:lvmeas691} & \CMPD\ at $(10,691)$ (v1.20.0-pre.5, slot-20 S2; predictions and the Python oracle committed first) & $\Delta$'s MTT measure $\equiv-s\,G_{a+1}G_{11-a}\pmod{691}$ on $343/343$ off-pole branches; $\lambda=1$ at $\{199,501\}$; pole points a unit off \\
\ref{rmk:lvcoeffring} & over $\mathcal O$ \PROVED\ (coefficient change); the cells \CMPD; Hida modulo $\m$ \CMPD\ at one component; F-P3-1 \CMPD\ as a census (v1.21.0-pre.1, slot-21 P3) & the measure over $W(\F_q)$ and totally ramified $\Z_\ell[x]/(F)$; refused: residual multiplicity $\ge2$, non-ordinary \\
\ref{chk:lvcoeffring} & \SPEC\ (v1.21.0-pre.1); first run $603$ checks, $0$ failing; ``pending P2/P4'' replaced by pointers (v1.21.0-pre.3) & P3's gate \texttt{HeckeLValueCoefficientRingValidation} (\texttt{standard}), arms (R0)--(R8) \\
\ref{rmk:lvtwovar} & assembly and P1 agreement \PROVED; items 1--3 \CONJ\ (rest on \Ucite{pr:u-msym}~(U16)), evidence \CMPD\ as a run; item 4 \STD\ (recalled) and \CMPD\ as a run (v1.21.0-pre.1, slot-21 P9) & the two-variable family by interpolation; the $\Lambda$-adic Eisenstein congruence at $X=\alpha_a$ modulo $37^{3}$, $103^{2}$; weight two and the $\mathcal L$-invariant of 11a \\
\ref{prop:lveismeas} & (1), (2), (4)'s constant, (5), (6) \CMPD; (2)'s twisted masses \PROVED\ from the convolution form and \STD; (3) \PROVED\ given \STD\ (Washington \S12.2, recalled); (4)'s shape \PROVED\ for the branch $-1$; (6)'s valuation $d-1-v(r)$ \PROVED\ (v1.21.0-pre.3, slot-21 P2) & the Eisenstein measure: Dedekind--Rademacher symbol, $\mu^{\mathrm{Eis}}=-E_1\star\iota_\ast E_{k-1}$, branch series $-G_{a+1}G_{k-1-a}$ in $\Lambda\otimes\Q$, the polar term $\mathrm{Res}$, $v(\mathrm{Res})=d$; the cusp form's pole branches; the located $14+4$ derived \\
\ref{chk:lveismeas} & \SPEC\ (v1.21.0-pre.3); first run $225$ checks, $0$ failing (heavy $259$); with E8L2 / E9L2 $239$ ($273$) & P2's gate \texttt{HeckeLValueEisensteinMeasureValidation} (\texttt{standard}), arms (E0)--(E9) \\
\ref{rmk:lvwild} & (1), (2), (3), (5), (6), (7) \CMPD; (4) \PROVED\ (Weierstrass) and \CMPD; (3)'s proof not given (v1.21.0-pre.3, slot-21 P4) & wild characters $\omega^{i}\psi_t$ of conductor $\ell^{m}$: $\tau(\chi)$ against complex $L$-values, root number $\tau(\chi)^{2}/M$; Proposition~\ref{prop:lvmtt}(4); $\lambda_{\chi,j}=-\varepsilon\lambda_{\bar\chi,n-j}$; $v_\pi=e\mu+\lambda$; Kubota--Leopoldt at $\zeta^{t}u^{r}-1$; the polar term visible at every wild critical point, at $(226,103)$ too since v1.21.0-pre.7 (P4x, $128$-bit) \\
\ref{chk:lvwild} & \SPEC\ (v1.21.0-pre.3); first run $201$ checks, $0$ failing (heavy $238$); one prediction wrong (the root number); with P4x $222$ (heavy $259$) & P4's gate \texttt{HeckeLValueWildCharacterValidation} (\texttt{standard}), arms (W0)--(W7)  \\
\ref{rmk:lviwasawa} & (1), (2), (3), (4), (5) and (7)'s formula \PROVED{} ((2) given Proposition~\ref{prop:lvmtt}(4) and Remark~\ref{rmk:lvwild}(3); (7) given Claim~\ref{cl:lvmeascong}(3) and $\mu(G)=0$); (6) \PROVED{} / \STD{} (Ferrero--Washington recalled; \cite{hho} read); (8), (9) records; every run \CMPD{} (v1.21.0-pre.7, slot-21 P5); (9)'s forced-zero theorem and (S) \PROVED{} modulo \STD{} inputs at the $30$ computed cells with the project's own readings of (H1), $+2$ once (H1) is gated at $(66,37)$, $(102,37)$, given \CMPD{} L3 at the pair and \CMPD{} (H1) at the cell (v1.22.0-pre.2, corrected at pre.3, slot-22 P5$''$; per pair; uniformly still conditional on (S)) & the tame functional equation at conductor $1$, $\ell$; $\mathcal L_b(T)=-(-1)^{b}\mathcal L_{n-b}(u^{n}(1+T)^{-1}-1)$; (FE-inv), (FE-parity): $\lambda$ odd at even self-dual branches; the near-trivial zero; the values oracle; $\lambda(G_b)=1$ at the irregular pairs; level-$1$ Greenberg--Vatsal; the $(\mu,\lambda)$ tables; $\lambda^{+}=1$ at $\{-k,n+k\}$ observed \\
\ref{chk:lviwasawa} & \SPEC\ (v1.21.0-pre.7); first run $580$ checks, $0$ failing (heavy $658$); two negative controls repaired, no predicted value changed & P5's gate \texttt{HeckeLValueIwasawaInvariantsValidation} (\texttt{standard}), arms (I0)--(I5) \\
\ref{cl:lvmeascong} & \CMPD\ at $3$ cells (v1.20.0-pre.1, slot-20 R20-1); shape \PROVED\ given its inputs; the pole branches an exception, computed; unexplained through v1.21.0-pre.2 (R20-L), explained by P2's polar term since v1.21.0-pre.3 (Proposition~\ref{prop:lveismeas}(6), \CMPD) & $\mathcal L^{-}_{f,a}\equiv-s\,G_{a+1}G_{k-1-a}\pmod{\ell^{d}}$ off the pole branches, $r^{-}\equiv sP^{\mathrm{Eis}}$; $\mu=\lambda=0$ on the minus side \\
\ref{rmk:lvmeasplus} & (ii) \PROVED\ from (L1), \CMPD; (iii) an observation, no claim (v1.20.0-pre.1); its pairs $\{-k,n+k\}$ unconditional at the computed cells since v1.22.0-pre.2 (Remark~\ref{rmk:lviwasawa}(9)) & $\mu^{+}=d$; the reduced plus measure's $\lambda=1$ branch pairs \\
\ref{chk:lvmeas} & \SPEC\ (v1.20.0-pre.1) & WP-MEAS's gate \texttt{HeckeLValueMeasureValidation} (\texttt{standard}), arms (M0)--(M6); (M7) and the env arm (M691) (v1.20.0-pre.5, S1\,/\,S2) \\
\ref{cl:lvtriangle} & (C) \CMPD; (L), (H) and the triangle \CONJ\ (v1.20.0-pre.1); (L1), (H) and the triangle \CMPD\ at the two cells (v1.20.0-pre.6, WP-LV3), (L)'s item~3 and the general statement \CONJ & at $(282,37)$, $(226,103)$ the congruence, linear and bilinear depths are all $2$; $\m_{\mathrm{Eis}}$ the joint ideal in $T_2,T_3,T_5$ \\
\ref{def:lvlattice} & \SPEC; $c_f=0$ for $n<\ell$ \PROVED; the value open (v1.20.0-pre.1); well-definedness (the $f$-line carried onto the $f$-line) \PROVED{} from \HA{lem:symdual} (v1.20.0-pre.2); value by Proposition~\ref{prop:lvlattice} (v1.20.0-pre.6) & the change-of-lattice constant between $\theta_\ast H^{1}(\Mn)$ and $H^{1}(\Dn)$ on the $f$-line \\
\ref{prop:lvlattice} & \PROVED\ (v1.20.0-pre.6, slot-20 R20-3) & $\theta_\ast$ Hecke- and $\iota$-equivariant; $c_f^{\varepsilon}=e\,v_\ell(n!)-d_f^{\varepsilon}$, $0\le d_f^{\varepsilon}\le e\max_tv_\ell\binom nt$; $d_f^{\varepsilon}=0$ if $H^{1}(\Mn/V'_n)^{\varepsilon}_{\m_f}=0$ \\
\ref{rmk:lvlatticevalues} & \CMPD\ (v1.20.0-pre.6; R20-3 two routes, WP-LV3 C++); the $d_f=1$ twist rule \CONJ\ (\CMPD\ at three cells) & $c=7$ at $(282,37)$, $2$ at $(226,103)$, $6$ ($v_\pi=12$) at $(258,37)$, $0$ for $n<\ell$; $d_f=1$ on five lines at $(40,37)$ \\
\ref{rmk:lviotak} & \PROVED\ from \HA{prop:iotasplit}, \HA{prop:iotacommutes} (v1.20.0-pre.1) & $\iota$ over $\Z/\ell^{k}$, $\ell$ odd: nothing new mathematically; the code honours it (WP-LV3, built at v1.20.0-pre.6; \texttt{hecke\_involution\_ring/v1}) \\
\ref{chk:lvtriangle} & \SPEC\ (v1.20.0-pre.1); outcome \CMPD\ (v1.20.0-pre.6, every gate held) & WP-LV3's gates: generic reader, cross-kind gate, $c_f$, $\iota$ over $\Z/\ell^{k}$, the twisted rung $(37;62,6)$ \\
\ref{cl:lveis} & (1), (3) \PROVED\ from \Ucite{pr:u-U1}; (2) \STD\ for $n<\ell-1$ (Harder, \VERIFYloc), \CONJ\ beyond (v1.19.0-pre.1) & the Eisenstein congruence module of the cuspidal sign is $\Z_\ell/(B_k/k)$; modulo $\ell^{m}$ on a saturated block (\texttt{stable} at $K=1$; pre.2, F1 M-1), a two-sided bound otherwise \\
\ref{cl:lvtwist} & \PROVED\ given \STD\ (Hida's control at nebentypus points, ENS 1986 Cor.~3.2 (3.3b), verified on the page by F1 at pre.2) (v1.19.0-pre.1) & twisted: $\T^{-}/I_{\mathrm{Eis}}\cong\Z_\ell/(B_{k,\chi}/k)$ at $\Ga_1(\ell)$, $\chi=\omega^{j}$; v1.19.0-pre.5: \CMPD{} at $4$ cells ($\ell=37,59,67$, $v=1$) and $2$ controls (slot-19 WP-LV2); v1.20.0-pre.6: \CMPD{} at the depth-$2$ rung $(37;62,6)$, $n=60$ (slot-20 WP-LV3(e)) \\
\S\ref{sec:lvnf} & \SPEC\ (v1.19.0-pre.1); built (v1.19.0-pre.3, WP-NF) & the number-field layer (WP-NF): objects, the local test of R18-2 (\PROVED\ there), the gate \\
\ref{rmk:i8switch} & \PROVED\ (from \cite{bk}; v1.16.3, owner O4) & I-8's split\,/\,ramified switch is not Bella\"iche--Khare obstructedness: both cases unobstructed (Thm.~22 (i), (iii) with Vandiver) \\
\ref{pr:frobeniuslayer} & \CONJ\ (v1.16.3, owner O4); proof \PROVED\ given $D(n_V)=4$ and reduced $n_0$ blocks; the onset \CONJ & the Frobenius layer (F1, law L-F1): the $\bar\rho_f\otimes\omega^j$ block doubles to $D=4$ at $n_V=\ell k_0+j(\ell+1)-2$ and is hidden from there; WP-12 test $0$ failures at $17$, $19$; slot~16: $3/3$ hidden-chain cells; slot~19: $9/9$ irreducible rungs $\le1220$ hidden at $(1220,37)$, $0$ spurious (v1.19.0-pre.5); slot~20 (v1.20.0-pre.3): family-(c) onsets, $\theta^{j}S_{k_0}=\ker x^{2}$ on forms --- the transport step under a non-reduced source, outside the hypotheses; slot~20 (v1.20.0-pre.6): WP-SW's $1258$, $1296$, $1334$ at $37$, $9/9$ rungs hidden at each; one unpredicted hidden block at $1334$ from $k_0=\ell+1=38$, explained by the mechanism from $n=1296$ (the early $V$-copy); v1.20.0-pre.7: refuted as stated at $k_0=\ell+1$ (one class, $(1334,37)$), hypothesis restricted to $k_0\le\ell$ \\
\ref{rmk:frobkummer} & \CMPD\ at $(1220,37)$; \CONJ\ beyond (v1.19.0-pre.5; slot-19 sat1220) & the Frobenius-layer rung of a Kummer class is the swapped family's $D=8$ block, with $\T/\ell=\Fl[x]/x^{4}$, $A_1=\Fl[x]/x^{3}$, $I=\m^{3}$, $\T^{\mathrm{sat}}/\T\cong\Z/\ell$; v1.20.0-pre.1: $\theta V_{\mathrm{src}}$ is not the socle (\CMPD, one run); v1.20.0-pre.2: consistent with Prediction~\ref{pr:frobeniuslayer} ``Why (ii)'' (its reduced-source hypothesis fails here); v1.20.0-pre.6: the swapped cells $b=2,3,4$ at $37$ as stated (\CMPD, WP-SW) \\
\ref{sec:ring} (WP-6 paragraph) & definitional (v1.15.1-pre.2); the \texttt{stable} verdict certified by \HA{cor:stablecert} (B1 form) and the digit route beyond $2^{64}$ (v1.15.1-pre.5, slot-15 WP-5); \texttt{HeckeRingLoewyValidation} (\texttt{standard}, 181 checks) & $A_K$, $e_\m$, $\varepsilon_K$ over $\Z/\ell^K$ read mod~$\ell$; \texttt{stable} verdicts certified ($\mathbb{T}'_{\m'}/\ell$ for $\mathbb{T}'=\Z_\ell[T_p:p\in S]$; generation not certified) \\
\ref{sec:localcert} & \PROVED & locality certificate with the deterministic $P_j$ recipe; witness ranges; strict mode (Frobenius radical, Berlekamp subalgebra, witness over $\bar A$, terminates) specified v1.4.0, implemented slot-3 WP-3 \\
\ref{cert:cyclic} & \PROVED, implemented (slot-2 WP-4) & cyclicity $\Rightarrow\T_\m=A_S$; the reported generator is canonical --- the least standard basis vector outside $\m_SV$, \emph{stated} v1.4.1 \\
\ref{cert:comm} & \PROVED\ (statement and socle bound); the cost \VERIFY{} discharged v1.4.0 ($u\le mD$, slot-2 WP-5) & commutant bound; the socle bound $\dim C\ge\dim A_S+fts-\dim\soc A_S$ decides when (b) can succeed \\
\ref{cert:sturm} & \PROVED\ (given $B'(k,\ell)$; range of Claim~\ref{cl:sturmmodl}); implemented on real blocks, slot-3 WP-2 & Sturm sweep with the computed bound \\
\ref{cl:frank} & \PROVED & $\Flr$-ranks from $\Fl$-ranks \\
\ref{cl:split} & \PROVED & scalar extension commutes with all invariants; Jordan types in units of $r$ \\
\ref{cl:canon} & \PROVED & canonical defining polynomial, deterministic \\
\ref{cl:ramanujan} & \PROVED, verified numerically & $\T_\m=\F_{691}[\varepsilon]/\varepsilon^2$; the identities for $q=3,5,7$ and $U_{691}=I+221E_{31}$ hold on the stored matrices \\
\S\ref{sec:tests} Kummer cells & \textbf{refuted 2026-09-16}, corrected (v1.4.0) & ``$r^1=0$ on the Eisenstein class, so $H^1_!=H^1$'' is false on all ten cells: $\rank r^1=1$, $\dim H^1_!=\dim H^1-1$; what holds is $r^1=0$ on the socle $E$ \\
\ref{D6} & \STD\ (socle semantics kept) / sharpened 2026-09-16 & the joint eigenspace on the unsplit $H^1$ is two-dimensional at every Kummer cell; the Eisenstein line is per sign; $\Z/\ell^2$ is a refinement (E8; v1.9.0: computed since slot~8, in the validations and golden only, \S\ref{sec:ring}~(d)) \\
\S\ref{sec:tests} others & \VERIFY & expected outcomes of the remaining test cases \\
\bottomrule
\end{longtable}
\end{center}

%======================================================================
\section{Audit response}\label{app:audit}
%======================================================================
Items are numbered as in \texttt{hecke\_extensions\_audit.md} (\S1 claim table, \S3--\S4
findings, \S7 amendment list). Every numerical claim of the audit that this revision relies
on was recomputed independently (Miller-basis Hecke matrices, lattice indices, $B'(k,\ell)$,
$\tau(691)$, Bernoulli numerators, and the pipeline reruns at $(10,691)$, $(22,5)$, $(26,7)$,
$(32,131)$); discrepancies are noted.
\begin{center}\footnotesize
\begin{longtable}{@{}>{\raggedright\arraybackslash}p{1.6cm}>{\raggedright\arraybackslash}p{1.4cm}>{\raggedright\arraybackslash}p{2.3cm}>{\raggedright\arraybackslash}p{10.2cm}@{}}
\toprule
item & severity & status & where it landed / remarks \\
\midrule
\endhead
B1 Sturm & blocking & \textsc{adopted with change} & Claim~\ref{cl:sturm}(a)/(b), Definition~\ref{def:bprime}, Claims~\ref{cl:bprime}, \ref{cl:sturmmodl}, Certificate~\ref{cert:sturm}, schema. All indices ($3,9,27,400,2240,40,80,36680,45$), the cuspidal index $1$ and the $2$-of-$3$ spans reproduce. \emph{Change}: $B'$ is defined through the $\Fl$-\emph{algebra} generated by the $T_q$, $q\le B'$ prime (what the certificate tests), not the linear span of $T_m$, $m\le B'$; hence $B'(34,131)=2$, and $(32,131)$ breaks the $\Z$-module statement only, while $(22,5)$, $(26,7)$ break the certificate on the pipeline's own $H^1$ (\S\ref{sec:tests}). The audit's remark that $131$ is an Eisenstein-congruence prime for $k=34$ is not adopted: $131\nmid\num(B_{34}/34)$. \\
B2 tangent space & blocking & \textsc{adopted} & Remark~\ref{rmk:coef}, Claim~\ref{cl:resfield} ($\hat w=w^{\ell^N}$ with $f\mid N$ so that $\hat w\equiv w$), Algorithm~\ref{alg:stage3}(2)--(3), \ref{D4}, Certificate~\ref{sec:localcert} (the deterministic $P_j$ recipe; the witness-range remark is stated with the corrected condition $\min(f,f')\le\ell$ for separation, and ``at most $\ell$ prime factors of $f$'' for primitivity, in place of the audit's ``$f\le100$''), Algorithm~\ref{alg:stage2}(2) (bounded attempts, \texttt{"local": false}). \\
M1 boundary & major & \textsc{adopted} & \HA{cl:chi} ($\chi_b$, same multiset on $H^1(\partial)$, $\chi_0$ for $n<\ell$; the $H^0$ and $H^1$ computations re-derived by hand) and \ref{add:bdry} (left-coset representatives $\mathrm{diag}(1,q)$, $(\begin{smallmatrix}q&j\\0&1\end{smallmatrix})$). \\
M2 duality & major & \textsc{adopted} & \HA{cl:pd}, Prediction~\ref{pr:dual}, identity 5 of \S\ref{sec:identities}. \\
M3 Jordan units & major & \textsc{adopted} & Claim~\ref{cl:split} (statement and proof), \ref{D3}. \\
M4 Eisenstein class & major & \textsc{adopted} (option b) & \ref{D6}: the Eisenstein line is $\soc(X^+_{4,\m})$ when one-dimensional, undefined otherwise; Algorithm~\ref{alg:stage5}(4); schema field \texttt{eisenstein}. A mod-$\ell^2$ run is Open Question~18. \\
M5 schema & major & \textsc{adopted} & \S\ref{sec:schema}: \texttt{"691": "active"}, \texttt{"H1\_shriek": "all"}, \texttt{"eisenstein": "bottom layer"}; added \texttt{local}, \texttt{coefficient\_field} ($P_j$), structure-map \texttt{matrices}, \texttt{sturm} ($B'$), \texttt{commit}, \texttt{model}, \texttt{iota\_sign}. \\
Am.~8 $U_{691}$ & minor & \textsc{adopted} & Claim~\ref{cl:ramanujan} and its ``Verified'' paragraph: $c_{691}=542$, $U_{691}=I+221E_{31}$ reproduced on the pipeline in this revision. \\
minor 1 & minor & \textsc{adopted} & \HA{cl:bockstein}: $n>0$. \\
minor 2 & minor & \textsc{adopted} & Claim~\ref{cl:harder} and \S\ref{sec:tests}: $(24,103),(24,2294797)$; both divisibilities recomputed. \\
minor 3 & minor & \textsc{adopted} & \HA{cl:amalgam}: Brown VII.9; $12$-torsion argument for all $n$. \\
minor 4 & minor & \textsc{adopted} & Claim~\ref{cl:bk} (``for $n\gg\ell$''); $(160,691)$ tagged \VERIFY\ in \S\ref{sec:tests} and Appendix~\ref{app:claims}. \\
minor 5 & minor & \textsc{adopted} & Stage 6: parenthetical replaced by the $\De\bmod11$ counterexample. \\
minor 6 & minor & \textsc{adopted} & Certificate~\ref{cert:sturm}: $S_{\mathrm{sep}}(\ell)$ of Claim~\ref{cl:sep} in place of the undefined $S_{\mathrm{Eis}}(\ell)$. \\
minor 7 & minor & \textsc{adopted} & Stage 7: $d\approx n/3$ for all six groups and both signs. \\
minor 8 & minor & \textsc{adopted} & Claim~\ref{cl:gap}: both clauses required for the ``iff''. \\
minor 9 & minor & \textsc{adopted} & Remark~\ref{rmk:coef}, referenced from Claims~\ref{cl:gen}, \ref{cl:yoneda}, \ref{cl:split}. \\
minor 10 & minor & \textsc{adopted} & \S\ref{sec:classdata}: Drozd 1972, Ringel 1975. \\
minor 11 & minor & \textsc{adopted} & Stage 1: $\Theta$ cost plus the fixed cohomology cost. \\
\S2.2 items & --- & \textsc{adopted} & Stage 0: the 7-vertex complex named as the substitute for the $\Ga(2)$ model; the other items were already declared. \\
\S6 concision & --- & \textsc{adopted in part} & \S\ref{sec:symdiv} shortened; the decision box after Claim~\ref{cl:gap} reduced to a pointer; \S\ref{sec:tests} rows at $\ell=691$ merged into one table; the $T_\ell$ paragraph kept in Claim~\ref{cl:Tl} only. \\
OQ13 & --- & \textsc{settled} & $691\,\|\,\num(B_{200}/200)$. The audit's residue $259816$ is the \emph{signed} numerator modulo $691^2$ and ours, $217665$, its absolute value ($B_{200}<0$; $217665+259816=691^2$): the discrepancy is the sign, not a disagreement (v1.12.0; through v1.11.1 this read ``does not reproduce \ldots\ in both normalisations''). The uniqueness at $(198,691)$ is closed (\S\ref{sec:tests}, v1.12.0). \\
\bottomrule
\end{longtable}
\end{center}
No item is contested. The programme's certificates that changed: (c) now needs a recorded
$B'(k,\ell)$ and is otherwise unavailable; the locality certificate and Stage 3 need the
coefficient field $\hat w$; the locality loop has a bounded exit; Open Question~7 is moot.

%======================================================================
\section{OQ2 verdict response and the owner's decisions}\label{app:oq2}
%======================================================================
Open Question~2 --- the Hecke action on the mapping cone --- was reviewed independently; the verdict
is \texttt{docs/implementation/\allowbreak hecke\_extensions\_\allowbreak OQ2\_\allowbreak verdict.md}
(\textsc{confirm with amendments}), and this revision applies it. The binding statement of the
adopted formulation is amendment C-A3 of
\texttt{docs/implementation/\allowbreak hecke\_extensions\_\allowbreak implementation\_\allowbreak plan.md}
\S0.4, which is the mechanism by which [C] gains it; this document states it in \ref{add:cone}. \textbf{The
eight points of the verdict's \S8 were decided by the owner on 2026-09-15}, and v1.2.2 applies
those decisions: the rows below that v1.2.1 labelled \emph{working assumption} are now plain
\textsc{adopted}, with the decision and its date. The owner followed the verdict's recommendation on
every point, with three amendments --- \emph{both} cone cross-routes are built rather than one; the
boundary group is $\Ga_\infty=\langle T\rangle\times\langle-I\rangle$; and the sign convention is
confirmed as already applied. The decision record, one row per point with the site of each decision,
is \S11 of the implementation plan; the two amendments that bear on this document are written into
\ref{add:bdry} and \ref{add:cone}. The status word in the second column records what this document
\emph{states}; where a decision leaves a construction open, the row says so.
\begin{center}\footnotesize
\begin{longtable}{@{}>{\raggedright\arraybackslash}p{2.4cm}>{\raggedright\arraybackslash}p{2.3cm}>{\raggedright\arraybackslash}p{10.8cm}@{}}
\toprule
item & status & where it landed / remarks \\
\midrule
\endhead
formulation (b$^\ast$) & \textsc{adopted} & \ref{add:cone}: the datum is the triple $(\tau,\tau_\partial,\eta)$ over $\Z$, the cochain homotopy $h=\tilde\eta$ is derived, the cone map is $\Phi_g(a,b)=(\Theta_ga,\Theta^\partial_gb+ha)$. Formulations ``on the nose'' and ``strict map of pairs'' are the special case $\eta\equiv0$; the higher-coherence formulation is not needed (the deliverable is cohomology-level). \\
sign & \textsc{confirmed} & $dh+hd=\Theta^\partial_gr-r\Theta_g$ with $D(a,b)=(-da,db+ra)$, re-derived independently: $D\Phi-\Phi D=(0,(dh+hd+r\Theta-\Theta^\partial r)a)$. The plan audit's opposite sign was wrong. \\
$\Hom(\ker r^1,\coker r^0)$ & \textsc{new, stated} & \ref{add:cone} third bullet and Prediction~\ref{pr:cone}: real for a free cochain-level $h$ (\emph{every} element occurs over a field, \PROVED), eliminated by the $\eta$-construction (\PROVED\ over $\Z$ for all $\ell$). This closes the gap that v1.1 did not mention. \\
the entrywise test & \textsc{demoted} & \ref{add:cone} fourth bullet: $r\Theta_g=\Theta^\partial_gr$ is a recorded diagnostic, never a branch; $h=0$ is admissible iff $\tau j=j\tau_\partial$, which the shipped $\tau$ fails. \\
boundary datum & \textsc{adopted} (owner decision 2026-09-15, as recommended) & The normative object is the $\Ga$-level $\tau_\partial$ on $\partial X$ with the bulk's own $s_i,\delta_i$; the one-dimensional $\tau_\infty$ of \ref{add:bdry} is its Shapiro transcription and a validation route. \ref{add:bdry}'s formulas are unchanged and remain the oracle. \\
integrality & \textsc{adopted} (owner decision 2026-09-15, as recommended) & \ref{add:cone}: $\tau_\partial$ and $\eta$ are integral for every $\ell$; the bulk $\tau$ is integral exactly when no fill used the mod-$\ell$ fallback. When all three tables are integral, the matrices are their reductions, so the reduction maps and Bockstein connecting maps are equivariant on the nose. \emph{Policy}: the mod-$m$ fill fallback is \textbf{forbidden by default} --- the integer tier is widened instead; it stays reachable behind a setting and then records a non-integral cell, but that is not the default escape, and a record claiming integrality must have solved all three tables over $\Z$. \\
sign convention & \textsc{adopted} (owner decision 2026-09-15, confirming what was applied) & \ref{add:cone} second bullet: $\partial\eta+\eta\partial=j\tau_\partial-\tau j$, hence $h=\tilde\eta$ with no sign. The cochain identity and the cone map are fixed by $D$ and are unmoved by the choice. \\
cross-routes and the PD gate & \textsc{adopted with amendment} (owner decision 2026-09-15) & Poincar\'e adjointness (\HA{cl:pd}, once $D_n$ is computed) is the \emph{acceptance gate} for $X_3$ at $\ell\ge5$; and \textbf{both} cross-routes --- the strict-pair $\tau^s$ and the relative complex --- are built as validation arms, where the verdict recommended one of the two. They are cross-checks of the independence theorem, never definitions. \\
boundary group at $\ell=2$ & \textsc{adopted with amendment} (owner decision 2026-09-15; the verdict made no recommendation) & $\Ga_\infty=\langle T\rangle\times\langle-I\rangle$, i.e.\ \HA{sec:six}'s own definition $H^\bullet(\partial;M):=H^\bullet(\Ga_\infty;M)$; the circle model of \ref{add:bdry} is its specialisation at $2\in R^\times$. Consequence, stated in \ref{add:bdry} and \ref{add:cone}: at $\ell=2$ the boundary --- and with it the bulk and the cone --- is \emph{unbounded above in degree}, so \HA{eq:les} is a window of a non-terminating long exact sequence. The degree bound / truncation policy that needs was settled in the contract: [C] \HA{open:ell2}, decided by the owner on 2026-09-16 as \HA{oc1a} ($q_{\max}$ a recorded run parameter on the free $\Ga$-model; this row said ``open \ldots\ not decided'' until v1.3.3). \\
golden tiers & \textsc{adopted} (owner decision 2026-09-15, as recommended) & $\dim X_3$, $\dim X_6$ and the char-poly / rank / trace of $A_c^{(q)}$ in golden tier~A (canonical for $\ell\ge5$); $\Theta^\partial_{\mathrm{flat}}$, $h_{\mathrm{flat}}$, the cone bases and the three table hashes in tier~B. The cone bases sit in tier~B because they are stored-basis dependent, \emph{not} because a change of $\eta$ or $\tau_\partial$ moves them --- it does not. \\
LES in all degrees & \textsc{adopted}; restated in [C] & v1.2 stated the long exact sequence in all degrees and named $H^0_c=0$ and $H^2(\Ga;M)=0$ as the hypotheses under which it collapses to the six-term sequence. \HA{cl:les} keeps the sequence and the collapse but not that pair of hypotheses: $H^0_c=0$ is unconditional (\HA{lem:conerelative}(4)) and is dropped, while $H^2(\Ga;M)=0$ is kept and shown not removable (\HA{lem:h2torsion}(2)). See item~3 of \S\ref{sec:imported}. \\
torsion primes & \textsc{adopted} (owner decision 2026-09-15, as recommended) & \ref{add:cone}, last paragraph: the formulation transfers verbatim to a free $\Ga$-model of the pair (Borel construction), which is the orbifold cohomology at $\ell\in\{2,3\}$; the 7-vertex complex is its specialisation where $6$ is invertible. Moerdijk--Pronk's complex is a \emph{computational model} of it: their charts are effective and the independence proofs for their $\mu$ are omitted, so the citation for the general case is \CONJ\ and the safe normative definition is the Borel construction they themselves name. The general model enters the binding contract \emph{now, as definitions}, with the construction deferred. \\
\bottomrule
\end{longtable}
\end{center}

%======================================================================
\begingroup\small
\begin{thebibliography}{99}
\bibitem{ash_stevens} A.~Ash, G.~Stevens, \emph{Cohomology of arithmetic groups and congruences between systems of Hecke eigenvalues}, J.~reine angew.\ Math.\ 365 (1986), 192--220.
\bibitem{bc} J.~Bella\"iche, G.~Chenevier, \emph{Families of Galois representations and Selmer groups}, Ast\'erisque 324 (2009).
\bibitem{bk} J.~Bella\"iche, C.~Khare, \emph{Level 1 Hecke algebras of modular forms modulo $p$}, Compositio Math.\ 151 (2015), 397--415.
\bibitem{eigenbook} J.~Bella\"iche, \emph{The Eigenbook: Eigenvarieties, families of Galois representations, $p$-adic $L$-functions}, Birkh\"auser, 2021.
\bibitem{bass_ubiquity} H.~Bass, \emph{On the ubiquity of Gorenstein rings}, Math.\ Z.\ 82 (1963), 8--28 (v1.19.0-pre.1; \S6 (ideals over Gorenstein orders) and \S7
(torsion-free modules over orders whose ideals are $2$-generated) used in Claim~\ref{cl:258lattice}; \textbf{not read}, locators pending, rule~A2; not in the project library, slot-19 referee F4, m1 --- the claim's route does not need it).
\bibitem{benson} D.~J.~Benson, \emph{Representations and Cohomology I}, Cambridge Studies in Adv.\ Math.\ 30, CUP, 1991.
\bibitem{berger_li_zhu} L.~Berger, H.~Li, H.~J.~Zhu, \emph{Construction of some families of $2$-dimensional crystalline representations}, Math.\ Ann.\ 329 (2004),
365--377 (v1.18.0-pre.2; Thm.~4.5, p.~12, and Prop.~4.4(1), p.~11, read on the arXiv copy, slot-17 R17-4; the $\ell=37$ report's V4).
\bibitem{brown} K.~S.~Brown, \emph{Cohomology of Groups}, GTM 87, Springer, 1982.
\bibitem{citro_ghitza} C.~Citro, A.~Ghitza, \emph{Computing level one Hecke eigensystems (mod $p$)}, LMS J.\ Comput.\ Math.\ 16 (2013), 246--270; arXiv:1102.3321 (v2 read; journal data not checked against it).
\bibitem{coates_wiles} J.~Coates, A.~Wiles, \emph{On the conjecture of Birch and Swinnerton-Dyer}, Invent.\ Math.\ 39 (1977), 223--251 (v1.19.0-pre.3; p.~230, (8),
the homomorphisms $\varphi_k\colon U_0\to\mathcal O_p/p$, read on the page image by the slot-19 referee F4).
\bibitem{deligne_serre} P.~Deligne, J.-P.~Serre, \emph{Formes modulaires de poids 1}, Ann.\ Sci.\ ENS 7 (1974), 507--530.
\bibitem{deng} T.~Deng, \emph{Torsions in cohomology of $\SL_2(\Z)$ and congruence of modular forms}, J.~Number Theory 246 (2023).
\bibitem{deo} S.~V.~Deo, \emph{Structure of Hecke algebras of modular forms modulo $p$}, Algebra \& Number Theory 11 (2017), 1--38.
\bibitem{diamond_im} F.~Diamond, J.~Im, \emph{Modular forms and modular curves}, in: Seminar on Fermat's Last Theorem (Toronto, 1993--1994), CMS Conf.\ Proc.\ 17, AMS, 1995, 39--133 (v1.18.0-pre.1; Thm.~12.2.2, p.~108, and Prop.~12.4.10, p.~117, verified on page images of the library copy, slot-17 L17-5 repairs).
\bibitem{drozd} Ju.~A.~Drozd, \emph{Representations of commutative algebras}, Funct.\ Anal.\ Appl.\ 6 (1972), 286--288.
\bibitem{edixhoven} B.~Edixhoven, \emph{The weight in Serre's conjectures on modular forms}, Invent.\ Math.\ 109 (1992), 563--594 (v1.18.0-pre.1; Thm.~2.6 read from the text layer only; v1.18.0-pre.2: statement and attribution verified in a secondary source, Edixhoven's Boston
survey \emph{Serre's conjecture}, printed p.~19, slot-18 R18-4; the theorem number \textbf{not verified}, locator pending, rule~A2).
\bibitem{fukaya_kato} T.~Fukaya, K.~Kato, \emph{On conjectures of Sharifi}, Kyoto J.\ Math.\ 64 (2024), no.~2, 277--420; preprint 2012 (v1.21.0-pre.9: recalled; v1.21.0-pre.10, P5$'$ audit: journal data from the Project Euclid listing,
the Introduction (Thms.~0.14--0.16, 0.20, Rmk.~0.17) read on the page of the authors' preprint of 15~Aug~2012, \texttt{math.ucla.edu/\string~sharifi/sharificonj.pdf}; not held, the journal copy not read).
\bibitem{glover} D.~J.~Glover, \emph{A study of certain modular representations}, J.~Algebra 51 (1978), 425--475.
\bibitem{hho} W.~Hart, D.~Harvey, W.~Ong, \emph{Irregular primes to two billion}, Math.\ Comp.\ 86 (2017), no.~308, 3031--3049 (v1.16.3; data as verified in \texttt{docs/latex/loewy\_and\_galois\_refs.bib}).
\bibitem{harder3} G.~Harder, \emph{Cohomology of Arithmetic Groups}, Vol.~III, manuscript, 2022 (\texttt{docs/literature/}).
\bibitem{hida_cong} H.~Hida, \emph{Congruences of cusp forms and special values of their zeta functions}, Invent.\ Math.\ 63 (1981), 225--261 (v1.19.0-pre.1:
cited for the shape of Hida's congruence criterion only; v1.19.0-pre.5: Thm.~5.1, pp.~247--248, (3.2), p.~237, Thm.~2.5 and (2.8), p.~236, Thm.~6.1, p.~250, read on the page by
slot-19 R19-9; PDF page $=$ journal page $-224$).
\bibitem{hida_ens} H.~Hida, \emph{Iwasawa modules attached to congruences of cusp forms}, Ann.\ Sci.\ ENS (4) 19 (1986), 231--273 (v1.19.0-pre.1; Thm.~3.1 and
Cor.~3.2, p.~250, verified for [U] and [C]; their nebentypus form at weights outside the class, used in Claim~\ref{cl:lvtwist} --- Cor.~3.2 (3.3b) for every $k\ge2$ with
p.~249's $S_k(\Ga_1(N_0p),a;A)=S_k(\Phi,\omega^{a-k};A)$ --- verified on the page by the slot-19 referee F1 (v1.19.0-pre.2)).
\bibitem{jochnowitz} N.~Jochnowitz, \emph{Congruences between systems of eigenvalues of modular forms}, Trans.\ AMS 270 (1982), 269--285.
\bibitem{lubeck} F.~L\"ubeck, \emph{Conway polynomials for finite fields}, online tables; definition in D.~Holt, B.~Eick, E.~O'Brien, \emph{Handbook of Computational Group Theory}, \S7.5.
\bibitem{mazur} B.~Mazur, \emph{Modular curves and the Eisenstein ideal}, Publ.\ Math.\ IHES 47 (1977), 33--186.
\bibitem{mtt} B.~Mazur, J.~Tate, J.~Teitelbaum, \emph{On $p$-adic analogues of the conjectures of Birch and Swinnerton-Dyer}, Invent.\ Math.\ 84 (1986), 1--48
(v1.19.0-pre.1; \S8, the modular-symbol formula for $L(f,\chi,j+1)$, shape only; constants not checked, rule~A2).
\bibitem{mccallum_sharifi} W.~McCallum, R.~Sharifi, \emph{A cup product in the Galois cohomology of number fields}, Duke Math.\ J.\ 120 (2003), no.~2, 269--310 (v1.21.0-pre.9: page range recalled; v1.21.0-pre.10: confirmed on the Project Euclid listing;
\S5, the relations and Thm.~5.1, and the remark after Cor.~5.2 read on the page of the authors' web PDF \texttt{math.ucla.edu/\string~sharifi/pairing.pdf}, version of 5~Dec~2003, pp.~12--14; not in the project library).
\bibitem{merel} L.~Merel, \emph{Universal Fourier expansions of modular forms}, in: On Artin's conjecture for odd 2-dimensional representations, LNM 1585, Springer, 1994, 59--94.
\bibitem{ribet} K.~Ribet, \emph{A modular construction of unramified $p$-extensions of $\Q(\mu_p)$}, Invent.\ Math.\ 34 (1976), 151--162.
\bibitem{ringel} C.~M.~Ringel, \emph{The representation type of local algebras}, in: Representations of Algebras (Ottawa 1974), LNM 488, Springer, 1975, 282--305.
\bibitem{serre} J.-P.~Serre, \emph{Congruences et formes modulaires (d'apr\`es H.~P.~F.~Swinnerton-Dyer)}, S\'em.\ Bourbaki 416 (1971/72); and \emph{Formes modulaires et fonctions z\^eta $p$-adiques}, LNM 350 (1973).
\bibitem{sharifi_reciprocity} R.~Sharifi, \emph{A reciprocity map and the two-variable $p$-adic $L$-function}, Ann.\ of Math.\ (2) 173 (2011), 251--300 (v1.21.0-pre.10, P5$'$ audit: the Introduction, (1.1), and \S5 (the map $\varpi$, Prop.~5.7, Conj.~5.8)
read on the page of the arXiv version \texttt{arXiv:0709.3591}; the journal copy not read, not held).
\bibitem{sharifi_eisenstein} R.~Sharifi, \emph{Iwasawa theory and the Eisenstein ideal}, Duke Math.\ J.\ 137 (2007), no.~1, 63--101 (v1.21.0-pre.9: page range recalled; v1.21.0-pre.10: confirmed on the Project Euclid listing; Introduction, p.~4,
read on the page of the author's web PDF \texttt{math.ucla.edu/\string~sharifi/paireis.pdf}; not in the project library).
\bibitem{shimura} G.~Shimura, \emph{Introduction to the Arithmetic Theory of Automorphic Functions}, Princeton UP, 1971.
\bibitem{stein} W.~Stein, \emph{Modular Forms, a Computational Approach}, GSM 79, AMS, 2007.
\bibitem{sturm} J.~Sturm, \emph{On the congruence of modular forms}, in: Number Theory (New York 1984--85), LNM 1240, Springer, 1987, 275--280.
\bibitem{washington} L.~C.~Washington, \emph{Introduction to Cyclotomic Fields}, 2nd ed., GTM 83, Springer, 1997 (table of irregular pairs).
\bibitem{washington1st} L.~C.~Washington, \emph{Introduction to Cyclotomic Fields}, 1st ed., GTM 83, Springer, 1982 (v1.18.0-pre.1; Cor.~5.15, p.~61, verified on page images, slot-17 R17-3;
v1.19.0-pre.3: pp.~63--65 (Thm.~5.18 and its proof), Prop.~8.12, p.~156, read by slot-19 R19-4; \S13.7, pp.~299--300, read by the slot-19 referee F4).
\bibitem{wiles_ordinary} A.~Wiles, \emph{On ordinary $\lambda$-adic representations associated to modular forms}, Invent.\ Math.\ 94 (1988), 529--573 (v1.18.0-pre.1; Thm.~2.1.4, p.~561, verified on the page image, slot-16 L16-3).
\end{thebibliography}
\endgroup

%======================================================================
% Revision history (1.25.0-pre.2): the date field and the abstract of 1.25.0-pre.1, moved here verbatim.
%======================================================================
\raggedright
\section{Revision history}\label{app:history}
This document follows [C]'s practice: a published version's text never changes under its own number. This appendix is the
record of the versions. Its first subsection is the release record, one paragraph per version, newest first (up to
1.25.0-pre.1 it was the date field of the title page); the second quotes the title page of 1.25.0-pre.1.

\subsection{Release record}
Version 1.25.3 --- 2026-11-10 (hecke-extensions slot~26, \textbf{signed by the owner 2026-11-10}: the owner-signed release of 1.25.3-pre.1; no change of content --- header, date and this record only; no mathematical statement, label, tag, schema key, record key, golden line or digest changes. Companion pointers: [C] v1.26.4, [U] 0.7.3.)

Version 1.25.3-pre.1 --- 2026-11-10 (hecke-extensions slot~26, a \textbf{pre-release, not signed}: title page only, no change of mathematical content: the date line of the title page names only the version, the date and this appendix; the words ``signed by the owner 2026-11-09'' leave the title page, and the sign-off record of every version stays in this record. No mathematical statement, label, tag, schema key, record key, golden line or digest changes. Companion pointers: [C] v1.26.4, [U] 0.7.3 (the same title-page change; signed with 1.25.3).)

Version 1.25.2 --- 2026-11-09 (hecke-extensions slot~26, \textbf{signed by the owner 2026-11-09}: the owner-signed release of 1.25.2-pre.1; no change of content --- header, date and this record only. Companion pointers: [C] 1.26.2, [U] 0.7.1.)

Version 1.25.2-pre.1 --- 2026-11-09 (hecke-extensions slot~26, a \textbf{pre-release, not signed}: paths only, no change of content (repository layout, owner ruling 2026-11-09): the program's sources moved under \texttt{src/}; one line names \texttt{src/ModArithmetic/ENGINE\_API.md}. No mathematical statement, label, tag, schema key, record key, golden line or digest changes. Companion pointers: [C] v1.26.2-pre.1, [U] 0.7.1.)

Version 1.25.1 --- 2026-11-09 (hecke-extensions slot~26, \textbf{signed by the owner 2026-11-09}: the owner-signed release of 1.25.1-pre.1; no change of content --- header, date and this record only; no mathematical statement, label, tag, schema key, record key, golden line or digest changes. Companion pointers: [C] v1.26.1, [U] 0.7.1.)

Version 1.25.1-pre.1 --- 2026-11-09 (slot~26, a \textbf{pre-release, not signed}: no change of mathematical content: machine names and private paths removed for the public release (owner ruling 2026-11-09).)

Version 1.25.0 --- 2026-11-07 (hecke-extensions slot~25, \textbf{signed by the owner 2026-11-07}: the owner-signed release of 1.25.0-pre.2; no change of content --- header, date and this record only; no mathematical statement, label, tag, schema key, record key, golden line or digest changes. Companion pointers: [C] v1.25.0, [U] 0.7.0.)

Version 1.25.0-pre.2 --- 2026-11-07 (\textbf{a pre-release, not signed, not yet seen by the owner}: front matter, section titles and layout only. The title page has a one-line date, the uniform author block and a generic abstract; the date field of 1.25.0-pre.1, one entry per version from 1.25.0-pre.1 back to v1.0, is moved here verbatim as the release record below (every entry boundary became a paragraph break; whitespace only: in 1.25.0-pre.1 the first nine entries were separated by line breaks and the rest ran on), and the earlier author line and abstract, whose paragraphs narrated versions 1.1 to 1.4.0, are quoted in the last subsection. The version tags are removed from eleven section titles; their earlier titles were ``The dimension of a reducible block at every \texorpdfstring{$\ell$}{l} (v1.18.0-pre.1)'' (\S\ref{sec:dimblock}), ``The block \texorpdfstring{$(258,37)$}{(258,37)} exactly: the Hecke algebra, its over-orders, the lattice (v1.19.0-pre.1)'' (\S\ref{sec:block258}), ``The Dickson--\texorpdfstring{$\theta$}{theta} objects (v1.12.0-pre.1)'' (\S\ref{sec:thetaobj}), ``The coefficient ring and the E8 objects (v1.9.0-pre.1; v1.9.0; v1.10.0-pre.1; v1.10.0-pre.2; v1.10.0; v1.10.1; v1.11.0-pre.1; v1.11.0-pre.2; v1.11.0; v1.12.0-pre.1; v1.12.0-pre.2; v1.12.0; v1.12.1)'' (\S\ref{sec:ring}), ``The congruence-census catalogue \texttt{hecke\_census/v1} (v1.12.0-pre.1)'' (\S\ref{sec:census}), ``The \texorpdfstring{$L$}{L}-value interface, and its Eisenstein instance (v1.19.0-pre.1)'' (\S\ref{sec:lvalue}), ``Measures: the measure-valued \textsf{linear} evaluation (v1.20.0-pre.1)'' (\S\ref{sec:lvmeasure}), ``The depth-2 triangle, the change of lattice, and \texorpdfstring{$\iota$}{iota} over \texorpdfstring{$\Z/\ell^{k}$}{Z/l\^{}k} (v1.20.0-pre.1)'' (\S\ref{sec:lvdepth2}), ``The truncated one-variable \texorpdfstring{$\ell$}{l}-adic \texorpdfstring{$L$}{L}-function as a record (v1.20.0-pre.5)'' (\S\ref{sec:lvlfunction}), ``Coefficient rings, the two-variable family, the Eisenstein measure, wild characters and Iwasawa invariants (v1.21.0-pre.1, pre.3, pre.7)'' (\S\ref{sec:lvslot21}) and ``OQ2 verdict response (v1.2) and the owner's decisions (v1.2.2)'' (Appendix~\ref{app:oq2}). No statement, label, number or status tag changes; the provenance notes at the start of these sections are unchanged. Companion pointers: [C] v1.25.0-pre.7, [U] 0.7.0-pre.4.)

Version 1.25.0-pre.1 --- 2026-11-06 (hecke-extensions slot~25; \textbf{a pre-release, not signed}; on top of the signed 1.24.0; \textbf{owner ruling of 2026-11-06} on the verification \texttt{docs/implementation/slot25/R59\_verification.md} of the gallery referee's finding R1, \texttt{docs/implementation/slot25/PG140\_referee.md}). Prediction~\ref{pr:latticesplit}~(R) \textbf{restated a fifth time} (a statement change of a \CONJ, the tag kept): on irreducible classes the hypothesis ``not dihedral'' is replaced by ``\emph{of one weight}'' --- exactly one pair $(w,j)$ of the congruence-free test with a non-zero generalised eigenspace in $M_w(\Fl)$; reducible classes unchanged (congruence-free alone). Of one weight implies not dihedral (proved in the statement), and excludes also the non-ordinary and the companion classes. New paragraph ``Refuted on cohomology: the v1.22.0-pre.1 form, at a non-ordinary class'': the class of $f_{16}$ at $\ell=59$ (irreducible, not dihedral, congruence-free, $a_{59}\equiv0$, Serre weights $16$ and $46$), four census-4.0 rows $n=884,942,944,1000$ with $(T_2-a_2)^2\ne0$ on $V$ ($\dim V=6,8,6,8$), and the consistency of the new form with the $672$ Loewy-$\ge3$ rows and the $5805$ irreducible rows with $\dim V\ge6$ of census~4.0; the last sentence of the dihedral paragraph gains a pointer; the row of Appendix~A updated. No label added, renamed or removed; no computed value, schema key, record key, golden line, digest, validation or code path moves. \textbf{Not yet seen by the owner as text.} Companion pointers: [C] v1.25.0-pre.7, [U] 0.7.0-pre.3.)

Version 1.24.0 --- 2026-11-05 (hecke-extensions slot~25, \textbf{signed by the owner 2026-11-05}: the owner-signed release of 1.24.0-rc.2; no change of content --- header and date only; no mathematical statement, label, tag, schema key, record key, golden line or digest changes. Companion pointers: [C] v1.24.0, [U] 0.6.0.)

Version 1.24.0-rc.2 --- 2026-11-05 (hecke-extensions slot~25, RC124 rc.2; \textbf{the release candidate of 1.24.0: not signed}; rc.1 refereed ACCEPT\_WITH\_MINORS, \texttt{docs/implementation/slot25/RC124\_referee.md}): Prediction~\ref{pr:latticesplit}, item~2 of ``Refuted as stated'' --- ``every saturation defect measured so far has length $\le1$'' is scoped to the Kummer classes, with a pointer to the longer defects elsewhere (length $2$ at $(540,23)$ in the same item; $\delta$ up to $190$ in census~4.0); the referee's F5, \textbf{not yet seen by the owner}; the sentence supports an expectation, not a claim, so no tag or label moves. Pre.2 of Claim~\ref{cl:sturm} (rc.1 below) was re-audited after A7 by the HeckeZ re-audit A7b (\texttt{heckez/A7b-Z0-reaudit.md} \S3, ACCEPT; on the branch \texttt{s24-wip-heckez-z0}). Nothing else changes. Companion pointers: [C] v1.24.0-rc.2, [U] 0.6.0.

Version 1.24.0-rc.1 --- 2026-11-05 (hecke-extensions slot~25, RC124; \textbf{the release candidate of 1.24.0: not signed, not yet refereed}, referee \texttt{docs/implementation/slot25/RC124\_referee.md}): the text of Claim~\ref{cl:sturm}(a), its status line, Definition~\ref{def:bprime}'s parenthesis and the claims-table row exactly as in 1.24.0-pre.2 below, taken onto the main line from the branch \texttt{s24-wip-heckez-z0} (\texttt{hecke\_extensions.tex} only; the erratum of pre.1 accepted by the owner, HeckeZ question Q9, 2026-11-04; pre.2's two corrections, after the HeckeZ referee A7 --- whose record \texttt{heckez/A7-Z0-referee.md} is on that branch, not on the main line ---, not yet seen by the owner). Nothing else changes: no other statement, label, tag, schema key, record key, golden line, digest, validation or code path. Companion pointers: [C] v1.24.0-rc.1, [U] 0.6.0.

Version 1.24.0-pre.2 --- 2026-11-04 (hecke-extensions slot~24, HeckeZ Z0 fix round after the adversarial referee A7, \texttt{docs/implementation/slot24/heckez/A7-Z0-referee.md}; \textbf{a pre-release, not signed}): Claim~\ref{cl:sturm}(a), second sentence --- for $k\equiv2\bmod12$, where $B(k)=d+1$, the condition equivalent to ``$B=B(k)$ suffices'' is $(a_d(f_0),a_{d+1}(f_0))\ne(0,0)$; pre.1's ``$a_{d+1}(f_0)\ne0$'' is sufficient, not equivalent (A7-m6); the verified range is extended to even $4\le k\le70$ (A7-N8; both coefficients are non-zero there). The status (\CONJ) and every consequence are unchanged; no computed value, record, golden line, digest, validation or code path moves. Companion pointers: [C] v1.24.0-pre.1 (main line; the HeckeZ text is [C] v1.25.0-pre.0b on the branch \texttt{s24-wip-heckez-z0}), [U] 0.6.0.

Version 1.24.0-pre.1 --- 2026-11-04 (hecke-extensions slot~24, HeckeZ package Z0; \textbf{a pre-release, not signed}, on top of the signed 1.23.0): the erratum to Claim~\ref{cl:sturm}(a), second sentence, accepted by the owner (HeckeZ plan, question Q9, signed 2026-11-04; \texttt{docs/implementation/slot24/heckez/ERRATUM\_E\_clsturm.md}). The argument via $\De g$ assumed $a_0(g)=0$ and is withdrawn; $\T_\Z(M_k)\otimes\Q$ is spanned by $T_m$, $m\le B$, for some finite $B$ (proved), and $B=B(k)$ is \CONJ, verified for even $4\le k\le58$. Consequential edits: the claim's status line, Definition~\ref{def:bprime}'s parenthesis, the claims table row. A correction of a proof: no computed value, record, golden line, digest, validation or code path depends on the sentence. Companion pointers: [C] v1.24.0-pre.1 (main line; the HeckeZ text is [C] v1.25.0-pre.0 on the branch \texttt{s24-wip-heckez-z0}), [U] 0.6.0.

Version 1.23.0 --- 2026-11-03 (hecke-extensions slot~24, \textbf{signed by the owner 2026-11-03}: the owner-signed release of 1.23.0-pre.1; no change of content --- header and date only; no mathematical statement, label, tag, schema key, record key, golden line or digest changes. Companion pointers: [C] v1.23.0, [U] 0.6.0.)

Version 1.23.0-pre.1 --- 2026-10-26 (hecke-extensions slot~24, D6 code audit finding M2; \textbf{a pre-release, not signed}). \S\ref{sec:ring}~(a) ``Reserved values'' and the campaign paragraph of the as-built list now agree with \HA{sec:recordschema} item~7 and \HA{def:ringsemantics}(3): \texttt{Z} is no longer refused inside \texttt{MathSettings::set()}; it is stored and honoured by the integral descent only, every other consumer refuses it by name (\texttt{NotImplementedInThisBuild}, exit~$5$); \texttt{Z/<m>} stays reserved. Wording only: no mathematical statement, label, tag, schema key, record key, golden line or digest changes. Companion pointers: [C] v1.23.0-pre.37.

Version 1.22.0 --- 2026-10-12 (hecke-extensions slot~22, CLOSE-T: the \textbf{release candidate} from 1.22.0-pre.3; \textbf{owner sign-off pending}, not yet signed). The text of 1.22.0-pre.1, pre.2 and pre.3 as stated in their entries (referee records: \texttt{F33-2\_referee.md}, \texttt{P5pp\_referee.md}, \texttt{F-T1\_release\_referee.md}); the \texttt{pre.N} suffix dropped from the current version only; \textbf{no statement's mathematics, label, tag or status changes}; no schema key, record key, golden line, digest or census verdict changes. Companion pointers: [C] v1.22.0, [U] 0.6.0, [L37] 0.5.2, [RC] 0.2.0.

Version 1.22.0-pre.3 --- 2026-10-12 (hecke-extensions slot~22; \textbf{a pre-release, not signed}; CLOSE-T T1 repairs of the Fable release referee \texttt{docs/implementation/slot22/F-T1\_release\_referee.md}). MINOR-2: the cell count of Remark~\ref{rmk:lviwasawa}(9), the status row of Appendix~A and this date line corrected from $32$ to $30$ cells (``PROVED modulo STD inputs, given CMPD L3 at the pair and CMPD (H1) at the cell''), $+2$, $(66,37)$ and $(102,37)$, once (H1) is gated there (the referee's one-route computation), as the P5$''$ referee's \S8 states. NIT-11 (the parity clause) and NIT-12 (the $(540,23)$ rows) applied as wording; NIT-13 (a \texttt{.bib} entry for Fukaya--Kato--Sharifi) not applied. No statement, label, tag or value moves. Companion pointers: [C] v1.22.0-pre.17, [U] 0.6.0-pre.3, [L37] 0.5.1, [RC] 0.1.1-pre.2.

Version 1.22.0-pre.2 --- 2026-10-09 (hecke-extensions slot~22; \textbf{a pre-release, not signed}; slot-22 P5$''$, \texttt{docs/implementation/slot22/P5pp\_note.md} v0.2, and its Fable referee
\texttt{docs/implementation/slot22/P5pp\_referee.md}: ACCEPT WITH REPAIRS, $0$\,B\,/\,$2$\,M\,/\,$7$\,m, the repairs R1--R8 applied in the note; only the tag changes of the referee's \S8).
Remark~\ref{rmk:lviwasawa}(9) gains the paragraph ``The modular-symbol proof'': at the $30$ cells with the project's own readings of (H1) (and, once (H1) is gated there, at $(66,37)$ and $(102,37)$, so far the referee's one-route computation)
the forced-zero theorem and (S) are \PROVED{} modulo \STD{} inputs, given \CMPD{} L3 at the pair and \CMPD{} (H1) at the cell --- through pre.1 conditional on Sharifi's conjecture; per pair, nothing uniform in $(\ell,r)$;
``consistent with Sharifi's conjecture (1.1)'' kept, (1.1), the index map and the descent untouched; the open list updated. Remark~\ref{rmk:lvmeasplus}(iii): ``conditional on Sharifi's conjecture'' $\to$ ``unconditional at the
computed cells''. Remark~\ref{rmk:lvlegc}: (L1) at depth exactly~$1$ at $(66,37)$, $(102,37)$, the referee's Python, one route (the C++ gate has not read them). No label added, renamed or removed; no schema key,
record key, golden line, digest or census verdict changes. Companion pointers: [C] v1.22.0-pre.2, [U] 0.6.0-pre.2, [L37] 0.5.1, [RC] 0.1.1-pre.2.

Version 1.22.0-pre.1 --- 2026-10-09 (hecke-extensions slot~22; \textbf{a pre-release, not signed}; census~3.3 finding F33-2, the owner ruling of 2026-10-09,
\texttt{docs/implementation/slot22/owner\_ruling\_F33-2\_1009.md}, with the lead note; wording of the slot-22 Fable referee \texttt{docs/implementation/slot22/F33-2\_referee.md}~\S3.3).
Prediction~\ref{pr:latticesplit}~(R) \textbf{restated a fourth time} (a statement change of a \CONJ, the tag kept): the hypothesis is now ``congruence-free and \emph{not dihedral}'',
dihedral meaning $\bar\rho_\m$ irreducible and $\bar\rho_\m\cong\bar\rho_\m\otimes\omega^{(\ell-1)/2}$ (reducible self-twist classes stay inside). New paragraph ``Refuted on cohomology: the
v1.17.0-pre.4 form, at dihedral classes'': the $161$ Loewy-$3$ rows of census~3.3 at $\ell=23$ (twists of $\Delta$) and $\ell=31$ (the weight-$16$ form), $A\cong\Fl[x]/x^3$, a computed
counterexample to the v1.17.0-pre.4 form; the narrowing removes $929$ dihedral rows ($161$ violating, $768$ consistent); on the $100\,025$ non-Dickson, non-dihedral residue-degree-$1$ rows no
row has Loewy length~$3$. The evidence sentence ``No row in the table violates the restated form'' gains the census-3.3 sentence; the prediction's header note and Appendix~A's row updated.
No label added, renamed or removed; no schema key, record key, golden line, digest or census verdict changes (the $161$ rows are kept unchanged). Companion pointers: [C] v1.22.0-pre.2, [U] 0.6.0-pre.2, [L37] 0.5.1, [RC] 0.1.1-pre.1.

Version 1.21.0 --- 2026-10-06 (hecke-extensions slot~21, the close; \textbf{the release} from 1.21.0-pre.10 under the \textbf{owner's decision D16}
(release of the pre-release documents at the close, accepted as recommended, 2026-10-06; \texttt{docs/implementation/slot21/slot21\_status\_and\_close\_plan.md}~\S5)).
Slot-21 release: the text of 1.21.0-pre.10 as published, the \texttt{pre.N} suffix dropped from the current version only; \textbf{no statement's mathematics, label, tag or status changes}.
The release upgrades no referee status: the P2x\,/\,P4x\,/\,P5\,/\,P5$'$ additions (pre.6, pre.7, pre.9) stay as stated in their entries below (F4 of pre.8 on P5, P2x, P4x; the adversarial audit of P5$'$ at pre.10;
whatever is marked not refereed, \STD{} recalled or conditional stays so). No schema key, record key, golden line, digest or census verdict changes.
Companion pointers: [C] v1.21.0, [U] 0.5.0, [L37] 0.5.1, [RC] 0.1.0, [LG] 0.1.5, [EX] 1.4.1, [CLI] 1.3.0.

Version 1.21.0-pre.10 --- 2026-10-06 (hecke-extensions slot~21; \textbf{a pre-release, not signed}): the adversarial audit of P5$'$ (\texttt{docs/implementation/slot21/rc/P5prime\_audit\_report.md}: ACCEPT WITH REPAIRS),
its repairs to Remark~\ref{rmk:lviwasawa}(9): the conditional theorem's conclusion is $h(b)=0$, hence $\lambda^{+}_b\ge1$ \emph{wherever $\mu^{+}_b=d$} (the hypothesis was dropped in the statement); the forced zero is $e_{\ell-r,r}=0$ of the
cup-product pairing (McCallum--Sharifi under Vandiver, \cite{hho}), not of the computed vector; Sharifi's conjecture cited at its statement \cite{sharifi_reciprocity}~(1.1) (read on the page), the index $t\leftrightarrow\omega^{t-1}$ Sharifi's; Fukaya--Kato's
theorems and hypotheses read (Introduction of the preprint) and the remaining gap named (descent to level one in this normalisation); ``nullity $0$ at regular $r\pm2$'' restricted to the pairs tried ($(257,162)$: nullity $10$); no novelty claimed;
independent re-derivation at five cells (three new irregular pairs). Bibliography: \cite{sharifi_reciprocity} added, journal data of \cite{fukaya_kato}, \cite{mccallum_sharifi}, \cite{sharifi_eisenstein} confirmed (Project Euclid listings).

Version 1.21.0-pre.9 --- 2026-10-06 (hecke-extensions slot~21; \textbf{a pre-release, not signed, not refereed}): P5$'$ (\texttt{docs/implementation/slot21/P5prime\_note.md}, \texttt{08d83ab3}) ---
Remark~\ref{rmk:lviwasawa}(9): the plus-side pair $\{-k,n+k\}$ (through pre.8 ``recorded, not derived'') identified with the McCallum--Sharifi forced zero: $h(b)=\ell^{-d}\mathcal L^{+}_b(0)\bmod\ell$ equals one unit
times the McCallum--Sharifi pairing vector $x_{b+1}$ (\CMPD, $32$ cells, $28$ irregular pairs, $37\le\ell\le691$, $4961$ branch values, $0$ mismatches), and $\lambda^{+}\ge1$ at $\{-k,n+k\}$ at every $d=1$ cell is a theorem
\textbf{conditional on Sharifi's conjecture} (Fukaya--Kato, \STD, recalled, not read); the second index silent on the plus side, seen by the minus side at $157$; what stays open listed there.
Remark~\ref{rmk:lvmeasplus}(iii)'s pointer sentence updated. Three bibliography items (\cite{mccallum_sharifi}, \cite{sharifi_eisenstein}: read on the authors' web PDFs, not held; \cite{fukaya_kato}: recalled).
No label added, renamed or removed; no schema key, record key, golden line, digest or census verdict changes. Companion pointers: [C] v1.21.0-pre.3, [U] 0.5.0-pre.8, [L37] 0.5.1-pre.6, [RC] 0.1.0-pre.12.

Version 1.21.0-pre.8 --- 2026-10-06 (hecke-extensions slot~21; \textbf{a pre-release, not signed}; the Fable referee F4 of P5, P2x and P4x
(\texttt{docs/implementation/slot21/rc/F4\_referee\_p5\_report.md}: P5 ACCEPT WITH REPAIRS, P2x ACCEPT, P4x ACCEPT; 0\,B\,/\,1\,M\,/\,9\,m), its repairs applied).
\textbf{M1} --- Remark~\ref{rmk:lviwasawa}(7): the pole branches of the minus measure are \emph{derived}, no longer a record: $(\mu^{-},\lambda^{-})=(0,\lambda(G_{b_0})-1)$ on both pole branches,
$b_0$ the cell's own irregular index, from Proposition~\ref{prop:lveismeas}(4), (6) (the polar term) and $\mu(G)=0$ --- the Greenberg--Vatsal formula with $\lambda(G_0):=-1$ for the pole;
$=(0,0)$ at the nine cells by $\lambda(G_{b_0})=1$. Minors: $\varepsilon=\pm1$ (not $\mp1$) for $r^{\pm}$ in the notation line of Remark~\ref{rmk:lviwasawa}; (3) $\sigma\circ\sigma=\mathrm{id}$ written in, and no self-dual
branch is a pole branch; (6) ``$\lambda(G_b)=0$ iff $\ell\nmid B_b/b$'' is \PROVED{} \emph{given} $\mu(G_b)=0$, and the step $\lambda_p=\sum_b\lambda(G_b)$ behind \cite{hho} is \STD{} (Iwasawa's theorem for the
minus part, with Vandiver from \cite{hho} item~(2)), not part of what is read on the page; (8) the unexplained pairs occur at the rate of chance; (9) the ``natural candidate'' corrected: $p_0$ itself carries no branch
but $0$ and $n$ (Remark~\ref{rmk:lvlegc}~(L2), Proposition~\ref{prop:lvmtt}(6)), the object is $w=(r^{+}-u\,p_0)/\ell^{d}$ of Remark~\ref{rmk:lvmeasplus}(iii). Independent re-derivation by the
referee (own Python, exact rationals, no project code, \texttt{docs/implementation/slot21/rc/f4p5/}): the Kubota--Leopoldt $(\mu,\lambda)$ at the seven irregular pairs and at every even class of the $16$ primes
(tangent-number Bernoulli numbers, divided differences); $\Delta$'s tables at $(10,11)$, $(10,23)$, $(10,37)$ by the values oracle, both signs, with $r^{+}$ as the $T_2$-eigenvector of Merel's Heilbronn matrices
(every entry of (8) reproduced; the parity rule fixes $\varepsilon(r^{-})=-1$); the wild pole-branch identity of Remark~\ref{rmk:lvwild}(7) at $(30,37)$, $t=1,2$ ($0/62$, the controls $62/62$, $62/62$, $0$).
No label added, renamed or removed; no schema key, record key, golden line, digest or census verdict changes. Companion pointers: [C] v1.21.0-pre.3, [U] 0.5.0-pre.8, [L37] 0.5.1-pre.6, [RC] 0.1.0-pre.11.

Version 1.21.0-pre.7 --- 2026-10-06 (hecke-extensions slot~21; \textbf{a pre-release, not signed, not refereed}): P4x and P5 (b), (c).
P4x --- Remark~\ref{rmk:lvwild}(7) at $(226,103)$ in $128$-bit $\mathcal O$-arithmetic ($0/454$; the ``not covered: $(226,103)$'' caveat removed; \CMPD) and Check~\ref{chk:lvwild}'s counts.
P5 --- new Remark~\ref{rmk:lviwasawa} (the tame functional equation $\lambda_{\chi,j}=-\chi(-1)(-1)^{j}\lambda_{\bar\chi,n-j}$ at conductor $1$ and $\ell$, \PROVED; its branch form and the
corollaries (FE-inv), (FE-parity), the near-trivial zero, \PROVED; the level-$1$ Greenberg--Vatsal formula off the pole branches, \PROVED{} given its inputs; the Kubota--Leopoldt invariants; the
$(\mu,\lambda)$ tables as a record; every run \CMPD; one plus-side observation recorded, not derived) and new Check~\ref{chk:lviwasawa}; pointer sentences in Claim~\ref{cl:lvmeascong}(5) and
Remark~\ref{rmk:lvmeasplus}(iii); Appendix~A rows. Two labels added (\texttt{rmk:lviwasawa}, \texttt{chk:lviwasawa}), none renamed or removed; no tag of an existing statement moved.
Companion pointers: [C] v1.21.0-pre.3, [U] 0.5.0-pre.7, [L37] 0.5.1-pre.6, [RC] 0.1.0-pre.10.

Version 1.21.0-pre.6 --- 2026-10-06 (hecke-extensions slot~21; \textbf{a pre-release, not signed, not refereed}): P2x --- at $(10,691)$ the Eisenstein congruence of measures is checked on $\Delta$'s own eigensymbol at level $2$
(Proposition~\ref{prop:lveismeas}(6), Claim~\ref{cl:lvmeascong}(4) and Check~\ref{chk:lveismeas}: arms E8L2 / E9L2, $14$ checks added; the ``level~$1$, modulo $691$'' caveat removed; \CMPD). No label added, renamed or removed; no other statement changed.
Companion pointers: [C] v1.21.0-pre.3, [U] 0.5.0-pre.6, [L37] 0.5.1-pre.6, [RC] 0.1.0-pre.9.

Version 1.21.0-pre.5 --- 2026-10-06 (hecke-extensions slot~21; \textbf{a pre-release, not signed, not refereed}): Claim~\ref{cl:glover}(3) \STD{} $\to$ \PROVED{} --- a
proof by Brauer characters is written in (the Brauer character of $\mathrm{Ind}_B^G(\chi_1^{s}\chi_2^{t})$ from Frobenius' formula equals that of
$\Sym^{r}\otimes\det^{t}\oplus\Sym^{\ell-1-r}\otimes\det^{s}$ on every semisimple class), since \cite{glover} has no induced-module statement (its (4.1), (4.2), (5.7), (6.4) are explicit
bases); \cite{glover} now cited for item~(1) by locator ((3.1), p.~433; (6.2), p.~459). Tag line and the summary row now \STD{} (1) / \PROVED{} (2)--(4). Checked exactly (own Python,
$\Z[x]/\Phi_{\ell^2-1}$, no project code) for every $(s,t)$ at $\ell\le13$, and the multiplicities of $\Sym^{n}$ against the recursion of (4) for $n\le3\ell(\ell+1)$.
No label added, renamed or removed; no other statement changed. Companion pointers: [C] v1.21.0-pre.3, [U] 0.5.0-pre.6, [L37] 0.5.1-pre.6, [RC] 0.1.0-pre.8.

Version 1.21.0-pre.4 --- 2026-10-06 (hecke-extensions slot~21; \textbf{a pre-release, not signed}; the Fable referee F4m of the measure side (P2, P3, P4;
\texttt{docs/implementation/slot21/rc/F4\_referee\_measure\_report.md}: ACCEPT WITH REPAIRS, 0\,B\,/\,1\,M\,/\,12\,m), its repairs applied). \textbf{M1} --- the literature leg of
Proposition~\ref{prop:lveismeas}(3) corrected: $E_1=-x\,\zeta_\ell$ (not $x\,\zeta_\ell$), so each of the two $G$-factors is \emph{minus} the branch series of $E_1$, resp.\ $\iota_\ast E_{k-1}$
(the product unchanged), and $dE_r=x^{r-1}dE_1$ in this text's normalisation (Washington's $E_k$ carries no $1/k$); Washington \S12.1 p.~232 and Thm.~12.2 p.~239 read on the page
(\STD{} recalled $\to$ read on the page). \textbf{Tag moves, each with its proof or its locator written in}: Zagier's $p_k^{-}\in\hat W_k^{-}$ read on the page (Prop.~(i), eq.~(11), p.~453; Claim~\ref{cl:lvmeascong}(1),
Proposition~\ref{prop:lveismeas}(1)); the residue constant of Proposition~\ref{prop:lveismeas}(4) \PROVED{} given (3) and $\zeta_\ell$'s residue; in Remark~\ref{rmk:lvwild}: the root number
$w=\tau(\chi)^{2}/M$ (1), the symmetry $\lambda_{\chi,j}=-\chi(-1)(-1)^{j}\lambda_{\bar\chi,n-j}$ for every $P\in W$ (3), the wild Kubota--Leopoldt convention and the pole-class valuation (5),
all \PROVED{} (elementary). Also: the title of \S\ref{sec:lvslot21}; the level-$1$ remark at $(10,691)$ and the $(66,37)$ remark in Proposition~\ref{prop:lveismeas}(6). Independent
re-derivation by the referee (own Python, exact rationals, no project code): Proposition~\ref{prop:lveismeas}(1)--(6) at $(12,5)$ $m\le4$, $(14,7)$ $m\le3$, $(30,37)$ $m\le2$, $(66,37)$,
$(282,37)$ $m\le2$ (the located $14$ reproduced), $(10,691)$; Remark~\ref{rmk:lvwild}(3), (5), (6) at $(10,5)$, $(10,11)$, $(5,2)$, $(7,2)$, $(12,5)$. No label added, renamed or removed;
no statement changed; no schema key, record key, golden line, digest or census verdict changes. Companion pointers: [C] v1.21.0-pre.2, [U] 0.5.0-pre.6, [L37] 0.5.1-pre.5, [RC] 0.1.0-pre.8.

Version 1.21.0-pre.3 --- 2026-10-06 (hecke-extensions slot~21; \textbf{a pre-release, not signed, not refereed}; the results of the $\ell$-adic packages P2 and P4 written in).
\textbf{New} at the end of \S\ref{sec:lvslot21} (so that no statement number moves): Proposition~\ref{prop:lveismeas} (slot-21 P2, \texttt{HeckeLValueEisensteinMeasureValidation}: Zagier's
Laurent polynomial through the Manin-path formula is a Dedekind--Rademacher sum exactly, the Eisenstein measure $\mu^{\mathrm{Eis}}=-E_1\star\iota_\ast E_{k-1}$, its branch series
$-G_{a+1}G_{k-1-a}$ in $\Lambda\otimes\Q$, the polar term on the two pole branches, the comparison with WP-MEAS's polynomial symbol, and the consequence for the cusp form's pole branches),
Check~\ref{chk:lveismeas}, Remark~\ref{rmk:lvwild} (slot-21 P4, \texttt{HeckeLValueWildCharacterValidation}: characters of conductor $\ell^{m}$, $m=2,3$, the constant $\tau(\chi)$ against complex
$L$-values, Proposition~\ref{prop:lvmtt}(4) at wild $\chi$, the symmetry $\lambda_{\chi,j}=-\varepsilon\lambda_{\bar\chi,n-j}$, the reading $v_\pi=e\mu+\lambda$, the pole branches at wild points),
Check~\ref{chk:lvwild}. \textbf{Added text} (no statement changed; tags as in the sources, none upgraded): Claim~\ref{cl:lvmeascong}(4) --- the pole branches \emph{explained} by P2's polar term, the
located $14+4$ derived (the points with $\ell\mid r$), R20-L no longer needed for it; Claim~\ref{cl:lvmeascong}(3) --- a second derivation through P2; Proposition~\ref{prop:lvmtt}, analytic form, and
Remark~\ref{rmk:lvshapes}~(B) --- $\tau(\chi)$ for $\ell$-power conductor \CMPD{} (P4; ``derived, not tested'' through pre.2); the ``Pending P2/P4'' sentences of Check~\ref{chk:lvcoeffring} and
Remark~\ref{rmk:lvtwovar}; \S\ref{sec:lvnorm} row~(M); Appendix~A (four rows added, three amended). Four labels added (\texttt{prop:lveismeas}, \texttt{chk:lveismeas}, \texttt{rmk:lvwild},
\texttt{chk:lvwild}), none renamed or removed. No schema key, record key, golden line, digest or census verdict changes. Companion pointers: [C] v1.21.0-pre.2, [U] 0.5.0-pre.5, [L37] 0.5.1-pre.4.

Version 1.21.0-pre.2 --- 2026-10-05 (hecke-extensions slot~21; \textbf{a pre-release, not signed, not refereed}): one sentence in Remark~\ref{rmk:lvshapes}~(B) --- the torsion of
$W(\Z/\ell^{R})$ and P8's reading of the $\m$-blocks through $W(\Z/\ell^{R+d})$ on both signs (slot-21 F4 on [U] part~4, S2, S3; \SPEC, no statement changes). Companion pointers: [C]
v1.21.0-pre.2, [U] 0.5.0-pre.4, [L37] 0.5.1-pre.3.

Version 1.21.0-pre.1 --- 2026-10-05 (hecke-extensions slot~21; \textbf{a pre-release, not signed, not refereed}; the slot-21 results written in). \textbf{New} at the end of
\S\ref{sec:lvalue}, so that no statement number moves: \S\ref{sec:lvslot21} --- Remark~\ref{rmk:lvcoeffring} (slot-21 P3: the MTT measure over the Hecke order $\mathcal O$ of the eigen-system,
$W(\F_q)$ or totally ramified, Proposition~\ref{prop:lvmtt} verbatim over $\mathcal O$, \PROVED; the cells \CMPD; Hida control modulo $\m$ at one residually non-split component, \CMPD; finding F-P3-1
as the scope sentence), Check~\ref{chk:lvcoeffring} (its gate \texttt{HeckeLValueCoefficientRingValidation}), Remark~\ref{rmk:lvtwovar} (slot-21 P9, \texttt{HeckeLambdaTwoVarValidation}: the
two-variable family by interpolation, the $\Lambda$-adic Eisenstein congruence at $X=\alpha_a$ modulo $37^{3}$ and $103^{2}$, weight two and the $\mathcal L$-invariant of 11a; everything resting on
\Ucite{pr:u-msym}~(U16) stays \CONJ; P1's two limits recorded for \S\ref{sec:lvlfunction}). \textbf{Added text} (no statement changed): Remark~\ref{rmk:lvshapes}~(M) (the coefficient
sentence), \S\ref{sec:lvnorm} row~(M) (values in $\mathcal O$; pointers to P3, P9), a reading of Claim~\ref{cl:lveis}(2) on the ordinary block after its proof (slot-21 F4 M3: [U] (U17)(a)),
\S\ref{sec:lvnf} (slot-21 T1: LLL and POLRED, \texttt{NumberFieldLllValidation}; T2: the degree-$\ell$ Ribet field, \texttt{NumberFieldRibetFieldValidation}, with the referee F4's verdicts on
its claims), Appendix~A (three rows). \textbf{Stale status sentences repaired} (the library inventory of 2026-10-05, \texttt{docs/implementation/slot21/library\_inventory.md}): Mazur and
Greenberg--Vatsal are \emph{held}, not yet read for Claim~\ref{cl:lvmeascong}, Remark~\ref{rmk:lvmeasplus}, Claim~\ref{cl:lvmeas691} and row~(M) (five sites; ``not in the library'' and ``deferred
to R20-L on the owner's scans'' were false); Glover is held and read (the remark after the proof of Claim~\ref{cl:dimformula}: the recursion, not a closed multiplicity formula, is in the paper).
Four labels added (\texttt{sec:lvslot21}, \texttt{rmk:lvcoeffring}, \texttt{chk:lvcoeffring}, \texttt{rmk:lvtwovar}), none renamed or removed; no earlier tag moves. No schema key, record key,
golden line, digest or census verdict changes. \emph{Pending P2/P4}: the exact $\Lambda$-adic Eisenstein measure and wild characters. Companion pointers: [C] v1.21.0-pre.2, [U] 0.5.0-pre.3, [L37] 0.5.1-pre.2.

Version 1.20.0 --- 2026-10-05 (hecke-extensions slot~20, the close, WP-18; \textbf{the release} from 1.20.0-pre.7 under the \textbf{owner's ruling O6} of 2026-10-04
(releases after the Fable windows), on the slot-20 referees' verdicts: F1 ACCEPT WITH REPAIRS (repairs at pre.2, \texttt{docs/implementation/slot20/WP-3\_repairs.md}),
F4 ACCEPT WITH REPAIRS (the research results; repairs at pre.4, \texttt{docs/implementation/slot20/R20-6\_repairs.md}) and F1$'$ ACCEPT WITH REPAIRS (the pre.2 $\to$ pre.6 delta;
repairs at pre.7, \texttt{docs/implementation/slot20/WP-17\_repairs.md}). Over pre.7: Remark~\ref{rmk:frobkummer} gains the result at $(1372,37)$, $b=5$ (slot-20 WP-SW follow-up, \texttt{SW1372\_report.md}): the
swapped-family statement \CMPD{} at all five cells of its range at $37$ (an evidence line; no tag moves, the statement stays \CONJ{} beyond $37$). No other statement,
tag or label changes; no schema key, record key, golden line, digest or census verdict changes. Companion pointers: [C] v1.20.0 (the milestone-M4
\texttt{\textbackslash LEAN} markers), [U] 0.4.0, [L37] 0.5.0, [EX] 1.4.0 (unchanged).

Version 1.20.0-pre.7 --- 2026-10-05 (hecke-extensions slot~20, WP-17: the repairs of the Fable referee F1$'$, \texttt{docs/implementation/slot20/F1prime\_review.md},
ACCEPT WITH REPAIRS, 0 B / 2 M / 6 m, all landed; \textbf{a pre-release, not signed}; record \texttt{docs/implementation/slot20/WP-17\_repairs.md}): M-E1 --- Proposition~\ref{prop:lvlattice}'s
lattice is the diagonal lattice $V'_n=\bigoplus_t\binom nt\Z_\ell u_t=n!\,\Phi_n^{-1}(\Dn)$, $\De$-stable through $\Phi_n$; the identification with the span of the $n$-th powers (false for
$n\ge\ell+1$) is removed (no value moves; WP-LV3's code reads the diagonal lattice); M-E2 --- \textbf{a hypothesis change} (no tag move): Prediction~\ref{pr:frobeniuslayer} is restricted to
$k_0\le\ell$, the boundary weight $k_0=\ell+1$ recorded as refuted as stated at one class, $(1334,37)$; m-E1 --- Remark~\ref{rmk:frobkummer}'s cell list ($J(32)=5$ at $37$, $J(24)=79$, $J(12)=679$;
\CMPD{} at the first four of five cells at $37$, $n=1372$ running); m-E2 --- Check~\ref{chk:lvmeas}(M7) under the heavy arm; m-E3 --- Claim~\ref{cl:lvtriangle}(L) and Remark~\ref{rmk:lvlegc}:
item~3 \CONJ{} as an index statement, \CMPD{} as a run; m-E4 --- Appendix~A row \texttt{chk:lvmeas}; m-E5 --- Proposition~\ref{prop:lvlattice}(3) names the reason for $H^{2}=0$;
m-E6 --- $n_0$ no longer re-used for $1296$. No label added, renamed or removed. Companion pointers: [C] v1.20.0-pre.2, [U] 0.4.0-pre.3, [L37] 0.5.0-pre.6.

Version 1.20.0-pre.6 --- 2026-10-04 (hecke-extensions slot~20, the close delta for the Fable referee F1$'$; \textbf{a pre-release, not signed}; list of changed labels
\texttt{docs/implementation/slot20/close\_delta.md}): the text brought into line with slot-20 WP-LV3 (merged \texttt{c1fda051}) and R20-3. \textbf{New}: Proposition~\ref{prop:lvlattice} (the
change-of-lattice constant $c_f^{\varepsilon}=e\,v_\ell(n!)-d_f^{\varepsilon}$, its bounds and a vanishing criterion; \PROVED, R20-3) and Remark~\ref{rmk:lvlatticevalues} (the values, \CMPD{} by
R20-3's two routes and WP-LV3's C++ record; the $d_f=1$ twist rule \CONJ). \textbf{Changed}: Definition~\ref{def:lvlattice} (the value no longer open); Claim~\ref{cl:lvtriangle} --- (L)'s (L1),
(H) and the cross-kind equality at depth~$2$ move from \CONJ{} to \CMPD{} (WP-LV3), $\m_{\mathrm{Eis}}$ cut jointly by $T_2,T_3,T_5$; Check~\ref{chk:lvtriangle} (the reader's scope: tame, totally
ramified orders, an unramified residue extension refused by name; (d)'s console script reads $(30,37)$, $(38,37)$; an outcome paragraph); Remark~\ref{rmk:lviotak} and \S\ref{sec:ring}'s
console list~(h) ($\iota$ honours $\Z/\ell^{k}$, schema \texttt{hecke\_involution\_ring/v1}); Remarks~\ref{rmk:lvlegc}, \ref{rmk:lvshapes}, \ref{rmk:lvgates} (the twisted depth-$2$ rung run, \CMPD),
\S\ref{sec:lvnorm} rows (B), (H), the code map, Appendix~A. Prediction~\ref{pr:frobeniuslayer} and Remark~\ref{rmk:frobkummer} gain slot-20 WP-SW's evidence --- the swapped cells $b=2,3,4$ \CMPD, and an observation at the boundary weight $k_0=\ell+1$
(an unpredicted hidden block at $(1334,37)$, explained by the mechanism, not by the onset formula). S1\,/\,S2's text of pre.5 re-read against \texttt{docs/implementation/slot20/S1\_S2\_report.md}: no change. Two labels added
(\texttt{prop:lvlattice}, \texttt{rmk:lvlatticevalues}); none renamed or removed. Tags moved: Claim~\ref{cl:lvtriangle}(L1), (H), the triangle (\CONJ{} $\to$ \CMPD); Remark~\ref{rmk:lvlegc}(L1) at
depth~$2$ (\CONJ{} $\to$ \CMPD); Definition~\ref{def:lvlattice}'s value (open $\to$ Proposition~\ref{prop:lvlattice}). No schema key of \texttt{hecke\_ext/v5}, golden line, digest or census verdict
changes. Companion pointers: [C] v1.20.0-pre.2 (code-map text only), [U] 0.4.0-pre.3, [L37] 0.5.0-pre.5.

Version 1.20.0-pre.5 --- 2026-10-04 (hecke-extensions slot~20, stretch packages S1 and S2, plan v1.1 \S0.6; \textbf{a pre-release, not signed}): S1 --- \S\ref{sec:lvlfunction},
the truncated one-variable $\ell$-adic $L$-function as the record \texttt{lfunction/v1} (Definition~\ref{def:lvlfun}, \SPEC; Lemma~\ref{lem:lvlfunprec}, the certified precision of each
$T$-coefficient, \PROVED, elementary; Remark~\ref{rmk:lvlfunschema}, the schema and the console verb \texttt{lfunction}); Check~\ref{chk:lvmeas} gains arm (M7). S2 --- Claim~\ref{cl:lvmeas691}, Ramanujan's congruence $\Delta\equiv E_{12}\pmod{691}$ as a congruence of
measures against Kubota--Leopoldt (\CMPD{} at $(10,691)$; $\lambda=1$ at the branches $199$, $501$ predicted and found; the $22$ pole points with a unit difference), and the env arm (M691) of
Check~\ref{chk:lvmeas}. Placed at the end of
\S\ref{sec:lvalue}, so that no statement number moves. Five labels added (\texttt{sec:lvlfunction}, \texttt{def:lvlfun}, \texttt{lem:lvlfunprec}, \texttt{rmk:lvlfunschema}, \texttt{cl:lvmeas691}); none renamed or
removed; no earlier claim's tag moves; a new record schema (not golden); no key of \texttt{hecke\_ext/v5}, golden line, digest or census verdict changes. Companion pointer: [P] v3.11.1.

Version 1.20.0-pre.4 --- 2026-10-04 (hecke-extensions slot~20, R20-6 repairs after the Fable referee F4, \texttt{docs/implementation/slot20/R20-6\_F4\_review.md},
ACCEPT WITH REPAIRS for this text, 0\,B\,/\,1\,M\,/\,1\,m; repairs listed in \texttt{docs/implementation/slot20/R20-6\_repairs.md}; \textbf{a pre-release, not signed}): M-E1 --- the
slot-20 evidence paragraph of Prediction~\ref{pr:frobeniuslayer} is qualified: the family-(c) seed is non-reduced and lies outside the statement's hypotheses, so the paragraph supports
the transport step~(ii) of \emph{Why} only; m-E1 --- the status-table row of Prediction~\ref{pr:frobeniuslayer} carries the slot-20 entry. The statement and its tag do not change.
No label added, renamed or removed; no schema key, record key, golden line, digest or census verdict changes. Companion pointer: [L37] 0.5.0-pre.4.

Version 1.20.0-pre.3 --- 2026-10-04 (hecke-extensions slot~20, R20-5; \textbf{a pre-release, not signed}, for the slot-20 Fable referee F4, R20-6): one evidence paragraph in
Prediction~\ref{pr:frobeniuslayer} (slot-20 R20-2~\S3, the family-(c) onsets on the forms side; \PROVED{} given curvilinearity, \CMPD{} at four cells); the statement and its tag do not
change. No label added, renamed or removed; no schema key, record key, golden line, digest or census verdict changes. Companion pointer: [L37] 0.5.0-pre.3.

Version 1.20.0-pre.2 --- 2026-10-04 (hecke-extensions slot~20, WP-3 repairs after the Fable referee F1, \texttt{docs/implementation/slot20/WP-3\_F1\_review.md},
ACCEPT WITH REPAIRS, 0\,B\,/\,1\,M\,/\,10\,m; repairs listed in \texttt{docs/implementation/slot20/WP-3\_repairs.md}; \textbf{a pre-release, not signed}): M-1 --- Proposition~\ref{prop:lvmtt}(6)
pins its object ($D^{(\alpha)}_Q$, the formula of Definition~\ref{def:lvmtt} on an arbitrary symbol with a fixed unit $\alpha$), Remark~\ref{rmk:lvmeasplus}(ii) cites it with linearity;
m-E1 --- Claim~\ref{cl:lvmeascong}: the reduction to the Eisenstein symbol \PROVED{} given its inputs, its Kubota--Leopoldt form \CMPD; m-E2 --- synced with WP-MEAS's first run (the $14$
pole points at $(282,37)$ listed and located, the two Kubota--Leopoldt routes agree $17/17$ and $50/50$, $178$ checks; no \texttt{CwLayer} addition); m-E3 --- the polar-term valuation's $r$ is the
weight index; m-E4 --- a cross-reference; m-E5 --- Definition~\ref{def:lvlattice}: $\theta=\Phi_n$ of \HA{lem:symdual}, and the carrying of the $f$-line, an ``input'' at pre.1, is \PROVED{} from
\HA{lem:symdual}; m-E6 --- Remark~\ref{rmk:frobkummer}: the $\theta$-socle run is consistent with Prediction~\ref{pr:frobeniuslayer} ``Why (ii)''; m-E7 --- the wall budget of
Check~\ref{chk:lvmeas} is not a gate. No label added, renamed or removed; no earlier claim's tag moves (the well-definedness clause of Definition~\ref{def:lvlattice} is new text with
its own tag). No schema key, record key, golden line, digest or census verdict changes. Companion pointers: [C] v1.20.0-pre.1 (unchanged), [U] 0.4.0-pre.2, [L37] 0.5.0-pre.2.

Version 1.20.0-pre.1 --- 2026-10-04 (hecke-extensions slot~20, WP-2; \textbf{a pre-release, not signed}, for the slot-20 Fable referee F1 (WP-3); release at the
slot-20 close only on F1's and F1$'$'s acceptance and the owner's yes; list of new and changed labels: \texttt{docs/implementation/slot20/WP-2\_drafts.md}). \textbf{New} at the end of
\S\ref{sec:lvalue}, so that no statement number moves: \S\ref{sec:lvmeasure}, the measure-valued \textsf{linear} evaluation (owner ruling O5 of 2026-10-04: no fourth kind) ---
Definitions~\ref{def:lvmeasure} (\SPEC), \ref{def:lvmtt} (the Mazur--Tate--Teitelbaum measure, with the factor $\ell^{k-2}$ the slot-20 plan's formula lacked), Proposition~\ref{prop:lvmtt}
(measure, moments, interpolation; \PROVED, slot-20 R20-1), Claim~\ref{cl:lvmeascong} (the Eisenstein congruence of measures against Kubota--Leopoldt, \CMPD{} at $(30,37)$,
$(282,37)$, $(226,103)$; the two pole branches an explicit exception with their computed failure points, the explanation deferred to R20-L), Remark~\ref{rmk:lvmeasplus} (the plus side;
the $\lambda=1$ branch pairs an observation, not a claim), Check~\ref{chk:lvmeas} (WP-MEAS's gate, arm (M0) in the \texttt{standard} tier); \S\ref{sec:lvdepth2} --- Claim~\ref{cl:lvtriangle}
(the three kinds at depth~$2$, \CONJ{} with the slot-20 evidence), Definition~\ref{def:lvlattice} (the change-of-lattice constant, \SPEC, value open until R20-3),
Remark~\ref{rmk:lviotak} ($\iota$ over $\Z/\ell^{k}$: nothing new mathematically, \PROVED{} from [C]), Check~\ref{chk:lvtriangle} (WP-LV3's gates). \textbf{Changed}: Remark~\ref{rmk:lvlegc}
(items 2--3 at depth~$2$ stay \CONJ, the new computation cited), Remark~\ref{rmk:lvshapes} (pointers to (M), to Claim~\ref{cl:lvtriangle} and to Definition~\ref{def:lvlattice}),
Remark~\ref{rmk:lvgates} (the twisted depth-$2$ rung's plan), Remark~\ref{rmk:frobkummer} (the $\theta$-socle run at $(1220,37)$: $\theta V_{\mathrm{src}}$ is not the socle, \CMPD), \S\ref{sec:lvnorm}
row~(M), the code map, Appendix~A. Twelve labels added (\texttt{sec:lvmeasure}, \texttt{def:lvmeasure}, \texttt{def:lvmtt}, \texttt{prop:lvmtt}, \texttt{cl:lvmeascong}, \texttt{rmk:lvmeasplus},
\texttt{chk:lvmeas}, \texttt{sec:lvdepth2}, \texttt{cl:lvtriangle}, \texttt{def:lvlattice}, \texttt{rmk:lviotak}, \texttt{chk:lvtriangle}); none renamed or removed; no earlier claim's tag moves. No
schema key, record key, golden line, digest or census verdict changes. Companion pointers: [C] v1.20.0-pre.1, [U] 0.4.0-pre.1, [L37] 0.5.0-pre.1.

Version 1.19.0 --- 2026-10-04 (hecke-extensions slot~19, the close, WP-18; \textbf{the release} from 1.19.0-pre.7 under the \textbf{owner's ruling O3},
on the slot-19 referees' verdicts: F1 ACCEPT WITH REPAIRS (repairs at pre.2), F4 ACCEPT WITH REPAIRS (the research results; repairs at pre.3) and
F1-delta ACCEPT WITH REPAIRS (the pre.3 $\to$ pre.6 delta; repairs at pre.7, \texttt{docs/implementation/slot19/F1delta\_repairs.md}). Release lines only: no statement,
tag or label changes over pre.7; no schema key, record key, golden line, digest or census verdict changes. Companion pointers: [C] v1.19.0 (the milestone-M3
\texttt{\textbackslash LEAN} markers), [U] 0.3.0, [L37] 0.4.0, [EX] 1.4.0 (unchanged).

Version 1.19.0-pre.7 --- 2026-10-04 (F1-delta repairs, textual only: the (H) constant retagged \STD+\CMPD, no tag otherwise moves; before that: hecke-extensions slot~19, WP-2 continued, \textbf{a pre-release, not signed}): the \textsf{bilinear} and \textsf{linear} instances of
Definition~\ref{def:lvroute} \textbf{implemented} --- Remark~\ref{rmk:lvshapes}~(H) (slot-19 WP-BIL, \texttt{HidaPairingRoute}: G1 on the code's \texttt{PoincareDualityPairing}, $T_p^{\ast}=T_p$
for $p=2,3,5$, R19-9's valuations reproduced, \CMPD; (H2) at depth~$2$ stays \CONJ; the change-of-lattice constant for $n\ge\ell$ open) and~(B) (slot-19 WP-LIN,
\texttt{BirchModularSymbolRoute} on the C++ cone's $H^{1}_c$: $\Psi\colon H^{1}_c\cong W$, [C]'s $T_p$ is the symbol $T_p$, a geometric-even symbol is [C]-minus, the $(258,37)$ sign facts,
\CMPD); a \textbf{new} Remark~\ref{rmk:lvlegc} (leg~(c) in modular-symbol form: $r^{+}\equiv u\,p_0$, $r^{-}\equiv s\,\mathrm{Bern}_k$ modulo~$\ell$ to depth $v(B_k/k)$, twisted values of
$r^{+}$ $\equiv0$; \CMPD{} at $8$ cells, \CONJ{} at depth $\ge2$); \S\ref{sec:lvnorm} rows~(B), (H), the code map and Appendix~A. One label added (\texttt{rmk:lvlegc}); none renamed
or removed; no earlier claim's tag moves. Companion pointers: [C] v1.19.0-pre.5, [U] 0.3.0-pre.4, [L37] 0.4.0-pre.3.

Version 1.19.0-pre.5 --- 2026-10-04 (hecke-extensions slot~19, WP-2 continued: the slot-19 results that landed after pre.3/pre.4 folded in with their own tags, \textbf{a pre-release, not signed};
change table \texttt{docs/implementation/slot19/WP-2\_pre5\_changes.md}). Remark~\ref{rmk:lvgates}: the twisted Claim~\ref{cl:lvtwist} \CMPD{} at four cells and two controls (slot-19 WP-LV2),
the deep rungs $(282,37)$, $(226,103)$ \CMPD{} (torsion-free, rank~$1$, saturated, $\#Q_m=\ell^{\min(m,2)}$; WP-LV2 and, by an independent reading, R19-7), the two-sided bound
tested on planted algebras; ``the twisted gate is not run'' deleted. \S\ref{sec:lvnorm} row~(a): the level-$m$ Coates--Wiles measure from $\Q_\ell(\zeta_{\ell^{m}})$ (slot-19 R19-7;
\PROVED{} identity, \CMPD{} values, the bridge to the unramified Kummer lift \PROVED{} from \cite{washington1st} Thm.~13.54, Prop.~13.51, Lemma~13.52); rows~(B), (H) and
Remark~\ref{rmk:lvshapes}: the research prototypes of the \textsf{linear} kind (slot-19 R19-8, Birch\,/\,Manin--Stevens) and the \textsf{bilinear} kind (slot-19 R19-9, Hida's
pairing) with their slot-20 specifications; (H)'s constant from \cite{hida_cong} Thm.~5.1, pp.~247--248, read on the page, no Euler factor removed (the \VERIFYloc{} of row~(H) is
settled). Prediction~\ref{pr:frobeniuslayer}: slot-19 evidence at $(1220,37)$ and a \textbf{new} Remark~\ref{rmk:frobkummer} (the Frobenius layer on a Kummer class: $I=\m^{3}$,
\CMPD{} at one cell, \CONJ{} beyond); Prediction~\ref{pr:latticesplit}: one evidence sentence, no change. The code map of \S\ref{sec:lvalue}: the twisted route \textbf{built}
(WP-LV2). One label added (\texttt{rmk:frobkummer}); none renamed or removed; no claim's tag moves. Companion pointers: [C] v1.19.0-pre.4, [U] 0.3.0-pre.4, [L37] 0.4.0-pre.2.)

Version 1.19.0-pre.4 --- 2026-10-04 (hecke-extensions slot~19, isign-fix: Claim~\ref{cl:258lattice} states that on $H^{1}_c$ the free sign is the other one ($\m$ residually Eisenstein; $\Lambda_\m^{+}\cong\T_\m$ on $H^{1}$ stands), its proof reads ``in characteristic~$0$''; Remark~\ref{rmk:lvshapes}(B): the normalisation fold-in of R19-8 \S1.4 (\texttt{docs/implementation/slot19/isign\_resolution.md} \S\S4--5). No label, tag, code-side value, golden line, digest or census verdict changes.)

Version 1.19.0-pre.3 --- 2026-10-04 (hecke-extensions slot~19, WP-2 continued after the Fable referee F4 of the slot-19 research results,
\texttt{docs/implementation/slot19/R19-6\_F4\_referee\_report.md} (ACCEPT WITH REPAIRS, $0$ blocking, $1$ major, $18$ minor); \textbf{a pre-release, not signed};
change table \texttt{docs/implementation/slot19/WP-2\_pre3\_changes.md}). The accepted results enter with F4's tags (its \S3): Claim~\ref{cl:258lattice} gains the
measured sign, $\Lambda_\m^{+}\cong\T_\m$, $\Lambda_\m^{-}\cong\m$ (\CMPD, WP-LC), and F4 m1's wording on \cite{bass_ubiquity}; Check~\ref{chk:latcmp}'s run (held at
$M=1,\dots,4$, \CMPD); a \textbf{new} Remark~\ref{rmk:latcmpseven}, the $\ell=7$ controls of WP-LC (\CMPD; F4 m5--m7: the over-order chain at $(48,7)$ told apart by
$U_7$, the scope of the per-sign freeness failure, freeness by \HA{lem:h2torsion}); Remark~\ref{rmk:258next}'s pointer; Remark~\ref{rmk:lvgates}: WP-LV's gate cells
(\CMPD, F4 m18), the Bernoulli routes as built, and the cells $(282,37)$, $(226,103)$ after [U] 0.3.0-pre.3's (U13); \S\ref{sec:lvnorm} rows~(a), (b): the
\SPEC{} placeholders replaced by slot-19 R19-4's text, with F4 m15's naming of $\varphi_k$ and its locators (\cite{washington1st} \S13.7, \cite{coates_wiles} p.~230);
\S\ref{sec:lvnf} and the code maps of \S\ref{sec:block258} and \S\ref{sec:lvalue}: \textbf{built} (WP-LC, WP-LV, WP-NF; the twisted Claim~\ref{cl:lvtwist} not built).
One label added (\texttt{rmk:latcmpseven}); none renamed or removed; tags: rows~(a), (b) of \S\ref{sec:lvnorm} move from \SPEC{} to their own tags, \S\ref{sec:lvnf} is
built; no claim's tag moves. Bibliography: \cite{coates_wiles}. No code-side value, golden line, digest, record key or census verdict changes. Companion pointers: [C]
v1.19.0-pre.3, [U] 0.3.0-pre.3, [L37] 0.4.0-pre.1.

Version 1.19.0-pre.2 --- 2026-10-03 (hecke-extensions slot~19, WP-3 repairs after the Fable referee F1, \texttt{docs/implementation/slot19/WP-3\_F1\_referee\_report.md};
\textbf{a pre-release, not signed}). M-2: Definition~\ref{def:lvroute}(2) gains the kind \textsf{congruence} (a Hecke character $e$ and its congruence module
$Q_e=\T(H')_{\mathcal O}/(I_e+\mathrm{Ann}\,I_e)$), which carries the instance (E) (Definition~\ref{def:lveis}); (B) stays \textsf{linear}, (H) \textsf{bilinear};
Remark~\ref{rmk:lvshapes}, the code map and Appendix~A updated. M-1: Claim~\ref{cl:lveis}(3)'s hypothesis is \emph{saturation} of the $\m_{\mathrm{Eis}}$-block's
Hecke algebra (the \texttt{stable} verdict at $K=1$; automatic on a rank-$1$ block), with a two-sided bound otherwise, and its inequality runs over $p\neq\ell$;
Remark~\ref{rmk:lvgates} says the hypothesis holds at the four untwisted cells and is open at $(282,37)$, $(226,103)$. Minor (F1 m-E1 \dots m-E8): Definition~\ref{def:lveis}'s
record reads $L=\zeta$ at $s_0=1-k$; the proof of Claim~\ref{cl:lveis}(1) matches $I_{\mathrm{Eis}}$ with [U]'s $I^{0}$; Claim~\ref{cl:lvtwist}'s Hida locator verified on the
page (ENS 1986 pp.~249--250, Cor.~3.2 (3.3b)) with the character-convention check, and \cite{hida_ens}'s note updated; Remark~\ref{rmk:lvshapes}: the congruence modules of
the $(258,37)$ block, and $P_{j,a}$ of degree $n$; Claim~\ref{cl:258alg}(5): the unique intermediate $\T_\m$-lattice (F1); Claim~\ref{cl:258lattice}'s proof: Jordan type in
place of Krull--Schmidt, the Gorenstein input proved inline, \HA{prop:cohforms}(5) cited; Remark~\ref{rmk:lvgates}: the $691$ rungs $342252$, $342440$ computed (F1, tracked
in \texttt{docs/implementation/slot19/wp3/}); a macro comment moved. Tags unchanged; no label added, renamed or removed; Appendix~A rows of
Claims~\ref{cl:258lattice}, \ref{cl:lveis}, \ref{cl:lvtwist} updated. Companion pointers: [C] v1.19.0-pre.2, [U] 0.3.0-pre.2.

Version 1.19.0-pre.1 --- 2026-10-03 (hecke-extensions slot~19, WP-2; \textbf{a pre-release, not signed}, for the slot-19 Fable referee F1 (WP-3);
release at the slot-19 close only on its acceptance and the owner's yes). \textbf{New}: \S\ref{sec:block258} at the end of \S\ref{sec:ext} --- the block
$(258,37)$ exactly (slot-19 R19-1): Claim~\ref{cl:258alg} (\CMPD, with (2) $\T_\m\cong\T^{\mathrm f}_\m$ \PROVED{} from \HA{prop:cohforms}), Claim~\ref{cl:258lattice}
($\Lambda_\m\cong\T_\m\oplus\m$, \PROVED{} given the \CMPD{} over-order data, [L37]'s $V$ and Gorenstein ideal theory --- a weaker input than R19-1's Bass structure theorem,
through the $\iota$-split), Remark~\ref{rmk:258Urank}, Check~\ref{chk:latcmp} (the gate of WP-LC), Remark~\ref{rmk:258next}; a new \S\ref{sec:lvalue} before the
appendices (owner ruling O6) --- the $L$-value interface (Definition~\ref{def:lvroute}, \SPEC; the Birch and Hida shapes, Remark~\ref{rmk:lvshapes}), the
Eisenstein\,/\,Bernoulli instance (Definition~\ref{def:lveis}; Claim~\ref{cl:lveis}: ring form \PROVED{} from [U]'s (U1), module form \STD{} for $n<\ell-1$; the twisted
Claim~\ref{cl:lvtwist}, \PROVED{} given Hida's control at the nebentypus points), the normalisation table (\S\ref{sec:lvnorm}; legs (a), (b) \SPEC, R19-4), the
number-field layer (\S\ref{sec:lvnf}, \SPEC), and the planned code maps of WP-LC, WP-LV, WP-NF. Harder's boundary-spectral-sequence route is not used. New tags
\CMPD, \SPEC; new environment \emph{Check}; bibliography: \cite{bass_ubiquity}, \cite{hida_cong}, \cite{hida_ens}, \cite{mtt}. Appendix~A: seven rows added. No earlier
statement, tag or label changed; no schema key, record key, golden line, digest or census verdict changes. Companion pointers: [C] v1.19.0-pre.1, [U]
0.3.0-pre.1, [EX] 1.4.0 (unchanged), [L37] 0.3.0 (unchanged).

Version 1.18.0 --- 2026-10-03 (hecke-extensions slot~18, the close, WP-18; \textbf{the release} from 1.18.0-pre.2 under the \textbf{owner's ruling O3
of 2026-10-03}, on the slot-18 referee F1's ACCEPT WITH REPAIRS, repairs at pre.2). Over pre.2, two pointer texts of the slot's stretch packages, no statement
or tag change: Stage~4's (\S\ref{sec:certgen}) ``Scope and order'' paragraph now cites census editions~1--3 and the mixed-finality scans, including slot-18 WP-19c's $64$-cell
probe ($\ell\in\{17,\dots,43\}$ on the Dickson cells, $n\le200$: no block with $\ker r^1|_V\ne0\ne S_V$), and states that Proposition~M stays specified only
(certificate (d2) lands with the next record event that needs it; \texttt{docs/implementation/slot18/WP-19c\_report.md} \S5); the digit route's item records
slot-18 WP-19b's cut at the Hecke step (\texttt{DigitCellSession}, value-free, \texttt{HeckeLadicSplitValidation}). Census edition~3.2 stays staged, not published
(owner ruling O7: published on the return that completes it). No label added, renamed or removed; no schema key, record key, golden line, digest or census
verdict changes. Companion pointers: [C] v1.18.0, [EX] 1.4.0, [U] 0.2.1 (unchanged).

Version 1.18.0-pre.2 --- 2026-10-03 (hecke-extensions slot~18, WP-3 repairs; \textbf{a pre-release, not signed}; release as 1.18.0 at the slot-18 close
on the owner's yes, ruling~O3). The slot-18 Fable referee F1 (\texttt{docs/implementation/slot18/WP-3\_F1\_referee\_report.md}) was \textbf{ACCEPT WITH REPAIRS},
0 blocking; its major M-2 and minors m-E1--m-E7 are applied here (table: \texttt{docs/implementation/slot18/WP-3\_repairs.md}). \textbf{Evidence, not statements}
(the referee's nine forms-side cells, predicted before each run, $9/9$ held; tracked at \texttt{docs/implementation/slot18/wp3/}): Prediction~\ref{pr:depth}'s
``untested'' $(m,r_0)=(2,3)$ now has a datum, $(34460,547)$, $\F_{547}[x]/x^6$, so the hypothesis $m\le2$ (not $m r_0\le4$) is the one the data support;
Prediction~\ref{pr:kfonset}~(C)'s first window sampled at its last and middle cells, (B) at five more cells; (C) still asserts nothing at the meeting.
\textbf{Wording} (no tag change): Claim~\ref{cl:dimformula} --- the Edixhoven fallback (a lower bound; m-E1), $d\in\{1,\ell-2\}$ and the semisimplicity /
multiplicity-one input named (m-E2), the parity corollary $m_O\equiv m_K$; the Edixhoven status ``attribution verified in a secondary source, locator pending''
(slot-18 R18-4), Berger--Li--Zhu as a second locator for $w\le\ell$ (new bibliography item \cite{berger_li_zhu}); Remark~\ref{rmk:dimwindows} cites the exact
event-sequence check (m-E3); Prediction~\ref{pr:depth}~(b)'s torsion-freeness (m-E5); census~3.2's ranges (m-E6); ``at the same twist'' (m-E7). The
\texttt{conj:law} pointer: R18-1's final figures, nothing taken in (M-2(d), m-E4; first written with F1's preliminary figures, replaced by R18-5). Appendix~A: four rows updated. No label added,
renamed or removed (one bibliography key added); no schema key, record key, golden line, digest or census verdict changes. Companion pointers: [C] v1.18.0-pre.2,
[EX] 1.4.0-pre.2, [U] 0.2.1 (unchanged).

Version 1.18.0-pre.1 --- 2026-10-03 (hecke-extensions slot~18, WP-2; \textbf{a pre-release, not signed}, for the slot-18 Fable referee F1 (WP-3);
release at the slot-18 close only on its acceptance and the owner's yes, ruling~O3 of 2026-10-03). \textbf{New statements}: \S\ref{sec:dimblock} at the end of
\S\ref{sec:ext} --- Claim~\ref{cl:seedtwist} (\PROVED; slot-17 R17-1 Lemma~3.1, with the standard\,/\,swapped partition), Claim~\ref{cl:selfmirror} (\PROVED{}
given one \STD; R17-1 Lemma~3.2), Claim~\ref{cl:dimformula} (\PROVED{} given \STD{} inputs, Edixhoven's locator not verified; R17-1 Theorem~4.1, the referee
F1's m-E5), Remark~\ref{rmk:dimwindows} (the windows on the multiplicity side) and Prediction~\ref{pr:depth} (\CONJ; the $\ell=37$ report's \texttt{conj:depth}).
\textbf{Prediction~\ref{pr:kfonset} generalised} (the label kept, the title changed): (A) the v1.17.0 statement; (B) the swapped family at $g=1$, same window
form; (C) $g=2$, $\Fl[x]/x^{3}$ in the first window; nothing asserted where no datum exists (the $g=2$ meeting, $g\ge3$, past the windows); new forms-side
evidence by slot-18 WP-2 ($11/11$ predictions held, \texttt{docs/implementation/slot18/wp2/}) and the slot-17 LAT37 lattice cells $(696,37)$--$(700,37)$.
\textbf{Corrections}: Prediction~\ref{pr:kfonset}~(c)'s ``$2092$ of $2092$'' standard classes is $1492$ of $1492$. \textbf{Pointers}: \S\ref{sec:census}
edition~3.2 (staged, not published); Prediction~\ref{pr:latticesplit}'s ``Observed on the Kummer locus'' and ``Where the excluded locus sits'' point to the new
statements; the report's \texttt{conj:law} is \emph{not} taken in (slot-18 R18-1 has not reported). Bibliography: \cite{diamond_im}, \cite{edixhoven},
\cite{washington1st}, \cite{wiles_ordinary}. Appendix~A: one row restated, five added. No label renamed or removed; no schema key, record key, golden line,
digest or census verdict changes. Companion pointers: [C] v1.18.0-pre.1, [EX] 1.4.0-pre.1, [U] 0.2.1 (unchanged).

Version 1.17.0 --- 2026-10-02 (hecke-extensions slot~17, the close, WP-18; \textbf{the release} from 1.17.0-pre.5 under the \textbf{owner's
ruling O3 of 2026-10-01}, which releases 1.17.0 on the slot-17 referee F1's acceptance). F1 (\texttt{docs/implementation/slot17/WP-3\_referee\_F1.md}) was
ACCEPT WITH REPAIRS, 2 blocking (B-1, B-2), 6 minor: B-2 and the minors applied at pre.3, B-1 by the owner's ruling of 2026-10-02 at pre.4 ((R) restated on
congruence-free classes); Prediction~\ref{pr:kfonset}'s window bound stands as a restatement of the \CONJ{} (the unbounded ``from $k$ on'' refuted on the forms side
at $\ell=37$, $k=1412$, $1450$). Over pre.5 nothing changes but this line: no statement, tag or label added, renamed or removed; no schema key, record key, golden
line, digest or census verdict changes. Census edition~3.2 is staged at the slot-17 close and \textbf{not} published (owner-gated); this version cites
edition~3.1 only. Companion pointers: [C] v1.17.0 (the first \texttt{\textbackslash LEAN} markers), [EX] 1.3.0, [U] 0.2.1.

Version 1.17.0-pre.5 --- 2026-10-02 (hecke-extensions slot~17, WP-2; \textbf{a pre-release, not signed}; pointers only, no statement or tag change):
Claim~\ref{cl:sigma} gains the verified locators \cite[Thm.~4.1, p.~276]{jochnowitz} (``weight $\le\ell+1$'') and \cite[Thm.~2.2, p.~204]{ash_stevens}
(finiteness), slot-17 R17-4; Prediction~\ref{pr:kfonset}~(c) records the window's last cell off~$37$ on the forms side --- $(3496,59)$, $(4498,67)$ still
curvilinear of rank~$4$ (the $\ell=37$ report's referee F4). No label added, renamed or removed. Companion pointers: [C] v1.17.0-pre.5, [EX] 1.3.0-pre.4,
[U] 0.2.1-pre.2.

Version 1.17.0-pre.4 --- 2026-10-02 (hecke-extensions slot~17, WP-2; \textbf{a pre-release, not signed}; \textbf{owner ruling of 2026-10-02 on
finding~B-1} of the referee F1). Prediction~\ref{pr:latticesplit} \textbf{restated a third time} (a statement change of a \CONJ, the tag kept): (R) is now asserted
on \emph{congruence-free} classes (for every even $w\le\ell+1$ and every twist $j$, the generalised eigenspace of $(p^{-j}a_p)_{p\ne\ell}$ in $M_w(\Fl)$ has
dimension $\le1$; on reducible non-Dickson classes this equals Kummer-free), as the referee proposed; the sentence ``(R) at every regular~$\ell$'' is deleted. The
v1.17.0-pre.2 Kummer-free form is recorded as \textbf{contradicted on the forms side} by the six cells $(2210/2278,67)$, $(3798/3930,131)$, $(4030/4170,139)$
($\T^{\mathrm f}_\m/\ell\cong\Fl[x]/x^4$ at irreducible classes of two congruent eigenforms of weight $\le\ell+1$; reproduced with tracked tooling, slot-17
R17-2~\S7; congruences confirmed over~$\Z$; $139$ regular; six further pairs the same), beside the v1.16.0 form and the v1.14.0 statement, both kept. Followers:
the prediction's evidence paragraph (no evidence lost), \S\ref{sec:census}'s clause, Appendix~A's row. No label added, renamed or removed; no schema key, record
key, golden line, digest or census verdict changes. Companion pointers: [C] v1.17.0-pre.4, [EX] 1.3.0-pre.4, [U] 0.2.1-pre.2.

Version 1.17.0-pre.3 --- 2026-10-02 (hecke-extensions slot~17, WP-2 resumed; \textbf{a pre-release, not signed}: the repairs of the slot-17 Fable
referee F1, \texttt{docs/implementation/slot17/WP-3\_referee\_F1.md}, [E] ACCEPT WITH REPAIRS, 2 blocking, 6 minor). \textbf{B-2} applied: Prediction~\ref{pr:kfonset}
gains the hypothesis $g_{\ell+1-w}=1$ (refuted as stated at $g=2$ by the document's own $(608,547)$: Loewy length~$3$ below the onset, $\dim V=12$ in the window),
``nothing asserted at $g\ge2$'', item~(d) widened to every class at a pair with $g\ge2$, Appendix~A's row; a restatement of a \CONJ{} under its label.
\textbf{B-1} \emph{not applied, awaiting the owner's ruling}: (R) of Prediction~\ref{pr:latticesplit} is contradicted on the forms side at irreducible classes
($\ell=67$, $131$ and the regular $139$; reproduced by slot-17 R17-2~\S7 with tracked tooling, six further pairs the same); its statement, the regular-prime sentence
and the ``Refuted'' list are unchanged at this version. Minors: m-E1 (item~(c): the window's upper end computed on the multiplicity side at $13$ pairs, $2092/2092$
classes; only the algebra's shape extrapolated), m-E2 (the embedding-dimension sentence located), m-E3 (one sentence in (R)'s evidence: $\dim V\le4$ forces
$\m_A^2=0$), m-E4 (item~2's ``$\dim I\ge2$''), m-E5 (R17-1's Theorem~4.1: its inputs VERIFY there, a slot-18 statement); m-E6 needs no edit. No label added,
renamed or removed; no schema key, record key, golden line, digest or census verdict changes. Companion pointers: [C] v1.17.0-pre.3, [EX] 1.3.0-pre.3, [U]
0.2.1-pre.2.

Version 1.17.0-pre.2 --- 2026-10-02 (hecke-extensions slot~17, WP-2, second pass; \textbf{a pre-release, not signed}, for the same referee F1
(WP-3); \textbf{owner ruling O-R17-2a of 2026-10-02}: restate (R) now, before the referee). (1)~Prediction~\ref{pr:latticesplit} \textbf{restated a second
time} (a statement change of a \CONJ): (R) is now asserted on \emph{Kummer-free} classes only (residual representation irreducible, or reducible with both seeds
$\ell$-regular), in place of ``not Kummer--Frobenius''. The v1.16.0 form is \textbf{refuted on cohomology} at $(608,547)$ (the $g=2$ first window:
$A\cong\F_{547}[x]/x^3$, saturated, \texttt{stable} certified; slot-17 R17-2) and contradicted on the forms side at the swapped classes at $\ell=37$, $103$,
$691$ ($\Fl[x]/x^4$) and at $(1156,547)$; it is kept below the new form, marked refuted, with its evidence; new paragraphs on where the excluded locus sits, what
held, and what is observed (not asserted) on the Kummer locus. (2)~Prediction~\ref{pr:kfonset}: the slot-17 evidence at the window's lower end off $37$
($(736,67)$ on cohomology; $67$, $59$, $101$ on the forms side), a scope item~(d), one parenthetical reworded (the $t=J(w)$ class is the Eisenstein end $b=0$ of
the swapped family), and a correction: the first Loewy-$4$ cell at $103$ is the swapped $(2574,103)$, not $(8342,103)$. (3)~\S\ref{sec:census}: one clause on the
second restatement; edition~3.1's \texttt{sat\_cells\_done} quoted at publication ($897$ of $1459$; the brief's E-2, mechanical). Appendix~A: two rows. No label
added, renamed or removed; no schema key, record key, golden line, digest or census verdict changes. Companion pointers: [C] v1.17.0-pre.2 (one citation repair
in \HA{lem:heckerankone}), [EX] 1.3.0-pre.2.

Version 1.17.0-pre.1 --- 2026-10-01 (hecke-extensions slot~17, WP-2, drafted during the slot-16 close; \textbf{a pre-release, not signed},
for the slot-17 Fable referee (WP-3) and the owner's ruling O3 of 2026-10-01, which releases 1.17.0 at the slot-17 close on the referee's acceptance).
(1)~Prediction~\ref{pr:kfonset} \textbf{restated} (a statement change of a \CONJ, hence the $X$-type version): the window
$\ell w+t(\ell+1)\le k<w+t(\ell+1)+(\ell^{2}-1)$ for $1\le t\le J(w)-1$ in place of ``from $k=\ell w+t(\ell+1)$ on'', which the forms side refutes at
$\ell=37$, $k=1412$, $1450$ (the $\ell=37$ report \texttt{docs/latex/l37\_loewy4\_example.tex}, \texttt{rmk:kfonset}, \texttt{conj:window}); the
superseded statement is kept below it, its evidence unchanged; the mechanism paragraph gains one sentence; Appendix~A's row restated. (2)~\S\ref{sec:census}:
a paragraph on \textbf{census edition~3.1} (add-only over edition~3; evidence and pointer only). No label added, renamed or removed; no schema key, record
key, golden line, digest or census verdict changes. Not in this version: slot-17 R17-1's double-divisibility observation (not merged when this was
drafted). Companion pointers: [C] v1.17.0-pre.1 (in the tree; v1.16.0 released), [EX] 1.3.0-pre.1, [U] 0.2.1-pre.1, [X] v1.21.

Version 1.16.3 --- 2026-10-01 (hecke-extensions slot~16, \textbf{owner ruling O4, 2026-10-01}: a \emph{patch} over 1.16.2 applying two
proposals of \texttt{docs/implementation/slot15/WP-12\_U\_E\_proposals.md}; \textbf{no existing statement, tag, label, schema key, record key,
golden line, digest or census verdict changes}). (1)~(E-3) \emph{accepted}: Remark~\ref{rmk:i8switch} (\PROVED{} from \cite{bk}) --- the
split\,/\,ramified switch of slot-14 I-8 is not Bella\"iche--Khare obstructedness; with it the bibliography entry \cite{hho}.
(2)~(E-2) \emph{accepted as a} \CONJ: Prediction~\ref{pr:frobeniuslayer}, the Frobenius layer (F1, WP-12's law L-F1), with its
proof given two per-cell facts, its WP-12 evidence and a dated slot-16 evidence remark (WP-5, WP-11: consistent, no qualification).
Appendix~A: two rows added. Prediction~\ref{pr:kfonset} is untouched (its window-bound change is deferred to slot~17).

Version 1.16.2 --- 2026-10-01 (hecke-extensions slot~16, the close: a \emph{patch} over 1.16.1, one locator only, \textbf{no
statement, tag, label, schema key, record key, golden line, digest or census verdict changes}). Claim~\ref{cl:mkfaithful}'s proof: Serre's
weight congruence is cited as LNM~350, n$^\circ$~1.3, Th\'eor\`eme~1, p.~198, verified on a public scan by slot-16 WP-8
(\texttt{docs/implementation/slot16/WP-8\_citations.md}~\S2), in place of ``\S1''; the paraphrase now gives the source's own hypotheses
($p\ge3$, rational coefficients, $v_p(f-f')\ge v_p(f)+m$), which contain both uses. Companion pointers: [C] v1.16.0, [EX] 1.2.2, [X] v1.20.2.

Version 1.16.1 --- 2026-10-01 (hecke-extensions slot~15, post-close documentation package after census edition~3; a
\emph{patch} release over the owner-signed 1.16.0: evidence and pointers only, \textbf{no statement, tag, label, schema key, record key,
golden line, digest or census verdict changes}; every addition is marked ``v1.16.1''). (1)~\S\ref{sec:census}: a paragraph on
\textbf{census edition~3} (slot-15 WP-15, \texttt{docs/matrix-engine/results/congruence\_census/edition3/}, published beside edition~2,
$80\,293$ rows, complete at $\ell=61\ldots103$, partial at $17\ldots59$), with its coverage and what it adds as evidence.
(2)~Prediction~\ref{pr:latticesplit}: one evidence row (the complete C1b block, $2165$ non-reduced rows at $\ell=61\ldots103$, Loewy
$\le2$) and the edition-3 pointer. (3)~Prediction~\ref{pr:kfonset}: one sentence --- at $\ell=61\ldots103$ the onset law predicts no
Loewy-$4$ cell inside $n\le700$ and none is observed, so the law's positive half is \emph{not} tested there.
Companion pointers: [C] v1.15.0, [EX] 1.2.1, [X] v1.20.2.

Version 1.16.0 --- 2026-10-01, \textbf{owner-signed release} (hecke-extensions slot~15, post-close; owner ruling Q13 of the
slot-15 plan --- release [E] 1.16.0 and [EX] 1.2.0 at the close --- unblocked by the ruling S15-O2 of 2026-10-01, which accepted every open item
and recommendation of \texttt{docs/implementation/slot15/WP-18\_close\_report.md}). \emph{Released} from v1.15.1-pre.7 with no further edit of
the text: v1.15.1 was never released, and because pre.7 changes the statement of a Prediction the release is an $X$-type version, 1.16.0. The
sites keep their ``v1.15.1-pre.$N$'' markers as provenance. What v1.16.0 contains over v1.15.0, gathered in one place:
(pre.1, slot~14) Citro--Ghitza's count of the level-one eigensystems mod~$p$ checked against $\Sigma(5)$, $\Sigma(7)$, $\Sigma(11)$,
$\Sigma(13)$, $\Sigma(691)$ (\S\ref{sec:atlas}), and its bibliography entry;
(pre.2, slot-14 audit, owner rulings Q-C, Q-D) the definitional paragraph ``The Hecke algebra over $\Z/\ell^{K}$, read mod~$\ell$'' at the end of
\S\ref{sec:ring}, and the sentence keeping the frozen $u_\ell$ RSS model \texttt{u\_ell\_rss/v2};
(pre.3, slot~14, P12's cluster half) the first-cluster evidence sentence of Prediction~\ref{pr:latticesplit} (job \texttt{34551683}, $1233$ blocks,
\texttt{noctua2\_section\_m/});
(pre.4, slot-15 WP-19, owner Q15) Claims~\ref{cl:mkfaithful} and~\ref{cl:formsF9}, and the existence of $B'(k,\ell)$ at every $\ell\ge5$;
(pre.5, slot-15 WP-5, owner Q4) the \emph{certified} \texttt{stable} verdict (\HA{cor:stablecert} in the referee's B1 form) and the digit route
beyond $2^{64}$ in that paragraph;
(pre.6, slot-15 WP-5 Fable audit M1--M3) $\varepsilon_K$ as the code computes it, the eigenspace comparison as a consistency check, one scope
sentence;
(pre.7, owner S15-O2) Prediction~\ref{pr:latticesplit} \textbf{refuted as stated} at $\ell=37$ and restated as~(R) with the Kummer--Frobenius
exclusion, and the new Prediction~\ref{pr:kfonset}.
No schema key, record key, golden line, digest or census verdict moves in any of the seven. Companion pointers: [C] v1.15.0, [EX] 1.2.0,
[U] 0.1.1-pre.6 in the tree (unchanged; binding v0.1.0), [X] v1.20.1.

v1.15.1-pre.7 --- 2026-10-01 (slot-15 post-close, \textbf{a pre-release}; owner ruling S15-O2 of 2026-10-01, which accepted
proposal~(E-1) of \texttt{docs/implementation/slot15/WP-12\_U\_E\_proposals.md} and the Fable referee
\texttt{docs/implementation/slot15/WP-4\_referee.md}~\S8): Prediction~\ref{pr:latticesplit} is \textbf{refuted as stated} --- at
$\ell=37$ the $18$ certified C1a records $\operatorname{eis}(a,a+5)$, $n=654\ldots700$, have $\texttt{loewy\_A}=4$, Hilbert function
$[1,1,1,1]$, module Loewy $[3,2,2,1]$, confirmed on the forms side and on the lattice mod~$37^2$ (slot-15 WP-13: $(654,37)$
ideal~$8$, \texttt{CENSUS\_LATTICE} \texttt{status=failed}, \texttt{first\_fail=(2,2)}) --- and restated, under its label, as the
conjecture~(R) with the Kummer--Frobenius exclusion past the onset weight $\ell w+t(\ell+1)$, with its evidence table; the v1.14.0
statement and its evidence stay in the text, marked refuted. New: Prediction~\ref{pr:kfonset} (\CONJ, the onset). Consistent with
[C] v1.15.0 \HA{prop:cohforms}(5). Appendix~A: one row changed, one added. The tag on Prediction~\ref{pr:latticesplit} stays \CONJ\
(for the restated form); no schema key, golden line, digest or verdict moves; the C++ readings keep reporting, never gating.

v1.15.1-pre.6 --- 2026-09-30 (slot-15 WP-5 Fable audit repairs, \textbf{a pre-release, not signed};
\texttt{docs/implementation/slot15/WP-5\_fable\_audit.md} \S3, M1--M3): in \S\ref{sec:ring}'s paragraph ``The Hecke algebra over $\Z/\ell^{K}$, read
mod~$\ell$'' --- (i) $\varepsilon_K$ defined as the code computes it, $K\cdot\dim_{\Fl}(A_K/\ell)-\mathrm{length}(A_K)$ (M2); (ii) the eigenspace-dimension
comparison stated as a consistency check, the identification $e_\m=e_{\m'}$ resting on $e_\m\in\Z[T_p]$, $e_\m^2=e_\m\ne0$ and the nilpotency (M1); one
scope sentence, the printed Hilbert function a reading, the certificate one of the isomorphism (M3); (iv) the check count. No statement, label,
STATUS tag, schema key, golden line, digest or verdict moves.

v1.15.1-pre.5 --- 2026-09-30 (slot-15 WP-5, \textbf{a pre-release, not signed}; plan \texttt{docs/implementation/slot15/00-slot15-plan.md}
\S WP-5, owner ruling Q4 of 2026-09-30; the Fable referee \texttt{docs/implementation/slot15/WP-4\_referee.md} accepted P15 and P17 and
\HA{cor:stablecert} with the statement-scope repair B1): the paragraph ``The Hecke algebra over $\Z/\ell^{K}$, read mod~$\ell$'' at the end of
\S\ref{sec:ring} --- its item~(ii) now states the \emph{certified} verdict (\HA{cor:stablecert} in B1's form: the block is that of a maximal
ideal of $\mathbb{T}'=\Z_\ell[T_p:p\in S]$, possibly a sum of blocks of the full algebra), replacing ``readings, not certificates'' \emph{for this verdict
only}; item~(i) and a new item~(v) take the route beyond $2^{64}$ (the digit route); items (iii), (iv) and the Appendix-A row follow. Placed after
WP-19's text (plan \S6 document order). No statement of another claim, no STATUS tag, schema key, record key, golden line, digest or census verdict
moves; the \texttt{LOEWY\_RING} lines gain one appended field \texttt{stable\_cert}, their old text byte-identical.

v1.15.1-pre.4 --- 2026-09-30 (slot-15 WP-19, \textbf{a pre-release, not signed}; plan \texttt{docs/implementation/slot15/00-slot15-plan.md}
\S WP-19, owner ruling Q15 of 2026-09-30): two new claims after Claim~\ref{cl:sturmmodl} --- Claim~\ref{cl:mkfaithful} (the forms' Hecke
algebra mod~$\ell$ has dimension $\dim M_k$ on $M_k(\Fl)$ at every $\ell\ge5$, including $(\ell-1)\mid k$; the $M_k$ claim is \emph{proved}, not
weakened) and Claim~\ref{cl:formsF9} (the F9 identity with its exact torsion term, which gives the torsion guard of plan WP-20) --- and one
paragraph after Claim~\ref{cl:sturmmodl} recording that $B'(k,\ell)$ exists at every $\ell\ge5$. Both proofs are refereed in the WP-4 window.
Placed before WP-5's text on the certified verdict (plan \S6 document order). No statement of another claim, no schema key, golden line,
digest or verdict moves. Appendix~A gains two rows.

v1.15.1-pre.3 --- 2026-09-30 (slot-14, \textbf{a pre-release, not signed}; the owner signed off proposal P12 on
2026-09-30, \texttt{docs/implementation/slot14/WP-D1\_binding\_doc\_proposals.md}): P12's cluster half --- one evidence sentence in
Prediction~\ref{pr:latticesplit} (the WP-M1 route on the first cluster, job \texttt{34551683}: $\m_A^2=0$ on all $1233$ non-reduced
$\Fl$-rational blocks beyond the census window, citing the data set \texttt{docs/matrix-engine/results/noctua2\_section\_m/}, published
2026-09-30), with the forms-side sentence of the WP-7 addendum, and the Appendix-A row. No statement, label, STATUS tag, schema key, golden
line, digest or verdict moves; the tag on Prediction~\ref{pr:latticesplit} stays \CONJ.

v1.15.1-pre.2 --- 2026-09-29 (slot-14 audit fixes, \textbf{a pre-release, not signed}; owner rulings Q-C and Q-D of 2026-09-29,
\texttt{docs/implementation/slot14/AUDIT\_contract\_currency\_2026-09-29.md} \S3.2): one definitional paragraph, ``The Hecke algebra over
$\Z/\ell^{K}$, read mod~$\ell$'', at the end of \S\ref{sec:ring} (E3 of the audit --- the objects, refusals and surface of
slot-14 WP-6's route, whose \texttt{stable} verdicts are \emph{readings, not certificates} until the statements P15 and P17 are refereed in
slot~15), one sentence at the $(d1)$ cost cap recording the owner's ruling that the frozen $u_\ell$ RSS model \texttt{u\_ell\_rss/v2} is
\emph{kept}, and one Appendix-A row. No statement of \S\ref{sec:census}, no label of another section, no STATUS tag, schema key, golden line,
digest or verdict moves; the tag on Prediction~\ref{pr:latticesplit} stays \CONJ.

v1.15.1-pre.1 --- 2026-09-28 (slot-14 references edit, \textbf{a pre-release, not signed}): one citation and
one check, at \S\ref{sec:atlas} --- Citro--Ghitza's count of the level-one eigensystems mod~$p$ against $\Sigma(5)$,
$\Sigma(7)$, $\Sigma(11)$, $\Sigma(13)$ and the stored $\Sigma(691)$, which agree --- and its bibliography entry. No statement,
label, STATUS tag, schema key, golden line, digest or verdict moves;

v1.15.0 --- 2026-09-28, \textbf{owner-signed release} (hecke-extensions slot 14, package WP-13, owner decision
2026-09-28 ``Go for publication and release'', slot-14 plan \S9 Q2 and Q3 with the recommended options). It applies the routing of
slot-14 plan \S8 (\texttt{docs/implementation/slot14/WP-D1\_binding\_doc\_proposals.md}) to this document. Every change is evidence, a
pointer or provenance; each is marked ``v1.15.0'' at its site. \textbf{P5}, with the local half of \textbf{P12}: Prediction~\ref{pr:latticesplit}'s
evidence gains a third, independent route (slot-14 WP-M1: $1229$ blocks in $250$ cells, all $259$ non-reduced with $\m_A^2=0$, and
the $\Delta$-block of $(400,7)$), slot-14 WP-6's C++ route over $\Z/\ell^{K}$ ($452/452$ on $348$ cells), and the sentence ``the prediction
concerns the reduction of the module, not of the algebra'' ($A_K\otimes\Fl$ reaches Loewy length $3$--$6$ on census rows and exactly
$8$ at $(400,7)$). \textbf{P13}, which supersedes P6: the $206$ Dickson-type blocks read by WP-6 are all reduced, so a non-reduced
Dickson-type block has not occurred and the extrapolation stays untested. \textbf{P11}: edition~2's split lines, the $\theta$-image
table and the forms census were reproduced byte for byte on a second machine. \textbf{P4}: a pointer from (d5) to the $\Z/\ell^2$
depth window beyond $n=64$, in the extract \texttt{docs/matrix-engine/results/noctua2\_run1/}, published on the owner's sign-off.
\emph{Not applied here}: P9, which goes to the \texttt{mixed} proposal document; P10 and P16, which go to \texttt{ENGINE\_API.md}; and
P12's cluster half (the section-M data, job \texttt{34551683}, unpublished), deferred to slot~15. No statement, label, STATUS tag,
schema key, golden line, digest or verdict moves, and the tag on Prediction~\ref{pr:latticesplit} stays \CONJ. Companion pointers:
[C] v1.14.0, [P] v3.9.3. Release note \texttt{docs/implementation/slot14/WP-13\_release\_note.md});

v1.14.0: 2026-09-27, \textbf{owner-signed release} (hecke-extensions slot 13, package WP-15
close, owner item~2 of the slot-13 close: \emph{released} from v1.14.0-pre.3 with no further edit --- the three
pre-releases' text stands as published, the \texttt{pre.N} suffix dropped from the version string at every site
it appears. What v1.14.0 therefore contains, gathered in one place: (pre.1, slot-13 WP-7, FIN\_audit G2 closed)
Certificate~\ref{cert:boundary}'s per-ideal \texttt{boundary\_data} (\texttt{rank\_r1}, \texttt{r1}, \texttt{t})
present iff (d1) passed, \texttt{verify()} code~$32$ re-deriving (H1) and (H2) from the record, and the schema
\texttt{hecke\_ext/v4}~$\to$~\texttt{v5} (golden R2: the $62$ \texttt{ext\_json\_*} lines, the schema string
only); (pre.2, slot-13 WP-9) Prediction~\ref{pr:latticesplit} (\CONJ, candidate C-F9): $\m_A^2=0$ on every
$\Fl$-rational $H^1$ block at $\ell\ge5$, held on all $654$ non-reduced census rows of edition~2, with the
Serre-weight count recorded as \emph{not} explaining it; (pre.3, slot-13 WP-11, answering the referee's
finding~M4) the evidence sentence of Prediction~\ref{pr:latticesplit} corrected: the evidence is entirely
non-Dickson (all $654$ tested rows \texttt{mixed} or \texttt{cuspidal}), so the extension to a non-reduced
Dickson-type $\m$ is an untested extrapolation, while the \emph{equivalence} $\m_A^2=0 \Leftrightarrow$ nearly
split itself needs no such hypothesis. No statement, label, STATUS tag, schema key, golden line, digest or
verdict moves beyond what pre.1--pre.3 already record; the tag on Prediction~\ref{pr:latticesplit} stays \CONJ.
Built clean: \texttt{pdflatex}~$\times\,3$, $0$ undefined references, $0$ undefined citations, $0$ multiply-defined
labels. Companion pointers: [C] v1.13.0, [P] v3.9.3, [X] v1.18.1. Report
\texttt{docs/implementation/slot13/WP-15\_report.md} (the slot close) and this release note);

v1.14.0-pre.3: 2026-09-27 (hecke-extensions slot 13, package WP-11, \textbf{the referee's
repair of Prediction~\ref{pr:latticesplit}}: finding~M4 of \texttt{docs/implementation/slot13/WP-11\_referee.md}).
One sentence of Prediction~\ref{pr:latticesplit} is replaced (marked ``v1.14.0-pre.3'' at its site), and the
\S\ref{sec:census} pointer and the Appendix-A row say the same: the \emph{equivalence} with $\m_A^2=0$ needs no
hypothesis ``$\m$ not of Dickson type'', but the \emph{evidence} is entirely non-Dickson --- all $654$ tested rows are
\texttt{mixed} ($567$) or \texttt{cuspidal} ($87$), with a free lattice and \texttt{sturm\_proof\_range = true} ---
and every Dickson-type block computed so far is reduced ($\m_A=0$: $340$ census rows and the $100$ Dickson-type
blocks of slot-13 WP-11's scan), so the extension to a non-reduced Dickson-type $\m$ has \emph{zero} evidence and
is an untested extrapolation. The tag stays \CONJ. No claim, certificate, schema key, golden line, digest or
verdict moves; no label or STATUS tag is removed or renamed (\texttt{docs/implementation/slot13/WP-11\_response.md});

v1.14.0-pre.2: 2026-09-27 (hecke-extensions slot 13, package WP-11, \textbf{the C-F9
prediction}, owner answer Q8 of the slot-13 plan: ``an [E] \CONJ{} remark if it holds on all rows'' --- it held on
all $654$; a non-normative addition to the $X$-type pre-release, marked ``v1.14.0-pre.2'' at its site).
(1)~\S\ref{sec:expect}: Prediction~\ref{pr:latticesplit} (\CONJ): at $\ell\ge5$, on every local block $V=H^1(\Ga;\Mn)_\m$
with residue field $\Fl$, the local Hecke algebra has $\m_A^2=0$ --- equivalently the lattice
$H^1(\Ga;\Sym^n\Z_\ell)_\m$ is ``nearly split'' --- with its evidence (census edition~2, $654/654$ non-reduced
$H^1$ rows, slot-13 WP-9, \texttt{HeckeCensusLatticeValidation}) and the record that the Serre-weight count
does \emph{not} explain it (Loewy length~$2$ on every row; the bound by the number of distinct weights fails on
$89$ rows). (2)~\S\ref{sec:census}: a pointer to it. No claim, certificate, schema key, golden line, digest or
verdict moves; no label or STATUS tag is removed or renamed. Report
\texttt{docs/implementation/slot13/WP-9\_report.md} \S5, WP-11 in \texttt{docs/implementation/slot13/});

v1.14.0-pre.1: 2026-09-26 (hecke-extensions slot 13, package WP-7, \textbf{spec first:
(d1) re-derivable from the record} (slot-12 FIN audit G2's repair, owner-approved 2026-09-26, slot-13 Q3), an
$X$-type change marked ``v1.14.0-pre.1'' at its sites. (1)~Certificate~\ref{cert:boundary}: on an ideal with
\texttt{d1 = passed} the record now carries the data (H1) and (H2) are decided on --- $\rank(r^1|_V)$, the
restriction $r^1|_V$ and the $\ell-1$ restrictions $t_p|_V$, $p\in R(\ell)$ --- and \texttt{verify()} re-derives
(d1) from the record as it re-derives (d0) (new code~32); the v1.13.1 sentence ``(H1) and (H2) in-process
only'' is superseded. (2)~\S\ref{sec:schema}: the record becomes \texttt{hecke\_ext/v5} (item~(j), the one
new key \texttt{boundary\_data}, present iff \texttt{d1 = passed}). (3)~The golden event (slot-13 R2): the
$62$ \texttt{ext\_json\_*} lines (the schema string only --- under the harness's \texttt{u\_ell = off} no (d1)
passes, so no \texttt{boundary\_data} is emitted there);
\texttt{docs/matrix-engine/rebaseline-2026-09-26-hecke-ext-v5-d1-data.md}. No claim, label or STATUS tag is
removed or renamed; no verdict moves. Report \texttt{docs/implementation/slot13/WP-7\_report.md});

v1.13.1: 2026-09-26 (hecke-extensions slot 12, S12-FIN, \textbf{the Fable audit's repairs}
(\texttt{docs/implementation/slot12/FIN\_audit.md}, 1 M / 3 m), a $Y$-type version marked ``v1.13.1''
at its sites: Certificate~\ref{cert:boundary} gains the representative stream's fixed cost cap (G1; the new
closed-list value \texttt{not\_run:budget}, tied to the \texttt{u\_ell} budget, no settings key) and states what
\texttt{verify()} re-derives for (d1) (G2: structure only; (H1), (H2) in-process); \texttt{boundary\_reps}'
condition corrected to the implemented one. No golden line, digest or census verdict moves);

v1.13.0: 2026-09-26, \textbf{owner-signed 2026-09-26} (hecke-extensions slot 12,
package S12-FIN: \textbf{the finality certificates land}, an $X$-type version, every site marked
``v1.13.0''; the owner-approved golden changed-line event and record schema event of 2026-09-26).
(1)~\S\ref{sec:certgen}: Certificates~\ref{cert:dickson} (d0) ``dickson'' on $X_1$ (Proposition~D of the
slot-12 finality proposal v2, from \HA{prop:heckeH0}(1)) and~\ref{cert:boundary} (d1) ``boundary'' on $X_4$
(Proposition~F$'$, from \HA{prop:bdrytranscription} and \HA{lem:bdrytorus}, which [C] v1.12.0 proves ---
so (d1) is \emph{unconditional} and may make a record \texttt{final = true} in the default
\texttt{closure} mode), with statements and proofs, run only where (a)--(c) do not pass; the practical
order extended. (2)~\S\ref{sec:schema}: the record becomes \texttt{hecke\_ext/v4} --- item~(i):
\texttt{certificate\_routes.d0} / \texttt{.d1}, the per-ideal \texttt{boundary\_reps}, the certificate names
\texttt{"dickson"} / \texttt{"boundary"} (also admitted after \texttt{certificate\_failed:}); the definition
of \texttt{final = true} extended to the two routes; \texttt{verify()} codes~30--31. (3)~The golden event:
the $62$ \texttt{ext\_json\_*} lines (schema and the new keys; under the harness's \texttt{u\_ell = off}
no (d0) / (d1) can pass, so no \texttt{final} and no raw object moved;
\texttt{docs/matrix-engine/rebaseline-2026-09-26-hecke-ext-v4-finality.md}). No claim, label or STATUS tag
is removed or renamed. Companion pointers: [C] v1.12.0, [P] v3.9.3, [X] v1.17. Report
\texttt{docs/implementation/slot12/FIN\_report.md});

v1.12.1: 2026-09-25 (hecke-extensions slot 11, package WP-17, the slot close: \textbf{the
WP-15 Fable referee's repairs} (\texttt{docs/implementation/slot11/WP-15\_referee.md}, 0 B / 2 M / 6 m), a
$Y$-type version, every site marked ``v1.12.1''. (1)~\S\ref{sec:ring}, ``Published v1.12.0'': the
\texttt{hida} record object's sixth key \texttt{"e"} named (M-2). (2)~\S\ref{sec:thetaobj}~(h): \textbf{the
specification stands} --- $\ell<5$ is a mathematical precondition, exit~$3$ --- and the build now follows it
(\texttt{DicksonThetaCell}'s \texttt{ENGINE\_CHECK\_MSG} naming \HA{lem:h2torsion}, slot-11 WP-17; the ruling
on F-2, m-1). (3)~Prediction~\ref{pr:eiscusp}: which ``cuspidal'' the reworded statement means (m-6). No
claim, label, STATUS tag, record member (the key \texttt{"e"} was already emitted), golden entry or digest
moves. Companion pointers: [C] v1.11.1, [P] v3.9.1, [X] v1.15. Report
\texttt{docs/implementation/slot11/WP-17\_report.md});

v1.12.0: 2026-09-25 (hecke-extensions slot 11, package WP-15: \textbf{v1.12.0
published}, an $X$-type version, from v1.12.0-pre.2 with what slot~11 built, every site marked
``v1.12.0''. (1)~\textbf{\S\ref{sec:ring}}: the paragraph ``Published v1.12.0'' --- the pre-release
paragraphs as built, member by member: the $\theta$ and census members of~(b1)'s closure with the
guards they actually carry, the console's $\theta$ verbs \emph{not built} (slot-11 WP-16 gave way),
the native \texttt{u128} layer as the third width and oracle, the digit route with its caps checked at
$K=k+K_w-1$, the refusal at $\ell\ge2^{31}$, the \texttt{hida} record object, and (h)'s field members
exiting~$3$ on a digit handle (slot-11 WP-14 pass~2, M-1); the console's \texttt{uint32\_t} gap
below $2^{64}$ recorded as an owner item. (2)~\textbf{\S\ref{sec:thetaobj}} and item~6 of
\S\ref{sec:identities}: as built, with the refusal at $\ell\in\{2,3\}$ being exit~$5$, not~$3$.
(3)~\textbf{\S\ref{sec:census}}: the catalogue as built --- the \texttt{t\_p=<fnv64>} field of
\texttt{CENSUS\_SPLIT}, the cell-exact acceptance, the published edition~1 and the W2 data set
\texttt{hecke\_census\_dataset/v1} in place of edition~2 (slot~12). (4)~\textbf{Prediction~\ref{pr:eiscusp}
reworded}: read literally, its ``iff'' is false at $(10,7)$ (Ramanujan's $\tau(n)\equiv n\sigma_3(n)\bmod7$);
it is now stated for $E_k$'s own class, with the exceptional residues named ($w\equiv0,2$). (5)~\S\ref{sec:tests}:
the uniqueness item at $(198,691)$ is \textbf{closed} (one congruent eigenform; degrees $1,1,2,2,10$), and
OQ13's residue discrepancy is the sign ($|\num|\equiv217665$, signed $\equiv259816$). No claim, label or
STATUS tag is removed or renamed; \ref{pr:eiscusp} keeps \STD{} for the restricted statement. No record
member, golden entry or digest moves. Companion pointers: [C] v1.11.0, [P] v3.9, [X] v1.14.1. Report
\texttt{docs/implementation/slot11/WP-15\_report.md});

v1.12.0-pre.2: 2026-09-25 (hecke-extensions slot 11, package WP-18c, \textbf{spec
first}, drafted by WP-18c's author for WP-4's editor row: the second pre-release of v1.12.0, an
$X$-type change marked ``v1.12.0-pre.2'' at its sites, under the owner's clarification of
2026-09-25 (latest; slot-11 plan \S10 last row): \emph{arbitrary $k$ by the (C2) digit induction,
in slot~11}. It supersedes pre.1's sentence ``(a)'s grammar $k\le63$ does not move in slot~11''.
\textbf{\S\ref{sec:ring}}: (a)'s grammar cap $63\to4096$; the per-cell rule (native below $2^{64}$,
the digit route at and above it for $\ell\ge5$, $\ell\in\{2,3\}$ refused there by name, a
$8192$-bit and a memory cap); the digit route's member list, its per-cell certified precision
bound $K=k+e_{\max}$ with the record keys that state it, the ordinary idempotent as a certified
datum, and the engine-family switch \texttt{ENGINE\_LADIC\_WIDTH}; (g)'s digit jobs; (h)'s digit
handles. No existing value, record member, golden entry or digest moves: every $\ell^k<2^{64}$
keeps the shipped route. Report \texttt{docs/implementation/slot11/WP-18c\_report.md});

v1.12.0-pre.1: 2026-09-25 (hecke-extensions slot 11, package WP-4, \textbf{spec
first}: the first pre-release of v1.12.0, with its own identifier so that no version number is
published twice under different text ([C] V1\,/\,V2), written before any slot-11 code package
commits a claim (slot-11 plan \texttt{docs/implementation/hecke\_extensions\_slot11\_plan.md}
\S3 ``WP-4'', rule~21) and under the owner's rulings of 2026-09-25 on every slot-11 decision ---
all at recommendation, decision~3 at option~(A$'$) (that plan's \S10); nothing below is
provisional. An $X$-type change; every addition is marked ``v1.12.0-pre.1'' at its site.
(1)~\textbf{\S\ref{sec:thetaobj}} (new): the Dickson--$\theta$ objects slot-11 WP-5 and WP-6
build --- the indexing, the quotient local system $Q_n$, the six terms and five maps of the long
exact sequence with their code names, what is certified and what is asserted, the hand-off file
to the census, the refusals, and what they are \emph{not} (no record member, no settings key, no
golden line, no certificate item code). (2)~\textbf{\S\ref{sec:identities} item~6}: its $L\cdot$
half (E1 identity class~5) extended by the identities (6a)--(6h) of that long exact sequence.
(3)~\textbf{Claim~\ref{cl:Tlnil}}: split into its two sentences; the first, $U_\ell\circ\theta=0$
over $\Fl$ (already on cochains), upgraded from \STD{} to a pointer at [C]'s proof,
\HA{prop:thetatwist} in \HA{sec:dicksontheta} (slot-11 WP-2), the second (the Jordan type) kept
\STD; and, with [C]'s correction \HA{rmk:heckeH0plan}, (6f) no longer says ``$U_\ell=0$ on all of
$H^0$'': $U_\ell$ has rank $\le1$ and square zero there, and is zero only for $n<\ell^3-\ell$
(\HA{prop:heckeH0}; $U_5D^6=-L^{20}$ at $(120,5)$); ``Dickson and $\theta$-images are
non-ordinary'' stands (\HA{cor:thetanonordinary}). (4)~\textbf{\S\ref{sec:census}}
(new): the congruence-census catalogue, a derived artefact with its own format tag
\texttt{hecke\_census/v1} --- its inputs and supersession rule, its rows, its columns each with
its source, the boundary split rule (Eisenstein\,/\,cuspidal\,/\,mixed by $\rank r^1$ on the
ideal's generalised eigenspace), the $\pm$, $\theta$-image, Glover-reconciliation and
$\Z/\ell^k$-depth rules, the extremal-example criteria [EX] selects from, the line formats of the
two computed columns, and the \texttt{results/} layout. (5)~\textbf{\S\ref{sec:ring}}: the
pre-release paragraph ``v1.12.0-pre.1'' --- the slot-11 members of~(b1)'s closure, all refusing a
non-default ring by name ($\theta$ and the census are $\Fl$-only in slot~11), and L1-X under
decision~3~(A$'$): step~1 (the $128$-bit engine element type with its oracle rule, and F-1, F-2) is
built in slot~11 and is value-neutral; \textbf{(a)'s grammar $k\le63$ does not move in slot~11};
the amendment $k\le126$ with the cell refusal at $\ell^k\ge2^{127}$ is slot~12's. No claim,
hypothesis, formula or label is removed or renamed and no STATUS tag moves except
Claim~\ref{cl:Tlnil}'s; no member of \texttt{hecke\_ext/v3}, \texttt{hecke\_report/v2} or
\texttt{hecke\_ring\_report/v1} moves; no settings key is added (eleven stay eleven); no gate, cell
list, golden entry or digest moves. Companion pointers: [C] v1.11.0-pre.1 (slot-11 WP-2), [X]
v1.14.1, [P] v3.9-pre.1. Report \texttt{docs/implementation/slot11/WP-4\_report.md});

v1.11.1: 2026-09-25 (hecke-extensions slot 10, packages WP-4 and WP-5, window
\textbf{W2}: a $Y$-type status change. The owner lent the box for W2 after the close window W1
(consent of 2026-09-25), and the two optional publications of decision~8~(A) ran there: slot-10
WP-4's $\ell$-adic survey under $\Z/\ell^2$ (\texttt{docs/matrix-engine/results/ring\_survey\_l2/},
$462$ jobs, $0$ failed) and slot-10 WP-5's second extension campaign
(\texttt{docs/matrix-engine/results/ext\_u691/}, the $45$ cells $\ell=691$, $n=112,\dots,200$ that
\texttt{ext\_first} recorded \texttt{u\_ell\_over\_budget}, $0$ failed). (1)~\textbf{The survey}
(\S\ref{sec:ring}, the status paragraph, and (g)): v1.11.0's ``not run in slot~10 (deferred)''
becomes ``run in W2'', with what it found; (g) still moves by nothing --- the survey wrote
\texttt{hecke\_ring\_report/v1} exactly as frozen, and no finding needs a member. Every change is
marked ``v1.11.1'' at its site. No claim, hypothesis, formula, STATUS tag, label, gate, cell list,
schema member or constant moves; no golden entry and no digest moved. Companion pointers: [C]
v1.10.1, [X] v1.14.1, [P] v3.8. Reports \texttt{docs/implementation/slot10/WP-4\_report.md},
\texttt{WP-5\_report.md});

v1.11.0: 2026-09-25 (hecke-extensions slot 10, package WP-15,
\textbf{phase~2}: publishes v1.11.0, the first text under this number. It is v1.11.0-pre.2 with
every site phase~1 marked filled from the reports of slot-10 WP-7 and WP-8, and with the phase-1
Fable referee's repairs that reach this document
(\texttt{docs/implementation/slot10/WP-15\_referee\_phase1.md}: M-1 and m-5 in (b3)). An $X$-type
change, as the pre-releases announced; every addition is marked ``v1.11.0'' at its site. No claim,
hypothesis, formula, STATUS tag or label is added, renamed or removed; no gate, cell list, schema
member or constant of an earlier version moves. Base \texttt{47e3ff9}: every slot-10 code package
merged. (1)~\textbf{Lane~B} (\S\ref{sec:ring}, the status paragraph): slot-10 WP-7's two levers ---
the $\iota$ sign split's independent checks and splits as host tasks, and L4, route~2's rows of
\texttt{HeckeOperatorOnK} on the consumer pool --- recorded as value-preserving with their evidence,
and its cone-certificate lever recorded for slot~11; slot-10 WP-8's ninth fingerprint bit (a cache
key) and its validation carries. (2)~\textbf{(b1)}: no lever in the cone classes or in
\texttt{HeckeSignSplit}; L4 and the $\iota$ lever located in the closure. (3)~\textbf{(f)}: route~3's
non-vacuity asserted (slot-10 WP-8, review m-12), \texttt{HeckeConeCrossRouteRingValidation}'s count
$386\to431$ by the source, re-derived at the close. (4)~\textbf{The survey} (status paragraph and
(g)): slot-10 WP-4 not run in slot~10 (the optional W2; deferred). (5)~\textbf{(b3)}: ``has always
computed'' becomes ``computes'', with the canonical $\tau$ dated to [C] v1.1.0 and the
legacy--canonical observation of \HA{rmk:tauexistsscope}(1) cited (referee M-1, m-5).
Companion pointers: [C] \textbf{v1.10.1}, [X] v1.13, [P] v3.8. Report
\texttt{docs/implementation/slot10/WP-15\_report.md});

v1.11.0-pre.2: 2026-09-25 (hecke-extensions slot 10, package WP-15,
\textbf{phase~1}: the second pre-release of v1.11.0, which records at the sites of
v1.11.0-pre.1 what slot-10 WP-3, WP-10, WP-11, WP-12 and WP-13 built (slot-10 plan \S3 ``WP-15'').
It carries its own identifier, so that no version number is published twice under different text
([C] V1\,/\,V2; slot-10 plan \S4.3): the packages of the slot's lane~B (WP-7, the cone and
$\iota$ profile with its optional levers, and WP-8, the carries) were still running when it was
written, and the sites they may complete were marked in its text, in italics, as bracketed
phase-2 sites (the markers are gone since v1.11.0, which filled them); phase~2 fills them and
publishes v1.11.0. An $X$-type change,
as the pre-release announced; every addition is marked ``v1.11.0-pre.2'' at its site and names the
validation that measures it. No claim, hypothesis, formula, STATUS tag or label is added, renamed
or removed; no gate, cell list, schema member or constant of an earlier version moves. Base
\texttt{f0640d0}: WP-2, WP-3, WP-6, WP-9, WP-10, WP-11, WP-12 and WP-13 merged, and the slot's
independent review (\texttt{docs/implementation/slot10/WP-14\_review.md}, passes 1 and~2).
(1)~\textbf{Status} (\S\ref{sec:ring}, new paragraph ``Status at v1.11.0-pre.2''): the three
members of \HA{obl:coneexact}(vii) over $\Z/\ell^k$ are built as specified (slot-10 WP-3), and
\textbf{E8's declaration is written, scoped by the review's finding M-1}: (vii)'s first two
sentences discharged at every E8 point ($\ell^k<2^{32}$), its third --- the Poincar\'e-duality
adjointness gate --- recorded unavailable over the ring, since the E7 pairing and the contragredient
refuse a composite ring; the quoted declaration of v1.11.0-pre.1 is amended accordingly.
(2)~\textbf{(b1) as built}: the cross-routes' guard and regime refusal, the changed rows of
arms~(X)/(X4b); line~2's free model (WP-10, WP-13, WP-11) outside the closure as specified, the
read-only use of \texttt{HeckeOperatorOnK} on the strict table at $k\ge2$ confirmed, and the
functoriality argument stated with its three inputs (review finding m-1); the console members.
(3)~\textbf{(b3)}: the decision-4 block as built --- nothing lifted, no line-2 package exhibited two
$\tau$ with different operators, and [C] v1.10.0's \HA{rmk:bredonhecke} carries the reason.
(4)~\textbf{(f) as built}: route~3's transport by an explicit inverse on classes instead of
column-by-column solves (a recorded deviation, verified by the review), the acceptance as measured,
the negative control's scope, the regime arm; the pointer into [C] closed.
(5)~\textbf{(h) as built} (slot-10 WP-12): the five members, the mismatch rule pinned both ways,
the member-getter position, and one exit code that moves by the Core exception's own rule --- an
unwrapped getter raising \texttt{NotImplementedInThisBuild} is now exit~$5$, not~$3$.
(6)~The two pointers into [C] of v1.11.0-pre.1 are closed by [C] \textbf{v1.10.0}.
Companion pointers: [C] \textbf{v1.10.0}, [X] v1.13, [P] \textbf{v3.8}. Report
\texttt{docs/implementation/slot10/WP-15\_report.md});

v1.11.0-pre.1: 2026-09-25 (hecke-extensions slot 10, package WP-2, \textbf{spec
first}: the pre-release of v1.11.0 that slot-10 WP-15 completes. It carries its own identifier
\texttt{1.11.0-pre.1}, so that no version number is published twice under different text ([C]
V1\,/\,V2; slot-10 plan \S4.3). An $X$-type change, written \emph{before} the code packages that
honour or refuse what it specifies (slot-10 WP-3, WP-10, WP-12, WP-13), under the owner's rulings
of 2026-09-25 on the slot-10 decisions, every one at recommendation~(A) (slot-10 plan \S10);
decisions 3, 4, 6 and~8 are the ones used here, and \textbf{nothing below is provisional}. No claim,
hypothesis, formula, STATUS tag or label is added, renamed or removed; no gate, cell list, schema
member or constant of an earlier version moves. Every addition is marked ``v1.11.0-pre.1'' at its
site, in \S\ref{sec:ring}; the paragraph ``Pre-release v1.11.0-pre.1'' there is its index.
(1)~\textbf{(b1) and (f): \HA{obl:coneexact}(vii) over $\Z/\ell^k$} (decision~3~(A)). Three members
move from ``keeps refusing'' to ``honours'' at $\ell\ge5$, $\ell^k<2^{32}$, from slot-10 WP-3's
lift commit: \texttt{HeckeConeStrictPairRoute}, \texttt{HeckeConeRelativeRoute} and
\texttt{EquivariantCohomologyOnK}'s relative sub-complex constructor, on the dense module route
of~(f) with~(c1)--(c5) verbatim; the acceptance is entrywise equality of each cross-route's
$A_c^{(q)}$ with route~1's in the cyclic-summand basis of~(f), at the seven E8 points of~(e); above
$2^{32}$ both cross-routes keep the ring context's named regime refusal (the \texttt{uint32\_t}
regime rule of~(b1)), and route~1 alone runs every regime; $\ell\in\{2,3\}$ keeps refusing. The
declaration E8 reaches on acceptance is stated, for slot-10 WP-15 to write.
(2)~\textbf{(b3) and (c1$'$): the Hecke action on the Bredon object at $\ell\in\{2,3\}$ is
declined} (decision~4~(A)), for two separately named objects: the cochain map $\Theta$ on
$C^{\mathrm{eq}}(\Z/\ell^k)$ is neither built nor needed (line~2 reads $\widehat\Theta_g$ on
$\widehat C$, \HA{prop:isomodel}(3)); the operator on cohomology is \textbf{not shown} independent
of $\tau$ (\HA{lem:tauexists}(2), \HA{prop:strictlift}(4)), which is the absence of a proof, not a
proof of dependence. \texttt{HeckeOperatorOnK}'s refusal there stands.
(3)~\textbf{(b1): line~2's free model is outside the closure}: the objects of slot-10 WP-10 and
WP-13 (\texttt{Integral/Model/StrictLift}, \texttt{PeriodicityClass}) take their ring as an explicit
\texttt{ModCtx} and read no setting, as slot~9's \texttt{IsotropyModel} and \texttt{IntegralGroups}
do; the read-only use of \texttt{HeckeOperatorOnK} on the strict table inside a validation is an
honouring use of an honouring member, and a refusal it meets stands.
(4)~\textbf{(h): the structural members} (decision~6~(ii)~(A)): \texttt{LocalSystem.n},
\texttt{.d\_V} and a new \texttt{.coefficient\_ring}, and \texttt{HeckeOperator.p} and \texttt{.r},
honour on a ring handle from slot-10 WP-12's lift commit; every other unlisted member still
exits~$5$; slot-9 WP-12's finding ``a ring mismatch is exit~$5$'' is recorded as the rule.
(5)~\textbf{(g): no schema change} (decision~8): slot-10 WP-4's survey writes
\texttt{hecke\_ring\_report/v1} as frozen in v1.10.0.
(6)~\textbf{Pointers}: the staleness of [C] \HA{obl:coneexact}(vii)'s ``its composite half is
open'' and of \HA{sec:notimplemented}'s matching sentence from slot-10 WP-3's merge until [C]'s next
edit (slot-10 WP-15, [C] v1.10.0), recorded and not patched (slot-10 plan \S4.3).
Companion pointers: [C] v1.9.1, [X] v1.13, [P] \textbf{v3.8-pre.1}. Report
\texttt{docs/implementation/slot10/WP-2\_report.md});

v1.10.1: 2026-09-25 (hecke-extensions slot 9, package WP-17, the slot close:
\textbf{the Fable referee's repairs to v1.10.0}, a $Y$-type change --- no claim, hypothesis,
formula, STATUS tag, label, schema member, constant or rule moves, and every addition is marked
``v1.10.1'' at its site. (1)~\textbf{M-1}: the slot's independent review (WP-14, pass~1), which
v1.10.0 cites in (b1), (c1$'$) and (f), is now in the tree as \texttt{docs/implementation/slot9/WP-14\_review.md} (reconstructed from the
reviewer's returned summaries after the original file was lost on 2026-09-25; it keeps its
provenance note), and each of those three sites names it. (2)~\textbf{m-7}: (b1)'s scoping row says
that the scope guard also exempts its own translation unit, \texttt{Registry/Slot9/WP9.cpp}.
(3)~\textbf{m-8}: (f)'s length identity for $\beta^1_c$ is asserted at the \texttt{uint32\_t} cells,
not at the two word-regime cells. (4)~\textbf{m-10}: v1.10.0's entry below said ``nothing else of
v1.10.0-pre.2 changed''; it also carried the (g7), (g9) and (g11) pointers and the slot-9 WP-9 sites
its items~(2)--(5) list, which v1.10.0-pre.2 had announced only as ``any site WP-6/7/9 touch'' ---
the sentence is read with that clause (the entry itself is left as written). Companion pointers:
[C] \textbf{v1.9.1}, [X] \textbf{v1.13}, [P] \textbf{v3.7.1}. Report
\texttt{docs/implementation/slot9/WP-17\_report.md});

v1.10.0: 2026-09-25 (hecke-extensions slot 9, package WP-15,
\textbf{phase~2}: publishes v1.10.0, the first text under this number. It is v1.10.0-pre.2 with
every site phase~1 left open filled from the reports of slot-9 WP-6, WP-7, WP-8, WP-9 and WP-4t,
and nothing else of v1.10.0-pre.2 changed. An $X$-type change, as the pre-releases announced
(one schema member is amended before the schema is first published, item~(2)); every addition is
marked ``v1.10.0'' at its site and names the validation that measures it. No claim, hypothesis,
formula, STATUS tag or label is renamed or removed, and no label is added. Base \texttt{9b004f2}:
every slot-9 code package merged. The Fable referee reads this version; its repairs go to v1.10.1.
(1)~\textbf{Status} (\S\ref{sec:ring}, new paragraph ``Status at v1.10.0''): the campaign (C5c,
slot-9 WP-8) built as specified; the merge fix \texttt{158c0be}; the carries of slot-9 WP-9; the
trimming of slot-9 WP-4t; E8's declaration unchanged.
(2)~\textbf{(g) as built} (slot-9 WP-8): the as-built paragraph --- the dispatch, the element
width, the routes at $\ell\ge5$ and at $\ell\in\{2,3\}$, \texttt{CampaignRingValidation}'s arms
and the fingerprint-$0$ byte identity --- with \textbf{four recorded deviations}: \texttt{kmode}
\texttt{not\_applicable} also at $\ell\ge5$ off the $K$-mode's range (noted at (g7)); the
uncertified state's \texttt{need} (noted at (g9)); \texttt{hecke} absent and not refused at
$\ell\ge5$ outside (g11)'s condition; each refusal's \texttt{who} and \texttt{message} from the
real guard, a guard that stops refusing failing the job. \textbf{(g11) amended} by the lead's
decision of 2026-09-24: the free \texttt{charpoly}'s coefficients are \emph{decimal strings},
since at a \texttt{uint64\_t} cell they exceed $2^{53}$; no other member can.
(3)~\textbf{(b1) and (b4): the campaign rows as built} --- the job loop honours $\Z/\ell^k$ at every
$\ell$ with $\ell^k<2^{64}$; the ext pass and a cell with $\ell^k\ge2^{64}$ refused at spec
validation, before any output exists.
(4)~\textbf{(b1): review finding m-1 and the \texttt{RunPipelineModEll} guard closed} (slot-9
WP-9, \texttt{SmallEllPolicyScopeValidation}; the merge fix \texttt{158c0be}); the
constant-fibre constructor's composite caller of \texttt{get\_prime()} re-pointed; the item-4b
tiers as trimmed by slot-9 WP-4t ($(60,5)$ in the \texttt{slow} sibling; $458\to306$ checks).
(5)~\textbf{(c1$'$): review finding m-2 closed} (slot-9 WP-9,
\texttt{HeckeBredonU64FlatRouteValidation}, one honouring \texttt{uint64\_t} cell at $\Z/2^{33}$).
Companion pointers: [C] v1.9.0 (unchanged; its pointer already names this version), [X] v1.12,
[P] \textbf{v3.7}. Report \texttt{docs/implementation/slot9/WP-15\_report.md});

v1.10.0-pre.2: 2026-09-24 (hecke-extensions slot 9, package WP-15,
\textbf{phase~1}: the second pre-release of v1.10.0, which records at the sites of
v1.10.0-pre.1 what slot-9 WP-3, WP-4, WP-12 and WP-13 built (slot-9 plan \S3 ``WP-15'', its
two-phase start of \S4.1 Rule~C). It carries its own identifier, so that no version number is
published twice under different text ([C] V1\,/\,V2; slot-9 plan \S4.3): the packages of the
slot's lane~B (WP-6 \ldots\ WP-9, in particular WP-8's campaign support, C5c) were still running
when it was written, and the sites they complete were marked in its text for phase~2, which
filled them and published v1.10.0 (above); the Fable referee reads
v1.10.0. An $X$-type change, as the pre-release announced; every addition is marked
``v1.10.0-pre.2'' at its site and names the validation that measures it. No claim, hypothesis,
formula, STATUS tag or label is renamed or removed, and no label is added. Base \texttt{7ddb555}:
WP-2, WP-3, WP-4 (with its review repairs and the golden event R1), WP-10, WP-11, WP-12 and WP-13
merged.
(1)~\textbf{Status} (\S\ref{sec:ring}, new paragraph ``Status at v1.10.0-pre.2''): E8 item~4b is
built by route~1, and \textbf{E8 is complete as specified} --- items 1--4b as specified, 4b over
$\Z/\ell^k$ by route~1, \HA{obl:coneexact}(vii)'s composite half open (the two cross-routes and
the relative constructor refuse, with arms); R1 has landed ($126$ entries, $0$ changed, $0$
removed; tier~A $11\,666$, tier~B $5901$).
(2)~\textbf{(a)}: the \texttt{--set} exit~$5$ and the \texttt{set}-statement position, as built
by WP-12 (\texttt{ConsoleErrorValidation} arm~(K), $1\to5$).
(3)~\textbf{(b1), (b3), (b4) as built}: the scoped lift at $\ell\in\{2,3\}$; \textbf{the $K$-mode's
refusal kind there is \texttt{NotImplementedInThisBuild}} (exit~$5$), recorded as standing; the
4b closure as WP-4 re-derived it, with the \texttt{uint32\_t} width of the four classes; the
console surface (WP-12, WP-13).
(4)~\textbf{(c1$'$)}: its ``the pipeline's $H^2=0$ is asserted'' is corrected --- the route does
not assert it; \texttt{HeckeBredonTreeRingValidation} checks it (review finding m-3).
(5)~\textbf{(f) as built}: \HA{obl:coneexact}(vi) \textbf{narrowed} to what is built --- the
commutator asserted, the $\nu=2$ recursion not computed over the ring (review finding M-W4-1);
item~14's completeness argument (m-W4-2); the measured values ($m=1$ at $(16,5)$, the $c_2$);
the landed set.
(6)~\textbf{(h) as built}: the honouring surface, and WP-12's four findings --- member getters
carry no source position (a next Core batch), unlisted structural members exit~$5$ on a ring
handle, a ring mismatch is exit~$5$, and \texttt{.A(q)} over the ring is a List of rows.
(7)~The pointers into [C] of (b3) and (f) now name [C] v1.9.0, which repairs both clauses.
Companion pointers: [C] \textbf{v1.9.0}, [X] v1.12, [P] \textbf{v3.7-pre.1}. Report
\texttt{docs/implementation/slot9/WP-15\_report.md});

v1.10.0-pre.1: 2026-09-24 (hecke-extensions slot 9, package WP-2, \textbf{spec
first}: the pre-release of v1.10.0 that slot-9 WP-15 completes. It carries its own identifier
\texttt{1.10.0-pre.1}, so that no version number is published twice under different text ([C]
V1\,/\,V2; slot-9 plan \S4.3). An $X$-type change, written \emph{before} the code packages that
honour or refuse what it specifies (slot-9 WP-3, WP-4, WP-8, WP-12), under the owner's rulings of
2026-09-24 on the slot-9 decisions, every one at recommendation~(A) (slot-9 plan \S10); decisions
2, 3 (with its sub-item (iii)), 4 and 6 are the ones used here, and \textbf{nothing below is
provisional}. No claim, hypothesis, formula, STATUS tag or label is renamed or removed, and no label
is added; no gate, cell list or constant of an earlier version moves. Every addition is marked
``v1.10.0-pre.1'' at its site, in \S\ref{sec:ring} unless noted.
(1)~\textbf{(b3) and (c1) at $\ell\in\{2,3\}$} (decision~4; slot-9 plan audit B-1, M-4).
\texttt{EquivariantCohomologyOnK} honours $\Z/\ell^k$ at $2$ and $3$, through a
\emph{flat-coordinate} module route (new item~(c1$'$)), reporting lengths and invariant factors
of the Bredon object only; the certified $K$-mode, \texttt{HeckeOperatorOnK},
\texttt{HeckeReduction} and every field-only consumer keep refusing there. (c1)'s ``every pivot is
asserted to be a unit, so $C^q_{\mathrm{eq}}$ is free'' is restricted to $\ell\ge5$. The as-built
refusal arm ``through \texttt{CoefficientRingContext}, $\ell=2$'' of~(b3) is recorded as the arm
that \emph{flips} to honouring, and its sentence ``WP-2's choice \ldots\ not objected to'' is marked
superseded. A pointer records that [C] \HA{def:ringsemantics}(2)'s clause ``the implementation of
this build refuses it there'' is superseded by~(b3) until [C]'s next edit (slot-9 WP-15).
(2)~\textbf{The 4b objects} (decision~3, with~(iii)): new paragraph~(f), the boundary complex and
$\mathrm{Cone}(r)$ over $\Z/\ell^k$ --- their coordinates, their tier-A and tier-B objects with golden
names at the seven existing E8 points, the generating rule of slot-9 WP-4's pre-stated set
($126$ entries), predicted values where they follow from the mathematics, and
\HA{obl:coneexact}(vii) over the ring as route~1 only, its composite half recorded as open.
(3)~\textbf{The job-level schema} \texttt{hecke\_ring\_report/v1} (decision~6): new
paragraph~(g), emitted by the campaign's job loop only under a non-default ring; a
fingerprint-$0$ job emits \texttt{campaign\_job/v1} with \texttt{hecke\_report/v2} byte for byte
as before, and the ext pass refuses a non-default ring by name at spec validation.
(4)~\textbf{The console under a non-default ring} (C5a, decision~2): new paragraph~(h), the verbs
and members that honour $\Z/\ell^k$, and~(a)'s ``Registration'' sentence rewritten as ruled: a
legal but unexecutable \texttt{--set} value is exit~$5$ from slot-9 WP-12's commit.
(5)~\textbf{(b1): the lift list}, one row per member that moves in slot~9, with its package and
its arm, and the new ``honours at $2$ and $3$'' row; (b4)'s surfaces restated for slot~9.
Companion pointers: [C] \textbf{v1.8.1}, [X] \textbf{v1.12}, [P] v3.6.1,
\texttt{ENGINE\_API.md} \S13 (slot-8 close, [WP8-14] \S8 item~7). Report
\texttt{docs/implementation/slot9/WP-2\_report.md});

v1.9.1: 2026-09-24 (hecke-extensions slot 8, package WP-13, the repairs of its
Fable referee pass, \texttt{docs/implementation/slot8/WP-13\_report.md} \S7. A $Y$-type change:
\textbf{pointers and wording only}. No claim, hypothesis, formula, STATUS tag or label is added,
renamed or removed, and no gate, cell list, constant or rule moves. Every repair is marked
``v1.9.1'' at its site.
(1)~\textbf{\S\ref{sec:ring}~(b1), the ``Unreachable'' bullet} (finding M-2). Its reason for
\texttt{kernel.cu}'s legacy benchmark said that no command-line switch sets a mathematical key
outside \texttt{--console} and \texttt{--campaign}. That is false: \texttt{--bench ext --u-ell}
installs \texttt{u\_ell}. The reason now given is that no switch sets \texttt{coefficient\_ring},
and that \texttt{--bench} refuses a non-default ring at \texttt{RunMode}'s entry.
(2)~\textbf{(a) and item~(2) of v1.9.0's entry: the \texttt{--set} exit code} (findings M-1,
m-2). Exit~$1$ for \texttt{--set coefficient\_ring=Z} was attributed to arms that did not test
it. It is now pinned by one check added to arm~(K) of \texttt{ConsoleErrorValidation}, and was
re-measured with a correct command line.
(3)~\textbf{(a): the emission and transcript facts} (finding m-1) are attributed to slot-8 WP-3's
gate runs, not to the arms.
(4)~\textbf{(c1): \HA{prop:coneintegral}(1)--(2)}, not~(2) alone (finding m-3).
(5)~\textbf{[IB] is defined} in \S\ref{sec:ring}'s preamble (finding m-4).
(6)~\textbf{(b3): ``the two refusals'' listed four}; it now reads ``the honouring object's
refusals'' (finding m-7).
Companion pointers: [C] v1.8.0, [X] v1.11, [P] \textbf{v3.6.1});

v1.9.0: 2026-09-24 (hecke-extensions slot 8, package WP-13. This revision
completes v1.9.0-pre.1 and publishes v1.9.0. It is an $X$-type change, as the pre-release
announced. No claim, hypothesis, formula, STATUS tag or label is added, renamed or removed. No
gate, cell list, constant or rule moves. One predicted literal is restated in model-free form,
item~(7). \S\ref{sec:ring} now records what was built against it,
at the sites it specifies. Each addition is marked ``v1.9.0'' and names the validation that
measures it. Base \texttt{5c70e95}: every slot-8 code package merged. Report
\texttt{docs/implementation/slot8/WP-13\_report.md}.
(1)~\textbf{E8 status} (\S\ref{sec:ring}, new paragraph ``Status at v1.9.0''). Items~2 and~3 are
implemented (slot-8 WP-1, WP-3), with the repairs from the independent review (WP-12 pass~2;
WP-3r). Item~4a is accepted (WP-4). R1 has landed: $124$ entries at the seven points, $0$
changed, $0$ removed, tier~A $11\,575$, tier~B $5866$, both digests unchanged. Item~4b is slot~9's.
Measured by \texttt{HeckeReductionValidation}, \texttt{HeckeReductionCoverageSlowValidation} and
\texttt{GOLDEN\_STRICT=1 GoldenRegressionValidation}.
(2)~\textbf{(a) as built.} The key is registered as specified, the eleventh key. At the console
the reserved values give exit~$5$ on the \texttt{set} statement (arm~(K) of
\texttt{ConsoleErrorValidation}; arm~(X)) and exit~$1$ as a \texttt{--set} argument (v1.9.1:
pinned by arm~(K); this entry first attributed it to both arms, which did not test it). Line~2's record keys are emitted
only under \texttt{Z}; the owner confirmed this (slot-8 ruling~6d).
(3)~\textbf{(b1): the list.} The closure is copied from slot-8 WP-3's derivation and re-derived
twice: by the review (WP-12), which added
\texttt{Integral/IntegralCharpolyCrt} (finding m-5; it refuses structurally), and at this
revision's base, which added the campaign's job loop, covered by its surface.
(4)~\textbf{(b3) and (b4) as built.} No \texttt{--bench} switch was added. The U64\_Wide range is
honoured and exercised (WP-3r, finding M-2). The refusal at $\ell\in\{2,3\}$ now points at [C]
\HA{def:ringsemantics}(2) and \HA{prop:certifiede}(5). The converse caller error is a named
\texttt{logic\_error} (finding m-6). Each surface's arm is named.
(5)~\textbf{(c): the coordinates of $(G,R_l)$ and the direct test.} (c1) now states that the input
is read in $C^q_{\mathrm{eq}}$ coordinates (finding m-9). WP-3r's arm~(M) tests the rule directly,
with mutations (finding M-1).
(6)~\textbf{(d) as built.} The $K$-mode; the Bockstein, whose meaning now points at [C]
\HA{rmk:bocksteinexp}(1)--(2) (proved since [C] v1.8.0; the F-A pointer the slot-8 plan asked
for); the free \texttt{charpoly}, asserted at weight $12$ (finding m-7); \texttt{red\_H\textasciicircum q}
with its arms~(R1)--(R5); and what the \ref{D6} arm measured at $(10,691)$.
(7)~\textbf{(e): footnote~$b$ restated, and the landed values.} ``$51\ne52$'' is [L]'s
amalgam-model coordinates. In the shipped model it is $234\ne235$. The footnote now states the
model-free content, a deficit of exactly one (WP-3's finding, confirmed by finding m-8; for the
owner to ratify, WP-13 report). The frozen values are listed after the object list.
(8)~\textbf{\ref{D6}}, the item and its row in \S\ref{app:claims}: an E8 run exists, in the
validations and golden only, so the published semantics is unchanged.
Companion pointers: [C] \textbf{v1.8.0}, [X] v1.11, [P] \textbf{v3.6},
\texttt{ENGINE\_API.md} \S13, whose pointer repairs (WP-12 pass~1, m-1--m-3) are the slot close's);

v1.9.0-pre.1: 2026-09-24 (hecke-extensions slot 8, package WP-2, \textbf{spec
first}: the pre-release of v1.9.0 that slot-8 WP-13 completes. It carries its own identifier
\texttt{1.9.0-pre.1}, so that no version number is published twice under different text ([C]
V1\,/\,V2; slot-8 plan \S4.3). An $X$-type change: a mathematical setting gains its grammar,
default and refusal rule, and the E8 objects gain their names and meaning, all \emph{before} slot-8
WP-3 registers the key and WP-4 emits the objects, so that their content gates are checks and not
stops. No claim, hypothesis, formula, STATUS tag or existing label is renamed or removed; nothing
here is emitted by a record in slot~8. Written under the owner's rulings of 2026-09-24 on the
slot-8 decisions 2, 3, 4 and 5, each at recommendation~(A) (slot-8 plan \S10); nothing below is
provisional. New \S\ref{sec:ring} at the end of \S\ref{sec:schema}:
(1)~\textbf{\texttt{coefficient\_ring}}: its grammar \texttt{F\_ell} $\mid$
\texttt{Z/ell\textasciicircum<k>}, the reservation of \texttt{Z} (line~2, [IB] OD-10~(a)) and of
\texttt{Z/<m>}, its default \texttt{F\_ell}, its canonical text, its fingerprint and golden cost
(the grammar and the default cost no golden line; a record field added for a ring costs golden
lines unless it is emitted only under a non-default ring), and where the ring is read;
(2)~\textbf{the refusal rule} (decision~4): every construction site and every field-only consumer
that does not honour a non-default ring raises \texttt{hecke::NotImplementedInThisBuild} naming
itself (console exit~5); the seven user-facing surfaces by name; the closure as a rule;
(3)~\textbf{the E8 objects} (decisions 2 and 3): \texttt{lengths\_H\textasciicircum q},
\texttt{invariant\_factors\_H\textasciicircum q}, \texttt{ladic\_type\textasciicircum q} with
\texttt{ladic\_type\_certified} and the two counts $r_K(d^0)$, $r_K(d^1)$, the free-quotient
\texttt{charpoly} with its precondition and precision $K-e_{\max}$, \texttt{bockstein\_rank\textasciicircum 0}
and \texttt{higher\_torsion\_count} with the bookkeeping identity and its meaning, and the
reduction maps (tier~B);
(4)~\textbf{the cyclic-summand basis rule} ([X] v1.11 E8 finding M3), stated once;
(5)~\textbf{the E8 golden point name} \texttt{E8 ring=Z/<ell>\textasciicircum<k> n=<n> ell=<ell>}
and the seven points with their primes (decision~5), with the predicted values.
Pointers added in \S\ref{sec:field} and at the two older \texttt{coefficient\_ring} sentences of
\S\ref{sec:schema}, each marked ``v1.9.0-pre.1''. Companion pointers unchanged: [C] v1.7.4,
[X] v1.11, [P] v3.5, \texttt{ENGINE\_API.md} \S13);

v1.8.2: 2026-09-23 (hecke-extensions slot 7, package WP-14, the slot close:
\textbf{status only}. No claim, hypothesis, formula, STATUS tag or label moves; no gate, cell
list, constant or rule moves. Close report \texttt{docs/implementation/slot7/WP-14\_report.md}.
(1)~\textbf{E7 is accepted} (slot-7 WP-1): \HA{cl:pd}(1)--(2) are tested on all $54$ cells of
the E2 case list at $\ell\ge5$ and fail on none, and the E7 objects are golden (add-only). The
two ``E7, scheduled'' status phrases (\S\ref{sec:identities} item~5, the Prediction table's
row for Prediction~\ref{pr:dual}) now say so --- finding m-7 of the slot-7 referee
(\texttt{docs/implementation/slot7/WP-13\_referee.md}), outside v1.8.0's diff and left to the close.
(2)~\textbf{E6 is declared} at the slot close. \S\ref{sec:field}'s cell list and failure rule
are now binding in [X] \textbf{v1.11} \S6.2~E6~(d), which pins the threshold $r_{\max}\ge10$;
the box's ``until then, the slot-7 plan's \S10 and \S2.4 row T3 are binding'' is marked
superseded. The \texttt{conway} fall-back is \textbf{routed}: every production entry point
raises \texttt{hecke::NotImplementedInThisBuild} (slot-7 WP-4).
(3)~\textbf{\S\ref{sec:identities} item~6: status at the close.} $\Sigma(691)$ is built (slot-7
WP-5: $1\,002\,225$ systems, byte-identical across two builds) and published by hash and build
command only; the identity is still not evaluated at $\ell=691$, and the record still carries
\texttt{atlas\_unavailable} there (ruling~3b~(B)).
(4)~\textbf{The first ext campaign is published} (\texttt{docs/matrix-engine/results/ext\_first/},
$505$ cells, \texttt{hecke\_ext/v3}, with \texttt{analysis.txt}); the record is
\texttt{hecke\_ext/v3} since slot-7 stage~3, exactly as \S\ref{sec:schema} (h) specified.
Companion pointers: [C] \textbf{v1.7.4}, [X] \textbf{v1.11}, [P] v3.5);

v1.8.1: 2026-09-23 (hecke-extensions slot 7, package WP-13, the independent
referee pass on v1.8.0; report \texttt{docs/implementation/slot7/WP-13\_referee.md}). A $Y$-type
change ([C] V2; [OD] slot-4 close-out~4): \textbf{pointers and wording only}. No claim,
hypothesis, formula, STATUS tag or label is added, renamed or removed, and no gate, cell list,
constant or rule moves. Every repair is marked ``v1.8.1'' at its site.
(1)~\textbf{Algorithm~\ref{alg:split}'s cost note: which run of arm~H measures what} (finding
M-1). v1.8.0 attributed the pre-fix fitted exponent ``$4.8$--$5.6$'' to
\texttt{HeckeCanonicalFieldValidation} arm~H. At this revision arm~H counts the \emph{fixed}
route and prints its fit ($4.89$, $4.23$, $3.87$ at $\ell=5$, $691$, $134217757$, over
$r=10..20$), and it counts the pre-fix route only for $r\le10$ (no fit, no $r=19$ count). The
pre-fix figures, $4.8$--$5.6$ and $144.1$M at $r=19$, are arm~H's output at commit
\texttt{fa829c7}. The note now names the commit for each figure.
(2)~\textbf{The two high-$r$ cells, framed once} (finding M-2). The note said that the bound failed
on ``both high-$r$ cells measured'' and then that E6~(d) is met on the $41$ gating cells, without
saying that the two are not among the $41$. Both notes now say that neither is a campaign cell
($n=210>200$; $\ell=134217757$ is not a campaign prime), that both are reported only, and what
W1 measured for them after the fix ($0.0062$ and $0.0151$).
(3)~\textbf{``Refuted at $r\ge11$''} (finding m-1) becomes ``refuted on the two high-$r$ cells
measured, at $r=14$ and $r=19$'': no cell at $r\in\{11,12,13\}$ was measured before the fix.
(4)~\textbf{The $r_{\max}\ge10$ threshold} (finding m-2) is marked as WP-4's reading of ruling~4's
``highest-$r$ cells'' (\texttt{e6\_measure.sh}'s \texttt{--r-min} default), to be pinned in [X]
v1.11.
(5)~\textbf{\S\ref{sec:field}: the superseded automatic flip} (finding m-3) is marked
superseded where it stands.
(6)~\textbf{\S\ref{sec:schema}, ``The stage-2 keys''~(e)} (finding m-4): which arm asserts which
of the three statements, and that the golden \texttt{ext\_json\_*} texts are text hashes.
Companion pointers unchanged: [C] v1.7.3, [X] v1.10, [P] v3.5);

v1.8.0: 2026-09-23 (hecke-extensions slot 7, package WP-12. This revision completes
v1.8.0-pre.1 and publishes v1.8.0. It is an $X$-type change, as the pre-release announced; no
claim, hypothesis, formula, STATUS tag or label is added, renamed or removed. Each rewritten claim
names the validation or window that measures it.
(1)~\textbf{Algorithm~\ref{alg:split}'s cost sentence} (slot-7 WP-4, commits \texttt{fa829c7},
\texttt{cab4e38}, verdict~(ii)). The sentence now gives a unit: $O(r^3\log\ell)$ $\F$-operations
$=O(r^5\log\ell)$ $\Fl$-multiplications with schoolbook arithmetic. ``Negligible'' is struck:
before the slot-7 fix it was refuted at $r\ge11$ on the high-$r$ cells. The fix and its counts are
recorded, measured by \texttt{HeckeCanonicalFieldValidation} arm~H. After the fix, window W1
measured the worst gating cell, $(170,691)$, at ratio $0.0104$ over the $41$ gating cells:
E6~(d) MET. The pre-release's wording is quoted in the note, marked as refuted.
(2)~\textbf{\S\ref{sec:field}'s decision box: E6~(d)'s cell list and failure rule}, as a pointer
(slot-7 owner ruling~4, cells~C1 and rule~Z). The binding text is [X] \S6.2~E6~(d) from [X] v1.11,
and the slot-7 plan \S10 until then. The W1 status is added. Nothing in this box moves the gate's
constant, numerator or denominator.
(3)~\textbf{\S\ref{sec:schema}, ``The stage-2 keys''~(e): \texttt{presentation} and the hashes}
(slot-7 WP-2 finding). The pre-release's ``leaves every hash of the record unchanged'' is true of
the record's \texttt{hashes} and of the golden \texttt{ext\_ideals} / \texttt{ext\_certify} hashes.
It is false of the record text, which moves through the \texttt{settings} block and its
fingerprint. The item now states both, as \texttt{HeckePresentationValidation} arms (k3) and (k4)
assert them.
(4)~\textbf{\S\ref{sec:identities} item~6: status}. The identity is asserted on $108$ cells at
$\ell\in\{5,7,11\}$ (slot-6 WP-9, \texttt{HeckeExtEndToEndValidation}). It is not evaluated at
$\ell=691$: $\Sigma(691)$ is slot-7 WP-5's and is not built at this revision, and under ruling~3b~(B)
the record carries \texttt{atlas\_unavailable} there either way. The body's ``not thereby asserted
--- that needs the atlas'' was stale since v1.7.2's changelog; it is kept, and followed by the
status.
(5)~\textbf{Addenda requested by slot-7 WP-3's content check}: none. WP-3 has not run at this
revision's base (\texttt{4b6e88a}). A request that arrives later is a v1.8.1 item for WP-13 or a
status item for WP-14.
(6)~\textbf{E8 findings checked, none contradicted here} (slot-7 WP-9's F1--F4 and WP-10's M1--M3,
\texttt{docs/implementation/slot7/WP-10\_review.md}). This document states nothing about
\texttt{ModCtx::inv}, invariant factors, basis reduction or $\rho_n$ over $\Z/\ell^k$. Its E8
sentences say three things: \ref{D6} says that the $\Z/\ell^2$ Eisenstein datum is a refinement;
the cone paragraph says that a claim is \emph{not} made over $\Z/\ell^k$; and \S\ref{sec:schema}
says that \texttt{coefficient\_ring} names the ring. F3's universal-coefficient point (a
$\Z/\ell^2$ cohomology reduces injectively into the $\Fl$ one, with a Tor cokernel) agrees with
\ref{D6}'s ``refinement, never a replacement''. The findings are [X]'s and E8's, for WP-14.
The companion pointer stays at [C] \textbf{v1.7.3});

v1.8.0-pre.1: 2026-09-23 (hecke-extensions slot 7, package WP-12a, \textbf{spec
first}: the pre-release of v1.8.0 that WP-12 completes. It carries its own identifier
\texttt{1.8.0-pre.1}, so that no version number is published twice under different text ([C]
V1\,/\,V2; slot-7 plan \S4.3): v1.8.0 will be WP-12's text, a different identifier. An $X$-type
change, since the record schema gains specified keys; no claim, hypothesis, formula, STATUS tag
or label is added, renamed or removed, and nothing here is emitted yet --- every key below is
written \emph{before} slot-7 WP-3 emits it, so that WP-3's content gate is a check and not a stop
(audit finding~M-3 of \texttt{docs/implementation/hecke\_extensions\_slot7\_plan\_audit.md}).
In \S\ref{sec:schema}, new paragraph ``The stage-2 keys'' and the corrected example:
(1)~\textbf{\texttt{sturm.proof\_range} and \texttt{sturm.proof\_range\_source}}: type, presence
rule, the closed value list \texttt{\{cell\_local, caller\_input, method\_undefined\}}, the
retired v2 value \texttt{dickson\_count}, and the default \texttt{cell\_local} at every $\ell$
(owner ruling~3a of 2026-09-23, slot-7 plan \S10), with the record-level
\texttt{anemic\_pool} the decision is made over;
(2)~\textbf{the atlas label and \texttt{atlas\_unavailable}}: the per-ideal \texttt{label} and
the record-level \texttt{s\_sep}, each either computed from the in-process $\Sigma(\ell)$ at
$\ell\in\{5,7,11,13\}$ or the placeholder \texttt{\{"status": "atlas\_unavailable"\}} at every
other $\ell$ (owner ruling~3b~(B)); Algorithm~\ref{alg:stage2}(3) says what the default run does
at an $\ell$ without an atlas;
(3)~\textbf{the split's location, stated once}: the per-ideal \texttt{split} is the
\texttt{hecke\_canonical\_split/v1} fragment verbatim, present iff \texttt{local} and
\texttt{split\_field\_report = canonical}; under \texttt{own\_field} the record carries
\texttt{coefficient\_field} only. The example's \texttt{field} and \texttt{char\_split}, which
no writer ever emitted, are retired;
(4)~\textbf{\texttt{predictions\_mode} and the prediction objects}: the record-level
\texttt{predictions\_mode} and \texttt{predictions} (the Kummer verdicts and the Glover factors),
and the per-ideal \texttt{predicted} / \texttt{confirmed} / \texttt{unpredicted} with a closed
prediction-name list;
(5)~\textbf{the \texttt{presentation} object's encoding} (audit finding~m-5; slot-7 owner
confirmation~15): generator indices are indices into the block basis, and the object names that
basis; present only under \texttt{presentation = on};
(6)~\textbf{the \texttt{settings} shape reconciled to the writer} (audit finding~m-12): each key
is an object \texttt{\{value, source\}}, the block holds exactly the registered keys in registry
order, and the fingerprint is sixteen hex digits. After slot-7 stage~1~(1) the golden texts see
only \texttt{\{fingerprint\}};
(7)~\textbf{\texttt{source} $\in\{$\texttt{default}, \texttt{spec}, \texttt{console}$\}$ is
binding and closed}, and the golden harness stamps \texttt{spec} (slot-7 owner confirmation~13).
In \S\ref{sec:field}, one pointer: tier~(ii)'s record placement is \S\ref{sec:schema}'s.
The companion pointer stays at [C] \textbf{v1.7.3}; [C] \HA{sec:recordschema} item~6(f) leaves the
field inventory to this document, so no [C] edit is forced);

v1.7.2: 2026-09-22 (hecke-extensions slot 6, package WP-14, the slot close:
\textbf{status only}. No claim, hypothesis, formula, STATUS tag or label moves.
(1)~\textbf{\S\ref{sec:field}'s pointer reduced} (slot-6 WP-12 finding~M-2, item~5 of its list
for the close): [X] \S6.2~E6~(d) now exists --- [X] v1.10 wrote the text of \texttt{e6\_gate\_brief.md} \S6.1 into it ---
so the interim pointer to \texttt{docs/implementation/e6\_gate\_brief.md} \S6.1 is marked
superseded. The gate's verdict is not this document's: it is [X] v1.10's and the slot-6 close
report's (\texttt{docs/implementation/slot6/WP-14\_report.md}), which records the two high-$r$
cells above $0.10$ and the owner's ruling on them ([X] v1.10, \texttt{[OD]} E6-3: a slot-7 fix of
the root finder, the default and this box's gate text unchanged, the verdict open, a high-$r$
cell added to the gate's list) and does not declare E6. Finding~B-1 of v1.7.1 (the reworded
$(160,691)$ sentence) is \textbf{accepted by the owner}, its departure from the slot-6 plan's
prescribed wording owner-approved.
(2)~\textbf{\texttt{conway}'s status wording} (finding~m-4): the value is parsed, validated and
moves the fingerprint; its consumer \texttt{FieldPolynomialFromSettings} raises
\texttt{hecke::NotImplementedInThisBuild} (console exit~5 since slot-6 WP-13). ``Rejected''
now says which. Slot-6 WP-13 found that the two production entry points
\texttt{CanonicalField(ell, r)} and \texttt{SplitInCanonicalField} bypass that consumer, so under
\texttt{conway} they silently use \texttt{least} --- latent, since no report consumer exists; an
open item of [X] v1.10, recorded here and not ruled.
(3)~\textbf{Status of the E7 and E6 items this document states}: $\Dn$ is computable since
slot-6 WP-1 (so the slot-3 addendum's last entry is no longer true); the pairings,
perfectness and PD adjointness of [C] \HA{cl:pd} are \emph{tested} on $20$ cells (WP-3),
socle $\leftrightarrow$ Loewy (Prediction~4.6, identity~7 of the certificate) is asserted
(WP-3); identity~6 (\S\ref{sec:alg} item~6, the residue atlas) is asserted on $108$ cells at
$\ell\in\{5,7,11\}$ under the standard-module convention of Claim~\ref{cl:glover}, and the
contragredient reading fails on all $63$ cells where it is distinguishable (WP-9) --- the
``a green suite cannot detect the discrepancy'' of that claim is now backed by a test for
identity~6; the atlas $\Sigma(\ell)$ exists for $\ell\le13$ (WP-8).
The companion pointer moves to [C] \textbf{v1.7.3});

v1.7.1: 2026-09-22 (hecke-extensions slot 6, package WP-12, the independent
referee pass on v1.7.0, report \texttt{docs/implementation/slot6/WP-12\_referee.md}: a $Y$ bump.
No claim, hypothesis, formula, STATUS tag or label moves. Three repairs, at two sites of
Claim~\ref{cl:bk} / \S\ref{sec:tests} and one of \S\ref{sec:field}.
(1)~\textbf{The $(160,691)$ sentence} (finding~B-1; \emph{not} pointers-only). v1.7.0's
``no single $T_q$ certifies the residue degrees that $\{T_2,T_3\}$ does'' contradicts the
sentence before it and the \texttt{SEPARATION\_SINGLE} line of
\texttt{HeckeExtSeparationValidation} that it cites: $T_2$ alone and $T_3$ alone each give the
family's $(\dim{:}f)$ list exactly, residue degrees included, and only $T_5$ alone under-reports
one ($f=1$ on the four-dimensional block). Both sites now read ``the block partition at
$(160,691)$, with its residue degrees, is already $T_2$'s (and $T_3$'s); $T_5$ alone gives the
same partition and under-reports one residue degree''; the v1.7.0 wording is kept in place,
marked as refuted. v1.7.0's item~(2) below, which quotes it, is superseded by this item; the
Appendix~\ref{app:claims} row made no residue-degree statement and is unchanged.
(2)~\textbf{The same two sites: where the partition identity comes from} (finding~M-1; one
clause, \emph{not} pointers-only). The validation compares $(\dim{:}f)$ lists in record order,
block counts and dimension multisets, not subspaces; that $T_5$'s partition \emph{is} the
family's, and the ``eight blocks'' argument, rest on the family's decomposition refining each
single operator's together with every run having seven blocks. The clause says so.
(3)~\textbf{\S\ref{sec:field}'s pointer} (finding~M-2; pointers-only). ``[X] \S6.2~E6~(d)''
names an item that [X] v1.9 does not yet contain; the pointer now says that it is written at
[X] v1.10 and that until then \texttt{docs/implementation/e6\_gate\_brief.md} \S6.1 is binding
(slot-6 plan \S4.3). v1.7.0's item~(1) is to be read with this pointer.
The companion pointer stays at [C] \textbf{v1.7.2});

v1.7.0: 2026-09-22 (hecke-extensions slot 6, package WP-11, on slot-6 WP-7's
code side and slot-6 WP-10's ratio-criterion audit: an $X$ bump, since one claim's verification
paragraph changes its conclusion and one gate changes its denominator. No claim is added,
renamed or removed, no STATUS tag moves, and no label is added or removed.
(1)~\textbf{\S\ref{sec:field}'s decision box: the E6 gate restated and reduced to a
pointer} (owner decisions \texttt{[OD]} E6-1 and E6-2 of 2026-09-22,
\texttt{docs/implementation/e6\_gate\_brief.md} \S6.3). The gate is now
$\texttt{canonical\_field\_s}\le0.10\times\texttt{total\_s}$ on the same 23 cells, with
numerator, denominator and protocol pinned by [X] \S6.2~E6~(d), the one governing statement;
the $\le2\times\texttt{own\_field\_s}$ half is deleted and \texttt{own\_field\_s} is a reported
column. The history of both halves is kept in the box. This matches the code comments slot-6
WP-7 rewrote (\texttt{Bench/BenchDriver.cpp}, \texttt{Hecke/HeckeSettings.h},
\texttt{HeckeExt/HeckeCanonicalField.h}) and the \texttt{[E6 gate]} line \texttt{--bench ext}
prints. Nothing here declares E6 or states a verdict: E6 acceptance is slot~7's
([X] \S6.2~E6~(c)), and the high-$r$ cells measured outside the 23 are an owner item, not text
of this document.
(2)~\textbf{Claim~\ref{cl:bk}'s verification paragraph and \S\ref{sec:tests}'s $(160,691)$
row repaired normatively} (v1.6.2 item~(1); slot-5 WP-8 finding~1). The measurement of
\texttt{HeckeExtSeparationValidation}'s \texttt{SEPARATION\_SINGLE} line shows that $T_2$ alone
already gives the family's seven blocks with their residue degrees, and that $T_5$'s $4{:}1$
\emph{is} the family's $4{:}2$ block --- the second of Claim~\ref{cl:bk}'s two alternatives, not
the first, which would give eight blocks. The closing sentence now reads ``no single $T_q$
certifies the residue degrees that $\{T_2,T_3\}$ does; the block partition at $(160,691)$ is
already $T_2$'s''; the refuted wording is kept in place, marked as such; the
Appendix~\ref{app:claims} row follows. Claim~\ref{cl:bk}'s general ``for $n\gg\ell$ no single
$T_q$ suffices'' is a counting statement and is unchanged; the $(160,691)$ row is no longer
cited as an instance of it.
(3)~\textbf{\texttt{sec:reportfield}: recorded as never having existed, and not added.}
\texttt{git log -S} finds no revision of this document that defined it; the text its citing
sites meant is \S\ref{sec:field}'s decision box with its tiers (i)/(ii). An alias would give
that one text a second name, which is what \texttt{[OD]} E6-2's one-authority rule was made to
remove, and the code sites already cite \texttt{sec:field} (slot-6 WP-7). The remaining
documentation sites are outside this package ([X] \S0 changelog, \texttt{docs/architecture-detail.md},
\texttt{HeckeExt/README.md}); an \texttt{[E]~<label>} citation is not scanned by the label census,
so the dangling name was invisible to it.
(4)~\textbf{From slot-6 WP-10's audit} (\texttt{docs/implementation/slot6/WP-10\_report.md},
rows~2 and~27): the \S\ref{sec:field} gate was the only cost bound in this document, and it is
the one (1) restates --- every other cost statement is asymptotic or a record; and Stage~1's
quoted $n=1000$ pair ($14.4$\,s against $14.8$\,s) has the whole below its part, so it now says
it supports ``dominates'' and is not a split.
(5)~\textbf{A status correction of v1.6.2 item~(3)}, which said ``with E2, E4b and E5 all
accepted in slots 4--5, E7's gates are all met'': E4b and E5 are accepted, E2 has not been
declared accepted (two of its acceptance items were open at the slot-5 close;
\texttt{docs/implementation/hecke\_extensions\_slot6\_plan.md} \S2.3 item~1). E7 starts on E2's
substantive state --- the record and the Loewy / socle data, which exist --- by the owner's
ruling of 2026-09-22 (that plan, \S10 ruling~1); the \texttt{det}-twist gate of
Claim~\ref{cl:glover} is met, as v1.6.2 said. The companion pointer stays at [C]
\textbf{v1.7.2}. No separate report: the commit message \texttt{hecke-ext(S6-WP-11)} carries the package record);

v1.6.2: 2026-09-22 (hecke-extensions slot 5, package WP-13, the slot close:
\textbf{status only --- no claim, hypothesis, formula, tag or label moves.} Three things are
recorded. (1)~\S\ref{sec:sep}'s closing separation sentence is \textbf{refuted by
measurement}: slot-5 WP-8 measured the claim on the shipped \texttt{ComputeMultiOp} at
$(160,691)$ and $(198,691)$, and \texttt{HeckeExtSeparationValidation} now records rather than
asserts it. The repair is a normative question for the next $X$ and is booked there, not taken
here. (2)~\textbf{\S\ref{sec:field}'s performance gate is defective as written} and is an
owner question: \texttt{canonical\_field\_s} $\le 2\times$ \texttt{own\_field\_s} cannot be met
by any implementation, since \texttt{own\_field\_s} is, by this document's own words, ``zero new
arithmetic'' --- 1--6\,$\mu$s of string formatting --- and the other half is a ratio whose
denominator is itself microseconds on 20 of the 23 atlas cells. Against \texttt{total\_s} at the
same $0.10$ threshold every one of the 23 cells passes, the worst at $2.9\,\%$. Slot-5 WP-7
measured the table and did \emph{not} flip the default; E6 acceptance is slot~6's. (3)~The
\texttt{det}-twist convention of Claim~\ref{cl:glover}, this document's gate on E7, is
\textbf{met}, and with E2, E4b and E5 all accepted in slots 4--5, E7's gates are all met. The
companion pointer moves to [C] \textbf{v1.7.2}. Close report:
\texttt{docs/implementation/slot5/WP-13\_report.md});

v1.6.1: 2026-09-22 (hecke-extensions slot 5, package WP-12, the independent
referee pass on v1.6.0: \textbf{a $Y$ bump --- pointers only. No claim, hypothesis, formula,
tag or label moves.} Two repairs, both of v1.6.0's own text, and one of a reading v1.6.0 left
standing. (1)~The version block said the \texttt{det}-twist convention is restated ``in the
one-sentence form an implementation quotes verbatim (\texttt{HeckeExt/GloverFactors.h} and the
console's \texttt{glover\_factors} both do)''. \emph{Neither quotes it}: both state the
convention correctly, in their own words --- \texttt{GloverFactors.h} as ``this is
\texttt{Sym\^{}n} of the STANDARD module, not of its dual: the two differ by
\texttt{det\^{}\{-n\}}, which would shift every \texttt{j\_i}'' and script
\texttt{017\_predictions.oc} likewise. The sentence is drafted so that they \emph{can} quote it
unchanged, and now says so; the mathematics, in all three places, agrees. (2)~The last bullet
of the slot-3 addendum, ``Still not computable: the production $\iota$ $\ldots$ the record
\texttt{hecke\_ext/v1}'', is correct \emph{of its own date} but is the last addendum in that
section, so it reads as the present state; it now points at what slots~4 and~5 built, leaving
$\Dn$ (E7) as the only entry still true. The companion pointer moves to [C] \textbf{v1.7.1}.
Report: \texttt{docs/implementation/referee\_report\_2026-09-22\_contract\_v1\_7\_0.md});

v1.6.0: 2026-09-22 (hecke-extensions slot 5, package WP-11: \textbf{the
convention that gates E7, and \S\ref{sec:schema} at \texttt{hecke\_ext/v2}} --- an $X$ bump,
since one claim's preamble acquires a consequence and one schema key changes its value. No
claim is added, renamed or removed and no tag moves.
(1)~\textbf{Claim~\ref{cl:glover}'s \texttt{det}-twist convention is recorded as E7's
gate.} The convention itself --- $\Sym^n$ is the \emph{standard} symmetric power of
\HA{def:rhon} and not its contragredient, the two readings shifting every $j_i$ by $-n$ ---
has been normative since v1.5.0; what this version adds is the consequence [WP4-7] did not
draw: E7's subject \emph{is} $\Dn$ and $\rho_n^{\vee}$, its first use of the factor multiset
is item~6 of \S\ref{sec:identities}, and under the two readings that comparison differs by
$\theta^{-n}$ on every factor. Both readings being internally consistent, \textbf{no green
validation suite can detect the discrepancy} --- every checkable invariant is
convention-independent or self-consistent --- so a wrong reading surfaces as an E7 finding
indistinguishable from a result. \textbf{This document at v1.6.0 is therefore a precondition
of E7's start}, and the convention is restated in the one-sentence form an implementation
can quote verbatim --- \texttt{HeckeExt/GloverFactors.h} and script
\texttt{017\_predictions.oc} both state the convention, in their own words, and can now
quote this sentence unchanged (v1.6.1: v1.6.0 said ``an implementation quotes verbatim
$\ldots$ both do'', which no file in the tree did).
(2)~\textbf{\S\ref{sec:schema} moves to \texttt{hecke\_ext/v2}}, making the \emph{two}
repairs v1.5.0 booked for ``the same slot-5 revision'' (owner ruling \texttt{[OD]} slot-4
close-out 5 of 2026-09-22; implemented by slot-5 WP-1): \texttt{u\_ell} is carried at
\textbf{record level with a single \texttt{resumable}}, so \texttt{resumable\_items} counts
\emph{decisions} and no \texttt{ideals[]} entry carries the object --- a change to an
existing key's \emph{value}, hence a version bump, made before the first emission --- and an
\texttt{ideals[]} entry gains the per-ideal \texttt{groups} list carrying the
\texttt{SameMaximalIdeal} witness, \texttt{modules} keeping its ``one entry per ambient
group'' shape exactly as v1.5.0 predicted. Three further schema corrections from the same
package: \texttt{canonical} is keyed by the ambient group at \emph{ideal} level (the example
had it nested inside \texttt{modules}); a closed-list scalar a record has not computed is
\texttt{null} and never dropped; and \texttt{les.hecke\_equivariant} is tri-state, a
cone-only record not having evaluated it. The example record is updated to all of these.
(3)~Claim~\ref{cl:harder}'s hypothesis-versus-domain paragraph and Claim~\ref{cl:glover}'s
recursion are \emph{unchanged}: both landed at v1.5.0/v1.5.1, so the slot-5 plan's reading
that they were outstanding was stale --- recorded here so the next reader does not look for
them. The companion pointer moves to [C] \textbf{v1.7.0}
(\HA{obl:coneexact}(iii) at every $\ell$ with its proof, \HA{sec:recordschema} item~6's nine
rules; v1.6.1 moved this pointer to v1.7.1). Report:
\texttt{docs/implementation/slot5/WP-11\_report.md});

v1.5.2: 2026-09-22 (hecke-extensions slot 4, package WP-11, the slot close:
\textbf{status only --- no claim, hypothesis, formula, tag or label moves}. What the slot
built of this document's programme: \textbf{E2} --- the extension record \texttt{hecke\_ext/v1}
of \S\ref{sec:schema} is written, for the bulk complex, by a \emph{separate} writer
(\texttt{HeckeExt/HeckeExtRecord} + \texttt{HeckeExtReport}, never an extension of
\texttt{Hecke/HeckeReport}), with the one tri-state \texttt{Final}, the settings block and
its fingerprint, and the resumable sidecar (slot-4 WP-1); and
\texttt{HeckeExt/HeckeExtDriver} turns a cell into that record through Stages~1--5, with a
streamed operator for a Sturm prime outside $S_0$ and the $U_\ell$ resume (WP-2).
\textbf{Claim~\ref{cl:ramanujan} (8.1) is verified end to end through that driver} at
$(10,691)$ --- the four exact identities in $M_3(\F_{691})$, \emph{including}
$3(U_{691}-1)=-542(T_2+24)$, which exercises $U_\ell$ and $\tau(691)$ --- with the
basis-independent content separated from the stored-basis pins, as the claim's
``Verified'' paragraph now says. \textbf{E3} --- a \emph{production} $\iota^{\ast}$ exists
(\texttt{Hecke/HeckeInvolution}, \texttt{HeckeExt/HeckeSignSplit}), so the sign split of
every stored $H^q$ and of every computed Hecke matrix, and the Eisenstein line of D6, are
pipeline objects and no longer validation-local (WP-3). \textbf{E6} --- the first two
pre-record units ship: \texttt{HeckeExt/HeckePredictions} (Claim~\ref{cl:harder}'s
predicate) and \texttt{HeckeExt/GloverFactors} (Claim~\ref{cl:glover}'s recursion), both
gated by \texttt{HeckePredictionsValidation} against independent formulas --- Glover
against the \emph{Brauer character} of the real $\rho_n$ on every $\ell$-regular class,
which is a proof cell by cell. Still open, and unchanged: the atlas, $S_{\mathrm{sep}}(\ell)$,
the canonical field, and the boundary / cone / per-sign record objects (slot 5). The [C]
pointer moves to v1.6.2. Report \texttt{docs/implementation/slot4/WP-11\_report.md});

v1.5.1: 2026-09-22 (hecke-extensions slot 4, package WP-9, the independent
referee pass on v1.5.0: \textbf{four mathematical repairs, two of them to false statements},
no claim added, renamed or removed and no tag moved.
(1)~Claim~\ref{cl:harder}'s new predicate paragraph is rewritten: v1.5.0 named only the two
instances $k=\ell-1$ and $k=\ell$ of the branch $(\ell-1)\mid k$ and so left every even $k$
with $(\ell-1)\mid k$ and $k>\ell-1$ --- at $\ell\in\{2,3\}$ \emph{every} even $k$, the
Kummer branch being vacuous there --- with no branch at all, and it answered $k=\ell$ with a
fact about \emph{odd} $\ell$ although the only in-range instance is $(n,\ell)=(0,2)$, where
the verdict is \texttt{false} and not the ``degenerate \texttt{true}'' the sentence gave.
The general von Staudt--Clausen branch is now stated, and the $B_k=0$ remark moved to odd
$k$, where it is true. The shipped \texttt{kummer\_verdict} already decided all of these
correctly; the document did not.
(2)~Claim~\ref{cl:glover}(4)'s counts are corrected to $\lfloor(n+1)/(\ell+1)\rfloor$ steps
and $2\lfloor(n+1)/(\ell+1)\rfloor+1$ raw factors: v1.5.0's bound undercounts at $n=\ell$
($2$ factors against $1$) and at $n=2\ell+1$ ($4$ against $3$).
(3)~Claim~\ref{cl:sharpbound}'s irreducibility step is moved to $A_S$, where the idempotent
a reducible $\min_V(t_p)$ produces actually lives; inside the residue field the argument
assumed what it proves.
(4)~Claim~\ref{cl:sharpbound}'s ``lower bound in general'' is replaced by the
surviving-block invariant --- every $\pi_p$ irreducible, hence $\deg\pi_p$ divides every
residue degree --- v1.5.0's reason being false in both directions on a general non-local
block. The companion pointer stays at [C] \textbf{v1.6.1}. Report:
\texttt{docs/implementation/referee\_report\_2026-09-22\_contract\_v1\_6\_0.md});

v1.5.0: 2026-09-21 (hecke-extensions slot 4, package WP-8: \textbf{five
open items of this document resolved, and the findings of the slot-4 units it owns} --- an
$X$ bump, since one claim is added and one certificate gains a hypothesis.
Certificate~\ref{sec:localcert} states what a \emph{reported} witness must satisfy
($w$ commuting with the family, the hypothesis a matrix witness admits and a coefficient
list makes automatic) and that the sub-block order after a split is deterministic
\emph{per mode} and not mode-independent --- the partition is, the list is not (owner
confirmation of 2026-09-21, item~1(c): per mode, explicitly not an intrinsic order);
Claim~\ref{cl:sharpbound} is new, $\dim_{\Fl}\bar A=\operatorname{lcm}_p\deg\pi_p$ on a
local block, which is why budget mode's \texttt{residue\_degree\_lower\_bound} is sharp
there; \S\ref{sec:schema} carries \texttt{phase\_reached} per block \emph{and} the run's
high-water mark under separate keys, the settings under their \emph{registered} names
(\texttt{sturm}, \texttt{sturm\_unavailable}, and the new \texttt{commutant\_budget}), the
two sidecar rules (a record-relative path, an injective name), the closed name list behind
\texttt{certificate\_failed:<code or certificate name>} and the four scalars that sit
inside \texttt{groups}. From the slot-4 units: Claim~\ref{cl:glover} is new --- the Glover
factors of $\Sym^n$ by the $\theta$-recursion, \textbf{with the standard-module convention
that fixes the $\det$-twists}, which item~6 of \S\ref{sec:identities} had left free --- and
Claim~\ref{cl:harder} separates the criterion's hypothesis $n<\ell-1$ from the
divisibility predicate's domain. From the E2b driver (slot-4 WP-2), three further
\S\ref{sec:schema} rules and one wording repair: a \emph{resumed} run writes the settings
block a \emph{fresh} run at the same effective values would write, which is what makes
``resume to the default budget equals a fresh default run'' true byte for byte and is
otherwise a fingerprint two implementations could disagree on; \texttt{u\_ell} is carried
\emph{per ideal} although the verdict is per \emph{cell}, so one decision is emitted once
per ideal and \texttt{resumable\_items} counts entries; an \texttt{ideals[]} entry is
\emph{per group}, the cross-group matching of Stage~2 being a driver diagnostic that no key
reports; and Claim~\ref{cl:ramanujan}'s $c'_q$ line is marked \textbf{basis-dependent},
the rest of that paragraph not being. The companion pointer moves to [C] \textbf{v1.6.0}
(slot-4 WP-8, the eight items v1.5.2 deferred; no label this document cites moved).
Report: \texttt{docs/implementation/slot4/WP-8\_report.md});

v1.4.2: 2026-09-17 (hecke-extensions slot 3, package WP-11, the slot close:
\textbf{status only}, a $Y$ bump --- no claim, certificate, tag or label moves. Strict
locality mode (Certificate~\ref{sec:localcert}) is \emph{implemented} (slot-3 WP-3), E1b's
Sturm sweep with both \texttt{unavailable} modes (Certificate~\ref{cert:sturm}) is
implemented on real blocks (WP-2), and the equivariant cone with its Hecke action exists
(WP-5, WP-6), so $X_3,X_6$ are computable on the model complexes; the addendum of
\S\ref{sec:today} records it. The companion pointer moves to [C] \textbf{v1.5.2} (slot-3
WP-11, status only; no label this document cites moved). Report:
\texttt{docs/implementation/slot3/WP-11\_report.md});

v1.4.1: 2026-09-17 (hecke-extensions slot 3, package WP-8: \textbf{one
certificate made precise}, and the companion pointer moved. Certificate~\ref{cert:cyclic}
now \emph{states} that the reported generator is the least standard basis vector outside
$\m_SV$, in a paragraph of its own, in place of v1.4.0's implementation note that the code
fixes that choice while the certificate says ``any $v\notin\m_SV$'' --- the gap slot-2
WP-4 \S7.2 reported and slot-2 WP-11 did not take up. \textbf{Why a $Y$ bump and not an
$X$}: the mathematical assertion of the certificate is unchanged --- every $v\notin\m_SV$
still generates, and the certificate is true for each --- what is added is a
\emph{conformance} requirement on what a record may report, which no claim of this
document depends on; nothing else in this revision. The companion pointer moves to
[C] \textbf{v1.5.0} (slot-3 WP-8, the seven items v1.4.2 deferred; no label this document
cites moved, so the citations are unaffected). Report:
\texttt{docs/implementation/slot3/WP-8\_report.md});

v1.4.0: 2026-09-16 (hecke-extensions slot 2, package WP-11: \textbf{the
consequential revision}, the first since v1.1 in which claims of this document change; it cited
the contract at [C] \textbf{v1.4.1}, and with it this document agrees with the implementation plan
[X] v1.5.1 and is the freeze candidate of [X] \S8 (OQ17's $(160,691)$ verdict being in since
v1.3.3). \emph{Two predictions refuted by computation}, each kept with its old text, the evidence
and the corrected statement rather than rewritten into its negation: \S\ref{sec:tests} /
\S\ref{sec:schema}'s ``at a Kummer prime $r^1=0$ on the Eisenstein class, so $H^1_!=H^1$'' is
\emph{false} on all ten Kummer cells --- $\rank r^1=1$, $\dim H^1_!=\dim H^1-1$
(\texttt{HeckeBoundaryParabolicValidation}, slot-2 WP-7, 2026-09-16) --- and \ref{D6}'s ``both
reduce to the same line'' is \emph{sharpened}: on the unsplit $H^1$ the joint $1+q^{n+1}$
eigenspace is two-dimensional at every Kummer cell, the socle of the block; the socle semantics
is kept. \emph{Made precise} from [X] v1.5.1 and the slot-1/slot-2 package reports:
\S\ref{sec:schema} --- the tri-state \texttt{final} with R3.2's reason list, the keys
\texttt{settings}, \texttt{resumable}, \texttt{resumable\_items}, \texttt{sturm.mode},
\texttt{sturm\_primes}, \texttt{coefficient\_ring}, the \emph{per-complex} \texttt{q\_max} /
\texttt{truncated} / \texttt{q\_max\_source} / \texttt{degrees}, \texttt{P} keyed by prime,
\texttt{hecke\_report/v2}; Definition~\ref{def:bprime}'s cap as the setting \texttt{sturm} and
the record-level \texttt{not\_applicable} at $n=0$; Certificate~\ref{cert:sturm}'s
\texttt{unavailable} branch (\texttt{nonfinal} default / \texttt{heuristic});
Algorithm~\ref{alg:stage2}(2) / Certificate~\ref{sec:localcert} with the \textbf{strict mode}
(Frobenius radical, Berlekamp subalgebra, residue witness over $\bar A$, termination, the
$(2,30)$ example) and the budget-mode reason \texttt{budget\_exhausted:locality};
Certificate~\ref{cert:cyclic}'s canonical generator; Algorithm~\ref{alg:stage3} as implemented;
Certificate~\ref{cert:comm}'s \textbf{socle bound} and the sharp unknown count $u\le mD$ (its
\VERIFY{} on the cost discharged); the canonical field as the default report form under the E6
gate (\S\ref{sec:field}); \S\ref{sec:tests}'s $(198,691)$ row (``possible ambiguity documented,
LMFDB to-do''); Stage~7's reach reframed, the first campaign at $n\le200$; $\Dn$ and the pairing
\emph{computed} (E7, scheduled) at \S\ref{sec:symdiv}, Prediction~\ref{pr:dual} and identity~5;
\ref{D6}'s $\Z/\ell^2$ refinement as a refinement. \emph{The boundary chapter of [C] v1.4.1
folded in} at \ref{add:bdry}, \ref{add:cone}, \ref{add:iota}: \HA{obl:taubdry}(v) is
(v.a)--(v.c) and degree~$0$ alone at $\nu\ge2$; the circle model carries the $-I$
self-relations (\HA{prop:circlemodel}(4)); the coordinate $\Phi$ of \HA{prop:bdrychar}(3);
$M_\beta$ on $M^T$ only; \HA{alg:taubdry} is the canonical $\tau_\partial$ and
\HA{rmk:taubdrynoiota} says it is not $\iota$-symmetric; (vii) at $\ell\in\{2,3\}$ a
\textsc{comp} record. \S\ref{sec:today} gains the slot-2 addendum; \HA{cl:chi}'s \VERIFY{} stays;
Appendix~\ref{app:claims} follows. Not changed: every \STD{} claim is recalled, not re-proved;
no claim of \S\ref{sec:hecke}--\S\ref{sec:primary} moves);

v1.3.3: 2026-09-16 (hecke-extensions slot 1 close; a status pass, no claim of
this document changed: the $(160,691)$ \VERIFY{} of Claim~\ref{cl:bk} and \S\ref{sec:tests} is
\emph{resolved} by the slot-1 probe --- all seven blocks local, the shipped per-operator
refinement already separates $T_5$'s $4{:}1$ block, the first of the two alternatives holds ---
and the row says so; Claim~\ref{cl:bprime}'s computed values gain $B'(258,691)=2$ and the values
the implementation \texttt{HeckeExt/HeckeSturmBound} now pins; the character $P_j$ is written
as the reduced representative $X\bmod\pi_w$ (constant at $r=1$) in Claim~\ref{cl:resfield}
and \S\ref{sec:schema}; \S\ref{sec:today} gains a dated addendum on what slot 1 built
(the boundary Hecke operator on $X_2,X_5$, the read-only locality certificate with its
all-local grid verdict, the Sturm bound, the identity certificate); the Appendix~\ref{app:oq2}
row ``boundary group at $\ell=2$'' no longer calls the truncation policy open --- it is
[C] \HA{open:ell2} / \HA{oc1a});

v1.3.2: 2026-09-16 (coherence pass: the (A2) closing sentence of
Stage~0 no longer calls the $\ell=2$ truncation policy ``open \ldots\ not decided
here'' --- it is decided in [C] v1.3.0, \HA{open:ell2} / \HA{oc1a}, $q_{\max}$ a recorded run
parameter; \ref{add:iota}'s ``not a pullback along any lift'' is qualified ``for $n$ even'';
and the involution is no longer called ``not implemented'' --- two validations build it. No
claim of this document changes);

v1.3.1: 2026-09-16 (owner decision OQ3 applied at its two sites,
\ref{add:iota} and the \texttt{iota\_sign} paragraph of \S\ref{sec:schema}: [C] v1.2.0
defines $\iota^{\ast}:=-\iota^{\ast}_{\mathrm{geo}}$, so that $H^+$ is by definition the
Eisenstein side and \texttt{iota\_sign} names that normalisation; no claim of this document
changes);

v1.3: 2026-09-16 (\textbf{the scope cut}: the nine statements of the
former \S1.1--\S1.2 and \S1.4 --- the six groups, the amalgam, the dimension of $H^1$, the
Dickson--Bockstein claim, the boundary groups, the six-term sequence and its display, the
boundary characters, Poincar\'e duality --- are removed from this document and cited from
\texttt{gl2\_contract.tex} instead, which is now [C]; the four of them that [C] corrects are
named as superseded in \S\ref{sec:imported}; \S\ref{sec:imported} replaces the removed text,
Appendix~\ref{app:claims} loses their STATUS rows, and the cone pointers into the frozen
\texttt{monodromy\_v5.tex} become pointers into [C]);

v1.2.2: 2026-09-15 (the owner's decisions on the OQ2 verdict's eight open points applied);

v1.2.1: 2026-09-15 (independent audit of the v1.2 revision applied);

v1.2: 2026-09-15 (OQ2 verdict applied, Appendix~\ref{app:oq2});

v1.1: 2026-09-14 (audit response applied, Appendix~\ref{app:audit});

v1.0: 2026-09-14, pre-audit

\subsection{The title page of 1.25.0-pre.1: author line and abstract, verbatim}
\begin{quote}
Orbifold Cohomology of $\GL(2)$ --- companion to \texttt{gl2\_contract.tex}
\end{quote}
\begin{quote}
This document specifies, in the notation and normalisations of the binding mathematical
contract \texttt{gl2\_contract.tex} (cited as [C]), what it means to compute and classify \emph{all}
extension classes of Hecke modules occurring in the cohomology groups the project can produce,
and how to do it with deterministic certificates at every step. It is the concise LaTeX
version of the owner's two notes \texttt{hecke\_extensions\_\allowbreak mathematics.md} ([M]) and
\texttt{hecke\_extensions\_\allowbreak computational\_\allowbreak strategy.md} ([S]), made rigorous and
computation-oriented. The answer to ``which $T_q$ are needed'' is reduced to a certified
local statement, one maximal ideal at a time: the operators computed must generate the local
Hecke algebra $\T_\m$, and the extension data are the isomorphism invariants of the
$\T_\m$-modules $H^\bullet_\m$ (Loewy and socle series, active tangent directions, Jordan
types, positions of the canonical submodules). Every claim carries a STATUS tag; the tags
are indexed in Appendix~\ref{app:claims}. The document is accurate about what the code
base implements today (the single- and multi-operator primary decomposition with certified
residue degrees, Stages S10--S11) and what is new. Version~1.1 applies the independent audit
of 2026-09-14 item by item (Appendix~\ref{app:audit}); the two blocking findings changed the
Sturm certificate (a per-$(k,\ell)$ computed bound $B'$, Claim~\ref{cl:bprime}) and the
tangent-space test (the coefficient field, Remark~\ref{rmk:coef}). Version~1.2 applies the
independent verdict on Open Question~2 (the Hecke action on the mapping cone,
Appendix~\ref{app:oq2}): the cone action is defined from a certified triple
$(\tau,\tau_\partial,\eta)$ over $\Z$ with the cochain homotopy \emph{derived}, which removes the
$\Hom(\ker r^1,\coker r^0)$ ambiguity that v1.1 did not state (\ref{add:cone}, Prediction~\ref{pr:cone}).
Version~1.2.1 applies the independent audit of that revision, all of it inside \ref{add:cone}: the
existence statement is corrected ($\tau_\partial$ and $\eta$ over $\Z$ for every $\ell$; the bulk
$\tau$ is the canonical table of \HA{alg:tau}, integral when no fill used the mod-$\ell$ fallback),
the realisation of the $\Hom$ ambiguity is marked as proved \emph{over a field}, the non-vanishing
of $h_\eta$ on $X_3$ is marked undecided rather than proved, the contractibility and
two-dimensionality of $X$ are stated where they are used, and the two
Appendix~\ref{app:oq2} rows that record undecided owner points are relabelled as working
assumptions. Version~1.2.2 applies the owner's decisions of 2026-09-15 on all eight open points of
that verdict: the working assumptions become settled text and Appendix~\ref{app:oq2}'s rows become
plain \textsc{adopted}. Three of the eight were decided as \emph{amendments} to the verdict's
recommendation; the one with mathematical content here is the boundary group, fixed as
$\Ga_\infty=\langle T\rangle\times\langle-I\rangle$ (the Borel reading, which is \HA{sec:six}'s own
definition), so that the circle model of \ref{add:bdry} is its specialisation at $2\in R^\times$ and
$H^\bullet(\partial;M)$ is \emph{unbounded above in degree} at $\ell=2$ (\ref{add:bdry}, \ref{add:cone}).
Version~1.3 carries out the scope cut: [C] is now the unified contract
\texttt{gl2\_contract.tex}, which supersedes \texttt{hecke\_action.tex} and
\texttt{monodromy\_v5.tex}, and the nine statements about the six groups that this document
used to state are stated there instead and cited here (\S\ref{sec:imported}). This document is
what the cut left it as: the conjectural and predictive classification specification --- the
Hecke algebra acting, the separation and generation bounds, the local algebras, the
classification data \ref{D1}--\ref{D7}, the predictions, the stage-by-stage algorithms with
their certificates, the scalar extension and the output schema. Version~1.4.0 (slot 2 of the
extension stage) is the first revision since v1.1 in which claims change: two predictions were
refuted by the boundary computation of slot 2 and are corrected with the evidence in place
(\S\ref{sec:tests}, \ref{D6}); the record schema, the Sturm, locality and generation certificates
and the local-algebra stage are made precise from what slots 1--2 built; and the boundary chapter
of [C] v1.4.1 is folded in.
\end{quote}

\end{document}
